\documentclass[11pt,oneside]{book}
\usepackage[margin=1in]{geometry}
\usepackage{amsmath,amssymb,amsthm,mathtools}
\usepackage{booktabs,longtable,array}
\usepackage{enumitem}
\usepackage[colorlinks=true,linkcolor=blue,urlcolor=blue,citecolor=blue]{hyperref}
\usepackage[nameinlink,noabbrev]{cleveref}

\newtheorem{theorem}{Theorem}[chapter]
\newtheorem{lemma}[theorem]{Lemma}
\newtheorem{proposition}[theorem]{Proposition}
\newtheorem{corollary}[theorem]{Corollary}
\newtheorem{assessment}[theorem]{Assessment}
\newtheorem{firewall}[theorem]{Claim firewall}
\newtheorem{openproblem}[theorem]{Open problem}
\newtheorem{record}[theorem]{Record}
\newtheorem{programme}[theorem]{Programme}
\newtheorem{audit}[theorem]{Audit}
\newtheorem{correction}[theorem]{Correction}
\newtheorem{warning}[theorem]{Warning}
\newtheorem{decision}[theorem]{Decision}
\newtheorem{scope}[theorem]{Scope}
\theoremstyle{definition}
\newtheorem{definition}[theorem]{Definition}
\theoremstyle{remark}
\newtheorem{remark}[theorem]{Remark}

\newcommand{\MGclosed}{\textsc{closed}}
\newcommand{\MGopen}{\textsc{open}}
\newcommand{\MGpartial}{\textsc{partial}}
\newcommand{\Tr}{\operatorname{Tr}}
\newcommand{\dist}{\operatorname{dist}}
\newcommand{\supp}{\operatorname{supp}}

\title{Renewal Yang--Mills Mass-Gap Research Notes\\
Compact Continuation, Version V1}
\author{Proof-indexed continuation after archive \texttt{V65\_f1}}
\date{8 September 2026}

\begin{document}
\maketitle
\tableofcontents

\chapter{Context and cumulative theorem ledger}
\label{ch:n1v1-context}

\section{Purpose of the compact continuation}
\label{sec:n1v1-purpose}

The objective is a rigorous construction of four-dimensional quantum
Yang--Mills theory for a compact nonabelian gauge group, followed by
Osterwalder--Schrader reconstruction and a proof that the physical
Hamiltonian has a strictly positive spectral gap above its vacuum.
The programme is constructive and multiscale.  It seeks exact
gauge-compatible block pushforwards at finite cutoff, a separation of
gauge, toron, retained, and eliminated modes before inversion, a
volume-uniform expansion of the nonlinear and bad-field sectors, a
controlled marginal trajectory, and eventual entry into a terminal
mixing domain whose Euclidean decay yields the Hamiltonian gap.

This file is a size-control rollover, not a mathematical restart.  It
records the major results that remain live, identifies the present
dependency boundary, and then continues the proof programme.  Detailed
proofs through the preceding active V65 remain in the frozen archive.

\begin{record}[Immediate archive]
\label{rec:n1v1-archive}
The immediate predecessor has been frozen as
\[
\texttt{Renewal\_Yang\_Mills\_Mass\_Gap\_Research\_Notes\_V65\_f1.tex}.
\]
Its validated record is
\[
\begin{array}{c|c}
\text{source lines}&27874\\
\text{bytes}&982002\\
\text{SHA--256}&
\texttt{7A5AC49398BF5FE00E454FA469459BC6216E6E1109C79FC376A15576C72C858C}.
\end{array}
\]
The present compact run restarts at V1.  Each successor is to append
one proof note to the complete active source and replace the preceding
active version only after validation.  Every fifth version inspects the
then-current Standard-Model--Einstein and Navier--Stokes notes.  At
V100 the active file is frozen under the next archive suffix and a new
contextual V1 begins.
\end{record}

\begin{firewall}[No mass-gap claim]
\label{fire:n1v1-no-gap}
The full Yang--Mills mass gap is not proved.  In particular, the
notes have not yet constructed the complete ultraviolet-to-terminal
renormalization trajectory, removed all regulators, constructed the
continuum Osterwalder--Schrader measure, or transferred a uniform
Euclidean clustering exponent to the spectrum of the continuum
Hamiltonian.
\end{firewall}

\section{Major results obtained so far}
\label{sec:n1v1-results}

\begin{record}[Finite-cutoff foundations and permanent no-go results]
\label{rec:n1v1-foundations}
The archive chain proves finite-cutoff Peter--Weyl and heat-kernel
identities, exact Duhamel and forest interpolation formulae,
endpoint-factorizing BKAR representations, rooted-tree estimates,
moving-conditional score identities, and exact ground-transform
projection ledgers on their declared domains.  It also proves that
the following shortcuts are invalid:
\begin{enumerate}[label=\textup{(\roman*)},leftmargin=3.5em]
\item an unconditioned small bad-set probability does not imply a
uniform conditional influence row;
\item pair covariance does not control the higher conditional
cumulants required by differentiated currents;
\item reflection positivity or static decay alone does not identify a
Hamiltonian gap;
\item gauge gradients, torons, retained modes, and other massless
directions cannot be hidden inside a uniformly gapped determinant; and
\item the exact RG differential contains a unit retained-pullback
sector, so contraction is available only after calibrated marginal and
redundant directions have been removed.
\end{enumerate}
\end{record}

\begin{record}[Exact block geometry and Gaussian sector]
\label{rec:n1v1-block-geometry}
The programme constructs exact retained/eliminated block
disintegrations, finite-range cochain liftings, reduced parity-tree
fibres, and associative fixed-factor pushforwards.  At a stationary
conditional centre the physical Wilson Hessian is
\[
                         D^2w(0)=d_1^*d_1,                 \tag{1.1}
\]
with gauge gradients and flat cohomological modes retained rather than
inverted.  On the reduced transverse space it has a positive
fixed-block floor.  The minimum-energy interpolation
\[
 J=L^{-1}Q^*(QL^{-1}Q^*)^{-1}                              \tag{1.2}
\]
gives an exact energy-orthogonal Gaussian disintegration and zero
Gaussian retained feedback on the conditionally centred detail space.
Canonical four-dimensional dilation leaves
\(\Tr F^2\) and \(\Tr F\widetilde F\) marginal and contracts declared
local curvature jets of dimension at least six by at most \(b^{-2}\).
\end{record}

\begin{record}[Determinant, Ward, and current ledgers]
\label{rec:n1v1-response}
For a finite-range uniformly gapped eliminated Hessian, the first two
log-determinant derivatives
\[
 D\log\det L=\Tr(L^{-1}DL),                                \tag{1.3}
\]
\[
 D_2D_1\log\det L
 =
 \Tr\!\left(L^{-1}D_2D_1L
 -L^{-1}D_2L\,L^{-1}D_1L\right)                            \tag{1.4}
\]
have exponentially local fixed-dictionary response rows.  The notes
separate the perimeter contact, the universal transverse moment, and
the two-dimensional hypercubic quartic Ward space.  They construct
explicit third and fourth Wilson jets, identify the redundant
equation-of-motion direction, and prove that a third differentiated
conditional current genuinely requires a fourth connected cumulant.
The pushed Jacobian Hessian has consequently been corrected to retain
the third current jet rather than being closed from only pair data.
\end{record}

\begin{record}[Reduced fibre, local vacuum, and good/bad control]
\label{rec:n1v1-fibre}
For every retained background \(W\), the reduced filling satisfies the
localized coercivity estimate
\[
 \sum_{y\in Y}D_y(z;W)
 \le20\sum_{p\in{\cal P}(Y)}\|U_p(z,W)-I\|_F^2.             \tag{1.5}
\]
For a coarse-flat background, including a fixed commuting toron, its
zero set is the unique moving section; for a nonflat background the
zero set is empty.  Thus an \(r\)-bad connected block set pays Wilson
action at least \(r^2|Y|/120\).

With
\[
 t=\frac{\beta}{120},\qquad r_\beta=\beta^{-2/5},           \tag{1.6}
\]
the dominated product-Haar reference has a nonnegative residual.  The
far-bad primitive is bounded by
\[
 q_{\rm far}(\beta)
 =
 C_-^{-45}\left(\frac{\beta}{120}\right)^{180}
 \exp\!\left\{-\frac{\beta^{1/5}}{120}\right\},             \tag{1.7}
\]
whereas the flat or commuting-toron good-chart Wilson remainder obeys
\[
 \eta_{\rm W}(\beta)=7680\pi^3\beta^{-1/5}.                 \tag{1.8}
\]
The exact hard-core four-mark polymer theorem closes once the total
owned activity \(u\) satisfies
\[
                  6401e^{1+\mu_*}u\le1.                    \tag{1.9}
\]
Moving Haar coordinates eliminate artificial derivatives of the
selector, but not the physical stationary-centre and current scores.
\end{record}

\begin{record}[Latest Jacobian and soft-selector results]
\label{rec:n1v1-latest}
The normalized score calculus, Haar Casimir identity, and the local
cubic compensation computation show that the fine compensated owner
has no \(F^2\) Hessian and that the local Jacobian marginal can be
removed as a coordinate coefficient.  A positive Gaussian forest
formula has been transferred to the reduced Yang--Mills covariance.
The compact smooth selector was proved unsuitable for an all-order
forest because a star tree differentiates one vertex arbitrarily many
times.  It was replaced by the exact analytic split
\[
 g_{\beta,y}=e^{-D_y/r_\beta^2},\qquad
 b_{\beta,y}=1-g_{\beta,y}.                                \tag{1.10}
\]
The preceding note bounded the soft-bad primitive by
\(O(\beta^{-1/5})\), but still treated the all-good factor as an
analytic quadratic vertex.  The first new theorem below removes that
unnecessary vertex exactly.
\end{record}

\section{Live dependency chain}
\label{sec:n1v1-frontier}

\begin{assessment}[Current frontier]
\label{ass:n1v1-frontier}
The shortest presently justified route is:
\begin{enumerate}[label=\textup{(\arabic*)},leftmargin=3.5em]
\item close the source-marked nonlinear cluster card with a
volume-independent strong-to-weak analytic estimate;
\item finish the raw pushed-Jacobian coefficient and compare the
compensated owner with the marginal Ward basis;
\item prove the calibrated stable-slice contraction, including the
anisotropic and topological rows;
\item iterate the exact fixed-factor cocycle into a terminal
Dobrushin or bridge domain;
\item construct the regulator-independent reflection-positive
continuum measure and perform OS reconstruction; and
\item transfer terminal Euclidean decay to a positive continuum
Hamiltonian gap.
\end{enumerate}
\end{assessment}

\chapter{V1: exact soft-reference rebasing and degree-faithful forests}
\label{ch:n1v1-new}

\section{Exact absorption of the all-good soft selector}
\label{sec:n1v1-rebase}

Let \(\Lambda\) be a finite block set.  The reduced links are
partitioned into disjoint sets \({\cal C}_y\), \(y\in\Lambda\), with
\(|{\cal C}_y|=45\).  In moving coordinates put
\[
 d_e(z;W)=\|z_eH_e(W)^{-1}-I\|_F^2,\qquad
 D_y=\sum_{e\in{\cal C}_y}d_e.                              \tag{1.11}
\]
For normalized Haar measure \(dU\), define
\[
 I(u)=\int_{SU(3)}e^{-u\|U-I\|_F^2}\,dU,                   \tag{1.12}
\]
and let \(\gamma_{u,W}\) be the product probability on all reduced
links with one-link density \(I(u)^{-1}e^{-u d_e}\).

Keep \(t=\beta/120\), set
\[
 s=r_\beta^{-2}=\beta^{4/5},                               \tag{1.13}
\]
and write \(g_y=e^{-sD_y}\), \(b_y=1-g_y\).

\begin{theorem}[Exact soft-pattern rebase]
\label{thm:n1v1-rebase}
For every block pattern \(B\subset\Lambda\) and every integrable
function \(F\),
\[
\begin{split}
 &{\mathbb E}_{\gamma_{t,W}}\!
 \left[
 F\prod_{y\notin B}g_y\prod_{y\in B}b_y
 \right]                                                     \\
 &\quad=
 \left(\frac{I(t+s)}{I(t)}\right)^{45|\Lambda|}
 {\mathbb E}_{\gamma_{t+s,W}}\!
 \left[
 F\prod_{y\in B}\bigl(e^{sD_y}-1\bigr)
 \right].
                                                               \tag{1.14}
\end{split}
\]
The prefactor is independent of \(B\) and \(W\).  Consequently it
cancels from every normalized conditional expectation after summing
over \(B\).
\end{theorem}

\begin{proof}
The elementary identity
\[
 \prod_{y\notin B}e^{-sD_y}
 \prod_{y\in B}(1-e^{-sD_y})
 =
 e^{-s\sum_yD_y}
 \prod_{y\in B}(e^{sD_y}-1)                                \tag{1.15}
\]
first extracts the same \(e^{-sd_e}\) from every reduced link.
For one link,
\[
 d\gamma_{t,W}(z_e)\,e^{-sd_e(z;W)}
 =
 \frac{I(t+s)}{I(t)}\,d\gamma_{t+s,W}(z_e).                 \tag{1.16}
\]
Multiply (1.16) over the \(45|\Lambda|\) disjoint links and use
(1.15).  This proves (1.14).  Both \(I(u)\) and the change to moving
coordinates are independent of \(W\) by Haar invariance.

Finally, the prefactor in (1.14) is common to all patterns.  The
numerator and denominator of a normalized conditional expectation
contain the same factor, so it cancels before division.
\end{proof}

\begin{corollary}[The soft-good factor is not a forest vertex]
\label{cor:n1v1-no-soft-good-vertex}
In the complete subset expansion, the product of all soft-good
factors may be absorbed exactly into the stronger product reference
\(\gamma_{t+s,W}\).  The only selector insertion attached to a
nonempty bad pattern is
\[
                         h_y=e^{sD_y}-1\ge0.                \tag{1.17}
\]
Thus no derivative of \(g_y\), no selector Cauchy radius, and no
soft quadratic vertex is required in the Gaussian forest bank.
\end{corollary}

\begin{proof}
This is exactly the pattern-independent change of reference in
Theorem~\ref{thm:n1v1-rebase}.  It is important that the
\({\cal C}_y\) form a disjoint partition; otherwise the exponent in
(1.15) would carry link multiplicities.
\end{proof}

\section{A quantitative activity for the rebased bad insertion}
\label{sec:n1v1-soft-activity}

Retain the one-link constants
\[
 C_-=\frac{3e^{-43/32}}{1024\pi^8},\qquad
 C_+=\frac{25\pi^7}{512},\qquad
 A_1=\frac{32C_+}{eC_-}.                                   \tag{1.18}
\]
The archived one-link moment estimate is
\[
 {\mathbb E}_{\gamma_u}\|U-I\|_F^2\le\frac{A_1}{u},
 \qquad u\ge2.                                             \tag{1.19}
\]

\begin{lemma}[Normalizer-ratio bound]
\label{lem:n1v1-ratio}
For \(t\ge2\) and \(s\ge0\),
\[
 1\le\frac{I(t)}{I(t+s)}
 \le\left(1+\frac{s}{t}\right)^{A_1}
 \le e^{A_1s/t}.                                          \tag{1.20}
\]
\end{lemma}

\begin{proof}
Differentiation under the compact Haar integral gives
\[
 -\frac{d}{du}\log I(u)
 ={\mathbb E}_{\gamma_u}\|U-I\|_F^2.                       \tag{1.21}
\]
The integrand in (1.12) is decreasing in \(u\), giving the left
inequality.  Integrate (1.21) from \(t\) to \(t+s\) and use (1.19):
\[
 \log\frac{I(t)}{I(t+s)}
 \le A_1\int_t^{t+s}\frac{du}{u}
 =A_1\log\left(1+\frac{s}{t}\right)
 \le\frac{A_1s}{t}.
\]
Exponentiation proves (1.20).
\end{proof}

\begin{theorem}[Rebased soft-bad primitive]
\label{thm:n1v1-rebased-bad}
Let \(R_{\beta,W}\ge0\) be the exact dominated residual.  For every
block set \(Y\subset\Lambda\), \(t=\beta/120\), \(s=\beta^{4/5}\),
and \(\beta\ge240\),
\[
 {\mathbb E}_{\gamma_{t+s,W}}
 \left[
 e^{-R_{\beta,W}}\prod_{y\in Y}(e^{sD_y}-1)
 \right]
 \le q_{\rm soft}(\beta)^{|Y|},                            \tag{1.22}
\]
where
\[
\begin{split}
 q_{\rm soft}(\beta)
 &:=
 \left(1+120\beta^{-1/5}\right)^{45A_1}-1                 \\
 &\le
 \exp\!\left\{
 \frac{172800C_+}{eC_-}\,\beta^{-1/5}
 \right\}-1.
                                                               \tag{1.23}
\end{split}
\]
In particular \(q_{\rm soft}(\beta)=O(\beta^{-1/5})\).
\end{theorem}

\begin{proof}
Because \(R_{\beta,W}\ge0\) and every factor in (1.17) is nonnegative
on the compact real fibre, discard \(e^{-R_{\beta,W}}\).
The sets \({\cal C}_y\) are disjoint, so the remaining expectation
factorizes over \(y\in Y\).  For one block,
\[
\begin{split}
 {\mathbb E}_{\gamma_{t+s,W}}e^{sD_y}
 &=
 \left(\frac{I(t)}{I(t+s)}\right)^{45},                    \\
 {\mathbb E}_{\gamma_{t+s,W}}(e^{sD_y}-1)
 &=
 \left(\frac{I(t)}{I(t+s)}\right)^{45}-1.                  \tag{1.24}
\end{split}
\]
Apply (1.20), insert \(s/t=120\beta^{-1/5}\), and use
\(45\cdot120A_1=172800C_+/(eC_-)\).  This proves (1.22)--(1.23).
\end{proof}

\begin{remark}[Two genuinely different bad primitives]
\label{rem:n1v1-two-bad}
The soft-transition activity (1.23) is polynomially small and
analytic after rebasing.  The far-bad activity (1.7) is
stretched-exponentially small and comes from geometric coercivity.
They are not interchangeable owner types and must remain separate in
the final activity ledger.
\end{remark}

\begin{firewall}[What rebasing does not prove]
\label{fire:n1v1-rebase-limit}
Theorem~\ref{thm:n1v1-rebase} removes the selector from the forest
bank, but it does not by itself make the Wilson residual a small
interaction.  The residual still contains the physical coupled
quadratic form.  It must be recentered at the exact stationary point,
its quadratic part must be absorbed into the reduced Gaussian
precision, and the remaining analytic vertex must be estimated in a
volume-uniform norm.
\end{firewall}

\section{A one-time analytic-radius estimate}
\label{sec:n1v1-radius}

The earlier conditional fixed-tree estimate omitted a degree
factorial.  A star tree shows that this cannot be ignored in an
ordinary analytic norm.  The factorial is harmless only after the
complete labelled-tree sum is evaluated.

\begin{definition}[Uniform Taylor--Wiener norm]
\label{def:n1v1-wiener}
Let \(F\) be holomorphic in an \(R\)-polytube around
\(\mathbb R^q\).  Define
\[
 \|F\|_R
 :=
 \sup_{\varphi\in\mathbb R^q}
 \sum_{\alpha\in\mathbb N^q}
 \frac{R^{|\alpha|}}{\alpha!}
 |\partial^\alpha F(\varphi)|.                             \tag{1.25}
\]
For \(0<r<R\), write \(\delta=R-r\).
\end{definition}

\begin{lemma}[All incident derivatives spend the radius once]
\label{lem:n1v1-one-radius}
For every multi-index \(\nu\),
\[
                    \|\partial^\nu F\|_r
 \le\frac{\nu!}{\delta^{|\nu|}}\|F\|_R
 \le\frac{|\nu|!}{\delta^{|\nu|}}\|F\|_R.                  \tag{1.26}
\]
\end{lemma}

\begin{proof}
Set \(\gamma=\alpha+\nu\).  The summand in
\(\|\partial^\nu F\|_r\) is
\[
 \frac{r^{|\alpha|}}{\alpha!}
 |\partial^{\alpha+\nu}F|
 =
 \nu!\binom{\gamma}{\nu}
 \frac{r^{|\gamma-\nu|}}{R^{|\gamma|}}
 \frac{R^{|\gamma|}}{\gamma!}|\partial^\gamma F|.          \tag{1.27}
\]
The multi-index binomial theorem with \(R=r+\delta\) gives
\[
 \binom{\gamma}{\nu}
 r^{|\gamma-\nu|}\delta^{|\nu|}
 \le R^{|\gamma|}.                                        \tag{1.28}
\]
Insert (1.28) into (1.27), sum over \(\alpha\), and take the
supremum over the real centre.  Since \(\nu!\le|\nu|!\), (1.26)
follows.  All derivatives at one vertex were estimated together;
there is no iterative chain \(R>R-\delta>R-2\delta>\cdots\).
\end{proof}

\section{The correct degree-weighted forest theorem}
\label{sec:n1v1-degree-forest}

Let \(C_{\xi\eta}\) be the reduced positive Gaussian covariance, with
coordinate labels \(\xi,\eta\) carrying their block positions.  Put
\[
 G_\mu
 =
 \sup_\xi\sum_\eta
 e^{\mu d(\xi,\eta)}|C_{\xi\eta}|<\infty.                  \tag{1.29}
\]
Let \(h_*\) be the finite maximum number of owner-coordinate and
orientation choices incident to one covariance endpoint, and define
\[
                         K_\mu=\frac{h_*G_\mu}{\delta^2}.   \tag{1.30}
\]

\begin{definition}[Anchored absolute tree norm]
\label{def:n1v1-tree-norm}
For a tree \(T\), choose one owner as root.  Expand every edge
contraction into its two coordinate endpoints
\((\xi_e,\eta_e)\), and let
\[
                         \ell_T=\sum_{e\in E(T)}
                         d(\xi_e,\eta_e).
\]
The norm \(\|{\cal A}_T\|_{\mu,r}\) means: take the absolute value of
the differentiated integrand before the endpoint sums; sum all
nonroot owner positions, finite owner types, orientations, and edge
endpoints; insert \(e^{\mu\ell_T}\); integrate the forest parameters
and the positive weakened Gaussian probability; and finally take the
supremum over the root position and over the real field centres in the
weak Taylor--Wiener norm.  This is an absolute majorant norm; no
cancellation is used in it.
\end{definition}

\begin{lemma}[Fixed tree with its necessary degree weight]
\label{lem:n1v1-fixed-tree}
Let \(T\) be a labelled tree on \(n\ge2\) local vertices
\(F_1,\ldots,F_n\), each satisfying \(\|F_i\|_R\le z\).
In the norm of Definition~\ref{def:n1v1-tree-norm}, the positive
Gaussian forest amplitude obeys
\[
 \|{\cal A}_T(F_1,\ldots,F_n)\|_{\mu,r}
 \le
 z^nK_\mu^{\,n-1}\prod_{i=1}^nd_T(i)!.                    \tag{1.31}
\]
\end{lemma}

\begin{proof}
Every forest edge contributes one covariance entry and one derivative
at each endpoint.  If \(d_i=d_T(i)\), all derivatives landing at
vertex \(i\) form a multi-index \(\nu_i\) of total degree \(d_i\).
Lemma~\ref{lem:n1v1-one-radius} bounds that vertex by
\[
                         z\,d_i!\,\delta^{-d_i}.            \tag{1.32}
\]
Since \(\sum_i d_i=2(n-1)\), the product of (1.32) supplies
\(\delta^{-2(n-1)}\prod_i d_i!\).

Root \(T\) at the anchored owner and remove terminal vertices one at a
time.  At each removal, the sum over the terminal position and its
finite endpoint types is at most \(h_*G_\mu\) by (1.29); the
exponential edge weights dominate the distance from the anchor to the
removed support.  There are \(n-1\) removals.  Positivity makes every
interpolated Gaussian measure a probability, so its expectation
cannot enlarge the uniform real-centre derivative bound.  Multiplying
the vertex, radius, and edge factors gives (1.31).
\end{proof}

\begin{lemma}[Exact degree-weighted Cayley sum]
\label{lem:n1v1-prufer}
For \(n\ge2\),
\[
 \sum_{T\ {\rm tree\ on}\ [n]}
 \prod_{i=1}^nd_T(i)!
 =
 (n-2)!\binom{3n-3}{n-2}.                                 \tag{1.33}
\]
\end{lemma}

\begin{proof}
A labelled tree corresponds to a Pruefer word of length \(n-2\).
If the label \(i\) occurs \(m_i\) times, then
\(d_T(i)=m_i+1\), and the number of words with this occurrence vector
is \((n-2)!/\prod_i m_i!\).  Hence the left side of (1.33) is
\[
 (n-2)!
 \sum_{\substack{m_1,\ldots,m_n\ge0\\
                  \sum_i m_i=n-2}}
 \prod_{i=1}^n(m_i+1).                                    \tag{1.34}
\]
Since
\[
 \sum_{m\ge0}(m+1)x^m=(1-x)^{-2},
\]
the sum in (1.34) is the coefficient of \(x^{n-2}\) in
\((1-x)^{-2n}\), namely \(\binom{3n-3}{n-2}\).
\end{proof}

\begin{theorem}[Connected forest sum with analytic degree loss]
\label{thm:n1v1-connected}
Under the hypotheses of Lemma~\ref{lem:n1v1-fixed-tree}, if
\[
                              8zK_\mu<1,                   \tag{1.35}
\]
then the complete connected Gaussian exponential converges absolutely
in the weak norm and satisfies
\[
 \left\|
 \sum_{n\ge1}\frac1{n!}\kappa_C(F,\ldots,F)
 \right\|_{\mu,r}
 \le
 z+\frac{4z^2K_\mu}{1-8zK_\mu}.                            \tag{1.36}
\]
\end{theorem}

\begin{proof}
For \(n\ge2\), combine (1.31) and (1.33).  The \(n\)-th connected
coefficient is at most
\[
 z^nK_\mu^{n-1}
 \frac1{n(n-1)}
 \binom{3n-3}{n-2}.                                       \tag{1.37}
\]
Using
\[
 \binom{3n-3}{n-2}\le2^{3n-3},
 \qquad n(n-1)\ge2,
\]
the expression in (1.37) is at most
\[
                         \frac1{16K_\mu}(8zK_\mu)^n.       \tag{1.38}
\]
Summing (1.38) from \(n=2\) to infinity gives the second term in
(1.36); the \(n=1\) term is bounded by \(z\).
\end{proof}

\begin{corollary}[Four marked sources]
\label{cor:n1v1-four-marks}
Suppose the vertex bank is holomorphic in four source variables on
\(|s_j|<\rho_j\), and the strong norm \(z\) is uniform on that
polydisc.  The connected coefficient containing all four marks obeys
the right side of (1.36), multiplied by
\[
                         \prod_{j=0}^3\rho_j^{-1}.          \tag{1.39}
\]
No source-radius power is absorbed into \(K_\mu\).
\end{corollary}

\begin{proof}
Theorem~\ref{thm:n1v1-connected} converges normally and uniformly on
the source polydisc.  Apply the four-variable Cauchy estimate once to
the complete connected sum.
\end{proof}

\begin{correction}[Correction to the V65 fixed-tree card]
\label{correction:n1v1-v65}
The V65 conditional estimate without
\(\prod_i d_T(i)!\) is stronger than the ordinary analytic
Taylor--Wiener norm proves.  The correct fixed-tree statement is
(1.31).  The exact Pruefer calculation (1.33) shows that the omitted
factor costs only an exponential tree constant, changing the
sufficient smallness condition from \(ezK<1\) to \(8zK<1\), not
destroying convergence.
\end{correction}

\section{Consequences for the Yang--Mills gate}
\label{sec:n1v1-gate}

\begin{proposition}[Corrected conditional closure of the good branch]
\label{prop:n1v1-good-gate}
Assume that, after stationary recentering and absorption of the exact
quadratic form, every unmarked Yang--Mills remainder owner satisfies
the uniform strong card
\[
 z_\beta
 \le C_{\rm an}
 \left\{
 e^{7680\pi^3\beta^{-1/5}}-1
 \right\},                                                 \tag{1.40}
\]
with fixed \(C_{\rm an}\), \(R>r>0\), \(h_*\), and \(G_\mu\).
Then the good connected activity satisfies
\[
 u_{\rm good}(\beta)
 \le C_{\rm own}
 \left\{
 z_\beta+
 \frac{4z_\beta^2K_\mu}{1-8z_\beta K_\mu}
 \right\}                                                  \tag{1.41}
\]
whenever \(8z_\beta K_\mu<1\).  In particular it tends to zero as
\(\beta\to\infty\).
\end{proposition}

\begin{proof}
Theorem~\ref{thm:n1v1-connected} gives the expression in braces.
The fixed owner conversion contributes \(C_{\rm own}\).  Equation
(1.40) implies \(z_\beta=O(\beta^{-1/5})\), so both the denominator
condition and convergence to zero hold for sufficiently large
\(\beta\).
\end{proof}

\begin{corollary}[Asymptotic numerical polymer gate]
\label{cor:n1v1-asymptotic-gate}
Under the hypothesis of Proposition~\ref{prop:n1v1-good-gate}, for
every fixed \(\mu_*>0\) there exists \(\beta_0<\infty\) such that
\[
\begin{split}
 6401e^{1+\mu_*}\Bigg[
 &q_{\rm far}(\beta)+q_{\rm soft}(\beta)\\
 &+
 C_{\rm own}\left\{
 z_\beta+
 \frac{4z_\beta^2K_\mu}{1-8z_\beta K_\mu}
 \right\}
 \Bigg]\le1                                                \tag{1.42}
\end{split}
\]
for all \(\beta\ge\beta_0\).
\end{corollary}

\begin{proof}
The far-bad term (1.7) tends to zero stretched exponentially.
Theorem~\ref{thm:n1v1-rebased-bad} gives
\(q_{\rm soft}(\beta)\to0\), and
Proposition~\ref{prop:n1v1-good-gate} gives
\(u_{\rm good}(\beta)\to0\).  The remaining factor in (1.42) is
fixed.
\end{proof}

\begin{assessment}[Progress in compact V1]
\label{ass:n1v1-status}
This note makes two rigorous advances at the V65 bottleneck.
First, it exactly absorbs the all-good soft selector into a stronger
product-Haar reference and proves the resulting bad insertion has
activity \(O(\beta^{-1/5})\); selector derivatives therefore disappear
from the forest bank.  Second, it supplies the missing one-time
strong-to-weak radius estimate, retains the unavoidable vertex-degree
factorials, sums them exactly by Pruefer words, and proves the
connected bound (1.36).  Thus arbitrary star degrees no longer cause
a circular analytic-radius argument.
\end{assessment}

\begin{decision}[V2 target]
\label{dec:n1v1-next}
V2 should populate the remaining hypothesis (1.40) for the exact
Yang--Mills remainder bank.  It should:
\begin{enumerate}[label=\textup{(\roman*)},leftmargin=3.5em]
\item rechart at the analytic stationary centre for a nonflat retained
source;
\item absorb the complete coupled quadratic form into the reduced
Gaussian precision;
\item estimate the Wilson, Haar, current, and three score owners in
the uniform Taylor--Wiener norm (1.25); and
\item display every source radius and covariance row constant.
\end{enumerate}
If a uniform real-centre norm fails, the proof must go upstream to a
Gaussian-regulated analytic norm or to a different exact reference;
it must not reintroduce a compact selector into an all-order forest.
\end{decision}

\begin{firewall}[Claim boundary after compact V1]
\label{fire:n1v1-claim}
The exact rebase and degree-weighted connected theorem are proved.
The concrete Yang--Mills strong card (1.40) remains an assumption, so
the numerical gate (1.42) is conditional.  This note does not evaluate
the remaining pushed-Jacobian coefficient, close the marginal and
anisotropic RG rows, construct the complete RG trajectory, remove the
regulators, construct the continuum OS theory, or prove the
Yang--Mills mass gap.
\end{firewall}

% =============================================================================
% BEGIN APPENDED COMPACT-SERIES V2 MATERIAL
% =============================================================================

\chapter{V2: the global-chart no-go and a compact heat-kernel forest}
\label{ch:n1v2}

\section{Testing the remaining V1 hypothesis}
\label{sec:n1v2-test}

The strong card (1.40) was deliberately left as a hypothesis.  Its
uniform Taylor--Wiener norm takes a supremum over every real
Lie-algebra centre.  That is compatible with an ordinary entire
perturbation of a Gaussian, but it is not compatible with subtracting
an unbounded quadratic form from an exact action living on a compact
group.  The obstruction occurs already in the zeroth derivative.

Let \(E\) be the finite-dimensional reduced real Lie algebra in one
finite-cutoff fibre, let
\[
 \mathcal E_g:E\longrightarrow{\cal M},\qquad
 \mathcal E_g(a)=\exp(ga)H(W),\qquad g=\beta^{-1/2},         \tag{2.1}
\]
denote the componentwise exponential followed by the moving section,
and let \({\cal M}\) be the compact reduced-link manifold.  Write
\[
 A_\beta(a)=\beta V_M(\mathcal E_g(a),W),\qquad
 Q(a)=\frac12\langle a,Ha\rangle,                           \tag{2.2}
\]
where \(H>0\) is the exact reduced physical Hessian.

\begin{theorem}[Global quadratic-subtraction no-go]
\label{thm:n1v2-global-nogo}
Set
\[
                   F_\beta(a)=e^{-[A_\beta(a)-Q(a)]}-1.    \tag{2.3}
\]
For every fixed \(\beta<\infty\),
\[
                         \sup_{a\in E}|F_\beta(a)|=\infty.  \tag{2.4}
\]
Consequently \(F_\beta\) has infinite uniform
Taylor--Wiener norm (1.25) for every \(R\ge0\), and the proposed
global interpretation of (1.40) is false.
\end{theorem}

\begin{proof}
Compactness of \({\cal M}\) and continuity of the Wilson action give
\[
                         0\le A_\beta(a)\le\beta V_{\max}  \tag{2.5}
\]
for all \(a\in E\).  Choose a unit eigenvector \(v\) of \(H\) with
eigenvalue \(\lambda>0\).  Along \(a=rv\),
\[
 A_\beta(rv)-Q(rv)
 \le\beta V_{\max}-\frac{\lambda r^2}{2}.                  \tag{2.6}
\]
Therefore
\[
 F_\beta(rv)
 \ge
 \exp\!\left\{\frac{\lambda r^2}{2}-\beta V_{\max}\right\}-1
 \longrightarrow\infty.                                   \tag{2.7}
\]
The zeroth summand in (1.25) is already unbounded.
\end{proof}

\begin{remark}[The obstruction is compactness, not a weak Taylor
estimate]
\label{rem:n1v2-compactness}
No improvement of the cubic constant in (1.8) repairs (2.4).
The local Taylor estimate is correct on
\(\|a\|=O(\beta^{1/10})\).  The failure comes from extending the
unbounded tangent quadratic \(Q\) through all logarithmic sheets of a
compact group.
\end{remark}

\section{Why a scalar Gaussian regulator is still insufficient}
\label{sec:n1v2-regulator-nogo}

One might try to replace (1.25) by a weighted norm.  The following
matrix statement shows that the most direct Gaussian weight has no
admissible parameter.

\begin{theorem}[Regulator versus integrability obstruction]
\label{thm:n1v2-regulator-nogo}
Let \(d\gamma_H\) be the centred Gaussian probability on \(E\) with
covariance \(H^{-1}\).  For a symmetric matrix \(B\), consider a
zeroth-order regulator
\[
 \|F\|_{0,B}
 =
 \sup_{a\in E}e^{-\langle a,Ba\rangle}|F(a)|.              \tag{2.8}
\]
If \(\|F_\beta\|_{0,B}<\infty\), then
\[
                              B\ge\frac12H.                \tag{2.9}
\]
On the other hand,
\[
 \int_Ee^{\langle a,Ba\rangle}\,d\gamma_H(a)<\infty
 \quad\Longleftrightarrow\quad
 I-2H^{-1/2}BH^{-1/2}>0.                                  \tag{2.10}
\]
Conditions (2.9) and (2.10) are incompatible.  This remains true
before any support-overlap loss is included.
\end{theorem}

\begin{proof}
Fix a unit vector \(v\).  From (2.5),
\[
 e^{-\langle rv,Brv\rangle}|F_\beta(rv)|
 \ge
 \exp\!\left\{
 r^2\left(\frac12\langle v,Hv\rangle-\langle v,Bv\rangle\right)
 -\beta V_{\max}
 \right\}-e^{-r^2\langle v,Bv\rangle}.                     \tag{2.11}
\]
Boundedness for every \(v\) forces (2.9).

Diagonalize
\[
                         S=H^{-1/2}BH^{-1/2}.
\]
After the change \(x=H^{1/2}a\), the integral in (2.10) is a product
of one-dimensional Gaussian integrals
\[
 \int_{\mathbb R}
 \exp\!\left\{-\frac12(1-2s_j)x_j^2\right\}
 \frac{dx_j}{\sqrt{2\pi}}.                                \tag{2.12}
\]
It is finite exactly when every eigenvalue \(s_j<1/2\), proving
(2.10).  But (2.9) gives \(S\ge I/2\).
\end{proof}

\begin{correction}[Correct scope of the V1 forest theorem]
\label{cor:n1v2-v1-scope}
The degree-weighted forest Theorem~\ref{thm:n1v1-connected} is valid.
Its proposed Yang--Mills input (1.40), however, cannot be the global
norm of the exact quadratic-subtracted compact action.  It can only
be a chart-local card accompanied by a separate exact large-field
extraction, or a global card relative to a compact reference which
already contains the logarithmic sheets.
\end{correction}

\section{A compact reference which contains all logarithmic sheets}
\label{sec:n1v2-heat-reference}

The no-go suggests moving the solvable reference from the tangent
space to the compact reduced fibre before taking a forest.  This
section constructs the exact positive interpolation required for
that route.

Let \(K=SU(3)^m\) with normalized Haar probability \(dm\), where
\(m\) is the number of reduced links in the finite fibre.  Let
\(\{X_\xi\}_{\xi\in{\cal I}}\) be a real left-invariant orthonormal
frame and let \(C=(C_{\xi\eta})\) be a strictly positive symmetric
matrix.  Define on \(C^\infty(K)\)
\[
 {\cal L}_C
 =
 \frac12\sum_{\xi,\eta\in{\cal I}}
 C_{\xi\eta}X_\xi X_\eta.                                 \tag{2.13}
\]
The fields \(X_\xi\) are skew-adjoint in Haar measure, so
\({\cal L}_C\) is nonpositive.

\begin{proposition}[Compact covariance reference]
\label{prop:n1v2-heat-reference}
For every \(g>0\), the closure of \({\cal L}_C\) generates a
positivity-preserving, constant-preserving, self-adjoint contraction
semigroup
\[
                         P_{g,C}=e^{g^2{\cal L}_C}.         \tag{2.14}
\]
There is a smooth strictly positive density \(k_{g,C}\) such that
\[
 [P_{g,C}f](I)=\int_Kf(U)k_{g,C}(U)\,dm(U)
 =:{\mathbb E}_{\nu_{g,C}}f.                              \tag{2.15}
\]
If a gauge, lattice, or reflection automorphism preserves \(C\), it
also preserves \(\nu_{g,C}\).
\end{proposition}

\begin{proof}
Choose a real matrix \(S\) with \(C=SS^T\), and put
\[
                         Y_r=\sum_\xi S_{\xi r}X_\xi.
\]
Then
\[
                         {\cal L}_C=\frac12\sum_rY_r^2,     \tag{2.16}
\]
and its Dirichlet form is
\[
 -\langle f,{\cal L}_Cf\rangle_{L^2(dm)}
 =
 \frac12\sum_r\|Y_rf\|_2^2\ge0.                            \tag{2.17}
\]
The Friedrichs closure therefore gives a self-adjoint contraction
semigroup.  The sum-of-squares diffusion construction gives
positivity and preservation of constants.  Strict positivity of
\(C\) makes (2.16) elliptic; on the compact connected manifold \(K\),
elliptic regularity and the strong maximum principle give a smooth
strictly positive heat kernel.  An automorphism preserving \(C\)
commutes with (2.13), hence with its semigroup and its kernel based at
the identity.
\end{proof}

\begin{remark}[Why this reference evades Theorem
\ref{thm:n1v2-global-nogo}]
\label{rem:n1v2-heat-wrap}
The density \(k_{g,C}\) is a function on \(K\), not
\(\exp[-Q(a)]\) continued through every \(a\in E\).  All logarithmic
sheets and the conjugate locus are already summed by the heat
semigroup.  No unbounded quadratic is subtracted from a bounded
compact action.
\end{remark}

\section{Positive replica interpolation on the compact group}
\label{sec:n1v2-replica}

For \(n\) replicas, let \(x=(x_{ij})\) be a real positive semidefinite
matrix with \(x_{ii}=1\).  On \(K^n\), set
\[
 {\cal L}^{(n)}_x
 =
 \frac12\sum_{i,j=1}^nx_{ij}
 \sum_{\xi,\eta}C_{\xi\eta}
 X_{i,\xi}X_{j,\eta}.                                     \tag{2.18}
\]
For \(F_1,\ldots,F_n\in C^\infty(K)\), define
\[
 \Phi_n(x)
 =
 \left[
 e^{g^2{\cal L}^{(n)}_x}
 \prod_{i=1}^nF_i(U_i)
 \right]_{U_1=\cdots=U_n=I}.                              \tag{2.19}
\]

\begin{lemma}[Positivity and endpoint factorization]
\label{lem:n1v2-replica-positive}
For every such \(x\), (2.18) generates a Markov probability on
\(K^n\).  If \(x\) is block diagonal according to a partition
\({\cal P}\) of \([n]\), then \(\Phi_n(x)\) factorizes over the blocks
of \({\cal P}\).  At the two endpoint matrices,
\[
 \Phi_n(I_n)=\prod_{i=1}^n{\mathbb E}_{\nu_{g,C}}F_i,\qquad
 \Phi_n({\bf1}{\bf1}^T)
 =
 {\mathbb E}_{\nu_{g,C}}\prod_{i=1}^nF_i.                  \tag{2.20}
\]
\end{lemma}

\begin{proof}
The coefficient matrix of the vector fields in (2.18) is
\(x\otimes C\ge0\).  Factoring it as \(RR^T\) writes the generator as
a sum of squares, exactly as in (2.16), so it is Markov.  A block
diagonal \(x\) has no mixed differential operator between distinct
blocks, and the semigroup tensor-factorizes.

At \(x=I_n\), the replicas are independent copies of (2.14).  At
\(x={\bf1}{\bf1}^T\), the sum-of-squares representation uses the same
driving diffusion in every replica.  Starting from the diagonal
identity, all endpoints remain equal, which gives the second identity
in (2.20).
\end{proof}

For \(i<j\), introduce the cross generator
\[
 {\cal D}_{ij}
 =
 \sum_{\xi,\eta}C_{\xi\eta}
 X_{i,\xi}X_{j,\eta}.                                     \tag{2.21}
\]

\begin{lemma}[Exact Duhamel edge derivative]
\label{lem:n1v2-duhamel-edge}
For every off-diagonal interpolation variable \(x_{ij}\),
\[
\begin{split}
 \partial_{x_{ij}}e^{g^2{\cal L}^{(n)}_x}
 =
 g^2\int_0^1
 e^{(1-s)g^2{\cal L}^{(n)}_x}
 {\cal D}_{ij}
 e^{sg^2{\cal L}^{(n)}_x}\,ds.                            \tag{2.22}
\end{split}
\]
For distinct formal edge derivatives \(e_1,\ldots,e_k\),
\[
\begin{split}
 \partial_{e_1}\cdots\partial_{e_k}e^{g^2{\cal L}^{(n)}_x}
 &=
 g^{2k}\sum_{\pi\in S_k}
 \int_{0<s_1<\cdots<s_k<1}
 e^{(1-s_k)g^2{\cal L}_x}{\cal D}_{e_{\pi(k)}}             \\
 &\quad\cdots
 e^{(s_2-s_1)g^2{\cal L}_x}{\cal D}_{e_{\pi(1)}}
 e^{s_1g^2{\cal L}_x}\,ds_1\cdots ds_k.                   \tag{2.23}
\end{split}
\]
The sum of the volumes of the \(k!\) ordered simplices in (2.23) is
one.
\end{lemma}

\begin{proof}
Since \({\cal L}^{(n)}_x\) depends affinely on \(x_{ij}\) and
\(\partial_{x_{ij}}{\cal L}^{(n)}_x={\cal D}_{ij}\), the
variation-of-constants formula gives (2.22).  Iterating (2.22) and
ordering the insertion times gives (2.23).  Every simplex has volume
\(1/k!\).
\end{proof}

\begin{theorem}[Compact heat-kernel forest identity]
\label{thm:n1v2-compact-forest}
Let \(x^T(w)\) be the positive BKAR weakening matrix associated with
a labelled tree \(T\).  Then
\[
 \boxed{
 \kappa_{\nu_{g,C}}(F_1,\ldots,F_n)
 =
 \sum_{T\ {\rm tree\ on}\ [n]}
 \int_{[0,1]^{E(T)}}
 \left[
 \partial_T\Phi_n(x)
 \right]_{x=x^T(w)}\,dw.}                                 \tag{2.24}
\]
Every interpolated object in (2.24) is a positive compact-group
Markov probability, and each tree derivative has the exact
time-ordered representation (2.23).
\end{theorem}

\begin{proof}
Apply the BKAR forest formula to the smooth function \(\Phi_n(x)\)
of the off-diagonal entries.  Positivity of every weakening matrix
\(x^F(w)\), together with Lemma~\ref{lem:n1v2-replica-positive},
gives a Markov interpolation.  When a forest has connected components
\({\cal P}\), endpoint factorization gives the product of the
\(\Phi\)-functions on those components.  Möbius inversion on set
partitions selects the connected forests, which are precisely
spanning trees, and yields (2.24).  Formula (2.23) supplies the stated
edge representation.  Degenerate weakening matrices follow by adding
\(\varepsilon I\), applying the identity to the elliptic
regularization, and using strong semigroup convergence as
\(\varepsilon\downarrow0\).
\end{proof}

\section{Tangent matching without choosing a logarithmic sheet}
\label{sec:n1v2-tangent}

\begin{proposition}[Quantitative tangent covariance]
\label{prop:n1v2-tangent}
For \(f\in C^\infty(K)\),
\[
 {\mathbb E}_{\nu_{g,C}}f
 =
 f(I)+g^2{\cal L}_Cf(I)+{\cal R}_{g,C}(f),                 \tag{2.25}
\]
with
\[
 |{\cal R}_{g,C}(f)|
 \le\frac{g^4}{2}\|{\cal L}_C^2f\|_\infty.                 \tag{2.26}
\]
Thus \(\nu_{g,C}\) has tangent covariance \(g^2C\) at the identity
with a uniform fourth-order semigroup remainder.
\end{proposition}

\begin{proof}
Twice integrating
\[
 \frac{d}{dt}e^{t{\cal L}_C}f
 =e^{t{\cal L}_C}{\cal L}_Cf
\]
from \(t=0\) to \(t=g^2\) gives
\[
 e^{g^2{\cal L}_C}f
 =
 f+g^2{\cal L}_Cf+
 \int_0^{g^2}(g^2-s)e^{s{\cal L}_C}{\cal L}_C^2f\,ds.
                                                               \tag{2.27}
\]
The Markov semigroup is an \(L^\infty\) contraction.  Evaluation at
the identity and integration of \(g^2-s\) prove (2.25)--(2.26).
\end{proof}

\begin{assessment}[Progress after compact V2]
\label{ass:n1v2-status}
V2 proves that the remaining V1 global norm hypothesis cannot hold
for an exact compact action and that a direct Gaussian regulator is
at the sharp nonintegrable threshold.  Going upstream, it constructs
a compact heat-kernel reference carrying the physical tangent
covariance and proves an exact positive replica forest formula with
time-ordered Duhamel edge insertions.  This removes the logarithmic-
sheet obstruction without sacrificing positivity or exact connected
cumulants.
\end{assessment}

\begin{decision}[V3 target]
\label{dec:n1v2-next}
V3 should build the strong analytic norm for the compact forest
(2.24).  The required estimate is a volume-uniform bound on the
time-ordered products (2.23), with:
\begin{enumerate}[label=\textup{(\roman*)},leftmargin=3.5em]
\item rescaled invariant derivatives so the factors \(g^2\) per edge
cancel exactly against the two endpoint derivatives;
\item the weighted covariance row \(G_\mu\);
\item a single analytic-radius loss per vertex, including commutators
between the invariant derivatives and the heat semigroup; and
\item four marked source radii kept explicit.
\end{enumerate}
After that norm is available, the next comparison is between the
exact Wilson conditional density and the compact heat-kernel reference.
If that ratio is not uniformly small, its marginal quadratic part
must be transported by a covariance path rather than declared an
interaction.
\end{decision}

\begin{firewall}[Claim boundary after compact V2]
\label{fire:n1v2-claim}
The compact heat-kernel reference and forest identity are exact at
finite cutoff.  V2 does not yet prove a tree norm for the Duhamel
insertions, a small Wilson-to-heat-kernel Radon--Nikodym activity, or
the numerical polymer gate.  It also does not close the pushed
Jacobian, marginal trajectory, stable RG contraction, continuum
limit, OS reconstruction, or Yang--Mills mass gap.
\end{firewall}

\subsection{Parent record}

This V2 note was formed by appending the preceding material to the
validated compact V1, whose record was
\[
\begin{array}{c|c}
\text{lines}&784\\
\text{bytes}&27964\\
\text{SHA--256}&
\texttt{E54AC40849413570B3EBD015FB8F18DADBB2A0C0392A22CA330552DEEB1A441A}.
\end{array}
\]

% =============================================================================
% END APPENDED COMPACT-SERIES V2 MATERIAL
% =============================================================================

% =============================================================================
% BEGIN APPENDED COMPACT-SERIES V3 MATERIAL
% =============================================================================

\chapter{V3: an invariant analytic norm for the compact forest}
\label{ch:n1v3}

\section{Opposite invariant fields remove the commutator loss}
\label{sec:n1v3-opposite}

The Duhamel products (2.23) place heat semigroups between cross
generators.  Estimating them with the same invariant fields used in
the generator creates unnecessary commutators.  On a Lie group the
opposite invariant fields commute exactly, which gives the required
one-time radius estimate.

Let \(X_\xi^L\) be the left-invariant fields used in (2.13), and let
\(X_\xi^R\) be the corresponding right-invariant fields.  On every
group factor,
\[
 [X_\xi^L,X_\eta^R]=0.                                    \tag{3.1}
\]
Introduce the rescaled fields
\[
 {\mathscr X}_\xi^L=gX_\xi^L,\qquad
 {\mathscr X}_\xi^R=gX_\xi^R.                              \tag{3.2}
\]
Then one Duhamel edge is
\[
 g^2{\cal D}_{ij}
 =
 \sum_{\xi,\eta}C_{\xi\eta}
 {\mathscr X}_{i,\xi}^L{\mathscr X}_{j,\eta}^L.            \tag{3.3}
\]
Thus the two physical derivatives consume the prefactor \(g^2\)
exactly; no power of \(g^{-1}\) remains in an edge constant.

\begin{lemma}[Right derivatives commute with every weakened
semigroup]
\label{lem:n1v3-commute}
For every replica \(i\), coordinate \(\xi\), positive weakening
matrix \(x\), and \(t\ge0\),
\[
 {\mathscr X}_{i,\xi}^R e^{tg^2{\cal L}^{(n)}_x}
 =
 e^{tg^2{\cal L}^{(n)}_x}{\mathscr X}_{i,\xi}^R.           \tag{3.4}
\]
They also commute with every cross generator \({\cal D}_{jk}\).
\end{lemma}

\begin{proof}
Every term of \({\cal L}^{(n)}_x\) and every
\({\cal D}_{jk}\) is a polynomial in left-invariant fields.
Equation (3.1) therefore gives commutation on
\(C^\infty(K^n)\).  Commutation with the semigroup follows first on
this invariant core from its power series and then on smooth
functions by semigroup closure.
\end{proof}

\section{The compact scaled Taylor--Wiener norm}
\label{sec:n1v3-norm}

For an ordered word
\(\boldsymbol\xi=(\xi_1,\ldots,\xi_m)\), write
\[
 {\mathscr X}_{\boldsymbol\xi}^R
 =
 {\mathscr X}_{\xi_1}^R\cdots{\mathscr X}_{\xi_m}^R.
\]
If \(F\) depends on the finite coordinate set \(S\), put
\[
 {\mathfrak d}_m(F)
 =
 \sum_{\boldsymbol\xi\in S^m}
 \|{\mathscr X}_{\boldsymbol\xi}^RF\|_\infty,\qquad
 {\mathfrak d}_0(F)=\|F\|_\infty,                          \tag{3.5}
\]
and define
\[
 \boxed{
 \|F\|^{\rm HK}_{R,g}
 =
 \sum_{m\ge0}\frac{R^m}{m!}{\mathfrak d}_m(F).}            \tag{3.6}
\]
This is a global norm on the compact group; it does not choose a
logarithmic sheet.

\begin{lemma}[Markov contraction in the opposite-invariant norm]
\label{lem:n1v3-markov-contraction}
For every positive weakening matrix \(x\), \(t\ge0\), and \(R>0\),
\[
 \left\|e^{tg^2{\cal L}^{(n)}_x}F\right\|^{\rm HK}_{R,g}
 \le\|F\|^{\rm HK}_{R,g}.                                 \tag{3.7}
\]
\end{lemma}

\begin{proof}
The weakened semigroup is Markov by
Lemma~\ref{lem:n1v2-replica-positive}, hence is an
\(L^\infty\) contraction.  Commute every ordered right-invariant word
through the semigroup using Lemma~\ref{lem:n1v3-commute}, apply the
\(L^\infty\) contraction to each term of (3.5), and sum (3.6).
\end{proof}

\begin{lemma}[One radius loss for a complete ordered word]
\label{lem:n1v3-radius}
Let \(0<r<R\), \(\delta=R-r\), and let
\({\boldsymbol\eta}\) be any fixed ordered word of length \(d\).
Then
\[
 \|{\mathscr X}_{\boldsymbol\eta}^RF\|^{\rm HK}_{r,g}
 \le\frac{d!}{\delta^d}\|F\|^{\rm HK}_{R,g}.               \tag{3.8}
\]
\end{lemma}

\begin{proof}
For each \(m\), the words counted in
\({\mathfrak d}_m({\mathscr X}_{\boldsymbol\eta}^RF)\) form a subset
of all words of length \(m+d\).  Hence
\[
 \|{\mathscr X}_{\boldsymbol\eta}^RF\|^{\rm HK}_{r,g}
 \le
 \sum_{m\ge0}\frac{r^m}{m!}{\mathfrak d}_{m+d}(F).         \tag{3.9}
\]
For \(k=m+d\), the scalar binomial theorem gives
\[
 \binom{k}{d}r^{k-d}\delta^d\le R^k,
\]
or
\[
 \frac{r^m}{m!}
 \le\frac{d!}{\delta^d}\frac{R^{m+d}}{(m+d)!}.             \tag{3.10}
\]
Insert (3.10) into (3.9) and enlarge the remaining sum to all
\(k\ge0\).
\end{proof}

\section{Conversion of a compact Duhamel edge}
\label{sec:n1v3-edge}

For \(U\in SU(3)\), left and right frames are related by
\[
 X_a^L(U)=\sum_{b=1}^8
 [\operatorname{Ad}_U]_{ab}X_b^R(U).                      \tag{3.11}
\]
With the invariant Euclidean colour metric,
\(\operatorname{Ad}_U\) is orthogonal.  Therefore
\[
 a_G:=
 \sup_{U\in SU(3)}\max_a
 \sum_{b=1}^8|[\operatorname{Ad}_U]_{ab}|
 \le\sqrt8.                                               \tag{3.12}
\]
Let \(G_\mu\) be the weighted absolute covariance row (1.29), and
retain the finite endpoint-incidence constant \(h_*\).

\begin{lemma}[Uniform edge row after left--right conversion]
\label{lem:n1v3-edge-row}
After rewriting both fields in (3.3) by (3.11), the transformed
covariance has weighted absolute row at most
\[
                         a_G^2G_\mu\le8G_\mu.              \tag{3.13}
\]
Including owner endpoints and the radius loss at both ends gives the
edge constant
\[
 \boxed{
 K_{\rm HK}(\mu;R,r)
 =
 \frac{8h_*G_\mu}{(R-r)^2}.}                              \tag{3.14}
\]
\end{lemma}

\begin{proof}
The adjoint matrices act only on colour and do not change the block
distance.  For a fixed transformed row, take absolute values, sum the
right colour index of the second adjoint matrix, apply the original
row bound \(G_\mu\), and then sum the first adjoint row.  Each adjoint
matrix costs at most \(a_G\), proving (3.13).  The finite endpoint
dictionary costs \(h_*\).  Lemma~\ref{lem:n1v3-radius} contributes
one factor \((R-r)^{-1}\) at each endpoint.
\end{proof}

\section{The compact fixed-tree estimate}
\label{sec:n1v3-fixed-tree}

Use the anchored absolute tree norm of
Definition~\ref{def:n1v1-tree-norm}, with group positions in place
of real field centres and with (3.6) as the vertex norm.

\begin{theorem}[Time-ordered compact fixed-tree bound]
\label{thm:n1v3-fixed-tree}
Let \(T\) be a labelled tree on \(n\ge2\) local compact-group
vertices satisfying
\[
                         \|F_i\|^{\rm HK}_{R,g}\le z.       \tag{3.15}
\]
Then the tree contribution in the compact forest identity (2.24)
obeys
\[
 \boxed{
 \|{\cal A}^{\rm HK}_T(F_1,\ldots,F_n)\|_{\mu,r}
 \le
 z^nK_{\rm HK}^{\,n-1}
 \prod_{i=1}^nd_T(i)!.}                                   \tag{3.16}
\]
The constant is independent of the volume, \(n\), the forest
parameters, and \(g\).
\end{theorem}

\begin{proof}
Fix one of the \((n-1)!\) Duhamel orderings in (2.23).
Starting with the outermost cross generator, use (3.11) and take its
absolute covariance row.  The resulting two right derivatives commute
through every inner heat semigroup and every inner left cross
generator by Lemma~\ref{lem:n1v3-commute}.  Repeat for every edge.
No right derivative ever differentiates an adjoint coefficient:
at each stage the pointwise left--right relation is used only for the
current outer cross generator, while the produced right derivatives
are moved inward by the exact operator commutator (3.1).

At vertex \(i\), the procedure leaves one ordered right derivative
for every incident edge, hence a word of length \(d_T(i)\).
Lemma~\ref{lem:n1v3-radius} bounds it by
\[
                         z\,d_T(i)!\,(R-r)^{-d_T(i)}.       \tag{3.17}
\]
The identity \(\sum_i d_T(i)=2(n-1)\), the edge row
Lemma~\ref{lem:n1v3-edge-row}, and leafwise summation of the weighted
spatial rows give the right side of (3.16).

Every heat segment is an \(L^\infty\) contraction by
Lemma~\ref{lem:n1v3-markov-contraction}.  Finally, each Duhamel
ordering has the same bound, and the total volume of its
\((n-1)!\) ordered simplices is one by
Lemma~\ref{lem:n1v2-duhamel-edge}.  Hence no ordering factorial
remains.
\end{proof}

\begin{theorem}[Complete compact connected bound]
\label{thm:n1v3-connected}
If
\[
                              8zK_{\rm HK}<1,              \tag{3.18}
\]
then the connected compact heat-kernel expansion converges absolutely
and
\[
 \left\|
 \sum_{n\ge1}\frac1{n!}
 \kappa_{\nu_{g,C}}(F,\ldots,F)
 \right\|_{\mu,r}
 \le
 z+\frac{4z^2K_{\rm HK}}{1-8zK_{\rm HK}}.                 \tag{3.19}
\]
\end{theorem}

\begin{proof}
Insert (3.16) into (2.24) and use the exact degree-weighted tree sum
(1.33).  The numerical summation is exactly the proof of
Theorem~\ref{thm:n1v1-connected}, with \(K_\mu\) replaced by
\(K_{\rm HK}\).
\end{proof}

\begin{corollary}[Explicit four-source form]
\label{cor:n1v3-four-source}
If the compact vertex bank is holomorphic for
\(|s_j|<\rho_j\), \(0\le j\le3\), and (3.15) holds uniformly on that
polydisc, the coefficient carrying all four marked sources is bounded
by the right side of (3.19) times
\[
                              \prod_{j=0}^3\rho_j^{-1}.     \tag{3.20}
\]
\end{corollary}

\begin{proof}
Normal convergence in (3.19) permits four-variable Cauchy
differentiation after the complete connected sum.  Applying Cauchy
once in each source gives (3.20).
\end{proof}

\section{What remains of the nonlinear cluster gate}
\label{sec:n1v3-gate}

\begin{assessment}[The forest constant is now concrete]
\label{ass:n1v3-forest-closed}
For the compact heat-kernel reference, the abstract forest constant is
now
\[
                         K_{\rm HK}
 =\frac{8h_*G_\mu}{(R-r)^2}.                               \tag{3.21}
\]
There is no hidden \(g^{-1}\), time-ordering factorial, repeated
analytic-radius expenditure, or semigroup commutator constant.
Arbitrary star degrees are already included through (1.33).
\end{assessment}

The remaining activity is not the impossible tangent-space function
(2.3).  It is the normalized compact density ratio
\[
 {\cal W}_{\beta,W}(U)
 =
 \frac{e^{-\beta V_M(U,W)}}
 {Z_{\rm W}(\beta,W)\,k_{\beta^{-1/2},C,W}(U)},             \tag{3.22}
\]
where
\(Z_{\rm W}(\beta,W)=\int e^{-\beta V_M(U,W)}\,dm(U)\).
Since the heat kernel is already normalized,
\({\mathbb E}_{\nu_{\beta^{-1/2},C,W}}{\cal W}_{\beta,W}=1\).
Its constant mode is removed before forming local owners.  Both
factors in (3.22) live on the compact fibre, so there is no
logarithmic-sheet growth.  The next question is whether its
nonconstant local owners have
\[
 \|{\cal W}_{\beta,W}-1\|^{\rm HK}_{R,g}
 =O(\beta^{-1/5})                                         \tag{3.23}
\]
after the exact marginal quadratic and determinant rows are transported
separately.

\begin{firewall}[A compact ratio is finite but not automatically
small]
\label{fire:n1v3-ratio}
Compactness and strict heat-kernel positivity make (3.22) finite at
each fixed cutoff.  They do not prove the uniform smallness (3.23).
A mismatch between the Wilson large-deviation rate and the heat-kernel
distance can make the ratio exponentially large as
\(\beta\to\infty\).  V3 therefore closes the forest norm, not the
Wilson-to-reference comparison.
\end{firewall}

\begin{assessment}[Progress after compact V3]
\label{ass:n1v3-status}
V3 proves the volume-uniform analytic estimate requested by V2.
The use of opposite invariant fields makes every weakened heat
semigroup contractive in the same analytic norm, converts each
Duhamel edge with the explicit row \(8h_*G_\mu\), and proves the
complete four-source connected bound.  The live bottleneck has moved
from forest combinatorics to the compact Radon--Nikodym comparison
(3.23).
\end{assessment}

\begin{decision}[V4 target]
\label{dec:n1v3-next}
V4 should compare the exact Wilson conditional density with the
compact heat-kernel reference.  It must first compare their global
large-deviation rates.  If the physical covariance heat kernel is not
globally dominated by the Wilson action, introduce a covariance path
from a globally dominated compact heat kernel to the physical
covariance and transport its quadratic determinant exactly.  Only the
remaining normalized nonconstant owner may enter \(z\) in (3.18).
\end{decision}

\begin{firewall}[Claim boundary after compact V3]
\label{fire:n1v3-claim}
V3 proves the compact fixed-tree and connected analytic estimates.
It does not prove (3.23), close the numerical polymer gate, finish the
pushed Jacobian or marginal trajectory, construct the stable RG
orbit, remove the regulators, construct the continuum OS measure, or
prove the Yang--Mills mass gap.
\end{firewall}

\subsection{Parent record}

This V3 note was formed by appending the preceding material to the
validated compact V2, whose record was
\[
\begin{array}{c|c}
\text{lines}&1230\\
\text{bytes}&43268\\
\text{SHA--256}&
\texttt{3AC8088548584F53005284AD5DF8B27EC352ADEEAF8D1D97B62F103CC95A72BE}.
\end{array}
\]

% =============================================================================
% END APPENDED COMPACT-SERIES V3 MATERIAL
% =============================================================================

% =============================================================================
% BEGIN APPENDED COMPACT-SERIES V4 MATERIAL
% =============================================================================

\chapter{V4: global rate matching and exact covariance transport}
\label{ch:n1v4}

\section{The Wilson and heat-kernel rate functions}
\label{sec:n1v4-rates}

Fix a coarse-flat retained background, use the moving section as the
identity of the reduced compact fibre \(K\), and abbreviate the exact
Wilson action by \(V(U)\).  Thus
\[
                         V(U)\ge0,\qquad V(U)=0
                         \Longleftrightarrow U=I.           \tag{4.1}
\]
Let \(H=D^2V(I)>0\) on the reduced physical tangent space and set
\(C=H^{-1}\).  The left-invariant cometric \(C\) defines a
Riemannian distance \(d_C\) on \(K\).  Near the identity,
\[
\begin{split}
 V(\exp x)
 &=\frac12\langle x,Hx\rangle+O(\|x\|^3),\\
 d_C(I,\exp x)^2
 &=\langle x,Hx\rangle+O(\|x\|^3).                         \tag{4.2}
\end{split}
\]

\begin{lemma}[Finite global rate ratio]
\label{lem:n1v4-rate-ratio}
The number
\[
 \sigma_*:=
 \sup_{U\ne I}\frac{d_C(I,U)^2}{2V(U)}                     \tag{4.3}
\]
is finite and satisfies \(\sigma_*\ge1\).
\end{lemma}

\begin{proof}
Equation (4.2) shows that the quotient in (4.3) tends to one as
\(U\to I\), uniformly in tangent directions.  Remove a sufficiently
small identity neighbourhood.  On the remaining compact set,
\(V\) has a positive minimum by (4.1), while \(d_C\) is bounded by
the compact diameter.  Hence the quotient is bounded there.  Its
limit one at the omitted point also proves \(\sigma_*\ge1\).
\end{proof}

The preceding coercivity theorem gives a more directly auditable
upper bound.  Let
\[
 D(U)=\sum_e\|U_e-I\|_F^2,\qquad
 L_C:=\sup_{U\ne I}\frac{d_C(I,U)^2}{D(U)}.                \tag{4.4}
\]
The same compactness and quadratic-comparability proof shows
\(L_C<\infty\).

\begin{lemma}[Volume-uniform \(SU(3)\) chord--metric comparison]
\label{lem:n1v4-chord-metric}
Let \(\|H\|_{\rm op}\) be the operator norm of the reduced Hessian in
the product Hilbert--Schmidt metric.  Then
\[
                         L_C\le
                         \frac{16\pi^2}{27}\|H\|_{\rm op}. \tag{4.4a}
\]
The bound is independent of the number of reduced links.
\end{lemma}

\begin{proof}
For \(U_e\in SU(3)\), choose a traceless logarithm \(X_e\) whose
eigenangles lie in \([-4\pi/3,4\pi/3]\).  Such a choice follows by
taking principal eigenangles: if their sum is \(2\pi\) or \(-2\pi\),
subtract or add \(2\pi\) to an extreme eigenangle.  On
\(|\theta|\le4\pi/3\),
\[
 \theta^2
 \le\frac{16\pi^2}{27}|e^{i\theta}-1|^2.                  \tag{4.4b}
\]
Indeed
\(\theta^2/|e^{i\theta}-1|^2
=(\theta/2\sin(\theta/2))^2\) is increasing in
\(|\theta|\) on this interval and has the displayed endpoint value.
Summing eigenangles gives
\[
 \|X_e\|_F^2
 \le\frac{16\pi^2}{27}\|U_e-I\|_F^2.                      \tag{4.4c}
\]

Join \(I\) to \(U\) by the simultaneous path
\(U_e(q)=\exp(qX_e)\), \(0\le q\le1\).  Its left-trivialized velocity
is the constant vector \(X=(X_e)_e\).  Since the metric tensor is
\(C^{-1}=H\), its squared length is at most
\[
 \langle X,HX\rangle
 \le\|H\|_{\rm op}\sum_e\|X_e\|_F^2.
\]
The geodesic distance is no larger than the length of this path.
Use (4.4c), divide by \(D(U)\), and take the supremum.
\end{proof}

\begin{corollary}[Coercive upper bound]
\label{cor:n1v4-sigma-bound}
The reduced-fibre inequality \(V\ge D/120\) implies
\[
 \sigma_*\le60L_C
 \le\frac{320\pi^2}{9}\|H\|_{\rm op}.                      \tag{4.5}
\]
\end{corollary}

\begin{proof}
For \(U\ne I\),
\[
 \frac{d_C(I,U)^2}{2V(U)}
 \le60\frac{d_C(I,U)^2}{D(U)}.
\]
Take the supremum.
\end{proof}

\section{The exact obstruction detected by Varadhan scaling}
\label{sec:n1v4-varadhan}

For clarity, write \(k_{g,C}(I,U)\) for the heat-kernel density in
(2.15), and
\[
 Z_{\rm W}(g)=\int_K e^{-V(U)/g^2}\,dm(U).                 \tag{4.6}
\]

\begin{lemma}[Uniform compact small-time rates]
\label{lem:n1v4-small-time}
On the fixed compact fibre,
\[
 \lim_{g\downarrow0}
 \sup_{U\in K}
 \left|
 g^2\log k_{g,C}(I,U)+\frac12d_C(I,U)^2
 \right|=0,                                                \tag{4.7}
\]
and
\[
                              \lim_{g\downarrow0}
                              g^2\log Z_{\rm W}(g)=0.       \tag{4.8}
\]
\end{lemma}

\begin{proof}
For (4.7), cover \(K\) by finitely many normal-coordinate balls.
On balls not meeting the cut locus, the heat parametrix has the form
\[
 k_{g,C}(I,U)
 =
 g^{-\dim K}
 e^{-d_C(I,U)^2/(2g^2)}
 \{a_0(I,U)+O(g^2)\},                                     \tag{4.9}
\]
with \(a_0>0\) on a smaller ball.  Across the cut locus, subdivide a
minimizing geodesic into finitely many such balls and use the
semigroup property.  Upper bounds follow by the same subdivision and
compact Gaussian bounds.  The finite cover makes the error in
\(g^2\log k+d_C^2/2\) uniform; powers of \(g\) disappear after
multiplication by \(g^2\).  This proves (4.7).

For (4.8), \(Z_{\rm W}(g)\le1\).  Conversely, (4.2) gives
\(V(\exp x)\le c\|x\|^2\) in a fixed identity chart.  Integrating over
\(\|x\|\le g\) yields
\[
                         Z_{\rm W}(g)\ge c_0g^{\dim K}
                         e^{-c_1}.
\]
Taking \(g^2\log\) proves (4.8).
\end{proof}

\begin{theorem}[Necessary global heat-reference condition]
\label{thm:n1v4-rate-obstruction}
For the normalized density ratio
\[
 {\cal W}_{g,C}(U)
 =
 \frac{e^{-V(U)/g^2}}
 {Z_{\rm W}(g)\,k_{g,C}(I,U)},                             \tag{4.10}
\]
one has
\[
 \lim_{g\downarrow0}g^2\log{\cal W}_{g,C}(U)
 =
 \frac12d_C(I,U)^2-V(U)                                   \tag{4.11}
\]
uniformly on \(K\).  If
\[
 \Delta_C:=
 \sup_U\left\{\frac12d_C(I,U)^2-V(U)\right\}>0,            \tag{4.12}
\]
then
\[
 \|{\cal W}_{g,C}-1\|_\infty
 \ge
 \exp\!\left\{\frac{\Delta_C+o(1)}{g^2}\right\}-1.          \tag{4.13}
\]
In that case the global compact activity (3.23) is impossible even
in zeroth order.
\end{theorem}

\begin{proof}
Take logarithms in (4.10), multiply by \(g^2\), and apply
Lemma~\ref{lem:n1v4-small-time}.  This proves the uniform limit
(4.11).  If \(\Delta_C>0\), choose a maximizer in the compact set and
evaluate (4.11) there to obtain (4.13).
\end{proof}

\begin{firewall}[Local Hessian matching does not imply global rate
matching]
\label{fire:n1v4-local-global}
Equation (4.2) makes the two rate functions agree to quadratic order
at the vacuum.  It gives no sign for (4.12) away from the vacuum.
Therefore the physical covariance \(C=H^{-1}\) cannot be used in a
global supremum activity until \(\Delta_C\le0\) has been proved.
\end{firewall}

\section{A globally dominated broad compact reference}
\label{sec:n1v4-broad}

For \(s>0\), replace \(C\) by \(sC\).  Scaling the cometric scales
the distance by
\[
                         d_{sC}(I,U)^2=s^{-1}d_C(I,U)^2.   \tag{4.14}
\]

\begin{theorem}[Existence of a globally rate-dominated reference]
\label{thm:n1v4-broad}
For every \(s>\sigma_*\),
\[
 V(U)-\frac12d_{sC}(I,U)^2
 \ge
 \left(1-\frac{\sigma_*}{s}\right)V(U)\ge0.                \tag{4.15}
\]
In particular, the Wilson-to-\(sC\) heat-kernel ratio has no positive
exponential large-deviation defect.  It is enough to choose the
auditable value
\[
 s>60L_C,
 \qquad\hbox{or more explicitly}\qquad
 s>\frac{320\pi^2}{9}\|H\|_{\rm op}.                       \tag{4.16}
\]
\end{theorem}

\begin{proof}
Definition (4.3) gives
\[
 \frac12d_C(I,U)^2\le\sigma_*V(U).
\]
Divide by \(s\) and use (4.14).  Equation (4.16) follows from (4.5).
\end{proof}

\begin{remark}[Why the broad ratio is not declared small]
\label{rem:n1v4-broad-not-small}
For \(s>1\), the broad heat kernel has tangent covariance \(sC\),
whereas the Wilson saddle has covariance \(C\).  Their likelihood
ratio therefore contains an order-one quadratic form on typical
rescaled fields.  Theorem~\ref{thm:n1v4-broad} removes exponential
large-field growth; it does not turn the covariance mismatch into an
irrelevant vertex.
\end{remark}

\section{Exact transport from the broad to the physical covariance}
\label{sec:n1v4-path}

Put
\[
 s_\tau=s-(s-1)\tau,\qquad C_\tau=s_\tau C,
 \qquad0\le\tau\le1.                                      \tag{4.17}
\]
Thus \(C_0=sC\) and \(C_1=C\).

\begin{theorem}[Scalar compact covariance path]
\label{thm:n1v4-covariance-path}
For every \(F\in C^\infty(K)\),
\[
 \frac{d}{d\tau}{\mathbb E}_{\nu_{g,C_\tau}}F
 =
 -(s-1)g^2
 {\mathbb E}_{\nu_{g,C_\tau}}{\cal L}_CF.                 \tag{4.18}
\]
Consequently
\[
 {\mathbb E}_{\nu_{g,C}}F-
 {\mathbb E}_{\nu_{g,sC}}F
 =
 -(s-1)g^2\int_0^1
 {\mathbb E}_{\nu_{g,C_\tau}}{\cal L}_CF\,d\tau.           \tag{4.19}
\]
This transports the entire quadratic mismatch as an exact two-leg
semigroup insertion; it is not entered into the polymer activity.
\end{theorem}

\begin{proof}
Because \({\cal L}_{C_\tau}=s_\tau{\cal L}_C\), all generators on the
path commute.  Differentiate
\[
 {\mathbb E}_{\nu_{g,C_\tau}}F
 =
 [e^{g^2s_\tau{\cal L}_C}F](I)
\]
and use \(s_\tau'=-(s-1)\).  Integration from zero to one gives
(4.19).
\end{proof}

\begin{corollary}[Uniform forest stock on the covariance path]
\label{cor:n1v4-path-stock}
If \(G_\mu(C)\) is the physical weighted covariance row, then
\[
 G_\mu(C_\tau)=s_\tau G_\mu(C)\le sG_\mu(C),               \tag{4.20}
\]
and every compact forest estimate of V3 remains valid along the path
with
\[
 K_{\rm HK}^{\rm path}
 =
 \frac{8h_*sG_\mu(C)}{(R-r)^2}.                            \tag{4.21}
\]
\end{corollary}

\begin{proof}
The weighted absolute row is homogeneous under positive scalar
multiplication.  Insert (4.20) into (3.14).
\end{proof}

\begin{record}[Owned Gaussian determinant]
\label{rec:n1v4-determinant}
In tangent variables of dimension \(D\), the covariance change
\(sC\to C\) carries the constant Gaussian normalization
\[
                              \frac{D}{2}\log s.            \tag{4.22}
\]
It is extensive when summed over blocks but cancels from normalized
conditional expectations.  Its retained-background derivatives,
when \(C=C(W)\), belong to the already separated log-determinant
response row (1.3)--(1.4), not to the small connected activity.
\end{record}

\section{The corrected reference factorization}
\label{sec:n1v4-factorization}

At the level of exact compact densities,
\[
 \frac{d\mu_{\rm W}}{d\nu_{g,sC}}
 =
 \frac{d\mu_{\rm W}}{d\nu_{g,C}}
 \frac{d\nu_{g,C}}{d\nu_{g,sC}}.                           \tag{4.23}
\]
The second factor is transported by
Theorem~\ref{thm:n1v4-covariance-path}.  The first factor is
quadratically matched at the vacuum and is the only candidate for a
small good-chart activity.  It must not be placed in a global
supremum norm when \(\Delta_C>0\); its complement is controlled by
the original Wilson coercivity and the far-bad primitive (1.7).

\begin{assessment}[Progress after compact V4]
\label{ass:n1v4-status}
V4 identifies the precise global obstruction to the physical
heat-kernel ratio: the sign of the rate defect (4.12).  It proves that
a broad covariance \(sC\), with the explicit sufficient choice
\(s>60L_C\), is globally rate dominated by the Wilson action.  It
then transports the order-one covariance mismatch exactly along
(4.17), with the uniform forest constant (4.21) and a separately
owned determinant.  Thus neither an exponentially growing density
ratio nor a marginal quadratic form is hidden in the small activity.
\end{assessment}

\begin{decision}[Mandatory V5 audit and next target]
\label{dec:n1v4-next}
V5 is the compulsory companion audit.  It should inspect the newest
active Standard-Model--Einstein and Navier--Stokes notes for a
large-deviation matched reference, covariance-path determinant, or
chart-local analytic extension theorem.  After the audit, the
Yang--Mills task is to split the physical ratio in (4.23) into:
\begin{enumerate}[label=\textup{(\roman*)},leftmargin=3.5em]
\item a quadratically matched principal-chart owner with a proved
\(O(\beta^{-1/5})\) compact analytic norm;
\item a nonprincipal-chart defect bounded by the Wilson far-bad
activity; and
\item the exact covariance path and determinant already isolated in
V4.
\end{enumerate}
\end{decision}

\begin{firewall}[Claim boundary after compact V4]
\label{fire:n1v4-claim}
V4 proves the rate criterion, existence of a globally dominated broad
reference, and exact covariance transport.  It does not yet construct
the exact chart-local/far factorization of the physical density ratio
or close the numerical polymer gate.  It also does not finish the
pushed Jacobian, marginal trajectory, stable RG contraction,
regulator removal, OS reconstruction, or Yang--Mills mass gap.
\end{firewall}

\subsection{Parent record}

This V4 note was formed by appending the preceding material to the
validated compact V3, whose record was
\[
\begin{array}{c|c}
\text{lines}&1609\\
\text{bytes}&56168\\
\text{SHA--256}&
\texttt{F3CBA422F9088589A5E306F24CF0779BD94505D4E377BAE7A322E9DCF2F60CDC}.
\end{array}
\]

% =============================================================================
% END APPENDED COMPACT-SERIES V4 MATERIAL
% =============================================================================

% =============================================================================
% BEGIN APPENDED COMPACT-SERIES V5 MATERIAL
% =============================================================================

\chapter{V5: companion audit and an exact analytic Morse chart}
\label{ch:n1v5}

\section{Mandatory five-version companion audit}
\label{sec:n1v5-audit}

The active companion files inspected before choosing the V5 argument
were
\[
\begin{array}{c|r|r|l}
\text{source}&\text{lines}&\text{bytes}&\text{SHA--256}\\ \hline
\text{SM--Einstein V73}&25889&955007&
\texttt{47D994C004A433D95698BE67D91BE8FD291DFB53438E1F9EE66FCBBED5F545B4}\\
\text{Navier--Stokes V26}&5319&278283&
\texttt{82C887F8C2C69E4F28FAD7C3DC8DE5B9099CB5A4BF431EBA1FFF3E91935978AD}.
\end{array}
\]

\begin{audit}[Transfer decision]
\label{aud:n1v5-transfer}
The SM V70--V71 degree-weighted forest and exact rebase results were
transferred from the compact YM V1 and add no new YM theorem.  SM V73
does add a useful general mechanism: a scalar coefficient majorant
turns an exact residual equation into a uniform analytic chart, with
all star derivatives controlled before the forest sum.  That analytic
implicit-function mechanism transfers to the reduced Wilson stationary
equation because its eliminated Hessian is invertible on the already
declared physical complement.

The Navier--Stokes file is unchanged.  Its rank-one Oseen estimates do
not control the full reduced colour-link Hessian and no constant is
imported from it.
\end{audit}

\begin{correction}[Canonical archive suffix]
\label{cor:n1v5-archive-name}
The validated 982002-byte predecessor recorded in compact V1 was
subsequently canonicalized in the shared folder as
\[
\texttt{Renewal\_Yang\_Mills\_Mass\_Gap\_Research\_Notes\_V65\_f8.tex},
\]
with unchanged SHA--256
\[
\texttt{7A5AC49398BF5FE00E454FA469459BC6216E6E1109C79FC376A15576C72C858C}.
\]
Earlier mentions of the temporary suffix
\(\texttt{\_V65\_f1.tex}\) refer to this same content.
\end{correction}

\begin{decision}[Audit-adjusted direction]
\label{dec:n1v5-audit-direction}
Rather than estimate every order of the Wilson remainder separately,
V5 applies the V73 majorant architecture upstream to an exact
parameter-dependent analytic Morse coordinate.  On its principal
chart the Wilson action becomes exactly quadratic.  The only
nonconstant good-field interaction is then the Haar/Morse Jacobian.
\end{decision}

\section{The normalized stationary family}
\label{sec:n1v5-stationary}

Let \(w\) denote retained variables in a fixed complex polydisc.
Assume the previously constructed stationary section has been used to
translate the reduced fine coordinate so that
\[
 D_xV(0,w)=0,\qquad H(w):=D_x^2V(0,w)>0                  \tag{5.1}
\]
on the real slice.  All gauge, toron, and cohomological zero modes
remain in \(w\); \(x\) contains only the finite-dimensional physical
complement.  Shrink the retained polydisc, if necessary, so that
\(H(w)\) and its analytic square roots are invertible throughout it.

Put
\[
 u=H(w)^{1/2}x,\qquad
 \Phi(u,w)
 =
 V(H(w)^{-1/2}u,w)-V(0,w).                                \tag{5.2}
\]
Then
\[
 \Phi(0,w)=0,\quad D_u\Phi(0,w)=0,\quad
 D_u^2\Phi(0,w)=I.                                        \tag{5.3}
\]

\begin{lemma}[Integral Hessian factorization]
\label{lem:n1v5-integral-hessian}
Define
\[
 A(u,w)
 =
 2\int_0^1(1-q)D_u^2\Phi(qu,w)\,dq.                       \tag{5.4}
\]
Then \(A\) is analytic and symmetric,
\[
                         A(0,w)=I,                         \tag{5.5}
\]
and
\[
                         \Phi(u,w)
 =\frac12\langle u,A(u,w)u\rangle.                        \tag{5.6}
\]
\end{lemma}

\begin{proof}
Analyticity permits differentiation under the finite integral, and
symmetry follows from the Hessian.  Equation (5.5) follows from
(5.3) and \(2\int_0^1(1-q)dq=1\).  Taylor's formula with integral
remainder, using the first two identities in (5.3), gives
\[
 \Phi(u,w)
 =
 \int_0^1(1-q)D_u^2\Phi(qu,w)[u,u]\,dq,
\]
which is (5.6).
\end{proof}

\section{A quantitative analytic square-root majorant}
\label{sec:n1v5-majorant}

Use the weighted support analytic algebra of the preceding programme.
Suppose that on \(\|u\|\le t\),
\[
 \|A-I\|_t\le{\mathfrak m}(t),\qquad
 \|D_uA\|_t\le{\mathfrak m}'(t),                           \tag{5.7}
\]
where \({\mathfrak m}\) is a convergent nonnegative scalar series with
\({\mathfrak m}(0)=0\).  For \(0\le q<1\), set
\[
\begin{split}
 a_{\sqrt{\ }}(q)
 &:=
 \sum_{n\ge1}\left|\binom{1/2}{n}\right|q^n,\\
 b_{\sqrt{\ }}(q)
 &:=
 \sum_{n\ge1}n\left|\binom{1/2}{n}\right|q^{n-1}.
                                                               \tag{5.8}
\end{split}
\]
Both are finite for \(q<1\).

\begin{lemma}[Square-root stocks]
\label{lem:n1v5-square-root}
If \(q_t={\mathfrak m}(t)<1\), the principal binomial series
\[
 B(u,w):=A(u,w)^{1/2}
 =
 I+\sum_{n\ge1}\binom{1/2}{n}(A-I)^n                     \tag{5.9}
\]
converges normally in the analytic algebra and satisfies
\[
\begin{split}
 \|B-I\|_t
 &\le a_{\sqrt{\ }}(q_t),\\
 \|D_uB\|_t
 &\le b_{\sqrt{\ }}(q_t){\mathfrak m}'(t).                \tag{5.10}
\end{split}
\]
Moreover \(B\) is symmetric and \(B^2=A\).
\end{lemma}

\begin{proof}
Submultiplicativity and (5.7) majorize (5.9) term by term by the
first series in (5.8).  Differentiating a word \((A-I)^n\) produces
\(n\) terms, each bounded by
\(q_t^{n-1}{\mathfrak m}'(t)\), which gives the second bound.
Normal convergence permits multiplication and differentiation.
Because every power of the symmetric matrix \(A-I\) is symmetric,
\(B\) is symmetric.  The scalar binomial identity, valid under normal
convergence, gives \(B^2=A\).
\end{proof}

\section{The exact analytic Morse coordinate}
\label{sec:n1v5-morse}

Define
\[
                         y={\cal Y}(u,w):=B(u,w)u.          \tag{5.11}
\]
Let
\[
 \chi(t)
 :=
 a_{\sqrt{\ }}({\mathfrak m}(t))
 +
 t\,b_{\sqrt{\ }}({\mathfrak m}(t)){\mathfrak m}'(t).
                                                               \tag{5.12}
\]

\begin{theorem}[Quantitative parameter-dependent Morse chart]
\label{thm:n1v5-morse}
Suppose
\[
                         {\mathfrak m}(t)<1,\qquad
                         \chi(t)<1.                        \tag{5.13}
\]
Then \({\cal Y}(\,\cdot\,,w)\) is injective on the \(u\)-ball of
radius \(t\), its image contains the \(y\)-ball of radius
\[
                         t_y=(1-\chi(t))t,                 \tag{5.14}
\]
and it has an analytic inverse \(u={\cal U}(y,w)\) there.  On this
chart the action has the exact normal form
\[
 \boxed{
 V(H^{-1/2}{\cal U}(y,w),w)
 =
 V(0,w)+\frac12\langle y,y\rangle.}                        \tag{5.15}
\]
\end{theorem}

\begin{proof}
Differentiate (5.11):
\[
 D_u{\cal Y}
 =
 B+(D_uB)[\,\cdot\,]u.
\]
Equations (5.10)--(5.12) give
\[
                         \|D_u{\cal Y}-I\|_t\le\chi(t)<1.  \tag{5.16}
\]
Writing \({\cal Y}(u)=u+N(u)\), (5.16) makes \(N\) Lipschitz with
constant \(\chi(t)\).  For \(\|y\|\le(1-\chi(t))t\), the map
\[
                         u\longmapsto y-N(u)
\]
is a contraction of the closed \(t\)-ball into itself.  Banach's
theorem gives a unique inverse.  Uniform convergence of its analytic
Picard iterates on smaller polydiscs proves analytic dependence on
\(y,w\).  Injectivity also follows from
\[
 \|{\cal Y}(u_1)-{\cal Y}(u_2)\|
 \ge(1-\chi(t))\|u_1-u_2\|.
\]

Finally, Lemma~\ref{lem:n1v5-integral-hessian} and
\(B^TB=B^2=A\) give
\[
 \Phi(u,w)=\frac12\langle Bu,Bu\rangle
 =\frac12\langle y,y\rangle,
\]
which is (5.15).
\end{proof}

\begin{remark}[Relation to the SM V73 transfer]
\label{rem:n1v5-v73}
The contraction and coefficient-majorant logic is the transferable
part of SM V73.  The explicit integral Hessian (5.4) and square-root
map (5.11) are Yang--Mills-specific improvements: they solve the
stationary residual and quadratic normalization simultaneously.
\end{remark}

\section{The only principal-chart interaction is a Jacobian}
\label{sec:n1v5-jacobian}

Write the reduced Haar density in the \(x\)-chart as
\[
                         dm_{\rm red}=J_{\rm H}(x,w)\,dx.
\]
Normalize its nonconstant logarithm after the linear change (5.2) by
\[
 {\mathfrak h}(u,w)
 :=
 \log\frac{J_{\rm H}(H^{-1/2}u,w)}{J_{\rm H}(0,w)},         \tag{5.17}
\]
and suppose its coefficient majorant on the \(t\)-ball is
\[
                         \|{\mathfrak h}\|_t
                         \le{\mathfrak j}(t),\qquad
                         {\mathfrak j}(0)=0.               \tag{5.18}
\]
Let \(d\) be the fixed local reduced dimension and define
\[
 \epsilon_{\rm M}(t)
 :=
 {\mathfrak j}(t)
 +d\{-\log(1-\chi(t))\}.                                  \tag{5.19}
\]

\begin{theorem}[Exact Gaussianization and Jacobian activity]
\label{thm:n1v5-gaussianization}
On the Morse chart, set \(y=gb\), \(g=\beta^{-1/2}\).  Then
\[
\begin{split}
 e^{-\beta V(x,w)}dm_{\rm red}(x)
 &=
 {\cal N}_{\beta}(w)
 e^{-\langle b,b\rangle/2}
 e^{\ell_\beta(b,w)}\,db,                                 \tag{5.20}\\
 {\cal N}_{\beta}(w)
 &:=
 e^{-\beta V(0,w)}g^d
 (\det H(w))^{-1/2}J_{\rm H}(0,w),                         \tag{5.21}
\end{split}
\]
where \(\ell_\beta(0,w)=0\).  On
\(\|gb\|\le t_y\), its coefficient-majorant norm satisfies
\[
                         \|\ell_\beta\|_t
                         \le\epsilon_{\rm M}(t),            \tag{5.22}
\]
and hence
\[
                         \|e^{\ell_\beta}-1\|_t
 \le e^{\epsilon_{\rm M}(t)}-1.                            \tag{5.23}
\]
\end{theorem}

\begin{proof}
The successive changes \(x=H^{-1/2}u\), \(y={\cal Y}(u,w)\),
and \(y=gb\), together with (5.15), give the Gaussian exponential
and the constant factor (5.21).  The nonconstant logarithm is
\[
 \ell_\beta(b,w)
 =
 {\mathfrak h}({\cal U}(gb,w),w)
 -
 \Tr\log D_u{\cal Y}({\cal U}(gb,w),w).                    \tag{5.24}
\]
The constant at \(b=0\) vanishes because
\({\cal U}(0,w)=0\) and \(D_u{\cal Y}(0,w)=I\).

Equation (5.16), normal convergence of
\(\Tr\log(I+E)=\sum_{n\ge1}(-1)^{n+1}\Tr(E^n)/n\), and the local trace
bound give
\[
 \|\Tr\log D_u{\cal Y}\|_t
 \le d\sum_{n\ge1}\frac{\chi(t)^n}{n}
 =d\{-\log(1-\chi(t))\}.                                  \tag{5.25}
\]
Composition with the inverse stays in the \(t\)-ball by
Theorem~\ref{thm:n1v5-morse}.  Equations (5.18), (5.24), and (5.25)
prove (5.22).  The Banach-algebra exponential estimate gives (5.23).
\end{proof}

\begin{corollary}[Running-window smallness]
\label{cor:n1v5-running}
Assume the majorants are uniform in the retained source polydisc and
\[
 {\mathfrak m}(t)\le M_1t,\qquad
 {\mathfrak m}'(t)\le M_1,\qquad
 {\mathfrak j}(t)\le J_1t                                  \tag{5.26}
\]
for \(0\le t\le t_0\).  With \(r_\beta=\beta^{-2/5}\), there are
fixed \(c_0,C_{\rm M},\beta_0>0\) such that on the physical
Morse ball
\[
                         \|y\|\le c_0r_\beta               \tag{5.27}
\]
and for \(\beta\ge\beta_0\),
\[
 \|e^{\ell_\beta}-1\|
 \le e^{C_{\rm M}\beta^{-2/5}}-1
 =O(\beta^{-2/5}).                                        \tag{5.28}
\]
Four retained derivatives are obtained by Cauchy with their four
source radii displayed exactly as in (3.20).
\end{corollary}

\begin{proof}
The power series in (5.8) vanish linearly at zero.  Equations
(5.12) and (5.26) therefore give
\(\chi(t)\le C_1t\) for sufficiently small \(t\).  Choose
\(c_0\) and \(\beta_0\) so that (5.13)--(5.14) hold with
\(t=C_2r_\beta\) and the ball (5.27) lies in the image.
Equations (5.19) and (5.26) then give
\(\epsilon_{\rm M}(t)\le C_{\rm M}r_\beta\).
Apply (5.23).  Source differentiation follows from normal convergence
and multivariable Cauchy.
\end{proof}

\begin{record}[Ownership after exact Gaussianization]
\label{rec:n1v5-ownership}
The constant factor (5.21) has three separate owners:
\[
 -\beta V(0,w),\qquad
 -\frac12\log\det H(w),\qquad
 \log J_{\rm H}(0,w).                                     \tag{5.29}
\]
They are respectively the retained classical action, physical
Gaussian determinant, and constant Haar/coarea density.  None is
charged again in the small interaction (5.23).  Its only owner is the
nonconstant Haar/Morse Jacobian (5.24).
\end{record}

\begin{assessment}[Progress after compact V5]
\label{ass:n1v5-status}
The compulsory audit imports the SM V73 coefficient-majorant
architecture and leads to a stronger YM construction.  V5 proves a
quantitative parameter-dependent analytic Morse theorem, exactly
Gaussianizes the Wilson action on the principal stationary chart, and
bounds the sole nonconstant Jacobian interaction by
\(O(\beta^{-2/5})\) on the running window.  This removes the
term-by-term Wilson remainder from the good-field forest.
\end{assessment}

\begin{decision}[V6 target]
\label{dec:n1v5-next}
V6 should glue the exact Morse chart to the far-bad sector without
differentiating a compact cutoff.  The required identity must compare
the truncated Morse Gaussian integral with a full Gaussian forest and
assign the difference to an explicitly bounded exit polymer.  It must
keep three distinct objects: the soft-transition activity (1.23), the
geometric far-bad activity (1.7), and the Gaussian/Morse chart-exit
tail.  Only after proving that gluing identity may (5.28) be inserted
as \(z\) in the connected bound (3.19).
\end{decision}

\begin{firewall}[Claim boundary after compact V5]
\label{fire:n1v5-claim}
V5 proves the local exact Morse normal form and its analytic Jacobian
bound.  It does not yet prove the cutoff-free chart/far gluing, close
the numerical polymer gate, finish the pushed Jacobian and marginal
trajectory, construct the stable RG orbit, remove regulators, perform
OS reconstruction, or prove the Yang--Mills mass gap.
\end{firewall}

\subsection{Parent record}

This V5 note was formed by appending the preceding material to the
validated compact V4, whose record was
\[
\begin{array}{c|c}
\text{lines}&2020\\
\text{bytes}&69371\\
\text{SHA--256}&
\texttt{B581452CFFBA4B2D8305C7B6E81E3AF5BFC59B3C0AF94738DF894139E94809A2}.
\end{array}
\]

% =============================================================================
% END APPENDED COMPACT-SERIES V5 MATERIAL
% =============================================================================

% =============================================================================
% BEGIN APPENDED COMPACT-SERIES V6 MATERIAL
% =============================================================================

\chapter{V6: exact Morse--Gaussian gluing without cutoff derivatives}
\label{ch:n1v6}

\section{Local scope of the Morse construction}
\label{sec:n1v6-scope}

\begin{correction}[No global square-root locality claim]
\label{cor:n1v6-local-morse}
The V5 Morse theorem is a theorem on one fixed-dimensional local
conditional fibre, or on a fixed filling neighbourhood whose
dimension and Hessian bounds are uniform.  It does not prove that the
square root of the complete lattice Hessian is exponentially local.
The archived square-root locality warning remains in force.  Global
assembly must use local conditional charts, energy-form coordinates,
or a separately proved functional-calculus locality theorem.
\end{correction}

This correction does not change (5.15): on every admitted local fibre
the action is exactly quadratic.  It changes only how those charts may
be assembled across the lattice.

\section{Local source-generating integral}
\label{sec:n1v6-local-integral}

Fix one admitted local fibre of real dimension \(d\), one retained
background \(w\), and a source vector
\(\boldsymbol s=(s_0,\ldots,s_3)\) in
\[
                         |s_j|<\rho_j.                     \tag{6.1}
\]
Let \(\Omega_0\) be the fixed Morse neighbourhood from V5 and let
\[
 {\cal J}(y,\boldsymbol s,w)
\]
denote the complete nonconstant local density in Morse coordinates:
the Haar/Morse Jacobian, the four bounded source insertions, and any
already-owned local chart score.  Normalize
\[
                         {\cal J}(0,\boldsymbol0,w)=1.      \tag{6.2}
\]
Assume it is analytic for \(\|y\|<t_0\) and on the source polydisc,
with coefficient-majorant bound
\[
 \|{\cal J}\|_{t_0,\boldsymbol\rho}\le M_{\cal J}.          \tag{6.3}
\]

Write
\[
 \ell(y,\boldsymbol s,w)=\log{\cal J}(y,\boldsymbol s,w)
 \]
in the fixed determinant-line branch and split its field Taylor
series exactly as
\[
\begin{split}
 \ell(y,\boldsymbol s,w)
 &=
 c(\boldsymbol s,w)
 +\langle a(\boldsymbol s,w),y\rangle
 +\frac12\langle y,B(\boldsymbol s,w)y\rangle\\
 &\quad+{\cal R}_3(y,\boldsymbol s,w).                     \tag{6.4}
\end{split}
\]
At zero source, \(c(\boldsymbol0,w)=0\).  The constant, linear, and
quadratic terms in (6.4) are not treated perturbatively below.

Choose
\[
 r_\beta=\beta^{-2/5},\qquad g=\beta^{-1/2},\qquad
 R_\beta=\frac{c_0r_\beta}{g}=c_0\beta^{1/10},             \tag{6.5}
\]
where \(c_0>0\) is fixed so that the physical ball
\(\|y\|\le c_0r_\beta\) lies in \(\Omega_0\) for large \(\beta\).

\section{Exact absorption of the two-jet}
\label{sec:n1v6-two-jet}

Put
\[
 A_g(\boldsymbol s,w)=I-g^2B(\boldsymbol s,w),\qquad
 m_g=gA_g^{-1}a.                                          \tag{6.6}
\]

\begin{lemma}[Shifted Gaussian two-jet]
\label{lem:n1v6-two-jet}
Suppose
\[
 \sup_{\boldsymbol s,w}g^2\|\operatorname{Re}B\|_{\rm op}
 \le\frac12.                                               \tag{6.7}
\]
Then \(\operatorname{Re}A_g\ge I/2\), and
\[
\begin{split}
 Z_2(\boldsymbol s,w)
 &:=
 {\mathbb E}_{\gamma_d}
 \exp\!\left\{
 c+g\langle a,b\rangle
 +\frac{g^2}{2}\langle b,Bb\rangle
 \right\}\\
 &=
 e^c\det(A_g)^{-1/2}
 \exp\!\left\{
 \frac{g^2}{2}\langle a,A_g^{-1}a\rangle
 \right\}.                                                \tag{6.8}
\end{split}
\]
Here \(\gamma_d\) is standard centred Gaussian probability and the
square-root branch is the one continued from \(g=0\).
\end{lemma}

\begin{proof}
The real-part bound in (6.7) gives
\(\operatorname{Re}\langle b,A_gb\rangle\ge\|b\|^2/2\), so
the integral is absolutely convergent throughout the source
polydisc.  Complete the square,
\[
 -\frac12\langle b,A_gb\rangle+g\langle a,b\rangle
 =
 -\frac12\langle b-m_g,A_g(b-m_g)\rangle
 +\frac{g^2}{2}\langle a,A_g^{-1}a\rangle,
\]
and evaluate the finite-dimensional Gaussian determinant.
\end{proof}

\begin{record}[Two-jet ownership]
\label{rec:n1v6-two-jet-owners}
The three terms
\[
 c,\qquad-\frac12\log\det A_g,\qquad
 \frac{g^2}{2}\langle a,A_g^{-1}a\rangle                  \tag{6.9}
\]
are respectively the constant source owner, the local covariance
determinant, and the shifted-current contraction.  They are transported
exactly and are not also charged as nonlinear activities.
\end{record}

\section{The cutoff-free add--subtract identity}
\label{sec:n1v6-gluing}

Let \({\cal N}_\beta(w)\) be the constant in (5.21), and let
\({\cal Z}_{\rm ext}(\boldsymbol s,w)\) be the exact original compact
integral over the complement of
\[
 \Omega_\beta:=\{U\in\Omega_0:\|y(U,w)\|\le c_0r_\beta\}.
                                                               \tag{6.10}
\]
Inside \(\Omega_\beta\), (5.20) and (6.4) are exact.
Define the quadratic exponential
\[
 {\cal Q}_2(b,\boldsymbol s,w)
 :=
 \exp\!\left\{
 c+g\langle a,b\rangle
 +\frac{g^2}{2}\langle b,Bb\rangle
 \right\}.                                                \tag{6.11}
\]

\begin{theorem}[Exact Morse--Gaussian gluing]
\label{thm:n1v6-gluing}
The complete local source-generating integral satisfies
\[
\begin{split}
 \frac{{\cal Z}_\beta(\boldsymbol s,w)}
 {{\cal N}_\beta(w)(2\pi)^{d/2}}
 &=
 \underbrace{{\mathbb E}_{\gamma_d}{\cal Q}_2}_{\rm two\!-\!jet}
 \\
 &\quad+
 \underbrace{{\mathbb E}_{\gamma_d}
 \left[
 {\bf1}_{\{\|b\|\le R_\beta\}}
 {\cal Q}_2
 \left(e^{{\cal R}_3(gb,\boldsymbol s,w)}-1\right)
 \right]}_{\rm interior\ Taylor\ defect}
 \\
 &\quad-
 \underbrace{{\mathbb E}_{\gamma_d}
 \left[
 {\bf1}_{\{\|b\|>R_\beta\}}{\cal Q}_2
 \right]}_{\rm Gaussian\ chart\ exit}
 \\
 &\quad+
 \underbrace{
 \frac{{\cal Z}_{\rm ext}(\boldsymbol s,w)}
 {{\cal N}_\beta(w)(2\pi)^{d/2}}
 }_{\rm Wilson\ exterior}.
                                                               \tag{6.12}
\end{split}
\]
No characteristic function occurs in the all-space two-jet term.
Every characteristic function in (6.12) belongs to a defect estimated
in total variation, not differentiated by a forest.
\end{theorem}

\begin{proof}
Split the original compact integral into \(\Omega_\beta\) and its
complement.  In \(\Omega_\beta\), substitute (5.20), (6.4), and
\(y=gb\); this gives
\[
 {\mathbb E}_{\gamma_d}
 \left[
 {\bf1}_{\{\|b\|\le R_\beta\}}
 {\cal Q}_2e^{{\cal R}_3(gb)}
 \right]
\]
after division by the constant normalization.  Add and subtract
\({\mathbb E}_{\gamma_d}{\cal Q}_2\), and split that full Gaussian
integral at \(R_\beta\).  The resulting four terms are exactly
(6.12).
\end{proof}

\section{Quantitative bounds for all three defects}
\label{sec:n1v6-defects}

Assume the logarithm has a coefficient majorant
\[
                         \|\ell\|_{t_0,\boldsymbol\rho}
                         \le M_\ell.                       \tag{6.13}
\]

\begin{lemma}[Interior third-order defect]
\label{lem:n1v6-interior}
On \(\|y\|\le c_0r_\beta<t_0\),
\[
 \|{\cal R}_3(y,\boldsymbol s,w)\|
 \le
 M_\ell
 \left(\frac{c_0r_\beta}{t_0}\right)^3
 =:\epsilon_3(\beta),                                     \tag{6.14}
\]
and hence
\[
 \left|e^{{\cal R}_3}-1\right|
 \le q_3(\beta):=e^{\epsilon_3(\beta)}-1
 =O(\beta^{-6/5}).                                        \tag{6.15}
\]
\end{lemma}

\begin{proof}
In the coefficient-majorant norm, deleting degrees zero, one, and two
leaves
\[
 \sum_{k\ge3}\|\ell_k\|(c_0r_\beta)^k
 \le
 \left(\frac{c_0r_\beta}{t_0}\right)^3
 \sum_{k\ge3}\|\ell_k\|t_0^k.
\]
Use (6.13) and the exponential estimate.
\end{proof}

\begin{lemma}[Uniform shifted-Gaussian exit]
\label{lem:n1v6-gaussian-exit}
Suppose, in addition to (6.7), that
\[
 A_{\rm abs}:=I-g^2\operatorname{Re}B,\qquad
 m_{\rm abs}:=gA_{\rm abs}^{-1}\operatorname{Re}a,
 \qquad
 \sup_{\boldsymbol s,w}\|m_{\rm abs}\|\le M_mg.            \tag{6.16}
\]
Put
\[
 Z_{2,\rm abs}:=
 {\mathbb E}_{\gamma_d}|{\cal Q}_2|.
\]
For \(R_\beta\ge2M_mg\),
\[
 \left|
 {\mathbb E}_{\gamma_d}
 \left[{\bf1}_{\{\|b\|>R_\beta\}}{\cal Q}_2\right]
 \right|
 \le
 Z_{2,\rm abs}(\boldsymbol s,w)\,
 2^{d/2}\exp\!\left\{-\frac{R_\beta^2}{32}\right\}.        \tag{6.17}
\]
\end{lemma}

\begin{proof}
Taking the absolute value replaces \(B,a,c\) by their real parts in
the exponent.  Completing that real square shows that the normalized
absolute majorant has covariance \(A_{\rm abs}^{-1}\le2I\), mean
\(m_{\rm abs}\), and total mass \(Z_{2,\rm abs}\).  On
\(\|b\|>R_\beta\), one has
\(\|b-m_{\rm abs}\|>R_\beta/2\).  For a centred Gaussian \(X\) with
covariance at most \(2I\), exponential Markov with parameter \(1/8\)
gives
\[
 {\mathbb P}\{\|X\|>R_\beta/2\}
 \le
 e^{-R_\beta^2/32}{\mathbb E}e^{\|X\|^2/8}
 \le2^{d/2}e^{-R_\beta^2/32}.
\]
Multiplication by the complete Gaussian normalization gives (6.17).
\end{proof}

\begin{lemma}[Wilson exterior tension]
\label{lem:n1v6-Wilson-exit}
Assume \(V(0,w)=0\), the local zero is unique, and
\({\cal N}_\beta(w)\ge n_-g^d\) uniformly on the retained real
domain.  Then for sufficiently large \(\beta\),
\[
 \left|
 \frac{{\cal Z}_{\rm ext}(\boldsymbol s,w)}
 {{\cal N}_\beta(w)(2\pi)^{d/2}}
 \right|
 \le
 C_{\rm ext}\beta^{d/2}
 \exp\!\left\{-\frac{c_0^2}{2}\beta^{1/5}\right\},         \tag{6.18}
\]
where bounded sources have been absorbed into \(C_{\rm ext}\).
\end{lemma}

\begin{proof}
Inside the fixed Morse neighbourhood but outside \(\Omega_\beta\),
the exact identity (5.15) gives
\[
                         V(U,w)\ge\frac12c_0^2r_\beta^2.
\]
On the complement of the fixed neighbourhood, compactness and
uniqueness of the zero give a positive minimum \(v_0\).  For large
\(\beta\), \(v_0\ge c_0^2r_\beta^2/2\).  Normalized Haar mass is at
most one, so the exterior compact integral is at most a fixed source
factor times
\[
 \exp\!\left\{-\frac12\beta c_0^2r_\beta^2\right\}.
\]
Divide by \(n_-g^d(2\pi)^{d/2}\), use
\(g^{-d}=\beta^{d/2}\), and
\(\beta r_\beta^2=\beta^{1/5}\).
\end{proof}

\begin{corollary}[One-block cutoff-free defect ledger]
\label{cor:n1v6-ledger}
Under the stated hypotheses, the exact local generating integral is a
shifted Gaussian two-jet plus three defects with respective primitive
sizes
\[
\boxed{
\begin{array}{c|c}
\text{defect}&\text{size}\\ \hline
\text{interior Taylor}&
C_2q_3(\beta)=O(\beta^{-6/5})\\
\text{Gaussian chart exit}&
C_2\,2^{d/2}e^{-c_0^2\beta^{1/5}/32}\\
\text{Wilson exterior}&
C_{\rm ext}\beta^{d/2}e^{-c_0^2\beta^{1/5}/2}.
\end{array}}                                               \tag{6.19}
\]
The soft-transition primitive (1.23) and geometric far-bad primitive
(1.7) remain separate and are not added twice to (6.19).
\end{corollary}

\begin{proof}
Combine Lemmas~\ref{lem:n1v6-interior},
\ref{lem:n1v6-gaussian-exit}, and
\ref{lem:n1v6-Wilson-exit}, using
\(R_\beta=c_0\beta^{1/10}\).  Completing the real square used in
Lemma~\ref{lem:n1v6-gaussian-exit} and shrinking the source polydisc
give a uniform constant
\[
                         Z_{2,\rm abs}\le C_2.
\]
The interior Taylor term is bounded by the same absolute normalization
times \(q_3(\beta)\).
\end{proof}

\section{Multi-block extension and its exact remaining hypothesis}
\label{sec:n1v6-multiblock}

For a set \(Y\) of disjoint Gaussian coordinate blocks with covariance
operator at most \(C_0I\), the same Chernoff calculation gives the
joint exit estimate
\[
 {\mathbb P}
 \left\{\|b_y\|>R_\beta\ \hbox{for every }y\in Y\right\}
 \le
 \left[
 2^{d_*/2}
 e^{-R_\beta^2/(4C_0)}
 \right]^{|Y|},                                           \tag{6.20}
\]
where \(d_*\) is the fixed block dimension.  Correlations do not spoil
the exponent because the restriction of the covariance to the union
of the disjoint coordinate sets still has operator norm at most
\(C_0\).

\begin{proof}
On the joint event,
\(\sum_{y\in Y}\|b_y\|^2\ge R_\beta^2|Y|\).  Apply exponential
Markov with parameter \(1/(4C_0)\).  The Gaussian determinant formula
in dimension \(d_*|Y|\) bounds the exponential moment by
\(2^{d_*|Y|/2}\), proving (6.20).
\end{proof}

\begin{firewall}[Local defects still need an owner product theorem]
\label{fire:n1v6-owner-product}
Equation (6.12) is exact and (6.19)--(6.20) are volume-free primitive
bounds.  To insert them into the KP gate one must still prove that the
local conditional Morse charts and their signed defect measures admit
a nonduplicating bounded-overlap owner product.  V5 does not supply
global locality of \(H^{1/2}\), and V6 does not infer it.  This owner
product must be built from fixed filling-neighbourhood charts or from
the archived energy-form localization, not from a global Hessian
square root.
\end{firewall}

\begin{assessment}[Progress after compact V6]
\label{ass:n1v6-status}
V6 corrects the scope of the Morse chart and proves the requested
cutoff-free local gluing.  The quadratic, linear, and constant
Jacobian jets are absorbed exactly into a shifted Gaussian.  Adding
and subtracting that full Gaussian produces three explicit signed
defects, none of whose indicators is differentiated.  Their sizes are
\(O(\beta^{-6/5})\) or stretched exponentially small, and correlated
Gaussian exits satisfy the joint bound (6.20).
\end{assessment}

\begin{decision}[V7 target]
\label{dec:n1v6-next}
V7 should construct the nonduplicating multi-block owner product for
(6.12).  It must use the forty-five-link conditional fibres and the
fixed filling-neighbourhood ownership already proved in the archive.
The decisive checks are:
\begin{enumerate}[label=\textup{(\roman*)},leftmargin=3.5em]
\item exterior conditioning does not create an unowned competing
stationary centre;
\item local Morse Jacobians overlap only through a fixed enlargement;
\item signed Taylor and Gaussian-exit defects satisfy a joint
conditional product bound; and
\item the shifted two-jet determinant is entered once in the marginal
ledger.
\end{enumerate}
If the first check fails for arbitrary exterior fields, the programme
must saturate that block into the bad component rather than choose a
nonuniform chart.
\end{decision}

\begin{firewall}[Claim boundary after compact V6]
\label{fire:n1v6-claim}
V6 proves a one-block exact gluing identity and joint Gaussian exit
bound, not the multi-block owner product or the numerical KP gate.
It does not finish the pushed Jacobian, marginal trajectory, stable RG
contraction, regulator removal, OS reconstruction, or Yang--Mills
mass gap.
\end{firewall}

\subsection{Parent record}

This V6 note was formed by appending the preceding material to the
validated compact V5, whose record was
\[
\begin{array}{c|c}
\text{lines}&2459\\
\text{bytes}&83549\\
\text{SHA--256}&
\texttt{E21004DB5552FF78FC35E71614733A0A5C47C7A3EC60F8E82BBD853A6F895961}.
\end{array}
\]

% =============================================================================
% END APPENDED COMPACT-SERIES V6 MATERIAL
% =============================================================================

% =============================================================================
% BEGIN APPENDED COMPACT-SERIES V7 MATERIAL
% =============================================================================

\chapter{V7: regular collars and a nonduplicating conditional-kernel product}
\label{ch:n1v7}

\section{Fixed local fibres and their collar witnesses}
\label{sec:n1v7-collars}

For a coarse block \(y\), let \({\cal C}_y\) be its forty-five-link
reduced fine fibre and let \({\cal N}_y\) be the fixed filling
neighbourhood containing every plaquette and retained datum on which
the conditional action depends.  Let \({\cal F}_y\) denote the compact
family of flat collar data, including the retained commuting toron
coordinate.  Define the gauge-invariant collar curvature
\[
 \kappa_y(W)^2
 =
 \sum_{p\subset{\cal N}_y\setminus{\cal C}_y}
 \|W_p-I\|_F^2.                                           \tag{7.1}
\]

\begin{lemma}[Fixed-collar reconstruction]
\label{lem:n1v7-collar-reconstruction}
There is a finite constant \(C_{\rm fill}\), independent of the
volume and of \(y\), such that
\[
 \dist_{\rm q}(W,{\cal F}_y)
 \le C_{\rm fill}\kappa_y(W)                              \tag{7.2}
\]
whenever the right side lies in the declared collar chart.  The
distance is taken after tree gauge and after retaining the toron
coordinates.
\end{lemma}

\begin{proof}
Gauge every link of a fixed collar spanning tree to the identity.
Every remaining contractible chord is the ordered product of
plaquette holonomies over one of the fixed filling surfaces.  If
\(A_1,\ldots,A_m\) are unitary, telescoping gives
\[
 \|A_1\cdots A_m-I\|_F
 \le\sum_{j=1}^m\|A_j-I\|_F.                              \tag{7.3}
\]
The filling length, the number of chords, and the incidence of a
plaquette in the filling table are fixed.  Cauchy--Schwarz and a sum
over this finite table yield (7.2).  Noncontractible chords are not
estimated by plaquettes; they are precisely the retained toron
coordinates in \({\cal F}_y\).
\end{proof}

\section{A uniform local stationary-centre card}
\label{sec:n1v7-centre}

In one toron patch, let \(x=0\) be the moving flat filling and write
\[
                         {\cal E}_y(x,W)=D_xV_y(x,W).       \tag{7.4}
\]
Compactness of the finite local dictionary supplies constants
\(\lambda_0,G_0,L_x,L_W>0\) such that, for \(W_0\in{\cal F}_y\),
\[
\begin{split}
 D_x{\cal E}_y(0,W_0)&\ge\lambda_0I,\\
 \|{\cal E}_y(0,W)\|&\le G_0\kappa_y(W),\\
 \|D_x{\cal E}_y(x,W)-D_x{\cal E}_y(0,W_0)\|
 &\le L_x\|x\|+L_W\kappa_y(W).                            \tag{7.5}
\end{split}
\]
The first line is the archived reduced vertical Hessian floor; the
other two are finite analytic derivative maxima combined with
Lemma~\ref{lem:n1v7-collar-reconstruction}.

Choose fixed \(r_0,\epsilon_0>0\) so that
\[
 \frac{G_0\epsilon_0}{\lambda_0}\le\frac{r_0}{4},
 \qquad
 \frac{L_xr_0+L_W\epsilon_0}{\lambda_0}\le\frac12.         \tag{7.6}
\]

\begin{theorem}[Regular collar gives one owned centre]
\label{thm:n1v7-regular-centre}
If \(\kappa_y(W)\le\epsilon_0\), then (7.4) has a unique solution
\(x_*(W)\) in \(\|x\|\le r_0\).  It depends analytically on \(W\) in
the smaller complex collar polydisc and
\[
 \|x_*(W)\|\le\frac{2G_0}{\lambda_0}\kappa_y(W),\qquad
 D_x^2V_y(x,W)\ge\frac{\lambda_0}{2}I                     \tag{7.7}
\]
throughout the real \(r_0\)-ball.
\end{theorem}

\begin{proof}
Let \(H_0=D_x{\cal E}_y(0,W_0)\) and define
\[
 {\cal T}_W(x)=x-H_0^{-1}{\cal E}_y(x,W).                  \tag{7.8}
\]
Equations (7.5)--(7.6) show that
\[
 \|D_x{\cal T}_W\|\le\frac12,\qquad
 \|{\cal T}_W(0)\|\le\frac{r_0}{4}.                        \tag{7.9}
\]
Thus \({\cal T}_W\) maps the \(r_0\)-ball into itself and is a
contraction.  Its unique fixed point is \(x_*(W)\).  The sharper first
bound in (7.7) follows by applying the contraction estimate at zero.
The Hessian bound follows from the last line of (7.5).  Uniform
analytic Picard iteration proves analytic dependence on \(W\).
\end{proof}

\begin{lemma}[Any competing centre is behind a fixed barrier]
\label{lem:n1v7-centre-barrier}
After possibly decreasing \(\epsilon_0\), there is \(v_*>0\) such
that, for every regular collar,
\[
 V_y(x,W)-V_y(x_*(W),W)\ge v_*                            \tag{7.10}
\]
whenever \(x\) is outside the \(r_0\)-ball but remains in the compact
reduced fibre.
\end{lemma}

\begin{proof}
For flat collar data, the archived local-vacuum theorem gives one
reduced zero, uniformly over the compact retained toron family.
The continuous action difference therefore has a positive minimum on
the compact complement of the \(r_0\)-ball.  Analytic dependence of
the action and Theorem~\ref{thm:n1v7-regular-centre} preserve half of
that minimum for sufficiently small \(\epsilon_0\).  Rename the
halved minimum \(v_*\).
\end{proof}

\begin{definition}[Regular and saturated-bad blocks]
\label{def:n1v7-regular}
A block is collar-regular if
\(\kappa_y(W)\le\epsilon_0\).  On a regular block, only the running
Morse ball around \(x_*(W)\) is assigned to the good kernel.  Its
complement is a chart-exit defect.  A collar-irregular block and every
block whose filling neighbourhood intersects it are placed in the
same saturated bad component.
\end{definition}

\begin{assessment}[No unowned competing centre]
\label{ass:n1v7-no-unowned-centre}
Two centres cannot occur inside a regular \(r_0\)-ball by
Theorem~\ref{thm:n1v7-regular-centre}.  Any centre outside that ball
pays the fixed barrier (7.10) and is a chart-exit defect.  If the
regularity hypothesis fails, (7.1) is an owned collar-curvature
witness.  Thus arbitrary exterior conditioning creates no unclassified
stationary branch.
\end{assessment}

\section{The collar-bad Peierls activity}
\label{sec:n1v7-collar-bad}

Let \(M_\kappa\) be the maximum number of collar sums (7.1) containing
one plaquette.  It is finite because \({\cal N}_y\) has fixed range.

\begin{lemma}[Joint collar witness]
\label{lem:n1v7-joint-collar}
For every set \(Y\) of collar-irregular blocks,
\[
 \sum_{p\subset\bigcup_{y\in Y}{\cal N}_y}
 \|U_p-I\|_F^2
 \ge\frac{\epsilon_0^2}{M_\kappa}|Y|.                     \tag{7.11}
\]
Consequently their Wilson action is at least
\[
                         \frac{\epsilon_0^2}{6M_\kappa}|Y|. \tag{7.12}
\]
\end{lemma}

\begin{proof}
Sum \(\kappa_y(W)^2>\epsilon_0^2\) over \(y\in Y\).
Every plaquette term appears at most \(M_\kappa\) times, which proves
(7.11).  For \(SU(3)\),
\[
 1-\frac13\operatorname{ReTr}U_p
 =\frac16\|U_p-I\|_F^2,
\]
giving (7.12).
\end{proof}

After division by a local Laplace normalization, (7.12) gives a
primitive of the form
\[
 q_{\rm collar}(\beta)
 =
 C_{\rm collar}\beta^{d_*/2}
 \exp\!\left\{
 -\frac{\beta\epsilon_0^2}{6M_\kappa}
 \right\}.                                                \tag{7.13}
\]
It is exponentially smaller than every running-window activity.

\section{Exact regular conditional kernels}
\label{sec:n1v7-kernels}

Let \(P_y\) be the exact normalized conditional Gibbs kernel on
\({\cal C}_y\), at fixed exterior field.  On a regular collar, let
\(\widehat{\cal U}_y(\,\cdot\,,W):\mathbb R^{d_*}\to{\cal C}_y\)
be any Borel extension of the inverse Morse map which agrees with it
on \(\|b\|\le R_\beta\) and maps the complement into the same local
fibre.  Let \(G_y\) be the pushforward by
\(\widehat{\cal U}_y\) of the normalized real shifted two-jet Gaussian
of V6.  Thus \(G_y\) is a probability kernel on the same compact
fibre as \(P_y\).  Define the signed
kernel
\[
                         \Delta_y=P_y-G_y.                 \tag{7.14}
\]
For a signed kernel \(K\), use the variation operator norm
\[
 \|K\|_{\rm var}
 =
 \sup_\omega\sup_{\|f\|_\infty\le1}|Kf(\omega)|.           \tag{7.15}
\]

Let
\[
\begin{split}
 \varepsilon_6(\beta):={}&
 C_2q_3(\beta)
 +C_2\,2^{d_*/2}e^{-c_0^2\beta^{1/5}/32}\\
 &+
 C_{\rm ext}\beta^{d_*/2}
 e^{-c_0^2\beta^{1/5}/2}.                                 \tag{7.16}
\end{split}
\]

\begin{lemma}[Normalization converts the V6 defects to a kernel
bound]
\label{lem:n1v7-kernel-bound}
If \(\varepsilon_6(\beta)<1\), then uniformly over regular exterior
fields,
\[
 \boxed{
 \|\Delta_y\|_{\rm var}
 \le
 q_{\rm ker}(\beta):=
 \frac{2\varepsilon_6(\beta)}
 {1-\varepsilon_6(\beta)}.}                               \tag{7.17}
\]
\end{lemma}

\begin{proof}
Apply the V6 add--subtract identity with an arbitrary measurable test
function \(\|f\|_\infty\le1\).  On the running ball, the Borel
extension agrees with the exact inverse Morse map.  Outside that ball,
its entire contribution is bounded by the Gaussian-exit term, so no
regularity of the extension is used.

Write the unnormalized exact conditional measure as
\[
                         Z_2G_y+\eta_y,
\]
where \(G_y\) is a probability and V6 gives
\(\|\eta_y\|_{\rm TV}\le Z_2\varepsilon_6\).
Its total mass is
\(Z_y=Z_2+\eta_y(1)\ge Z_2(1-\varepsilon_6)\).  Therefore
\[
 P_y-G_y
 =
 \frac{\eta_y}{Z_y}
 +\left(\frac{Z_2}{Z_y}-1\right)G_y.
\]
The variation norm is at most
\[
 \frac{\|\eta_y\|_{\rm TV}+|\eta_y(1)|}{Z_y}
 \le\frac{2\varepsilon_6}{1-\varepsilon_6}.
\]
\end{proof}

\section{The exact one-color owner product}
\label{sec:n1v7-product}

Colour the fixed-range overlap graph of the neighbourhoods
\({\cal N}_y\).  A finite number \(\chi_*\le\Delta_*+1\) of colours
suffices, and neighbourhoods belonging to one colour are disjoint.
Fix one colour \(C\) and freeze its exterior.  No Wilson interaction
meets two fibres of \(C\), so their exact conditional kernel
factorizes:
\[
                         P_C=\bigotimes_{y\in C}P_y.        \tag{7.18}
\]
Let \(C_{\rm reg}\) and \(C_{\rm bad}\) be the regular and
collar-irregular subsets determined by this frozen exterior.

\begin{theorem}[Exact frozen-colour defect expansion]
\label{thm:n1v7-colour-product}
\[
\boxed{
 P_C
 =
 \sum_{S\subset C_{\rm reg}}{\cal T}_{C,S},}               \tag{7.19}
\]
where
\[
 {\cal T}_{C,S}
 =
 \left(\bigotimes_{y\in C_{\rm bad}}P_y\right)
 \otimes
 \left(\bigotimes_{y\in S}\Delta_y\right)
 \otimes
 \left(\bigotimes_{y\in C_{\rm reg}\setminus S}G_y\right).
                                                               \tag{7.19a}
\]
Every colour block occurs exactly once.  The regular signed part obeys
\[
 \|{\cal T}_{C,S}\|_{\infty\to\infty}
 \le q_{\rm ker}(\beta)^{|S|}.                             \tag{7.20}
\]
\end{theorem}

\begin{proof}
Use the conditional factorization (7.18), substitute
\(P_y=G_y+\Delta_y\) for \(y\in C_{\rm reg}\), and distribute the
finite tensor product.  This proves (7.19)--(7.19a).
Every exact or Gaussian conditional probability has \(L^\infty\)
operator norm one, while (7.17) bounds every signed factor.  Tensor
submultiplicativity gives (7.20).
\end{proof}

\begin{remark}[Why overlap does not duplicate a determinant]
\label{rem:n1v7-determinant}
The normalized \(G_y\) contains no scalar determinant.  The two-jet
normalization (6.8) is assigned separately to the lexicographic owner
\(y\) when its conditional kernel is formed.  Equation (7.19) expands
normalized kernels only, so the colour product cannot charge that
determinant again.
\end{remark}

\section{Support graph and connected collection}
\label{sec:n1v7-support}

Join \(y\) and \(z\) when
\({\cal N}_y\cap{\cal N}_z\ne\varnothing\), and let
\(\Delta_*\) be the maximum degree of this fixed-range graph.  The
frozen-colour products above have disjoint supports.  Collection across
different colours requires an ordered schedule because updating one
colour can change the collar data seen by the next.

\begin{lemma}[Connected-word count]
\label{lem:n1v7-connected-count}
The number of connected block sets of size \(n\) containing a fixed
root is at most
\[
                         \{e(\Delta_*+1)\}^{n-1}.           \tag{7.21}
\]
\end{lemma}

\begin{proof}
Every connected set has a rooted spanning tree.  Explore that tree in
depth-first order and encode each new vertex by one of at most
\(\Delta_*+1\) neighbour choices.  The standard exponential
generating bound for rooted trees supplies \(e^{n-1}\); multiplying
the choices gives (7.21).
\end{proof}

Let \(M_{\rm sat}\) be the maximum number of original blocks in one
fixed saturation neighbourhood and put
\[
 q_{\rm sat}(\beta)
 =
 q_{\rm collar}(\beta)^{1/M_{\rm sat}}.                   \tag{7.22}
\]
Selecting a disjoint witness subfamily of size at least
\(|Y|/M_{\rm sat}\) and using (7.11) yields this per-saturated-block
activity.

\begin{theorem}[Polymer gate under a colour-compatible schedule]
\label{thm:n1v7-kernel-gate}
Define
\[
                         u_7(\beta)
 =q_{\rm ker}(\beta)+q_{\rm sat}(\beta).                  \tag{7.23}
\]
Assume an all-colour schedule has the following property: once an
owner is declared regular or is placed in a saturated bad component,
later colour kernels can change that declaration only by enlarging
the same saturated component, and every original block is assigned
once.  If
\[
 e^{1+\mu}(\Delta_*+1)u_7(\beta)<1,                       \tag{7.24}
\]
then the connected scheduled-kernel defect expansion converges
absolutely in the block weight \(e^{\mu|Y|}\), uniformly in the
volume.  Since
\[
 q_{\rm ker}(\beta)=O(\beta^{-6/5}),\qquad
 q_{\rm sat}(\beta)=
 O\!\left(
 \beta^{d_*/(2M_{\rm sat})}e^{-c\beta}
 \right),                                                 \tag{7.25}
\]
(7.24) holds for all sufficiently large \(\beta\).
\end{theorem}

\begin{proof}
For a regular connected scheduled word, (7.20) pays
\(q_{\rm ker}^{|Y|}\).  Saturated collar-bad words pay
\(q_{\rm sat}^{|Y|}\) by the witness selection preceding the theorem.
The schedule hypothesis prevents a later colour from creating a
second owner for either payment.
Mixed words are bounded by expanding which type owns every block and
using the binomial theorem, giving \(u_7^{|Y|}\).  Multiply by the
connected-set count (7.21) and by \(e^{\mu|Y|}\); the resulting
geometric series has ratio bounded by the left side of (7.24).
Equation (7.25) follows from (7.13), (7.16), and (7.17).
\end{proof}

\begin{firewall}[The all-colour schedule is not yet proved]
\label{fire:n1v7-kernel-vs-rg}
Theorem~\ref{thm:n1v7-colour-product} is unconditional for one frozen
colour.  Theorem~\ref{thm:n1v7-kernel-gate} is conditional on the
stated all-colour schedule.  V7 does not prove that successive colour
updates preserve earlier collar declarations.  Nor does it prove that
one sweep has a uniform physical-sector contraction or that iterating
the reference sweep reproduces the exact coarse pushforward.
\end{firewall}

\begin{assessment}[Progress after compact V7]
\label{ass:n1v7-status}
V7 constructs the requested nonduplicating owner product for every
frozen colour layer.  Small collar curvature gives one uniformly
nondegenerate local centre; failed collars and remote competing
centres have explicit owned action witnesses.  The ordered
colour expansion (7.19) pays every regular defect once and the
two-jet determinant is separately owned once.  The numerical
kernel-polymer gate closes at sufficiently large \(\beta\) once the
explicit all-colour schedule card is established.
\end{assessment}

\begin{decision}[V8 target]
\label{dec:n1v7-next}
V8 should first prove the all-colour saturation schedule required by
Theorem~\ref{thm:n1v7-kernel-gate}.  It should then compare the exact
heat-bath sweep with the shifted-Gaussian reference sweep in an
influence norm.  It must prove either:
\begin{enumerate}[label=\textup{(\roman*)},leftmargin=3.5em]
\item a strict reference-sweep contraction on the reduced physical
sector, stable under the defect norm \(u_7\); or
\item an exact resolvent/Poisson comparison showing that their
invariant conditional measures differ by a connected kernel polymer
of size \(O(u_7)\).
\end{enumerate}
Retained slow and toron variables must remain parameters; they may not
be included in the contracted norm.
\end{decision}

\begin{firewall}[Claim boundary after compact V7]
\label{fire:n1v7-claim}
V7 closes the frozen-colour conditional-kernel owner product.  Its
large-\(\beta\) all-colour connected gate remains conditional on the
schedule card.  It does not prove a physical-sector sweep contraction,
identify the exact invariant coarse kernel,
complete the pushed Jacobian or marginal flow, construct the full RG
trajectory, remove regulators, perform OS reconstruction, or prove
the Yang--Mills mass gap.
\end{firewall}

\subsection{Parent record}

This V7 note was formed by appending the preceding material to the
validated compact V6, whose record was
\[
\begin{array}{c|c}
\text{lines}&2925\\
\text{bytes}&98414\\
\text{SHA--256}&
\texttt{14B8E1D7A4E622034681495B9A94F68581D8EB43E1F1904755CE7DD12795198F}.
\end{array}
\]

% =============================================================================
% END APPENDED COMPACT-SERIES V7 MATERIAL
% =============================================================================

% =============================================================================
% BEGIN APPENDED COMPACT-SERIES V8 MATERIAL
% =============================================================================

\chapter{V8: a triangular all-colour saturation schedule}
\label{ch:n1v8}

\section{What has to be scheduled}
\label{sec:n1v8-setup}

The missing hypothesis in
Theorem~\ref{thm:n1v7-kernel-gate} is not a harmless choice of
notation.  Conditional kernels belonging to different colours see
different exterior data, and multiplying all of them as if those data
were frozen would in general be false.  The correct object is a
triangular disintegration: integrate one colour, retain its exact
normalization in the marginal action, and form the next conditional
kernel from that marginal action.  An integrated fibre is never
updated or classified again.

Let \({\cal B}\) be the finite block set and retain the disjoint
reduced coordinate fibres
\[
             {\cal X}=\prod_{y\in{\cal B}}{\cal X}_y .
                                                               \tag{8.1}
\]
The sets \({\cal X}_y\) are disjoint even though their filling
neighbourhoods \({\cal N}_y\) overlap.  Fix once and for all a proper
colouring
\[
 {\cal B}=C_1\mathbin{\dot\cup}\cdots
 \mathbin{\dot\cup}C_{\chi_*},
 \qquad
 \chi_*\le\Delta_*+1,                                      \tag{8.2}
\]
of the overlap graph of the \({\cal N}_y\).  Put
\[
 I_a=\bigcup_{b\le a}C_b,\qquad
 R_a={\cal B}\setminus I_a,\qquad I_0=\varnothing.          \tag{8.3}
\]
Slow coarse fields, torons, and the four source variables are denoted
collectively by \(s\); they are parameters and are never members of
\(I_a\).

\begin{definition}[Exact triangular colour pushforward]
\label{def:n1v8-pushforward}
Let \(W_{a-1}(x_{C_a},x_{R_a};s)\) be the exact Boltzmann density after
the fibres in \(I_{a-1}\) have been integrated.  Define
\[
 W_a(x_{R_a};s)
 =
 \int_{{\cal X}_{C_a}}
 W_{a-1}(x_{C_a},x_{R_a};s)\,dH_{C_a}(x_{C_a}),             \tag{8.4}
\]
where \(dH_{C_a}\) is product Haar measure on the reduced fibres of
that colour.  At zero source, whenever \(W_a>0\), set
\[
 P_a(dx_{C_a}\mid x_{R_a};s)
 =
 \frac{W_{a-1}(x_{C_a},x_{R_a};s)}
      {W_a(x_{R_a};s)}\,dH_{C_a}(x_{C_a}).                  \tag{8.5}
\]
For complex sources, (8.4), rather than a logarithm or a choice of
phase, is the definition; identities first proved at real source are
continued analytically.
\end{definition}

\begin{lemma}[Exactness and absence of repeated updates]
\label{lem:n1v8-fubini}
The maps (8.4) satisfy
\[
 W_a
 =
 \int_{{\cal X}_{I_a}}W_0\,dH_{I_a},\qquad
 W_{\chi_*}
 =
 \int_{\cal X}W_0\,dH_{\cal B}.                             \tag{8.6}
\]
At physical source, the joint conditional law admits the triangular
factorization
\[
 dP(x_{\cal B}\mid s)
 =
 P_1(dx_{C_1}\mid x_{R_1};s)\cdots
 P_{\chi_*}(dx_{C_{\chi_*}}\mid s).                         \tag{8.7}
\]
In particular, a variable in \(C_a\) occurs in one and only one
conditional factor.
\end{lemma}

\begin{proof}
Equation (8.6) follows by induction from (8.4) and Tonelli's theorem
at physical source.  The spaces are compact and all finite-volume
densities are continuous, so the integrals are finite.  Successive
substitution of (8.5) into
\[
 W_{a-1}\,dH_{C_a}=P_a\,W_a
\]
telescopes all intermediate \(W_a\)'s and proves (8.7).  Fubini's
theorem gives (8.6) for complex sources as well.  Both sides are
analytic on the common source polydisc, so (8.7), interpreted as the
corresponding identity between unnormalized analytic kernels,
continues there.  Disjointness in (8.2)--(8.3) proves the last
assertion.
\end{proof}

\begin{remark}[Why this is not an ordinary heat-bath sweep]
\label{rem:n1v8-not-sampler}
If one repeatedly applied full conditionals of the original action,
earlier variables would remain present and their collar data could be
changed by later updates.  That product need not equal the original
joint law.  In (8.4), earlier variables have instead been eliminated
and their effect is carried exactly by \(W_a\).  Lemma
\ref{lem:n1v8-fubini}, not a stationarity assertion for a Gibbs
sampler, is the algebraic basis of the schedule below.
\end{remark}

\section{Calibrated rows and the stable owner norm}
\label{sec:n1v8-norm}

At every stage split the effective local action into three rows:
\[
 {\cal S}_a
 =
 {\cal S}^{\rm cal}_a+{\cal J}_a+\Psi_a.                   \tag{8.8}
\]
Here \({\cal S}^{\rm cal}_a\) belongs to the fixed compact dictionary
of reduced Wilson, covariance, slow, and toron backgrounds;
\({\cal J}_a\) contains constants, currents, quadratic forms,
determinants, and the declared marginal source rows; and \(\Psi_a\)
is the connected stable owner family.  The extraction of
\({\cal J}_a\) is the exact finite Taylor-jet projection used in V6.
No determinant is put into a normalized kernel.

There is one important bookkeeping convention.  The initial
calibrated rows may be of order one.  Their archived uniform Hessian
and analytic bounds are part of \({\cal S}^{\rm cal}_0\), not part of
a small norm.  Every increment generated during the present sweep,
including its extracted linear and quadratic jet, is counted in the
increment norm before it is transported into
\({\cal S}^{\rm cal}_a+{\cal J}_a\).  Thus a smallness statement below
never silently discards a newly generated Hessian.

For a connected owner family \(\Phi=\{\Phi(Y)\}\), define on a
strictly smaller source polydisc
\[
 \|\Phi\|_{\mu,\rho}
 =
 \sup_{z\in{\cal B}}
 \sum_{\substack{Y\ni z\\Y\ {\rm connected}}}
 e^{\mu|Y|}\|\Phi(Y)\|_\rho.                               \tag{8.9}
\]
The local analytic norm \(\|\cdot\|_\rho\) is the compact
right-invariant Taylor--Wiener norm of V3, augmented by the four
source radii.  Let
\[
 \eta_a
 =
 \left\|
   \hbox{all generated increments through colour \(a\)}
 \right\|_{\mu,\rho},\qquad \eta_0=0.                       \tag{8.10}
\]

\begin{lemma}[Uniform local derivative embeddings]
\label{lem:n1v8-embeddings}
After one fixed restriction of the field and source radii, there are
volume-independent constants \(C_1,C_2<\infty\) such that the
generated perturbation seen by any active fibre satisfies
\[
 \sup_y\|D_x\Psi^{\rm loc}_{a,y}\|
 \le C_1\eta_a,\qquad
 \sup_y\|D_x^2\Psi^{\rm loc}_{a,y}\|
 \le C_2\eta_a.                                            \tag{8.11}
\]
The same inequalities hold for the sum of transported generated
jet increments.
\end{lemma}

\begin{proof}
Only owners meeting \({\cal N}_y\) enter
\(\Psi^{\rm loc}_{a,y}\).  Apply the Cauchy estimate on the fixed
difference between the original and restricted field radii.  For
each owner, the first or second derivative costs the corresponding
inverse radius difference.  The number of endpoint coordinates in
one \({\cal N}_y\), and the number of \(y\)'s whose neighbourhood can
contain a fixed endpoint, are bounded by the fixed local dictionary.
Summing with (8.9) proves (8.11).  The finite Taylor-jet projections
and their inclusions have bounded operator norm on the same smaller
polydisc, so the identical argument includes the transported
generated jets.
\end{proof}

\section{A quantitative stability reserve for later colours}
\label{sec:n1v8-stability}

Shrink the V7 constants, if necessary, to
\(\bar r_0\le r_0\) and \(\bar\epsilon_0\le\epsilon_0\) such that
\[
 \frac{G_0\bar\epsilon_0}{\lambda_0}
 \le\frac{\bar r_0}{8},\qquad
 \frac{L_x\bar r_0+L_W\bar\epsilon_0}{\lambda_0}
 \le\frac14.                                               \tag{8.12}
\]
This is possible because all four constants are fixed and
\(\lambda_0>0\).  Let \(\eta_{\rm an}>0\) be the distance, in the
norm (8.9), from the smaller V5 Morse majorant ball to the boundary of
the larger analytic ball.  Compactness of the local dictionary and
the strict radius restriction make this distance positive.  Put
\[
 \eta_*=
 \min\left\{
 \frac{\lambda_0\bar r_0}{8C_1},
 \frac{\lambda_0}{4C_2},
 \eta_{\rm an}
 \right\}.                                                 \tag{8.13}
\]

\begin{theorem}[Morse centre stability under generated owners]
\label{thm:n1v8-morse-stability}
Suppose the current collar satisfies
\(\kappa_y(W)\le\bar\epsilon_0\) and the cumulative generated
increment satisfies \(\eta_a\le\eta_*\).  Then the perturbed local
stationary equation
\[
 D_x\{V_y(x,W)+\Psi^{\rm loc}_{a,y}(x,W)\}=0               \tag{8.14}
\]
has exactly one solution in \(\|x\|\le\bar r_0\).  Throughout that
real ball its reduced Hessian is at least
\(\lambda_0 I/2\), and the solution and its local Morse map are
analytic on the restricted complex polydisc.
\end{theorem}

\begin{proof}
Let \(H_0=D_x{\cal E}_y(0,W_0)\) be the V7 flat-collar Hessian and set
\[
 {\cal T}^{\Psi}_W(x)
 =
 x-H_0^{-1}
 \bigl\{
 {\cal E}_y(x,W)+D_x\Psi^{\rm loc}_{a,y}(x,W)
 \bigr\}.                                                  \tag{8.15}
\]
Equations (7.5), (8.11), (8.12), and (8.13) give
\[
 \|{\cal T}^{\Psi}_W(0)\|
 \le\frac{\bar r_0}{8}+\frac{\bar r_0}{8}
 =\frac{\bar r_0}{4},\qquad
 \|D_x{\cal T}^{\Psi}_W\|
 \le\frac14+\frac14=\frac12.                               \tag{8.16}
\]
Hence \({\cal T}^{\Psi}_W\) maps the closed
\(\bar r_0\)-ball into itself and is a contraction.  Its fixed point
is the unique solution of (8.14) in the ball.  Moreover,
\[
 D_x^2(V_y+\Psi^{\rm loc}_{a,y})
 \ge
 \lambda_0I-\frac{\lambda_0}{4}I-\frac{\lambda_0}{4}I
 =\frac{\lambda_0}{2}I.                                   \tag{8.17}
\]
Uniform analytic Picard iteration gives analytic dependence of the
centre.  The last entry in (8.13), together with the V5
square-root majorant, keeps the analytic Hessian square root and the
inverse Morse map inside their convergence balls.
\end{proof}

\begin{corollary}[The V6 kernel estimate survives every admissible
stage]
\label{cor:n1v8-kernel-survives}
On any history for which \(\eta_a\le\eta_*\), every later
collar-regular fibre has the same form of exact
Morse--Gaussian decomposition as in V6--V7.  After increasing its
fixed prefactor, but not changing its \(\beta\)-power,
\[
 \|\Delta_{a,y}\|_{\rm TV}\le q_{\rm ker}(\beta)
 =O(\beta^{-6/5}).                                        \tag{8.18}
\]
Collar-irregular histories remain owned by the V7 Wilson-action
witness and are not required to satisfy an analytic Morse estimate.
\end{corollary}

\begin{proof}
Theorem~\ref{thm:n1v8-morse-stability} supplies the uniform positive
Hessian, analytic chart, and inverse-chart radii used in the V6
proof.  All altered constants range over a compact subset of the
strict inequalities (8.12)--(8.13); their maxima can therefore be
absorbed into one fixed prefactor.  The Taylor defect is still cubic
after exact two-jet extraction and the running radius is unchanged,
so its leading power remains \(\beta^{-6/5}\).  The chart-exit terms
remain exponentially small.  The collar-irregular statement is
exactly Lemma~\ref{lem:n1v7-joint-collar}.
\end{proof}

\section{The finite connected-output compiler}
\label{sec:n1v8-compiler}

We record all harmless but necessary conversion constants rather than
calling them one.  Let
\[
\begin{split}
 c_{\rm own}&=\hbox{operator norm of the canonical owner
 projection},\\
 c_{\rm jet}&=\max\{\|{\cal J}\|,\|1-{\cal J}\|\},\\
 c_{\rm src}&=\hbox{the fixed four-source Cauchy restriction
 constant},\\
 c_{\rm loc}&=\hbox{the maximum local endpoint-incidence
 constant},
\end{split}
\]
and define
\[
 c_{\rm cmp}
 =
 \max\{1,c_{\rm own}c_{\rm jet}c_{\rm src}c_{\rm loc}\}.   \tag{8.19}
\]
Every entry is a maximum over a finite local dictionary and strict
radius inclusions; hence \(c_{\rm cmp}<\infty\) and is independent of
the volume and \(\beta\).  Normalized probability kernels cost one in
supremum norm.  Their scalar normalizations and determinants are
transported in \({\cal J}_a\) and are not charged a second time in
\(c_{\rm cmp}\).

Set
\[
 A_\mu=e^{1+\mu}(\Delta_*+1),\qquad
 {\mathfrak C}_\mu(v)=\frac{e^\mu v}{1-A_\mu v}.            \tag{8.20}
\]

\begin{lemma}[Rooted connected-output bound]
\label{lem:n1v8-compiler}
Suppose the total primitive strength incident on any block is at most
\(v\) after the conversions in (8.19).  If \(A_\mu v<1\), the
connected outputs containing a fixed root have total stable norm at
most
\[
                         {\mathfrak C}_\mu(v).              \tag{8.21}
\]
\end{lemma}

\begin{proof}
Expand each connected output into its primitive owner word and choose
the least block as root.  A word on \(n\) distinct blocks contributes
at most \(v^n\) after (8.19).  By
Lemma~\ref{lem:n1v7-connected-count}, its connected support has at
most
\(\{e(\Delta_*+1)\}^{n-1}\) rooted choices.  The weight in (8.9)
contributes \(e^{\mu n}\).  Therefore the root sum is bounded by
\[
 \sum_{n\ge1}
 e^{\mu n}\{e(\Delta_*+1)\}^{n-1}v^n
 =
 \frac{e^\mu v}{1-e^{1+\mu}(\Delta_*+1)v},
\]
which is (8.21).  Canonical rooting removes any multiplicity from
possible spanning trees; the spanning-tree count is used only as an
upper bound.
\end{proof}

\begin{definition}[Pass, hit, and exact jet transport]
\label{def:n1v8-pass-hit}
Immediately before colour \(C_a\), split every current owner
canonically as
\[
 \Psi_{a-1}=\Psi^{\rm pass}_{a-1}+\Psi^{\rm hit}_{a-1},    \tag{8.22}
\]
where a pass owner has support disjoint from every
\({\cal N}_y\), \(y\in C_a\), and a hit owner is every other owner.
Pass owners are copied unchanged.  Hit owners are consumed in the
colour integral and replaced by its connected outputs.  The finite
jet of each output is extracted exactly once and transported into
\({\cal J}_a\); its stable complement becomes a new owner rooted at
the least block of its support.
\end{definition}

\begin{lemma}[Finite-colour recursion]
\label{lem:n1v8-recursion}
Assume \(\eta_{a-1}\le\eta_*\), and let
\[
                         v_a
 =c_{\rm cmp}\{\eta_{a-1}+u_7(\beta)\}.                    \tag{8.23}
\]
If \(A_\mu v_a<1\), then the exact \(a\)-th colour pushforward is
well-defined in the owner algebra and
\[
 \eta_a
 \le
 \eta_{a-1}+{\mathfrak C}_\mu(v_a).                        \tag{8.24}
\]
\end{lemma}

\begin{proof}
On the zero-defect branch,
Corollary~\ref{cor:n1v8-kernel-survives} constructs the normalized
reference kernels.  Every nontrivial primitive meeting the current
colour is either a previously generated hit owner, a regular signed
kernel \(\Delta_{a,y}\), or a newly born saturated witness.  Their
root activity is bounded respectively by
\(\eta_{a-1}\), \(q_{\rm ker}\), and \(q_{\rm sat}\).
Thus (8.19) bounds the total primitive strength by (8.23).
Lemma~\ref{lem:n1v8-compiler} bounds all connected replacements by
\({\mathfrak C}_\mu(v_a)\).  Pass owners are unchanged and contribute
at most \(\eta_{a-1}\).  Exact jet extraction cannot increase the
bound beyond the factor already included in \(c_{\rm cmp}\), proving
(8.24).  Absolute convergence permits the integrations and connected
collection to be interchanged.
\end{proof}

\section{Monotone saturation and unique payment}
\label{sec:n1v8-saturation}

Fix a total order on \({\cal B}\).  For \(S\subset{\cal B}\), define
its one-step filling saturation by
\[
 {\rm Sat}(S)
 =
 \{z\in{\cal B}:
 {\cal N}_z\cap{\cal N}_y\ne\varnothing
 \hbox{ for some }y\in S\}.                                \tag{8.25}
\]
Its cardinality for a singleton is at most \(M_{\rm sat}\), after
enlarging the V7 constant by a fixed amount if necessary.

\begin{definition}[Historical saturation algorithm]
\label{def:n1v8-history}
Start with no assigned blocks and empty saturated set \(B_0\).
At stage \(a\), process the blocks of \(C_a\) in the fixed order.
\begin{enumerate}[label=\textup{(\roman*)},leftmargin=3.5em]
\item If \(y\) was assigned at an earlier birth, integrate its fibre
with the exact residual kernel but create no new mark.
\item If \(y\) is unassigned and its current collar is regular, assign
the permanent mark \((y,{\rm reg},a)\).
\item If \(y\) is unassigned and its current collar is irregular, set
\(H_y={\rm Sat}(\{y\})\), enlarge the saturated set by \(H_y\), and
assign every still-unassigned \(z\in H_y\) the mark
\((y,{\rm bad},a)\).
\end{enumerate}
Whenever two saturated sets meet, replace them by their union.  The
owner labels on blocks already assigned are never altered.
\end{definition}

\begin{lemma}[Monotonicity and exhaustion]
\label{lem:n1v8-history}
The historical algorithm has the following properties.
\begin{enumerate}[label=\textup{(\alph*)},leftmargin=3.5em]
\item The saturated set and each of its connected components can only
grow or merge.
\item Every original block receives exactly one mark by the end of
stage \(\chi_*\).
\item A previously regular mark swallowed by a later saturated
component is not charged again.
\item Every bad birth owns at least one and at most \(M_{\rm sat}\)
previously unassigned blocks.
\end{enumerate}
\end{lemma}

\begin{proof}
Only unions are performed, proving (a).  The assignment rule is
applied only to unassigned blocks and never changes a label, proving
uniqueness.  If a block is still unassigned when its unique colour is
processed, either (ii) or (iii) assigns it; hence all blocks are
assigned, proving (b).  The immutable label rule gives (c).  At a bad
birth the centre \(y\) is itself unassigned, while (8.25) and the
definition of \(M_{\rm sat}\) give the upper bound in (d).
\end{proof}

\begin{lemma}[Bad-witness cost can be distributed without
duplication]
\label{lem:n1v8-bad-distribution}
Assume \(q_{\rm collar}(\beta)\le1\).  If a bad birth owns \(m\)
new blocks, then
\[
 q_{\rm collar}(\beta)
 \le q_{\rm sat}(\beta)^m,\qquad 1\le m\le M_{\rm sat}.     \tag{8.26}
\]
For any family of bad births, the product of their witness weights is
therefore bounded by \(q_{\rm sat}\) to the number of bad-assigned
blocks.
\end{lemma}

\begin{proof}
By (7.22),
\[
 q_{\rm sat}^m=q_{\rm collar}^{m/M_{\rm sat}}.
\]
Since \(0<q_{\rm collar}\le1\) and \(m/M_{\rm sat}\le1\), this is at
least \(q_{\rm collar}\), proving (8.26).  Distinct births own
disjoint sets of newly assigned blocks by
Definition~\ref{def:n1v8-history}.  Multiplication proves the final
claim.  The joint action estimate
Lemma~\ref{lem:n1v7-joint-collar} supplies the product witness even
when the filling neighbourhoods of different birth blocks overlap.
\end{proof}

\section{Closure of the all-colour schedule card}
\label{sec:n1v8-closure}

Put
\[
 B_\mu=2e^\mu c_{\rm cmp},\qquad
 D_\mu=1+B_\mu,\qquad
 K_a=D_\mu^a-1.                                            \tag{8.27}
\]

\begin{theorem}[Triangular all-colour schedule]
\label{thm:n1v8-all-colour}
Suppose \(\beta\) is large enough that
\[
\boxed{
 u_7(\beta)
 \le
 \min\left\{
 \frac{\eta_*}{K_{\chi_*}},
 \frac{1}
 {2A_\mu c_{\rm cmp}(K_{\chi_*}+1)}
 \right\}.}                                                \tag{8.28}
\]
Then the exact colour pushforwards (8.4), equipped with
Definition~\ref{def:n1v8-history}, satisfy all of the following,
uniformly in the finite volume:
\begin{enumerate}[label=\textup{(\roman*)},leftmargin=3.5em]
\item every active regular fibre has a unique analytic Morse centre
and the V7 signed-kernel estimate;
\item every original block is paid exactly once, either by a regular
kernel mark or by one canonical bad birth;
\item later colours can only enlarge or merge saturated components
and cannot change a historical owner;
\item all connected colour outputs converge absolutely in the
\(e^{\mu|Y|}\) owner norm;
\item
\[
                         \eta_a\le K_a u_7(\beta)
 \quad(0\le a\le\chi_*).                                  \tag{8.29}
\]
\end{enumerate}
Consequently the schedule hypothesis of
Theorem~\ref{thm:n1v7-kernel-gate} is proved for this triangular
disintegration, and its all-colour connected gate is unconditional
for all sufficiently large \(\beta\).
\end{theorem}

\begin{proof}
We use induction on \(a\).  The assertion is trivial at \(a=0\).
Assume
\(\eta_{a-1}\le K_{a-1}u_7\).  Since
\(K_{a-1}\le K_{\chi_*}\), the second entry of (8.28) gives
\[
 A_\mu v_a
 \le
 A_\mu c_{\rm cmp}(K_{\chi_*}+1)u_7
 \le\frac12.                                               \tag{8.30}
\]
Lemma~\ref{lem:n1v8-recursion} and (8.20) therefore imply
\[
\begin{split}
 \eta_a
 &\le\eta_{a-1}+2e^\mu c_{\rm cmp}
       (\eta_{a-1}+u_7)\\
 &=D_\mu\eta_{a-1}+B_\mu u_7\\
 &\le\{D_\mu K_{a-1}+B_\mu\}u_7
   =(D_\mu^a-1)u_7=K_a u_7.                               \tag{8.31}
\end{split}
\]
The first entry of (8.28) now gives \(\eta_a\le\eta_*\), so
Theorem~\ref{thm:n1v8-morse-stability} and
Corollary~\ref{cor:n1v8-kernel-survives} apply at the next colour.
This closes the analytic induction.

The historical conclusions are
Lemma~\ref{lem:n1v8-history}.  A regular birth pays
\(q_{\rm ker}\); a bad birth pays its disjoint set of newly assigned
blocks by Lemma~\ref{lem:n1v8-bad-distribution}.  Thus mixed histories
pay \(u_7\) once per assigned block, with no second determinant or
collar charge.  Lemma~\ref{lem:n1v8-compiler} and (8.30) give absolute
connected convergence at every stage.  Lemma~\ref{lem:n1v8-fubini}
shows that the stages compose to the exact integral rather than to an
approximate update.  These statements are precisely the schedule
requirements in Theorem~\ref{thm:n1v7-kernel-gate}.

Finally, (7.25) says \(u_7(\beta)\to0\), while every constant in
(8.28) is fixed.  Hence (8.28) holds for all sufficiently large
\(\beta\).
\end{proof}

\begin{corollary}[Uniform all-colour connected activity]
\label{cor:n1v8-all-colour-gate}
Under (8.28), the complete triangular colour disintegration has a
connected owner expansion satisfying the V7 geometric gate
\[
 e^{1+\mu}(\Delta_*+1)u_7(\beta)<1                         \tag{8.32}
\]
after increasing the lower threshold for \(\beta\), and all four
source derivatives may be taken after the normally convergent
connected sum.
\end{corollary}

\begin{proof}
Theorem~\ref{thm:n1v8-all-colour} supplies the previously conditional
owner schedule.  Since \(u_7(\beta)\to0\), (8.32) can be imposed
together with (8.28).  The smaller source polydisc in
Lemma~\ref{lem:n1v8-embeddings} is fixed, so normal convergence and
the four-variable Cauchy theorem justify differentiation term by
term.
\end{proof}

\begin{firewall}[Finite-scale meaning of the closure]
\label{fire:n1v8-finite-scale}
Theorem~\ref{thm:n1v8-all-colour} closes the all-colour owner and
saturation card for one fixed-factor, finite-scale triangular
integration.  It does not assert that a repeated Gibbs sampler
contracts, that the resulting coarse action returns to the same
calibrated family at every renormalization scale, or that the
constants in (8.28) remain uniform along an infinite RG trajectory.
Those are separate statements.
\end{firewall}

\section{The exact next bridge: physical-sector contraction}
\label{sec:n1v8-bridge}

The schedule proof identifies the next genuine numerical bottleneck.
Let \({\mathscr X}_{\rm ph}\) be an influence space after quotienting
constants and functions of retained slow or toron variables.  A
reference coarse sweep \(G\) and its exact counterpart \(P\) must act
on this same quotient.  The following elementary resolvent card fixes
the target constants.

\begin{lemma}[Contraction transfer and invariant-state comparison]
\label{lem:n1v8-resolvent}
Suppose \(G1=P1=1\) and, on \({\mathscr X}_{\rm ph}\),
\[
 \|G\|\le\rho<1,\qquad \|P-G\|\le\delta,\qquad
 \rho+\delta<1.                                            \tag{8.33}
\]
Then \(P\) is a strict contraction on
\({\mathscr X}_{\rm ph}\), both \(I-G\) and \(I-P\) have norm
convergent Neumann inverses there, and
\[
 \|(I-P)^{-1}-(I-G)^{-1}\|
 \le
 \frac{\delta}{(1-\rho-\delta)(1-\rho)}.                   \tag{8.34}
\]
If the corresponding dual invariant states \(\nu_P,\nu_G\) are
normalized and the dual norm of a state is at most one, then
\[
 \|\nu_P-\nu_G\|_{{\mathscr X}_{\rm ph}^*}
 \le\frac{\delta}{1-\rho}.                                 \tag{8.35}
\]
\end{lemma}

\begin{proof}
The triangle inequality gives
\(\|P\|\le\rho+\delta<1\) on the quotient.  Hence both Neumann series
converge.  The resolvent identity
\[
 (I-P)^{-1}-(I-G)^{-1}
 =(I-P)^{-1}(P-G)(I-G)^{-1}
\]
proves (8.34).  Invariance gives, on the quotient,
\[
 \nu_P-\nu_G
 =(\nu_P-\nu_G)G+\nu_P(P-G).
\]
Right multiplication by \((I-G)^{-1}\) and
\(\|\nu_P\|\le1\) prove (8.35).
\end{proof}

\begin{openproblem}[Numerical sweep card]
\label{open:n1v8-sweep-card}
V8 does not yet prove either inequality in (8.33).  The connected
kernel expansion gives a candidate
\(\delta\le C_{\rm swp}u_7(\beta)\), but the conversion constant
\(C_{\rm swp}\) must be derived in one explicit local influence norm.
More importantly, the reduced shifted-Gaussian reference sweep needs
a volume-independent \(\rho<1\) after slow and toron directions are
removed.  Treating those retained directions as contractive would be
false.
\end{openproblem}

\begin{decision}[V9 target]
\label{dec:n1v8-next}
Construct the physical influence seminorm from the archived reduced
Hessian and Schur-complement decomposition.  Compute the block
influence matrix of the shifted-Gaussian reference disintegration and
test its weighted row sum.  If the raw row sum reaches one, go
upstream: enlarge the retained slow sector until the complementary
Schur operator has a strict gap.  Then prove the connected-kernel
conversion
\(\|P-G\|\le C_{\rm swp}u_7(\beta)\) and apply
Lemma~\ref{lem:n1v8-resolvent}.
\end{decision}

\begin{assessment}[Progress after compact V8]
\label{ass:n1v8-status}
The all-colour hypothesis left open in V7 is now discharged for the
exact triangular pushforward.  The proof combines a monotone
historical saturation rule, one-time owner payment, a quantitative
Morse stability reserve, and a finite connected-output recursion.
The remaining obstruction has moved downstream to a strict
physical-sector contraction of the reference coarse sweep and its
stability under the \(O(u_7)\) exact-kernel defect.
\end{assessment}

\begin{firewall}[Claim boundary after compact V8]
\label{fire:n1v8-claim}
V8 proves a one-scale all-colour disintegration and its connected
kernel gate at sufficiently large \(\beta\).  It does not yet prove
the reference influence contraction, exact invariant coarse-kernel
identification, a scale-uniform RG invariant domain, regulator
removal, Osterwalder--Schrader reconstruction, or the Yang--Mills
mass gap.
\end{firewall}

\subsection{Parent record}

This V8 note was formed by appending the preceding material to the
validated compact V7, whose record was
\[
\begin{array}{c|c}
\text{lines}&3401\\
\text{bytes}&115294\\
\text{SHA--256}&
\texttt{09B911B8C7EBCD784059FEDB8743150C38459CEBB654DC3A42339294106206E1}.
\end{array}
\]

% =============================================================================
% END APPENDED COMPACT-SERIES V8 MATERIAL
% =============================================================================

% =============================================================================
% BEGIN APPENDED COMPACT-SERIES V9 MATERIAL
% =============================================================================

\chapter{V9: an exact measured-coordinate slice for the Jacobian
marginal}
\label{ch:n1v9}

\section{Upstream correction to the V8 acquisition order}
\label{sec:n1v9-correction}

\begin{correction}[Do not rebuild the obstructed influence norm]
\label{cor:n1v9-influence-route}
The V8 decision proposed constructing an exponentially weighted
physical influence norm directly from the reduced Schur operator.
That construction would duplicate an archived route whose endpoint
is already known to fail: total-variation influence of a coupled
unbounded Gaussian saturates at one, and the normalized physical
curvature multiplier has a direction-dependent Hodge limit at zero
momentum.  The archive already replaces exponential localization by
dyadic annular packet localization and proves the corresponding
principal contraction and packet product estimates.  Those results
remain live.  V9 therefore moves to the earlier unresolved quotient
card rather than redoing them.
\end{correction}

The open issue is the marginal Hessian of the Haar divergence of the
raw pushed conditional current.  The archive proves that its value is
a coordinate coefficient and can be shifted by a cubic compact
vector field, but it stops short of turning that observation into an
exact analytic section compatible with a pushforward.  We now do so.

\section{Exact action of measured coordinate flows}
\label{sec:n1v9-flows}

Let \(X\) be a finite product of \(SU(N)\) link fibres with product
Haar measure \(dH\).  Let \(S\) be a real analytic action and let
\({\cal W}\) be a real analytic gauge-equivariant vector field on
\(X\).  Compactness makes its real flow \(\varphi_t\) complete.
For sufficiently small complex \(t\), the flow is analytic in every
fixed complex chart used by the owner norm.

\begin{definition}[Measured coordinate action]
\label{def:n1v9-coordinate-action}
Define \(S^{[t]}\) exactly by
\[
 e^{-S^{[t]}}\,dH
 =
 (\varphi_{-t})_*(e^{-S}\,dH).                              \tag{9.1}
\]
Equivalently,
\[
 S^{[t]}(y)
 =
 S(\varphi_t(y))
 -\log\operatorname{Jac}_H\varphi_t(y).                    \tag{9.2}
\]
\end{definition}

\begin{lemma}[The infinitesimal owner includes the Jacobian]
\label{lem:n1v9-infinitesimal}
The derivative of (9.2) at \(t=0\) is
\[
 \left.\frac d{dt}\right|_{t=0}S^{[t]}
 =
 DS[{\cal W}]-\operatorname{div}_H{\cal W}
 =:{\cal O}_{\cal W}.                                      \tag{9.3}
\]
For every integrable observable \(F\),
\[
 \int_XF\,e^{-S^{[t]}}dH
 =
 \int_XF\circ\varphi_{-t}\,e^{-S}dH.                        \tag{9.4}
\]
\end{lemma}

\begin{proof}
The standard Jacobian equation for a flow is
\[
 \left.\frac d{dt}\right|_{0}
 \log\operatorname{Jac}_H\varphi_t
 =\operatorname{div}_H{\cal W}.
\]
Differentiate (9.2) to obtain (9.3).  Equation (9.4) is the defining
pushforward identity (9.1).
\end{proof}

\begin{corollary}[Measured redundant directions are exact coordinate
orbits]
\label{cor:n1v9-exact-orbit}
The infinitesimal owner \({\cal O}_{\cal W}\) is tangent to an exact
family of probability measures related by a change of variables.
It is therefore redundant only when observables are transported by
(9.4); deleting the divergence while keeping observables fixed is
not the same operation.
\end{corollary}

\begin{proof}
Normalization and equality of all transported expectations follow
from (9.4).  Formula (9.3) shows that the action and Jacobian terms
are the two inseparable parts of the tangent.
\end{proof}

\begin{lemma}[Compatibility with an exact block pushforward]
\label{lem:n1v9-block-compatibility}
Let \({\cal B}:X_{\rm f}\to X_{\rm c}\) be a measurable block map and
let \(\nu_\varepsilon\) be a differentiable family of fine measures.
If
\(\rho_\varepsilon dH_{\rm c}={\cal B}_*\nu_\varepsilon\), then for
any retained vector field \({\cal W}\)
\[
 \widetilde\rho_\varepsilon dH_{\rm c}
 :=
 (\varphi_{-\varepsilon\lambda})_*
 (\rho_\varepsilon dH_{\rm c})                              \tag{9.5}
\]
is the pushforward under the postcomposed block coordinate
\(\varphi_{-\varepsilon\lambda}\circ{\cal B}\).
If the original logarithmic response is
\({\cal O}_{\overline V}\), the response of (9.5) is
\[
                         {\cal O}_{\overline V+\lambda{\cal W}}.
                                                               \tag{9.6}
\]
\end{lemma}

\begin{proof}
Functoriality of pushforwards gives
\[
 (\varphi_{-\varepsilon\lambda})_*{\cal B}_*\nu_\varepsilon
 =
 (\varphi_{-\varepsilon\lambda}\circ{\cal B})_*
 \nu_\varepsilon.
\]
Apply Lemma~\ref{lem:n1v9-infinitesimal} to the second pushforward.
Both the differential of the action and Haar divergence are linear
in the infinitesimal velocity, which proves (9.6).
\end{proof}

\section{The canonical compact compensation direction}
\label{sec:n1v9-compensation}

On a retained lattice put
\[
 r_p(U)=1-\frac1N\operatorname{ReTr}U_p,\qquad
 V_{\rm W}=\sum_pr_p,\qquad
 \Phi=\sum_pr_p^2,\qquad {\cal W}_\Phi=\nabla\Phi.           \tag{9.7}
\]
The archive's exact Haar--Casimir calculation gives
\[
 \left(D^2\operatorname{div}_H{\cal W}_\Phi\right)_I
 =
 8\left(C_F+\frac2N\right)D^2V_{\rm W}(I),
 \qquad C_F=\frac{N^2-1}{N}.                                \tag{9.8}
\]
The vector field \({\cal W}_\Phi\) is cubic in exponential link
coordinates and hence has zero vacuum linearization.

Let \({\mathfrak m}\) be the fixed Wilson marginal extractor and put
\[
 K_{\rm W}=D^2V_{\rm W}(I),\qquad
 \omega={\mathfrak m}(K_{\rm W})\ne0,\qquad
 c_H=8\left(C_F+\frac2N\right)\omega.                        \tag{9.9}
\]
The normalization of \({\mathfrak m}\) may of course be chosen so
that \(\omega=1\).

\begin{lemma}[Stationarity removes the action-gradient debit]
\label{lem:n1v9-gradient-debit}
Let \(S\) be any analytic gauge-invariant action stationary at the
flat vacuum.  Then
\[
 D^2\bigl(DS[{\cal W}_\Phi]\bigr)(I)=0.                     \tag{9.10}
\]
Consequently, for an analytic velocity \(U\),
\[
 \zeta_{\rm Jac}(U+\lambda{\cal W}_\Phi)
 =
 \zeta_{\rm Jac}(U)-c_H\lambda,                             \tag{9.11}
\]
while the marginal Hessian of the action-gradient part is unchanged.
\end{lemma}

\begin{proof}
In exponential coordinates \(a\), stationarity gives
\(DS(a)=O(\|a\|)\), while
\({\cal W}_\Phi(a)=O(\|a\|^3)\).  Their contraction is
\(O(\|a\|^4)\), proving (9.10).  The Haar divergence is linear in the
vector field.  Insert (9.8), apply the minus sign in
\(\zeta_{\rm Jac}={\mathfrak m}
(-D^2\operatorname{div}_HU)\), and use (9.9) to obtain (9.11).
\end{proof}

\section{Analytic zero-Jacobian section}
\label{sec:n1v9-section}

For an ownerwise analytic velocity \(U\), write
\(\|U\|_{(3),\mu,\rho}\) for the V8 owner norm including background
jets through order three.  The third-jet divergence formula from the
archive and the finite marginal extractor imply a fixed bound
\[
 |\zeta_{\rm Jac}(U)|
 \le C_{\rm J}\|U\|_{(3),\mu,\rho}.                         \tag{9.12}
\]
The constant is a finite contraction of the first and third
background jets with the Haar connection at the identity.  The V8
four-source connected gate supplies precisely the four marked
observables needed to make the right side volume independent.

\begin{definition}[Compensated pushed current]
\label{def:n1v9-compensated-current}
For every current owner set
\[
 \lambda(U):=\frac{\zeta_{\rm Jac}(U)}{c_H},\qquad
 U^\circ:=U+\lambda(U){\cal W}_\Phi.                        \tag{9.13}
\]
\end{definition}

\begin{theorem}[Exact analytic measured-coordinate slice]
\label{thm:n1v9-coordinate-slice}
On the small-activity domain of
Theorem~\ref{thm:n1v8-all-colour}, (9.13) has the following
properties.
\begin{enumerate}[label=\textup{(\roman*)},leftmargin=3.5em]
\item \(U^\circ\) has the same vacuum linearization as \(U\).
\item
\[
                         \zeta_{\rm Jac}(U^\circ)=0.        \tag{9.14}
\]
\item The map \(U\mapsto U^\circ\) is analytic, local,
gauge-equivariant, reflection and hypercubic equivariant, and
idempotent.
\item There is a volume-independent bound
\[
 \|U^\circ\|_{(3),\mu,\rho}
 \le
 \left(
 1+\frac{C_{\rm J}\|{\cal W}_\Phi\|_{(3),\mu,\rho}}{|c_H|}
 \right)
 \|U\|_{(3),\mu,\rho}.                                      \tag{9.15}
\]
\item If \(U\) is the raw coarse current produced by an exact
pushforward, \(U^\circ\) is the current produced after the exact
postcomposition (9.5).  Hence it represents the same measured
field-redefinition class.
\end{enumerate}
The section is the unique affine correction of \(U\) in the
\({\cal W}_\Phi\) direction satisfying (9.14).
\end{theorem}

\begin{proof}
The zero linearization of \({\cal W}_\Phi\) proves (i).
Equations (9.11) and (9.13) give
\[
 \zeta_{\rm Jac}(U^\circ)
 =\zeta_{\rm Jac}(U)-c_H\frac{\zeta_{\rm Jac}(U)}{c_H}=0.
\]
The conditional current and its first three jets are normally
convergent analytic owner sums by the V8 gate.  The marginal
extractor is a finite linear map, so \(\lambda(U)\) is analytic.
The vector field in (9.7) has fixed finite range and all the listed
symmetries.  Thus the correction has the same properties.

If \(\zeta_{\rm Jac}(U)=0\), then \(\lambda(U)=0\); applying the map
again changes nothing, which proves idempotence.  Equation (9.12)
and the triangle inequality give (9.15).  Item (v) is
Lemma~\ref{lem:n1v9-block-compatibility}.  Finally, (9.11) has
nonzero slope \(-c_H\), so its unique zero on the affine line
\(U+\mathbb R{\cal W}_\Phi\) is (9.13).
\end{proof}

\begin{corollary}[The raw coefficient is not a physical prerequisite]
\label{cor:n1v9-raw-not-needed}
In the quotient by exact local measured coordinate changes, the raw
number \(\zeta_{\rm Jac}(U)\) need not be evaluated.  The canonical
representative \(U^\circ\) has zero Jacobian marginal, no changed
linearization, and no action-gradient \(F^2\) debit.  Any fixed raw
coordinate calculation may still report the original number, but it
cannot alter a quotient observable or a response computed with the
transport rule (9.4).
\end{corollary}

\begin{proof}
Theorem~\ref{thm:n1v9-coordinate-slice}(i)--(ii) and
Lemma~\ref{lem:n1v9-gradient-debit} prove the local statements.
Items (iv)--(v) and (9.4) show that the replacement is an exact
bounded coordinate postcomposition, not deletion of a term.
\end{proof}

\begin{firewall}[Scope of the quotient]
\label{fire:n1v9-section-scope}
The coordinate slice is used inside the RG comparison and all
observables are transported back by (9.4).  V9 does not replace the
original reflection-positive lattice measure by an unrelated
measure.  Nor does it assert that a Jacobian term vanishes in an
arbitrarily fixed raw coordinate system.
\end{firewall}

\section{Exact comparison with the isotropic dimension-six
completion}
\label{sec:n1v9-o1}

Let \(\widehat{\cal O}_1\) denote the fixed positive normalization of
the archived compact completion
\[
 {\cal O}_1(U)
 =
 \frac12\sum_{x,\mu<\nu}\sum_{\rho=1}^4
 \|{\cal D}_\rho P_{\mu\nu}(x)\|_{\rm HS}^2                \tag{9.16}
\]
for which
\[
                         D^2\widehat{\cal O}_1(I)
 ={\cal P}^2.                                               \tag{9.17}
\]
Here \({\cal P}=d_1^*d_1\) is the flat Wilson precision.
For
\[
 {\cal V}_*
 =
 \frac12\nabla V_{\rm W}
 +\lambda_*\nabla\Phi,\qquad
 \lambda_*=\frac{N^2-1}{4(N^2+1)},                          \tag{9.18}
\]
define the vacuum-subtracted measured owner
\[
\begin{split}
 {\cal O}_{\rm R}^{(\beta)}
 :={}&
 \beta^{-1}\left\{
 D(\beta V_{\rm W})[{\cal V}_*]
 -\operatorname{div}_H{\cal V}_*
 \right\}\\
 &-\left.
 \beta^{-1}\left\{
 D(\beta V_{\rm W})[{\cal V}_*]
 -\operatorname{div}_H{\cal V}_*
 \right\}\right|_{I}.                                      \tag{9.19}
\end{split}
\]
It is the measured redundant owner generated by
\(\beta^{-1}{\cal V}_*\).

\begin{theorem}[Two-jet exact completion comparison]
\label{thm:n1v9-completion-comparison}
The local analytic functional
\[
 {\cal E}_{\rm comp}^{(\beta)}
 :=
 \widehat{\cal O}_1-{\cal O}_{\rm R}^{(\beta)}              \tag{9.20}
\]
has vanishing vacuum two-jet:
\[
 {\cal E}_{\rm comp}^{(\beta)}(I)=0,\qquad
 D{\cal E}_{\rm comp}^{(\beta)}(I)=0,\qquad
 D^2{\cal E}_{\rm comp}^{(\beta)}(I)=0.                    \tag{9.21}
\]
In the measured field-redefinition quotient,
\[
                         [\widehat{\cal O}_1]
 = [{\cal E}_{\rm comp}^{(\beta)}].                        \tag{9.22}
\]
\end{theorem}

\begin{proof}
The vacuum values agree by the subtraction in (9.19) and
\(\widehat{\cal O}_1(I)=0\).  Gauge, reflection, and translation
symmetry make both first derivatives vanish at the flat vacuum.
The archived Haar--Casimir computation at the value (9.18) gives
\[
 D^2\left\{
 D(\beta V_{\rm W})[{\cal V}_*]
 -\operatorname{div}_H{\cal V}_*
 \right\}(I)
 =\beta{\cal P}^2.
\]
After division by \(\beta\), this equals (9.17), proving the third
identity in (9.21).  The second term in (9.20) is an exact measured
redundant owner by (9.19) and
Corollary~\ref{cor:n1v9-exact-orbit}; it vanishes in the quotient,
which proves (9.22).
\end{proof}

\begin{corollary}[Running-window size of the completion defect]
\label{cor:n1v9-completion-size}
Choose a fixed logarithm ball on which the third analytic derivative
of (9.20) is bounded by \(M_3\), uniformly for
\(\beta\ge\beta_0>0\).  In weak-field variables \(U=e^{ga}\),
\(g=\beta^{-1/2}\),
\[
 |{\cal E}_{\rm comp}^{(\beta)}(e^{ga})|
 \le\frac{M_3}{6}g^3\|a\|_{\rm loc}^3.                     \tag{9.23}
\]
On the inherited running good window
\[
 \|a\|_{\rm loc}
 \le\sqrt{720}\,\pi\,\beta^{1/10},
\]
one obtains
\[
\boxed{
 |{\cal E}_{\rm comp}^{(\beta)}(e^{ga})|
 \le
 C_{\rm comp}\beta^{-6/5},\qquad
 C_{\rm comp}:=
 \frac{M_3}{6}(\sqrt{720}\,\pi)^3.}                         \tag{9.24}
\]
The same estimate holds in the restricted analytic owner norm after
one fixed Cauchy-radius loss.
\end{corollary}

\begin{proof}
Taylor's integral formula and (9.21) give
\[
 {\cal E}_{\rm comp}^{(\beta)}(e^{ga})
 =
 \frac12\int_0^1(1-t)^2
 D^3{\cal E}_{\rm comp}^{(\beta)}(e^{tga})
 [ga,ga,ga]\,dt,
\]
which proves (9.23) because
\(\frac12\int_0^1(1-t)^2dt=1/6\).
The third derivatives of the first term in (9.20) are fixed.
Those of the action-gradient part of (9.19) are also fixed after the
factor \(\beta^{-1}\) cancels the Wilson coefficient \(\beta\);
the remaining divergence derivatives carry the smaller factor
\(\beta^{-1}\).  Thus one \(M_3\) works for
\(\beta\ge\beta_0\).  Inserting the running-window bound gives
\[
 g^3\|a\|_{\rm loc}^3
 \le(\sqrt{720}\,\pi)^3
 \beta^{-3/2+3/10}
 =(\sqrt{720}\,\pi)^3\beta^{-6/5},
\]
which is (9.24).  Cauchy's theorem on nested fixed analytic balls
gives the norm version.
\end{proof}

\begin{corollary}[A stable primitive, not a new marginal row]
\label{cor:n1v9-completion-primitive}
For every fixed \(|\tau|\le\tau_*\), the good-window exponential
factor satisfies
\[
 \left|
 e^{-\tau{\cal E}_{\rm comp}^{(\beta)}}-1
 \right|
 \le
 \exp\{\tau_*C_{\rm comp}\beta^{-6/5}\}-1
 =O(\beta^{-6/5}).                                         \tag{9.25}
\]
Its complement is assigned to the existing collar/far-bad owner.
Thus comparison of \({\cal O}_1\) with the compensated measured
owner creates no \(F^2\) marginal and fits the V8 connected compiler
as a stable primitive.
\end{corollary}

\begin{proof}
Use \(|e^{-z}-1|\le e^{|z|}-1\) and (9.24).  Equation (9.21)
places the local factor in the stable two-jet kernel, while the V7
saturation rule owns the complement.
\end{proof}

\begin{assessment}[What V9 removes from the dependency chain]
\label{ass:n1v9-progress}
V9 proves that the unevaluated raw
\(\zeta_{\rm Jac}\) is not an obstruction to the physical quotient:
there is a unique, bounded, analytic, symmetry-compatible
postcomposition of the exact block coordinate that sets it to zero
without changing the current linearization.  It also completes the
structural comparison with \({\cal O}_1\): after subtracting the
measured redundant representative, the completion defect has zero
two-jet and size \(O(\beta^{-6/5})\) on the running window.
\end{assessment}

\begin{decision}[Mandatory V10 audit and next target]
\label{dec:n1v9-next}
V10 must inspect the then-current Standard-Model--Einstein and
Navier--Stokes notes.  The audit should look specifically for a
cofinal symmetry-restoration estimate, a quotient-compatible
anisotropy recursion, or a terminal bridge usable for the remaining
hypercubic owner
\[
                         {\cal O}_{\rm A}
 ={\cal O}_2-\frac23{\cal O}_1.                             \tag{9.26}
\]
The Yang--Mills continuation should then prove an actual contraction
or an explicit obstruction for this anisotropic row in the archived
annular packet space.
\end{decision}

\begin{firewall}[Claim boundary after compact V9]
\label{fire:n1v9-claim}
V9 removes the raw Jacobian marginal from the quotient dependency and
turns the isotropic completion difference into a small stable owner.
It does not prove the hypercubic anisotropy contraction, the complete
nonlinear packet-space invariant ball, the universal marginal
coefficient, a scale-uniform RG trajectory, regulator removal,
Osterwalder--Schrader reconstruction, or the Yang--Mills mass gap.
\end{firewall}

\subsection{Parent record}

This V9 note was formed by appending the preceding material to the
validated compact V8, whose record was
\[
\begin{array}{c|c}
\text{lines}&4128\\
\text{bytes}&142140\\
\text{SHA--256}&
\texttt{90D1A84C251F1A51F2E4E22B0C208B149131275A64911D882406D592BD7571F0}.
\end{array}
\]

% =============================================================================
% END APPENDED COMPACT-SERIES V9 MATERIAL
% =============================================================================

% =============================================================================
% BEGIN APPENDED COMPACT-SERIES V10 MATERIAL
% =============================================================================

\chapter{V10: companion audit and strict contraction of the
hypercubic row}
\label{ch:n1v10}

\section{Mandatory five-version companion audit}
\label{sec:n1v10-audit}

\begin{audit}[Files inspected]
\label{aud:n1v10-files}
The live companion snapshots inspected before choosing the V10
argument were
\[
\begin{array}{c|r|r|l}
\text{file}&\text{lines}&\text{bytes}&\text{SHA--256}\\ \hline
\texttt{Renewal\_SM\_Einstein\_Continuum\_Research\_Notes\_V87.tex}
&30681&1117435&
\texttt{2EA56F3C6A620C7C51157C6BFDF175B2AB5942A8ADC9295B792DB9CBB2FE62CE}
\\
\texttt{Renewal\_NS\_Stability\_Research\_Notes\_V26.tex}
&5319&278283&
\texttt{82C887F8C2C69E4F28FAD7C3DC8DE5B9099CB5A4BF431EBA1FFF3E91935978AD}.
\end{array}                                                     \tag{10.1}
\]
The SM file advanced during the audit, so (10.1) records the later
V87 file actually read.  The NS file is unchanged.
\end{audit}

SM V86 proves abstract finite-colour telescoping and invariant-kernel
resolvent estimates.  SM V87 then computes the Wasserstein influence
matrix of a block Gaussian conditional:
\[
 c_{ij}=\|A_{ii}^{-1}A_{ij}\|,\qquad
 \kappa_w=\sup_iw_i^{-1}\sum_{j\ne i}c_{ij}w_j.             \tag{10.2}
\]
Its proof gives contraction if \(\kappa_w<1\), but its massless
nearest-neighbour example gives
\[
 \kappa_h=\frac{2D}{2D+m^2h^2},                             \tag{10.3}
\]
which equals one at \(m=0\) and approaches one at fixed physical
mass as \(h\downarrow0\).  This confirms the V9 route correction:
a raw single-site influence gap is not uniform on an unprojected
massless continuum sector.

\begin{audit}[Transfer decision]
\label{aud:n1v10-transfer}
The SM precision proof transfers as a useful diagnostic, not as a
new Yang--Mills contraction.  Gauge, toron, and infrared variables
must remain outside the contracted sector.  The older Yang--Mills
annular packet construction already implements that separation and
is retained.  SM V87 supplies no hypercubic-response coefficient.
The NS rank-one Oseen estimates concern a deterministic
three-dimensional parabolic kernel and supply neither a gauge Ward
projection nor a four-dimensional anisotropy row.  Accordingly V10
works only with the finite local hypercubic coordinate after the V9
measured quotient.
\end{audit}

\section{The calibrated anisotropy coordinate}
\label{sec:n1v10-coordinate}

Recall the two normalized dimension-six completions
\(\widehat{\cal O}_1,\widehat{\cal O}_2\).  Their quadratic Ward
tensors span the two-dimensional local hypercubic space.  Define
\[
 {\cal O}_{\rm A}
 :=
 \widehat{\cal O}_2-\frac23\widehat{\cal O}_1.              \tag{10.4}
\]
Its rotational average vanishes, whereas its flat quadratic symbol
is fourth order in momentum.  The isotropic redundant direction is
represented by \({\cal O}_{\rm R}^{(\beta)}\) from (9.19) and is
removed by the V9 measured-coordinate section.

\begin{definition}[Anisotropy extractor]
\label{def:n1v10-extractor}
On the finite local Ward-jet space choose the linear functional
\({\mathfrak a}\) by
\[
 {\mathfrak a}({\cal O}_{\rm A})=1,\qquad
 {\mathfrak a}({\cal O}_{\rm R}^{(\beta)})=0,\qquad
 {\mathfrak a}(V_{\rm W})=0,                                \tag{10.5}
\]
and let it vanish on the chosen stable complement.  Its operator norm
on the restricted analytic local norm is denoted \(C_{\mathfrak a}\).
\end{definition}

The functional exists and is bounded because the displayed rows are
linearly independent in a fixed finite-dimensional jet space.  It is
important that (10.5) is imposed after the measured redundant
representative has been fixed; otherwise a raw Jacobian coordinate
could be mistaken for anisotropy.

Let \(b\ge2\) be the fixed blocking factor and let
\({\cal R}_\beta^\circ\) denote one exact triangular pushforward,
followed by the V9 measured-coordinate section and the declared
constant, marginal, redundant, and stable projections.  Let
\({\cal L}_0\) be its calibrated Gaussian principal derivative.
The inherited canonical dilation theorem gives
\[
 \left|
 {\mathfrak a}\bigl({\cal L}_0{\cal O}_{\rm A}\bigr)
 \right|
 \le b^{-2}.                                                 \tag{10.6}
\]
This is the engineering contraction of a dimension-six local
curvature jet in four dimensions.  It does not involve an inverse on
the retained massless sector.

\section{One distinguished anisotropy mark}
\label{sec:n1v10-mark}

The V8 compiler was stated for unmarked activities.  We need its
one-source derivative.  Define the V9 completion activity
\[
 q_{\rm comp}(\beta)
 :=
 \exp\{\tau_*C_{\rm comp}\beta^{-6/5}\}-1                   \tag{10.7}
\]
and the complete nonreference strength
\[
 w_{10}(\beta)
 :=
 (K_{\chi_*}+1)u_7(\beta)+q_{\rm comp}(\beta).              \tag{10.8}
\]
The factor \(K_{\chi_*}+1\) includes all generated increments in
(8.29) and one fresh defect.  In particular,
\[
                         w_{10}(\beta)\longrightarrow0.     \tag{10.9}
\]

Let \(c_{10}\ge1\) contain the fixed owner, jet, packet-to-local, and
source-radius conversions for one insertion of
\({\cal O}_{\rm A}\), and put
\[
 x_{10}(\beta)=A_\mu c_{10}w_{10}(\beta).                  \tag{10.10}
\]
Let \(M_{\rm A}\) be the strong analytic norm of one normalized
\({\cal O}_{\rm A}\) insertion on the larger source polydisc.

\begin{lemma}[One-mark connected defect bound]
\label{lem:n1v10-one-mark}
If \(x_{10}(\beta)<1\), the sum of all connected outputs containing
one distinguished anisotropy insertion and at least one
nonreference primitive obeys
\[
\boxed{
 \|{\cal E}_{\beta,{\rm A}}\|_{\mu,\rho}
 \le
 \varepsilon_{\rm A}(\beta)
 :=
 \frac{e^\mu M_{\rm A}c_{10}w_{10}(\beta)}
      {(1-x_{10}(\beta))^2}.}                              \tag{10.11}
\]
The bound is uniform in the finite volume, in the colour stage, and
on the smaller anisotropy-source disc.
\end{lemma}

\begin{proof}
Subtract the Gaussian principal response before expanding.  Every
remaining connected word then contains the distinguished source mark
and at least one of: a regular signed kernel, a saturated witness, a
generated hit owner, or the completion defect (9.25).  Their total
incident strength is bounded by (10.8), and all deterministic norm
conversions are included in \(c_{10}\).

For a word with \(n\ge1\) nonreference primitives, root at the source
support and use the connected-set exploration of
Lemma~\ref{lem:n1v7-connected-count}.  There are at most \(n+1\)
places at which the distinguished mark can enter the rooted
exploration.  After the first primitive, every further primitive
costs at most \(x_{10}\).  Hence the marked sum is bounded by
\[
 e^\mu M_{\rm A}c_{10}w_{10}
 \sum_{n\ge0}(n+1)x_{10}^n
 =
 \frac{e^\mu M_{\rm A}c_{10}w_{10}}
      {(1-x_{10})^2}.
\]
The V8 normal convergence permits the source derivative to pass
through the connected sum.  The fixed source support and owner
incidence make every constant volume independent.
\end{proof}

\begin{corollary}[Marked remainder tends to zero]
\label{cor:n1v10-mark-small}
There is \(\beta_{10}<\infty\) such that, for
\(\beta\ge\beta_{10}\),
\[
 x_{10}(\beta)\le\frac12,\qquad
 \varepsilon_{\rm A}(\beta)
 \le4e^\mu M_{\rm A}c_{10}w_{10}(\beta)
 \longrightarrow0.                                         \tag{10.12}
\]
\end{corollary}

\begin{proof}
Use (10.9) in (10.10)--(10.11).
\end{proof}

\section{Strict one-step contraction}
\label{sec:n1v10-contraction}

Fix all calibrated marginal, topological, slow, and toron
coordinates in their declared polydiscs.  Let
\[
 F_{\beta,r}(\alpha)
 :=
 {\mathfrak a}\!\left(
 {\cal R}_\beta^\circ(S_r+\alpha{\cal O}_{\rm A})
 -{\cal R}_\beta^\circ(S_r)
 \right),                                                   \tag{10.13}
\]
where \(r\) denotes the retained parameters.  The subtraction makes
\(F_{\beta,r}(0)=0\) and removes any affine anisotropy forcing from
the Lipschitz calculation.

\begin{theorem}[Uniform hypercubic-row contraction]
\label{thm:n1v10-anisotropy-contraction}
On a fixed smaller source interval
\(|\alpha|\le\alpha_*\), and for all sufficiently large \(\beta\),
\[
 |F_{\beta,r}(\alpha)-F_{\beta,r}(\alpha')|
 \le q_{\rm A}(\beta)|\alpha-\alpha'|,                      \tag{10.14}
\]
uniformly in the volume and retained parameters, where
\[
 q_{\rm A}(\beta)
 :=
 b^{-2}+C_{\mathfrak a}\varepsilon_{\rm A}(\beta).          \tag{10.15}
\]
In particular, after increasing \(\beta_{10}\),
\[
\boxed{
                         q_{\rm A}(\beta)\le\frac12
 \qquad(b\ge2).}                                            \tag{10.16}
\]
\end{theorem}

\begin{proof}
Differentiate (10.13) with respect to \(\alpha\).
The Gaussian principal term is bounded by \(b^{-2}\) through (10.6).
After subtracting that term, the source derivative is precisely a
connected output with one distinguished anisotropy mark and at least
one nonreference primitive.  Lemma~\ref{lem:n1v10-one-mark} and the
extractor norm give
\[
                         |F_{\beta,r}'(\alpha)|
 \le b^{-2}+C_{\mathfrak a}\varepsilon_{\rm A}(\beta)
\]
uniformly on the smaller source interval.  Integrate along the real
segment from \(\alpha'\) to \(\alpha\) to prove (10.14).
For \(b\ge2\), \(b^{-2}\le1/4\).  By
Corollary~\ref{cor:n1v10-mark-small}, choose the threshold so that
\(C_{\mathfrak a}\varepsilon_{\rm A}\le1/4\); this proves (10.16).
\end{proof}

\begin{remark}[What is and is not contracted]
\label{rem:n1v10-sector}
Theorem~\ref{thm:n1v10-anisotropy-contraction} is a scalar
finite-jet statement after the redundant direction has been
sectioned out.  It neither inverts nor contracts the Yang--Mills
massless covariance.  Gauge, toron, topological, and infrared
coordinates remain parameters, consistently with the companion
audit.
\end{remark}

\section{Affine forcing and the calibrated graph}
\label{sec:n1v10-graph}

An exact lattice step may generate anisotropy even when the input
coordinate is zero.  Let
\[
 f_{\beta,r}
 :=
 {\mathfrak a}\!\left(
 {\cal R}_\beta^\circ(S_r)
 -{\cal R}_{0}^\circ(S_r)
 \right),                                                   \tag{10.17}
\]
where the exact Gaussian reference anisotropy has been included in
the calibrated base \(S_r\).  The same connected expansion, now with
the output extractor as the sole mark, gives a fixed constant
\(C_{{\rm A},0}\) such that
\[
                         |f_{\beta,r}|
 \le C_{{\rm A},0}w_{10}(\beta).                            \tag{10.18}
\]

\begin{proposition}[Unique one-step anisotropy graph]
\label{prop:n1v10-fixed-graph}
Suppose the retained parameters are fixed and \(\beta\) is large
enough for (10.16).  The scalar map
\[
                         T_{\beta,r}(\alpha)
 =f_{\beta,r}+F_{\beta,r}(\alpha)                           \tag{10.19}
\]
has a unique fixed point in every interval containing
\[
 |\alpha|\le
 \frac{C_{{\rm A},0}w_{10}(\beta)}
      {1-q_{\rm A}(\beta)}.                                 \tag{10.20}
\]
The fixed graph is analytic in the retained parameters and is
\(O(w_{10}(\beta))\).
\end{proposition}

\begin{proof}
Equation (10.14) makes \(T_{\beta,r}\) a contraction with factor
\(q_{\rm A}<1\).  Equations (10.18)--(10.19) show that the interval
in (10.20) is invariant.  Banach's fixed-point theorem gives
existence and uniqueness.  Uniform analytic contraction iteration
gives analytic dependence on \(r\).
\end{proof}

\begin{lemma}[Cofinal restoration under vanishing forcing]
\label{lem:n1v10-cofinal}
Let a sequence of anisotropy coordinates obey
\[
 |\alpha_{j+1}|
 \le q|\alpha_j|+f_j,\qquad 0\le q<1,\qquad f_j\longrightarrow0.
                                                               \tag{10.21}
\]
Then \(\alpha_j\to0\).  More precisely,
\[
 |\alpha_n|
 \le q^n|\alpha_0|
 +\sum_{j=0}^{n-1}q^{n-1-j}f_j.                            \tag{10.22}
\]
\end{lemma}

\begin{proof}
Induction gives (10.22).  Given \(\epsilon>0\), choose \(J\) so that
\(f_j\le\epsilon(1-q)/2\) for \(j\ge J\).  The part of the sum before
\(J\) tends to zero after multiplication by \(q^{n-J}\); the tail is
at most \(\epsilon/2\).  The initial term also tends to zero.
\end{proof}

\begin{corollary}[Conditional cofinal anisotropy restoration]
\label{cor:n1v10-cofinal}
Along any weak-coupling cofinal trajectory for which the other
calibrated coordinates remain in the V8 domain and
\[
                         w_{10}(\beta_j)\longrightarrow0,   \tag{10.23}
\]
the hypercubic coordinate relative to the calibrated graph tends to
zero.
\end{corollary}

\begin{proof}
Apply Theorem~\ref{thm:n1v10-anisotropy-contraction},
(10.18), and Lemma~\ref{lem:n1v10-cofinal}.
\end{proof}

\begin{firewall}[Trajectory hypothesis remains visible]
\label{fire:n1v10-trajectory}
Corollary~\ref{cor:n1v10-cofinal} does not construct the weak-coupling
trajectory or prove (10.23) at all scales.  It proves the anisotropy
row once such a trajectory remains in the already established
one-scale domain.  The physical Wilson and theta coordinates, and
the topological sectors, are not included in this contraction.
\end{firewall}

\begin{assessment}[Progress after compact V10]
\label{ass:n1v10-status}
The compulsory audit confirms that a massless single-site influence
route cannot be uniform and provides no missing anisotropy theorem.
V10 instead uses the finite hypercubic jet after the V9 quotient.
One marked V8 connected expansion makes the non-Gaussian derivative
an explicit \(o(1)\) correction to canonical \(b^{-2}\) scaling.
This proves a strict, volume-uniform one-step anisotropy contraction
and a unique small calibrated anisotropy graph.
\end{assessment}

\begin{decision}[V11 target]
\label{dec:n1v10-next}
The next unresolved finite row is the topological and winding sector.
V11 should separate:
\begin{enumerate}[label=\textup{(\roman*)},leftmargin=3.5em]
\item the local density \(\operatorname{Tr}F\widetilde F\), whose
coefficient is marginal and cannot be contracted by dimensional
scaling;
\item exact finite-volume flux and toron labels, which cannot be put
inside the local owner norm; and
\item contractible topological fluctuations, to which the V8 marked
compiler may apply.
\end{enumerate}
It should prove a sectorwise pushforward identity and determine
whether reflection positivity forces the physical theta coordinate
to the real zero section before any terminal mixing argument.
\end{decision}

\begin{firewall}[Claim boundary after compact V10]
\label{fire:n1v10-claim}
V10 proves the finite hypercubic-row contraction on the established
one-scale weak-coupling domain.  It does not construct the complete
calibrated RG trajectory, control the topological/flux sector, prove
the universal marginal coefficient, enter a terminal mixing domain,
remove regulators, perform Osterwalder--Schrader reconstruction, or
prove the Yang--Mills mass gap.
\end{firewall}

\subsection{Parent record}

This V10 note was formed by appending the preceding material to the
validated compact V9, whose record was
\[
\begin{array}{c|c}
\text{lines}&4645\\
\text{bytes}&160005\\
\text{SHA--256}&
\texttt{ABD29A58D14806C7705611E40EBC28BFB92EAED1A397BBE3FEC47FDA73A0F1CD}.
\end{array}
\]

% =============================================================================
% END APPENDED COMPACT-SERIES V10 MATERIAL
% =============================================================================

% =============================================================================
% BEGIN APPENDED COMPACT-SERIES V11 MATERIAL
% =============================================================================

\chapter{V11: sectorwise topological transport and reflection
positivity}
\label{ch:n1v11}

\section{Admissible sectors and retained global labels}
\label{sec:n1v11-sectors}

The term \emph{topological} has covered three mathematically different
objects in the earlier ledgers:
\begin{enumerate}[label=\textup{(\roman*)},leftmargin=3.5em]
\item the local pseudoscalar density whose weak-field limit is
\(\operatorname{Tr}F\widetilde F\);
\item an integer global topological charge; and
\item flat flux, winding, and toron data not reconstructed from a
contractible filling.
\end{enumerate}
They cannot share one local polymer coordinate.  The first is a local
source, the second is a sector label, and the third is retained
boundary/cohomology data.

Let \(X_\Lambda\) be the finite reduced-link configuration space.
Choose a gauge-invariant admissible open set
\({\cal A}_\Lambda\subset X_\Lambda\) on which the geometric
topological charge
\[
                         Q_\Lambda:{\cal A}_\Lambda\to\mathbb Z
                                                               \tag{11.1}
\]
is continuous.  Let
\(\tau_\Lambda:{\cal A}_\Lambda\to{\mathscr H}_\Lambda\)
record the declared flux and toron data.  Discrete components of
\(\tau_\Lambda\) are locally constant; continuous toron coordinates
are held fixed by every reduced conditional fibre.

\begin{lemma}[Integer charge is constant on a regular fibre]
\label{lem:n1v11-charge-constant}
If \(C\) is a connected subset of one conditional fibre and
\(C\subset{\cal A}_\Lambda\), then \(Q_\Lambda\) is constant on
\(C\).  If the fibre definition also fixes \(\tau_\Lambda\), the pair
\((\tau_\Lambda,Q_\Lambda)\) is constant on \(C\).
\end{lemma}

\begin{proof}
The continuous image of a connected set is connected.  The only
connected subsets of the discrete space \(\mathbb Z\) are
singletons, proving the assertion for \(Q_\Lambda\).  Discrete flux
labels obey the same argument, while fixed continuous torons are
constant by the definition of the fibre.
\end{proof}

For \(h\in{\mathscr H}_\Lambda\) and \(n\in\mathbb Z\), put
\[
 {\cal A}_{h,n}
 =
 \{U\in{\cal A}_\Lambda:
   \tau_\Lambda(U)=h,\ Q_\Lambda(U)=n\}.                    \tag{11.2}
\]
Only finitely many integers occur at fixed lattice volume when the
geometric charge is a finite sum of bounded cell densities.

\begin{definition}[Extended coarse block map]
\label{def:n1v11-extended-map}
For an ordinary gauge-covariant block map
\({\cal B}:X_\Lambda\to X_{\Lambda'}\), define on the admissible set
\[
 \widehat{\cal B}(U)
 =
 \bigl({\cal B}(U),\tau_\Lambda(U),Q_\Lambda(U)\bigr).
                                                               \tag{11.3}
\]
The last two entries are retained labels; they are not recomputed
from \({\cal B}(U)\).
\end{definition}

\section{Exact sectorwise pushforward}
\label{sec:n1v11-pushforward}

Let \(d\nu=e^{-S(U)}dH(U)\) be an unnormalized positive
finite-volume Wilson measure.  Define
\[
 d\nu_{h,n}
 =
 {\bf1}_{{\cal A}_{h,n}}e^{-S}dH,\qquad
 d\nu_{\rm bad}
 =
 {\bf1}_{{\cal A}_\Lambda^c}e^{-S}dH.                      \tag{11.4}
\]
For real \(\theta\), the admissible theta deformation is
\[
 d\nu_\theta^{\rm adm}
 =
 \sum_{h,n}e^{i\theta n}d\nu_{h,n}.                        \tag{11.5}
\]

\begin{theorem}[Sectorwise block identity]
\label{thm:n1v11-sector-pushforward}
The extended pushforward is
\[
\boxed{
 \widehat{\cal B}_*\nu_\theta^{\rm adm}
 =
 \sum_{h,n}e^{i\theta n}
 \bigl({\cal B}_*\nu_{h,n}\bigr)
 \otimes\delta_h\otimes\delta_n.}                           \tag{11.6}
\]
In particular, exact integration of any family of connected regular
conditional fibres is diagonal in \((h,n)\).  The theta coefficient
is transported unchanged modulo \(2\pi\).
\end{theorem}

\begin{proof}
For a bounded measurable test function \(F(W,h,n)\), the left side
of (11.6) evaluated at \(F\) is
\[
 \sum_{h,n}e^{i\theta n}
 \int_{{\cal A}_{h,n}}
 F({\cal B}(U),h,n)e^{-S(U)}dH(U),
\]
which is the right side.  Lemma
\ref{lem:n1v11-charge-constant} makes every connected admissible
conditional integration diagonal.  Since the only theta dependence
in the \(n\)-th summand is \(e^{i\theta n}\), no integration changes
its coefficient.  Integer \(n\) gives \(2\pi\)-periodicity.
\end{proof}

\begin{corollary}[Exact effective action on the extended state]
\label{cor:n1v11-effective-sector}
If
\[
 {\cal B}_*\nu_{h,n}(dW)
 =
 e^{-S'_{h,n}(W)}dH_{\Lambda'}(W)
\]
on its absolutely continuous part, then the theta-deformed effective
action on the extended state is
\[
                         S'_{h,n}(W)-i\theta n.             \tag{11.7}
\]
All sector-dependent free energies belong to
\(S'_{h,n}\); none is a renormalization of \(\theta\).
\end{corollary}

\begin{proof}
Insert the density in (11.6) and combine
\(e^{i\theta n}\) with \(e^{-S'_{h,n}}\).
\end{proof}

\begin{warning}[Forgetting the label destroys the local statement]
\label{warn:n1v11-forgetting}
If \(n\) is discarded after the pushforward, the coarse density is
\[
                         \sum_n e^{i\theta n}\rho_n(W),     \tag{11.8}
\]
which need not be the exponential of a local action with one
coarse topological density.  The exact nonrunning statement is
therefore a theorem on the extended state (11.3), not a license to
sum over sectors and then postulate locality.
\end{warning}

\section{Ownership of the nonadmissible remainder}
\label{sec:n1v11-bad}

Choose the admissibility threshold inside the V7 collar chart so that
every configuration satisfying all tightened collar bounds
\(\kappa_y\le\bar\epsilon_0\) lies in
\({\cal A}_\Lambda\).  Then
\[
 {\cal A}_\Lambda^c
 \subset
 \bigcup_y\{\kappa_y>\bar\epsilon_0\}.                      \tag{11.9}
\]
This is a choice of the regular domain, not a claim that every
Wilson configuration possesses a smooth geometric charge.

\begin{proposition}[Topological transitions are bad-owned]
\label{prop:n1v11-transition-owned}
Any continuous path that changes the integer charge, or leaves one
admissible component on which the charge is defined, meets a
collar-irregular block.  In the triangular expansion its first such
block is assigned to one canonical V8 bad birth.  Consequently every
off-admissible or sector-transition word carries at least one
\(q_{\rm sat}(\beta)\) primitive and is never charged to a regular
local theta vertex.
\end{proposition}

\begin{proof}
If the path remained in \({\cal A}_\Lambda\), its connected image
under the continuous integer-valued map \(Q_\Lambda\) would be one
integer by Lemma~\ref{lem:n1v11-charge-constant}.  Hence it leaves
\({\cal A}_\Lambda\).  Inclusion (11.9) gives a collar-irregular
block at the first exit.  Definition~\ref{def:n1v8-history} assigns
the first unassigned such block to a unique bad birth, and
Lemma~\ref{lem:n1v8-bad-distribution} supplies its
\(q_{\rm sat}\) payment.
\end{proof}

\begin{corollary}[Sectorwise V8 compiler]
\label{cor:n1v11-sector-compiler}
For each fixed retained pair \((h,n)\), the V8 connected expansion
and all its constants apply without modification to regular
contractible owners.  Words leaving that sector are in the
saturated-bad bank.  No local connected word receives both a
topological-transition charge and a collar-bad charge.
\end{corollary}

\begin{proof}
Regular fibre kernels are diagonal by
Theorem~\ref{thm:n1v11-sector-pushforward}.  The V8 historical
assignment is unique, and
Proposition~\ref{prop:n1v11-transition-owned} assigns every
nondiagonal word to that existing bad mark.
\end{proof}

\section{Reflection positivity with a theta phase}
\label{sec:n1v11-rp}

Let \({\cal A}_+\) be the algebra of observables supported in the
positive-time half lattice and let \(\Theta\) be the usual
antilinear time-reflection involution.  A functional
\(\langle\cdot\rangle_0\) is reflection positive if
\[
                         \langle(\Theta F)F\rangle_0\ge0
 \quad(F\in{\cal A}_+).                                    \tag{11.10}
\]

\begin{theorem}[Half-space phase lemma]
\label{thm:n1v11-half-phase}
Suppose \(\langle\cdot\rangle_0\) is reflection positive and
\(q_+\in{\cal A}_+\) is real.  For any real \(\theta\), define
\[
 h_\theta=e^{i\theta q_+},\qquad
 w_\theta=h_\theta\,\Theta h_\theta
 =e^{i\theta(q_+-\Theta q_+)}.                              \tag{11.11}
\]
Then the unnormalized functional
\[
                         \langle F\rangle_\theta
 :=\langle Fw_\theta\rangle_0                              \tag{11.12}
\]
is reflection positive.  Its partition value
\(\langle1\rangle_\theta\) is nonnegative; if it is positive, the
normalized functional is reflection positive as well.
\end{theorem}

\begin{proof}
Antilinearity of \(\Theta\) gives
\(\Theta h_\theta=e^{-i\theta\Theta q_+}\), proving (11.11).
For \(F\in{\cal A}_+\),
\[
\begin{split}
 \langle(\Theta F)F\rangle_\theta
 &=
 \left\langle
 (\Theta F)Fh_\theta\Theta h_\theta
 \right\rangle_0\\
 &=
 \left\langle
 \Theta(Fh_\theta)(Fh_\theta)
 \right\rangle_0\ge0.
\end{split}                                                 \tag{11.13}
\]
Taking \(F=1\) gives the partition statement.  Division by a positive
number preserves the inequality.
\end{proof}

\begin{corollary}[Reflection positivity does not force
\(\theta=0\)]
\label{cor:n1v11-rp-theta}
If the lattice topological charge has the exact reflection split
\[
                         Q_\Lambda=q_+-\Theta q_+,          \tag{11.14}
\]
then multiplication of a reflection-positive Yang--Mills weight by
\(e^{i\theta Q_\Lambda}\) preserves reflection positivity for every
real \(\theta\) for which the partition value is nonzero.  Therefore
reflection positivity alone does not select the zero-theta section.
\end{corollary}

\begin{proof}
Equations (11.11) and (11.14) identify the theta factor with
\(w_\theta\).  Apply
Theorem~\ref{thm:n1v11-half-phase}.
\end{proof}

\begin{warning}[Boundary residue]
\label{warn:n1v11-boundary}
If instead
\[
                         Q_\Lambda
 =q_+-\Theta q_++q_0,                                     \tag{11.15}
\]
where \(q_0\) is supported on the reflection plane, the theorem
applies only if \(e^{i\theta q_0}\) is itself a positive reflection
kernel, namely a limit of finite sums
\(\sum_j(\Theta b_j)b_j\).  Reflection invariance of \(q_0\) alone is
not enough.  The boundary residue must remain in the boundary ledger.
\end{warning}

\begin{proof}
The cone of reflection-positive insertions is generated and closed by
factors \(\sum_j(\Theta b_j)b_j\).  Multiplying (11.13) by such a
factor preserves its representation as a sum of nonnegative OS
forms.  Without that representation, no sign conclusion follows.
\end{proof}

\section{The zero-theta physical section}
\label{sec:n1v11-zero}

The Clay mass-gap problem for pure Yang--Mills is ordinarily posed at
zero theta.  This is a choice of theory, not a consequence of
Corollary~\ref{cor:n1v11-rp-theta}.  At \(\theta=0\), every sector
weight in (11.6) is positive and the phase and boundary-residue
questions disappear.

\begin{proposition}[Zero theta is exactly preserved]
\label{prop:n1v11-zero-preserved}
On the extended state, if the input theta coordinate is zero, every
exact triangular RG pushforward has theta coordinate zero.  No
regular connected owner can generate it, and every charge-changing
path is already in the saturated-bad bank.
\end{proposition}

\begin{proof}
At \(\theta=0\), (11.6) contains no sector phase.
Corollary~\ref{cor:n1v11-effective-sector} shows that an exact
pushforward cannot create one.  Regular fibre integrations preserve
the integer label, and
Proposition~\ref{prop:n1v11-transition-owned} classifies all exits.
\end{proof}

\begin{remark}[Local pseudoscalar marks]
\label{rem:n1v11-local-density}
A source coupled to a local representative of
\(\operatorname{Tr}F\widetilde F\) may still be differentiated to
compute topological correlations.  On a fixed sector its total sum
is the constant \(n\); local fluctuations are ordinary marked
connected owners, whereas the constant mode belongs to the retained
sector label.  This separates the local source calculus from the
global theta coordinate without setting either to zero by notation.
\end{remark}

\begin{assessment}[Progress after compact V11]
\label{ass:n1v11-status}
V11 gives an exact extended-state pushforward in which global charge,
flux, and toron labels are retained and the theta coefficient is
transported without renormalization.  Regular local fibres are
sector-diagonal; every admissibility or charge transition is paid by
the already existing saturated-bad owner.  The half-space phase lemma
also proves that reflection positivity does not by itself force
\(\theta=0\).  The zero-theta theory is nevertheless an exact
invariant section and is the section used below.
\end{assessment}

\begin{decision}[V12 target]
\label{dec:n1v11-next}
Combine the V8 all-colour compiler, the V9 redundant-coordinate
section, the V10 anisotropy contraction, and the V11 extended-sector
identity into one calibrated one-step domain.  The next note should
state a triangular majorant map for:
\[
 (\hbox{physical marginal},\hbox{theta label},
 \hbox{anisotropy},\hbox{stable packet owners},
 \hbox{bad activity}).                                     \tag{11.16}
\]
It must distinguish a one-step invariant domain from an infinite
trajectory and expose the still-missing marginal beta-function
coefficient and scale-to-scale parameter matching.
\end{decision}

\begin{firewall}[Claim boundary after compact V11]
\label{fire:n1v11-claim}
V11 proves sectorwise transport and the reflection-positive theta
phase lemma.  It does not prove a uniform sum over infinite-volume
topological sectors, boundary-residue positivity for an arbitrary
lattice charge, the marginal Yang--Mills flow, the full RG
trajectory, terminal mixing, regulator removal,
Osterwalder--Schrader reconstruction, or the Yang--Mills mass gap.
\end{firewall}

\subsection{Parent record}

This V11 note was formed by appending the preceding material to the
validated compact V10, whose record was
\[
\begin{array}{c|c}
\text{lines}&5068\\
\text{bytes}&175478\\
\text{SHA--256}&
\texttt{EB228387210698AF0D6AA82B4194A843FFE0B920033EDF201C9232AEC36A517D}.
\end{array}
\]

% =============================================================================
% END APPENDED COMPACT-SERIES V11 MATERIAL
% =============================================================================

% =============================================================================
% BEGIN APPENDED COMPACT-SERIES V12 MATERIAL
% =============================================================================

\section{Compact-series V12: the calibrated product domain and the
marginal two-jet firewall}
\label{sec:n1v12}

\subsection{Why the next step is an exponent calculation}

The legacy packet programme already proved a conditional aligned
stable contraction and a conditional cofinal trajectory theorem.
Repeating that graph transform would not advance the argument.
Compact V8--V11 instead discharged four pieces which were not
simultaneously present there:
\begin{enumerate}
\item the all-colour fibre is one exact triangular pushforward, with
      a unique historical owner for every saturated component;
\item the raw Jacobian marginal may be removed by an exact measured
      redundant coordinate;
\item the remaining anisotropy coordinate is strictly contractive
      and has a unique small calibrated graph; and
\item global flux, toron, and charge labels are transported exactly,
      while the zero-theta theory is an invariant section.
\end{enumerate}
The correct new question is therefore whether inserting those
calibrated lower coordinates changes the quadratic coefficient of
the physical marginal map.  This note proves the exact sufficient
exponent condition.  It also proves that norm-smallness of the lower
coordinates, without a response factor, cannot determine the
Yang--Mills beta coefficient.

\subsection{The triangular calibrated cone}

Write
\[
 {\bf z}=(z_1,\ldots,z_N)\in{\mathbb R}_+^N                 \tag{12.1}
\]
for majorants of all nonmarginal dynamical rows after the exact
redundant quotient.  The list includes the saturated-bad activity,
the stable packet owner norm, the stationary-reference mismatch,
and the anisotropy coefficient.  It does not include the retained
sector labels \((h,n)\), which are transported exactly, or theta,
which is fixed to zero.

Let \(u>0\) be a normalized physical marginal coordinate, with
\(u=0\) the ultraviolet fixed point.  Different components may have
different decay exponents \(p_i>0\).  Put
\[
 D(u;{\bf H})
 :=
 \left\{
 {\bf z}:0\le z_i\le H_i u^{p_i},\ 1\le i\le N
 \right\}.                                                   \tag{12.2}
\]
For a nonnegative matrix \(Q=(Q_{ij})\), set
\[
 (Q_{\bf p}(\sigma))_{ij}
 :=\sigma^{-p_i}Q_{ij},
 \qquad
 d_{\bf p}(\sigma)_i:=\sigma^{-p_i}d_i.                       \tag{12.3}
\]
This row rescaling records that the target radius is evaluated at
the new marginal coordinate.

\begin{theorem}[Triangular calibrated-cone theorem]
\label{thm:n1v12-cone}
Fix \(0<\sigma<1\).  Suppose a one-step map obeys
\[
             \frac{u^+}{u}\ge\sigma                            \tag{12.4}
\]
and, for every \({\bf z}\in D(u;{\bf H})\),
\[
 z_i^+
 \le
 \sum_{j=1}^N Q_{ij}z_j+d_i u^{p_i}.                           \tag{12.5}
\]
Assume \(Q_{ij}=0\) whenever \(p_j\ne p_i\), or more generally that
every nonzero entry satisfies \(p_j\ge p_i\) and \(0<u\le1\).
If
\[
          \rho\!\left(Q_{\bf p}(\sigma)\right)<1,              \tag{12.6}
\]
then
\[
 {\bf H}:=
 \left(I-Q_{\bf p}(\sigma)\right)^{-1}
 d_{\bf p}(\sigma)                                            \tag{12.7}
\]
is nonnegative and
\[
             {\bf z}\in D(u;{\bf H})
 \quad\Longrightarrow\quad
             {\bf z}^+\in D(u^+;{\bf H}).                      \tag{12.8}
\]
If \(Q\) is lower triangular, condition (12.6) is precisely
\[
                         Q_{ii}<\sigma^{p_i}
 \quad(1\le i\le N).                                          \tag{12.9}
\]
\end{theorem}

\begin{proof}
For a nonzero entry with \(p_j\ge p_i\) and \(u\le1\),
\[
 Q_{ij}H_j u^{p_j}\le Q_{ij}H_j u^{p_i}.
\]
Thus (12.5) gives
\[
 z_i^+
 \le u^{p_i}\{(Q{\bf H})_i+d_i\}.                              \tag{12.10}
\]
By (12.6), the Neumann series
\[
 (I-Q_{\bf p}(\sigma))^{-1}
 =
 \sum_{k\ge0}Q_{\bf p}(\sigma)^k
\]
converges entrywise and has nonnegative entries.  Equation (12.7)
is equivalent, row by row, to
\[
                   (Q{\bf H})_i+d_i
                   =\sigma^{p_i}H_i.                           \tag{12.11}
\]
Combining (12.4), (12.10), and (12.11) yields
\[
 z_i^+\le(\sigma u)^{p_i}H_i
        \le (u^+)^{p_i}H_i,
\]
which is (12.8).  A triangular matrix has its diagonal entries as
its eigenvalues; row rescaling changes them to
\(\sigma^{-p_i}Q_{ii}\), proving (12.9).
\end{proof}

\begin{remark}[Why a fixed \(\sigma\) is harmless cofinally]
\label{rem:n1v12-sigma}
If \(u^+=u+O(u^2)\), then \(u^+/u=1+O(u)\).  Hence every fixed
\(\sigma<1\) satisfies (12.4) for all sufficiently small \(u\).
The strict packet and anisotropy margins may therefore be spent
against \(\sigma^{-p_i}\) once, without imposing a false uniform
gap on the marginal direction.
\end{remark}

\begin{proposition}[Exact sector extension of the cone]
\label{prop:n1v12-sector-cone}
Let \(D_{h,n}(u;{\bf H})\) be (12.2) in the retained sector
\((h,n)\), with redundant coordinate zero and theta zero.  Under the
extended block map of V11,
\[
 D_{h,n}(u;{\bf H})
 \longrightarrow
 D_{h^+,n}(u^+;{\bf H})                                      \tag{12.12}
\]
whenever the hypotheses of
Theorem~\ref{thm:n1v12-cone} hold.  No sector sum enters the
majorant \(Q\), and every charge-changing exit is already charged to
the saturated-bad component.
\end{proposition}

\begin{proof}
V11 proves \(n^+=n\) on every regular admissible fibre and assigns
every exceptional transition to one historical bad owner.  It also
proves \(\theta^+=\theta\), so the zero-theta section remains zero.
V9 makes the redundant coordinate an exact section of the measured
quotient.  The only numerical coordinates left are therefore those
in (12.1), to which
Theorem~\ref{thm:n1v12-cone} applies.  The flux label \(h\) is
transported by the declared coarse sector map and is not summed in
any row of (12.5).
\end{proof}

\begin{assessment}[What the product-domain splice actually proves]
\label{ass:n1v12-product}
On every scale where the previously proved packet, forest, collar,
and analytic cards supply (12.5), compact V8--V11 turn the old
conditional stable ball into a single sectorwise calibrated product
domain.  The calibration order is noncircular: bad ownership first,
then stable packet contraction, then the exact measured quotient,
then anisotropy, while theta and the integer charge are retained
labels.  This is a one-step domain theorem; an infinite orbit still
requires the marginal recursion.
\end{assessment}

\subsection{The physical two-jet under calibration}

Let \(X_i\) be Banach spaces and write
\[
 {\cal D}(u)
 :=
 \left\{
 x=(x_1,\ldots,x_N):
 \|x_i\|_{X_i}\le H_i u^{p_i}
 \right\}.                                                   \tag{12.13}
\]
Let \(F(u,x)\) be the physical marginal component of the exact
one-step map after all normalizations have been declared.

\begin{theorem}[Marginal two-jet preservation]
\label{thm:n1v12-twojet}
Suppose, for some \(b>0\), \(C_0<\infty\), and
\(\delta_0>0\),
\[
 |F(u,0)-u+bu^2|
 \le C_0u^{2+\delta_0}.                                      \tag{12.14}
\]
Suppose also that on the segment from \(0\) to every
\(x\in{\cal D}(u)\),
\[
 \|D_{x_i}F(u,x)\|_{X_i^*}
 \le L_i u^{q_i},                                            \tag{12.15}
\]
where
\[
                         p_i+q_i>2
 \quad(1\le i\le N).                                         \tag{12.16}
\]
Then, uniformly for \(x\in{\cal D}(u)\),
\[
 F(u,x)
 =
 u-bu^2+
 O\!\left(u^{2+\delta}\right),                               \tag{12.17}
\]
where
\[
 \delta
 :=
 \min\left\{
 \delta_0,\min_i(p_i+q_i-2)
 \right\}>0.                                                 \tag{12.18}
\]
In particular the quadratic coefficient \(b\) is unchanged by the
complete calibrated lower graph.
\end{theorem}

\begin{proof}
The Banach-space fundamental theorem of calculus gives
\[
 F(u,x)-F(u,0)
 =
 \sum_{i=1}^N
 \int_0^1
 D_{x_i}F(u,tx)[x_i]\,dt.                                    \tag{12.19}
\]
Equations (12.13) and (12.15) bound the \(i\)-th term by
\[
                         L_iH_i u^{p_i+q_i}.                  \tag{12.20}
\]
Add (12.14), take the minimum exponent, and use the finiteness of
the coordinate list.  This proves (12.17)--(12.18).
\end{proof}

\begin{corollary}[The exponent required by the compact activities]
\label{cor:n1v12-exponent}
On the compact V6--V10 regular branch, the two-jet-subtracted good
activity, the compensated Jacobian defect, and the anisotropy graph
have the common conservative order
\[
             O(\beta^{-6/5})=O(u^{6/5}),
 \qquad u=\beta^{-1}.                                        \tag{12.21}
\]
For any such coordinate, the response estimate
\[
                         \|D_xF(u,x)\|\le Lu                  \tag{12.22}
\]
is sufficient to preserve the marginal quadratic coefficient,
because
\[
                         \frac65+1=\frac{11}{5}>2.            \tag{12.23}
\]
More sharply, any exponent \(q>4/5\) in (12.15) is sufficient.
\end{corollary}

\begin{proof}
Apply Theorem~\ref{thm:n1v12-twojet} with \(p=6/5\).
\end{proof}

\begin{warning}[Activity smallness alone does not identify the beta
coefficient]
\label{warn:n1v12-smallness}
The bound \(x=O(u^{6/5})\) cannot by itself be inserted into the
marginal row.  A response factor such as (12.22), or an exact
orthogonality which is even stronger, is indispensable.
\end{warning}

\begin{proof}
On scalar \(x\), take
\[
                         F(u,x)=u-bu^2+x,\qquad x(u)=u^{6/5}.
                                                                    \tag{12.24}
\]
Then \(x(u)\to0\), but
\[
 \frac{u-F(u,x(u))}{u^2}
 =
 b-u^{-4/5},
\]
which has no finite limit.  Thus absolute smallness of a calibrated
coordinate does not control its marginal projection at quadratic
order.
\end{proof}

\subsection{Coordinates which provably cannot alter the coefficient}

\begin{lemma}[Parabolic two-jet is invariant under normalized
reparametrization]
\label{lem:n1v12-reparam}
Let
\[
 f(u)=u-bu^2+O(u^{2+\delta}),\qquad
 \phi(u)=u+cu^2+O(u^{2+\delta})
                                                                    \tag{12.25}
\]
for some \(\delta>0\), with \(\phi'(0)=1\).  Then
\[
                  \phi\circ f\circ\phi^{-1}(v)
                  =v-bv^2+O(v^{2+\delta'}),                    \tag{12.26}
\]
where \(\delta'=\min\{\delta,1\}\).
\end{lemma}

\begin{proof}
Taylor inversion gives
\[
 \phi^{-1}(v)=v-cv^2+O(v^{2+\delta'}).
\]
Writing \(u=\phi^{-1}(v)\), one has
\[
 f(u)=v-(b+c)v^2+O(v^{2+\delta'}).
\]
Finally,
\[
 \phi(f(u))
 =
 f(u)+c f(u)^2+O(f(u)^{2+\delta'})
 =
 v-bv^2+O(v^{2+\delta'}).
\]
\end{proof}

\begin{proposition}[The V9 redundant correction has exactly zero
physical leakage]
\label{prop:n1v12-redundant-zero}
Let a physical marginal coordinate be defined from the normalized
measured pushforward, rather than from the coefficient of a chosen
action representative.  The V9 replacement
\[
                         {\cal U}\longmapsto{\cal U}^{\circ}
                                                                    \tag{12.27}
\]
which kills the raw Jacobian extractor changes that physical
marginal map by exactly zero.
\end{proposition}

\begin{proof}
V9 constructs \({\cal U}^{\circ}\) by adding the divergence of a
Haar vector field together with the corresponding action
derivative.  Its defining identity says that the two expressions
are the same measured coordinate class: their normalized integrals
and all pushed physical source derivatives agree.  A marginal
coordinate defined from that pushed measure is therefore identical
before and after (12.27).  This is stronger than
Lemma~\ref{lem:n1v12-reparam}; there is no remainder to estimate.
\end{proof}

\begin{corollary}[Only two leakage families remain]
\label{cor:n1v12-two-leaks}
On the zero-theta extended state:
\begin{enumerate}
\item theta and the integer sector label have zero marginal
      derivative because they are retained, not calibrated;
\item the V9 redundant direction has zero physical derivative by
      Proposition~\ref{prop:n1v12-redundant-zero}; and
\item the possible changes of the marginal two-jet come only from
      the stable/owner graph and the anisotropy graph, together with
      finite-regulator mismatch in the raw one-loop coefficient.
\end{enumerate}
\end{corollary}

\begin{proof}
The first item is the exact sectorwise transport of V11, the second
is the preceding proposition, and the list of numerical coordinates
in (12.1) gives the third.
\end{proof}

\subsection{The reduced marginal matching card}

Index a cofinal family of finite regulators by \(r\).  Let
\(F_r^0(u):=F_r(u,0)\), after the exact Haar determinant, gauge
quotient, ghost contribution, normalization, and block convention
have all been included.

\begin{definition}[Reduced marginal matching card]
\label{def:n1v12-RMMC}
RMMC consists of the following two estimates on a common
regulator-uniform neighbourhood:
\begin{align}
 |F_r^0(u)-u+b_{\rm YM}u^2|
 &\le
 C u^{2+\delta_0}+\epsilon_r u^2,
 \qquad \epsilon_r\longrightarrow0,                            \tag{12.28}\\
 \|D_{x_i}F_r(u,x)\|
 &\le L_i u^{q_i},
 \qquad p_i+q_i>2,                                             \tag{12.29}
\end{align}
for the stable/owner and anisotropy coordinates only.  Here
\(b_{\rm YM}>0\) is computed in the same normalized block
convention as \(F_r\).
\end{definition}

\begin{theorem}[RMMC closes the full calibrated marginal two-jet]
\label{thm:n1v12-RMMC}
Assume RMMC and the calibrated cone
Theorem~\ref{thm:n1v12-cone}.  Then, uniformly on that cone,
\[
 F_r(u,x)
 =
 u-b_{\rm YM}u^2+
 O(u^{2+\delta})+O(\epsilon_r u^2),                            \tag{12.30}
\]
with \(\delta>0\) given by (12.18).  Consequently
\[
 \lim_{\substack{u\downarrow0\\ r\to\infty}}
 \frac{u-F_r(u,x)}{u^2}
 =
 b_{\rm YM}                                                   \tag{12.31}
\]
for every joint limit along which the calibrated bounds hold.
\end{theorem}

\begin{proof}
Use (12.28) in place of (12.14) in the proof of
Theorem~\ref{thm:n1v12-twojet}.  Division of (12.30) by \(u^2\)
leaves errors \(O(u^\delta)+O(\epsilon_r)\), proving (12.31).
\end{proof}

\begin{assessment}[What has and has not been reduced]
\label{ass:n1v12-RMMC-status}
The continuum determinant calculation in the legacy notes fixes the
candidate universal coefficient, and the compact-link work fixes
the exact finite-regulator expression.  Those facts do not yet
prove (12.28) for the complete generated family.  Compact V9 removes
the measured redundant contamination exactly, V10 supplies the
\(u^{6/5}\) anisotropy graph, and V11 removes theta and topology from
the leakage calculation.  Thus the remaining new estimates have the
precise form:
\[
 \boxed{
 \text{finite-regulator raw matching (12.28)}
 \quad+\quad
 \text{an }O(u)\text{ stable/anisotropy response (12.29)}.}     \tag{12.32}
\]
Neither estimate may be replaced by the statement that the
activities tend to zero.
\end{assessment}

\begin{decision}[V13 upstream target]
\label{dec:n1v12-next}
Work upstream on the response half of RMMC.  Differentiate the
retained Wilson coefficient with respect to one stable or anisotropy
insertion in the exact background-field quotient.  The target is a
factorization
\[
 D_xF_r(u,x)
 =
 u\,{\cal K}_{r,x}(u,x),\qquad
 \sup_{r,u,x}\|{\cal K}_{r,x}(u,x)\|<\infty,                    \tag{12.33}
\]
obtained from the Ward/background-field normalization, not from a
bare activity estimate.  If (12.33) fails for an allowed owner, that
owner must be promoted into the retained coordinate list before the
beta coefficient can be defined.
\end{decision}

\begin{firewall}[Claim boundary after compact V12]
\label{fire:n1v12-claim}
V12 proves the sectorwise calibrated-cone theorem, the marginal
two-jet preservation criterion, the insufficiency of ungraded
activity smallness, and exact zero leakage of the measured redundant
coordinate.  It does not prove RMMC (12.28)--(12.29), the universal
finite-regulator marginal recursion, the infinite RG orbit,
regulator removal, terminal mixing, Osterwalder--Schrader
reconstruction, or the Yang--Mills mass gap.
\end{firewall}

\subsection{Parent record}

This V12 note was formed by appending the preceding material to the
validated compact V11, whose record was
\[
\begin{array}{c|c}
\text{lines}&5468\\
\text{bytes}&190050\\
\text{SHA--256}&
\texttt{723BBD8CAD9141B52825EEDAFC7CD89DBBDB2D323CD3EE5174907A682B092D42}.
\end{array}
\]

% =============================================================================
% END APPENDED COMPACT-SERIES V12 MATERIAL
% =============================================================================

% =============================================================================
% BEGIN APPENDED COMPACT-SERIES V13 MATERIAL
% =============================================================================

\section{Compact-series V13: reciprocal-coupling shielding and
closure of the calibrated-response row}
\label{sec:n1v13}

\subsection{The correct variable for the marked estimate}

V12 asked for an \(O(u)\) bound on the response of the physical
coupling \(u\) to a stable or anisotropy insertion.  There is a
stronger route.  The local Wilson extractor naturally produces the
inverse coupling
\[
                              \beta=u^{-1}.                    \tag{13.1}
\]
A uniformly bounded marked response of the one-step increment in
\(\beta\) automatically acquires two powers of \(u\) after inversion.
This observation moves the difficult factor out of the forest
estimate and into an exact scalar identity.

For one finite regulator \(r\), write the inverse-coupling row as
\[
 B_r(\beta,x)=\beta+c_r(\beta,x),                              \tag{13.2}
\]
where \(x\) denotes the complete calibrated stable/owner and
anisotropy variables.  Its physical-coupling form is
\[
 F_r(u,x)
 :=
 \frac{1}{B_r(u^{-1},x)}
 =
 \frac{u}{1+u c_r(u^{-1},x)}.                                 \tag{13.3}
\]
The orientation is the ultraviolet one: a positive limiting
increment \(c_r\) makes \(u\) decrease.

\begin{lemma}[Reciprocal shielding]
\label{lem:n1v13-shield}
Suppose, on a domain \({\cal X}\),
\[
 |c_r(\beta,x)|\le M,\qquad
 \|D_xc_r(\beta,x)\|_{{\cal X}^*}\le L                         \tag{13.4}
\]
uniformly in \(r,\beta,x\).  If \(0<u\le(2M)^{-1}\), then
\[
                         \|D_xF_r(u,x)\|
                         \le4Lu^2.                             \tag{13.5}
\]
\end{lemma}

\begin{proof}
Differentiate (13.3) at fixed \(u\):
\[
 D_xF_r(u,x)
 =
 -\frac{u^2D_xc_r(u^{-1},x)}
        {(1+u c_r(u^{-1},x))^2}.                               \tag{13.6}
\]
The first bound in (13.4) gives
\(|1+uc_r|\ge1-uM\ge1/2\).  Insert the second bound into
(13.6).
\end{proof}

\begin{corollary}[V12 response exponent with room to spare]
\label{cor:n1v13-v12-response}
For every calibrated coordinate of size \(O(u^{6/5})\), (13.5)
gives
\[
                    F_r(u,x)-F_r(u,0)
                    =O(u^{16/5}).                              \tag{13.7}
\]
Thus the V12 condition \(p+q>2\) holds with
\[
                         p=\frac65,\qquad q=2.                 \tag{13.8}
\]
\end{corollary}

\begin{proof}
Integrate (13.5) on the segment from \(0\) to \(x\) and use
\(\|x\|=O(u^{6/5})\).
\end{proof}

\subsection{A marked bound for the inverse-coupling increment}

Let \(\mathfrak b\) be the bounded local extractor of the normalized
isotropic Wilson coefficient in the exact output action.  Thus,
schematically but with a fixed declared normalization,
\[
 c_r(\beta,x)
 =
 \mathfrak b\!\left({\cal R}_{\beta,r}^{\circ}(S_\beta+x)\right)
 -\beta.                                                       \tag{13.9}
\]
Here \({\cal R}_{\beta,r}^{\circ}\) is the V8 exact all-colour
pushforward followed by the V9 measured-coordinate section.  Sector
labels are held fixed as in V11.

Let \(L_{\rm G}\) be the operator norm of the Gaussian-principal
stable/anisotropy-to-\(\mathfrak b\) response.  Let \(M_x\) be the
strong analytic norm of one unit input mark, and let
\(C_{\mathfrak b}\) contain the fixed output extractor, packet-to-
local, owner, and source-radius embeddings.  Retain the compact V10
quantities
\[
 w_{10}(\beta),\qquad
 x_{10}(\beta)=A_\mu c_{10}w_{10}(\beta).                     \tag{13.10}
\]

\begin{theorem}[Uniform marked inverse-coupling response]
\label{thm:n1v13-marked-inverse}
On every compact all-colour domain for which \(x_{10}(\beta)<1\),
\[
\boxed{
 \|D_xc_r(\beta,x)\|
 \le
 L_{\rm inv}(\beta)
 :=
 L_{\rm G}
 \frac{e^\mu C_{\mathfrak b}M_xc_{10}w_{10}(\beta)}
      {(1-x_{10}(\beta))^2}.}                                 \tag{13.11}
\]
The bound is uniform in finite volume and in the colour stage.  If
the fixed local extractor and input-mark embeddings are uniform in
the regulator family, it is uniform in \(r\) as well.
\end{theorem}

\begin{proof}
Differentiate (13.9) in a unit direction of \(x\).  Exact
differentiation of a normalized conditional logarithm gives one
distinguished marked insertion.  Split its connected expansion into
the Gaussian-principal word and the words containing at least one
nonreference primitive.  The first has norm at most \(L_{\rm G}\)
by definition.

For the second family, root the selected connected exploration at
the distinguished input mark.  A word with \(n\ge1\)
nonreference primitives has its first incident strength bounded by
\[
                     e^\mu C_{\mathfrak b}M_xc_{10}w_{10},
                                                                    \tag{13.12}
\]
and each subsequent primitive costs at most \(x_{10}\).
There are at most \(n+1\) positions of the distinguished mark in the
rooted exploration.  Therefore the complete marked remainder is at
most
\[
 e^\mu C_{\mathfrak b}M_xc_{10}w_{10}
 \sum_{n\ge0}(n+1)x_{10}^n
 =
 \frac{e^\mu C_{\mathfrak b}M_xc_{10}w_{10}}
      {(1-x_{10})^2}.                                         \tag{13.13}
\]
This is the same marked resolvent mechanism proved for the
anisotropy output in
Lemma~\ref{lem:n1v10-one-mark}; only the bounded output extractor is
changed.  Normal convergence in
Theorem~\ref{thm:n1v8-all-colour} justifies differentiation and
summation.  Historical saturation assigns every bad word once, so
no additional sector or colour multiplicity appears.
\end{proof}

\begin{corollary}[Cofinal boundedness of the response]
\label{cor:n1v13-response-bounded}
There is a finite \(\beta_{13}\) such that, for
\(\beta\ge\beta_{13}\),
\[
 x_{10}(\beta)\le\frac12,\qquad
 L_{\rm inv}(\beta)
 \le
 L_{\rm G}
 +4e^\mu C_{\mathfrak b}M_xc_{10}w_{10}(\beta).                \tag{13.14}
\]
In particular
\[
                    \sup_{\beta\ge\beta_{13}}L_{\rm inv}(\beta)
                    <\infty.                                  \tag{13.15}
\]
\end{corollary}

\begin{proof}
Compact V10 proves \(w_{10}(\beta)\to0\), hence
\(x_{10}(\beta)\to0\).  Choose \(\beta_{13}\) so that
\(x_{10}\le1/2\) and use (13.11).
\end{proof}

\begin{proposition}[Closure of the RMMC response half]
\label{prop:n1v13-RMMC-response}
Assume the inverse-coupling increment itself is uniformly bounded,
\[
             |c_r(\beta,x)|\le M_c                             \tag{13.16}
\]
on the compact calibrated domain.  Then the RMMC response estimate
(12.29) holds with \(q_i=2\) for every stable/owner and anisotropy
coordinate controlled by
Theorem~\ref{thm:n1v13-marked-inverse}.
\end{proposition}

\begin{proof}
Combine Lemma~\ref{lem:n1v13-shield}, with
\(M=M_c\), and
Corollary~\ref{cor:n1v13-response-bounded}.  The resulting constant
\(4\sup L_{\rm inv}\) is uniform on a sufficiently small
\(u\)-interval.
\end{proof}

\begin{remark}[The bounded-increment hypothesis is the natural
normalization card]
\label{rem:n1v13-bounded-increment}
The quantity \(c_r\) is one finite-scale change of the coefficient
of the normalized Wilson operator.  Its Gaussian-principal part is
the one-loop number, while its connected remainder tends to zero on
the compact weak-coupling domain.  Thus (13.16) is strictly weaker
than identifying its limit.  Nevertheless it must be checked in the
same finite-regulator convention; compactness of each individual
fibre does not by itself make the bound uniform in \(r\).
\end{remark}

\subsection{Inverse-coupling matching is equivalent to RMMC}

\begin{theorem}[Scalar inversion with a uniform remainder]
\label{thm:n1v13-inversion}
Suppose
\[
 c_r(\beta,0)
 =
 b_{\rm YM}+e_r(\beta),\qquad
 |e_r(\beta)|
 \le C\beta^{-\delta}+\epsilon_r,                              \tag{13.17}
\]
where \(\delta>0\), \(\epsilon_r\to0\), and
\(|c_r(\beta,0)|\le M_c\).  Then, for
\(0<u\le(2M_c)^{-1}\),
\[
 \left|
 F_r(u,0)-u+b_{\rm YM}u^2
 \right|
 \le
 C u^{2+\delta}
 +\epsilon_r u^2
 +2M_c^2u^3.                                                   \tag{13.18}
\]
Consequently (13.17) implies the raw RMMC estimate (12.28).
\end{theorem}

\begin{proof}
The exact algebraic identity
\[
 \frac{u}{1+uc}
 =
 u-cu^2+\frac{c^2u^3}{1+uc}                                   \tag{13.19}
\]
with \(c=c_r(u^{-1},0)\) gives
\[
 F_r(u,0)-u+b_{\rm YM}u^2
 =
 -e_r(u^{-1})u^2+
 \frac{c_r(u^{-1},0)^2u^3}
      {1+u c_r(u^{-1},0)}.                                    \tag{13.20}
\]
Use (13.17), \(|c_r|\le M_c\), and
\(|1+uc_r|\ge1/2\).
\end{proof}

\begin{corollary}[Reduced remaining marginal card]
\label{cor:n1v13-reduced-card}
Subject to the already declared uniform local extractor and marked-
tree embeddings, the full RMMC of V12 reduces to the single
inverse-coupling limit
\[
\boxed{
 c_r(\beta,0)
 =
 b_{\rm YM}
 +O(\beta^{-\delta})
 +O(\epsilon_r),
 \qquad
 \epsilon_r\longrightarrow0.}                                \tag{13.21}
\]
\end{corollary}

\begin{proof}
Theorem~\ref{thm:n1v13-inversion} gives the raw half of RMMC.
Proposition~\ref{prop:n1v13-RMMC-response} gives the calibrated
response half.  V12 removes theta, topological labels, and the
measured redundant coordinate from the leakage list.
\end{proof}

\subsection{The scalar cofinal orbit once (13.21) is known}

\begin{theorem}[Asymptotic reciprocal trajectory]
\label{thm:n1v13-reciprocal-trajectory}
Let
\[
 \beta_{j+1}
 =
 \beta_j+b_{\rm YM}+e_j,\qquad
 |e_j|\le C\beta_j^{-\delta}+\epsilon_j,                       \tag{13.22}
\]
where \(b_{\rm YM}>0\), \(\delta>0\), and
\(\epsilon_j\to0\).  For sufficiently large \(\beta_0\), the
sequence is strictly increasing and
\[
                         \frac{\beta_j}{j}\longrightarrow
                         b_{\rm YM},\qquad
                         u_j=\beta_j^{-1}\sim
                         \frac1{b_{\rm YM}j}.                  \tag{13.23}
\]
\end{theorem}

\begin{proof}
Since \(\beta^{-\delta}\to0\) and \(\epsilon_j\to0\), choose the
initial cofinal index so that \(|e_j|\le b_{\rm YM}/2\) as long as
\(\beta_j\ge\beta_0\).  Then
\[
 \frac12b_{\rm YM}
 \le\beta_{j+1}-\beta_j
 \le\frac32b_{\rm YM},                                        \tag{13.24}
\]
so induction preserves the domain and gives
\(\beta_j\ge\beta_0+\frac12b_{\rm YM}j\).
Summing (13.22) yields
\[
 \frac{\beta_j}{j}
 =
 \frac{\beta_0}{j}+b_{\rm YM}
 +\frac1j\sum_{k<j}e_k.                                      \tag{13.25}
\]
The last term tends to zero by the Cesaro theorem because
\[
 |e_k|
 \le
 C\left(\beta_0+\frac12b_{\rm YM}k\right)^{-\delta}
 +\epsilon_k
 \longrightarrow0.
\]
This proves the first limit in (13.23); taking reciprocals proves
the second.
\end{proof}

\begin{warning}[A scalar orbit is not yet the Yang--Mills measure]
\label{warn:n1v13-scalar-only}
Theorem~\ref{thm:n1v13-reciprocal-trajectory} controls only the
retained coupling coordinate.  Coupling it to the calibrated cone
requires the regulator-uniform hypotheses in (13.21), the sectorwise
lower-row bounds, and the moving-reference compatibility already
listed in the programme.  It neither constructs an infinite-volume
measure nor supplies terminal exponential clustering.
\end{warning}

\begin{assessment}[Progress after compact V13]
\label{ass:n1v13-status}
V13 closes the response half of the reduced marginal matching card
at the level of the established compact marked-tree estimates.
The key gain is reciprocal shielding: a bounded response of the
inverse-coupling increment becomes an \(O(u^2)\) response of the
physical coupling, so every \(O(u^{6/5})\) calibrated defect enters
only at \(O(u^{16/5})\).  The marginal bottleneck is now the one
scalar, regulator-matched limit (13.21), rather than an uncontrolled
coupling to the entire stable Banach space.
\end{assessment}

\begin{decision}[V14 upstream target]
\label{dec:n1v13-next}
Return to the finite-regulator background-field calculation and
write \(c_r(\beta,0)\) as
\[
 \text{Gaussian determinant coefficient}
 +\text{compact/nonlinear remainder}
 +\text{reset mismatch}.                                     \tag{13.26}
\]
Use the exact V9 measured quotient before taking the extractor.
The next note should prove a telescoping criterion for (13.21) and
identify, without a placeholder constant, which of the compact
kernel difference and reset mismatch still lacks a summable bound.
\end{decision}

\begin{firewall}[Claim boundary after compact V13]
\label{fire:n1v13-claim}
V13 proves reciprocal shielding, a uniform one-mark response bound
from the existing compact compiler, the reduction of RMMC to
(13.21), and the scalar asymptotic trajectory conditional on that
limit.  It does not prove the finite-regulator limit (13.21), the
complete coupled RG orbit, regulator removal, terminal mixing,
Osterwalder--Schrader reconstruction, or the Yang--Mills mass gap.
\end{firewall}

\subsection{Parent record}

This V13 note was formed by appending the preceding material to the
validated compact V12, whose record was
\[
\begin{array}{c|c}
\text{lines}&5970\\
\text{bytes}&206747\\
\text{SHA--256}&
\texttt{9041D8F8330E8D0E9B1BE4CF771E5EFF02EDBB359D74441B181E694AFF16AB2E}.
\end{array}
\]

% =============================================================================
% END APPENDED COMPACT-SERIES V13 MATERIAL
% =============================================================================

% =============================================================================
% BEGIN APPENDED COMPACT-SERIES V14 MATERIAL
% =============================================================================

\section{Compact-series V14: averaged determinant matching and the
sublinear reset criterion}
\label{sec:n1v14}

\subsection{Why absolute Cauchy control is stronger than necessary}

The legacy one-loop analysis proved a contact-null Cauchy telescope:
absolute summability of adjacent weighted-kernel differences implies
convergence of the one-step beta coefficient.  That theorem remains
useful for full locality, but the scalar ultraviolet law needs less.
Indeed, if
\[
                         \beta_{j+1}=\beta_j+c_j,              \tag{14.1}
\]
then only the average of the increments enters \(\beta_N/N\).
V14 records the exact weaker condition and applies the reciprocal
shielding of V13 to the nonlinear calibrated rows.

\begin{lemma}[Cesaro marginal criterion]
\label{lem:n1v14-cesaro}
Let \(b>0\), let \(\beta_0>0\), and suppose
\[
                         \beta_{j+1}
                         =\beta_j+b+d_j.                       \tag{14.2}
\]
If
\[
                         \frac1N\sum_{j=0}^{N-1}d_j
                         \longrightarrow0,                    \tag{14.3}
\]
and \(\beta_j>0\), then
\[
                         \frac{\beta_N}{N}\longrightarrow b,
 \qquad
                         u_N=\beta_N^{-1}\sim\frac1{bN}.        \tag{14.4}
\]
\end{lemma}

\begin{proof}
Summing (14.2) gives
\[
 \frac{\beta_N}{N}
 =
 \frac{\beta_0}{N}+b+
 \frac1N\sum_{j=0}^{N-1}d_j.                                  \tag{14.5}
\]
The first and last terms tend to zero.  The limit is positive, so
taking reciprocals gives the second assertion.
\end{proof}

\begin{corollary}[Pointwise decay is sufficient, not necessary]
\label{cor:n1v14-pointwise}
Condition \(d_j\to0\) implies (14.3).  Absolute summability
\(\sum_j|d_j|<\infty\) also implies it, but is not necessary.
\end{corollary}

\begin{proof}
The first implication is the Cesaro theorem.  For the second,
\[
 \left|\frac1N\sum_{j<N}d_j\right|
 \le\frac1N\sum_{j\ge0}|d_j|\longrightarrow0.
\]
For nonnecessity, \(d_j=(j+1)^{-1/2}\) has averages tending to zero
but is not summable.
\end{proof}

\subsection{Exact cumulative Gaussian ownership}

Let \(\Gamma_{{\rm G},j}(W)\) be the measured one-loop Gaussian
shell functional at scale \(j\), including the Haar density or its
equivalent compensated-Haar vertices.  Let
\(\mathfrak b\) denote the normalized transverse \(F^2\) extractor.
Put
\[
                         b_{{\rm G},j}
                         :=\mathfrak b\Gamma_{{\rm G},j}.       \tag{14.6}
\]
For \(N\) nested exact shell eliminations, let
\(\Gamma_{\rm G}^{[0,N)}\) be the fine-minus-final-coarse measured
Gaussian determinant.

\begin{proposition}[Exact determinant sum before reference reset]
\label{prop:n1v14-det-sum}
If all \(N\) eliminations use the same exact nested Gaussian
quadratic form and its induced measured quotient, then
\[
 \Gamma_{\rm G}^{[0,N)}
 =
 \sum_{j=0}^{N-1}\Gamma_{{\rm G},j}+
 C_N,                                                         \tag{14.7}
\]
where \(C_N\) is independent of the retained background.  Hence
\[
 \boxed{
 \mathfrak b\Gamma_{\rm G}^{[0,N)}
 =
 \sum_{j=0}^{N-1}b_{{\rm G},j}.}                              \tag{14.8}
\]
\end{proposition}

\begin{proof}
Successive Schur complements factor the determinant of a positive
finite-dimensional quadratic form into the product of its vertical
shell determinants and the determinant of the final coarse Schur
form.  Taking logarithms gives (14.7), up to normalization constants.
The measured-slice Jacobian factors in the same order because it is
the Jacobian of the composed triangular change of variables.
Every term in \(C_N\) is background independent and is annihilated
by \(\mathfrak b\), proving (14.8).
\end{proof}

\begin{remark}[No new Cauchy theorem is being asserted]
\label{rem:n1v14-no-cauchy}
Equation (14.8) is the exact determinant cocycle already established
in the legacy programme, used here at the scalar level.  It does not
say that \(b_{{\rm G},j}\) converges.  Its new use is that a linear
asymptotic for the left side suffices for the averaged coefficient.
\end{remark}

Now insert a reference reset after each exact shell.  Let
\(r_{{\rm reset},j}\) be the scalar change in the next shell
increment caused by that reset, after the same-scale moment has been
preserved.  Define the cumulative reset memory
\[
                         R_N^{\rm reset}
                         :=\sum_{j=0}^{N-1}r_{{\rm reset},j}.   \tag{14.9}
\]
Let \(r_{{\rm nl},j}\) contain the non-Gaussian compact remainder
and the displacement from the exact redundant and anisotropy
calibrations.  Then the exact inverse-coupling ledger has the form
\[
 \beta_N-\beta_0
 =
 \mathfrak b\Gamma_{\rm G}^{[0,N)}
 +R_N^{\rm reset}
 +\sum_{j=0}^{N-1}r_{{\rm nl},j}
 +R_N^{\rm bdry}.                                             \tag{14.10}
\]
The last term retains genuine spatial boundary, toron-transition,
or sector-boundary effects not already assigned to a local
saturated owner.

\begin{warning}[Every term in (14.10) has a distinct meaning]
\label{warn:n1v14-ledger}
A moment-null contact is absent from
\(\mathfrak b\Gamma_{\rm G}^{[0,N)}\) exactly.  A same-scale moment-
preserving reset may still affect a later shell and belongs to
\(R_N^{\rm reset}\).  A topology-changing collar exit belongs to a
saturated activity and hence to \(r_{{\rm nl},j}\); a global
boundary condition not represented by that local owner remains in
\(R_N^{\rm bdry}\).  Combining these terms before proving their
ownership would make the telescope circular.
\end{warning}

\subsection{The calibrated nonlinear average is already negligible}

On the compact V8--V10 domain, let
\[
 w_j:=w_{10}(\beta_j),\qquad w_j\longrightarrow0.             \tag{14.11}
\]
The unmarked connected compiler and bounded marginal extractor give
a finite constant \(C_{\rm out}\) such that the direct non-Gaussian
increment obeys
\[
                         |r_{{\rm dir},j}|
                         \le C_{\rm out}w_j.                   \tag{14.12}
\]
Let \(x_j\) denote the calibrated stable/owner and anisotropy graph.
Its conservative bound is
\[
                         \|x_j\|\le H u_j^{6/5}.               \tag{14.13}
\]

\begin{proposition}[Vanishing of calibrated inverse-coupling
contamination]
\label{prop:n1v14-nonlinear-average}
Assume the marked inverse-coupling response of V13 is uniformly
bounded by \(L_{\rm inv}\).  Then
\[
 |r_{{\rm nl},j}|
 \le
 C_{\rm out}w_j+
 L_{\rm inv}H u_j^{6/5},                                     \tag{14.14}
\]
and consequently
\[
                         r_{{\rm nl},j}\longrightarrow0,
 \qquad
 \frac1N\sum_{j=0}^{N-1}r_{{\rm nl},j}
 \longrightarrow0.                                          \tag{14.15}
\]
\end{proposition}

\begin{proof}
Equation (14.12) bounds the direct connected correction.  Integrate
the marked derivative of the inverse-coupling increment from zero
to \(x_j\) and use (14.13), giving the second term in (14.14).
Both terms tend to zero by (14.11) and \(u_j\to0\).  The second
claim in (14.15) is the Cesaro theorem.
\end{proof}

\begin{corollary}[The old dimension-six DQC values are not needed
for the leading coefficient]
\label{cor:n1v14-no-DQC}
The exact numerical irrelevant-to-marginal mixing values of the two
local dimension-six reset directions are unnecessary for the
leading beta coefficient, provided their marked response is
uniformly bounded and their calibrated coefficients tend to zero.
They remain relevant to sharper finite-scale errors and to full
action locality.
\end{corollary}

\begin{proof}
Their complete contribution is included in the second term of
(14.14), hence has zero Cesaro average.  No assertion that the
individual mixing values vanish is used.
\end{proof}

\subsection{The averaged determinant matching card}

\begin{definition}[Averaged determinant matching card]
\label{def:n1v14-ADMC}
ADMC on a cofinal regulator/volume diagonal is the pair of
sublinear estimates
\begin{align}
 \mathfrak b\Gamma_{\rm G}^{[0,N)}
 &=
 Nb_{\rm YM}+o(N),                                            \tag{14.16}\\
 R_N^{\rm reset}+R_N^{\rm bdry}
 &=o(N).                                                       \tag{14.17}
\end{align}
All determinant, Haar, quotient, toron, and normalization
conventions in (14.16) must be the ones used by the exact compact
block.
\end{definition}

\begin{theorem}[ADMC is sufficient for the asymptotically free
marginal orbit]
\label{thm:n1v14-ADMC}
Assume ADMC, the calibrated compact domain, and the bounded marked
response of V13.  Then
\[
                         \frac{\beta_N}{N}\longrightarrow
                         b_{\rm YM},\qquad
                         u_N\sim\frac1{b_{\rm YM}N}.            \tag{14.18}
\]
\end{theorem}

\begin{proof}
Divide (14.10) by \(N\).  Equation (14.16) gives
\(b_{\rm YM}+o(1)\), (14.17) gives \(o(1)\), and
Proposition~\ref{prop:n1v14-nonlinear-average} makes the nonlinear
average \(o(1)\).  Also \(\beta_0/N\to0\).  The conclusion follows
as in Lemma~\ref{lem:n1v14-cesaro}.
\end{proof}

\begin{corollary}[Relation to the legacy Cauchy card]
\label{cor:n1v14-old-stronger}
Absolute summability of every scalar adjacent-scale reset and
kernel difference implies ADMC, but ADMC permits nonsummable errors.
For example, a reset memory
\[
                         r_{{\rm reset},j}=O(j^{-1/2})         \tag{14.19}
\]
has cumulative size \(O(N^{1/2})=o(N)\) and is admissible for
(14.18).
\end{corollary}

\begin{proof}
Absolute summability gives bounded cumulative errors, hence
sublinear ones.  Summing (14.19) gives \(O(N^{1/2})\) by comparison
with its integral.
\end{proof}

\subsection{The exact surviving coefficient obstruction}

The prior compact finite-card theorem supplies an absolutely
convergent one-shell pure-Wilson moment
\[
 b_{\rm tree}
 =
 -\frac32\sum_{x\in{\mathbb Z}^4}x_1^2j(x),                   \tag{14.20}
\]
and proves convergence of finite quotient cards to this number.
The continuum background-field calculation supplies the candidate
\[
 b_{\rm YM}
 =
 \frac{22C_A}{3(4\pi)^2}\log2                                \tag{14.21}
\]
in the upward ultraviolet orientation.  What has not been proved is
the equality of the compact measured determinant anomaly with
(14.21), either one shell at a time or in the weaker averaged form
(14.16).

\begin{proposition}[Minimal Gaussian alternative]
\label{prop:n1v14-gaussian-alternative}
Any one of the following is sufficient for (14.16):
\begin{enumerate}
\item \(b_{{\rm G},j}\to b_{\rm YM}\);
\item
      \[
       \sum_{j<N}(b_{{\rm G},j}-b_{\rm YM})=o(N);
                                                                    \tag{14.22}
      \]
\item a decomposition
      \[
       b_{{\rm G},j}-b_{\rm YM}
       =a_{j+1}-a_j+e_j,                                      \tag{14.23}
      \]
      with \(a_N=o(N)\) and
      \(N^{-1}\sum_{j<N}e_j\to0\).
\end{enumerate}
\end{proposition}

\begin{proof}
The first implies the second by Cesaro.  The second is (14.16)
after (14.8).  Under the third,
\[
 \sum_{j<N}(b_{{\rm G},j}-b_{\rm YM})
 =
 a_N-a_0+\sum_{j<N}e_j=o(N),
\]
which is the second condition.
\end{proof}

\begin{assessment}[The named remaining terms]
\label{ass:n1v14-remaining}
The marginal programme no longer needs a summable norm for every
local generated reset.  Two precise limits remain:
\begin{enumerate}
\item the \emph{compact Gaussian anomaly}: prove (14.16), equivalently
      one of the alternatives in
      Proposition~\ref{prop:n1v14-gaussian-alternative}, for the
      parity-tree/Haar-compensated determinant on the physical
      dyadic diagonal; and
\item the \emph{propagated global reset boundary}: prove (14.17) for
      the part not already represented by the vanishing local
      calibrated graph, including any toron/flux or chart-overlap
      boundary response.
\end{enumerate}
The first is an ultraviolet anomaly-matching calculation.  The
second is a sublinear boundary theorem.  Neither may be replaced by
the same-scale statement that the reset preserves the classical
quadratic moment.
\end{assessment}

\begin{decision}[Mandatory V15 audit and direction]
\label{dec:n1v14-next}
Before the next calculation, inspect the current SM and
Navier--Stokes companion notes as required at every fifth compact
version.  Look specifically for:
\begin{enumerate}
\item a determinant-cocycle theorem producing a boundary coboundary
      of the form (14.23);
\item a sublinear rather than summable reference-memory estimate;
      or
\item a uniform global-sector boundary estimate.
\end{enumerate}
After the audit, attack the first transferable item; otherwise
return to the compact Gaussian anomaly and build an explicit
low/edge homotopy whose derivative is a moment-null contact plus a
Cesaro-small boundary term.
\end{decision}

\begin{firewall}[Claim boundary after compact V14]
\label{fire:n1v14-claim}
V14 proves that averaged, rather than absolutely summable,
determinant matching is sufficient for the asymptotically free
marginal orbit.  It proves that all locally calibrated nonlinear
rows have zero leading Cesaro contribution under the V13 marked
bound.  It does not prove ADMC (14.16)--(14.17), construct the full
coupled orbit, remove regulators, establish terminal mixing or
Osterwalder--Schrader reconstruction, or prove the Yang--Mills mass
gap.
\end{firewall}

\subsection{Parent record}

This V14 note was formed by appending the preceding material to the
validated compact V13, whose record was
\[
\begin{array}{c|c}
\text{lines}&6375\\
\text{bytes}&220149\\
\text{SHA--256}&
\texttt{C550C93F0488FBB5D364F9E512245EB3941AAB59130F94B4FB591FAF40A4EF06}.
\end{array}
\]

% =============================================================================
% END APPENDED COMPACT-SERIES V14 MATERIAL
% =============================================================================

% =============================================================================
% BEGIN APPENDED COMPACT-SERIES V15 MATERIAL
% =============================================================================

\section{Compact-series V15: compulsory companion audit and
sublinear measured-conjugacy invariance}
\label{sec:n1v15}

\subsection{Companion files inspected}

The live unarchived companions at this fifth compact version were
\[
\begin{array}{c|r|r|l}
\text{file}&\text{lines}&\text{bytes}&\text{SHA--256}\\ \hline
\texttt{Renewal\_SM\_Einstein\_Continuum\_Research\_Notes\_V95.tex}
&33270&1206780&
\texttt{F38AE6C577B8A31B907AB1611173B5E1D024301E0415C370E1731696FBBA43BD}
\\
\texttt{Renewal\_NS\_Stability\_Research\_Notes\_V26.tex}
&5319&278283&
\texttt{82C887F8C2C69E4F28FAD7C3DC8DE5B9099CB5A4BF431EBA1FFF3E91935978AD}.
\end{array}                                                     \tag{15.1}
\]

The relevant SM material is its V90--V95 block.  It independently
confirms three structural points used here: an exact all-colour
disintegration is not a repeated Gibbs sweep; measured-coordinate
directions must be removed before contracting irrelevant jets; and
topological labels require sectorwise projective limits plus
discrete-label tightness.  It also excludes anomalies, determinant
lines, and topological classes from the redundant quotient unless
their exact triviality has been proved.  These are consistency
checks, not new numerical Yang--Mills estimates.

The NS file remains at V26.  Its rank-one Oseen-kernel calculation
has no compact determinant, reference reset, topological sector, or
reflection-positive quantum measure.  No constant or theorem from
that calculation transfers to the V14 anomaly problem.

\begin{assessment}[V15 audit decision]
\label{ass:n1v15-audit}
Neither companion supplies the determinant coboundary, a sublinear
reference-memory estimate, or a uniform Yang--Mills sector-boundary
bound requested in V14.  The SM note strengthens the firewall:
topological and anomalous directions must remain retained while the
comparison is made.  The Yang--Mills direction therefore stays on
the compact Gaussian anomaly, now at the weaker averaged level.
\end{assessment}

\subsection{One-loop cocycles and changes of regulator coordinates}

Let \({\cal G}\) be a real vector space of background-dependent
quadratic germs and let
\[
                         \mathfrak b:{\cal G}\longrightarrow
                         {\mathbb R}                            \tag{15.2}
\]
be the normalized transverse \(F^2\) extractor.  Constants
independent of the retained background lie in
\(\ker\mathfrak b\).

Consider two regulator towers, with measured one-loop shell
functionals
\[
                         \gamma_j,\widetilde\gamma_j\in{\cal G}.
                                                                    \tag{15.3}
\]
The towers need not have convergent individual shell coefficients.

\begin{definition}[Measured comparison cochain]
\label{def:n1v15-cochain}
A sequence \(A_j\in{\cal G}\) is a measured comparison cochain if
\[
 \widetilde\gamma_j-\gamma_j
 =
 A_{j+1}-A_j+E_j+C_j,\qquad
                         \mathfrak b(C_j)=0.                   \tag{15.4}
\]
Here \(E_j\) is the genuine failure of the finite-scale
intertwining; \(A_{j+1}-A_j\) is an endpoint rechart, and \(C_j\)
is a background-independent or exactly moment-null contact.
\end{definition}

\begin{theorem}[Sublinear measured-conjugacy invariance]
\label{thm:n1v15-sublinear-conjugacy}
Suppose (15.4) holds and
\begin{align}
                         \mathfrak b(A_N)&=o(N),               \tag{15.5}\\
 \frac1N\sum_{j=0}^{N-1}\mathfrak b(E_j)&\longrightarrow0.     \tag{15.6}
\end{align}
Then
\[
 \frac1N\sum_{j=0}^{N-1}
 \mathfrak b(\widetilde\gamma_j-\gamma_j)
 \longrightarrow0.                                          \tag{15.7}
\]
Consequently the two towers have the same averaged one-loop
coefficient whenever either average exists.
\end{theorem}

\begin{proof}
Apply \(\mathfrak b\) to (15.4) and sum:
\[
 \sum_{j=0}^{N-1}
 \mathfrak b(\widetilde\gamma_j-\gamma_j)
 =
 \mathfrak b(A_N)-\mathfrak b(A_0)
 +\sum_{j=0}^{N-1}\mathfrak b(E_j).                            \tag{15.8}
\]
The \(C_j\) terms vanish exactly.  Divide by \(N\) and use
(15.5)--(15.6).
\end{proof}

\begin{corollary}[Bounded and logarithmic endpoint recharting]
\label{cor:n1v15-endpoint-growth}
Condition (15.5) holds if either
\[
 \sup_N|\mathfrak b(A_N)|<\infty
 \quad\hbox{or}\quad
 |\mathfrak b(A_N)|\le C\log(2+N).                             \tag{15.9}
\]
No adjacent-scale absolute summability is required.
\end{corollary}

\begin{proof}
Both bounds in (15.9) are \(o(N)\).
\end{proof}

\begin{corollary}[Coboundaries may have nonsummable increments]
\label{cor:n1v15-nonsummable}
Let \(\mathfrak b(A_N)=N^{1/2}\) and \(E_j=0\).  Then the increments
\[
 \mathfrak b(A_{j+1}-A_j)
 =
 (j+1)^{1/2}-j^{1/2}                                         \tag{15.10}
\]
are not absolutely summable, but (15.7) still holds.
\end{corollary}

\begin{proof}
The partial sum of (15.10) is \(N^{1/2}\), whose quotient by \(N\)
tends to zero.
\end{proof}

\subsection{Derivation from an exact finite-scale intertwiner}

Let \(T_j\) and \(\widetilde T_j\) be exact Gaussian shell
eliminations.  Suppose \(C_j\) is an invertible background-dependent
linear change of the complete fine-plus-retained variables and
\[
                         \widetilde T_j C_j
                         =
                         C_{j+1}T_j                            \tag{15.11}
\]
on their common finite-dimensional measured domain.  Let
\({\cal A}_j\) be the measured one-loop functional of \(C_j\):
the half logarithm of its induced Hessian determinant minus the
logarithm of its measure Jacobian.

\begin{theorem}[Exact intertwining produces a coboundary]
\label{thm:n1v15-intertwiner}
Under (15.11), the shell functionals satisfy
\[
                         \widetilde\gamma_j-\gamma_j
                         =
                         {\cal A}_{j+1}-{\cal A}_j+C_j^{0},    \tag{15.12}
\]
where \(C_j^0\) is independent of the retained background.
Therefore \({\cal A}_j\) is a measured comparison cochain with
\(E_j=0\).
\end{theorem}

\begin{proof}
Compute the Gaussian integral around the retained background in two
orders.  Along the lower route in the commuting square (15.11),
first change variables by \(C_j\) and then eliminate the shell.
Along the upper route, first eliminate the shell and then change the
retained variables by \(C_{j+1}\).  Equality of the resulting
finite-dimensional integrals implies equality of their one-loop
coefficients.  Moving the two change-of-variable contributions to
opposite sides gives (15.12).  Normalizing each Gaussian integral
may add a background-independent constant \(C_j^0\), but no other
term.
\end{proof}

\begin{corollary}[Pure measured-slice changes have zero cochain]
\label{cor:n1v15-slice-zero}
If \(C_j\) merely changes a gauge complement and its exact orbit
Jacobian is included, then
\[
                         {\cal A}_j=0                          \tag{15.13}
\]
for every \(j\).  Hence all background jets, including the averaged
anomaly, are exactly slice independent.
\end{corollary}

\begin{proof}
On a complement, the Hessians obey
\[
                         \widetilde H=C_j^*HC_j.
\]
Thus
\[
 \frac12\log\det\widetilde H
 =
 \frac12\log\det H+\log|\det C_j|.
\]
The induced measure Jacobian changes its logarithm by the same
\(\log|\det C_j|\).  Their difference is zero, proving (15.13).
\end{proof}

\begin{corollary}[The V9 measured redundant section is anomaly
neutral]
\label{cor:n1v15-v9-neutral}
The exact V9 measured-coordinate correction cannot contribute to
\(A_N\), \(E_j\), or the averaged anomaly.
\end{corollary}

\begin{proof}
V9 identifies the original and corrected actions as the same
measured pushforward, including the Haar divergence.  It is
therefore an exact measured intertwiner with zero physical
functional, as already proved at the marginal level in
Proposition~\ref{prop:n1v12-redundant-zero}.
\end{proof}

\subsection{Approximate intertwining and a usable norm bound}

Frequently the two towers agree only up to a local remainder.
Let \(\|\cdot\|_{\cal G}\) be any germ norm for which
\[
                         |\mathfrak b K|
                         \le C_{\mathfrak b}\|K\|_{\cal G}.    \tag{15.14}
\]

\begin{proposition}[Norm-small intertwining defect]
\label{prop:n1v15-defect}
If (15.4) holds and
\[
 \frac1N\sum_{j=0}^{N-1}\|E_j\|_{\cal G}
 \longrightarrow0,                                          \tag{15.15}
\]
then (15.6) holds.
\end{proposition}

\begin{proof}
By (15.14),
\[
 \left|
 \frac1N\sum_{j<N}\mathfrak b(E_j)
 \right|
 \le
 \frac{C_{\mathfrak b}}N\sum_{j<N}\|E_j\|_{\cal G},
\]
and the right side tends to zero.
\end{proof}

\begin{corollary}[Vanishing local mismatch is enough]
\label{cor:n1v15-vanishing-defect}
A uniform estimate
\[
                         \|E_j\|_{\cal G}
                         \le e_j,\qquad e_j\longrightarrow0   \tag{15.16}
\]
implies (15.15), even when \(\sum_je_j=\infty\).
\end{corollary}

\begin{proof}
Apply the Cesaro theorem to \(e_j\).
\end{proof}

\subsection{Sectorwise form}

Let \(\gamma_{j,h,n}\) and \(\widetilde\gamma_{j,h,n}\) be the
shell functionals on a regular retained flux/charge sector.

\begin{theorem}[Uniform sectorwise conjugacy]
\label{thm:n1v15-sectorwise}
Suppose (15.4)--(15.6) hold in every retained sector, uniformly in
all labels occurring on the cofinal domain:
\[
 \sup_{h,n}\frac{|\mathfrak b(A_{N,h,n})|}{N}\longrightarrow0,
                                                                    \tag{15.17}
\]
\[
 \sup_{h,n}
 \left|
 \frac1N\sum_{j<N}\mathfrak b(E_{j,h,n})
 \right|\longrightarrow0.                                    \tag{15.18}
\]
Then (15.7) holds uniformly in \((h,n)\).  At theta zero, any
positive sector mixture whose total mass is finite has the same
averaged anomaly, provided the normalized sector weights are used.
\end{theorem}

\begin{proof}
The proof of Theorem~\ref{thm:n1v15-sublinear-conjugacy} is uniform
under (15.17)--(15.18).  If \(p_{h,n}\ge0\) and
\(\sum p_{h,n}=1\), the absolute value of the weighted average is
bounded by the supremum over sectors.  This proves the last claim.
\end{proof}

\begin{warning}[Sector summability remains an independent card]
\label{warn:n1v15-sector-sum}
Theorem~\ref{thm:n1v15-sectorwise} does not prove that the infinite-
volume sector weights have finite total mass or are tight in
\((h,n)\).  It says only that, if a normalized positive zero-theta
mixture exists, a uniform conjugacy estimate passes to that mixture.
This agrees with the SM V95 audit.
\end{warning}

\subsection{The compact-to-continuum comparison card}

Let \(\gamma_j^{\rm cmp}\) be the exact parity-tree,
Haar-compensated compact Gaussian shell and let
\(\gamma_j^{\rm cov}\) be a covariant background regulator for
which the continuum calculation gives
\[
                         \mathfrak b\gamma_j^{\rm cov}
                         =b_{\rm YM}+e_j^{\rm cov},\qquad
 \frac1N\sum_{j<N}e_j^{\rm cov}\longrightarrow0.              \tag{15.19}
\]

\begin{definition}[Sublinear compact--continuum intertwiner card]
\label{def:n1v15-SCIC}
SCIC consists of a measured comparison cochain \(A_j\), defects
\(E_j\), and moment-null contacts \(C_j\) such that
\[
 \gamma_j^{\rm cmp}-\gamma_j^{\rm cov}
 =
 A_{j+1}-A_j+E_j+C_j,                                        \tag{15.20}
\]
with (15.5)--(15.6), uniformly on the physical dyadic diagonal and
on every retained sector admitted to the final mixture.
\end{definition}

\begin{theorem}[SCIC closes the compact Gaussian anomaly]
\label{thm:n1v15-SCIC}
If SCIC and (15.19) hold, then
\[
 \frac1N\sum_{j=0}^{N-1}
 \mathfrak b\gamma_j^{\rm cmp}
 \longrightarrow b_{\rm YM}.                                 \tag{15.21}
\]
Thus the Gaussian half of ADMC (14.16) holds.
\end{theorem}

\begin{proof}
Apply Theorem~\ref{thm:n1v15-sublinear-conjugacy} to (15.20).
It identifies the compact and covariant averages.  Equation (15.19)
evaluates the latter.
\end{proof}

\begin{assessment}[The post-audit marginal frontier]
\label{ass:n1v15-frontier}
The measured gauge-slice change and the V9 redundant section are
already exact zero-cochain transformations.  V13--V14 remove
vanishing local stable and anisotropy graphs from the leading
average.  The unresolved Gaussian problem is therefore not a
generic scheme-independence slogan: it is the construction of
(15.20) between the exact compact parity-tree determinant and a
covariant continuum-computable tower, with sublinear endpoint
cochain and Cesaro-small defect.  Any leftover topological boundary
must satisfy the uniform sector version (15.17)--(15.18).
\end{assessment}

\begin{decision}[V16 target]
\label{dec:n1v15-next}
Construct SCIC first on the contractible flat sector.  Use a
low/edge partition of the exact \(45\times45\) compact symbol, but
apply the Ward extractor only after recombining the partition.
The target is:
\[
 \gamma_j^{\rm cmp}-\gamma_j^{\rm cov}
 =
 A_{j+1}-A_j+C_j+E_j,\qquad
 \mathfrak b(C_j)=0,                                         \tag{15.22}
\]
with \(A_j=O(1+j^\alpha)\) for some \(\alpha<1\) and
\(\|E_j\|_{\cal G}\to0\).  If the Brillouin edge gives an
order-one noncoboundary term, retain its exact scalar rather than
calling it a contact.
\end{decision}

\begin{firewall}[Claim boundary after compact V15]
\label{fire:n1v15-claim}
V15 performs the required companion audit and proves invariance of
the averaged anomaly under exact or Cesaro-accurate measured
conjugacies with sublinear endpoint cochain.  It reduces compact
Gaussian matching to SCIC.  It does not construct SCIC, prove
sector-label tightness, complete ADMC, construct the coupled RG
orbit or continuum measure, prove terminal clustering or
Osterwalder--Schrader reconstruction, or establish the
Yang--Mills mass gap.
\end{firewall}

\subsection{Parent record}

This V15 note was formed by appending the preceding material to the
validated compact V14, whose record was
\[
\begin{array}{c|c}
\text{lines}&6773\\
\text{bytes}&234163\\
\text{SHA--256}&
\texttt{CF0973C459AB1535394F0FBC5018FC285962BAD8A696500C0DADED2BA4285B57}.
\end{array}
\]

% =============================================================================
% END APPENDED COMPACT-SERIES V15 MATERIAL
% =============================================================================

% =============================================================================
% BEGIN APPENDED COMPACT-SERIES V16 MATERIAL
% =============================================================================

\section{Compact-series V16: a Ward-recombined low/edge compiler
for the averaged anomaly}
\label{sec:n1v16}

\subsection{The split must be made after the Ward sum}

Let \(h_j=2^{-j}\) and use physical loop momentum on the periodic
torus
\[
 {\mathbb T}_{h_j}^4
 :=
 \left(-\frac{\pi}{h_j},\frac{\pi}{h_j}\right]^4.             \tag{16.1}
\]
All opposite faces are identified.  Let
\(\chi_{{\rm L},j}+\chi_{{\rm E},j}=1\) be a smooth periodic
partition into a low region and its Brillouin-edge complement.

\begin{lemma}[Periodic Ward divergence]
\label{lem:n1v16-periodic-div}
If \(V\in W^{1,1}({\mathbb T}_{h_j}^4;{\mathbb R}^4)\) is
periodic, then
\[
                         \int_{{\mathbb T}_{h_j}^4}
                         \nabla_k\mathbin{\cdot}V(k)\,dk=0.    \tag{16.2}
\]
\end{lemma}

\begin{proof}
Integrate each \(\partial_{k_\nu}V_\nu\) along the \(\nu\)-th
coordinate.  The two endpoint traces agree by periodicity and
cancel.  Fubini's theorem applies to the \(W^{1,1}\) representative.
Sum over \(\nu\).
\end{proof}

\begin{lemma}[Exact cancellation of the artificial low/edge
boundary]
\label{lem:n1v16-low-edge-cancel}
Under the hypotheses of
Lemma~\ref{lem:n1v16-periodic-div},
\[
\begin{split}
 \int\chi_{{\rm L},j}\nabla_k\mathbin{\cdot}V
 &=-\int(\nabla_k\chi_{{\rm L},j})\mathbin{\cdot}V,\\
 \int\chi_{{\rm E},j}\nabla_k\mathbin{\cdot}V
 &=+\int(\nabla_k\chi_{{\rm L},j})\mathbin{\cdot}V.            \tag{16.3}
\end{split}
\]
Thus the two cutoff-boundary terms cancel exactly after
recombination.
\end{lemma}

\begin{proof}
Periodic integration by parts gives the first identity.  Since
\(\chi_{{\rm E},j}=1-\chi_{{\rm L},j}\), its gradient is the
negative of the low gradient, giving the second.  Their sum is
(16.2).
\end{proof}

\begin{warning}[Separate low and edge absolute values can create a
false anomaly]
\label{warn:n1v16-false-edge}
Neither integral in (16.3) need be small.  Estimating them
separately by absolute value loses their exact cancellation and can
manufacture an order-one edge term.  The complete tadpole, bubble,
Haar, router, toron, and overlap response must first be placed in
one periodic Ward sum.
\end{warning}

\subsection{The exact compact response density}

Along a \(C^1\) regulator homotopy \(t\in[0,1]\), let
\(A_{j,t}(k)>0\) be the \(45\times45\) vertical symbol and
\[
                         G_{j,t}(k)=A_{j,t}(k)^{-1}.           \tag{16.4}
\]
Let \(Q_{j,t}(k,q)\) and \(R_{j,t}(k,q)\) be the complete
Haar-compensated quadratic and linear background vertices,
including the total harmonic-router derivative.  Define
\[
\begin{split}
 \Omega_{j,t}(k,q)
 :={}&
 \frac12\operatorname{Tr}\!\left[
 G_{j,t}(k)Q_{j,t}(k,q)\right]\\
 &-\frac13\operatorname{Tr}\!\left[
 G_{j,t}(k+q\widehat e_1)R_{j,t}(k,q)
 G_{j,t}(k)R_{j,t}(k,q)^*\right].
                                                               \tag{16.5}
\end{split}
\]
All finite Haar-density terms not absorbed into \(Q\) are appended
to the first line of (16.5).  The Ward-recombined anomaly density is
\[
 W_{j,t}(k)
 :=
 \frac32
 \left.\partial_q^2\right|_{q=0}\Omega_{j,t}(k,q).             \tag{16.6}
\]

\begin{proposition}[The response card is a classical periodic
integrand]
\label{prop:n1v16-response-card}
Suppose the vertical spectral floor is uniform along the homotopy
and all incidence and vertex symbols are trigonometric
polynomials.  Then \(W_{j,t}\) and \(\partial_tW_{j,t}\) are
smooth periodic functions of \(k\), the toron value \(k=0\) is
included, and
\[
 \partial_t\int_{{\mathbb T}_{h_j}^4}W_{j,t}(k)\,dk
 =
 \int_{{\mathbb T}_{h_j}^4}\partial_tW_{j,t}(k)\,dk.           \tag{16.7}
\]
\end{proposition}

\begin{proof}
The spectral floor makes matrix inversion smooth with
\[
                         \partial_tG=-G(\partial_tA)G.         \tag{16.8}
\]
Products, adjoints, traces, and the two external derivatives
preserve smooth periodicity.  The domain is compact for fixed
\(j\), so \(\partial_tW\) is bounded on every closed homotopy
interval.  Differentiation under the finite-volume integral follows
from dominated convergence.  No deletion of \(k=0\) occurs in this
argument.
\end{proof}

\begin{remark}[Six response words, one Ward sum]
\label{rem:n1v16-six-words}
Before the two external derivatives are taken, differentiating
(16.5) produces
\[
\begin{split}
 \partial_t\Omega={}&
 \frac12\operatorname{Tr}[\dot GQ+G\dot Q]\\
 &-\frac13\operatorname{Tr}[
 \dot G_+RGR^*
G_+\dot RGR^*
G_+R\dot G R^*
G_+RG\dot R^*].                                               \tag{16.9}
\end{split}
\]
The dot on a vertex is total: it includes the router and stationary-
centre derivatives.  Formula (16.9) is not split into six positive
majorants before applying the finite-lattice Ward identity.
\end{remark}

\subsection{A sufficient Ward-homotopy decomposition}

\begin{definition}[Ward-recombined homotopy card]
\label{def:n1v16-WRHC}
WRHC at scale \(j\) is a decomposition
\[
 \partial_tW_{j,t}(k)
 =
 \nabla_k\mathbin{\cdot}V_{j,t}(k)
 +D_{j+1,t}(k)-D_{j,t}(k)
 +R_{j,t}^{\rm hom}(k)
 +C_{j,t}(k),                                                 \tag{16.10}
\]
where:
\begin{enumerate}
\item \(V_{j,t}\) is periodic and belongs to \(W^{1,1}\);
\item \(D_{j,t}\) is an integrable endpoint-rechart density;
\item
      \[
       \int C_{j,t}(k)\,dk=0                                  \tag{16.11}
      \]
      by an exact moment-null contact or finite Ward identity; and
\item \(R_{j,t}^{\rm hom}\) is the genuine comparison remainder.
\end{enumerate}
\end{definition}

\begin{theorem}[WRHC compiles to SCIC]
\label{thm:n1v16-WRHC}
Suppose WRHC holds between a covariant endpoint \(t=0\) and the
compact endpoint \(t=1\).  Define
\[
 a_j
 :=
 \int_0^1\int_{{\mathbb T}_{h_j}^4}
 D_{j,t}(k)\,dk\,dt,                                         \tag{16.12}
\]
\[
 e_j
 :=
 \int_0^1\int_{{\mathbb T}_{h_j}^4}
 R_{j,t}^{\rm hom}(k)\,dk\,dt.                               \tag{16.13}
\]
Then their one-loop coefficients obey
\[
 b_j^{\rm cmp}-b_j^{\rm cov}
 =
 a_{j+1}-a_j+e_j.                                             \tag{16.14}
\]
If \(a_N=o(N)\) and
\[
                         \frac1N\sum_{j<N}e_j\longrightarrow0,
                                                                    \tag{16.15}
\]
then the compact and covariant averaged anomalies agree.
\end{theorem}

\begin{proof}
Integrate (16.10) over the periodic torus.  The divergence vanishes
by Lemma~\ref{lem:n1v16-periodic-div}; the contact vanishes by
(16.11).  Then integrate in \(t\) and use (16.7), which gives
(16.14).  Sum over \(j\), divide by \(N\), and apply
Theorem~\ref{thm:n1v15-sublinear-conjugacy}.
\end{proof}

\begin{corollary}[Cutoff partitions do not alter WRHC]
\label{cor:n1v16-partition}
One may verify (16.10) separately on the low and edge supports
provided the two pieces use the same partition of the already
recombined Ward density.  The artificial boundary flux cancels
exactly and does not enter \(R_{j,t}^{\rm hom}\).
\end{corollary}

\begin{proof}
Multiply (16.10) by
\(\chi_{{\rm L},j}\) and \(\chi_{{\rm E},j}\), add the two
identities, and use
Lemma~\ref{lem:n1v16-low-edge-cancel}.
\end{proof}

\subsection{A physical scale window for the low remainder}

Choose \(0<\alpha<1\) and
\[
                         \Lambda_j:=h_j^{-\alpha}.             \tag{16.16}
\]
Thus
\[
                         \Lambda_j\longrightarrow\infty,
 \qquad h_j\Lambda_j\longrightarrow0.                         \tag{16.17}
\]
This permits an expanding physical low region while every lattice
argument \(h_jk\) in that region tends uniformly to zero.

\begin{lemma}[Four-dimensional low-region estimate]
\label{lem:n1v16-low-estimate}
Suppose, after the divergence, rechart, and contact pieces in
(16.10) have been removed, that for some \(s>0\)
\[
 |R_{j,t}^{\rm hom}(k)|
 \le
 \begin{cases}
 C h_j^s,&|k|\le1,\\
 C h_j^s|k|^{s-4},&1<|k|\le\Lambda_j,
 \end{cases}                                                  \tag{16.18}
\]
uniformly in \(t\).  Then
\[
 \int_0^1\int_{|k|\le\Lambda_j}
 |R_{j,t}^{\rm hom}(k)|\,dk\,dt
 \le C_s\left(h_j^s+
                    (h_j\Lambda_j)^s\right)
 \longrightarrow0.                                          \tag{16.19}
\]
\end{lemma}

\begin{proof}
The unit ball contributes at most a fixed volume times \(Ch_j^s\).
In polar coordinates in four dimensions, the annulus contributes
\[
 C|{\mathbb S}^3|h_j^s
 \int_1^{\Lambda_j}r^{s-4}r^3\,dr
 =
 \frac{C|{\mathbb S}^3|}{s}
 h_j^s(\Lambda_j^s-1).                                       \tag{16.20}
\]
Use (16.17).
\end{proof}

\begin{corollary}[Any strict physical window closes the low row]
\label{cor:n1v16-low-window}
With \(\Lambda_j=h_j^{-\alpha}\), the right side of (16.19) is
\[
                         O(h_j^s)+
                         O(h_j^{s(1-\alpha)}).                 \tag{16.21}
\]
It is summable in \(j\), although only its convergence to zero is
needed for the averaged anomaly.
\end{corollary}

\begin{proof}
Insert (16.16) into (16.19).  Both powers are positive and
\(h_j=2^{-j}\).
\end{proof}

\begin{warning}[Low Taylor control cannot decide the edge mean]
\label{warn:n1v16-edge-mean}
Lemma~\ref{lem:n1v16-low-estimate} says nothing about
\(\Lambda_j<|k|\le\pi/h_j\).  Periodicity kills a proved total
derivative there, but it does not kill an arbitrary order-one
residual.  Moreover, writing an arbitrary zero-mean function as a
divergence would be circular: verifying zero mean is precisely the
anomaly calculation.
\end{warning}

\subsection{The exact remaining edge scalar}

Define, after a chosen algebraic Ward divergence and endpoint
rechart have been extracted,
\[
 e_j^{\rm edge}
 :=
 \int_0^1
 \int_{\Lambda_j<|k|\le\pi/h_j}
 R_{j,t}^{\rm hom}(k)\,dk\,dt.                               \tag{16.22}
\]

\begin{proposition}[Three sufficient edge alternatives]
\label{prop:n1v16-edge-alternatives}
Each of the following closes (16.15):
\begin{enumerate}
\item \(e_j^{\rm edge}\to0\);
\item \(N^{-1}\sum_{j<N}e_j^{\rm edge}\to0\); or
\item
      \[
       e_j^{\rm edge}=d_{j+1}-d_j+\varepsilon_j,              \tag{16.23}
      \]
      where \(d_N=o(N)\) and
      \(N^{-1}\sum_{j<N}\varepsilon_j\to0\).
\end{enumerate}
\end{proposition}

\begin{proof}
The low contribution tends to zero by
Lemma~\ref{lem:n1v16-low-estimate}, hence has zero Cesaro average.
The first alternative implies the second.  The second is (16.15).
Under the third, the partial sum telescopes to
\(d_N-d_0+\sum_{j<N}\varepsilon_j=o(N)\).
\end{proof}

\begin{definition}[Actual compact edge card]
\label{def:n1v16-CEC}
For the \(45\times45\) response (16.9), CEC consists of:
\begin{enumerate}
\item the complete periodic Ward sum, including \(k=0\), Haar,
      router, stationary-centre, and cutoff-overlap terms;
\item an algebraically derived \(V_{j,t}\), not one obtained by
      subtracting the unknown mean;
\item a moment-null contact \(C_{j,t}\) verified before taking
      absolute values;
\item the low estimate (16.18) on one window (16.16); and
\item one of the edge alternatives in
      Proposition~\ref{prop:n1v16-edge-alternatives}.
\end{enumerate}
\end{definition}

\begin{assessment}[What V16 resolves]
\label{ass:n1v16-status}
V16 proves that a low/edge cutoff introduces no anomaly when applied
after the full Ward recombination.  It converts a differentiated
compact-to-covariant symbol homotopy into the SCIC coboundary and
gives an explicit four-dimensional low-window estimate.  The
unresolved Gaussian datum is now the scalar edge mean (16.22), not
a full weighted kernel norm and not any of the six response words
separately.
\end{assessment}

\begin{decision}[V17 target]
\label{dec:n1v16-next}
Attack CEC item 5.  Rescale the edge variable
\(\kappa=h_jk\in[-\pi,\pi]^4\) in the full Ward-recombined
expression (16.9).  Separate:
\begin{enumerate}
\item the homogeneous continuum logarithmic density;
\item periodic \(\kappa\)-divergences;
\item moment-null external contacts; and
\item the residual edge scalar.
\end{enumerate}
Determine whether the residual is an exact scale coboundary under
the dyadic alias map.  If it is not, its nonzero mean is the actual
regulator anomaly and must be evaluated rather than bounded away.
\end{decision}

\begin{firewall}[Claim boundary after compact V16]
\label{fire:n1v16-claim}
V16 proves the Ward-recombined low/edge compiler and closes its low
row under the explicit symbol estimate (16.18).  It does not verify
CEC for the actual \(45\times45\) compact symbols, evaluate the
edge scalar, complete SCIC or ADMC, construct the coupled RG orbit
or continuum measure, prove terminal clustering or
Osterwalder--Schrader reconstruction, or establish the
Yang--Mills mass gap.
\end{firewall}

\subsection{Parent record}

This V16 note was formed by appending the preceding material to the
validated compact V15, whose record was
\[
\begin{array}{c|c}
\text{lines}&7186\\
\text{bytes}&248613\\
\text{SHA--256}&
\texttt{E7A5A89DC77F506EAAE52F57AC1B9D97C31077563D9D8DE69DAE877014D9B2F2}.
\end{array}
\]

% =============================================================================
% END APPENDED COMPACT-SERIES V16 MATERIAL
% =============================================================================

% =============================================================================
% BEGIN APPENDED COMPACT-SERIES V17 MATERIAL
% =============================================================================

\section{Compact-series V17: dyadic alias cohomology and the unique
edge-anomaly scalar}
\label{sec:n1v17}

\subsection{The normalized alias transfer}

Write \({\mathbb T}^4=({\mathbb R}/2\pi{\mathbb Z})^4\) and use
normalized Haar average
\[
                         \langle f\rangle
 :=
 \frac1{(2\pi)^4}\int_{{\mathbb T}^4}f(k)\,dk.                \tag{17.1}
\]
Define the dyadic inverse-branch transfer
\[
 ({\cal P}f)(q)
 :=
 \frac1{16}
 \sum_{\sigma\in\{0,1\}^4}
 f\!\left(\frac q2+\pi\sigma\right).                         \tag{17.2}
\]
All arguments in (17.2) are read modulo \(2\pi\).

\begin{lemma}[Haar invariance of the alias transfer]
\label{lem:n1v17-Haar}
For every \(f\in L^1({\mathbb T}^4)\),
\[
                         \langle{\cal P}f\rangle
                         =\langle f\rangle.                   \tag{17.3}
\]
Moreover \({\cal P}\) is a positive contraction on
\(L^\infty\).
\end{lemma}

\begin{proof}
For a fixed \(\sigma\), put
\(k=q/2+\pi\sigma\).  The sixteen images, modulo opposite faces,
partition \({\mathbb T}^4\), and \(dq=16\,dk\) on each image.
The factor \(1/16\) in (17.2) therefore gives (17.3).
Positivity is immediate, and an average of sixteen values has
absolute value at most their largest absolute value.
\end{proof}

For
\[
                         f(k)=
 \sum_{m\in{\mathbb Z}^4}\widehat f(m)e^{im\cdot k},           \tag{17.4}
\]
the transfer has an exact arithmetic action.

\begin{lemma}[Fourier decimation]
\label{lem:n1v17-Fourier}
For every trigonometric polynomial, and hence by density whenever
the Fourier series converges absolutely,
\[
                         \widehat{{\cal P}f}(m)
                         =\widehat f(2m).                      \tag{17.5}
\]
Consequently
\[
                         \widehat{{\cal P}^nf}(m)
                         =\widehat f(2^nm).                    \tag{17.6}
\]
\end{lemma}

\begin{proof}
Apply (17.2) to \(e^{im\cdot k}\).  The branch average factors into
four sums
\[
 \frac12\sum_{\sigma_\nu=0}^1e^{i\pi m_\nu\sigma_\nu}.
\]
It is zero if some \(m_\nu\) is odd and is one if every component is
even.  In the latter case the surviving exponential is
\(e^{i(m/2)\cdot q}\), proving (17.5).  Iterate for (17.6).
\end{proof}

\subsection{Analytic contraction off the constant mode}

For \(\rho>0\), let \({\cal A}_\rho\) be the Wiener algebra with
\[
 \|f\|_\rho
 :=
 \sum_{m\in{\mathbb Z}^4}
 e^{\rho|m|_1}|\widehat f(m)|.                                \tag{17.7}
\]
Matrix-valued symbols use any fixed finite-dimensional entry norm
inside the sum.

\begin{theorem}[Supergeometric alias contraction]
\label{thm:n1v17-contraction}
Let \(f\in{\cal A}_\rho\) and \(\langle f\rangle=0\).  For every
\(0<\rho'\le\rho\) and every \(n\ge1\),
\[
 \boxed{
 \|{\cal P}^nf\|_{\rho'}
 \le
 e^{-(\rho2^n-\rho')}\|f\|_\rho.}                             \tag{17.8}
\]
\end{theorem}

\begin{proof}
The zero-average condition is \(\widehat f(0)=0\).
By (17.6),
\[
\begin{split}
 \|{\cal P}^nf\|_{\rho'}
 &=
 \sum_{m\ne0}e^{\rho'|m|_1}
                  |\widehat f(2^nm)|\\
 &\le
 \sup_{m\ne0}
 e^{\rho'|m|_1-\rho2^n|m|_1}
 \sum_{m\ne0}e^{\rho|2^nm|_1}
                  |\widehat f(2^nm)|.
\end{split}
\]
The last sum is bounded by \(\|f\|_\rho\), and
\(|m|_1\ge1\) bounds the supremum by the factor in (17.8).
\end{proof}

\begin{corollary}[Exact analytic Poisson solution]
\label{cor:n1v17-Poisson}
If \(f\in{\cal A}_\rho\) and \(\langle f\rangle=0\), then for every
\(0<\rho'\le\rho\) the series
\[
                         a:=\sum_{n=0}^\infty{\cal P}^nf       \tag{17.9}
\]
converges in \({\cal A}_{\rho'}\) and solves
\[
                         (I-{\cal P})a=f.                     \tag{17.10}
\]
The solution is unique modulo an additive constant.
\end{corollary}

\begin{proof}
The \(n=0\) term belongs to \({\cal A}_{\rho'}\), and (17.8)
makes the remaining series absolutely convergent.  Partial sums
obey
\[
 (I-{\cal P})\sum_{n=0}^N{\cal P}^nf
 =
 f-{\cal P}^{N+1}f.
\]
The last term tends to zero in \({\cal A}_{\rho'}\), proving
(17.10).

If \((I-{\cal P})g=0\), then (17.5) gives
\(\widehat g(m)=\widehat g(2m)\).  Absolute Fourier summability
forces this common value to tend to zero along
\(m,2m,4m,\ldots\) for every \(m\ne0\).  Thus only
\(\widehat g(0)\) can remain.
\end{proof}

\begin{theorem}[One-dimensional analytic alias cohomology]
\label{thm:n1v17-cohomology}
Every \(f\in{\cal A}_\rho\) has the unique decomposition
\[
\boxed{
 f=\langle f\rangle+(I-{\cal P})a,}                           \tag{17.11}
\]
where \(a\in{\cal A}_{\rho'}\) is unique after the normalization
\(\langle a\rangle=0\).  Hence
\[
 {\cal A}_\rho/
 \operatorname{ran}(I-{\cal P})
 \cong{\mathbb R}                                             \tag{17.12}
\]
through the Haar mean.
\end{theorem}

\begin{proof}
Apply Corollary~\ref{cor:n1v17-Poisson} to
\(f-\langle f\rangle\).  Haar invariance (17.3) shows that every
element of \(\operatorname{ran}(I-{\cal P})\) has zero mean.
Therefore the constant in (17.11) is forced.  The normalized
uniqueness follows from the last statement of the corollary.
\end{proof}

\begin{warning}[The decomposition does not evaluate its constant]
\label{warn:n1v17-mean}
Theorem~\ref{thm:n1v17-cohomology} proves that the mean is the only
analytic obstruction to an alias coboundary.  It does not prove
that the mean is zero.  Constructing \(a\) after first subtracting
the unknown mean cannot be cited as a proof of anomaly
cancellation.
\end{warning}

\subsection{Application to the Ward-recombined edge response}

Rescale the physical edge variable by
\[
                         \kappa=h_jk\in{\mathbb T}^4.          \tag{17.13}
\]
Let \(f_{j,t}(\kappa)\) be the fully recombined scalar response
obtained from (16.9) after the low continuum contribution,
algebraic periodic divergences, endpoint rechart, and moment-null
external contacts have been removed.  Assume
\[
                         \sup_{j,t}\|f_{j,t}\|_\rho<\infty.    \tag{17.14}
\]

\begin{proposition}[Canonical edge split]
\label{prop:n1v17-edge-split}
For every \(j,t\) and \(0<\rho'\le\rho\),
\[
 f_{j,t}(\kappa)
 =
 m_{j,t}+(I-{\cal P})a_{j,t}(\kappa),                         \tag{17.15}
\]
where
\[
                         m_{j,t}:=\langle f_{j,t}\rangle,      \tag{17.16}
\]
\[
 a_{j,t}
 =
 \sum_{n\ge0}{\cal P}^n(f_{j,t}-m_{j,t}),                     \tag{17.17}
\]
and
\[
 \sup_{j,t}\|a_{j,t}\|_{\rho'}<\infty.                        \tag{17.18}
\]
\end{proposition}

\begin{proof}
Equations (17.15)--(17.17) are
Theorem~\ref{thm:n1v17-cohomology}.  Uniformity in (17.18) follows
by summing the right side of (17.8) against the uniform stock in
(17.14).
\end{proof}

\begin{corollary}[The bounded coboundary is automatically
sublinear]
\label{cor:n1v17-sublinear}
The alias-coboundary component in (17.15) supplies a uniformly
bounded endpoint cochain in V15 and hence contributes zero to the
averaged anomaly.  The complete edge obstruction is
\[
                         \overline m_N
 :=
 \frac1N\sum_{j=0}^{N-1}\int_0^1m_{j,t}\,dt.                  \tag{17.19}
\]
\end{corollary}

\begin{proof}
The transfer potential has the uniform bound (17.18), so its
endpoint value divided by \(N\) tends to zero.  The constant mode is
fixed by \({\cal P}\) and cannot be written as a coboundary.
Integrating it in the homotopy parameter and averaging over scales
gives (17.19).
\end{proof}

\begin{theorem}[Exact edge criterion for SCIC]
\label{thm:n1v17-edge-SCIC}
Under (17.14), the edge row of SCIC holds if and only if
\[
                         \overline m_N\longrightarrow0.       \tag{17.20}
\]
If \(m_{j,t}\to m_{\infty,t}\) in \(L^1(dt)\), then this is
equivalent to
\[
                         \int_0^1m_{\infty,t}\,dt=0.           \tag{17.21}
\]
\end{theorem}

\begin{proof}
By
Proposition~\ref{prop:n1v17-edge-split} and
Corollary~\ref{cor:n1v17-sublinear}, every nonconstant analytic
mode is a bounded coboundary and has zero averaged contribution.
Thus (17.20) is necessary and sufficient.  Under the stated
\(L^1\) convergence,
\[
 \int_0^1m_{j,t}\,dt
 \longrightarrow
 \int_0^1m_{\infty,t}\,dt.
\]
Its Cesaro average has the same limit, proving the last assertion.
\end{proof}

\begin{assessment}[The edge problem is exactly one scalar]
\label{ass:n1v17-edge}
After Ward recombination and analytic rescaling, there is no
remaining infinite-dimensional edge cohomology.  Periodic
divergences, alias coboundaries, and moment-null contacts are exact
zero classes.  The sole class is the Haar mean \(m_{j,t}\).
Consequently an attempted full weighted-kernel contraction is
unnecessary for anomaly matching, although it remains relevant to
locality of the complete effective action.
\end{assessment}

\subsection{A certified finite-grid mean}

For an integer \(M\ge1\), let
\[
 {\cal K}_M
 :=
 \left\{\frac{2\pi}{M}\ell:
 \ell\in\{0,\ldots,M-1\}^4\right\}                            \tag{17.22}
\]
and define the trapezoidal mean
\[
                         \langle f\rangle_M
 :=
 \frac1{M^4}\sum_{\kappa\in{\cal K}_M}f(\kappa).               \tag{17.23}
\]

\begin{theorem}[Exponential finite-card certificate]
\label{thm:n1v17-grid}
For \(f\in{\cal A}_\rho\),
\[
 \left|\langle f\rangle_M-\langle f\rangle\right|
 \le
 \|f\|_\rho
 \left[
 \left(\frac{1+e^{-\rho M}}{1-e^{-\rho M}}\right)^4-1
 \right].                                                     \tag{17.24}
\]
In particular the error is \(O_\rho(e^{-\rho M})\).
\end{theorem}

\begin{proof}
The grid average of \(e^{im\cdot\kappa}\) is one precisely when
\(m\in M{\mathbb Z}^4\), and is zero otherwise.  Hence
\[
 \langle f\rangle_M-\langle f\rangle
 =
 \sum_{\ell\in{\mathbb Z}^4\setminus\{0\}}
 \widehat f(M\ell).                                           \tag{17.25}
\]
Since
\[
 |\widehat f(M\ell)|
 \le e^{-\rho M|\ell|_1}\|f\|_\rho,
\]
sum over \(\ell\ne0\).  The one-dimensional geometric sum is
\[
 \sum_{n\in{\mathbb Z}}e^{-\rho M|n|}
 =
 \frac{1+e^{-\rho M}}{1-e^{-\rho M}}.
\]
Taking its fourth power and deleting the zero vector gives
(17.24).
\end{proof}

\begin{corollary}[Finite matrices can decide a nonzero edge class]
\label{cor:n1v17-decision}
Suppose interval evaluation of the complete finite-grid cards gives
\[
 |\langle f_{j,t}\rangle_M|\ge\delta
                                                                    \tag{17.26}
\]
uniformly on a homotopy subinterval, while the right side of
(17.24) is at most \(\delta/2\) with a certified common analytic
stock.  Then the edge mean is nonzero there and cannot be discarded
as a contact, divergence, or alias coboundary.
\end{corollary}

\begin{proof}
The triangle inequality and (17.24) give
\(|\langle f_{j,t}\rangle|\ge\delta/2\).
Theorem~\ref{thm:n1v17-cohomology} makes every listed discarded
class mean zero.
\end{proof}

\begin{decision}[V18 target]
\label{dec:n1v17-next}
Populate the scalar card \(m_{j,t}\) from the complete
\(45\times45\) expression.  First exploit colour trace, parity,
hypercubic symmetry, and orientation reversal to reduce the owner
bank.  Then either:
\begin{enumerate}
\item prove the mean vanishes algebraically along the compact-to-
      covariant homotopy; or
\item evaluate it by the finite-grid certificate (17.24), with
      interval bounds on the common analytic stock.
\end{enumerate}
Only after this scalar is known should the programme return to the
global sector-boundary part of ADMC.
\end{decision}

\begin{firewall}[Claim boundary after compact V17]
\label{fire:n1v17-claim}
V17 proves the analytic dyadic alias cohomology, identifies the edge
obstruction with one constant Fourier mode, and gives an
exponentially convergent finite-grid certificate.  It does not
evaluate the actual \(45\times45\) edge mean, complete SCIC or
ADMC, prove sector tightness, construct the coupled RG orbit or
continuum measure, establish terminal clustering or
Osterwalder--Schrader reconstruction, or prove the
Yang--Mills mass gap.
\end{firewall}

\subsection{Parent record}

This V17 note was formed by appending the preceding material to the
validated compact V16, whose record was
\[
\begin{array}{c|c}
\text{lines}&7597\\
\text{bytes}&262005\\
\text{SHA--256}&
\texttt{73B345ADD69F25F71FFE922569949ED0E9FD34164D2FFF3A77A1C56A25C6EDFF}.
\end{array}
\]

% =============================================================================
% END APPENDED COMPACT-SERIES V17 MATERIAL
% =============================================================================

% =============================================================================
% BEGIN APPENDED COMPACT-SERIES V18 MATERIAL
% =============================================================================

\section{Compact-series V18: an explicit complex strip and a
terminating interval card for the edge mean}
\label{sec:n1v18}

\subsection{Exact finite-incidence input}

The attached reduced-fibre source constructs the parity-tree curl
matrix \(D_R(k)\), with \(96\) plaquette rows and \(45\) reduced
link columns, and proves for
\[
                         A(k)=\frac13D_R(k)^*D_R(k)            \tag{18.1}
\]
the exact real-torus floor
\[
 A(k)\succeq\frac{\lambda_*}{3}I,\qquad
 \frac{\lambda_*}{3}>
 0.039201395042156420883.                                    \tag{18.2}
\]
It also proves that the coarse-cell range of \(A\) is at most two.
For a rational stock below, fix
\[
                         \underline\lambda:=\frac{49}{1250}
                         =0.0392.                              \tag{18.3}
\]

Write the finite Laurent expansion
\[
                         A(k)=\sum_{n\in{\mathbb Z}^4}
                         A_ne^{in\cdot k}                      \tag{18.4}
\]
and define
\[
                         M_A:=\sum_n\|A_n\|_{\rm op}.          \tag{18.5}
\]
There are at most \(96\cdot4=384\) signed unit incidence entries
before multiplication.  The triangle and Frobenius bounds therefore
give the deliberately crude but literal estimate
\[
                         M_A\le\frac{384^2}{3}=49152.          \tag{18.6}
\]

\begin{lemma}[Complex perturbation of a finite Laurent symbol]
\label{lem:n1v18-Laurent}
Suppose \(A_n=0\) for \(|n|_1>R_A\).  If
\(z=k+iy\) and \(\|y\|_\infty\le\rho\), then
\[
 \|A(z)-A(k)\|_{\rm op}
 \le
 M_A\left(e^{R_A\rho}-1\right).                               \tag{18.7}
\]
\end{lemma}

\begin{proof}
For every occurring \(n\),
\[
 |e^{-n\cdot y}-1|
 \le e^{|n\cdot y|}-1
 \le e^{R_A\rho}-1.
\]
Insert this in (18.4), subtract the real value, and sum the operator
norms of the coefficients.
\end{proof}

\begin{theorem}[Explicit parity-tree inverse strip]
\label{thm:n1v18-strip}
Set
\[
 \boxed{
 \rho_A:=
 \frac12\log\left(
 1+\frac{\underline\lambda}{98304}
 \right)>0.}                                                  \tag{18.8}
\]
For every \(z=k+iy\) with
\(\|y\|_\infty\le\rho_A\), the matrix \(A(z)\) is invertible and
\[
 \boxed{
                         \|A(z)^{-1}\|_{\rm op}
                         \le\frac{2}{\underline\lambda}
                         =\frac{2500}{49}.}                   \tag{18.9}
\]
\end{theorem}

\begin{proof}
Use \(R_A=2\), (18.6), and (18.8) in (18.7):
\[
 \|A(z)-A(k)\|
 \le49152(e^{2\rho_A}-1)
 =\frac{\underline\lambda}{2}.                               \tag{18.10}
\]
The real inverse obeys
\(\|A(k)^{-1}\|\le\underline\lambda^{-1}\) by
(18.2)--(18.3).  Hence
\[
 A(z)
 =
 A(k)\left[
 I+A(k)^{-1}(A(z)-A(k))
 \right],
\]
and the norm of the perturbation in brackets is at most \(1/2\).
The Neumann series gives inverse norm at most
\[
 2\|A(k)^{-1}\|\le2\underline\lambda^{-1},
\]
which is (18.9).
\end{proof}

\begin{remark}[Size of the strip]
\label{rem:n1v18-strip-size}
The radius (18.8) is extremely conservative because (18.6) discards
all incidence orthogonality.  Its only purpose is to be a proved
positive number obtained from the literal \(96\)-by-\(45\) grammar.
A sharper coefficient sum may replace \(49152\) without changing
any theorem below.
\end{remark}

\subsection{A strip stock for the differentiated one-loop card}

For a matrix Laurent polynomial \(P(k,q)\), define its coefficient
stock and range by
\[
                         {\cal M}(P):=\sum_{m,n}\|P_{m,n}\|_{\rm op},
                                                                    \tag{18.11}
\]
\[
 R(P):=\max\{|m|_1+|n|:(m,n)\ \hbox{occurs}\}.                \tag{18.12}
\]
Then on
\(\|\operatorname{Im}k\|_\infty\le r_k\),
\(|\operatorname{Im}q|\le r_q\),
\[
                         \|P(k,q)\|
 \le{\cal M}(P)e^{R(P)(r_k+r_q)}.                             \tag{18.13}
\]
Every quantity on the right is a finite sum of norms of explicit
incidence coefficients.

Along a homotopy, define the strip stocks
\[
\begin{array}{lll}
 a_1:=\sup\|\dot A\|,&
 q_0:=\sup\|Q\|,&q_1:=\sup\|\dot Q\|,\\
 &&\\[-1.2em]
 r_0:=\sup\|R\|,&r_1:=\sup\|\dot R\|,&
 g_0:=2500/49,
\end{array}                                                    \tag{18.14}
\]
where the suprema are bounded by (18.13) on
\[
                         r_k=r_q=\frac{\rho_A}{4}.            \tag{18.15}
\]
The shifted inverse \(G(k+q\widehat e_1)\) remains within the
\(\rho_A/2\) strip, so (18.9) applies.

\begin{lemma}[Explicit supremum for the six-word response]
\label{lem:n1v18-six-bound}
Let \(\dot\Omega\) be (16.9), including any Haar term in
\(\dot Q\).  On the strip (18.15),
\[
                         |\dot\Omega(k,q)|\le B_{\dot\Omega}, \tag{18.16}
\]
where
\[
\boxed{
\begin{split}
 B_{\dot\Omega}
 :={}&
 \frac{45}{2}
 \left(g_0^2a_1q_0+g_0q_1\right)\\
 &+15\left(
 2g_0^3a_1r_0^2+
 2g_0^2r_0r_1
 \right).
\end{split}}                                                  \tag{18.17}
\]
\end{lemma}

\begin{proof}
Equation (16.8) and (18.9) give
\[
                         \|\dot G\|\le g_0^2a_1.              \tag{18.18}
\]
For \(45\times45\) matrices,
\[
                         |\operatorname{Tr}(XY\cdots Z)|
 \le45\|X\|\|Y\|\cdots\|Z\|.                                 \tag{18.19}
\]
Apply (18.18)--(18.19) to the two tadpole words and four bubble
words in (16.9).  The two differentiated inverse slots have the
same bound, as do the two differentiated vertex slots.  The
coefficients \(45/2\) and \(45/3=15\) give (18.17).
\end{proof}

\begin{theorem}[Explicit analytic stock for the edge response]
\label{thm:n1v18-response-stock}
Let
\[
 \dot W(k)
 :=
 \frac32\left.\partial_q^2\right|_{q=0}\dot\Omega(k,q).
                                                                    \tag{18.20}
\]
Then \(\dot W\) is analytic for
\(\|\operatorname{Im}k\|_\infty\le\rho_A/4\) and
\[
                         \sup|\dot W|
 \le
 B_W:=
 \frac{48B_{\dot\Omega}}{\rho_A^2}.                           \tag{18.21}
\]
Moreover, with
\[
                         \rho_W:=\frac{\rho_A}{8},             \tag{18.22}
\]
its Wiener norm obeys
\[
\boxed{
 \|\dot W\|_{\rho_W}
 \le
 B_W
 \left(
 \frac{1+e^{-\rho_W}}{1-e^{-\rho_W}}
 \right)^4.}                                                  \tag{18.23}
\]
\end{theorem}

\begin{proof}
For fixed complex \(k\) in the declared strip, use Cauchy's
second-derivative estimate on the \(q\)-circle of radius
\(\rho_A/4\):
\[
 |\partial_q^2\dot\Omega(k,0)|
 \le
 \frac{2B_{\dot\Omega}}{(\rho_A/4)^2}
 =
 \frac{32B_{\dot\Omega}}{\rho_A^2}.
\]
Multiplication by \(3/2\) gives (18.21).

Shift each Fourier coefficient of \(\dot W\) to the boundary of
the \(k\)-strip with signs opposite to its multi-index.  This gives
\[
                         |\widehat{\dot W}(m)|
 \le B_We^{-(\rho_A/4)|m|_1}.                                \tag{18.24}
\]
Multiply by \(e^{\rho_W|m|_1}\), sum, and use
\(\rho_A/4-\rho_W=\rho_W\).  The product of the four
one-dimensional geometric sums is (18.23).
\end{proof}

\begin{firewall}[Which finite stocks are still data]
\label{fire:n1v18-finite-stocks}
The inverse radius and bound (18.8)--(18.9) are already populated
from the attached exact incidence theorem.  The numbers
\(a_1,q_0,q_1,r_0,r_1\) in (18.14) are not analytic-existence
hypotheses: they are finite coefficient sums of the chosen
compact-to-covariant homotopy matrices.  That homotopy matrix bank
has not been supplied or constructed, so V18 does not assign
numerical values to those five sums.
\end{firewall}

\subsection{A terminating interval enclosure}

For a fixed homotopy parameter \(t\), evaluate the complete
Ward-recombined response \(\dot W_t\) on the \(M^4\) grid
of V17 using outward-rounded interval arithmetic.  Let
\[
                         I_{M,t}=[L_{M,t},U_{M,t}]             \tag{18.25}
\]
contain the grid mean.  Put
\[
 {\cal E}_M
 :=
 B_W
 \left(
 \frac{1+e^{-\rho_W}}{1-e^{-\rho_W}}
 \right)^4
 \left[
 \left(
 \frac{1+e^{-\rho_WM}}{1-e^{-\rho_WM}}
 \right)^4-1
 \right].                                                     \tag{18.26}
\]

\begin{theorem}[Certified edge-mean interval]
\label{thm:n1v18-interval}
Under the stocks above,
\[
\boxed{
 \langle\dot W_t\rangle
 \in
 [L_{M,t}-{\cal E}_M,\,
  U_{M,t}+{\cal E}_M].}                                       \tag{18.27}
\]
The deterministic error \({\cal E}_M\) tends exponentially to
zero as \(M\to\infty\).
\end{theorem}

\begin{proof}
The interval \(I_{M,t}\) encloses the finite grid average by
construction.  Apply
Theorem~\ref{thm:n1v17-grid} with
\(\rho=\rho_W\) and the norm bound (18.23).  The resulting error is
exactly (18.26).
\end{proof}

For the homotopy integral, partition
\[
                         0=t_0<t_1<\cdots<t_K=1.               \tag{18.28}
\]
Suppose interval evaluation on the whole panel
\([t_{\ell-1},t_\ell]\) gives a common enclosure
\[
 \langle\dot W_t\rangle
 \in J_{\ell,M}
 \quad(t_{\ell-1}\le t\le t_\ell),                            \tag{18.29}
\]
after adding (18.26).

\begin{corollary}[Certified homotopy scalar]
\label{cor:n1v18-homotopy}
The unique edge class of V17 satisfies
\[
\boxed{
 \int_0^1\langle\dot W_t\rangle\,dt
 \in
 \sum_{\ell=1}^K
 (t_\ell-t_{\ell-1})J_{\ell,M}.}                              \tag{18.30}
\]
If the interval on the right excludes zero, the proposed homotopy
does not prove compact-to-covariant anomaly equality.  If it is
contained in \([-\varepsilon,\varepsilon]\) along a cofinal family
with \(\varepsilon\to0\), the edge row of SCIC closes.
\end{corollary}

\begin{proof}
Integrate the interval inclusion (18.29) on each panel and add
Minkowski intervals.  The last two statements are respectively the
necessity and sufficiency of the constant-mode criterion
Theorem~\ref{thm:n1v17-edge-SCIC}.
\end{proof}

\begin{assessment}[What became executable]
\label{ass:n1v18-status}
V18 replaces the unquantified analytic stock in V17 by an explicit
positive strip derived from the exact parity-tree gap and a literal
incidence count.  Once a compact-to-covariant Laurent homotopy is
listed, five finite coefficient sums populate the response bound,
and (18.30) gives a terminating rigorous interval calculation of
the only edge cohomology class.  The attached source contains the
compact endpoint matrices but not this homotopy derivative bank.
\end{assessment}

\begin{decision}[V19 target]
\label{dec:n1v18-next}
Construct the missing homotopy bank without changing dimension:
interpolate the compact vertical Hessian and vertices to their
second-order covariant Taylor symbols on the same \(45\)-component
parity-tree fibre, retaining a stabilizing positive irrelevant
term so that the real spectral floor persists.  Differentiate the
Schur router and Haar-compensated vertices explicitly.  The output
must be the five Laurent coefficient sums in (18.14), together with
the exact Ward test at \(q=0\); no numerical anomaly value is to be
reported before those checks pass.
\end{decision}

\begin{firewall}[Claim boundary after compact V18]
\label{fire:n1v18-claim}
V18 proves an explicit complex inverse strip, a quantitative
analytic response norm, and a terminating interval enclosure for
the edge mean.  It does not construct the comparison homotopy,
populate its five coefficient stocks, evaluate the edge class,
complete SCIC or ADMC, prove sector tightness, construct the
continuum measure or terminal clustering, or establish the
Yang--Mills mass gap.
\end{firewall}

\subsection{Parent record}

This V18 note was formed by appending the preceding material to the
validated compact V17, whose record was
\[
\begin{array}{c|c}
\text{lines}&8011\\
\text{bytes}&274649\\
\text{SHA--256}&
\texttt{E8ED68538DE49A4F657A2DE6041F32FAB1406ACDD4A6FAE43D166D752848AA0C}.
\end{array}
\]

% =============================================================================
% END APPENDED COMPACT-SERIES V18 MATERIAL
% =============================================================================

% =============================================================================
% BEGIN APPENDED COMPACT-SERIES V19 MATERIAL
% =============================================================================

\section{Compact-series V19: a positive covariant stabilizer for
the comparison homotopy}
\label{sec:n1v19}

\subsection{A compact gauge-invariant fourth-order stabilizer}

Let \(P_{x,\mu\nu}(U)\) be a plaquette holonomy based at \(x\).
For a matrix field \(F_x\) transforming by conjugation at \(x\),
define the covariant lattice Laplacian
\[
\begin{split}
 ({\cal L}_UF)_x
 :=
 \sum_{\rho=1}^4\bigl(
 &U_{x,\rho}F_{x+\widehat e_\rho}U_{x,\rho}^{-1}\\
 &+
 U_{x-\widehat e_\rho,\rho}^{-1}
 F_{x-\widehat e_\rho}U_{x-\widehat e_\rho,\rho}
 -2F_x\bigr).                                                \tag{19.1}
\end{split}
\]
Put
\[
 {\cal V}_{\rm stab}(U)
 :=
 c_{\rm stab}
 \sum_{x,\mu<\nu}
 \left\|
 {\cal L}_U(P_{\cdot,\mu\nu}-I)_x
 \right\|_F^2,                                                \tag{19.2}
\]
where the positive constant \(c_{\rm stab}\) is fixed below.

\begin{lemma}[Symmetry and positivity of the stabilizer]
\label{lem:n1v19-stab-symmetry}
The functional (19.2) is nonnegative, gauge invariant, translation
invariant, reflection invariant, hypercubic invariant, and has
finite range.
\end{lemma}

\begin{proof}
Every term is a squared Frobenius norm.  Under a gauge
transformation, both \(P_{x,\mu\nu}-I\) and its covariant Laplacian
transform by conjugation at \(x\); the Frobenius norm is conjugation
invariant.  Summation over all sites, plaquette orientations, and
four displacement directions gives the remaining symmetries.
Only links in a fixed two-neighbourhood of a plaquette occur, so the
range is finite.
\end{proof}

At the flat connection, let \(d_1\) be the linearized plaquette curl
and let
\[
                         \omega(k):=
 4\sum_{\rho=1}^4\sin^2(k_\rho/2).                            \tag{19.3}
\]

\begin{proposition}[Exact flat stabilizer symbol]
\label{prop:n1v19-stab-symbol}
Choose \(c_{\rm stab}\) so that the unit Wilson vertical Hessian is
\[
                         A(k)=\frac13D_R(k)^*D_R(k).
\]
Then the reduced vertical Hessian of (19.2) is
\[
\boxed{
                         S(k)=\omega(k)^2A(k).}                \tag{19.4}
\]
In particular,
\[
                         S(k)\succeq0,\qquad
                         S(k)=O(|k|^4),                        \tag{19.5}
\]
so adding \(s{\cal V}_{\rm stab}\) changes no constant, linear, or
quadratic momentum jet of the reference precision.
\end{proposition}

\begin{proof}
The first variation of \(P-I\) at the flat connection is \(d_1a\).
The flat Fourier symbol of (19.1) is
\[
 \sum_\rho(e^{ik_\rho}+e^{-ik_\rho}-2)
 =-\omega(k).
\]
Therefore the quadratic form obtained by twice differentiating
(19.2) is a positive constant times
\[
                         \|\omega(k)D_R(k)a_R\|_2^2.
\]
Fixing \(c_{\rm stab}\) to match the \(1/3\) Wilson normalization
gives (19.4).  Since
\(\omega(k)=|k|^2+O(|k|^4)\), (19.5) follows.
\end{proof}

\begin{warning}[Irrelevant at the classical two-jet is not
one-loop invisible]
\label{warn:n1v19-loop-visible}
Equation (19.5) permits the stabilizer in a regulator homotopy
without changing the classical continuum kinetic term.  It does not
say that its one-loop background response has zero constant Fourier
mode.  That response is part of the V17--V18 edge scalar and must be
evaluated.
\end{warning}

\subsection{A positive endpoint from an arbitrary covariant
two-jet representative}

Let \(H(k)\) be a finite Laurent, Hermitian vertical symbol obtained
from a gauge-covariant finite-range reference action.  Assume its
flat two-jet has the desired covariant continuum normalization and
\[
                         H(0)=A(0).                            \tag{19.6}
\]
Define literal coefficient stocks
\[
 M_H:=\sum_n\|H_n\|_{\rm op},\qquad
 L_H:=\sum_n|n|_1\|H_n\|_{\rm op}.                            \tag{19.7}
\]
With \(\underline\lambda=49/1250\) from V18, put
\[
 r_0:=
 \min\left\{
 \frac{\pi}{2},
 \frac{\underline\lambda}
      {2\max\{1,L_H\}}
 \right\},                                                    \tag{19.8}
\]
\[
 \omega_0:=4\sin^2(r_0/2),\qquad
 s_*:=
 \frac{M_H+\underline\lambda/2}
      {\omega_0^2\underline\lambda}.                          \tag{19.9}
\]

\begin{theorem}[Uniformly positive stabilized endpoint]
\label{thm:n1v19-positive-endpoint}
The symbol
\[
                         H_*(k):=H(k)+s_*S(k)                 \tag{19.10}
\]
satisfies
\[
\boxed{
                         H_*(k)\succeq
                         \frac{\underline\lambda}{2}I
 \qquad(k\in{\mathbb T}^4).}                                 \tag{19.11}
\]
It has the same momentum two-jet as \(H\).
\end{theorem}

\begin{proof}
First suppose \(\|k\|_\infty\le r_0\).  From (19.7),
\[
 \|H(k)-H(0)\|
 \le
 \sum_n\|H_n\||e^{in\cdot k}-1|
 \le L_H\|k\|_\infty
 \le\frac{\underline\lambda}{2}.                             \tag{19.12}
\]
Equations (18.2)--(18.3) and (19.6) give
\(H(0)\succeq\underline\lambda I\).  Hence
\(H(k)\succeq\underline\lambda I/2\), and adding the positive
\(s_*S(k)\) preserves this bound.

Now suppose \(\|k\|_\infty\ge r_0\).  At least one component has
absolute value at least \(r_0\), so
\[
                         \omega(k)\ge\omega_0.                \tag{19.13}
\]
Also \(H(k)\succeq-M_HI\), while (18.2)--(18.3) and (19.4) give
\[
                         S(k)\succeq
                         \omega_0^2\underline\lambda I.
\]
The definition (19.9) yields (19.11).  Finally
Proposition~\ref{prop:n1v19-stab-symbol} shows that the added term
begins at fourth order in momentum.
\end{proof}

\begin{remark}[The normalization at zero may be reached by
congruence]
\label{rem:n1v19-congruence}
If a proposed covariant endpoint has \(H(0)>0\) but not (19.6), a
constant measured change of vertical variables can normalize it:
\[
 C=H(0)^{-1/2}A(0)^{1/2},\qquad
 \widetilde H(k)=C^*H(k)C.                                   \tag{19.14}
\]
The exact measure Jacobian makes this congruence anomaly neutral by
Corollary~\ref{cor:n1v15-slice-zero}.
\end{remark}

\subsection{The gapped convex homotopy}

Define
\[
                         A_t(k):=(1-t)A(k)+tH_*(k),
 \qquad0\le t\le1.                                           \tag{19.15}
\]

\begin{theorem}[Uniform real gap along the homotopy]
\label{thm:n1v19-homotopy-gap}
For all \(k\in{\mathbb T}^4\) and \(0\le t\le1\),
\[
\boxed{
                         A_t(k)\succeq
                         \frac{\underline\lambda}{2}I.}       \tag{19.16}
\]
Moreover
\[
                         \dot A_t=H+s_*S-A                    \tag{19.17}
\]
is a finite Laurent symbol with the explicit coefficient bound
\[
\boxed{
 {\cal M}(\dot A_t)
 \le
 M_H+49152+
 12582912\,s_*.}                                              \tag{19.18}
\]
\end{theorem}

\begin{proof}
The compact endpoint obeys
\(A\succeq\underline\lambda I\), and the stabilized endpoint obeys
(19.11).  Their convex combination gives (19.16).

The coefficient sum of \(\omega\) in (19.3) is sixteen, so the
coefficient sum of \(\omega^2\) is at most \(16^2=256\).
Equations (18.5)--(18.6) and (19.4) give
\[
                         {\cal M}(S)\le256\cdot49152
                         =12582912.
\]
Apply the triangle inequality to (19.17).
\end{proof}

\begin{corollary}[The first V18 homotopy stock is populated]
\label{cor:n1v19-a1}
For the homotopy (19.15), the coefficient \(a_1\) in (18.14) may
be taken to be the right side of (19.18), multiplied only by the
explicit strip exponential from (18.13).
\end{corollary}

\begin{proof}
Equation (19.17) is independent of \(t\).  Apply (18.13) with its
finite Laurent range.
\end{proof}

\subsection{Persistence on a background tube}

Let \(W\) be a retained background in a fixed chord chart and let
\(A_t(k;W)\) be the exact background vertical Hessian of the
homotopy action, including the stabilizer (19.2).  Suppose the
finite word differentiation gives
\[
 \|A_t(k;W)-A_t(k;I)\|
 \le L_Wd_{\rm ch}(W,I)                                      \tag{19.19}
\]
uniformly in \(t,k\).

\begin{proposition}[Uniform background gap]
\label{prop:n1v19-background-gap}
On
\[
                         d_{\rm ch}(W,I)
 \le\frac{\underline\lambda}{4\max\{1,L_W\}},                 \tag{19.20}
\]
one has
\[
\boxed{
                         A_t(k;W)\succeq
                         \frac{\underline\lambda}{4}I.}       \tag{19.21}
\]
\end{proposition}

\begin{proof}
The perturbation in (19.19) is at most
\(\underline\lambda/4\).  Subtract it from the flat floor
(19.16).
\end{proof}

\begin{corollary}[A common complex strip exists for the full
homotopy]
\label{cor:n1v19-common-strip}
Let \(M_{A,*}\) be the maximum Laurent coefficient sum of
\(A_t(\,\cdot\,;W)\) on (19.20), and let \(R_{A,*}\) be its
maximum range.  Then
\[
 \rho_*:=
 \frac1{R_{A,*}}
 \log\left(
 1+\frac{\underline\lambda}{8M_{A,*}}
 \right)                                                     \tag{19.22}
\]
is a common complex momentum strip, with
\[
                         \|A_t(z;W)^{-1}\|
                         \le\frac8{\underline\lambda}.        \tag{19.23}
\]
\end{corollary}

\begin{proof}
Lemma~\ref{lem:n1v18-Laurent} bounds the complex perturbation by
\(\underline\lambda/8\).  Relative to the real floor
\(\underline\lambda/4\), its inverse-relative norm is at most
\(1/2\).  The Neumann series doubles the real inverse bound
\(4/\underline\lambda\), giving (19.23).
\end{proof}

\subsection{Ward compatibility of the homotopy}

\begin{proposition}[No gauge-breaking stabilizer row]
\label{prop:n1v19-Ward}
If \(H\) is the vertical restriction of a gauge-invariant
finite-range background action, then every action in (19.15) has an
exact gauge-invariant nonlinear completion.  Its full background
Hessian obeys the finite-lattice Ward identity, and the stabilizer
adds no gauge-boson mass term.
\end{proposition}

\begin{proof}
The compact endpoint is the Wilson action and is gauge invariant.
The proposed covariant endpoint is gauge invariant by hypothesis,
and Lemma~\ref{lem:n1v19-stab-symmetry} proves the same for
(19.2).  Linear combinations preserve gauge invariance.  Twice
differentiating the exact gauge invariance gives the background Ward
identity.  At zero momentum the stabilizer Hessian vanishes by
(19.4), so it cannot add a mass term.
\end{proof}

\begin{warning}[A quadratic symbol is not yet its vertex bank]
\label{warn:n1v19-vertices}
Theorem~\ref{thm:n1v19-homotopy-gap} populates
\(\dot A\), but the V18 stocks \(q_0,q_1,r_0,r_1\) require first
and second background derivatives of the chosen nonlinear
completion.  Positivity of \(H_*\) does not determine those
vertices.  The covariant endpoint action must be declared before
the four remaining finite coefficient sums can be evaluated.
\end{warning}

\begin{assessment}[Progress after compact V19]
\label{ass:n1v19-status}
V19 constructs an explicit compact, positive, gauge-invariant
irrelevant action whose vertical symbol is
\(\omega^2A\).  It proves that any finite-range covariant two-jet
representative can be stabilized and joined to the compact
parity-tree symbol through a uniformly gapped convex homotopy.
The first of the five V18 Laurent stocks now has the explicit bound
(19.18).  The remaining acquisition is not positivity or
analyticity: it is the nonlinear covariant endpoint action and its
four background-vertex coefficient sums.
\end{assessment}

\begin{decision}[Mandatory V20 audit and next acquisition]
\label{dec:n1v19-next}
Inspect the current SM and Navier--Stokes companion notes before
continuing.  Look specifically for a finite-word construction of a
gauge-covariant higher-derivative endpoint or a reusable bound on
its first two background vertices.  If none is present, use the
already defined covariant plaquette-difference operators from the
legacy Yang--Mills dimension-six bank to instantiate \(H\), \(Q\),
and \(R\), and compute the four coefficient sums in (18.14).
\end{decision}

\begin{firewall}[Claim boundary after compact V19]
\label{fire:n1v19-claim}
V19 proves the positive stabilized homotopy and its uniform real and
complex vertical gaps.  It does not populate the four
background-vertex stocks, evaluate the edge mean, complete SCIC or
ADMC, prove sector tightness, construct the continuum measure or
terminal clustering, or establish the Yang--Mills mass gap.
\end{firewall}

\subsection{Parent record}

This V19 note was formed by appending the preceding material to the
validated compact V18, whose record was
\[
\begin{array}{c|c}
\text{lines}&8410\\
\text{bytes}&286768\\
\text{SHA--256}&
\texttt{AD01B9B698D24C0147DD11D31A43B3DFA28C6EB76F6973BB3198F0A5184ECB17}.
\end{array}
\]

% =============================================================================
% END APPENDED COMPACT-SERIES V19 MATERIAL
% =============================================================================

% =============================================================================
% BEGIN APPENDED COMPACT-SERIES V20 MATERIAL
% =============================================================================

\section{Compact-series V20: compulsory companion audit and a
closed stationary-vertex majorant for the comparison endpoint}
\label{sec:n1v20}

\subsection{The V20 companion audit}

The mandatory five-version audit used the newest complete
companions available in the common notes directory:
\[
\begin{array}{c|r|r|l}
\text{programme and file}&\text{lines}&\text{bytes}&
\text{SHA--256}\\ \hline
\text{SM--Einstein V100\_f10}
&34908&1268894&
\texttt{B6F956D0CEC6B68E0CB6198AC27BA42E66650C4F5663A8F8E3BEA287BAFD3720}\\
\text{Navier--Stokes V26}
&5319&278283&
\texttt{82C887F8C2C69E4F28FAD7C3DC8DE5B9099CB5A4BF431EBA1FFF3E91935978AD}.
\end{array}                                                    \tag{20.1}
\]
The SM stream has completed and archived its tenth compact cycle;
there is no newer unarchived SM note.  Its V100 audit independently
imports the Yang--Mills V14--V17 conclusions that an exact measured
intertwiner produces a scale coboundary and that Ward recombination
must precede a low/edge split.  It also supplies an important
restriction: a sublinear coboundary is sufficient for an averaged
beta coefficient but not for convergence of the cylinder
coefficients defining a projective continuum measure.  The
Navier--Stokes V26 note contains no determinant, gauge endpoint, or
background-vertex construction that transfers to the present
acquisition.

\begin{proposition}[The averaged and projective ledgers are
logically different]
\label{prop:n1v20-two-ledgers}
Let \(e_j\) be a scalar local comparison error.
\begin{enumerate}
\item If \(e_j=B_{j+1}-B_j\) and \(B_N=o(N)\), then
      \[
       \frac1N\sum_{j=0}^{N-1}e_j\longrightarrow0.            \tag{20.2}
      \]
\item The conclusion (20.2), even together with \(e_j\to0\), does
      not imply convergence of
      \(c_N=c_0+\sum_{j<N}e_j\).
\end{enumerate}
\end{proposition}

\begin{proof}
The first assertion follows from
\[
 \frac1N\sum_{j<N}e_j=\frac{B_N-B_0}{N}.
\]
For the second, take \(e_j=(j+1)^{-1}\).  Then \(e_j\to0\) and its
Cesaro mean tends to zero, whereas
\(\sum_{j<N}e_j\) diverges.  Thus the resulting \(c_N\) has no
finite limit.
\end{proof}

\begin{firewall}[Two independent quantitative requirements]
\label{fire:n1v20-two-ledgers}
SCIC and ADMC may use sublinear endpoint cochains and
Cesaro-small errors to identify the universal averaged inverse-
coupling coefficient.  Construction of the continuum functional
still requires summable errors on every fixed cylinder, together
with the weighted and sector-uniform variants needed for
tightness.  No later use of V14--V15 is allowed to replace that
projective summability card.
\end{firewall}

\subsection{A concrete positive covariant endpoint}

V19 left the nonlinear endpoint unspecified.  We now choose the
first covariant plaquette-difference completion from legacy
Yang--Mills V52--V54.  In brief, if
\[
 ({\cal D}_\rho P_{\mu\nu})(x)
 =
 U_\rho(x)P_{\mu\nu}(x+\widehat e_\rho)U_\rho(x)^{-1}
 -P_{\mu\nu}(x),                                             \tag{20.3}
\]
then, with the previously fixed positive normalization \(\nu_1\),
\[
 {\cal O}_1(U)
 =
 \frac{\nu_1}{2}
 \sum_{x,\mu<\nu}\sum_{\rho=1}^4
 \|{\cal D}_\rho P_{\mu\nu}(x)\|_{\rm HS}^2.                 \tag{20.4}
\]
The endpoint action and its flat reduced vertical Hessian are
\[
 {\cal V}_*={\cal V}_{\rm W}+{\cal O}_1,\qquad
 A_*(k)=A(k)+A_1(k),                                        \tag{20.5}
\]
\[
 A_1(k)=D_R(k)^*\Omega_1(k)D_R(k),\qquad
 \Omega_1(k)=\sum_{\rho=1}^4E_\rho(k)^*E_\rho(k).             \tag{20.6}
\]
Here \(E_\rho\) is the exact forward fine-parity difference.
Equations (20.3)--(20.6) import the already proved nonlinear
completion and factorization; their word differentiation is not
repeated here.

\begin{proposition}[Gap, Ward identity, and a literal Hessian
stock]
\label{prop:n1v20-endpoint}
For every real \(k\),
\[
 A_1(k)\succeq0,\qquad
 A_*(k)\succeq\underline\lambda I,\qquad
 \underline\lambda=\frac{49}{1250}.                          \tag{20.7}
\]
Moreover the complete link Hessian of \({\cal O}_1\) annihilates
the fine gauge-gradient incidence exactly.  Its Laurent
coefficient sum satisfies
\[
\boxed{
 {\cal M}(A_1)
 \le384^2\cdot16
 =2359296.}                                                  \tag{20.8}
\]
\end{proposition}

\begin{proof}
The first assertion follows from the factorization (20.6).
The compact parity-tree floor gives
\(A\succeq\underline\lambda I\), proving (20.7).
On the complete link space the plaquette incidence \(D\) obeys
\(Dd=0\) for the gauge-gradient incidence \(d\); hence
\(D^*\Omega_1Dd=0\).

For (20.8), the attached incidence stock is
\({\cal M}(D_R)\le384\).  Each \(E_\rho\) has coefficient sum at
most two, so
\[
 {\cal M}(\Omega_1)
 \le\sum_{\rho=1}^4{\cal M}(E_\rho)^2\le16.
\]
Submultiplicativity of the Laurent coefficient norm in (20.6)
gives (20.8).
\end{proof}

\begin{remark}[Why this endpoint is covariant and irrelevant]
\label{rem:n1v20-covariant-endpoint}
The long-wave quadratic germ of (20.4) is
\(\sum_{\rho,\mu<\nu}\|\nabla_\rho
F_{\mu\nu}^{\rm lin}\|^2\).  It is fourth order in physical
momentum and therefore leaves the Yang--Mills kinetic two-jet
unchanged.  Unlike an arbitrarily signed quartic symbol, (20.4)
is an exact nonnegative compact-link action.  The stronger V19
stabilizer remains available for other signed representatives but
is not needed for this endpoint.
\end{remark}

\subsection{A common complex strip for the direct endpoint path}

Use the direct flat path
\[
                         A_t=A+tA_1,\qquad0\le t\le1.         \tag{20.9}
\]
Put
\[
 M_{20}:=49152+2359296=2408448                              \tag{20.10}
\]
and let \(R_{20}\ge1\) be the maximum Laurent range of \(A\) and
\(A_1\).  Define
\[
 \rho_{20}:=
 \frac1{R_{20}}
 \log\left(
 1+\frac{\underline\lambda}{2M_{20}}
 \right).                                                    \tag{20.11}
\]

\begin{lemma}[Uniform inverse on the V20 strip]
\label{lem:n1v20-strip}
If \(k\in{\mathbb T}^4\), \(0\le t\le1\), and
\(\|\operatorname{Im}z\|_\infty\le\rho_{20}\), then
\[
\boxed{
 \|A_t(k+z)^{-1}\|\le
 \frac2{\underline\lambda}
 =\frac{2500}{49}=:g_0.}                                    \tag{20.12}
\]
\end{lemma}

\begin{proof}
The real matrix \(A_t(k)\) has floor
\(\underline\lambda\) by (20.7).  The finite Laurent estimate and
(20.10)--(20.11) give
\[
 \|A_t(k+z)-A_t(k)\|
 \le M_{20}(e^{R_{20}\rho_{20}}-1)
 =\frac{\underline\lambda}{2}.
\]
The inverse-relative perturbation has norm at most \(1/2\).
The Neumann series therefore doubles the real inverse bound
\(\underline\lambda^{-1}\), proving (20.12).
\end{proof}

\begin{corollary}[The first response stock is numerical]
\label{cor:n1v20-a1}
On the response strip
\[
 \|\operatorname{Im}k\|_\infty,\,
 |\operatorname{Im}q|
 \le\frac{\rho_{20}}4,                                      \tag{20.13}
\]
the V18 stock \(a_1\) may be taken to be
\[
\boxed{
 a_1=
 2359296\,
 \exp\left(\frac{R(A_1)\rho_{20}}4\right).}                  \tag{20.14}
\]
\end{corollary}

\begin{proof}
Here \(\dot A=A_1\).  Apply the finite Laurent estimate to
(20.8).  The shifted inverse at \(k+q\widehat e_1\) stays inside
half of the strip (20.11), so (20.12) applies.
\end{proof}

\subsection{Raw word stocks and the stationary-centre map}

We next turn the inherited finite word bank into total background
vertices.  Work in the fixed parity-cell chord coordinates before
stationary recentering.  Write \(x\) for the forty-five vertical
coordinates and \(v\) for a retained background polarization.
For endpoint \(e\in\{{\rm W},*\}\), set
\[
 J_{ab}^{e}:=
 \sup_{(20.13)}
 \left\|
 D_x^aD_v^b{\cal V}_e(0,0)
 \right\|,
 \qquad
 1\le a+b\le4.                                               \tag{20.15}
\]
The norm includes the finite parity offsets and the fixed colour
trace.  Thus (20.15) is a literal finite coefficient-sum norm,
not a thermodynamic-volume norm.

The legacy factorized word compiler gives, for one normalized
owner of (20.4), the order two, three, and four bounds
\[
 C_2=100,\qquad C_3=1560,\qquad C_4=19312.                   \tag{20.16}
\]
There are \(16\cdot6\cdot4=384\) owners per parity cell.  If
\(d_x,d_v\) are the exact coefficient norms of the fixed vertical
and retained link injections, then
\[
\boxed{
 J_{ab}^{*}
 \le
 J_{ab}^{\rm W}
 +384\,\nu_1 C_{a+b}d_x^ad_v^b
 \exp\left(\frac{R_{ab}\rho_{20}}2\right),
 \quad 2\le a+b\le4.}                                       \tag{20.17}
\]
All quantities on the right are already fixed finite incidence or
word data.  Numerically, before the injection and strip factors,
\[
 384C_2=38400,\quad
 384C_3=599040,\quad
 384C_4=7415808.                                             \tag{20.18}
\]

\begin{lemma}[No hidden fifth action derivative]
\label{lem:n1v20-four-jets}
The total first and second background derivatives of the stationary
vertical Hessian at \(v=0\) use only the raw stocks (20.15) of
total order at most four.
\end{lemma}

\begin{proof}
Let \(x_e(v)\) solve
\[
 D_x{\cal V}_e(x_e(v),v)=0,\qquad x_e(0)=0,                  \tag{20.19}
\]
and abbreviate
\[
 X_e=D_vx_e(0),\qquad X_{2,e}=D_v^2x_e(0).
\]
Twice differentiating (20.19) gives the exact identities
\[
 X_e=-A_e^{-1}D_xD_v{\cal V}_e,                              \tag{20.20}
\]
\[
 X_{2,e}
 =-A_e^{-1}\left(
 D_xD_v^2{\cal V}_e
 +2D_x^2D_v{\cal V}_e[X_e]
 +D_x^3{\cal V}_e[X_e,X_e]
 \right).                                                    \tag{20.21}
\]
For
\(\widehat A_e(v)=D_x^2{\cal V}_e(x_e(v),v)\), the chain rule gives
\[
 R_e:=D_v\widehat A_e(0)
 =
 D_x^2D_v{\cal V}_e
 +D_x^3{\cal V}_e[X_e],                                     \tag{20.22}
\]
\[
\begin{split}
 Q_e:=D_v^2\widehat A_e(0)
 ={}&
 D_x^2D_v^2{\cal V}_e
 +2D_x^3D_v{\cal V}_e[X_e]\\
 &+D_x^4{\cal V}_e[X_e,X_e]
 +D_x^3{\cal V}_e[X_{2,e}].                                 \tag{20.23}
\end{split}
\]
The largest derivative in (20.20)--(20.23) has total order four.
These are precisely the stationary score and Schur terms that
would be lost by freezing the centre.
\end{proof}

\begin{theorem}[Explicit stationary-vertex majorant]
\label{thm:n1v20-stationary-majorant}
For \(e\in\{{\rm W},*\}\), define
\[
 \ell_e:=g_0J_{11}^e,                                       \tag{20.24}
\]
\[
 \kappa_e:=
 g_0\left(
 J_{12}^e+2J_{21}^e\ell_e+J_{30}^e\ell_e^2
 \right),                                                    \tag{20.25}
\]
\[
 {\mathfrak r}_e:=
 J_{21}^e+J_{30}^e\ell_e,                                   \tag{20.26}
\]
\[
 {\mathfrak q}_e:=
 J_{22}^e+2J_{31}^e\ell_e
 +J_{40}^e\ell_e^2+J_{30}^e\kappa_e.                        \tag{20.27}
\]
Then on (20.13),
\[
 \|X_e\|\le\ell_e,\qquad
 \|X_{2,e}\|\le\kappa_e,\qquad
 \|R_e\|\le{\mathfrak r}_e,\qquad
 \|Q_e\|\le{\mathfrak q}_e.                                 \tag{20.28}
\]
\end{theorem}

\begin{proof}
Apply (20.12) and submultiplicativity successively to
(20.20), (20.21), (20.22), and (20.23).  The four resulting
inequalities are exactly (20.24)--(20.28).
\end{proof}

\subsection{Haar absorption and the four remaining V18 stocks}

Let \(h_e(v)=-\log J(x_e(v))\) be the reduced Haar density.  In the
fixed coordinates,
\[
 Dh(0)=0,\qquad \|D^2h(0)\|=\frac12.                         \tag{20.29}
\]
Consequently
\[
 |D_v^2h_e(0)|\le\frac12\ell_e^2.                            \tag{20.30}
\]
If \(M_e=\sup_{(20.13)}\|\widehat A_e\|\), define the canonically
Haar-absorbed quadratic vertex
\[
 Q_e^{\rm meas}
 :=
 Q_e+\frac{2D_v^2h_e(0)}{45}\widehat A_e.                    \tag{20.31}
\]
Then
\[
 \frac12\Tr(\widehat A_e^{-1}Q_e^{\rm meas})
 =
 \frac12\Tr(\widehat A_e^{-1}Q_e)+D_v^2h_e(0)                \tag{20.32}
\]
and
\[
 \|Q_e^{\rm meas}\|
 \le
 \widehat{\mathfrak q}_e
 :=
 {\mathfrak q}_e+\frac{M_e\ell_e^2}{45}.                     \tag{20.33}
\]

The one-loop comparison need not interpolate nonlinear actions.
It may interpolate the two exact positive stationary Hessian
families and their density scalars:
\[
 \widehat A_t(v)
 =(1-t)\widehat A_{\rm W}(v)+t\widehat A_*(v),\qquad
 h_t(v)=(1-t)h_{\rm W}(v)+th_*(v).                           \tag{20.34}
\]
This makes
\[
 R_t=(1-t)R_{\rm W}+tR_*,\qquad
 Q_t^{\rm meas}
 =(1-t)Q_{\rm W}^{\rm meas}+tQ_*^{\rm meas}.                 \tag{20.35}
\]

\begin{theorem}[Completion of the five-stock analyticity card]
\label{thm:n1v20-five-stocks}
For the Gaussian homotopy (20.34), the five stocks required in
V18 may be chosen as
\[
\boxed{\begin{aligned}
 a_1&=
 2359296
 \exp\left(\frac{R(A_1)\rho_{20}}4\right),\\
 r_0&=\max\{
 {\mathfrak r}_{\rm W},{\mathfrak r}_*\},&
 r_1&={\mathfrak r}_{\rm W}+{\mathfrak r}_*,\\
 q_0&=\max\{
 \widehat{\mathfrak q}_{\rm W},
 \widehat{\mathfrak q}_*\},&
 q_1&=
 \widehat{\mathfrak q}_{\rm W}
 +\widehat{\mathfrak q}_* .
\end{aligned}}                                               \tag{20.36}
\]
Every entry on the right is an explicit finite expression in the
attached Wilson incidence bank, the 384 six-word owners, the fixed
normalization \(\nu_1\), and the certified inverse bound
\(g_0=2500/49\).
\end{theorem}

\begin{proof}
The \(a_1\) row is Corollary~\ref{cor:n1v20-a1}.
Equation (20.35) and convexity of the norm give the two undotted
bounds.  Differentiation in \(t\) gives
\[
 \dot R_t=R_*-R_{\rm W},\qquad
 \dot Q_t^{\rm meas}
 =Q_*^{\rm meas}-Q_{\rm W}^{\rm meas}.
\]
The triangle inequality and (20.28), (20.33) give the dotted
bounds.
\end{proof}

\begin{proposition}[The Gaussian path is endpoint exact]
\label{prop:n1v20-path-exact}
For every retained background in the common positive tube,
\[
\begin{split}
 &\int_0^1\partial_t\left[
 \frac12\log\det\widehat A_t(v)+h_t(v)
 \right]\,dt\\
 &\qquad=
 \frac12\log\det\widehat A_*(v)+h_*(v)
 -\frac12\log\det\widehat A_{\rm W}(v)-h_{\rm W}(v).          \tag{20.37}
\end{split}
\]
Thus use of (20.34) changes no endpoint determinant response.
\end{proposition}

\begin{proof}
Positivity makes the logarithmic determinant a smooth real
function throughout the path.  The fundamental theorem of
calculus gives (20.37).  Background differentiation commutes with
the \(t\)-integral on the common compact tube by the uniform gap
and the bounds (20.36).
\end{proof}

\begin{proposition}[Ward compatibility of the direct Gaussian
path]
\label{prop:n1v20-Ward}
The complete Ward identity, including the Haar and stationary
terms, holds for every \(t\) in (20.34).  The endpoint deformation
adds no gauge-boson mass term.
\end{proposition}

\begin{proof}
Both endpoint Hessian families are obtained from exact
gauge-invariant compact-link actions after measured stationary
reduction.  They therefore transform covariantly under retained
background gauge transformations.  Convex combinations preserve
that covariance, and differentiating it gives the Ward identities
for (20.35).  The Haar scalar is interpolated together with its
endpoint density and hence is not omitted.  Finally the physical
long-wave Hessian of \({\cal O}_1\) is fourth order in momentum, so
its zero-momentum mass row vanishes.
\end{proof}

\begin{assessment}[Progress after compact V20]
\label{ass:n1v20-status}
The required companion audit is complete.  It adds the
projective-summability firewall (20.2) and prevents the averaged
anomaly criterion from being misused later.  The nonlinear
comparison endpoint is now the explicit positive action
\({\cal V}_{\rm W}+{\cal O}_1\).  Its flat Hessian has an exact
factorization, Ward zero, spectral floor, numerical Laurent stock,
and common complex strip.  The inherited order-four word bank,
together with the exact stationary recursion
(20.20)--(20.23), populates all five V18 response stocks by
(20.36).  The remaining edge acquisition is finite evaluation:
assemble the 384 owner arrays, verify their quadratic
factorization and Ward residual exactly, and run the outward-
rounded grid enclosure of the single alias mean.
\end{assessment}

\begin{decision}[V21 finite execution target]
\label{dec:n1v20-next}
Instantiate the factorized mixed-jet certificate for
\({\cal O}_1\) on one parity cell.  The execution must produce:
\begin{enumerate}
\item the exact quadratic residual against (20.6);
\item exact or outward-rounded arrays for
      \(R_*-R_{\rm W}\) and
      \(Q_*^{\rm meas}-Q_{\rm W}^{\rm meas}\);
\item the complete toron-inclusive Ward residual before any
      low/edge split; and
\item an interval for the constant alias class whose truncation
      error is bounded with (20.36).
\end{enumerate}
If the attachment lacks machine-readable owner arrays, the note
must expose that missing finite input rather than replacing it by
a universality assertion.
\end{decision}

\begin{firewall}[Claim boundary after compact V20]
\label{fire:n1v20-claim}
V20 closes the construction and analytic majorization of the
finite comparison bank.  It does not execute the mixed-jet array,
evaluate or prove vanishing of the edge mean, complete SCIC or
ADMC, prove summable cylinder comparison, establish sector
tightness, construct the continuum measure or terminal
clustering, or prove the Yang--Mills mass gap.
\end{firewall}

\subsection{Parent record}

This V20 note was formed by appending the preceding material to the
validated compact V19, whose record was
\[
\begin{array}{c|c}
\text{lines}&8805\\
\text{bytes}&299706\\
\text{SHA--256}&
\texttt{15D26459E9EEB32043639C8F16678B943011262E2AC397DBAED43C8412F09507}.
\end{array}
\]

% =============================================================================
% END APPENDED COMPACT-SERIES V20 MATERIAL
% =============================================================================

% =============================================================================
% BEGIN APPENDED COMPACT-SERIES V21 MATERIAL
% =============================================================================

\section{Compact-series V21: exact execution of the parity-cell
quadratic owner certificate}
\label{sec:n1v21}

\subsection{Purpose and certificate provenance}

V20 reduced the comparison to a finite owner bank but retained a
very crude coefficient bound.  This note executes the entire
quadratic part of that bank over the integer Laurent ring
\[
 {\cal L}:=
 {\mathbb Z}[z_1,z_1^{-1},\ldots,z_4,z_4^{-1}],\qquad
 z_\rho=e^{ik_\rho}.                                        \tag{21.1}
\]
The reproducible verifier is
\[
 \texttt{scripts/verify\_ym\_v21\_incidence.py}.              \tag{21.2}
\]
Its validated record is
\[
\begin{array}{c|c}
\text{lines}&311\\
\text{bytes}&10796\\
\text{SHA--256}&
\texttt{DE678E5CAE7671B90C8788AF9A489444D8A63B4F32EA78F08BC68AA0131B0AE0}.
\end{array}                                                   \tag{21.3}
\]
The calculation uses no floating-point arithmetic.  Every matrix
entry is a finite dictionary
\(\alpha\mapsto c_\alpha\in{\mathbb Z}\) representing
\(\sum_\alpha c_\alpha z^\alpha\).

\subsection{Two independent constructions of the owner rows}

Index fine edge types by
\[
 {\cal E}=\{(s,\mu):s\in\{0,1\}^4,\ 1\le\mu\le4\}            \tag{21.4}
\]
and plaquette rows by
\[
 {\cal P}=\{(s,\mu\nu):s\in\{0,1\}^4,\ 1\le\mu<\nu\le4\}.
                                                                    \tag{21.5}
\]
For an edge beginning at \(z=2y+s'\), its Laurent monomial is
\(z^{y}\) and its type is \((s',\mu)\).  Applying this rule to the
four oriented edges of every plaquette constructs
\[
                         D\in{\cal L}^{96\times64}.           \tag{21.6}
\]

The reduced set is generated rather than entered by hand.  With
\[
 m_{s,\mu}=
 \begin{cases}
 |\{\nu<\mu:s_\nu=1\}|,&s_\mu=0,\\
 2|\{\nu>\mu:s_\nu=1\}|
 +|\{\nu<\mu:s_\nu=1\}|,&s_\mu=1,
 \end{cases}                                                 \tag{21.7}
\]
put
\[
 {\cal R}=\{(s,\mu):m_{s,\mu}>0\},\quad
 {\cal C}=\{(\widehat e_\mu,\mu):1\le\mu\le4\},              \tag{21.8}
\]
and let \({\cal T}\) be the complement.

For each \((s,\mu\nu,\rho)\), the first construction forms the
linearized transported-plaquette difference directly:
\[
 K^{\rm dir}_{s,\mu\nu,\rho}
 :=
 D_{s+\widehat e_\rho,\mu\nu}-D_{s,\mu\nu},                  \tag{21.9}
\]
where a base point crossing the parity cell is converted to its
coarse Laurent monomial.  The second construction uses only the
stored rows of \(D\):
\[
 K^{\rm fac}_{s,\mu\nu,\rho}
 :=
 \begin{cases}
 D_{s+\widehat e_\rho,\mu\nu}-D_{s,\mu\nu},
 &s_\rho=0,\\
 z_\rho D_{s-\widehat e_\rho,\mu\nu}-D_{s,\mu\nu},
 &s_\rho=1.
 \end{cases}                                                 \tag{21.10}
\]
Thus \(K^{\rm fac}\) is the row realization of \(E_\rho D\),
whereas \(K^{\rm dir}\) is generated from the two actual
plaquette words.

\begin{theorem}[Exact owner and factorization residuals]
\label{thm:n1v21-residuals}
The certificate gives, as identities in \({\cal L}\),
\[
\boxed{
 K^{\rm dir}=K^{\rm fac},\qquad
 (K^{\rm dir}_R)^*K^{\rm dir}_R
 =
 (K^{\rm fac}_R)^*K^{\rm fac}_R.}                            \tag{21.11}
\]
The common Gram matrix is Hermitian.  The generated cardinalities
and filling checks are
\[
\begin{array}{c|rrrrrr}
 &|{\cal P}|&|{\cal E}|&|{\cal R}|&|{\cal C}|&|{\cal T}|
 &\text{owners}\\ \hline
\text{value}&96&64&45&4&15&384,
\end{array}                                                   \tag{21.12}
\]
\[
 \max m_{s,\mu}=6,\qquad
 \sum_{s,\mu}m_{s,\mu}=96,\qquad
 \sum_{s,\mu}m_{s,\mu}^2=272.                               \tag{21.13}
\]
\end{theorem}

\begin{proof}
The verifier constructs (21.9) by calling the plaquette-incidence
routine at the two integer base points independently.  It
constructs (21.10) from the parity cases after \(D\) is complete.
It compares every coefficient dictionary in all
\(384\cdot64\) entries.  Equality of the reduced Gram matrices is
then checked entry by entry.  Hermiticity is checked independently
by comparing entry \(ij\) with the Laurent adjoint of entry \(ji\).
All comparisons return true.  The counts in (21.12)--(21.13) are
obtained from the generated sets (21.4)--(21.8), and agree with
the legacy hand counts.
\end{proof}

\subsection{Exact Ward residual}

Let \(d\in{\cal L}^{64\times16}\) be the generated fine
site-to-edge gradient.  Its row for \((s,\mu)\) is
\[
 (d\phi)_{s,\mu}=
 \begin{cases}
 \phi_{s+\widehat e_\mu}-\phi_s,&s_\mu=0,\\
 z_\mu\phi_{s-\widehat e_\mu}-\phi_s,&s_\mu=1.
 \end{cases}                                                 \tag{21.14}
\]

\begin{theorem}[Complete-link Ward zero before reduction]
\label{thm:n1v21-Ward}
The exact residuals are
\[
\boxed{
                         Dd=0,\qquad
                         K^{\rm dir}d=0}                     \tag{21.15}
\]
in \({\cal L}\).  Consequently the full completion Hessian obeys
\[
                         (K^{\rm dir})^*K^{\rm dir}d=0.       \tag{21.16}
\]
\end{theorem}

\begin{proof}
The verifier multiplies the independently generated matrices in
(21.6), (21.9), and (21.14), collecting like Laurent monomials
with exact integer coefficients.  Every coefficient dictionary of
both products in (21.15) is empty after cancellation.  Equation
(21.16) follows algebraically.  The test is performed on the
complete sixty-four-column incidence before the tree and
second-constituent columns are removed.
\end{proof}

\begin{firewall}[What the exact Ward test does not test]
\label{fire:n1v21-Ward-scope}
Equation (21.15) is the quadratic curl--gradient identity.  It
does not yet test the background differentiated tadpole, bubble,
Haar, stationary-centre, and toron sum.  Those rows require the
third and fourth noncommutative word jets and must pass a second,
coordinate-complete Ward residual.
\end{firewall}

\subsection{Sharp Laurent inventories}

For a matrix Laurent polynomial \(P=(P_{ij})\), write
\[
 {\cal M}(P):=
 \sum_{i,j}\sum_{\alpha}
 |[z^\alpha]P_{ij}|.                                        \tag{21.17}
\]
The exact generated inventories are
\[
\begin{array}{c|rrrrr}
P&
\#\{P_{ij}\ne0\}&
\#\text{ Laurent terms}&
{\cal M}(P)&
\max|[z^\alpha]P_{ij}|&
\max|\alpha|_1\\ \hline
D_R&270&270&270&1&1\\
D_R^*D_R&561&673&898&6&2\\
A_1&1523&4067&13764&54&3.
\end{array}                                                   \tag{21.18}
\]
Here \(A_1=(K_R^{\rm dir})^*K_R^{\rm dir}\), with the
normalization of (20.6).  Its canonical ordered coefficient list
has SHA--256
\[
 \texttt{21465A8A21CC2E25A1ACF56F47093E5108DD11D0518DC67ED1D7002B523115E8}.
                                                                    \tag{21.19}
\]

\begin{remark}[Normalization reconciliation]
\label{rem:n1v21-normalization}
The factorization convention (20.6) fixes the \(\nu_1\) appearing
in (20.4) to one.  Thus (21.18) is the literal stock for the
chosen endpoint.  If a later convention rescales the owner by a
positive scalar \(\nu\), the \(A_1\) coefficient stock and all its
action jets are multiplied by \(\nu\); no combinatorial array is
changed.
\end{remark}

\begin{corollary}[Sharp replacement of the V20 Hessian bound]
\label{cor:n1v21-sharp-stock}
For the chosen normalization,
\[
\boxed{
 {\cal M}(A)=\frac{898}{3},\qquad
 {\cal M}(A_1)=13764,\qquad
 {\cal M}(A_t)\le\frac{42190}{3}.}                            \tag{21.20}
\]
In particular the V20 bound \(2359296\) for \(A_1\) may be
replaced by \(13764\).
\end{corollary}

\begin{proof}
Use \(A=\frac13D_R^*D_R\), the two exact rows of (21.18), and
subadditivity along \(A_t=A+tA_1\).
\end{proof}

\subsection{A sharper certified complex strip}

Because the maximum \(\ell^1\)-range in (21.18) is three, put
\[
 \rho_{21}:=
 \frac13\log\left(
 1+\frac{\underline\lambda}
 {2(42190/3)}
 \right)
 =
 \frac13\log\left(
 1+\frac{147}{105475000}
 \right).                                                    \tag{21.21}
\]

\begin{theorem}[Refined inverse and response stock]
\label{thm:n1v21-refined-strip}
For every real \(k\), \(0\le t\le1\), and
\(\|\operatorname{Im}z\|_\infty\le\rho_{21}\),
\[
 \|A_t(k+z)^{-1}\|\le\frac{2500}{49}.                        \tag{21.22}
\]
On
\[
 \|\operatorname{Im}k\|_\infty,\,
 |\operatorname{Im}q|\le\rho_{21}/4,
\]
the first V18 stock may be sharpened to
\[
\boxed{
 a_1=
 13764\exp(3\rho_{21}/4).}                                   \tag{21.23}
\]
\end{theorem}

\begin{proof}
The real floor remains \(\underline\lambda=49/1250\).
Equations (21.20)--(21.21) make the complex perturbation at most
\(\underline\lambda/2\).  The same Neumann argument as in
Lemma~\ref{lem:n1v20-strip} proves (21.22).  Apply the finite
Laurent bound to the exact range and stock of \(A_1\) in
(21.18) to obtain (21.23).
\end{proof}

\subsection{The precise nonlinear remainder}

The quadratic part of the V21 decision has now been executed.
The attachment contains the rule generating every six-link owner,
but it contains no machine-readable array for the labelled
noncommutative jets
\[
 Y_{ab},\qquad Y_{abc},                                     \tag{21.24}
\]
after insertion of the eight colour directions, the harmonic
background lift, its deformation, and the stationary recursion.
The absence of a stored array is not an existence obstruction:
the legacy product rule uniquely generates (21.24).  It is an
execution obstruction because a numerical edge interval cannot be
audited without preserving the exact basis order, owner offsets,
colour contractions, Haar row, and toron convention.

\begin{proposition}[Quadratic success cannot certify the
one-loop response]
\label{prop:n1v21-nonlinear-independence}
The identities (21.11), (21.15), and the stock (21.18) do not
determine \(R_*-R_{\rm W}\) or
\(Q_*^{\rm meas}-Q_{\rm W}^{\rm meas}\).
\end{proposition}

\begin{proof}
Add to the endpoint any gauge-invariant local action beginning at
fourth order in the link logarithms.  Its value, gradient, and
flat Hessian vanish, so it leaves every V21 quadratic identity and
inventory unchanged.  Its fourth derivative can be nonzero and
enters the two-background tadpole vertex.  Hence the nonlinear
response is independent data.
\end{proof}

\begin{assessment}[Progress after compact V21]
\label{ass:n1v21-status}
The first finite execution gate is complete and reproducible.
The direct 384-owner construction agrees exactly with the
factorized \(E_\rho D_R\) construction; the complete Ward
residuals vanish over the integer Laurent ring; and the crude
endpoint coefficient stock has fallen from \(2359296\) to the
exact value \(13764\).  This yields the sharper strip (21.21) and
stock (21.23).  The sole unfinished part of the finite comparison
bank is now sharply typed as the noncommutative third/fourth
owner-jet array with its coordinate-complete contraction, not the
incidence, gap, or quadratic Ward rows.
\end{assessment}

\begin{decision}[V22 nonlinear execution target]
\label{dec:n1v21-next}
Extend the verifier with an exact degree-four jet algebra for
ordered \(SU(3)\) link words.  First reproduce the quadratic Gram
matrix and hash (21.19) from the jet engine as an independent
normalization test.  Then emit the third and fourth mixed owners
in a fixed adjoint basis, apply the stationary formulas
(20.20)--(20.23), and test the complete background Ward sum.
Only after that residual is zero may the toron-inclusive interval
for the constant alias class be evaluated.
\end{decision}

\begin{firewall}[Claim boundary after compact V21]
\label{fire:n1v21-claim}
V21 executes and certifies the quadratic endpoint bank.  It does
not execute the noncommutative mixed jets, evaluate the edge mean,
complete SCIC or ADMC, prove summable cylinder comparison,
establish sector tightness, construct the continuum measure or
terminal clustering, or prove the Yang--Mills mass gap.
\end{firewall}

\subsection{Parent record}

This V21 note was formed by appending the preceding material to the
validated compact V20, whose record was
\[
\begin{array}{c|c}
\text{lines}&9357\\
\text{bytes}&317451\\
\text{SHA--256}&
\texttt{842D84EF6BE00B9C2DE525EF5CB389D0FF3BBAA46A8854F4D0CBFD596A9A021A}.
\end{array}
\]

% =============================================================================
% END APPENDED COMPACT-SERIES V21 MATERIAL
% =============================================================================

% =============================================================================
% BEGIN APPENDED COMPACT-SERIES V22 MATERIAL
% =============================================================================

\section{Compact-series V22: an exact noncommutative
\(SU(3)\) word-jet engine}
\label{sec:n1v22}

\subsection{Exact colour field and basis}

The nonlinear V21 remainder cannot be computed over scalar
incidence coefficients: ordered colour products first appear in
the cubic owner jet.  We therefore use the exact field
\({\mathbb Q}(i)\), not floating-point Gell--Mann matrices.
Let \(E_{ab}\) denote the elementary \(3\times3\) matrix and set
\[
\begin{array}{lll}
 T_1=E_{12}-E_{21},&
 T_2=i(E_{12}+E_{21}),&
 T_3=i\operatorname{diag}(1,-1,0),\\
 T_4=E_{13}-E_{31},&
 T_5=i(E_{13}+E_{31}),&
 T_6=E_{23}-E_{32},\\
 T_7=i(E_{23}+E_{32}),&
 T_8=i\operatorname{diag}(0,1,-1).
\end{array}                                                   \tag{22.1}
\]
Every \(T_a\) is anti-Hermitian and traceless.  In the real
Hilbert--Schmidt metric
\[
                         g_{ab}:=
 \Re\Tr(T_a^*T_b),                                          \tag{22.2}
\]
the exact Gram matrix is
\[
 g=
 \begin{pmatrix}
 2&0&0&0&0&0&0&0\\
 0&2&0&0&0&0&0&0\\
 0&0&2&0&0&0&0&-1\\
 0&0&0&2&0&0&0&0\\
 0&0&0&0&2&0&0&0\\
 0&0&0&0&0&2&0&0\\
 0&0&0&0&0&0&2&0\\
 0&0&-1&0&0&0&0&2
 \end{pmatrix}.                                              \tag{22.3}
\]
In particular \(g>0\); the nonorthogonal last Cartan direction
introduces only the rational inverse of
\(\left(\begin{smallmatrix}2&-1\\-1&2\end{smallmatrix}\right)\).

\begin{lemma}[Exact exponential jets]
\label{lem:n1v22-exp-jets}
For matrices \(H_1,\ldots,H_n\) and \(n\ge1\),
\[
 D^n\exp(0)[H_1,\ldots,H_n]
 =
 \frac1{n!}\sum_{\pi\in{\mathfrak S}_n}
 H_{\pi(1)}\cdots H_{\pi(n)}.                               \tag{22.4}
\]
For the inverse link \(\exp(-X)\), the right side is multiplied by
\((-1)^n\).
\end{lemma}

\begin{proof}
Extract the coefficient of \(t_1\cdots t_n\) from
\[
 \exp\left(\sum_{j=1}^nt_jH_j\right)
 =
 \sum_{m\ge0}\frac1{m!}
 \left(\sum_{j=1}^nt_jH_j\right)^m.
\]
Only \(m=n\) contributes, and its multilinear words are indexed by
the \(n!\) permutations.  Replacing every \(H_j\) by \(-H_j\)
proves the inverse formula.
\end{proof}

\subsection{The labelled finite-word algorithm}

For an oriented link factor \(g_r\), retain its fine-edge type,
coarse Laurent offset, and orientation sign.  If
\(W=g_1\cdots g_m\), assign every labelled derivative to one of
the \(m\) factors.  For a map
\(\sigma:\{1,\ldots,n\}\to\{1,\ldots,m\}\), put
\[
 I_r(\sigma):=\sigma^{-1}(r).
\]
The exact word jet used by the verifier is
\[
\boxed{
 W_{1\cdots n}
 =
 \sum_{\sigma}
 (g_1)_{I_1(\sigma)}
 \cdots
 (g_m)_{I_m(\sigma)},}                                      \tag{22.5}
\]
where each factor derivative is (22.4) with its Laurent
monomial.  Empty derivative sets contribute the identity.

For one covariant-difference owner, write
\[
 Y=CP_+C^{-1}-P_0,\qquad
 o(Y)=\frac12\Re\Tr(Y^*Y).                                  \tag{22.6}
\]
At the flat connection \(Y=0\).  Hence for
\(2\le n\le4\) and \(I=\{1,\ldots,n\}\),
\[
\boxed{
 D^no[h_I]
 =
 \frac12
 \sum_{\varnothing\ne J\subsetneq I}
 \left\langle
 Y_J,Y_{I\setminus J}
 \right\rangle_{\mathbb R}.}                                \tag{22.7}
\]
This single subset formula contains the three \(2+1\) terms at
order three and the three \(2+2\) plus four \(3+1\) terms at
order four.

\begin{proposition}[Exactness and permutation symmetry of the
jet algorithm]
\label{prop:n1v22-jet-exact}
Equations (22.4)--(22.7) compute the Fréchet derivatives of the
literal compact-link owner through order four.  The resulting
multilinear forms are symmetric in their labelled directions.
\end{proposition}

\begin{proof}
Repeated product differentiation gives (22.5), with each
assignment occurring exactly once.  Differentiating the real
quadratic pairing in (22.6) gives (22.7).  Both formulas are
coefficients of the ordinary multivariate Taylor expansion of a
smooth scalar function, so mixed derivatives commute.  Equivalently,
permuting the labels bijects the assignment and subset families.
\end{proof}

\subsection{Executable normalization replay}

The implementation is
\[
 \texttt{scripts/verify\_ym\_v22\_word\_jets.py}.             \tag{22.8}
\]
Its record is
\[
\begin{array}{c|c}
\text{lines}&544\\
\text{bytes}&17742\\
\text{SHA--256}&
\texttt{9FCEF587697796C5DC3AFE28B98AAB4B9705728F9B4003D9D391B060F4C1CA21}.
\end{array}                                                   \tag{22.9}
\]
The program represents every colour entry by two rational numbers,
every Laurent coefficient by such an exact Gaussian rational, and
every derivative coefficient in (22.4) by a rational number.

\begin{theorem}[All first word jets reproduce the V21 bank]
\label{thm:n1v22-replay}
For every one of the \(384\) owners and every one of the
forty-five reduced edge types, the first derivative obtained from
the literal six-factor word (22.5) equals the corresponding entry
of \(K^{\rm dir}\) times \(T_1\).  The Gram matrix reconstructed
from those word jets has canonical SHA--256
\[
 \boxed{
 \texttt{21465A8A21CC2E25A1ACF56F47093E5108DD11D0518DC67ED1D7002B523115E8},}
                                                                    \tag{22.10}
\]
identical to (21.19).
\end{theorem}

\begin{proof}
For each of the \(384\cdot45\) inputs, the verifier generates the
transported six-factor word and the four-factor local plaquette
independently of the V21 row.  It applies (22.4)--(22.5), compares
every colour Laurent coefficient with
\((K^{\rm dir})_{oe}T_1\), and obtains equality in every entry.
Since \(g_{11}=2\), division of the colour Gram by two gives the
scalar spatial Gram.  Its independently serialized coefficient
list has the hash (22.10).
\end{proof}

\begin{firewall}[Why the hash replay matters]
\label{fire:n1v22-replay}
A nonlinear word program with a reversed link order, wrong inverse
sign, shifted parity phase, or inconsistent colour normalization
can still produce finite symmetric tensors.  Equality of all
first jets and the full Gram hash anchors the nonlinear engine to
the exact quadratic endpoint that carries the certified gap.
\end{firewall}

\subsection{Nonzero exact cubic and quartic tests}

Use the owner based at parity \(0000\), plaquette directions
\((1,2)\), and transport direction \(3\).  Order its distinct
fine-edge types lexicographically.  The cubic test inserts
\[
\begin{array}{c|c}
\text{edge type}&\text{colour}\\ \hline
(0000,1)&T_1\\
(0000,2)&T_2\\
(0010,1)&T_3,
\end{array}                                                   \tag{22.11}
\]
where \(0010\) denotes the parity with its third component equal
to one.  The quartic test inserts \(T_1\) four times on
\((0000,1)\).

\begin{theorem}[Exact nonlinear sample values]
\label{thm:n1v22-samples}
For the unnormalized owner (22.6), the engine obtains
\[
\boxed{
 D^3o[(0000,1)T_1,(0000,2)T_2,(0010,1)T_3]=2,}              \tag{22.12}
\]
\[
\boxed{
 D^4o[(0000,1)T_1]^4=-2.}                                   \tag{22.13}
\]
The cubic value is unchanged under all six permutations, and the
quartic value under all twenty-four labelled permutations.  Their
canonical hashes are respectively
\[
 \texttt{661B1386C4E3D6654F77CC2F9C604BFC51B11B90535062A3C8963E806A0B19A9},
                                                                    \tag{22.14}
\]
\[
 \texttt{87270C18CBCD3079ED04326D82BA8C0A30E470DD1DBFA00319EA5A1859580397}.
                                                                    \tag{22.15}
\]
\end{theorem}

\begin{proof}
The verifier evaluates (22.5)--(22.7) over
\({\mathbb Q}(i)[z_1^{\pm1},\ldots,z_4^{\pm1}]\).  Each selected
answer has only the zero Laurent monomial, with rational
coefficients \(2\) and \(-2\).  It recomputes the scalar derivative
for every permutation of the labelled input tuple and compares the
full coefficient dictionaries.  All comparisons return true.
\end{proof}

\begin{corollary}[The nonlinear endpoint data are now
algorithmically populated]
\label{cor:n1v22-populated}
For any specified vertical and retained background directions,
every direct endpoint contribution to \(R_*-R_{\rm W}\) and
\(Q_*-Q_{\rm W}\) is a terminating exact calculation in the
algebra used by the V22 verifier.
\end{corollary}

\begin{proof}
The linear vertex requires two vertical and one background
direction, hence (22.7) with \(n=3\).  The quadratic vertex
requires two of each, hence \(n=4\).  There are finitely many
owners, parity offsets, edge types, and eight colour directions.
All arithmetic operations are exact and all assignment sums are
finite.
\end{proof}

\subsection{What remains before the full response array}

The engine deliberately stops before contracting the full owner
bank.  Three coordinate layers must be fixed in the same run:
\begin{enumerate}
\item the four retained link directions and their harmonic lift
      \(L=-A^{-1}B\) at each internal momentum;
\item the stationary-centre derivatives (20.20)--(20.23) and
      the Haar absorption (20.31); and
\item the dual-colour contraction using \(g^{-1}\), followed by
      the complete toron-inclusive tadpole--bubble Ward sum.
\end{enumerate}
The first layer contains the certified analytic matrix inverse and
therefore leaves the integer Laurent ring, but it remains a
rigorous complex interval operation on the strip (21.21).

\begin{assessment}[Progress after compact V22]
\label{ass:n1v22-status}
The missing noncommutative local calculus has been constructed and
executed.  It is anchored to the V21 endpoint by all
\(384\cdot45\) first jets and the complete quadratic hash, and it
produces exact nonzero symmetric cubic and quartic derivatives.
There is no longer an ambiguity about link order, inverse-link
differentiation, colour normalization, or owner Taylor
coefficients.  The remaining finite acquisition is the
momentum-dependent harmonic/stationary contraction of these jets,
followed by a Ward residual and interval integration.
\end{assessment}

\begin{decision}[V23 contraction target]
\label{dec:n1v22-next}
Add the rational dual colour basis \(g^{-1}\) and the four retained
coordinate columns to the verifier.  At a fixed rational torus
point, construct \(A,B,L=-A^{-1}B\) exactly, contract the cubic and
quartic owner tensors into \(R,Q\), and check the stationary
identities before interval momentum is introduced.  Use the
quadratic hash (22.10) and the sample values (22.12)--(22.13) as
hard regression tests.
\end{decision}

\begin{firewall}[Claim boundary after compact V22]
\label{fire:n1v22-claim}
V22 supplies the exact local third/fourth word-jet engine.  It does
not yet assemble the harmonic and stationary response matrices,
check the complete background Ward sum, evaluate the edge mean,
complete SCIC or ADMC, prove summable cylinder comparison,
establish sector tightness, construct the continuum measure or
terminal clustering, or prove the Yang--Mills mass gap.
\end{firewall}

\subsection{Parent record}

This V22 note was formed by appending the preceding material to the
validated compact V21, whose record was
\[
\begin{array}{c|c}
\text{lines}&9723\\
\text{bytes}&329896\\
\text{SHA--256}&
\texttt{CA468DD3181CAE731FC57635C8B69497B0D456C5E0407D6FEEFDEE6FC79DB9D9}.
\end{array}
\]

% =============================================================================
% END APPENDED COMPACT-SERIES V22 MATERIAL
% =============================================================================

% =============================================================================
% BEGIN APPENDED COMPACT-SERIES V23 MATERIAL
% =============================================================================

\section{Compact-series V23: exact harmonic routers at all
sixteen parity corners}
\label{sec:n1v23}

\subsection{The rational corner problem}

The first coordinate layer in the V22 decision is the harmonic
lift.  At a general momentum it is an analytic rational matrix
function.  At the sixteen parity corners
\[
 {\cal K}_{\rm c}:=
 \{k\in{\mathbb T}^4:e^{ik_\mu}\in\{1,-1\}
 \text{ for every }\mu\},                                   \tag{23.1}
\]
all Laurent entries are integers.  This permits an exact
large-matrix check before interval arithmetic is introduced.

Split the V21 incidence and owner-difference matrices into reduced
and second-constituent columns:
\[
 D=(D_R\ D_C\ D_T),\qquad
 K=(K_R\ K_C\ K_T).                                         \tag{23.2}
\]
The tree columns are fixed to zero.  After clearing the Wilson
factor \(1/3\), define
\[
\begin{array}{lll}
 A_{\rm W}^{c}=D_R^*D_R,&
 B_{\rm W}^{c}=D_R^*D_C,&
 C_{\rm W}^{c}=D_C^*D_C,\\
 A_*^{c}=A_{\rm W}^{c}+3K_R^*K_R,&
 B_*^{c}=B_{\rm W}^{c}+3K_R^*K_C,&
 C_*=\frac13C_{\rm W}^{c}+K_C^*K_C.
\end{array}                                                   \tag{23.3}
\]
The exact routers are
\[
 L_{\rm W}=-(A_{\rm W}^{c})^{-1}B_{\rm W}^{c},\qquad
 L_*=-(A_*^{c})^{-1}B_*^{c}.                                \tag{23.4}
\]

\begin{lemma}[Clearing the normalization is exact]
\label{lem:n1v23-clearing}
Equations (23.3)--(23.4) are precisely the Wilson and
\({\cal V}_{\rm W}+{\cal O}_1\) harmonic equations at every
corner.
\end{lemma}

\begin{proof}
The Wilson blocks are
\(\frac13D_R^*D_R\) and \(\frac13D_R^*D_C\); multiplication of
both sides of the router equation by three does not alter its
solution.  The endpoint adds the Gram blocks of \(K^*K\).
Multiplication of its equation by three yields the second row of
(23.3).
\end{proof}

\subsection{Executable exact solve}

The verifier is
\[
 \texttt{scripts/verify\_ym\_v23\_router.py}.                 \tag{23.5}
\]
Its record is
\[
\begin{array}{c|c}
\text{lines}&329\\
\text{bytes}&11204\\
\text{SHA--256}&
\texttt{079B932D30B1786C966DF941B656C2B331005A953EC3F502CBB6125CEF3A7F60}.
\end{array}                                                   \tag{23.6}
\]
For each corner the program evaluates the Laurent matrices
exactly, performs two \(45\times45\) Gauss--Jordan eliminations
over \({\mathbb Q}\), and constructs both four-by-four Schur
matrices.  Its canonical aggregate output hash is
\[
\boxed{
 \texttt{C4BF7BFB11888993FCD75C4D0E03D6FAD3CAF0DD067700F53FF9D34F2308BACD}.}
                                                                    \tag{23.7}
\]

\begin{theorem}[All corner router and Schur residuals pass]
\label{thm:n1v23-corner-residuals}
At every \(k\in{\cal K}_{\rm c}\), exact rational arithmetic gives
\[
\boxed{\begin{aligned}
 A_{\rm W}^{c}L_{\rm W}+B_{\rm W}^{c}&=0,\\
 A_*^{c}L_*+B_*^{c}&=0,
\end{aligned}}                                               \tag{23.8}
\]
and both Schur matrices
\[
 {\mathsf S}_{\rm W}
 =\frac13\left(
 C_{\rm W}^{c}+(B_{\rm W}^{c})^*L_{\rm W}
 \right),                                                    \tag{23.9}
\]
\[
 {\mathsf S}_*
 =C_*+
 \left(\frac13B_{\rm W}^{c}+K_R^*K_C\right)^*L_*             \tag{23.10}
\]
are exactly symmetric.  The Wilson and endpoint Schur ranks are
zero at \(k=0\) and three at each of the other fifteen corners.
\end{theorem}

\begin{proof}
The program substitutes each of the sixteen sign vectors in
(23.1), solves (23.4), and reduces every rational entry to lowest
terms.  It then computes all entries of the residuals (23.8) and
their Schur transposition defects.  Every residual numerator is
zero.  Exact row reduction gives the stated ranks.
\end{proof}

\begin{corollary}[The corner physical kernel has the expected
Ward rank]
\label{cor:n1v23-Ward-rank}
At every nonzero parity corner the retained four-direction kernel
has one Ward direction and three nonzero transverse directions,
for both endpoints.  At the toron corner the entire commuting
constant four-plane is flat.
\end{corollary}

\begin{proof}
The exact ranks are those of
Theorem~\ref{thm:n1v23-corner-residuals}.  Gauge covariance gives
the retained gradient null direction.  A four-by-four symmetric
matrix of rank three has no additional null direction.  At zero,
rank zero means all four constant directions vanish in the
quadratic form.
\end{proof}

\subsection{The toron corner}

At
\[
                         z=(1,1,1,1),                        \tag{23.11}
\]
the two routers coincide.  Their common \(45\times4\) matrix has
exactly twenty-eight entries equal to one and 152 entries equal to
zero.  Its canonical hash is
\[
 \texttt{C6ABDD14F92A50DD05FA306BAC032F2092662873ABEDE0218EB0D4AAB8AFFF4B}.
                                                                    \tag{23.12}
\]

\begin{theorem}[Exact toron preservation by the endpoint]
\label{thm:n1v23-toron}
At (23.11),
\[
\boxed{\begin{aligned}
 L_*&=L_{\rm W},\\
 D_RL_{\rm W}+D_C&=0,\\
 K_RL_*+K_C&=0,\\
 {\mathsf S}_{\rm W}&={\mathsf S}_*=0.
\end{aligned}}                                               \tag{23.13}
\]
\end{theorem}

\begin{proof}
All entries of the four displayed residuals are evaluated as
rational numbers by the independent incidence and router
constructions.  Their reduced numerators vanish identically.
The second equation says that the Wilson harmonic lift of a
constant retained link field has zero fine curl.  Applying the
plaquette-difference rows to that zero-curvature lift gives the
third equation.  The two Schur identities also follow by inserting
the corresponding zero vectors into the Gram forms.
\end{proof}

\begin{remark}[The toron was not deleted]
\label{rem:n1v23-toron-included}
The zero corner is solved with the same invertible forty-five
dimensional vertical matrix as every other corner.  Only the
retained Schur kernel vanishes.  Thus (23.13) is not produced by
removing an internal zero mode or by assigning a limiting value to
an undefined inverse.
\end{remark}

\subsection{Endpoint routing differs away from the toron}

The exact solves also show
\[
 L_*=L_{\rm W}
 \quad\Longleftrightarrow\quad
 z=(1,1,1,1)\ \hbox{or}\ z=(-1,-1,-1,-1)                    \tag{23.14}
\]
within the sixteen-corner set.  At the other fourteen corners the
two routers differ.  Across all thirty-two router matrices, the
largest reduced denominator is \(18172\), the largest absolute
numerator is \(13193\), and every tested equation remains exact.

\begin{warning}[Freezing the Wilson router would fail]
\label{warn:n1v23-router-change}
Although the two actions have the same physical kinetic two-jet,
their harmonic lifts differ at fourteen of the sixteen parity
corners.  Therefore the endpoint background vertices cannot be
formed by inserting the Wilson lift into the new owner words.
The router derivative is a genuine part of
\(R_*-R_{\rm W}\) and \(Q_*-Q_{\rm W}\).
\end{warning}

\begin{proof}
Equation (23.14) and the rational bounds are direct properties of
the serialized exact solutions.  The warning follows because a
background Hessian differentiated after substituting \(L(k)\)
contains the derivative of that substitution; replacing \(L_*\)
by \(L_{\rm W}\) changes the endpoint family at the displayed
corners.
\end{proof}

\begin{firewall}[Corner exactness is not a torus certificate]
\label{fire:n1v23-corners}
The sixteen exact solves test phase, normalization, toron
inclusion, and the physical Ward rank at the most important
rational points.  They do not bound the routers between corners,
certify a complete background Ward identity, or enclose the
Brillouin-zone mean.  Those require the complex strip and interval
matrix arithmetic.
\end{firewall}

\subsection{A verified interval entry point}

For a complex matrix \(A(k)\) with real-axis inverse bound \(g_0\),
let an interval box centered at \(k_c\) have radius \(\delta\).
If
\[
 \|A(k)-A(k_c)\|\le\eta,\qquad
 \|A(k_c)^{-1}\|\eta<1,                                     \tag{23.15}
\]
then
\[
 A(k)^{-1}
 =
 \left[I+A(k_c)^{-1}(A(k)-A(k_c))\right]^{-1}A(k_c)^{-1}     \tag{23.16}
\]
and
\[
 \|A(k)^{-1}\|
 \le
 \frac{\|A(k_c)^{-1}\|}
 {1-\|A(k_c)^{-1}\|\eta}.                                   \tag{23.17}
\]

\begin{proposition}[Finite interval cover criterion]
\label{prop:n1v23-cover}
Let \(M_A\) and \(R_A\) be coefficient sum and range stocks for
the endpoint path.  Every torus box satisfying
\[
 M_A(e^{R_A\delta}-1)<g_0^{-1}                              \tag{23.18}
\]
admits the certified inverse enclosure (23.16)--(23.17).
A finite uniform grid with such \(\delta\) covers
\({\mathbb T}^4\).
\end{proposition}

\begin{proof}
The Laurent estimate gives (23.15) with
\(\eta=M_A(e^{R_A\delta}-1)\).  Condition (23.18) makes the
Neumann series converge and proves (23.16)--(23.17).  Compactness
of the torus and a positive admissible \(\delta\) give a finite
grid cover.
\end{proof}

\begin{assessment}[Progress after compact V23]
\label{ass:n1v23-status}
The harmonic-coordinate layer is now exact at all sixteen parity
corners.  Both endpoint router equations and Schur symmetries pass,
the retained kernels have precisely the expected Ward ranks, and
the toron is included with an exact common lift and zero curvature.
The calculation also proves that endpoint router motion is
unavoidable away from two corners.  Together with the exact V22
word engine and the cover criterion (23.18), the next acquisition
is no longer algebraic ambiguity but interval contraction of the
mixed vertices on a finite momentum cover.
\end{assessment}

\begin{decision}[V24 mixed-vertex checkpoint]
\label{dec:n1v23-next}
At the toron corner, where the router is the integer matrix
(23.12), contract all 384 cubic and quartic owner jets against the
four retained directions and eight colour directions.  Include
the stationary-centre and Haar rows, then test the complete
tadpole--bubble Ward response at zero external momentum.  This
exact rational checkpoint must pass before the construction is
extended to interval boxes away from the corner.
\end{decision}

\begin{firewall}[Claim boundary after compact V23]
\label{fire:n1v23-claim}
V23 certifies exact harmonic routing at the parity corners.  It
does not yet assemble the complete mixed response, check its
background Ward sum, evaluate the edge mean, complete SCIC or
ADMC, prove summable cylinder comparison, establish sector
tightness, construct the continuum measure or terminal
clustering, or prove the Yang--Mills mass gap.
\end{firewall}

\subsection{Parent record}

This V23 note was formed by appending the preceding material to the
validated compact V22, whose record was
\[
\begin{array}{c|c}
\text{lines}&10050\\
\text{bytes}&341433\\
\text{SHA--256}&
\texttt{775B384B03A010065CB77DAC9AA98A74818FEE306FE37F29C2C42C42BE2402CE}.
\end{array}
\]

% =============================================================================
% END APPENDED COMPACT-SERIES V23 MATERIAL
% =============================================================================

% =============================================================================
% BEGIN APPENDED COMPACT-SERIES V24 MATERIAL
% =============================================================================

\section{Compact-series V24: exact dual-colour compression of
the mixed-vertex contraction}
\label{sec:n1v24}

\subsection{Why the colour trace is an upstream acquisition}

A literal vertical response matrix has
\(45\cdot8=360\) rows.  Forming its dense bubble product before
using adjoint invariance would repeat the same colour contraction
for every spatial entry and obscure exact cancellations.  V24
therefore evaluates the dual-basis identities first.  This is an
upstream move: it does not omit a vertex, but changes the order of
an exact finite contraction.

Let \(g_{ab}\) be the Gram matrix (22.3), and let \(g^{ab}\) denote
its inverse.  Direct rational inversion gives
\[
 (g^{ab})=
 \begin{pmatrix}
 \frac12&0&0&0&0&0&0&0\\
 0&\frac12&0&0&0&0&0&0\\
 0&0&\frac23&0&0&0&0&\frac13\\
 0&0&0&\frac12&0&0&0&0\\
 0&0&0&0&\frac12&0&0&0\\
 0&0&0&0&0&\frac12&0&0\\
 0&0&0&0&0&0&\frac12&0\\
 0&0&\frac13&0&0&0&0&\frac23
 \end{pmatrix}.                                              \tag{24.1}
\]

\subsection{Completeness and Casimir identities}

\begin{theorem}[Exact fundamental completeness]
\label{thm:n1v24-completeness}
For the basis (22.1) and its dual metric (24.1),
\[
\boxed{
 \sum_{a,b=1}^8g^{ab}T_aT_b=-\frac83I_3.}                   \tag{24.2}
\]
\end{theorem}

\begin{proof}
Multiply the eight explicit matrices (22.1), weight them by the
sixty-four rational entries of (24.1), and add.  All off-diagonal
entries and the unequal diagonal parts cancel.  Each diagonal
entry is \(-8/3\).  The computation is also certified
independently below.
\end{proof}

\begin{theorem}[Exact adjoint Casimir]
\label{thm:n1v24-Casimir}
For every \(X\in\mathfrak{su}(3)\),
\[
\boxed{
 \sum_{a,b=1}^8g^{ab}[T_a,[T_b,X]]=-6X.}                    \tag{24.3}
\]
\end{theorem}

\begin{proof}
The left side is an adjoint-equivariant endomorphism of the simple
real Lie algebra \(\mathfrak{su}(3)\), so it is scalar.  Exact
evaluation on each of the basis elements (22.1) gives
\(-6T_c\).  Linearity proves (24.3).  The sign is negative because
the \(T_a\) are anti-Hermitian; the positive quadratic Casimir is
\(-\sum g^{ab}\operatorname{ad}(T_a)\operatorname{ad}(T_b)\).
\end{proof}

\begin{corollary}[No dense colour matrix is needed in the bubble]
\label{cor:n1v24-bubble}
Whenever the two cubic background vertices reduce to adjoint
commutators after their ordered spatial coefficients are fixed,
the complete internal colour contraction in the bubble is the
scalar factor six times the corresponding spatial contraction.
\end{corollary}

\begin{proof}
The vertical covariance contracts internal colour indices with
\(g^{ab}\).  Two commutator vertices therefore produce the
negative of the operator in (24.3); its value is \(6I\).
\end{proof}

\subsection{Ordered triple and quartic traces}

\begin{theorem}[Real cubic trace is the lowered bracket]
\label{thm:n1v24-triple}
For all basis indices,
\[
\boxed{
 \Re\Tr(T_aT_bT_c)
 =
 -\frac12
 \left\langle T_a,[T_b,T_c]\right\rangle_{\rm HS}.}          \tag{24.4}
\]
Exactly sixty of the \(8^3\) ordered triples have nonzero real
part.
\end{theorem}

\begin{proof}
Cyclicity gives
\[
\begin{split}
 \left\langle T_a,[T_b,T_c]\right\rangle_{\rm HS}
 &=-\Re\Tr\left(T_aT_bT_c-T_aT_cT_b\right).
\end{split}
\]
For anti-Hermitian matrices,
\[
 \overline{\Tr(T_aT_cT_b)}
 =
 -\Tr(T_bT_cT_a)
 =
 -\Tr(T_aT_bT_c),
\]
where the last equality is cyclicity.  Hence the two real parts
in the bracket are opposites, proving (24.4).  Enumeration in the
exact basis gives the stated count.
\end{proof}

\begin{theorem}[One contracted quartic ordering]
\label{thm:n1v24-quartic}
For every \(c,d\),
\[
\boxed{
 \sum_{a,b=1}^8g^{ab}
 \Re\Tr(T_aT_bT_cT_d)
 =
 \frac83g_{cd}.}                                             \tag{24.5}
\]
\end{theorem}

\begin{proof}
Use (24.2):
\[
\begin{split}
 \sum_{a,b}g^{ab}\Re\Tr(T_aT_bT_cT_d)
 &=
 -\frac83\Re\Tr(T_cT_d)
 =\frac83g_{cd}.
\end{split}
\]
\end{proof}

\begin{warning}[Take the real invariant contraction at the right
stage]
\label{warn:n1v24-imaginary-order}
An individual ordered trace
\(\Tr(T_aT_bT_c)\) or
\(\Tr(T_aT_bT_cT_d)\) can have a nonzero imaginary part.  It is
incorrect to require each ordered word to be real before adding
the conjugate owner terms.  The real Hilbert--Schmidt action and
the complete dual-colour contraction are real; the intermediate
ordered traces need not be.
\end{warning}

\subsection{Executable certificate}

The exact verifier is
\[
 \texttt{scripts/verify\_ym\_v24\_colour.py}.                 \tag{24.6}
\]
Its record is
\[
\begin{array}{c|c}
\text{lines}&227\\
\text{bytes}&7407\\
\text{SHA--256}&
\texttt{71313AFECACFEBACD73BED6FA64FC56E6185DA2090D7B87CF0CDC1EB5562C63E}.
\end{array}                                                   \tag{24.7}
\]
The canonical payload containing \(g^{-1}\), every commutator
coefficient, the eight Casimir images, the completeness matrix,
and the contracted quartic matrix has SHA--256
\[
\boxed{
 \texttt{BC13B912E71E7EAD878B9A7C50E633C085541717E7DE1CD71E1861A264978360}.}
                                                                    \tag{24.8}
\]

\begin{proposition}[Exhaustive exact checks]
\label{prop:n1v24-checks}
The verifier returns true for all of the following:
\begin{enumerate}
\item \(gg^{-1}=I_8\);
\item closure of all sixty-four commutators in the basis (22.1);
\item all \(8^3=512\) Jacobi identities;
\item the eight basis instances of (24.3);
\item the matrix identity (24.2);
\item all 512 instances of (24.4); and
\item all sixty-four instances of (24.5).
\end{enumerate}
\end{proposition}

\begin{proof}
Every matrix entry is a Gaussian rational and every metric
coefficient is rational.  The verifier compares reduced rational
numerators after exact matrix multiplication.  No tolerance or
floating-point branch occurs.
\end{proof}

\subsection{Compression of the V22 owner bank}

For a fixed background colour \(c\), write the endpoint cubic
vertex schematically as
\[
 R^{(c)}
 =
 \sum_\alpha
 {\cal R}_\alpha\otimes{\cal C}^{(c)}_\alpha,                 \tag{24.9}
\]
where \({\cal R}_\alpha\) is a \(45\times45\) spatial Laurent
matrix and \({\cal C}^{(c)}_\alpha\) is an ordered colour
operator generated by (22.5).  Similarly write
\[
 Q^{(c,d)}
 =
 \sum_\beta
 {\cal Q}_\beta\otimes{\cal D}^{(c,d)}_\beta.                 \tag{24.10}
\]

\begin{proposition}[Finite colour-first response compiler]
\label{prop:n1v24-colour-first}
The complete tadpole and bubble colour traces formed from
(24.9)--(24.10) can be evaluated exactly before any spatial
matrix inverse:
\begin{enumerate}
\item replace every real cubic ordered trace using (24.4);
\item contract paired commutators using (24.3);
\item reduce adjacent internally contracted quartic words using
      (24.5); and
\item evaluate the finitely many remaining quartic orderings
      directly with (22.1) and (24.1).
\end{enumerate}
The output is a finite rational linear combination of spatial
\(45\times45\) trace words.
\end{proposition}

\begin{proof}
The V22 engine produces only products of at most four basis
matrices.  The covariance supplies only the dual metric
\(g^{ab}\).  Items one through three apply the proved identities.
There are finitely many permutations not in the adjacent form of
(24.5), and the exact eight-element table evaluates them.  Colour
coefficients are rational and are independent of the spatial
matrix inverse.  Thus the contraction terminates before spatial
interval arithmetic begins.
\end{proof}

\begin{corollary}[Dimension reduction]
\label{cor:n1v24-dimension}
The response calculation need not store or multiply a dense
\(360\times360\) matrix.  It is enough to retain the spatial
\(45\times45\) coefficient arrays and the finite rational colour
weights generated by
Proposition~\ref{prop:n1v24-colour-first}.
\end{corollary}

\begin{proof}
The tensor product decompositions (24.9)--(24.10), followed by the
complete colour contraction, leave only spatial indices.
\end{proof}

\begin{firewall}[Colour compression is not a Ward cancellation]
\label{fire:n1v24-not-Ward}
The identities above remove redundant colour arithmetic.  They do
not imply that the tadpole and bubble sum vanishes, that the edge
mean is zero, or that a toron contribution may be deleted.  The
spatial owner, router, stationary, Haar, and momentum contractions
must still be assembled in the declared order.
\end{firewall}

\begin{assessment}[Progress after compact V24]
\label{ass:n1v24-status}
The exact local jet engine and the harmonic router are now joined
by a certified dual-colour compiler.  Its completeness and
Casimir identities collapse the large colour trace to rational
weights on \(45\times45\) spatial words, while preserving every
ordered quartic term.  This materially reduces the remaining
toron checkpoint and removes a source of sign and normalization
errors.  The outstanding operation is the spatial contraction of
the 384 owners with the V23 router and the stationary/Haar rows.
\end{assessment}

\begin{decision}[Mandatory V25 audit and spatial contraction]
\label{dec:n1v24-next}
Before continuing, perform the five-version SM/NS companion
audit.  Then apply the colour-first compiler to the toron router
(23.12), assemble the exact spatial \(R,Q\) arrays, and test the
complete toron-inclusive Ward response.  If no companion supplies
a stronger contraction identity, retain the V22 owner ordering
and use sparse spatial arrays rather than a dense colour matrix.
\end{decision}

\begin{firewall}[Claim boundary after compact V24]
\label{fire:n1v24-claim}
V24 proves and certifies the exact dual-colour compression.  It
does not yet assemble the complete mixed response, check its
background Ward sum, evaluate the edge mean, complete SCIC or
ADMC, prove summable cylinder comparison, establish sector
tightness, construct the continuum measure or terminal
clustering, or prove the Yang--Mills mass gap.
\end{firewall}

\subsection{Parent record}

This V24 note was formed by appending the preceding material to the
validated compact V23, whose record was
\[
\begin{array}{c|c}
\text{lines}&10370\\
\text{bytes}&352619\\
\text{SHA--256}&
\texttt{A3068FC9B52381A599407F2C88476E4BCD9116CB95F2159A37039A2F581B17CA}.
\end{array}
\]

% =============================================================================
% END APPENDED COMPACT-SERIES V24 MATERIAL
% =============================================================================

% =============================================================================
% BEGIN APPENDED COMPACT-SERIES V25 MATERIAL
% =============================================================================

\section{Compact-series V25: compulsory companion audit and an
exact residual-inverse atlas seed}
\label{sec:n1v25}

\subsection{The V25 companion audit}

The current companion records are
\[
\begin{array}{c|r|r|l}
\text{programme and file}&\text{lines}&\text{bytes}&
\text{SHA--256}\\ \hline
\text{SM--Einstein V11}
&3323&107237&
\texttt{455C080F110A9352D5C593B14D97687C1A5CF6A9C1229C4406F51F93DEC47738}\\
\text{Navier--Stokes V26}
&5319&278283&
\texttt{82C887F8C2C69E4F28FAD7C3DC8DE5B9099CB5A4BF431EBA1FFF3E91935978AD}.
\end{array}                                                    \tag{25.1}
\]
The NS companion is unchanged and again supplies no determinant or
gauge-response construction.

The SM programme has restarted after its tenth archived compact
cycle and advanced to V11.  Its V10 audit imports the exact
Yang--Mills V21--V22 workflow: literal word generation, an
independent quadratic hash, a complete unreduced Ward check, and
separate determinant-measure rows.  Its new V11 result supplies
two transferable finite-gap certificates.  The relevant one here
is the residual inverse criterion
\[
 \|I-XH\|\le\eta<1
 \quad\Longrightarrow\quad
 \|H^{-1}\|\le\frac{\|X\|}{1-\eta}.                          \tag{25.2}
\]
It also records the finite-range weighted-conjugation estimate
needed to turn a certified local inverse into exponential kernel
decay.  No SM note supplies a Yang--Mills mixed owner array or edge
scalar.

\begin{audit}[Transfer and direction adjustment]
\label{audit:n1v25-transfer}
The V24 plan is adjusted as follows.
\begin{enumerate}
\item Preserve the exact colour-first spatial contraction.
\item Replace one globally crude inverse strip, when evaluating
      real momenta, by an adaptive residual-inverse atlas seeded
      with the exact V23 corner inverses.
\item Keep Haar, stationary-centre, toron, and determinant rows
      separate until the complete Ward sum is formed.
\item Use finite-range conjugation only after each local inverse
      residual has passed.
\end{enumerate}
This improves the numerical route but changes none of the open
physical claims.
\end{audit}

\subsection{Induced row stocks}

For a matrix Laurent polynomial \(P(k)=\sum_\alpha P_\alpha
e^{i\alpha\cdot k}\), define
\[
 {\cal M}_\infty(P)
 :=
 \max_i\sum_{j,\alpha}|(P_\alpha)_{ij}|,                     \tag{25.3}
\]
\[
 {\cal L}_\infty(P)
 :=
 \max_i\sum_{j,\alpha}
 |\alpha|_1|(P_\alpha)_{ij}|.                               \tag{25.4}
\]
Then
\[
 \|P(k)\|_\infty\le{\cal M}_\infty(P)                        \tag{25.5}
\]
and, for real \(k,k'\),
\[
 \|P(k)-P(k')\|_\infty
 \le{\cal L}_\infty(P)\|k-k'\|_\infty.                       \tag{25.6}
\]

\begin{proof}
Equation (25.5) is the induced row-sum estimate.  For each
monomial,
\[
 |e^{i\alpha\cdot k}-e^{i\alpha\cdot k'}|
 \le|\alpha|_1\|k-k'\|_\infty.
\]
Insert this in the row sum to obtain (25.6).
\end{proof}

The exact Laurent arrays from V21 give
\[
\begin{array}{c|ccc}
P&{\cal M}_\infty(P)&{\cal L}_\infty(P)&
\max_i\sum_{j,\alpha}|\alpha_\mu||(P_\alpha)_{ij}|\\ \hline
A=\frac13D_R^*D_R&\frac{23}{3}&\frac{10}{3}&\frac53
\quad(\mu=1,2,3,4)\\
A_1&348&210&(79,88,90,90)_\mu.
\end{array}                                                   \tag{25.7}
\]
In particular
\[
 {\cal L}_\infty(A_t)
 \le\frac{640}{3},\qquad
 \|\partial_tA_t\|_\infty
 \le348.                                                     \tag{25.8}
\]

\subsection{Exact corner inverse norms}

At every corner in (23.1), the V23 rational inverse matrices can
be measured in the induced infinity norm without interval
rounding.

\begin{proposition}[Endpoint corner bounds]
\label{prop:n1v25-corner-inverses}
For the sixteen exact corner inverses,
\[
\boxed{
 \max_{k_c\in{\cal K}_{\rm c}}\|A(k_c)^{-1}\|_\infty=53,}     \tag{25.9}
\]
\[
\boxed{
 \max_{k_c\in{\cal K}_{\rm c}}\|
 (A(k_c)+A_1(k_c))^{-1}\|_\infty
 =
 \frac{199572383}{72207277}<2.764.}                          \tag{25.10}
\]
\end{proposition}

\begin{proof}
The V25 verifier performs exact Gauss--Jordan inversion on the
same integer matrices used in V23, restores the factor three, and
takes exact rational row sums.  Maximization over the finite
sixteen-element set gives (25.9)--(25.10).
\end{proof}

\begin{remark}[A positive deformation can improve a nonunitarily
invariant norm]
\label{rem:n1v25-infinity}
The large difference between (25.9) and (25.10) is not itself a
spectral claim: the infinity norm depends on the chosen parity
basis.  It is nevertheless a valid and useful residual
certificate norm.
\end{remark}

\subsection{A certified two-parameter box}

Let \(t_c\in\{0,1\}\), \(k_c\in{\cal K}_{\rm c}\), and put
\[
 X_c=A_{t_c}(k_c)^{-1}.                                     \tag{25.11}
\]
For
\[
 \|k-k_c\|_\infty\le\delta_k,\qquad
 |t-t_c|\le\delta_t,                                        \tag{25.12}
\]
equations (25.6)--(25.8) give
\[
 \|I-X_cA_t(k)\|_\infty
 \le
 \|X_c\|_\infty
 \left(
 \frac{640}{3}\delta_k+348\delta_t
 \right).                                                    \tag{25.13}
\]

\begin{theorem}[Uniform half-residual seed boxes]
\label{thm:n1v25-seed-boxes}
The exact rational choices
\[
\boxed{
 \delta_k=\frac3{135680},\qquad
 \delta_t=\frac1{73776}}                                    \tag{25.14}
\]
make the right side of (25.13) at most \(1/2\) for every one of
the thirty-two endpoint-corner pairs.  Consequently
\[
\boxed{
 \|A_t(k)^{-1}\|_\infty
 \le2\|X_c\|_\infty\le106}                                  \tag{25.15}
\]
throughout each seed box.
\end{theorem}

\begin{proof}
Use the common corner bound
\(\|X_c\|_\infty\le53\).  The two terms in (25.13) become
\[
 53\frac{640}{3}\frac3{135680}=\frac14,\qquad
 53\cdot348\frac1{73776}=\frac14.
\]
Their sum is \(1/2\).  Apply (25.2) with
\(\eta=1/2\) to obtain (25.15).
\end{proof}

\begin{corollary}[Certified continuation rule]
\label{cor:n1v25-continuation}
Suppose a box centered at \((k_c,t_c)\) has a certified inverse
enclosure \(X_c\) and bound \(K_c\).  Any adjacent box satisfying
\[
 K_c\left(
 \frac{640}{3}\delta_k+348\delta_t
 \right)<1                                                  \tag{25.16}
\]
inherits an inverse certificate.  Re-centering with the newly
enclosed inverse gives a finite adaptive continuation algorithm
through any chain of boxes for which the strict residual margin is
maintained.
\end{corollary}

\begin{proof}
Equation (25.13) is local and does not require a parity-corner
centre.  The residual theorem (25.2) supplies the next inverse and
its bound.  Repetition proves the claim along a finite chain.
\end{proof}

\subsection{Weighted inverse after a residual pass}

Let \(A_B\) be a certified center matrix of range \(r_A\), with
\[
 \|A_B\|_{0,\infty}\le M_B,\qquad
 \|A_B^{-1}\|_\infty\le K_B.                                \tag{25.17}
\]

\begin{proposition}[Finite-range exponential kernel bound]
\label{prop:n1v25-weighted}
If
\[
 K_BM_B(e^{\mu r_A}-1)\le\frac12,                            \tag{25.18}
\]
then its position-space inverse obeys
\[
 \boxed{
 \|E_xA_B^{-1}E_y\|_\infty
 \le2K_Be^{-\mu d(x,y)}.}                                   \tag{25.19}
\]
\end{proposition}

\begin{proof}
Conjugate by the diagonal weight
\((W_yf)(x)=e^{\mu d(x,y)}f(x)\).  Finite range gives
\[
 \|W_yA_BW_y^{-1}-A_B\|_\infty
 \le M_B(e^{\mu r_A}-1).
\]
Condition (25.18) makes the inverse-relative perturbation at most
\(1/2\).  The Neumann series bounds the conjugated inverse by
\(2K_B\).  Its \(x,y\) block is
\(e^{\mu d(x,y)}E_xA_B^{-1}E_y\), proving (25.19).
\end{proof}

\begin{warning}[Seed boxes do not cover the comparison domain]
\label{warn:n1v25-no-cover}
The radii (25.14) certify nonempty neighborhoods and regression
tests.  Their union is far too small to cover
\({\mathbb T}^4\times[0,1]\).  A completed certificate must run
the adaptive continuation, record every accepted box and overlap,
and verify coverage.  The global spectral theorem proves that such
a cover exists, but it does not provide the interval response
values or their serialized atlas.
\end{warning}

\subsection{Executable record}

The stock and seed-box verifier is
\[
 \texttt{scripts/verify\_ym\_v25\_residual\_stocks.py}.       \tag{25.20}
\]
Its record is
\[
\begin{array}{c|c}
\text{lines}&145\\
\text{bytes}&4999\\
\text{SHA--256}&
\texttt{0C5467AEDC573430C2984CF327AAE0086209FCC367965CFA306A5D3D5E58148F}.
\end{array}                                                   \tag{25.21}
\]
It regenerates the V21 Laurent matrices rather than copying the
stocks (25.7), solves all endpoint-corner inverses exactly, and
outputs the rational radii (25.14).

\begin{assessment}[Progress after compact V25]
\label{ass:n1v25-status}
The required companion audit is complete and yields a concrete
improvement: exact corner inverses now seed a residual atlas, with
sharp induced row stocks and nonzero rational boxes.  The
finite-range conjugation theorem also connects each successful
inverse certificate to the weighted kernel decay needed in the
response and later projective estimates.  The colour-first V24
contraction remains the correct local algebra; it should now be
performed inside certified boxes rather than under a single very
crude global norm.
\end{assessment}

\begin{decision}[V26 sparse spatial contraction]
\label{dec:n1v25-next}
Contract the V22 cubic and quartic owners at the exact toron
router, using the V24 dual-colour identities and sparse spatial
indices.  Produce separate direct-owner, router,
stationary-centre, and Haar arrays, then verify their complete
pre-split Ward sum.  Use the V25 half-residual boxes as the first
interval regression cells after the exact toron checkpoint passes.
\end{decision}

\begin{firewall}[Claim boundary after compact V25]
\label{fire:n1v25-claim}
V25 performs the mandatory audit and certifies residual-inverse
atlas seeds.  It does not yet assemble the complete mixed response,
cover the full comparison domain, evaluate the edge mean, complete
SCIC or ADMC, prove summable cylinder comparison, establish sector
tightness, construct the continuum measure or terminal
clustering, or prove the Yang--Mills mass gap.
\end{firewall}

\subsection{Parent record}

This V25 note was formed by appending the preceding material to the
validated compact V24, whose record was
\[
\begin{array}{c|c}
\text{lines}&10693\\
\text{bytes}&363243\\
\text{SHA--256}&
\texttt{0162994668377362B5AA45715B0DB71C2C58141711CE9FA0E1F1CD626B400275}.
\end{array}
\]

% =============================================================================
% END APPENDED COMPACT-SERIES V25 MATERIAL
% =============================================================================

% =============================================================================
% BEGIN APPENDED COMPACT-SERIES V26 MATERIAL
% =============================================================================

\section{Compact-series V26: exact elimination of the
stationary and Haar rows on the commuting toron checkpoint}
\label{sec:n1v26}

\subsection{The common toron lift}

Let \(a\in{\mathbb R}^4\) be a retained constant one-form and
\(X\in\mathfrak{su}(3)\) a fixed colour direction.  At the zero
corner, form the fine scalar lift
\[
 b_R=L_{\rm W}a=L_*a,\qquad b_C=a,\qquad b_T=0,              \tag{26.1}
\]
using Theorem~\ref{thm:n1v23-toron}.  Thus
\[
                         Db=0,\qquad Kb=0.                  \tag{26.2}
\]
Put the compact links on the one-parameter background
\[
                         U_e(t)=\exp(t b_eX).                \tag{26.3}
\]

\begin{theorem}[The lifted toron is exactly flat, not merely
linearly flat]
\label{thm:n1v26-exact-flat}
Every fine plaquette of (26.3) is the identity for every real
\(t\).  Consequently
\[
 {\cal V}_{\rm W}(U(t))=0,\qquad
 {\cal O}_1(U(t))=0,\qquad
 {\cal V}_*(U(t))=0.                                        \tag{26.4}
\]
\end{theorem}

\begin{proof}
All link matrices in (26.3) belong to the same one-parameter
subgroup and therefore commute.  For a plaquette \(p\),
\[
 P_p(U(t))
 =
 \exp\left(
 tX\sum_{e\in\partial p}\epsilon_{p,e}b_e
 \right).
\]
The exponent is the \(p\)-th component of \(Db\), which is zero
by (26.2).  Hence \(P_p=I\).  Every Wilson plaquette defect
vanishes.  Both plaquettes in every covariant difference (20.3)
are the identity, so every \({\cal O}_1\) owner vanishes as well.
\end{proof}

\begin{corollary}[Exact endpoint preservation of the commuting
toron valley]
\label{cor:n1v26-toron-valley}
The endpoint deformation does not lift or tilt the one-colour
commuting toron family.  In particular, it adds neither a
classical toron potential nor a classical mass row.
\end{corollary}

\begin{proof}
Equation (26.4) is identically zero as a function of \(t\), so all
of its derivatives vanish.
\end{proof}

\subsection{Stationarity in translated fibre coordinates}

Let \(z=(z_e)_{e\in{\cal R}}\) denote reduced link variables around
the identity and translate them by
\[
                         (T_tz)_e=e^{tb_eX}z_e.              \tag{26.5}
\]
Use the retained links in (26.3) on the four \(C\)-columns and
keep tree links fixed.  Define
\[
 \widetilde{\cal V}_e(z,t)
 :=
 {\cal V}_e(T_tz,U_C(t)),
 \qquad e\in\{{\rm W},*\}.                                  \tag{26.6}
\]

\begin{theorem}[The translated stationary centre is identically
zero]
\label{thm:n1v26-centre-zero}
For both endpoints and all sufficiently small real \(t\),
\[
 D_z\widetilde{\cal V}_e(0,t)=0,\qquad
                         z_{*,e}(t)=0.                       \tag{26.7}
\]
Thus
\[
                         z_{*,e}'(0)=z_{*,e}''(0)=0.          \tag{26.8}
\]
\end{theorem}

\begin{proof}
At \(z=0\), the physical fine field in (26.6) is precisely the
flat connection (26.3).  Both actions are sums of nonnegative
squared plaquette defects or squared covariant plaquette
differences and vanish there.  Hence their first variation is
zero.  The vertical Hessians at \(t=0\) are positive by (20.7);
continuity preserves a positive vertical Hessian for small \(t\).
The implicit-function theorem gives a unique nearby stationary
centre.  Since \(z=0\) is stationary for every \(t\), uniqueness
gives (26.7), and differentiation gives (26.8).
\end{proof}

\begin{remark}[Why nonnegativity is used]
\label{rem:n1v26-nonnegative}
Gauge invariance alone would make the flat family critical for the
Wilson term, but the declared sum-of-squares completion makes the
same statement immediate for \({\cal O}_1\).  This is one reason
for choosing the positive endpoint in V20.
\end{remark}

\subsection{The Haar row vanishes in this chart}

\begin{theorem}[Exact product-Haar invariance]
\label{thm:n1v26-Haar}
The coordinate change (26.5) preserves the reduced product-Haar
measure for every \(t\).  Therefore its logarithmic Jacobian is
identically zero and contributes no first or second background
derivative to the toron response.
\end{theorem}

\begin{proof}
Each reduced factor is left multiplication by the fixed group
element \(e^{tb_eX}\).  Haar measure is left invariant on every
factor, and products of invariant measures are invariant.  Thus
the Jacobian equals one for every \(t\).
\end{proof}

\begin{corollary}[Coordinate-complete toron simplification]
\label{cor:n1v26-complete-simplification}
On the path (26.3), all stationary-centre chain-rule terms in
(20.22)--(20.23) and the Haar absorption (20.31) vanish for both
endpoints.  The complete toron background vertices are the direct
derivatives of the translated endpoint Hessians at \(z=0\).
\end{corollary}

\begin{proof}
The stationary terms contain \(z_*'\) or \(z_*''\), which vanish
by (26.8).  The measure is independent of \(t\) by
Theorem~\ref{thm:n1v26-Haar}.
\end{proof}

\subsection{Exact endpoint vertex arrays at the toron}

For reduced fluctuation directions \(u,v\), define
\[
 R_1(b)[u,v]
 :=
 D^3{\cal O}_1(0)[u,v,b],                                   \tag{26.9}
\]
\[
 Q_1(b)[u,v]
 :=
 D^4{\cal O}_1(0)[u,v,b,b],                                 \tag{26.10}
\]
where every insertion uses the literal translated link order
(26.5).  These are generated by (22.5)--(22.7).  Let
\(R_{\rm W}(b),Q_{\rm W}(b)\) denote the attached exact Wilson
suffix arrays in the same translated chart.

\begin{theorem}[No omitted toron rows]
\label{thm:n1v26-endpoint-vertices}
The complete endpoint vertices are exactly
\[
\boxed{
 R_*(b)=R_{\rm W}(b)+R_1(b),\qquad
 Q_*(b)=Q_{\rm W}(b)+Q_1(b).}                               \tag{26.11}
\]
No router, stationary, or Haar term must be added to (26.11).
\end{theorem}

\begin{proof}
Action differentiation is linear, giving the two direct sums.
The router is common by (26.1), the stationary derivatives vanish
by Theorem~\ref{thm:n1v26-centre-zero}, and the Haar row vanishes
by Theorem~\ref{thm:n1v26-Haar}.  These exhaust the
coordinate-complete rows identified in V20.
\end{proof}

\begin{corollary}[Endpoint-exact toron determinant response]
\label{cor:n1v26-response}
Let
\[
 {\mathbb A}_e=A_e\otimes g,\qquad
 {\mathbb G}_e={\mathbb A}_e^{-1},                           \tag{26.12}
\]
with the nonorthogonal colour metric \(g\) of (22.3).  For the
one-colour toron direction \(bX\),
\[
\boxed{
 \Gamma_e''(0)
 =
 \frac12\Tr\left(
 {\mathbb G}_e{\mathbb Q}_e(b)
 -
 {\mathbb G}_e{\mathbb R}_e(b)
 {\mathbb G}_e{\mathbb R}_e(b)
 \right),}                                                   \tag{26.13}
\]
and the endpoint comparison is
\[
 \Delta\Gamma''(0)=\Gamma_*''(0)-\Gamma_{\rm W}''(0),        \tag{26.14}
\]
with \(R_*,Q_*\) given by (26.11).
\end{corollary}

\begin{proof}
Twice differentiate
\(\frac12\log\det{\mathbb A}_e(t)\).
There is no density derivative by
Theorem~\ref{thm:n1v26-Haar}.  Equation (26.14) is subtraction of
the two exact endpoint identities.
\end{proof}

\begin{warning}[A finite toron determinant need not vanish]
\label{warn:n1v26-toron-potential}
Classical flatness (26.4) and the absence of a local mass row do
not force the finite-volume determinant response (26.13) to be
zero.  A global toron holonomy can have a finite-size effective
potential.  The toron value must be retained in the momentum sum
and separated from a local \(F^2\) coefficient only by the
declared Ward/moment extractor.
\end{warning}

\subsection{A sparse rather than dense execution}

For one owner, the transported word and local plaquette together
meet at most ten fine-edge occurrences.  Hence a direct cubic or
quartic contribution with two vertical insertions has at most
\(10^2\) ordered spatial pairs before coincident-edge
cancellations.  Across all 384 owners this gives the crude
preassembly bound
\[
                         384\cdot10^2=38400                 \tag{26.15}
\]
on owner-pair records for each fixed background direction.

\begin{proposition}[Sparse exact contraction suffices]
\label{prop:n1v26-sparse}
For a fixed toron background polarization and colour, the arrays
(26.9)--(26.10) can be assembled from at most 38400 owner-pair
records, each carrying an exact rational colour weight after V24.
Accumulating equal reduced-edge pairs produces the complete
\(45\times45\) spatial arrays without constructing a dense
\(360\times360\) matrix.
\end{proposition}

\begin{proof}
A derivative of a finite word is zero unless every vertical
insertion meets an occurrence of its edge type.  There are at most
ten occurrences in the difference of the six- and four-factor
words.  Choose the two vertical occurrences, then sum over the
384 owners.  V24 contracts all internal colour labels exactly,
leaving only the reduced spatial pair.  Addition of records with
the same pair is precisely matrix assembly.
\end{proof}

\begin{firewall}[The exact reduction is not the executed array]
\label{fire:n1v26-not-executed}
V26 proves that the toron checkpoint has no hidden
stationary-centre, router, or Haar correction and bounds its sparse
workload.  It does not serialize \(R_1,Q_1\), evaluate
(26.13), or test the complete momentum Ward extractor.
\end{firewall}

\begin{assessment}[Progress after compact V26]
\label{ass:n1v26-status}
The toron checkpoint is substantially simplified.  The exact V23
lift gives a commuting flat connection for both endpoints, the
translated stationary centre is identically zero, and product
Haar measure contributes no density row.  Therefore the endpoint
difference is determined solely by the direct V22 cubic/quartic
owners and the already attached Wilson suffix arrays, with colour
compressed by V24.  The remaining task is a bounded sparse
assembly rather than an unresolved coordinate-completion problem.
\end{assessment}

\begin{decision}[V27 serialized toron arrays]
\label{dec:n1v26-next}
Extend the V22 engine so one background direction may be a full
fine lift \(b\), not a single edge type.  Assemble and serialize
the sparse spatial arrays \(R_1(b)\) and \(Q_1(b)\) for each of the
four toron polarizations, verify their skew/symmetric types and
hypercubic orbit relations, and replay selected entries directly
from the word formula before evaluating (26.13).
\end{decision}

\begin{firewall}[Claim boundary after compact V26]
\label{fire:n1v26-claim}
V26 removes the toron coordinate-completion ambiguity.  It does
not execute the mixed arrays, evaluate the edge mean, complete
SCIC or ADMC, prove summable cylinder comparison, establish sector
tightness, construct the continuum measure or terminal
clustering, or prove the Yang--Mills mass gap.
\end{firewall}

\subsection{Parent record}

This V26 note was formed by appending the preceding material to the
validated compact V25, whose record was
\[
\begin{array}{c|c}
\text{lines}&11030\\
\text{bytes}&374092\\
\text{SHA--256}&
\texttt{8124B18E09A4A425B322A5F3283D024D8B6B4F3632B18A48DA972E5AD222188E}.
\end{array}
\]

% =============================================================================
% END APPENDED COMPACT-SERIES V26 MATERIAL
% =============================================================================

% =============================================================================
% BEGIN APPENDED COMPACT-SERIES V27 MATERIAL
% =============================================================================

\section{Compact-series V27: exact sparse assembly of one
covariant-endpoint toron vertex card}
\label{sec:n1v27}

\subsection{Declared polarization and normalization}

Take the first retained spatial polarization in (26.1) and the
background colour \(T_1\) of (22.1).  Its Hilbert--Schmidt norm is
\[
                         g_{11}=2.                           \tag{27.1}
\]
The exact toron lift has eight nonzero fine-edge types, all with
coefficient one; every other type has coefficient zero.  For
reduced spatial indices \(i,j\), define the cubic spatial
coefficient \({\cal R}_1\) by
\[
 D^3{\cal O}_1[e_iT_a,e_jT_b,bT_1]
 =
 ({\cal R}_1)_{ij}
 \left\langle T_a,[T_1,T_b]\right\rangle_{\rm HS}.           \tag{27.2}
\]
Define the unit-colour-traced quartic coefficient by
\[
 ({\cal Q}_1)_{ij}
 :=
 \frac1{g_{11}}
 \sum_{a,b=1}^8g^{ab}
 D^4{\cal O}_1[e_iT_a,e_jT_b,bT_1,bT_1].                    \tag{27.3}
\]
Thus (27.3) already includes the complete internal colour trace
and divides by the squared norm of the chosen background colour.

\subsection{Sparse exact algorithm}

The V22 engine is extended so that a derivative direction may be
a scalar field on all sixty-four fine-edge types rather than a
single edge type.  At a word factor based on edge type \(e\), such
a direction inserts its coefficient \(b_eT_a\).  The assignment
formula (22.5) is unchanged.

For each owner, only reduced edge types actually occurring in its
six- and four-factor words are retained.  Equal spatial pairs are
accumulated immediately.  The program visits
\[
                         12435                              \tag{27.4}
\]
owner-pair records, instead of forming a dense colour-spatial
matrix.

\begin{lemma}[Two independent cubic projections]
\label{lem:n1v27-two-projections}
For \(X=T_1\), the two exact lowered-bracket components
\[
 \left\langle T_2,[T_1,T_3]\right\rangle_{\rm HS}=-4,\qquad
 \left\langle T_2,[T_1,T_8]\right\rangle_{\rm HS}=2          \tag{27.5}
\]
are nonzero.  Dividing the corresponding complete cubic owner
sums by the two values in (27.5) gives identical
\(45\times45\) rational matrices.
\end{lemma}

\begin{proof}
The bracket values follow by direct multiplication of (22.1).
The verifier assembles (27.2) once from each component in (27.5)
without sharing the resulting spatial matrix.  Entrywise
comparison of the two rational arrays returns equality.
\end{proof}

\begin{theorem}[Executed toron endpoint arrays]
\label{thm:n1v27-arrays}
The matrices defined by (27.2)--(27.3) have integer entries and
obey
\[
\boxed{
 {\cal R}_1^T=-{\cal R}_1,\qquad
 {\cal Q}_1^T={\cal Q}_1.}                                  \tag{27.6}
\]
Their exact inventories are
\[
\begin{array}{c|rrr|l}
 &\text{nonzero entries}&
 \max|\text{numerator}|&
 \max\text{ denominator}&\text{SHA--256}\\ \hline
{\cal R}_1&706&14&1&
\texttt{9C376E5DC8C6AEC32877E23BDD32B08602D15B94C607B5F262291A93FDDD8593}\\
{\cal Q}_1&830&84&1&
\texttt{CD4DA2C234ADBD85DCDBA0FB76325BD4915FB24EF16B2CCB0BBE89683B26BD49}.
\end{array}                                                   \tag{27.7}
\]
\end{theorem}

\begin{proof}
Each owner derivative is evaluated by the exact rational
assignment and subset formulas (22.5)--(22.7).  The cubic
coefficient is independently projected through both nonzero
components (27.5); the arrays agree by
Lemma~\ref{lem:n1v27-two-projections}.  The quartic row sums the
ten nonzero entries of the dual metric (24.1) and divides by
\(g_{11}=2\).  Direct transposition tests give (27.6).
Canonical row-major serialization of reduced rational entries
gives the hashes in (27.7).
\end{proof}

\begin{corollary}[The direct endpoint toron bank is populated]
\label{cor:n1v27-populated}
For the first toron polarization, the direct completion terms
\(R_1(b)\) and \(Q_1(b)\) in (26.11) are no longer formal finite
sums; they are the explicit serialized arrays (27.7).
\end{corollary}

\subsection{Executable provenance}

The verifier is
\[
 \texttt{scripts/verify\_ym\_v27\_toron\_vertices.py}.        \tag{27.8}
\]
Its record is
\[
\begin{array}{c|c}
\text{lines}&344\\
\text{bytes}&11978\\
\text{SHA--256}&
\texttt{F453E4368E57D0433C4256446B872AFD9C26A1E0147B887EFC4DEA80F8AD766F}.
\end{array}                                                   \tag{27.9}
\]
The script regenerates the incidence, exact zero-corner Wilson
router, background lift, dual colour metric, owner words, and all
mixed derivatives.  It does not read a precomputed \(R\) or
\(Q\) matrix.

\begin{proposition}[Type checks follow from the scalar Hessian]
\label{prop:n1v27-types}
The skew type of \({\cal R}_1\) and symmetric type of
\({\cal Q}_1\) are necessary consequences of symmetry of the
full vertical Hessian in the two fluctuation slots.
\end{proposition}

\begin{proof}
Interchanging the fluctuation directions in the cubic derivative
does not change the scalar action derivative.  The colour tensor
\langle T_a,[T_1,T_b]\rangle\) is skew in \(a,b\), so its
spatial coefficient must also be skew.  In the quartic row the
two vertical colour indices are contracted with the symmetric
dual metric and the action derivative is symmetric in the
vertical slots, leaving a symmetric spatial coefficient.
\end{proof}

\begin{firewall}[One polarization is not yet the hypercubic
card]
\label{fire:n1v27-one-polarization}
V27 executes the first constant spatial polarization.  The
parity-tree coordinate section is not manifestly hypercubic, so
the other three arrays may not be declared to be literal
permutations until their restoring tree transformations or their
independent hashes are checked.
\end{firewall}

\subsection{Insertion into the endpoint trace}

Let \(R_{\rm W},Q_{\rm W}\) be the attached Wilson suffix arrays
for the same unit-colour toron lift.  The complete endpoint cubic
operator has spatial coefficient
\[
 {\cal P}_*:=\frac13R_{\rm W}+{\cal R}_1,                    \tag{27.10}
\]
because the Wilson cubic Hessian is
\(\frac13R_{\rm W}\otimes\operatorname{ad}X\).  The complete
colour-traced endpoint quartic coefficient is
\[
                         {\cal Q}_*:=Q_{\rm W}+{\cal Q}_1.    \tag{27.11}
\]

\begin{proposition}[Colour-reduced endpoint formula]
\label{prop:n1v27-endpoint-formula}
For a unit-norm background colour and the first toron
polarization,
\[
\boxed{
 \Gamma_*''(0)
 =
 \frac12\Tr(G_*{\cal Q}_*)
 +3\Tr(G_*{\cal P}_*G_*{\cal P}_*),}                         \tag{27.12}
\]
where
\[
                         G_*=(A+A_1)^{-1}.                   \tag{27.13}
\]
The Wilson endpoint is
\[
 \Gamma_{\rm W}''(0)
 =
 \frac12\Tr(G_{\rm W}Q_{\rm W})
 +\frac13\Tr(G_{\rm W}R_{\rm W}
                   G_{\rm W}R_{\rm W}).                     \tag{27.14}
\]
\end{proposition}

\begin{proof}
The tadpole colour trace has already been taken in
(27.3) and in the attached definition of \(Q_{\rm W}\).  For the
bubble, Theorem~\ref{thm:n1v24-Casimir} gives
\(\Tr_{\rm ad}(\operatorname{ad}X)^2=-6\) for unit \(X\).
The factor \(-1/2\) in the second derivative of the logarithmic
determinant therefore becomes \(+3\), proving (27.12).
Setting \({\cal R}_1={\cal Q}_1=0\) and using the Wilson factor
\(1/3\) gives (27.14).
\end{proof}

\begin{assessment}[Progress after compact V27]
\label{ass:n1v27-status}
The first completion toron card is now an executed exact object.
Two independent cubic colour projections agree, the full
dual-colour quartic contraction is complete, the required
skew/symmetric types pass, and canonical hashes fix both arrays.
Together with the exact Wilson suffix formulas and the V23
endpoint inverse, equation (27.12) reduces the first toron
determinant comparison to rational \(45\times45\) arithmetic.
\end{assessment}

\begin{decision}[V28 exact toron determinant difference]
\label{dec:n1v27-next}
Generate \(R_{\rm W},Q_{\rm W}\) from the attached suffix formulas
on the same eight-edge toron lift, verify their declared types,
and evaluate (27.12)--(27.14) exactly.  Record tadpole and bubble
parts separately, retain their signs, and do not infer a local
mass term from the finite toron value.
\end{decision}

\begin{firewall}[Claim boundary after compact V27]
\label{fire:n1v27-claim}
V27 executes one completion toron vertex card.  It does not yet
evaluate its determinant response, verify the other spatial
polarizations, evaluate the edge mean, complete SCIC or ADMC,
prove summable cylinder comparison, establish sector tightness,
construct the continuum measure or terminal clustering, or prove
the Yang--Mills mass gap.
\end{firewall}

\subsection{Parent record}

This V27 note was formed by appending the preceding material to the
validated compact V26, whose record was
\[
\begin{array}{c|c}
\text{lines}&11348\\
\text{bytes}&385324\\
\text{SHA--256}&
\texttt{FCFE0FB979D045C358449F2D207EB0416031C913DCE77FEA515B6B8355B8C19C}.
\end{array}
\]

% =============================================================================
% END APPENDED COMPACT-SERIES V27 MATERIAL
% =============================================================================

% =============================================================================
% BEGIN APPENDED COMPACT-SERIES V28 MATERIAL
% =============================================================================

\section{Compact-series V28: exact endpoint-minus-Wilson
determinant response at the first toron polarization}
\label{sec:n1v28}

\subsection{Wilson suffix arrays in the identical lift}

For each oriented Wilson plaquette with edge occurrences
\((e_i,\epsilon_i)\), \(1\le i\le4\), use the common toron lift
\(b\) from V27 and define
\[
 f=\sum_{i=1}^4\epsilon_ib_{e_i},                            \tag{28.1}
\]
\[
 s_i=\sum_{j>i}\epsilon_jb_{e_j}
 +{\bf1}_{\{\epsilon_i=-1\}}\epsilon_ib_{e_i},\qquad
 d_{ij}=s_i-s_j-\frac f2.                                   \tag{28.2}
\]
The exact attached Wilson matrices are generated occurrence by
occurrence:
\[
 (R_{\rm W})_{e_ie_j}
 =\epsilon_i\epsilon_jd_{ij},\qquad
 (R_{\rm W})_{e_je_i}
 =-\epsilon_i\epsilon_jd_{ij},                              \tag{28.3}
\]
\[
 (Q_{\rm W})_{e_ie_i}=-\frac89f^2,                          \tag{28.4}
\]
\[
 (Q_{\rm W})_{e_ie_j}
 =(Q_{\rm W})_{e_je_i}
 =
 \epsilon_i\epsilon_j
 \left(-2d_{ij}^2-\frac7{18}f^2\right).                     \tag{28.5}
\]
Repeated edge types are accumulated on their actual reduced
matrix indices.  On the exact toron lift, \(f=0\) on every
plaquette, but it is retained in the generator as an independent
residual check.

\begin{theorem}[Executed Wilson toron arrays]
\label{thm:n1v28-Wilson-arrays}
For the first toron polarization,
\[
 R_{\rm W}^T=-R_{\rm W},\qquad Q_{\rm W}^T=Q_{\rm W}.         \tag{28.6}
\]
Their inventories are
\[
\begin{array}{c|rrr|l}
 &\text{nonzero entries}&
 \max|\text{numerator}|&
 \max\text{ denominator}&\text{SHA--256}\\ \hline
R_{\rm W}&90&1&1&
\texttt{35AA5495F57CCADEA0DD19D828C81E1291015144F2A21AD4A01BDD2F827AF93C}\\
Q_{\rm W}&90&2&1&
\texttt{F80E3577CEAF38B13453762272B9AE63AF11892C1BFD3F73D72D1D0B60EC200F}.
\end{array}                                                   \tag{28.7}
\]
\end{theorem}

\begin{proof}
The generator applies (28.1)--(28.5) to the ninety-six literal
plaquettes constructed independently from the incidence routine.
It accumulates rational entries, tests transposition, and
serializes the full \(45\times45\) matrices.  The type identities
pass exactly and give the hashes in (28.7).
\end{proof}

\subsection{Exact rational inverse and trace evaluation}

At the zero corner define
\[
 G_{\rm W}=A^{-1},\qquad
 G_*=(A+A_1)^{-1},                                          \tag{28.8}
\]
\[
 {\cal P}_*=\frac13R_{\rm W}+{\cal R}_1,\qquad
 {\cal Q}_*=Q_{\rm W}+{\cal Q}_1,                            \tag{28.9}
\]
where the completion arrays are the independently regenerated
V27 arrays.  Their hashes are checked against (27.7) before any
trace is evaluated.

Separate the four rational response pieces:
\[
 T_{\rm W}:=\frac12\Tr(G_{\rm W}Q_{\rm W}),\qquad
 B_{\rm W}:=\frac13\Tr(G_{\rm W}R_{\rm W}
                            G_{\rm W}R_{\rm W}),             \tag{28.10}
\]
\[
 T_*:=\frac12\Tr(G_*{\cal Q}_*),\qquad
 B_*:=3\Tr(G_*{\cal P}_*G_*{\cal P}_*).                     \tag{28.11}
\]

\begin{theorem}[Exact first-polarization toron response]
\label{thm:n1v28-response}
The four contributions and their endpoint totals are
\[
\boxed{
 T_{\rm W}=\frac{441}{8},\qquad
 B_{\rm W}=-\frac{197}{4},\qquad
 \Gamma_{\rm W}''(0)=\frac{47}{8},}                         \tag{28.12}
\]
\[
\boxed{
 T_*=\frac{194816091}{4713800},\qquad
 B_*=-\frac{14704054060106}{347186100625},}                 \tag{28.13}
\]
\[
\boxed{
 \Gamma_*''(0)
 =
 -\frac{2841921261373}{2777488805000}.}                     \tag{28.14}
\]
Consequently
\[
\boxed{
 \Gamma_*''(0)-\Gamma_{\rm W}''(0)
 =
 -\frac{4789916997687}{694372201250}
 \in(-6.899,-6.898).}                                       \tag{28.15}
\]
\end{theorem}

\begin{proof}
The matrices in (28.8) are obtained by exact rational
Gauss--Jordan inversion, retaining the Wilson factor \(1/3\).
Equations (28.10)--(28.11) are then finite rational matrix
products and traces.  Direct reduction to lowest terms gives
(28.12)--(28.14).  Subtraction gives the exact fraction in
(28.15); rational comparison with the two displayed decimals
gives its interval.
\end{proof}

\begin{proposition}[Both bubble signs are certified]
\label{prop:n1v28-bubble-sign}
\[
                         B_{\rm W}\le0,\qquad B_*\le0.        \tag{28.16}
\]
\end{proposition}

\begin{proof}
For a positive symmetric \(G\) and a real skew matrix \(R\),
\[
 \Tr(GRGR)
 =
 -\|G^{1/2}RG^{1/2}\|_{\rm HS}^2\le0.                       \tag{28.17}
\]
Apply (28.17) to \(R_{\rm W}\) and \({\cal P}_*\), using the
positive coefficients \(1/3\) and \(3\) in
(28.10)--(28.11).  The exact negative values in
(28.12)--(28.13) independently confirm the sign.
\end{proof}

\subsection{Reproducible certificate}

The verifier is
\[
 \texttt{scripts/verify\_ym\_v28\_toron\_response.py}.        \tag{28.18}
\]
Its record is
\[
\begin{array}{c|c}
\text{lines}&340\\
\text{bytes}&11991\\
\text{SHA--256}&
\texttt{31F6B9E0156A6D7F9205B65931422DCE451C69C4E8AA535DA5015A34806F62BA}.
\end{array}                                                   \tag{28.19}
\]
The canonical payload containing all four \(45\times45\) vertex
arrays and all response fractions has SHA--256
\[
\boxed{
 \texttt{72A6BFC479DE1972DEFE37681DFF10FAA67C949E0E48928AB58DEA1389F49978}.}
                                                                    \tag{28.20}
\]
The following checks all return true:
\begin{enumerate}
\item both V27 completion hashes replay;
\item \(R_{\rm W},{\cal R}_1\) are skew;
\item \(Q_{\rm W},{\cal Q}_1\) are symmetric; and
\item both bubble contributions are nonpositive.
\end{enumerate}

\begin{firewall}[The nonzero toron number is not a mass]
\label{fire:n1v28-not-mass}
Equation (28.15) is the response of one finite parity-cell
internal toron mode to one constant commuting background
polarization.  It is not the external-momentum second moment, the
Brillouin-zone average, the constant alias class of V17, or a
gauge-boson mass.  It must enter the toron-inclusive finite sum
with its exact weight; it cannot be promoted to, or deleted as, a
local continuum coefficient.
\end{firewall}

\subsection{A toron consistency equation for the full grid}

Let \({\cal W}_M(q)\) be the complete finite-grid response before
the low/edge split, and let the internal momentum grid include
\(k=0\).  Write
\[
 {\cal W}_M(q)
 =
 \frac1{M^4}{\cal W}_{0,M}(q)
 +
 \frac1{M^4}\sum_{k\ne0}{\cal W}_{k,M}(q).                  \tag{28.21}
\]

\begin{proposition}[The exact toron row is a regression value]
\label{prop:n1v28-regression}
For the first constant polarization at \(q=0\), the endpoint
difference of the first term in (28.21), before any external
moment extraction, must equal
\[
 -\frac1{M^4}
 \frac{4789916997687}{694372201250}                         \tag{28.22}
\]
in the normalization of V28.  Any interval grid implementation
whose zero-cell enclosure excludes (28.22) has a router, colour,
Haar, or suffix-order error.
\end{proposition}

\begin{proof}
The zero-grid point carries the uniform Riemann-sum weight
\(M^{-4}\).  Its endpoint response is exactly (28.15), and V26
proves that no stationary or Haar term is missing in this chart.
\end{proof}

\begin{warning}[External differentiation remains]
\label{warn:n1v28-external-q}
The beta coefficient uses the second derivative in external
momentum \(q\), not merely the value at \(q=0\).  V28 validates the
zero-cell value of the response integrand.  The derivatives of
the background lift, phases, and shifted inverse
\(G(k+q\widehat e_1)\) must still be included when computing the
edge class.
\end{warning}

\begin{assessment}[Progress after compact V28]
\label{ass:n1v28-status}
The first toron endpoint comparison is fully executed in exact
rational arithmetic.  Wilson and completion arrays are separately
hashed, both bubble signs pass, and the endpoint-minus-Wilson
response is the nonzero fraction (28.15).  This supplies a hard
zero-cell regression value for the eventual interval grid and
confirms that toron deletion would change the finite comparison.
The next finite task is to differentiate the toron card in
external momentum and to verify the remaining spatial
polarizations.
\end{assessment}

\begin{decision}[V29 external-momentum toron jet]
\label{dec:n1v28-next}
Differentiate the generated Laurent incidence, both routers, and
the V22 mixed owner arrays with respect to one external axial
momentum at \(q=0\).  Use implicit differentiation of
\(AL+B=0\) rather than finite differences.  Produce exact first
and second \(q\)-jets of the zero-cell endpoint response and check
reflection forces the first derivative to vanish.  Keep the
finite-grid toron weight explicit.
\end{decision}

\begin{firewall}[Claim boundary after compact V28]
\label{fire:n1v28-claim}
V28 evaluates one exact toron endpoint response.  It does not
compute its external-momentum second derivative, verify all
polarizations, evaluate the edge mean, complete SCIC or ADMC,
prove summable cylinder comparison, establish sector tightness,
construct the continuum measure or terminal clustering, or prove
the Yang--Mills mass gap.
\end{firewall}

\subsection{Parent record}

This V28 note was formed by appending the preceding material to the
validated compact V27, whose record was
\[
\begin{array}{c|c}
\text{lines}&11607\\
\text{bytes}&394540\\
\text{SHA--256}&
\texttt{D441E998FB881E84239CD3C572446AA34CA2C5A6CA60C4B0C38F91A6C315594F}.
\end{array}
\]

% =============================================================================
% END APPENDED COMPACT-SERIES V28 MATERIAL
% =============================================================================

% =============================================================================
% BEGIN APPENDED COMPACT-SERIES V29 MATERIAL
% =============================================================================

\section{Compact-series V29: exact first and second axial jets
of both harmonic routers}
\label{sec:n1v29}

\subsection{Differentiating the router equation}

Let \(q\) be axial external momentum in the first coordinate and
write
\[
 A_e(q)L_e(q)+B_e(q)=0,\qquad
 e\in\{{\rm W},*\}.                                         \tag{29.1}
\]
For any matrix \(F(q)\), put
\[
                         F^{[n]}:=\partial_q^nF(0).           \tag{29.2}
\]
The zero-corner inverses are \(G_e=(A_e^{[0]})^{-1}\).

\begin{lemma}[Exact router recursion]
\label{lem:n1v29-router-recursion}
\[
\boxed{
 L_e^{[0]}=-G_eB_e^{[0]},}                                  \tag{29.3}
\]
\[
\boxed{
 L_e^{[1]}
 =-G_e\left(
 A_e^{[1]}L_e^{[0]}+B_e^{[1]}
 \right),}                                                   \tag{29.4}
\]
\[
\boxed{
 L_e^{[2]}
 =-G_e\left(
 A_e^{[2]}L_e^{[0]}
 +2A_e^{[1]}L_e^{[1]}
 +B_e^{[2]}
 \right).}                                                   \tag{29.5}
\]
\end{lemma}

\begin{proof}
Equation (29.3) is (29.1) at zero.  Differentiate (29.1) once
and twice, then multiply the resulting equations by \(G_e\).
\end{proof}

For a Laurent monomial \(z^\alpha=e^{i\alpha\cdot k}\), exact
axial differentiation gives
\[
 (z^\alpha)^{[0]}=1,\qquad
 (z^\alpha)^{[1]}=i\alpha_1,\qquad
 (z^\alpha)^{[2]}=-\alpha_1^2.                              \tag{29.6}
\]
Thus all input jets in (29.3)--(29.5) lie in
\({\mathbb Q}(i)\).

\subsection{Endpoint block derivatives}

Let
\[
 H_{RR}(q)=K_R(q)^*K_R(q),\qquad
 H_{RC}(q)=K_R(q)^*K_C(q).                                  \tag{29.7}
\]
After clearing the Wilson factor \(1/3\), the two router systems
are
\[
 A_{\rm W}^{c}=D_R^*D_R,\qquad
 B_{\rm W}^{c}=D_R^*D_C,                                    \tag{29.8}
\]
\[
 A_*^{c}=A_{\rm W}^{c}+3H_{RR},\qquad
 B_*^{c}=B_{\rm W}^{c}+3H_{RC}.                             \tag{29.9}
\]
Their derivatives are formed by the product rule before the
router recursion is applied.  For example,
\[
\begin{split}
 (D_R^*D_C)^{[2]}
 ={}&(D_R^{[2]})^*D_C^{[0]}
 +2(D_R^{[1]})^*D_C^{[1]}\\
 &+(D_R^{[0]})^*D_C^{[2]}.                                 \tag{29.10}
\end{split}
\]

\begin{theorem}[All differentiated router residuals vanish]
\label{thm:n1v29-residuals}
For both \(e={\rm W},*\), exact Gaussian-rational arithmetic gives
\[
\boxed{\begin{aligned}
 A_e^{[0]}L_e^{[0]}+B_e^{[0]}&=0,\\
 A_e^{[1]}L_e^{[0]}+A_e^{[0]}L_e^{[1]}+B_e^{[1]}&=0,\\
 A_e^{[2]}L_e^{[0]}+2A_e^{[1]}L_e^{[1]}
 +A_e^{[0]}L_e^{[2]}+B_e^{[2]}&=0.
\end{aligned}}                                               \tag{29.11}
\]
\end{theorem}

\begin{proof}
The verifier generates \(D,K\) over the integer Laurent ring,
applies (29.6) coefficient by coefficient, forms the product jets
(29.8)--(29.10), and evaluates (29.3)--(29.5) with the exact
zero-corner inverse.  Substitution into (29.11) leaves zero
Gaussian rationals in every entry.
\end{proof}

\subsection{Reflection types}

\begin{theorem}[Exact axial parity of the router jets]
\label{thm:n1v29-reflection}
For both endpoints,
\[
\boxed{
 L_e^{[0]},L_e^{[2]}\in{\mathbb Q}^{45\times4},\qquad
 L_e^{[1]}\in i{\mathbb Q}^{45\times4}.}                    \tag{29.12}
\]
\end{theorem}

\begin{proof}
The Laurent coefficients are real.  Equation (29.6) is real at
even order and purely imaginary at odd order.  Products in
(29.8)--(29.10) preserve this parity.  The zero-corner inverse is
real, so the recursion (29.3)--(29.5) preserves it.  The verifier
also checks every output entry directly.
\end{proof}

\begin{corollary}[First response derivatives can have a forced
imaginary representation]
\label{cor:n1v29-first-imaginary}
When a reflection-even scalar response combines the \(q\) and
\(-q\) contributions, every term linear in \(L^{[1]}\) cancels
with its reflected conjugate.  A nonzero real first derivative
would therefore be a phase or orientation residual.
\end{corollary}

\begin{proof}
Reflection sends \(q\) to \(-q\) and complex conjugates the
Fourier matrices.  By (29.12), the order-one coefficient changes
sign, while the scalar reflection-even sum is unchanged.
\end{proof}

\subsection{A new cancellation and the first endpoint difference}

\begin{theorem}[The endpoint first router jet is unchanged]
\label{thm:n1v29-first-equality}
\[
\boxed{
 L_*^{[0]}=L_{\rm W}^{[0]},\qquad
 L_*^{[1]}=L_{\rm W}^{[1]},\qquad
 L_*^{[2]}\ne L_{\rm W}^{[2]}.}                             \tag{29.13}
\]
\end{theorem}

\begin{proof}
The first two matrix coefficient lists agree entry by entry in
exact arithmetic and have identical hashes below.  The second-
order lists have different hashes and their difference contains a
nonzero rational entry.
\end{proof}

\begin{corollary}[The completion router response begins at
second axial order]
\label{cor:n1v29-router-order}
At the zero corner,
\[
                         L_*(q)-L_{\rm W}(q)=O(q^2).          \tag{29.14}
\]
Equivalently, the completion blocks obey the exact first-order
compatibility identity
\[
 H_{RR}^{[1]}L^{[0]}
 +H_{RR}^{[0]}L^{[1]}
 +H_{RC}^{[1]}=0,                                           \tag{29.15}
\]
where \(L^{[j]}\) is the common Wilson/endpoint jet for
\(j=0,1\).
\end{corollary}

\begin{proof}
Taylor's theorem and (29.13) give (29.14).  Subtract the first-
derivative Wilson router equation from the endpoint equation in
(29.11), use (29.9), and divide by three to obtain (29.15).
\end{proof}

\begin{remark}[Consistency with irrelevance]
\label{rem:n1v29-irrelevance}
Equation (29.14) is stronger than the classical statement that
\({\cal O}_1\) begins at fourth physical momentum order.  It says
that the harmonic coordinate change itself is identical through
first axial order at the toron.  The second-order difference is
genuine and must enter the external-momentum response.
\end{remark}

\subsection{Exact inventories and hashes}

The output inventories are
\[
\begin{array}{c|c|rrr|l}
e&\text{jet}&\text{nonzero}&
\max|\text{numerator}|&\max\text{ denominator}&
\text{SHA--256}\\ \hline
{\rm W}&L^{[0]}&28&1&1&
\texttt{FC03F50C84C78D8605F099A478D9DDA98F420C9567109540ADF283B4B7386FE3}\\
{\rm W}&L^{[1]}&36&1&4&
\texttt{DE6205C6D7DDDD73983C2E77EC04C36B1520AE30F55BE0AC02F23FBFD04DAA05}\\
{\rm W}&L^{[2]}&74&19&24&
\texttt{C41717DACBFA42FEABE44746733D0A38920DF5A0B0D1569349A1888077461791}\\
*&L^{[0]}&28&1&1&
\texttt{FC03F50C84C78D8605F099A478D9DDA98F420C9567109540ADF283B4B7386FE3}\\
*&L^{[1]}&36&1&4&
\texttt{DE6205C6D7DDDD73983C2E77EC04C36B1520AE30F55BE0AC02F23FBFD04DAA05}\\
*&L^{[2]}&74&47611&288600&
\texttt{D6E6F8021D84969AE23DE342B281445056813A4CA1AEA1FAE178048315916631}.
\end{array}                                                   \tag{29.16}
\]

The verifier is
\[
 \texttt{scripts/verify\_ym\_v29\_router\_jets.py},           \tag{29.17}
\]
with record
\[
\begin{array}{c|c}
\text{lines}&274\\
\text{bytes}&9179\\
\text{SHA--256}&
\texttt{AAB0E08B2C2C2B7A4B5BC96D6455FC19323AE2106AF6016EC7D0C278A954F233}.
\end{array}                                                   \tag{29.18}
\]
The canonical payload containing all six matrices in (29.16) has
SHA--256
\[
\boxed{
 \texttt{2E5C2E755499E8BBC8AEB22C0B9E931165A9FFA23C46B40392D182E8A70C112E}.}
                                                                    \tag{29.19}
\]

\begin{firewall}[Router jets are not response jets]
\label{fire:n1v29-router-only}
V29 differentiates the harmonic lift.  The external-momentum
derivative of (28.21) also differentiates owner phases, the
background insertions, the shifted internal inverse, and the
Wilson suffix quantities.  Equality of \(L^{[1]}\) does not make
the first response derivative vanish by itself; the complete
reflection sum must be formed.
\end{firewall}

\begin{assessment}[Progress after compact V29]
\label{ass:n1v29-status}
The external-momentum coordinate input at the toron is now exact
through second order.  Both endpoint recursions and all
differentiated residuals pass.  Reflection parity is explicit,
and the positive completion unexpectedly leaves the first router
jet unchanged, so the first coordinate difference occurs only in
the rational second jet.  The remaining toron calculation is to
propagate these six matrices through the sparse owner and
log-determinant response derivatives.
\end{assessment}

\begin{decision}[Mandatory V30 audit and response jet]
\label{dec:n1v29-next}
Perform the five-version SM/NS companion audit.  Then assemble the
first and second external-\(q\) derivatives of the Wilson and
endpoint toron responses.  Use (29.11) as the router regression,
prove the complete first derivative vanishes after reflection
recombination, and serialize the exact second derivative before
it is inserted with weight \(M^{-4}\) into the full momentum sum.
\end{decision}

\begin{firewall}[Claim boundary after compact V29]
\label{fire:n1v29-claim}
V29 executes the axial harmonic-router jets.  It does not yet
execute the external response derivative, verify all
polarizations, evaluate the edge mean, complete SCIC or ADMC,
prove summable cylinder comparison, establish sector tightness,
construct the continuum measure or terminal clustering, or prove
the Yang--Mills mass gap.
\end{firewall}

\subsection{Parent record}

This V29 note was formed by appending the preceding material to the
validated compact V28, whose record was
\[
\begin{array}{c|c}
\text{lines}&11900\\
\text{bytes}&404317\\
\text{SHA--256}&
\texttt{9F47D7238B580EF7BE78FF34B3B1733CE6314531F76FBE3434F0764E88232977}.
\end{array}
\]

% =============================================================================
% END APPENDED COMPACT-SERIES V29 MATERIAL
% =============================================================================

% =============================================================================
% BEGIN APPENDED COMPACT-SERIES V30 MATERIAL
% =============================================================================

\section{Compact-series V30: mandatory companion audit and an exact
independent-leg Laurent compiler}
\label{sec:n1v30}

\subsection{The five-version companion checkpoint}

The current active companion notes at this checkpoint are
\[
\begin{array}{c|r|r|l}
\text{programme}&\text{version}&\text{bytes}&\text{SHA--256}\\ \hline
\text{SM--Einstein}&19&181945&
\texttt{E797F3D40BF86D131DCD52FE67C8BD74258F4336CE35ADEED4D25C6CC278E39E}\\
\text{Navier--Stokes}&26&278283&
\texttt{82C887F8C2C69E4F28FAD7C3DC8DE5B9099CB5A4BF431EBA1FFF3E91935978AD}.
\end{array}                                                   \tag{30.1}
\]
They contain \(5856\) and \(5319\) lines respectively.

\begin{audit}[Mathematical transfer at V30]
\label{audit:n1v30-companions}
The SM note has advanced from V11 to V19.  Its relevant new results
are: an exact boundary zero mode invalidating a naive Dirichlet gap;
a tree-gauge periodic completion on good contractible blocks; retention
of finite flat-holonomy modes before alias estimation; complementary
minor formulas through massless determinant zeros; and a quaternionic
determinant frame on the commuting flat locus.  Its explicit firewall
is that the last frame does not extend to curved noncommuting fields.

The NS note remains at V26.  It computes exact rank-one norms of first
and second Oseen-kernel derivatives.  The rank-one restriction is exact
for the first derivative but gives only a localized modest improvement
for the second, and it does not close the global regularity estimate.
\end{audit}

\begin{proposition}[Consequences for the present route]
\label{prop:n1v30-audit-transfer}
The audit implies the following three rules for the Yang--Mills
calculation.
\begin{enumerate}
\item A toron or flat-sector simplification must be retained with its
finite weight before an edge or alias estimate; it cannot be deleted
because a cofinal mean is expected to vanish.
\item A one-polarization toron calculation is a regression certificate,
not a substitute for a uniform transverse-polarization estimate.
\item Restricted-locus symmetry may simplify a finite block, but every
phase and momentum label required off that locus must be kept before
differentiation.
\end{enumerate}
\end{proposition}

\begin{proof}
The first and third statements are exactly the order-of-operations and
restricted-locus firewalls in SM V16--V19.  The second is the logical
content of the NS comparison between unrestricted symmetric-traceless
and rank-one input norms.  None of the companion results supplies a
Yang--Mills cylinder comparison or a mass-gap estimate.
\end{proof}

\subsection{Why the response calculation must go one step upstream}

The V22 word engine attached a single Laurent variable to every
direction and V27 evaluated that variable at the identity.  This is
sufficient for the exact number in V28.  It is not sufficient for an
external-momentum derivative: the two internal legs and one or two
background legs carry different momenta.  Assigning one common
character before differentiating would identify distinct chain-rule
terms.

Let
\[
 \Phi_m:=({\mathbb Z}^4)^m,\qquad
 \boldsymbol\zeta^{\boldsymbol\alpha}
 :=
 \prod_{\ell=1}^m\zeta_\ell^{\alpha_\ell},
 \quad
 \boldsymbol\alpha=(\alpha_1,\ldots,\alpha_m)\in\Phi_m.       \tag{30.2}
\]
Every differentiated leg now has its own character \(\zeta_\ell\).

Consider an oriented link factor \(f\).  The parity decomposition gives
its edge type \(e_f\), coarse-cell offset \(r_f\in{\mathbb Z}^4\), and
orientation \(\epsilon_f\in\{-1,1\}\).  For labelled directions
\[
                         X_\ell=v_\ell T_{a_\ell},            \tag{30.3}
\]
and a set \(S\) of labels assigned to this factor, define
\[
\begin{split}
 {\cal E}_{f,S}
 :={}&
 \frac{\epsilon_f^{|S|}}{|S|!}
 \prod_{\ell\in S}v_\ell(e_f)
 \sum_{\pi\in{\mathfrak S}(S)}
 T_{a_{\pi(1)}}\cdots T_{a_{\pi(|S|)}}\\
 &\hspace{35mm}\times
 \prod_{\ell\in S}\zeta_\ell^{r_f}.                          \tag{30.4}
\end{split}
\]
For \(S=\varnothing\), this is the identity.  Formula (30.4) is the
mixed derivative of the oriented link exponential, including its
Fourier translation.

For an ordered word \(W=(f_1,\ldots,f_s)\), sum the ordered product of
(30.4) over all assignments of the requested labels to the factors.
Call the result \(J_{W,S}\).  For an owner difference
\[
                         Y=W_{\rm transported}-W_{\rm lower}, \tag{30.5}
\]
put \(J_{Y,S}=J_{W_{\rm transported},S}-J_{W_{\rm lower},S}\).
The \(m\)-linear derivative of the owner energy is then
\[
 {\cal V}_m(X_1,\ldots,X_m)
 =
 \frac12
 \sum_{\varnothing\ne S\subsetneq[m]}
 \operatorname{ReTr}\!
 \left(J_{Y,S}^{\dagger}J_{Y,S^c}\right).                   \tag{30.6}
\]
The adjoint in (30.6) reverses the colour word and sends the phase
\(\boldsymbol\alpha\) to \(-\boldsymbol\alpha\).

\begin{theorem}[Exact independent-leg word formula]
\label{thm:n1v30-word-formula}
Equation (30.6) is the mixed derivative
\[
 \partial_{t_1}\cdots\partial_{t_m}
 \frac12\|Y(\exp(\sum_\ell t_\ell X_\ell))\|_{\rm HS}^2
 \bigg|_{t=0},                                               \tag{30.7}
\]
as an element of the real Laurent ring
\({\mathbb Q}[\Phi_m]\).  It is equivariant under every simultaneous
permutation of direction labels and their four-dimensional phase
slots.
\end{theorem}

\begin{proof}
Repeated Leibniz differentiation of one exponential gives (30.4):
the factorial is the exponential coefficient and the permutation sum
is the polarization of a noncommutative power.  A second Leibniz
expansion over the factors of \(W\) gives \(J_{W,S}\).  Finally,
differentiating \(\frac12\operatorname{ReTr}(Y^\dagger Y)\) distributes
a nonempty proper subset of labels to the first factor and its
complement to the second, which is precisely (30.6).  Relabelling the
derivative variables permutes both the direction data and the
corresponding offsets, leaving the mixed derivative unchanged.
\end{proof}

\begin{corollary}[Exact zero-momentum specialization]
\label{cor:n1v30-zero-specialization}
The homomorphism
\[
 \varepsilon_0:{\mathbb Q}[\Phi_m]\longrightarrow{\mathbb Q},
 \qquad
 \varepsilon_0(\boldsymbol\zeta^{\boldsymbol\alpha})=1,      \tag{30.8}
\]
sends (30.6) to the collapsed zero-momentum derivative used in V27.
\end{corollary}

\begin{proof}
Under (30.8), every link translation character equals one.  The
remaining orientation, field, colour-word, factorial, partition, and
Hilbert--Schmidt factors are term-for-term those of the V27 formula.
\end{proof}

\subsection{The complete cubic symbol}

Use the first retained spatial polarization and background colour
\(T_1\) of V27.  For the two nonzero colour projections used there,
let
\[
 C^{(a,b)}_{ij}(\zeta_1,\zeta_2,\zeta_3)
 :=
 \frac{
 \sum_{\rm owners}{\cal V}_3(e_iT_a,e_jT_b,bT_1)}
 {\langle T_a,[T_1,T_b]\rangle}.                             \tag{30.9}
\]

\begin{theorem}[Full colour projection and leg-exchange certificate]
\label{thm:n1v30-cubic}
For the two projections
\[
 (a,b,\langle T_a,[T_1,T_b]\rangle)=(2,3,-4),(2,8,2),        \tag{30.10}
\]
with one-based colour labels, the matrices (30.9) agree coefficient by
coefficient in \({\mathbb Q}[\Phi_3]\).  Their common value \(C\)
satisfies
\[
 C_{ij}(\zeta_1,\zeta_2,\zeta_3)
 =
 -C_{ji}(\zeta_2,\zeta_1,\zeta_3),                          \tag{30.11}
\]
and
\[
 \varepsilon_0(C)={\cal R}_1,\qquad
 \operatorname{SHA256}(\varepsilon_0(C))
 =
 \texttt{9C376E5DC8C6AEC32877E23BDD32B08602D15B94C607B5F262291A93FDDD8593}.
                                                                    \tag{30.12}
\]
\end{theorem}

\begin{proof}
The verifier applies (30.4)--(30.6) to all \(384\) owner words and all
\(12435\) locally supported ordered edge pairs.  It compares the two
colour projections in the full Laurent ring, applies the internal-leg
transposition to every monomial, and finally applies (30.8).  All three
comparisons are exact rational equalities.  The last hash replays V27.
\end{proof}

The complete cubic inventory is
\[
\begin{array}{c|r}
\text{nonzero matrix entries}&974\\
\text{Laurent terms}&10380\\
\max|\text{numerator}|&27\\
\max\text{ denominator}&4\\
\max\text{ single-leg } \ell^1\text{ range}&2\\
\max\text{ total } \ell^1\text{ range}&6
\end{array}                                                    \tag{30.13}
\]
and its canonical SHA--256 is
\[
\boxed{
\texttt{97FE887E18F7296A4C65146D4477FC3F1F81C270BAC8A87275A0EBE017EF3F54}.}
                                                                    \tag{30.14}
\]
The rise from \(706\) nonzero entries after zero specialization to
\(974\) before specialization records genuine phase cancellations.

\subsection{A quartic independent regression panel}

The full zero-momentum quartic matrix was already certified in V27.
For the new phase engine, use the transpose-closed index set
\[
                         I=\{0,1,2,7,14,22,31,44\}.           \tag{30.15}
\]
For all \(64\) ordered pairs in \(I^2\), contract the two internal
colours with the exact inverse colour metric and use two copies of the
background direction.  Denote the resulting panel by \(Q_I\).

\begin{theorem}[Quartic phase regressions]
\label{thm:n1v30-quartic}
The panel satisfies
\[
 Q_{I,ij}(\zeta_1,\zeta_2,\zeta_3,\zeta_4)
 =
 Q_{I,ji}(\zeta_2,\zeta_1,\zeta_3,\zeta_4),                 \tag{30.16}
\]
\[
 Q_{I,ij}(\zeta_1,\zeta_2,\zeta_3,\zeta_4)
 =
 Q_{I,ij}(\zeta_1,\zeta_2,\zeta_4,\zeta_3).                 \tag{30.17}
\]
Moreover, every zero-momentum entry agrees exactly with an independent
evaluation by the collapsed V27 engine.
\end{theorem}

\begin{proof}
The program forms both sides of (30.16) and (30.17) as sparse rational
Laurent polynomials and compares their dictionaries.  It then sums the
coefficients of each entry and compares with the old scalar word
derivative.  All \(64\) comparisons pass.
\end{proof}

There are \(36\) nonzero panel entries and \(2881\) Laurent terms.  The
largest numerator and denominator are \(83\) and \(12\), the largest
single-leg range is \(2\), and the largest total range is \(8\).  The
panel and its zero specialization have hashes
\[
\begin{split}
 \operatorname{SHA256}(Q_I)
 &=
 \texttt{2E2B48251758BA216988D06BEE9F75C8E0FC19C521376E29363E801D4CBC2C69},\\
 \operatorname{SHA256}(\varepsilon_0Q_I)
 &=
 \texttt{A001DA6FC1BCB3D67AFE4333D9DA0E87B85835AAD4439CA50BB6AB3947EF69E8}.
                                                               \tag{30.18}
\end{split}
\]

\subsection{The exact axial specialization required by the toron card}

For \(\sigma=(\sigma_1,\ldots,\sigma_m)\in{\mathbb Z}^m\), define
\[
 {\mathfrak s}_\sigma
 \left(\boldsymbol\zeta^{\boldsymbol\alpha}\right)(q)
 :=
 \exp\left(
 iq\sum_{\ell=1}^m\sigma_\ell(\alpha_\ell)_1
 \right).                                                     \tag{30.19}
\]
If \(n_\sigma(\boldsymbol\alpha)\) denotes the integer in the sum, then
\[
 \partial_q{\mathfrak s}_\sigma
 \left(\boldsymbol\zeta^{\boldsymbol\alpha}\right)(0)
 =i n_\sigma(\boldsymbol\alpha),\qquad
 \partial_q^2{\mathfrak s}_\sigma
 \left(\boldsymbol\zeta^{\boldsymbol\alpha}\right)(0)
 =-n_\sigma(\boldsymbol\alpha)^2.                             \tag{30.20}
\]
Thus no finite differencing is needed.

In an all-incoming convention, the two cubic bubble vertices use
\[
 \sigma_+=(0,-1,1),\qquad \sigma_-=(0,1,-1),                 \tag{30.21}
\]
and the tadpole uses
\[
                         \sigma_Q=(0,0,1,-1).                \tag{30.22}
\]
The two background slots must additionally be composed with the V29
jets of the harmonic lift \(b_e(q)\).

Let \(A_e(q)\) be the reduced quadratic block and
\(G_e(q)=A_e(q)^{-1}\), with the endpoint normalization restored after
the cleared integer calculation.  Its exact inverse jets are
\[
 G_e^{[1]}=-G_eA_e^{[1]}G_e,                                \tag{30.23}
\]
\[
 G_e^{[2]}
 =
 2G_eA_e^{[1]}G_eA_e^{[1]}G_e-G_eA_e^{[2]}G_e.              \tag{30.24}
\]

\begin{proof}
Differentiate \(A_e(q)G_e(q)=I\) once and twice and solve on the left
with \(G_e(0)\).
\end{proof}

Write \(P_e^\pm(q)\) for the two specializations (30.21), including
the lifted background, and \(Q_e(q)\) for (30.22).  The toron card in
the normalization of V28 is
\[
\begin{split}
 {\cal W}_{0,e}(q)
 ={}&\frac12\Tr\!\left(G_e(0)Q_e(q)\right)\\
 &+c_e\Tr\!\left(
 G_e(0)P_e^+(q)G_e(q)P_e^-(q)
 \right),
 \qquad
 c_{\rm W}=\frac13,\quad c_*=3.                              \tag{30.25}
\end{split}
\]
At \(q=0\), this is exactly the V28 formula.

\begin{proposition}[Closed second-jet formula]
\label{prop:n1v30-second-jet}
With all unlabelled factors on the right evaluated at \(q=0\),
\[
\begin{split}
 {\cal W}_{0,e}^{[2]}
 ={}&\frac12\Tr(G_e Q_e^{[2]})\\
 &+c_e\Tr G_e\bigl(
 P_e^{+,[2]}G_eP_e^-
 +P_e^+G_e^{[2]}P_e^-
 +P_e^+G_eP_e^{-,[2]}\\
 &\hspace{19mm}
 +2P_e^{+,[1]}G_e^{[1]}P_e^-
 +2P_e^{+,[1]}G_eP_e^{-,[1]}
 +2P_e^+G_e^{[1]}P_e^{-,[1]}
 \bigr).                                                     \tag{30.26}
\end{split}
\]
The reflection-completed card
\[
 {\cal W}_{0,e}^{\rm sym}(q)
 :=
 \frac12\bigl({\cal W}_{0,e}(q)+{\cal W}_{0,e}(-q)\bigr)     \tag{30.27}
\]
obeys
\[
 ({\cal W}_{0,e}^{\rm sym})^{[1]}=0,\qquad
 ({\cal W}_{0,e}^{\rm sym})^{[2]}={\cal W}_{0,e}^{[2]}.      \tag{30.28}
\]
\end{proposition}

\begin{proof}
Apply the second-order Leibniz rule to the three varying factors in
the bubble trace; the first \(G_e(0)\) is fixed.  This gives the six
terms in (30.26).  Equation (30.27) is even by definition, so its first
derivative vanishes.  The two second derivatives at zero agree, which
gives (30.28).
\end{proof}

\begin{firewall}[Reflection is not a missing-term cure]
\label{fire:n1v30-reflection}
Equation (30.28) cancels the complete first derivative only after both
orientations have been assembled.  It does not permit deletion of an
individual imaginary first jet.  Formula (30.26) still requires every
owner phase, both background-lift derivatives, and the shifted inverse
derivatives.
\end{firewall}

\subsection{Executable certificate and status}

The verifier is
\[
             \texttt{scripts/verify\_ym\_v30\_multileg\_phases.py}. \tag{30.29}
\]
Its record is
\[
\begin{array}{c|c}
\text{lines}&523\\
\text{bytes}&17684\\
\text{SHA--256}&
\texttt{8D0E6907392D0DAF563FA64158DEA9B88101272A1B14CBD3D9D8950DE0741D4F}.
\end{array}                                                    \tag{30.30}
\]
The canonical payload containing the full cubic symbol and quartic
panel has SHA--256
\[
\boxed{
\texttt{6CAE6E15C8D18ED4D66DAD08E6283BAC259F91B62B385594DFCDADAA14AAAB74}.}
                                                                    \tag{30.31}
\]
The two full colour projections, the V27 cubic replay, cubic skew
exchange, quartic internal exchange, quartic background exchange, and
all quartic-panel zero specializations pass exactly.

\begin{assessment}[Progress after compact V30]
\label{ass:n1v30-status}
The mandatory companion audit is complete.  More importantly, the
phase ambiguity blocking the external toron derivative has been
removed: every differentiated leg now owns an independent Laurent
character, the complete cubic symbol is certified, a quartic panel
cross-checks the old engine, and the exact second-response formula is
fixed.  This is an upstream correction, not a numerical mass claim.
\end{assessment}

\begin{decision}[V31 direct axial response compiler]
\label{dec:n1v30-next}
Avoid serializing the entire four-leg quartic symbol.  Instead, apply
(30.20) during owner assembly, compose the two background slots with
the six V29 router matrices, compute (30.23)--(30.24), and evaluate
(30.26) in exact Gaussian-rational arithmetic.  Check (30.28)
independently and retain the factor \(M^{-4}\).
\end{decision}

\begin{firewall}[Claim boundary after compact V30]
\label{fire:n1v30-claim}
V30 certifies the multileg phase algebra and the response-jet formula.
It does not yet output the numerical second toron jet, verify all
external polarizations, evaluate the nonzero momentum edge mean,
complete SCIC or ADMC, prove summable cylinder comparison, establish
sector tightness, construct the continuum measure or terminal
clustering, or prove the Yang--Mills mass gap.
\end{firewall}

\subsection{Parent record}

This V30 note was formed by appending the preceding material to the
validated compact V29, whose record was
\[
\begin{array}{c|c}
\text{lines}&12200\\
\text{bytes}&414183\\
\text{SHA--256}&
\texttt{BED1ED939F606480DE9B341DD5BB6363275256150FC0A1753E355DA01E05B46F}.
\end{array}
\]

% =============================================================================
% END APPENDED COMPACT-SERIES V30 MATERIAL
% =============================================================================

% =============================================================================
% BEGIN APPENDED COMPACT-SERIES V31 MATERIAL
% =============================================================================

\section{Compact-series V31: translated-chart correction and exact
cubic bubble jets}
\label{sec:n1v31}

\subsection{A coordinate inconsistency exposed by the Wilson suffix}

V26 declared the reduced-fibre chart
\[
                         U_e(t,z)=e^{t b_eX}z_e.              \tag{31.1}
\]
For exponential vertical coordinates \(z_e=e^{Z_e}\), a positive link
factor is therefore \(e^{B_e}e^{Z_e}\), while its inverse is
\[
                         (e^{B_e}e^{Z_e})^{-1}
                         =e^{-Z_e}e^{-B_e}.                  \tag{31.2}
\]
The V27 engine instead placed background and vertical directions in a
single exponential and symmetrized all coincident insertions.  V30
faithfully retained the phases of that engine, but did not repair this
coordinate ordering.

\begin{lemma}[The two charts differ already in a mixed link jet]
\label{lem:n1v31-chart-difference}
For noncommuting \(B,Z\),
\[
 \partial_t\partial_s e^{tB+sZ}\big|_{0}
 =\frac12(BZ+ZB),\qquad
 \partial_t\partial_s(e^{tB}e^{sZ})\big|_{0}=BZ.             \tag{31.3}
\]
Their difference is \(\frac12[B,Z]\).  For the inverse translated
factor the ordered mixed coefficient is \(ZB\), with the appropriate
two negative signs.
\end{lemma}

\begin{proof}
Take the \(ts\)-coefficient in the three displayed exponential
products.  The exponential of a sum receives the two orderings with
coefficient \(1/2\); the product chart fixes their order.
\end{proof}

\begin{correction}[Scope of the V27--V30 coordinate correction]
\label{corr:n1v31-scope}
The following earlier claims are superseded as physical
left-translated-chart statements:
\begin{enumerate}
\item the V27 completion matrices \({\cal R}_1,{\cal Q}_1\);
\item the endpoint bubble, tadpole, total, and endpoint-minus-Wilson
numbers in V28; and
\item the identification of the V30 exponential-chart cubic payload
with the translated toron response.
\end{enumerate}
The V27 and V30 hashes remain valid certificates of the exponential
chart they actually compute.  The V28 Wilson suffix arrays and Wilson
response remain valid.  The V29 router jets and the V30 independent-leg
Laurent bookkeeping also remain valid.
\end{correction}

\begin{proof}
Lemma~\ref{lem:n1v31-chart-difference} affects precisely nonlinear word
derivatives having coincident noncommuting background and vertical
insertions.  It does not affect the linear incidence, quadratic
routers, their inverses, or a suffix computation already derived in
the product chart.  The exact replay and failure statements below
separate these cases computationally.
\end{proof}

\subsection{The correctly ordered factor compiler}

Partition the labels assigned to a link factor into background labels
\(S_B\) and vertical labels \(S_Z\).  Write
\(\operatorname{Sym}(S)\) for the sum of all ordered colour words
obtained by permuting the labels in \(S\).  The translated factor
coefficient replacing (30.4) is
\[
\begin{split}
 {\cal E}^{\rm tr}_{f,S}
 :={}&
 \frac{\epsilon_f^{|S|}}
 {|S_B|!\,|S_Z|!}
 \prod_{\ell\in S}v_\ell(e_f)
 \prod_{\ell\in S}\zeta_\ell^{r_f}\\
 &\quad\times
 \begin{cases}
 \operatorname{Sym}(S_B)\operatorname{Sym}(S_Z),
     &\epsilon_f=1,\\
 \operatorname{Sym}(S_Z)\operatorname{Sym}(S_B),
     &\epsilon_f=-1.
 \end{cases}                                                \tag{31.4}
\end{split}
\]

\begin{theorem}[Translated word derivative]
\label{thm:n1v31-translated-word}
Replacing (30.4) by (31.4), then applying the word assignment and
Hilbert--Schmidt partition formula (30.6), gives the exact mixed
derivative of every Wilson plaquette and every transported owner in
the chart (31.1).
\end{theorem}

\begin{proof}
On a positive physical edge, all background derivatives act on the
first exponential and all vertical derivatives on the second.  Their
internal repeated derivatives are separately symmetrized, giving the
first order in (31.4).  On an inverse occurrence, (31.2) reverses the
two exponential blocks and supplies one minus sign per derivative,
giving the second order.  Leibniz assignment over word factors and
over \(Y^\dagger Y/2\) is unchanged and is exact by the proof of
Theorem~\ref{thm:n1v30-word-formula}.
\end{proof}

\subsection{An independent Wilson regression}

Let \(C_{\rm W}^{\rm tr}\) be the cubic spatial symbol obtained by
differentiating \(-\operatorname{ReTr}U_p\) with (31.4), dividing by
the lowered colour bracket as in V27, and summing all \(96\)
plaquettes.

\begin{theorem}[The translated compiler replays the Wilson suffix]
\label{thm:n1v31-Wilson-replay}
No normalization correction is required:
\[
 \boxed{\text{raw word normalization}=1.}                   \tag{31.5}
\]
At zero momentum,
\[
 \varepsilon_0(C_{\rm W}^{\rm tr})=R_{\rm W},\qquad
 \operatorname{SHA256}(R_{\rm W})
 =
 \texttt{35AA5495F57CCADEA0DD19D828C81E1291015144F2A21AD4A01BDD2F827AF93C}.
                                                                    \tag{31.6}
\]
The Wilson zero-momentum bubble is consequently
\[
                         {\cal B}_{\rm W}(0)=-\frac{197}{4}, \tag{31.7}
\]
exactly as in V28.
\end{theorem}

\begin{proof}
The suffix array and the word array are generated by independent
algorithms.  The program first tests that their supports agree, then
divides every pair of corresponding nonzero entries.  All ratios equal
one.  It next performs the exact \(45\times45\) trace contraction,
which gives (31.7).
\end{proof}

\subsection{The corrected completion cubic at zero momentum}

Let \(\widetilde{\cal R}_1\) denote the completion cubic generated in
the translated chart.  Exact assembly over the \(12435\) owner-pair
records gives
\[
\begin{array}{c|c|c|l}
&\text{nonzero}&\max|\text{numerator}|&\text{SHA--256}\\ \hline
\widetilde{\cal R}_1&662&14&
\texttt{578910422F9F68F7EF82C12DF6F28BAEB721268B1ECA849A7F54271417BB0BBE}\\
{\cal P}_*^{[0]}=\frac13R_{\rm W}+\widetilde{\cal R}_1&662&43&
\texttt{E8D4E01B578A3B54EB0C0814DA8C7B4C5CE694ABFEB1D03794ADE64299F41FA8}.
\end{array}                                                    \tag{31.8}
\]
Both matrices are exactly skew.

\begin{corollary}[The old completion cubic was in the wrong chart]
\label{cor:n1v31-old-cubic}
\[
 \operatorname{SHA256}(\widetilde{\cal R}_1)
 \ne
 \texttt{9C376E5DC8C6AEC32877E23BDD32B08602D15B94C607B5F262291A93FDDD8593}.
                                                                    \tag{31.9}
\]
In particular, the difference is not a global normalization: the two
matrices have different supports, with \(662\) and \(706\) nonzero
entries.
\end{corollary}

\begin{proof}
Equation (31.9) and the support counts follow from exact sparse
serialization.  The old hash is the V27 exponential-chart hash.
\end{proof}

Using the unchanged endpoint covariance \(G_*(0)\), the corrected
endpoint bubble value is
\[
 \boxed{
 {\cal B}_*(0)
 =3\Tr\!\left(
 G_*(0){\cal P}_*^{[0]}G_*(0){\cal P}_*^{[0]}
 \right)
 =-\frac{74509879}{1178450}.}                               \tag{31.10}
\]
The change from the V28 exponential-chart bubble is
\[
 -\frac{74509879}{1178450}
 +\frac{14704054060106}{347186100625}
 =
 -\frac{14494975333563}{694372201250}.                       \tag{31.11}
\]

\begin{firewall}[No corrected endpoint total yet]
\label{fire:n1v31-no-total}
The endpoint quartic tadpole must be recomputed with (31.4).  Therefore
(31.10) is a corrected bubble, not a corrected total response.  The
V28 endpoint tadpole, endpoint total, and endpoint-minus-Wilson total
must not be used downstream.
\end{firewall}

\subsection{Lift and inverse jets in the corrected calculation}

Write the axial lifted background as
\[
 b_e(q)=b^{[0]}+iq\,\dot b_e+\frac{q^2}{2}b_e^{[2]}+O(q^3).
                                                               \tag{31.12}
\]
The first two fields are common to the endpoints, while the second
derivative distinguishes them:
\[
\begin{array}{c|r|r|r|l}
\text{field}&\text{nonzero}&\max|\text{numerator}|&
\max\text{ denominator}&\text{SHA--256}\\ \hline
b^{[0]}&8&1&1&
\texttt{1B97C1852385FC72F11FFB877B366AF36D6741683C02A3BD9AC9E852AF34BB79}\\
\dot b&12&1&4&
\texttt{089FB353F23741964B9432BA3122D28F0ABA0939FEA3038FC0282C4022940C6B}\\
b_{\rm W}^{[2]}&18&13&24&
\texttt{BDD6433C96BE2AC704D8F826F70635AB7A49E16AB0254112132DDAEABA3F9A9B}\\
b_*^{[2]}&18&20309&288600&
\texttt{592088C2A1610F133FE1F94ABF6B52E0E5C46079A4128532DE2B5D3DD81A6AEE}.
\end{array}                                                    \tag{31.13}
\]
These are the first-polarization columns of the V29 router jets,
embedded in the \(64\) fine-edge types.

For the shifted covariances, write
\[
 G_e(q)=G_e^{[0]}+iq\,\dot G_e+\frac{q^2}{2}G_e^{[2]}+O(q^3).
                                                               \tag{31.14}
\]
The inverse formulas (30.23)--(30.24) give the following canonical
hashes:
\[
\begin{array}{c|lll}
e&G_e^{[0]}&\dot G_e&G_e^{[2]}\\ \hline
{\rm W}&
\texttt{E92F5C77B2B507DF7169038260777A2BE56BC539D1068D2EB508274FA760448B}&
\texttt{6A0D92F9EEFC124A6DF279C8AEA845D6FA320B46E66B81CEC670BE7822BDFFF3}&
\texttt{34F39AF8E3B7171A215037D90452C7DA816B7AE6EB468D1C89A179429566EF9E}\\
*&
\texttt{811E3FEE872D3325202D8E7BC11940869825DF445C3EA1E5167367B4C2BAF65B}&
\texttt{87063A9C0EE0E2F7BA375262740F4357411C7A33899FC5329046911797D51468}&
\texttt{7E10E0FB6773AADB62C0EDB2AB780AB68A6C13890B6C8F068E80B3D1E0F672E9}.
\end{array}                                                    \tag{31.15}
\]

\begin{theorem}[All inverse-jet equations vanish exactly]
\label{thm:n1v31-inverse-residuals}
For \(e={\rm W},*\), substituting (31.14) in
\(A_e(q)G_e(q)=I\) makes the coefficients of orders zero, one, and two
equal respectively to \(I,0,0\) in exact Gaussian-rational arithmetic.
\end{theorem}

\begin{proof}
The program generates the Laurent derivatives of \(A_e\), applies
(30.23)--(30.24), and multiplies the full \(45\times45\) matrices.
Every real and imaginary rational numerator in the two residual
matrices is zero.
\end{proof}

\subsection{Translated cubic jets}

Let the two all-incoming bubble orientations be those of (30.21), but
evaluate their word factors with (31.4).  Their cubic matrices have the
form
\[
 P_e^+(q)=P_e^{[0]}+iq\,\dot P_e+\frac{q^2}{2}P_e^{[2]}+O(q^3),
                                                               \tag{31.16}
\]
\[
 P_e^-(q)=P_e^{[0]}-iq\,\dot P_e+\frac{q^2}{2}P_e^{[2]}+O(q^3).
                                                               \tag{31.17}
\]

\begin{proposition}[Exact phase and lift chain rule]
\label{prop:n1v31-cubic-chain}
For a three-leg Laurent coefficient with axial character integer
\[
 n=-\alpha_{2,1}+\alpha_{3,1},                              \tag{31.18}
\]
the three matrices in (31.16) receive respectively
\[
\begin{array}{c|c}
\text{order}&\text{coefficient}\\ \hline
0&C(b^{[0]})\\
1\text{ after removing }i&
nC(b^{[0]})+C(\dot b)\\
2&
-n^2C(b^{[0]})-2nC(\dot b)+C(b^{[2]}).
\end{array}                                                    \tag{31.19}
\]
The opposite orientation gives the signs in (31.17).
\end{proposition}

\begin{proof}
Apply (30.20) to the phase and differentiate the linear background
slot in (31.12).  The order-two cross term is
\(2(in)(iC(\dot b))=-2nC(\dot b)\).  Replacing \(q\) by \(-q\)
changes both first derivatives and preserves both second derivatives.
\end{proof}

The canonical cubic-jet records are
\[
\begin{array}{c|c|r|l}
e&\text{jet}&\text{nonzero}&\text{SHA--256}\\ \hline
{\rm W}&P^{[0]}&90&
\texttt{35AA5495F57CCADEA0DD19D828C81E1291015144F2A21AD4A01BDD2F827AF93C}\\
{\rm W}&\dot P&170&
\texttt{2B7A861DED30DE1B148B0D47D42B8EDAFDB43A4188C175BFF0EBFF4552CD8AF3}\\
{\rm W}&P^{[2]}&320&
\texttt{0C3C671FE30DFA1E0D1D9B453ABC753F053B3262D5D7F9E336CB17304FFEC410}\\
*&P^{[0]}&662&
\texttt{E8D4E01B578A3B54EB0C0814DA8C7B4C5CE694ABFEB1D03794ADE64299F41FA8}\\
*&\dot P&816&
\texttt{1E75907C72B47B3A435F03529B57E66A8FE361E4BCEA07D3128BF74A3D68F10F}\\
*&P^{[2]}&1239&
\texttt{E2F89E31F5E0E3CBA74AB0CC8CFEE74D7E9E0F88FCFB7FA4CA06BDBDE7DA41F5}.
\end{array}                                                    \tag{31.20}
\]
For both endpoints the program independently constructs the two
orientations and verifies (31.16)--(31.17).  It also verifies that all
three zero-order Wilson, completion, and combined endpoint matrices
are skew.

\subsection{The complete bubble second derivative}

Put \(G_0=G_e^{[0]}\), \(G_1=\dot G_e\), \(G_2=G_e^{[2]}\), and
similarly \(P_0,\ P_1,\ P_2\).  From
\[
 {\cal B}_e(q)
 =
 c_e\Tr\!\left(G_0P_e^+(q)G_e(q)P_e^-(q)\right),
 \qquad c_{\rm W}=\frac13,\quad c_*=3,                       \tag{31.21}
\]
one obtains
\[
\begin{split}
 {\cal B}_e'(0)
 ={}&ic_e\Tr G_0
 \left(P_1G_0P_0+P_0G_1P_0-P_0G_0P_1\right),               \tag{31.22}\\
 {\cal B}_e''(0)
 ={}&c_e\Tr G_0\bigl(
 P_2G_0P_0+P_0G_2P_0+P_0G_0P_2\\
 &\hspace{18mm}
 -2P_1G_1P_0+2P_1G_0P_1+2P_0G_1P_1
 \bigr).                                                     \tag{31.23}
\end{split}
\]

\begin{theorem}[Exact translated-chart bubble jets]
\label{thm:n1v31-bubble-jets}
For the first retained axial polarization,
\[
\boxed{
 {\cal B}_{\rm W}'(0)={\cal B}_*'(0)=0,}                    \tag{31.24}
\]
\[
\boxed{
 {\cal B}_{\rm W}''(0)=-\frac{5147}{1024},}                 \tag{31.25}
\]
\[
\boxed{
 {\cal B}_*''(0)
 =-\frac{740420392225351}{355518567040000}.}                \tag{31.26}
\]
Consequently the endpoint-minus-Wilson bubble second jet is
\[
\boxed{
 {\cal B}_*''(0)-{\cal B}_{\rm W}''(0)
 =
 \frac{261636616922881}{88879641760000}
 \approx2.943718175973.}                                    \tag{31.27}
\]
\end{theorem}

\begin{proof}
Equations (31.22)--(31.23) are the product rule applied to
(31.21), with the factors of \(i\) from (31.14),
(31.16), and (31.17) multiplied explicitly.  The verifier evaluates
all trace words over exact rationals.  Both rational coefficients in
(31.22) reduce to zero.  The order-two contractions reduce to
(31.25)--(31.27).  The order-zero Wilson and endpoint values reproduce
(31.7) and (31.10), providing independent normalizations.
\end{proof}

\begin{firewall}[The bubble jet is not the toron response jet]
\label{fire:n1v31-bubble-only}
The numbers (31.25)--(31.27) omit the translated-chart quartic
tadpole second derivative.  No sign or endpoint comparison may be
assigned to the complete toron response until that term is assembled.
They also concern one polarization at the zero internal momentum and
still carry the eventual finite-grid weight \(M^{-4}\).
\end{firewall}

\subsection{Executable record and next acquisition}

The exact verifier is
\[
       \texttt{scripts/verify\_ym\_v31\_cubic\_bubble\_jets.py}. \tag{31.28}
\]
Its record is
\[
\begin{array}{c|c}
\text{lines}&833\\
\text{bytes}&26834\\
\text{SHA--256}&
\texttt{0F8068E04A5DFA673E8F5E326D7665331EA8C0F6CAE5DFD76ED221BC002D588A}.
\end{array}                                                    \tag{31.29}
\]
The canonical payload containing the four lifted fields, all six
inverse matrices, twelve component cubic matrices, and both bubble
cards has SHA--256
\[
\boxed{
\texttt{E719818F389E7140B0838CDEF15CBB3C111B518910ABEA08E84F2D656AB75506}.}
                                                                    \tag{31.30}
\]

\begin{assessment}[Progress after compact V31]
\label{ass:n1v31-status}
V31 found and repaired a noncommutative coordinate error before it
entered the external response.  The corrected factor compiler exactly
replays the independently derived Wilson suffix, changes the completion
cubic as required, certifies both inverse jets, and evaluates the full
bubble second derivative.  Earlier endpoint response numbers are
quarantined rather than silently reused.
\end{assessment}

\begin{decision}[V32 translated quartic tadpole]
\label{dec:n1v31-next}
Apply (31.4) to the four-leg owner derivative with two internal colour
slots and the \(b(q),b(-q)\) background slots.  Contract the internal
colours first using V24, accumulate directly the zero and second axial
moments without serializing the full \(\Phi_4\) polynomial, and verify
the corrected Wilson quartic suffix.  Only then add the tadpole to
(31.25)--(31.26).
\end{decision}

\begin{firewall}[Claim boundary after compact V31]
\label{fire:n1v31-claim}
V31 corrects the translated cubic and computes the bubble jet.  It does
not yet compute the corrected quartic tadpole, the full toron response
jet, all polarizations, the nonzero-momentum edge mean, SCIC or ADMC,
summable cylinder comparison, sector tightness, a continuum measure,
terminal clustering, or the Yang--Mills mass gap.
\end{firewall}

\subsection{Parent record}

This V31 note was formed by appending the preceding material to the
validated compact V30, whose record was
\[
\begin{array}{c|c}
\text{lines}&12670\\
\text{bytes}&431067\\
\text{SHA--256}&
\texttt{9C12A3729B18AC29FC2737B7FD1FF6162DC1E8E4154762DC0EF0344F3553AA0A}.
\end{array}
\]

% =============================================================================
% END APPENDED COMPACT-SERIES V31 MATERIAL
% =============================================================================

% =============================================================================
% BEGIN APPENDED COMPACT-SERIES V32 MATERIAL
% =============================================================================

\section{Compact-series V32: corrected translated quartic and the
complete zero-toron card}
\label{sec:n1v32}

\subsection{Completion of the coordinate correction at zero momentum}

Correction~\ref{corr:n1v31-scope} also applies to the V30 quartic
regression panel when that panel is interpreted as a physical
left-translated-chart object.  Its hash remains a valid
exponential-chart regression.  The present note replaces the complete
quartic matrix at the toron by the correctly ordered factor rule
(31.4).

\begin{lemma}[Why the direct quartic is complete at the toron]
\label{lem:n1v32-zero-complete}
For the constant commuting lift \(b=b^{[0]}\), the stationary-centre
acceleration and the reduced Haar-density response vanish for both
endpoints.  Hence at internal and external momentum zero,
\[
 Q_e^{\rm meas}(0)
 =
 D^4{\cal V}_e(0)[\,\cdot,\cdot,b,b\,]                      \tag{32.1}
\]
in the translated chart.
\end{lemma}

\begin{proof}
The exact flatness theorem of V26 gives
\(U_e(t)=e^{tb_eX}\) with every plaquette and owner difference equal
to the identity or zero for all \(t\).  Left translation sends this
path to the zero vertical coordinate, so the stationary centre is
identically zero, not merely zero to first order.  Its second
derivative therefore vanishes.  Product Haar measure is exactly left
invariant, so its logarithmic Jacobian is also identically zero.
Substitution in (20.23) and (20.31) leaves (32.1).
\end{proof}

\subsection{Phase-free translated quartic compiler}

At \(q=0\), set every independent character in (31.4) to one before
assembly.  This produces a phase-free colour-word dictionary while
retaining the two ordered blocks
\[
 \operatorname{Sym}(S_B)\operatorname{Sym}(S_Z)
 \quad\hbox{and}\quad
 \operatorname{Sym}(S_Z)\operatorname{Sym}(S_B)              \tag{32.2}
\]
on positive and inverse occurrences respectively.  For reduced
spatial indices \(i,j\), define
\[
 (\widetilde{\cal Q}_1)_{ij}
 :=
 \frac1{g_{11}}\sum_{a,b=1}^8g^{ab}
 D^4{\cal O}_1^{\rm tr}
 [e_iT_a,e_jT_b,bT_1,bT_1].                                \tag{32.3}
\]
The program visits the same \(12435\) locally supported owner-pair
records as V27, but uses (32.2) at every coincident link.

\begin{theorem}[Independent Wilson quartic normalization]
\label{thm:n1v32-Wilson-quartic}
The raw translated word derivative of
\(-\operatorname{ReTr}U_p\), after dual-colour contraction, obeys
\[
                       Q_{\rm W}^{\rm raw}=3Q_{\rm W}.       \tag{32.4}
\]
After restoring the Wilson factor \(1/3\),
\[
 \operatorname{SHA256}(Q_{\rm W})
 =
 \texttt{F80E3577CEAF38B13453762272B9AE63AF11892C1BFD3F73D72D1D0B60EC200F},
                                                                    \tag{32.5}
\]
exactly replaying the independently derived V28 suffix matrix.
\end{theorem}

\begin{proof}
The translated word engine and the suffix formula are generated
independently.  Their supports agree.  For every corresponding nonzero
entry, exact division of the raw word value by the suffix value gives
the same rational number \(3\).  Division by three then reproduces
every suffix entry and the hash (32.5).
\end{proof}

\begin{theorem}[Corrected completion and endpoint quartics]
\label{thm:n1v32-quartics}
The three symmetric quartic matrices have records
\[
\begin{array}{c|r|r|r|l}
Q&\text{nonzero}&\max|\text{numerator}|&
\max\text{ denominator}&\text{SHA--256}\\ \hline
Q_{\rm W}&90&2&1&
\texttt{F80E3577CEAF38B13453762272B9AE63AF11892C1BFD3F73D72D1D0B60EC200F}\\
\widetilde{\cal Q}_1&693&84&1&
\texttt{63142B04120255711D0ECAE9BC6B40793A047EC459C737475E5F3A801B3F8140}\\
Q_*=Q_{\rm W}+\widetilde{\cal Q}_1&693&86&1&
\texttt{943C5F281836D8575475C0FDA29C280A302308E9518E0BAE232E392574DDFCA2}.
\end{array}                                                    \tag{32.6}
\]
In particular,
\[
 \operatorname{SHA256}(\widetilde{\cal Q}_1)
 \ne
 \texttt{CD4DA2C234ADBD85DCDBA0FB76325BD4915FB24EF16B2CCB0BBE89683B26BD49}.
                                                                    \tag{32.7}
\]
\end{theorem}

\begin{proof}
For each owner-pair record the program evaluates all ten nonzero
entries of \(g^{-1}\), divides by \(g_{11}=2\), and accumulates exact
rationals.  Direct transposition checks prove symmetry.  Sparse
row-major serialization gives (32.6).  The old V27
exponential-chart matrix had \(830\) nonzero entries and the hash on
the right of (32.7), so the translated correction is nontrivial.
\end{proof}

\subsection{The corrected complete zero response}

Define the tadpole at the toron by
\[
                         {\cal T}_e(0)
 :=\frac12\Tr(G_e(0)Q_e(0)).                                \tag{32.8}
\]
Exact contraction of (32.6) with the two V31 covariance matrices gives
\[
 \boxed{{\cal T}_{\rm W}(0)=\frac{441}{8},}                 \tag{32.9}
\]
\[
 \boxed{{\cal T}_*(0)=\frac{61348647}{942760}.}             \tag{32.10}
\]
The first value independently replays V28.

\begin{theorem}[Corrected zero-toron endpoint card]
\label{thm:n1v32-zero-card}
Combining (32.9)--(32.10) with the translated bubble values
(31.7) and (31.10) yields
\[
 \boxed{
 {\cal W}_{0,{\rm W}}(0)
 =
 \frac{441}{8}-\frac{197}{4}
 =
 \frac{47}{8},}                                             \tag{32.11}
\]
\[
 \boxed{
 {\cal W}_{0,*}(0)
 =
 \frac{61348647}{942760}
 -\frac{74509879}{1178450}
 =
 \frac{8703719}{4713800},}                                  \tag{32.12}
\]
and therefore
\[
\boxed{
 {\cal W}_{0,*}(0)-{\cal W}_{0,{\rm W}}(0)
 =
 -\frac{2373732}{589225}
 \approx-4.02856633713777.}                                 \tag{32.13}
\]
\end{theorem}

\begin{proof}
Lemma~\ref{lem:n1v32-zero-complete} proves that the translated quartic
and cubic trace terms exhaust the measured stationary response at this
point.  Equations (32.11)--(32.13) are exact additions and subtraction
of the four certified fractions.
\end{proof}

\begin{corollary}[Replacement of the V28 endpoint regression]
\label{cor:n1v32-regression}
In the complete finite-grid sum, the zero-cell endpoint difference for
the first polarization is
\[
                         -\frac1{M^4}
                         \frac{2373732}{589225}.             \tag{32.14}
\]
This replaces the V28 value (28.22), which used the
exponential-chart completion vertices.
\end{corollary}

\begin{proof}
The zero internal-momentum cell has weight \(M^{-4}\), and
Theorem~\ref{thm:n1v32-zero-card} supplies its corrected integrand.
\end{proof}

\begin{firewall}[The corrected zero card is not a mass]
\label{fire:n1v32-not-mass}
Equation (32.13) is a finite one-cell, one-polarization response value.
It is neither an external-momentum second derivative nor a gauge-boson
mass.  Gauge covariance requires the full measured response to have no
zeroth-order mass term after the complete momentum and Ward sum; a
nonzero individual cell is allowed and must be retained.
\end{firewall}

\subsection{Why the external second jet remains open}

At \(q=0\), exact flatness killed \(X_{2,e}\) and the Haar term.
For \(q\ne0\), the harmonic lift is not an exactly flat
one-parameter subgroup.  The full quadratic Hessian vertex is
\[
\begin{split}
 Q_e^{\rm meas}(q)
 ={}&
 D^4{\cal V}_e[\,\cdot,\cdot,b_e(q),b_e(-q)\,]
 +D^3{\cal V}_e[\,\cdot,\cdot,X_{2,e}(q)\,]\\
 &+\frac{2D_v^2h_e(q)}{45}\widehat A_e(q),                  \tag{32.15}
\end{split}
\]
where the first line abbreviates the full expansion of (20.23) along
the harmonic direction.  Although the last two terms vanish at
\(q=0\), their second axial derivatives need not vanish.

\begin{warning}[Direct quartic differentiation is insufficient]
\label{warn:n1v32-stationary}
Adding the direct quartic second jet to (31.25)--(31.26) before
differentiating \(X_{2,e}(q)\) and the Haar density would repeat the
omission already identified in V20.  The next computation must retain
all three rows of (32.15).
\end{warning}

\subsection{Executable record and status}

The exact verifier is
\[
     \texttt{scripts/verify\_ym\_v32\_translated\_quartic\_zero.py}. \tag{32.16}
\]
Its record is
\[
\begin{array}{c|c}
\text{lines}&463\\
\text{bytes}&15555\\
\text{SHA--256}&
\texttt{2397973F33CECC97C066EF9531C6486C0CA7AAEBBFB511884DBB29CE33E0682C}.
\end{array}                                                    \tag{32.17}
\]
The payload containing all three quartic matrices and both complete
zero cards has SHA--256
\[
\boxed{
\texttt{5B0676C79B7AF0BF989A264195B427BBCED5C1C50EBA80DFBCAFBEE5C9B09D08}.}
                                                                    \tag{32.18}
\]

\begin{assessment}[Progress after compact V32]
\label{ass:n1v32-status}
The translated-coordinate correction is now complete at the toron
point.  Both Wilson cubic and quartic word engines replay independent
suffix formulas, the corrected endpoint cubic and quartic arrays are
separately hashed, and the full measured zero response is the exact
fraction (32.12).  The external second derivative remains separated
from this zero value.
\end{assessment}

\begin{decision}[V33 stationary acceleration and Haar jets]
\label{dec:n1v32-next}
Differentiate the stationary equation (20.21) in axial momentum at
\(q=0\), using the translated three- and four-leg word compiler.
Compute the zero and second jets of
\(D^3{\cal V}_e[X_{2,e}(q)]\) and \(D_v^2h_e(q)\), and verify their
Ward recombination with the direct quartic row before forming the full
toron response second derivative.
\end{decision}

\begin{firewall}[Claim boundary after compact V32]
\label{fire:n1v32-claim}
V32 corrects the complete zero-toron card.  It does not yet compute the
full external second jet, all polarizations, the nonzero-momentum edge
mean, SCIC or ADMC, summable cylinder comparison, sector tightness, a
continuum measure, terminal clustering, or the Yang--Mills mass gap.
\end{firewall}

\subsection{Parent record}

This V32 note was formed by appending the preceding material to the
validated compact V31, whose record was
\[
\begin{array}{c|c}
\text{lines}&13145\\
\text{bytes}&448263\\
\text{SHA--256}&
\texttt{55CAEE957B49FF407D9B5B523D0D2B75182E1FBE1BF8277FED154AB27FD0A229}.
\end{array}
\]

% =============================================================================
% END APPENDED COMPACT-SERIES V32 MATERIAL
% =============================================================================

% =============================================================================
% BEGIN APPENDED COMPACT-SERIES V33 MATERIAL
% =============================================================================

\section{Compact-series V33: exact translated quartic jets and the
complete first toron second response}
\label{sec:n1v33}

\subsection{A phase-free exact jet ring}

To differentiate the translated factors without serializing a
four-leg Laurent polynomial, represent every coefficient by
\[
 u(q)=u_0+iq\,u_1+\frac{q^2}{2}u_2+O(q^3),
 \qquad (u_0,u_1,u_2)\in{\mathbb Q}^3.                       \tag{33.1}
\]
Multiplication and conjugation in this jet ring are
\[
\begin{split}
 (uv)_0&=u_0v_0,\\
 (uv)_1&=u_1v_0+u_0v_1,\\
 (uv)_2&=u_2v_0-2u_1v_1+u_0v_2,\\
 \overline{u}&=(u_0,-u_1,u_2).
                                                               \tag{33.2}
\end{split}
\]
The minus sign in the order-two cross term is \(i^2=-1\).

For a direction with axial momentum sign \(\sigma\) inserted at a
factor of offset \(r_f\), its character has jet
\[
                         (1,\sigma(r_f)_1,-(\sigma(r_f)_1)^2). \tag{33.3}
\]
For the two background slots use
\[
\begin{split}
 b_+(q)&=(b^{[0]},\dot b,b^{[2]}),\qquad \sigma=1,\\
 b_-(q)&=(b^{[0]},-\dot b,b^{[2]}),\qquad \sigma=-1,          \tag{33.4}
\end{split}
\]
in the convention (31.12).  The factor ordering remains (31.4).

\begin{theorem}[Jet-ring equivalence]
\label{thm:n1v33-jet-ring}
Applying (33.2)--(33.4) at every translated link factor gives exactly
the zero, first, and second axial derivatives obtained by first forming
the independent-leg Laurent expression and then applying
(30.19)--(30.20) together with the router chain rule.
\end{theorem}

\begin{proof}
For one direction, multiply its field jet by (33.3).  Formula (33.2)
is the second-order Leibniz rule after removing the common \(i\) from
odd derivatives.  Repeated application therefore gives the derivative
of every ordered factor product.  Addition over assignments, owner
differences, Hilbert--Schmidt partitions, and colour contractions is
linear and commutes with taking jets.
\end{proof}

\subsection{The stationary acceleration vanishes in the same-colour row}

For an arbitrary vertical direction \(w\), define the mixed stationary
score
\[
 {\cal S}_e(q)[w]
 :=
 D^3{\cal V}_e(0)
 [w,b_e(q)T_1,b_e(-q)T_1].                                  \tag{33.5}
\]
The residual centre acceleration is
\[
                         X_{2,e}(q)=-G_e(0){\cal S}_e(q).     \tag{33.6}
\]

\begin{theorem}[All stationary-score jets through order two vanish]
\label{thm:n1v33-score-zero}
For the Wilson plaquette action, for its evaluation on the endpoint
lift, and for the completion action separately,
\[
\boxed{
 {\cal S}_e^{[0]}={\cal S}_e^{[1]}={\cal S}_e^{[2]}=0}       \tag{33.7}
\]
as covectors on all \(45\) reduced spatial directions and all eight
colour directions.  Consequently
\[
\boxed{
 X_{2,e}^{[0]}=X_{2,e}^{[1]}=X_{2,e}^{[2]}=0.}              \tag{33.8}
\]
\end{theorem}

\begin{proof}
The real cubic colour tensor is the lowered bracket (24.4), and the
two background colours are equal, so its invariant colour factor
contains \([T_1,T_1]=0\).  As an exhaustive check independent of this
reduction, the verifier applies the ordered translated word engine to
every reduced edge, every colour basis vector, and every plaquette or
owner.  All three rational jet components vanish separately for the
Wilson, endpoint-Wilson, and completion scores.  Equation (33.8)
follows from (33.6).
\end{proof}

\subsection{Haar invariance is exact in the group chart}

\begin{theorem}[No Haar response at any external momentum]
\label{thm:n1v33-Haar}
Let \(g_e(v,q)\in SU(3)\) be any background-dependent recentering of
the reduced link \(z_e\).  The map
\[
                         z_e\longmapsto g_e(v,q)z_e           \tag{33.9}
\]
preserves the reduced product-Haar measure exactly.  Hence its
logarithmic Jacobian and all background and momentum derivatives
vanish.
\end{theorem}

\begin{proof}
Normalized Haar measure on each \(SU(3)\) factor is left invariant.
Taking the product over the \(45\) reduced link factors proves (33.9)
without a small-field expansion.  Dependence of \(g_e\) on \(v,q\)
does not alter the identity for each fixed parameter value, so every
parameter derivative is zero.
\end{proof}

\begin{correction}[Refinement of the V32 warning]
\label{corr:n1v33-Haar-score}
Warning~\ref{warn:n1v32-stationary} correctly required the missing rows
to be checked.  The present calculation closes them for this
same-colour toron card: the stationary row vanishes through the
required axial order by (33.7), and the Haar row vanishes identically
by Theorem~\ref{thm:n1v33-Haar}.  The nonzero direct quartic jet may
therefore be added to the V31 bubble jet.
\end{correction}

\subsection{Exact quartic matrices through second axial order}

Write
\[
 Q_e(q)=Q_e^{[0]}+iq\,\dot Q_e+\frac{q^2}{2}Q_e^{[2]}+O(q^3).
                                                               \tag{33.10}
\]
Both first-jet matrices vanish entry by entry.  Their common zero hash
is
\[
 \operatorname{SHA256}(0_{45\times45})
 =
 \texttt{D3CD5ACF94F846A7E5FA993AFE37299B2F1C1C94BBBAB29BB3063C56633E3622}.
                                                                    \tag{33.11}
\]
The nontrivial records are
\[
\begin{array}{c|c|r|r|r|l}
e&\text{jet}&\text{nonzero}&\max|\text{numerator}|&
\max\text{ denominator}&\text{SHA--256}\\ \hline
{\rm W}&Q^{[0]}&90&2&1&
\texttt{F80E3577CEAF38B13453762272B9AE63AF11892C1BFD3F73D72D1D0B60EC200F}\\
{\rm W}&Q^{[2]}&235&53&36&
\texttt{0D1BB2DD54FF597AEBED6BABDC070142F640429A321502FC0E6D0F21AC7349F8}\\
{\cal O}_1&Q^{[0]}&693&84&1&
\texttt{63142B04120255711D0ECAE9BC6B40793A047EC459C737475E5F3A801B3F8140}\\
{\cal O}_1&Q^{[2]}&1399&21078221&72150&
\texttt{98DA0BE0DAA80AFB87F3E172628051BF0382ED851CEDF8253AEEBF3B0EE7BED1}\\
*&Q^{[0]}&693&86&1&
\texttt{943C5F281836D8575475C0FDA29C280A302308E9518E0BAE232E392574DDFCA2}\\
*&Q^{[2]}&1399&63258713&432900&
\texttt{1FDF53E67FDB1AA5BE79745E3A4DD5AD2B82C99C979F286916538F9B3775E5B6}.
\end{array}                                                    \tag{33.12}
\]

\begin{proposition}[Quartic regression and reflection checks]
\label{prop:n1v33-quartic-checks}
The order-zero matrices in (33.12) reproduce all three V32 hashes.
The raw Wilson word normalization is \(3\) for both the Wilson and
endpoint lifts.  Both order-one matrices vanish before the trace is
taken.
\end{proposition}

\begin{proof}
The verifier compares the order-zero sparse matrices entry by entry
with V32, checks the two exact normalization ratios, and tests every
entry of the order-one matrices for zero.
\end{proof}

\subsection{Tadpole and complete response jets}

The tadpole jet is
\[
 {\cal T}_e^{[r]}
 =\frac12\Tr\!\left(G_e(0)Q_e^{[r]}\right),
 \qquad r=0,1,2.                                             \tag{33.13}
\]
Exact trace contraction of (33.12) gives
\[
\boxed{
 ({\cal T}_{\rm W}^{[0]},{\cal T}_{\rm W}^{[1]},
 {\cal T}_{\rm W}^{[2]})
 =
 \left(\frac{441}{8},0,-\frac{2809}{128}\right),}           \tag{33.14}
\]
\[
\boxed{
 ({\cal T}_*^{[0]},{\cal T}_*^{[1]},{\cal T}_*^{[2]})
 =
 \left(
 \frac{61348647}{942760},
 0,
 -\frac{136877990034143}{906935120000}
 \right).}                                                   \tag{33.15}
\]

\begin{theorem}[Complete first-polarization toron response jet]
\label{thm:n1v33-complete-response}
After adding the bubble jets of
Theorem~\ref{thm:n1v31-bubble-jets}, the complete measured stationary
toron cards are
\[
\boxed{
 ({\cal W}_{0,{\rm W}}^{[0]},
  {\cal W}_{0,{\rm W}}^{[1]},
  {\cal W}_{0,{\rm W}}^{[2]})
 =
 \left(
 \frac{47}{8},0,-\frac{27619}{1024}
 \right),}                                                   \tag{33.16}
\]
\[
\boxed{
 ({\cal W}_{0,*}^{[0]},
  {\cal W}_{0,*}^{[1]},
  {\cal W}_{0,*}^{[2]})
 =
 \left(
 \frac{8703719}{4713800},
 0,
 -\frac{54396592485609407}{355518567040000}
 \right).}                                                   \tag{33.17}
\]
The endpoint-minus-Wilson card is therefore
\[
\boxed{
\left(
-\frac{2373732}{589225},
0,
-\frac{11201914893111883}{88879641760000}
\right),}                                                    \tag{33.18}
\]
whose second component is approximately
\[
                         -126.034653957767.                  \tag{33.19}
\]
\end{theorem}

\begin{proof}
Theorems~\ref{thm:n1v33-score-zero} and
\ref{thm:n1v33-Haar} prove that (33.13) contains every tadpole-side
stationary and measure contribution through axial order two.  Add
(33.14) to the Wilson bubble triple
\[
 \left(-\frac{197}{4},0,-\frac{5147}{1024}\right)
\]
and (33.15) to the endpoint bubble triple
\[
 \left(
 -\frac{74509879}{1178450},
 0,
 -\frac{740420392225351}{355518567040000}
 \right).
\]
Reduction of the resulting rational fractions gives
(33.16)--(33.18).  Their order-zero entries replay V32.
\end{proof}

\begin{corollary}[Finite-grid toron second-jet regression]
\label{cor:n1v33-grid-regression}
For the first retained axial polarization, the endpoint-minus-Wilson
zero-cell contribution to the external second derivative of the
finite-grid response is
\[
 -\frac1{M^4}
 \frac{11201914893111883}{88879641760000}.                   \tag{33.20}
\]
\end{corollary}

\begin{proof}
The uniform internal-momentum grid assigns weight \(M^{-4}\) to
the toron.  Equation (33.18) supplies its exact second derivative.
\end{proof}

\begin{firewall}[A complete toron card is not the edge mean]
\label{fire:n1v33-toron-only}
Theorem~\ref{thm:n1v33-complete-response} is complete only for one
external polarization and the single internal momentum \(k=0\).
The number (33.20) is not the Brillouin-zone mean, a continuum beta
coefficient, a gauge-boson mass, a spectral gap, or a clustering rate.
It is a hard regression value that the full interval compiler must
contain with its exact weight.
\end{firewall}

\subsection{Executable certificate}

The verifier is
\[
 \texttt{scripts/verify\_ym\_v33\_translated\_quartic\_jets.py}. \tag{33.21}
\]
Its record is
\[
\begin{array}{c|c}
\text{lines}&836\\
\text{bytes}&26328\\
\text{SHA--256}&
\texttt{F3B696072C832D4B98652544A04754F5592AB620A4DE725BF134ACC14804B702}.
\end{array}                                                    \tag{33.22}
\]
The payload containing all nine quartic-jet matrices, both tadpole
triples, and both complete response triples has SHA--256
\[
\boxed{
\texttt{D85B44787BABB3CFDFD2C6583224E3A243E362B1E70B83B3780126B2D0B255C2}.}
                                                                    \tag{33.23}
\]
Fourteen checks pass: the same-colour commutator, three independent
stationary-score rows, two Wilson normalizations, three V32 zero-order
replays, two zero first-jet matrices, two zero-card replays, and the
vanishing of both total first derivatives.

\begin{assessment}[Progress after compact V33]
\label{ass:n1v33-status}
The first translated toron response is now exact through external
second order.  The stationary acceleration and Haar rows have been
closed rather than bounded, the direct quartic and bubble pieces are
separately serialized, and reflection is verified entrywise before
the trace.  This completes one finite regression card while preserving
the distinction between that card and the global coefficient.
\end{assessment}

\begin{decision}[V34 four-polarization toron orbit]
\label{dec:n1v33-next}
Repeat the V29 lift, V31 cubic, and V33 quartic jet construction for
the other three retained polarizations.  Determine the exact restoring
tree-coordinate transformations under axis permutations, or else
retain four independent hashes.  Only after all four cards pass may
the toron block enter the hypercubic polarization trace.
\end{decision}

\begin{firewall}[Claim boundary after compact V33]
\label{fire:n1v33-claim}
V33 completes one toron response second jet.  It does not yet verify
the other three polarizations, evaluate any nonzero internal momentum,
enclose the edge mean, complete SCIC or ADMC, prove summable cylinder
comparison, establish sector tightness, construct the continuum
measure or terminal clustering, or prove the Yang--Mills mass gap.
\end{firewall}

\subsection{Parent record}

This V33 note was formed by appending the preceding material to the
validated compact V32, whose record was
\[
\begin{array}{c|c}
\text{lines}&13440\\
\text{bytes}&458760\\
\text{SHA--256}&
\texttt{C2794F960F65463F5CFECE3D25EA6039612A8E2CEBD8AA3B2B6BE5870269B0F0}.
\end{array}
\]

% =============================================================================
% END APPENDED COMPACT-SERIES V33 MATERIAL
% =============================================================================

% =============================================================================
% BEGIN APPENDED COMPACT-SERIES V34 MATERIAL
% =============================================================================

\section{Compact-series V34: the transverse toron card and the complete
hypercubic polarization orbit}
\label{sec:n1v34}

\subsection{Longitudinal and transverse orbits}

Keep external axial momentum \(q e_0\).  The V33 background used the
retained polarization \(e_0\), which is longitudinal relative to this
axis.  The present calculation uses retained polarization \(e_1\).
The remaining \(e_2,e_3\) polarizations lie in the same orbit under
axis permutations fixing \(e_0\).

\begin{theorem}[Transverse quotient-measure symmetry]
\label{thm:n1v34-transverse-orbit}
For either endpoint and every finite parity-cell response coefficient,
\[
 {\cal W}_{\mu\mu}(q e_0)
 =
 {\cal W}_{\nu\nu}(q e_0),
 \qquad \mu,\nu\in\{1,2,3\}.                                \tag{34.1}
\]
Every off-diagonal component vanishes:
\[
                         {\cal W}_{\mu\nu}(q e_0)=0
                         \quad(\mu\ne\nu).                   \tag{34.2}
\]
\end{theorem}

\begin{proof}
Any permutation of axes \(1,2,3\) fixes \(q e_0\) and bijects the fine
edge, plaquette, and transported-owner sets.  It preserves orientation
up to loop reversal; \(\operatorname{ReTr}U_p\) and the squared
Hilbert--Schmidt owner norm are invariant under that reversal.  It also
preserves product Haar measure and the gauge-group action.  Therefore
it induces a measure-preserving isomorphism of the finite gauge
quotient carrying polarization \(\mu\) to polarization \(\nu\).
The quotient Hessians and their differentiated trace words are
conjugate, proving (34.1).

The rooted tree slice need not be fixed by this permutation, but tree
gauge has background-independent unit Faddeev--Popov determinant.
Changing the tree is only a coordinate change on the same quotient and
cannot change the scalar trace response.

For (34.2), reflect one transverse axis occurring in exactly one of the
two polarization labels.  The external momentum is unchanged, while
that background component changes sign; the mixed coefficient is
therefore its own negative.  The same argument handles a
longitudinal--transverse pair and two distinct transverse labels.
\end{proof}

\begin{corollary}[Only two diagonal cards are independent]
\label{cor:n1v34-two-cards}
The complete toron polarization tensor at axial momentum has the form
\[
 \operatorname{diag}({\cal W}_{L},{\cal W}_{T},
                     {\cal W}_{T},{\cal W}_{T}).             \tag{34.3}
\]
Thus V33 and one independently computed transverse card determine the
whole four-polarization toron block.
\end{corollary}

\subsection{Exact transverse cubic and quartic records}

The V34 verifier repeats the full translated construction with
zero-based retained polarization index \(1\).  It does not obtain the
transverse scalar by copying the longitudinal answer.  The combined
endpoint cubic jets have records
\[
\begin{array}{c|r|r|r|l}
\text{jet}&\text{nonzero}&\max|\text{numerator}|&
\max\text{ denominator}&\text{SHA--256}\\ \hline
P_*^{[0]}&664&43&3&
\texttt{D16EE39AA679356E5A99FF34A14870B57301F54F1D30C2FEA35C8EE7478885E2}\\
\dot P_*&818&79&24&
\texttt{3C09369E8D700F49EF3E75EA52D20EBA2D414706490DD7020BBCC7554CA6A91D}\\
P_*^{[2]}&1317&25546709&1731600&
\texttt{F63CA37AC2ED892ADEC0490C7E28FACF3D85599E77A602A942F46D7D6A816CBB}.
\end{array}                                                    \tag{34.4}
\]
The combined endpoint quartic jets are
\[
\begin{array}{c|r|r|r|l}
\text{jet}&\text{nonzero}&\max|\text{numerator}|&
\max\text{ denominator}&\text{SHA--256}\\ \hline
Q_*^{[0]}&696&86&1&
\texttt{D29EE00B1407FFD84A8FED0F878B9CACCC0B3CA76BF3AA1E36263B6C744F8069}\\
\dot Q_*&0&0&1&
\texttt{D3CD5ACF94F846A7E5FA993AFE37299B2F1C1C94BBBAB29BB3063C56633E3622}\\
Q_*^{[2]}&1407&63258713&432900&
\texttt{5CC967BA406A8ACFF46E65FAD476F729C978F6D5E7D8FF4CDA7546AF2708D768}.
\end{array}                                                    \tag{34.5}
\]
The different hashes from V33 are expected: the fixed rooted tree
coordinate matrices are not literally axis-permutation matrices.
Only their quotient trace scalars have the symmetry of
Theorem~\ref{thm:n1v34-transverse-orbit}.

\begin{proposition}[All transverse coordinate gates pass]
\label{prop:n1v34-gates}
For the transverse card:
\begin{enumerate}
\item the Wilson cubic word normalization is one;
\item both Wilson quartic word normalizations are three;
\item both cubic orientations have the required first/second parity;
\item the Wilson, endpoint-Wilson, and completion stationary-score
jets vanish through order two;
\item both quartic first-jet matrices vanish; and
\item both bubble and total first derivatives vanish.
\end{enumerate}
\end{proposition}

\begin{proof}
Each item is an exact rational comparison performed before the response
fractions are emitted.  No floating-point tolerance occurs.
\end{proof}

\subsection{The exact transverse response}

The transverse tadpole and bubble triples are
\[
\boxed{
 ({\cal T}_{\rm W,T}^{[0]},{\cal T}_{\rm W,T}^{[1]},
  {\cal T}_{\rm W,T}^{[2]})
 =
 \left(\frac{487}{8},0,-\frac{17009}{384}\right),}           \tag{34.6}
\]
\[
\boxed{
 ({\cal B}_{\rm W,T}^{[0]},{\cal B}_{\rm W,T}^{[1]},
  {\cal B}_{\rm W,T}^{[2]})
 =
 \left(-55,0,\frac{14713}{1024}\right),}                    \tag{34.7}
\]
\[
\boxed{
 ({\cal T}_{*,T}^{[0]},{\cal T}_{*,T}^{[1]},
  {\cal T}_{*,T}^{[2]})
 =
 \left(
 \frac{5132669}{72520},
 0,
 -\frac{60057901669949}{388686480000}
 \right),}                                                   \tag{34.8}
\]
\[
\boxed{
 ({\cal B}_{*,T}^{[0]},{\cal B}_{*,T}^{[1]},
  {\cal B}_{*,T}^{[2]})
 =
 \left(
 -\frac{162459883}{2356900},
 0,
 \frac{1809930195288949}{355518567040000}
 \right).}                                                   \tag{34.9}
\]

\begin{theorem}[Complete transverse toron response jet]
\label{thm:n1v34-transverse-card}
The complete measured stationary transverse cards are
\[
\boxed{
 ({\cal W}_{0,{\rm W},T}^{[0]},
  {\cal W}_{0,{\rm W},T}^{[1]},
  {\cal W}_{0,{\rm W},T}^{[2]})
 =
 \left(
 \frac{47}{8},0,-\frac{91933}{3072}
 \right),}                                                   \tag{34.10}
\]
\[
\boxed{
 ({\cal W}_{0,*,T}^{[0]},
  {\cal W}_{0,*,T}^{[1]},
  {\cal W}_{0,*,T}^{[2]})
 =
 \left(
 \frac{8703719}{4713800},
 0,
 -\frac{159369091596473209}{1066555701120000}
 \right).}                                                   \tag{34.11}
\]
Their endpoint difference is
\[
\boxed{
 \left(
 -\frac{2373732}{589225},
 0,
 -\frac{10620935983976257}{88879641760000}
 \right),}                                                   \tag{34.12}
\]
whose second derivative is approximately
\[
                         -119.497961216538.                  \tag{34.13}
\]
\end{theorem}

\begin{proof}
Add (34.6) to (34.7) and (34.8) to (34.9).  The stationary row
vanishes by the exhaustive score check, and the Haar row vanishes by
Theorem~\ref{thm:n1v33-Haar}.  Exact fraction reduction gives
(34.10)--(34.12).
\end{proof}

\subsection{The complete diagonal toron tensor}

Let \(\Delta_L^{[2]}\) be the endpoint difference in (33.18) and
\(\Delta_T^{[2]}\) that in (34.12).  Theorem
\ref{thm:n1v34-transverse-orbit} gives
\[
 \Delta_{\mu\nu}^{[2]}(0)
 =
 \operatorname{diag}\left(
 -\frac{11201914893111883}{88879641760000},
 -\frac{10620935983976257}{88879641760000},
 -\frac{10620935983976257}{88879641760000},
 -\frac{10620935983976257}{88879641760000}
 \right).                                                    \tag{34.14}
\]
Its polarization trace is
\[
\boxed{
 \Tr_{\rm pol}\Delta^{[2]}(0)
 =
 -\frac{581955714122171}{1201076240000}
 \approx-484.5285376073804.}                                \tag{34.15}
\]
Separately,
\[
 \Tr_{\rm pol}{\cal W}_{0,{\rm W}}^{[2]}=-\frac{467}{4},
 \qquad
 \Tr_{\rm pol}{\cal W}_{0,*}^{[2]}
 =-\frac{722181365142171}{1201076240000}.                   \tag{34.16}
\]

\begin{proof}
Equation (34.14) combines the longitudinal V33 entry with the
transverse orbit (34.1), while (34.2) removes all off-diagonal entries.
Adding one longitudinal and three transverse fractions gives
(34.15)--(34.16).
\end{proof}

\begin{corollary}[Full toron-block finite-grid regression]
\label{cor:n1v34-grid}
The endpoint-minus-Wilson polarization-traced toron contribution to the
finite-grid external second derivative is
\[
 -\frac1{M^4}
 \frac{581955714122171}{1201076240000}.                      \tag{34.17}
\]
\end{corollary}

\begin{proof}
Apply the toron grid weight \(M^{-4}\) to (34.15).
\end{proof}

\begin{firewall}[Orbit closure is still a one-cell result]
\label{fire:n1v34-one-cell}
V34 closes all four external polarization components only at internal
momentum zero.  It does not supply the nonzero-momentum integrand, its
interval enclosure, the Ward cancellation between internal cells, or
the Brillouin-zone mean.  In particular, the nonzero longitudinal
toron entry does not violate the full Ward identity; that identity
applies after the complete internal sum is assembled.
\end{firewall}

\subsection{Executable record and next acquisition}

The verifier is
\[
              \texttt{scripts/verify\_ym\_v34\_transverse\_toron.py}. \tag{34.18}
\]
Its record is
\[
\begin{array}{c|c}
\text{lines}&414\\
\text{bytes}&13331\\
\text{SHA--256}&
\texttt{A1A2B86A69AED5F3F9AC7EFFBF6B85EBEFEC7726AF50FCEB7EFC75767831D788}.
\end{array}                                                    \tag{34.19}
\]
The canonical transverse cubic, quartic, and response payload has
SHA--256
\[
\boxed{
\texttt{570DCE62A3067BBB8965B580D9CDC47EDCBDBB2AE37CB61975C910742977A386}.}
                                                                    \tag{34.20}
\]

\begin{assessment}[Progress after compact V34]
\label{ass:n1v34-status}
The toron block is now closed through external second order for the
longitudinal and transverse quotient orbits.  One independent
transverse calculation passes the same normalization, parity,
stationary, and Haar gates as V33; quotient symmetry supplies the other
two transverse axes; and the full polarization trace is the exact
fraction (34.15).
\end{assessment}

\begin{decision}[Mandatory V35 audit and first nonzero momentum cell]
\label{dec:n1v34-next}
Perform the required SM/NS companion audit.  Then move to the first
nonzero internal momentum orbit.  Retain independent internal and
external phase labels, certify both shifted covariance inverses, and
replay the toron limit as the orbit approaches zero.  Do not begin the
global interval mean until one nonzero orbit passes the complete
cubic, quartic, stationary-score, Haar, and Ward gates.
\end{decision}

\begin{firewall}[Claim boundary after compact V34]
\label{fire:n1v34-claim}
V34 completes the four-polarization toron block.  It does not evaluate
nonzero internal momenta, enclose the edge mean, complete SCIC or ADMC,
prove summable cylinder comparison, establish sector tightness,
construct the continuum measure or terminal clustering, or prove the
Yang--Mills mass gap.
\end{firewall}

\subsection{Parent record}

This V34 note was formed by appending the preceding material to the
validated compact V33, whose record was
\[
\begin{array}{c|c}
\text{lines}&13809\\
\text{bytes}&471453\\
\text{SHA--256}&
\texttt{ACC5EA1D076C855363B1317C668144914A65C5E245EC5B9D84B279D9BBBC344C}.
\end{array}
\]

% =============================================================================
% END APPENDED COMPACT-SERIES V34 MATERIAL
% =============================================================================

% =============================================================================
% BEGIN APPENDED COMPACT-SERIES V35 MATERIAL
% =============================================================================

\section{Compact-series V35: companion audit and an explicit
weak-coupling good-block theorem}
\label{sec:n1v35}

\subsection{Why the programme moves upstream here}

V31--V34 corrected the noncommutative background chart and completed
the zero-internal-momentum response through second external order.
The next algebraic acquisition remains the first nonzero internal
momentum orbit.  The compulsory five-version audit, however, exposes a
logically earlier missing input: local exact response calculations can
enter a continuum argument only on blocks on which their coordinate,
gap, and inverse estimates hold with summably high probability.

This note therefore proves that probability row directly for both the
Wilson measure and the positive comparison endpoint already fixed in
V20--V21.  It uses only positivity, normalized Haar measure, and a
small-link comparison event.  It does not assume a cluster expansion,
a thermodynamic pressure limit, or independence of plaquettes.

\subsection{The V35 compulsory companion audit}

The current companion files at the time of this audit are
\[
\begin{array}{c|r|r|l}
\text{programme and file}&\text{lines}&\text{bytes}&
\text{SHA--256}\\ \hline
\text{SM--Einstein V43}
&13619&401300&
\texttt{9D8FBDD37858BEEC785F16609CE9726175F8A0D1B07517B6D9EDF63EA164DB15}\\
\text{Navier--Stokes V26}
&5319&278283&
\texttt{82C887F8C2C69E4F28FAD7C3DC8DE5B9099CB5A4BF431EBA1FFF3E91935978AD}.
\end{array}                                                    \tag{35.1}
\]
The relevant SM conclusions are threefold.
\begin{enumerate}
\item V41 obtains a plaquette mean and chord tail from a global
      partition-function comparison, avoiding the volume loss in a
      conditional crossing argument.
\item V42 shows how a summable gauge failure intensity must be combined
      with every deterministic adjacent-error stock; neither row alone
      constructs the limiting cylinder functional.
\item V43 proves a scalar-shift small-singular-value estimate but also
      proves that its crude inverse bound cannot be used to remove a
      monotone ridge.  Thus positivity of the present endpoint action
      must not be confused with global convexity of its stationary
      Hessian away from the good chart.
\end{enumerate}
The NS V26 note sharpens rank-one tensor norms in a heat-kernel
stability problem.  It supplies no gauge-measure or noncommutative
endpoint identity that changes the present acquisition.

\subsection{Finite periodic endpoint measures}

Let
\[
 \Lambda_L=({\mathbb Z}/L{\mathbb Z})^4,\qquad
 |{\cal E}_L|=4L^4,\qquad |{\cal P}_L|=6L^4,                 \tag{35.2}
\]
where \({\cal P}_L\) contains one positively oriented plaquette for
each pair \(\mu<\nu\).  Put \(G=SU(3)\), give every link normalized
Haar measure, and write
\[
 z_p(U):=1-\frac13\operatorname{ReTr}U_p
 =\frac16\|I-U_p\|_{\rm HS}^2,\qquad
 V_{\rm W}:=\sum_{p\in{\cal P}_L}z_p.                       \tag{35.3}
\]
For \(p=(x,\mu\nu)\) and \(1\leq\rho\leq4\), define
\[
 Q_{p,\rho}:=
 U_{x,\rho}U_{p+\widehat e_\rho}U_{x,\rho}^{-1},
 \qquad
 {\cal O}_1:=
 \frac12\sum_{p,\rho}\|Q_{p,\rho}-U_p\|_{\rm HS}^2.         \tag{35.4}
\]
This is precisely the V20 endpoint owner with the normalization
\(\nu_1=1\) reconciled in V21.  For \(0\leq t\leq1\), set
\[
 V_t:=V_{\rm W}+t{\cal O}_1,\qquad
 d\nu_{\beta,L,t}:=
 Z_{\beta,L,t}^{-1}e^{-\beta V_t}\prod_{e\in{\cal E}_L}dH(U_e).
                                                                    \tag{35.5}
\]
Thus \(t=0\) is the Wilson ensemble and \(t=1\) is the positive
comparison endpoint.

For \(0<\delta\leq1\), let
\[
 B_G(\delta):=\{U\in SU(3):\|U-I\|_{\rm op}\leq\delta\},
 \qquad v_G(\delta):=H(B_G(\delta)),                         \tag{35.6}
\]
and let \({\cal A}_\delta\) be the event that every link belongs to
\(B_G(\delta)\).

\begin{lemma}[Exact small-link comparison constants]
\label{lem:n1v35-comparison}
On \({\cal A}_\delta\),
\[
 V_{\rm W}\leq8\delta^2|{\cal P}_L|,\qquad
 {\cal O}_1\leq384\delta^2|{\cal P}_L|,                     \tag{35.7}
\]
and consequently
\[
 \boxed{
 V_t\leq C_t\delta^2|{\cal P}_L|,
 \qquad C_t:=8+384t\leq392.}                                \tag{35.8}
\]
\end{lemma}

\begin{proof}
For unitary matrices \(A_1,\ldots,A_m\), telescoping gives
\[
 \|A_1\cdots A_m-I\|_{\rm op}
 \leq\sum_{j=1}^m\|A_j-I\|_{\rm op}.                        \tag{35.9}
\]
The same distance is unchanged by inversion.  Every plaquette word
has four link factors, so
\[
 \|U_p-I\|_{\rm op}\leq4\delta.
\]
Since \(\|M\|_{\rm HS}^2\leq3\|M\|_{\rm op}^2\), (35.3) gives
\[
 z_p\leq\frac16\,3(4\delta)^2=8\delta^2.                   \tag{35.10}
\]
Moreover \(Q_{p,\rho}\) is conjugate to
\(U_{p+\widehat e_\rho}\), whence
\[
 \|Q_{p,\rho}-I\|_{\rm op}\leq4\delta,\qquad
 \|Q_{p,\rho}-U_p\|_{\rm op}\leq8\delta.                   \tag{35.11}
\]
Each owner in (35.4) is therefore at most
\[
 \frac12\,3(8\delta)^2=96\delta^2.                         \tag{35.12}
\]
There are four owners per plaquette.  Summation proves the second
inequality in (35.7), and (35.8) follows.
\end{proof}

\begin{theorem}[Finite-volume plaquette large-deviation bound]
\label{thm:n1v35-tail}
Let
\[
 \overline z_L:=\frac1{|{\cal P}_L|}
 \sum_{p\in{\cal P}_L}z_p,\qquad
 r_L:=\frac{|{\cal E}_L|}{|{\cal P}_L|}=\frac23.            \tag{35.13}
\]
For every \(\beta>0\), \(0\leq t\leq1\),
\(0<\delta\leq1\), and \(\tau\geq0\),
\[
 \boxed{
 \nu_{\beta,L,t}\{\overline z_L\geq\tau\}
 \leq
 \min\left\{
 1,\,
 \exp\left[
 -\beta|{\cal P}_L|
 \left(\tau-a_{\beta,t}(\delta)\right)
 \right]
 \right\},}                                                 \tag{35.14}
\]
where
\[
 a_{\beta,t}(\delta):=
 C_t\delta^2+
 \frac{2}{3\beta}\log\frac1{v_G(\delta)}.                  \tag{35.15}
\]
No thermodynamic-volume constant occurs in \(a_{\beta,t}\).
\end{theorem}

\begin{proof}
Product Haar measure gives
\[
 H^{{\cal E}_L}({\cal A}_\delta)
 =v_G(\delta)^{|{\cal E}_L|}.
\]
Lemma~\ref{lem:n1v35-comparison} therefore yields
\[
 Z_{\beta,L,t}
 \geq
 v_G(\delta)^{|{\cal E}_L|}
 \exp\{-\beta C_t\delta^2|{\cal P}_L|\}.                   \tag{35.16}
\]
Both terms in (35.5) are nonnegative, so
\[
 V_t\geq V_{\rm W}
 =|{\cal P}_L|\overline z_L.                               \tag{35.17}
\]
The unnormalized mass of
\(\{\overline z_L\geq\tau\}\) is consequently at most
\(\exp\{-\beta\tau|{\cal P}_L|\}\), because normalized product
Haar measure has total mass one.  Divide by (35.16), use (35.13),
and also use the trivial probability bound one.
\end{proof}

\begin{corollary}[Uniform local mean]
\label{cor:n1v35-mean}
Hypercubic and translation invariance imply, for every plaquette \(p\),
\[
 {\mathbb E}_{\beta,L,t}z_p
 =
 {\mathbb E}_{\beta,L,t}\overline z_L.
\]
If \(a_{\beta,t}(\delta)\leq2\), then
\[
 \boxed{
 {\mathbb E}_{\beta,L,t}z_p
 \leq
 a_{\beta,t}(\delta)
 +\frac1{\beta|{\cal P}_L|}.}                               \tag{35.18}
\]
Every translation-invariant weak subsequential infinite-volume limit
satisfies
\[
 \boxed{
 {\mathbb E}_{\beta,\infty,t}z_p
 \leq a_{\beta,t}(\delta).}                                 \tag{35.19}
\]
\end{corollary}

\begin{proof}
The action and product Haar measure have all periodic translations and
hypercubic rotations as symmetries, so all expectations in the first
display are equal.  Since \(0\leq z_p\leq2\),
\[
 {\mathbb E}\overline z_L
 =
 \int_0^2\nu_{\beta,L,t}\{\overline z_L\geq s\}\,ds.
\]
Use the bound one up to \(a_{\beta,t}(\delta)\) and integrate the
exponential in (35.14) thereafter.  Extending its upper endpoint to
infinity gives (35.18).  Compactness supplies subsequential local weak
limits, and a one-plaquette observable is continuous; take the
limsup as \(L\to\infty\) to obtain (35.19).
\end{proof}

\begin{lemma}[The \(SU(3)\) Haar small ball]
\label{lem:n1v35-Haar-ball}
There are constants \(c_G>0\) and
\(0<\delta_G\leq1\), independent of \(L,\beta,t\), such that
\[
 \boxed{
 v_G(\delta)\geq c_G\delta^8
 \qquad(0<\delta\leq\delta_G).}                             \tag{35.20}
\]
\end{lemma}

\begin{proof}
Equip \(\mathfrak{su}(3)\), of real dimension eight, with any Euclidean
norm equivalent to the operator norm.  The exponential map is a
diffeomorphism from a sufficiently small ball about zero onto a
neighbourhood of the identity.  In these coordinates normalized Haar
measure has a smooth strictly positive density whose value at zero is
positive.  Shrink the coordinate ball so that the density is bounded
below by half its value at zero and so that a Euclidean ball of radius
\(c\delta\) maps inside \(B_G(\delta)\).  Its eight-dimensional
Lebesgue volume is a positive constant times \(\delta^8\), proving
(35.20).
\end{proof}

\begin{theorem}[Explicit weak-coupling plaquette bound]
\label{thm:n1v35-weak}
For \(\beta\geq\delta_G^{-2}\), choose
\(\delta=\beta^{-1/2}\).  Every translation-invariant infinite-volume
limit of (35.5) obeys
\[
 \boxed{
 {\mathbb E}_{\beta,\infty,t}z_p
 \leq
 \frac{
 C_t+\frac23\log(c_G^{-1})+\frac83\log\beta
 }{\beta}.}                                                 \tag{35.21}
\]
In particular the Wilson ensemble has the sharper value \(C_0=8\);
the estimate is uniform on the complete positive homotopy
\(0\leq t\leq1\) with \(C_t\leq392\).
\end{theorem}

\begin{proof}
Lemma~\ref{lem:n1v35-Haar-ball} gives
\[
 \log v_G(\beta^{-1/2})^{-1}
 \leq\log(c_G^{-1})+4\log\beta.
\]
Insert this and (35.15) into (35.19).
\end{proof}

\subsection{From plaquette probability to a local logarithm chart}

For a finite set \(B\) of plaquettes define the chord-bad event
\[
 {\cal F}_{B,\epsilon}
 :=
 \left\{
 \max_{p\in B}\|I-U_p\|_{\rm op}>\epsilon
 \right\}.                                                   \tag{35.22}
\]

\begin{corollary}[Block chord tail]
\label{cor:n1v35-block-tail}
Let
\[
 m_{\beta,t}:=
 \frac{
 C_t+\frac23\log(c_G^{-1})+\frac83\log\beta
 }{\beta}.
\]
For every \(\beta\geq\delta_G^{-2}\), every finite plaquette block
\(B\), and \(0<\epsilon\leq2\),
\[
 \boxed{
 \nu_{\beta,\infty,t}({\cal F}_{B,\epsilon})
 \leq
 \frac{6|B|m_{\beta,t}}{\epsilon^2}.}                       \tag{35.23}
\]
\end{corollary}

\begin{proof}
By (35.3),
\[
 \|I-U_p\|_{\rm op}>\epsilon
 \quad\Longrightarrow\quad
 z_p>\frac{\epsilon^2}{6}.
\]
Markov''s inequality and Theorem~\ref{thm:n1v35-weak} bound the
probability for one plaquette by \(6m_{\beta,t}/\epsilon^2\).
Take the union bound over \(B\).
\end{proof}

\begin{lemma}[Contractible block reconstruction]
\label{lem:n1v35-reconstruction}
Let \(K\) be a finite contractible cubical lattice complex with a fixed
spanning tree and a fixed cellular contraction of every fundamental
edge loop.  Let \(A_K\) be the largest number of plaquette faces used
by these contractions.  If
\[
 \max_{p\subset K}\|I-U_p\|_{\rm op}\leq\epsilon,            \tag{35.24}
\]
then there is a vertex gauge transformation, anchored at the tree
root, for which
\[
 \boxed{
 \max_{e\subset K}\|I-U_e^{\,g}\|_{\rm op}
 \leq A_K\epsilon.}                                         \tag{35.25}
\]
\end{lemma}

\begin{proof}
Gauge recursively along the spanning tree, making every tree link the
identity.  A nontree link then equals, up to orientation, the holonomy
of its fundamental loop.  The declared cellular contraction writes
that loop holonomy as an ordered product of at most \(A_K\) conjugates
of plaquette holonomies or their inverses.  Operator distance from the
identity is invariant under conjugation and inversion.  Applying the
telescoping inequality (35.9) to that product proves (35.25).
\end{proof}

\begin{corollary}[A summable local-chart schedule]
\label{cor:n1v35-schedule}
Let \((K_i,\beta_i,t_i,\epsilon_i)\) be any sequence of contractible
blocks and parameters, with \(B_i\) the plaquettes of \(K_i\).  If
\[
 \sum_i
 \frac{|B_i|\,[1+\log\beta_i]}
 {\beta_i\epsilon_i^2}<\infty,
 \qquad
 A_{K_i}\epsilon_i\leq\delta_{\rm ch},                       \tag{35.26}
\]
where \(\delta_{\rm ch}\) is a fixed principal-logarithm chart radius,
then the chart-failure probabilities are summable, uniformly for
\(0\leq t_i\leq1\).  On any compatible common coupling, only finitely
many of the blocks fail the chart condition almost surely.
\end{corollary}

\begin{proof}
The second condition and
Lemma~\ref{lem:n1v35-reconstruction} show that the complement of
\({\cal F}_{B_i,\epsilon_i}\) enters the declared link chart after
anchored tree gauge fixing.  Equation (35.23), with \(C_{t_i}\leq392\),
and the first condition make the failure probabilities summable.
The final statement is the first Borel--Cantelli lemma; it needs no
independence.
\end{proof}

\begin{warning}[Topology is not removed by a local plaquette bound]
\label{warn:n1v35-topology}
Lemma~\ref{lem:n1v35-reconstruction} is deliberately restricted to a
contractible block.  On a periodic torus, small plaquette curvature
does not force noncontractible Wilson lines or torons to be close to
the identity.  Those variables remain in the sector and holonomy
ledger.  V35 therefore supplies the local good-chart probability row,
not global sector tightness.
\end{warning}

\subsection{Consequences for the acquisition order}

\begin{assessment}[Progress after compact V35]
\label{ass:n1v35-status}
The positive endpoint introduced in V20 now has a volume-uniform
weak-coupling plaquette estimate, and the same proof applies directly
to the Wilson target with a much smaller comparison constant.  The
result is uniform along the positive homotopy, gives an explicit block
chord tail, and converts by tree gauge into a principal-link-chart
event on every contractible block.  Thus the deterministic local
response compiler can be paired with a summable probabilistic
good-block schedule rather than an assumed small-field restriction.
\end{assessment}

\begin{decision}[V36 returns to the first nonzero momentum orbit]
\label{dec:n1v35-next}
Construct the first nonzero internal momentum orbit using independent
internal and external Laurent phases.  Require both shifted covariance
inverses, the translated-chart cubic and quartic vertices, stationary
score and Haar rows, the complete Ward recombination, and an exact
toron-limit regression.  Carry (35.23)--(35.26) as the probability
gate for using that deterministic card; do not replace the unresolved
sector or adjacent-cylinder rows by the local chord estimate.
\end{decision}

\begin{firewall}[Claim boundary after compact V35]
\label{fire:n1v35-claim}
V35 proves a finite- and infinite-volume plaquette concentration bound
for the Wilson and positive endpoint measures and a local
contractible-block chart consequence.  It does not prove nonzero
momentum response, a Brillouin-zone mean, global stationary-Hessian
tails, comparison of the endpoint measure with a continuum target,
sector tightness, projective cylinder convergence, reflection
positivity of the limiting construction, clustering, or the
Yang--Mills mass gap.
\end{firewall}

\subsection{Parent record}

This V35 note was formed by appending the preceding material to the
validated compact V34, whose record was
\[
\begin{array}{c|c}
\text{lines}&14148\\
\text{bytes}&483056\\
\text{SHA--256}&
\texttt{E4AB11D1DCC9EC43955A92008AC9EFA838CA5FABA7BD3F92656762A6AFEE3C30}.
\end{array}
\]

% =============================================================================
% END APPENDED COMPACT-SERIES V35 MATERIAL
% =============================================================================

% =============================================================================
% BEGIN APPENDED COMPACT-SERIES V36 MATERIAL
% =============================================================================

\section{Compact-series V36: the first exact nonzero internal
momentum response orbit}
\label{sec:n1v36}

\subsection{Independent base and external phases}

V31--V34 evaluated the complete stationary response at internal
momentum zero.  The new acquisition is the self-inverse nonzero orbit
\[
 \kappa=\pi\widehat e_2,\qquad
 q=q\,\widehat e_1,\qquad
 \text{external polarization } \widehat e_2.                \tag{36.1}
\]
Coordinates in (36.1) are one-based: the executable record stores
internal axis one, external axis zero, and polarization one.
The internal and external axes are deliberately different.  Thus a
base-momentum sign can never be silently reused as an external
derivative phase.

For a Laurent monomial \(z^\alpha\), define its exact jet at (36.1) by
\[
 {\cal J}_{\kappa}^{[j]}(z^\alpha)
 :=
 (-1)^{\alpha_2}
 \left.\frac{d^j}{dq^j}e^{iq\alpha_1}\right|_{q=0},
 \qquad 0\leq j\leq2.                                      \tag{36.2}
\]
The real jet representation used by the verifier stores
\[
 {\cal J}_{\kappa}(z^\alpha)
 =
 (-1)^{\alpha_2}
 \left(1,\alpha_1,-\alpha_1^2\right),                       \tag{36.3}
\]
where the middle component is the coefficient of \(i\).
Multiplication is therefore
\[
 (a_0,a_1,a_2)(b_0,b_1,b_2)
 =
 (a_0b_0,\,
  a_1b_0+a_0b_1,\,
  a_2b_0-2a_1b_1+a_0b_2).                                  \tag{36.4}
\]

\begin{lemma}[Exactness of the orbit jet ring]
\label{lem:n1v36-jet-ring}
Equations (36.2)--(36.4) give the value, coefficient of \(i\) in
the first derivative, and real second derivative of every finite
Laurent polynomial at \(\kappa+q\widehat e_1\).  Conjugate transpose
acts by
\[
 (A_0,A_1,A_2)^*
 =
 (A_0^*,-A_1^*,A_2^*).                                     \tag{36.5}
\]
\end{lemma}

\begin{proof}
For a monomial,
\[
 e^{i\alpha\cdot(\kappa+q\widehat e_1)}
 =
 (-1)^{\alpha_2}e^{iq\alpha_1}.
\]
Differentiate twice.  The second derivative of a product contains
two first-derivative factors, each carrying \(i\), which gives the
minus sign in the middle term of (36.4).  Linearity proves the
polynomial statement.  Complex conjugation reverses the sign of the
coefficient of \(i\), proving (36.5).
\end{proof}

\subsection{The independent-leg vertex rule}

The cubic word bank has two vertical legs and one background leg.
For a phase tuple
\(\boldsymbol\alpha=(\alpha^{(0)},\alpha^{(1)},\alpha^{(2)})\),
the base and plus-oriented external weights are
\[
 \eta=(1,-1,0),\qquad
 \sigma_+=(0,-1,1),                                        \tag{36.6}
\]
\[
 w_\kappa(\boldsymbol\alpha)
 =
 (-1)^{
 \alpha^{(0)}_2-\alpha^{(1)}_2},
 \qquad
 m_+(\boldsymbol\alpha)
 =
 -\alpha^{(1)}_1+\alpha^{(2)}_1.                            \tag{36.7}
\]
The minus-oriented writer uses
\(\sigma_-=(0,1,-1)\).  Thus the phase contribution to order \(j\)
is exactly
\[
 w_\kappa(\boldsymbol\alpha)m_\pm(\boldsymbol\alpha)^j.     \tag{36.8}
\]
The background lift jets remain those of the external router at
momentum \(q\widehat e_1\); they are not evaluated at \(\kappa\).

For a quartic word, assign internal rates \(+1,-1\) to its two
vertical legs and zero to the two background legs, and external rates
\(0,0,+1,-1\), respectively.  At every occurrence with displacement
\(a\), its exact scalar factor is
\[
 (-1)^{\eta_\ell a_2}
 \left(1,\sigma_\ell a_1,-(\sigma_\ell a_1)^2\right).       \tag{36.9}
\]
This factor is inserted before the noncommutative colour word is
multiplied.  Consequently coincident translated background and
vertical insertions retain their declared left-translated order; they
are never symmetrized into an exponential-of-sum chart.

\begin{proposition}[Self-inverse plus/minus and Hermitian gates]
\label{prop:n1v36-gates}
For both the Wilson and endpoint banks, the exact computation gives
\[
 A^{[0]\,T}=A^{[0]},\qquad
 A^{[1]\,T}=-A^{[1]},\qquad
 A^{[2]\,T}=A^{[2]},                                       \tag{36.10}
\]
where \(A^{[1]}\) denotes the stored imaginary coefficient.  For the
Wilson cubic bank, the Wilson bank evaluated with the endpoint
router, the translated completion bank, and their endpoint sum,
\[
 P_-^{[j]}=(-1)^jP_+^{[j]},
 \qquad 0\leq j\leq2.                                      \tag{36.11}
\]
The Wilson and endpoint quartic first-jet contractions, bubble first
jets, and complete response first jets all vanish.
\end{proposition}

\begin{proof}
The generated incidence polynomials are evaluated coefficient by
coefficient using (36.2).  Equality (36.10) is checked in all
\(45^2\) entries for each of the six matrices.  The translated cubic
compiler retains a phase vector for every leg separately, applies
(36.6)--(36.8), and checks (36.11) entry by entry.  The quartic
compiler applies (36.9) before colour multiplication, contracts the
complete \(SU(3)\) dual-colour pair, and then tests the stated scalar
first jets.  Every residual is exactly zero in rational arithmetic.
\end{proof}

\subsection{Covariances and the complete cell response}

Let \(D_R(k)\) be the reduced plaquette incidence and \(K_R(k)\) the
384-row translated-plaquette-difference incidence.  In the established
normalization put
\[
 H_{\rm W}(k):=D_R(k)^*D_R(k),\qquad
 H_*(k):=H_{\rm W}(k)+3K_R(k)^*K_R(k),                       \tag{36.12}
\]
\[
 G_{\rm W}(k):=3H_{\rm W}(k)^{-1},\qquad
 G_*(k):=3H_*(k)^{-1}.                                     \tag{36.13}
\]
The V18--V21 positive floor makes both inverses well defined.

For \(e\in\{{\rm W},*\}\), let \(P_e^\pm(\kappa,q)\) be the complete
cubic writers and \(Q_e(\kappa,q)\) the complete quartic writer in the
translated chart.  The cell bubble and tadpole are
\[
 {\cal B}_e(q;\kappa)
 =
 c_e\Tr\!\left[
 G_e(\kappa)P_e^+(\kappa,q)
 G_e(\kappa+q\widehat e_1)P_e^-(\kappa,q)
 \right],                                                   \tag{36.14}
\]
\[
 {\cal T}_e(q;\kappa)
 =
 \frac12\Tr\!\left[G_e(\kappa)Q_e(\kappa,q)\right],
 \qquad
 c_{\rm W}=\frac13,\quad c_*=3.                             \tag{36.15}
\]
The stationary-score rows vanish through order two separately for
the Wilson, endpoint-routed Wilson, and completion actions.  The Haar
left-translation row vanishes identically.  Hence the complete
measured response on this declared cell is
\[
 {\cal W}_e^{[j]}(\kappa)
 =
 {\cal T}_e^{[j]}(\kappa)+{\cal B}_e^{[j]}(\kappa).          \tag{36.16}
\]

\begin{theorem}[First nonzero internal-momentum response card]
\label{thm:n1v36-nonzero-card}
At the orbit (36.1), the exact tadpole triples are
\[
 \boxed{
 ({\cal T}_{\rm W}^{[0]},{\cal T}_{\rm W}^{[1]},
  {\cal T}_{\rm W}^{[2]})
 =
 \left(\frac{302}{5},0,-\frac{7457}{168}\right),}           \tag{36.17}
\]
\[
 \boxed{
 ({\cal T}_{*}^{[0]},{\cal T}_{*}^{[1]},
  {\cal T}_{*}^{[2]})
 =
 \left(
 \frac{38665469995}{505450939},
 0,
 -\frac{28758791690948779}{182341426244250}
 \right).}                                                  \tag{36.18}
\]
The exact bubble triples are
\[
 \boxed{
 ({\cal B}_{\rm W}^{[0]},{\cal B}_{\rm W}^{[1]},
  {\cal B}_{\rm W}^{[2]})
 =
 \left(
 -\frac{14162}{245},
 0,
 \frac{101613881}{2881200}
 \right),}                                                  \tag{36.19}
\]
\[
 \boxed{
 ({\cal B}_{*}^{[0]},{\cal B}_{*}^{[1]},
  {\cal B}_{*}^{[2]})
 =
 \left(
 -\frac{183187480214}{2527254695},
 0,
 \frac{4599350785850746783639351}
 {737317160910043246806000}
 \right).}                                                  \tag{36.20}
\]
Consequently
\[
 \boxed{
 ({\cal W}_{\rm W}^{[0]},{\cal W}_{\rm W}^{[1]},
  {\cal W}_{\rm W}^{[2]})
 =
 \left(
 \frac{636}{245},
 0,
 -\frac{26273669}{2881200}
 \right),}                                                  \tag{36.21}
\]
\[
 \boxed{
 ({\cal W}_{*}^{[0]},{\cal W}_{*}^{[1]},
  {\cal W}_{*}^{[2]})
 =
 \left(
 \frac{1448552823}{361036385},
 0,
 -\frac{111689915331712918387988497}
 {737317160910043246806000}
 \right).}                                                  \tag{36.22}
\]
The endpoint-minus-Wilson cell difference is
\[
 \boxed{
 \Delta^{[0,1,2]}(\pi\widehat e_2)
 =
 \left(
 \frac{5011052031}{3538156573},
 0,
 -\frac{107153116857979395343397926}
 {752677935095669147781125}
 \right).}                                                  \tag{36.23}
\]
Numerically its zeroth and second components are
\[
 \Delta^{[0]}\approx1.41628894245094,\qquad
 \Delta^{[2]}\approx-142.362505743389.                      \tag{36.24}
\]
\end{theorem}

\begin{proof}
Differentiate (36.14) in the exact jet ring
of Lemma~\ref{lem:n1v36-jet-ring}, including the shifted inverse
through
\[
 G''=-GH''G,\qquad
 G''''=2GH''GH''G-GH''''G.                                       \tag{36.25}
\]
Differentiate (36.15), contract all \(45^2\) entries, and use the
stationary and Haar zeros preceding (36.16).  Exact fraction reduction
gives (36.17)--(36.23).
\end{proof}

\subsection{A sign change in the actual internal-momentum integrand}

\begin{corollary}[The response difference is not momentum constant]
\label{cor:n1v36-sign-change}
For the fixed external transverse channel of (36.1),
\[
 \Delta^{[0]}(0)
 =
 -\frac{2373732}{589225}<0,\qquad
 \Delta^{[0]}(\pi\widehat e_2)
 =
 \frac{5011052031}{3538156573}>0.                           \tag{36.26}
\]
There is therefore an \(s_*\in(0,\pi)\) such that
\[
 \Delta^{[0]}(s_*\widehat e_2)=0.                           \tag{36.27}
\]
\end{corollary}

\begin{proof}
The first value is the V34 transverse toron card and the second is
(36.23).  The incidence and vertex entries are finite Laurent
polynomials.  The uniform positive Hessian floor makes their inverses
continuous on the real momentum torus.  Hence
\(\Delta^{[0]}(s\widehat e_2)\) is continuous, and the intermediate
value theorem proves (36.27).
\end{proof}

\begin{warning}[One positive cell does not determine the mean]
\label{warn:n1v36-mean}
The sign change (36.26) rules out any argument that infers the
Brillouin-zone average from the toron value or from one nonzero sample.
The remaining edge scalar needs either exact orbit cancellation or a
certified global quadrature with an analytic-strip error.  V36 supplies
the first nontrivial regression point for that global certificate.
\end{warning}

\subsection{Executable record and next acquisition}

The exact verifier is
\[
 \texttt{scripts/verify\_ym\_v36\_first\_nonzero\_orbit.py}. \tag{36.28}
\]
Its source record is
\[
\begin{array}{c|c}
\text{lines}&1533\\
\text{bytes}&46742\\
\text{SHA--256}&
\texttt{586C31525D1A0133EBCF59EABE63EA545C6393F065F15847EC4FD5FE57F0020A}.
\end{array}                                                  \tag{36.29}
\]
The full generalized-code-path toron regression reproduces the V34
fractions exactly and has canonical payload SHA--256
\[
 \texttt{5C4214862089A555A9A2147BCED67D5F5AA3C1B9D6FA4CFD6711E3224BA3E029}.
                                                                    \tag{36.30}
\]
The nonzero-orbit payload has canonical SHA--256
\[
 \boxed{
 \texttt{437181372E25FCA3679A1A80BBC32408225B7820607970AE85C35C106C0A48D8}.}
                                                                    \tag{36.31}
\]

\begin{assessment}[Progress after compact V36]
\label{ass:n1v36-status}
The complete translated-chart response compiler now evaluates one
genuinely nonzero internal momentum with base and external phases
independent.  It passes a full toron regression, Hermitian covariance,
cubic plus/minus, stationary-score, Haar, and reflection first-jet
gates.  The exact nonzero response (36.23) also proves that the
zeroth-order endpoint difference changes sign along an internal
momentum axis.
\end{assessment}

\begin{decision}[V37 orbit pairing before global quadrature]
\label{dec:n1v36-next}
Derive the exact transformation of the complete cell response under
\(k\mapsto-k\), coordinate reflections, and the stabilizer of the
external axis.  Use it to identify the minimal fundamental momentum
domain and to state the Ward telescope on paired cells.  Only after
that algebraic reduction should the programme construct interval
quadrature for the remaining domain; otherwise the numerical card
would pay for cancellations that should be exact.
\end{decision}

\begin{firewall}[Claim boundary after compact V36]
\label{fire:n1v36-claim}
V36 proves one complete nonzero internal-momentum response card and
one sign-change consequence.  It does not compute the
Brillouin-zone mean, prove the cellwise or summed Ward telescope,
establish a universal beta coefficient for the endpoint homotopy,
complete SCIC or ADMC, prove sector tightness or projective continuum
convergence, establish terminal clustering, or prove the
Yang--Mills mass gap.
\end{firewall}

\subsection{Parent record}

This V36 note was formed by appending the preceding material to the
validated compact V35, whose record was
\[
\begin{array}{c|c}
\text{lines}&14600\\
\text{bytes}&498660\\
\text{SHA--256}&
\texttt{23F674C142CA9B298546508F7A3E9EAF2CE2E154078C32B0BA2FC8660AD25098}.
\end{array}
\]

% =============================================================================
% END APPENDED COMPACT-SERIES V36 MATERIAL
% =============================================================================

% =============================================================================
% BEGIN APPENDED COMPACT-SERIES V37 MATERIAL
% =============================================================================

\section{Compact-series V37: exact momentum-orbit reduction and
the finite-grid Ward telescope}
\label{sec:n1v37}

\subsection{Coordinate-free response notation}

Let
\[
 {\mathscr W}^{[j]}_{e;a,b}(k),\qquad
 e\in\{{\rm W},*\},\quad 0\leq j\leq2,                     \tag{37.1}
\]
denote the complete measured cell response in which the external
momentum is \(q\widehat e_a\), the retained background polarization is
\(\widehat e_b\), and \(j\) external momentum derivatives are taken at
\(q=0\).  Complete means that the cubic bubble, quartic tadpole,
stationary-centre, Haar-density, and quotient-coordinate terms are
recombined before the scalar trace is formed.  The endpoint difference
is
\[
 \Delta_{a,b}^{[j]}(k)
 :=
 {\mathscr W}^{[j]}_{*;a,b}(k)
 -
 {\mathscr W}^{[j]}_{{\rm W};a,b}(k).                       \tag{37.2}
\]

The raw parity-tree matrices are not tensors under every lattice
symmetry.  If \(R\) is a signed coordinate permutation, let
\[
 S_R(k):{\cal H}_k^{\rm red}\longrightarrow
 {\cal H}_{Rk}^{\rm red}                                   \tag{37.3}
\]
be the map induced by transforming the complete sixty-four-link
field and then applying the fixed tree quotient at \(Rk\).  Factors
\(S_R\), their inverses, and the corresponding density Jacobian cancel
only in the complete measured scalar (37.1).  This is why the symmetry
theorem below is not asserted entrywise for a raw \(45\times45\)
matrix.

\subsection{Signed hypercubic covariance}

Write a signed permutation as
\[
 R\widehat e_\mu
 =
 \epsilon_\mu\widehat e_{\pi(\mu)},
 \qquad \epsilon_\mu\in\{-1,1\}.                            \tag{37.4}
\]

\begin{theorem}[Exact covariance of the complete cell scalar]
\label{thm:n1v37-covariance}
For \(j=0,1,2\),
\[
 \boxed{
 {\mathscr W}^{[j]}_{e;\pi(a),\pi(b)}(Rk)
 =
 \epsilon_a^j
 {\mathscr W}^{[j]}_{e;a,b}(k).}                            \tag{37.5}
\]
The same identity holds for \(\Delta^{[j]}\).
\end{theorem}

\begin{proof}
Both \(V_{\rm W}\) and \({\cal O}_1\) are sums over all translated
oriented plaquettes and all four displacement directions.  A signed
coordinate permutation merely permutes these owners, possibly
reversing a plaquette orientation.  Real trace and Hilbert--Schmidt
norm are unchanged by inversion, so both actions are invariant.
Product Haar measure is invariant as well.

Transform the complete link field by \(R\).  On reduced coordinates,
the vertical Hessian, its inverse, and every differentiated vertex are
related by the maps (37.3).  Each \(S_R\) adjacent to a propagator
cancels with its inverse, and cyclicity removes the remaining
similarity factors in every bubble and tadpole trace.  The measured
stationary and density terms transform as scalars because they are
obtained by differentiating the same invariant finite-dimensional
density.  Hence the undifferentiated response with external vector
\(q\widehat e_a\) equals the transformed response with external vector
\(q\epsilon_a\widehat e_{\pi(a)}\).

Changing the sign of the background polarization contributes twice
and therefore gives \(\epsilon_b^2=1\).  Taking \(j\) derivatives with
respect to the positively oriented scalar momentum at the transformed
point contributes \(\epsilon_a^j\).  This proves (37.5), and
subtraction proves the endpoint statement.
\end{proof}

\begin{corollary}[Internal inversion parity]
\label{cor:n1v37-inversion}
For every \(a,b\),
\[
 \boxed{
 {\mathscr W}^{[j]}_{e;a,b}(-k)
 =
 (-1)^j{\mathscr W}^{[j]}_{e;a,b}(k),\qquad
 \Delta_{a,b}^{[j]}(-k)
 =
 (-1)^j\Delta_{a,b}^{[j]}(k).}                              \tag{37.6}
\]
In particular every odd jet vanishes at a self-inverse momentum
\(k=-k\) modulo \(2\pi{\mathbb Z}^4\).
\end{corollary}

\begin{proof}
Take \(R=-I\) in Theorem~\ref{thm:n1v37-covariance}.  Then
\(\pi\) is the identity and every \(\epsilon_\mu=-1\).  At a
self-inverse momentum the two sides refer to the same cell, so an odd
jet equals its negative.
\end{proof}

This proves conceptually the zero first jets certified independently
at \(k=0\) in V33--V34 and at \(k=\pi\widehat e_2\) in V36.

\subsection{The minimal domain for the V36 channel}

Fix \(a=1\) and \(b=2\), as in V36, and let
\[
 {\cal G}_{1,2}
 :=
 \{R:\ R\widehat e_1=\pm\widehat e_1,\ 
        R\widehat e_2=\pm\widehat e_2,\ 
        R\{\widehat e_3,\widehat e_4\}
        =\{\pm\widehat e_3,\pm\widehat e_4\}\}.              \tag{37.7}
\]
This group has \(2^4\cdot2=32\) elements.  For even \(j\), (37.5)
shows that the V36 scalar is constant on every
\({\cal G}_{1,2}\)-orbit.

Let
\[
 {\cal D}_{1,2}
 :=
 \left\{
 k\in[0,\pi]^4:\ 0\leq k_4\leq k_3\leq\pi
 \right\}.                                                   \tag{37.8}
\]

\begin{theorem}[Exact fundamental-domain reduction]
\label{thm:n1v37-domain}
Every momentum in \({\mathbb T}^4\) has a representative in
\({\cal D}_{1,2}\), and two interior points of
\({\cal D}_{1,2}\) lie in the same \({\cal G}_{1,2}\)-orbit only if
they are equal.  Therefore, for every integrable even response jet,
\[
 \boxed{
 \int_{{\mathbb T}^4}
 {\mathscr W}^{[j]}_{e;1,2}(k)\,dk
 =
 32\int_{{\cal D}_{1,2}}
 {\mathscr W}^{[j]}_{e;1,2}(k)\,dk,}                         \tag{37.9}
\]
where boundary sets have Lebesgue measure zero.  The same formula
holds for \(\Delta^{[j]}\).
\end{theorem}

\begin{proof}
Independent sign changes fold each coordinate into \([0,\pi]\).
The remaining transposition of coordinates three and four orders
their folded values.  Away from the hyperplanes
\(k_\mu\in\{0,\pi\}\) and \(k_3=k_4\), no nonidentity group element
fixes a point; hence the orbit has 32 elements and the interior
representative is unique.  Partition the torus into these images and
use the even-\(j\) invariance from (37.5).  Boundaries have measure
zero, giving (37.9).
\end{proof}

For the uniform grid
\[
 \Gamma_M
 :=
 \left\{\frac{2\pi}{M}n:\ n\in({\mathbb Z}/M{\mathbb Z})^4\right\},
                                                                    \tag{37.10}
\]
let \({\cal R}_M\) be one representative of every
\({\cal G}_{1,2}\)-orbit and let
\[
 m_M(r):=\frac{32}{|\operatorname{Stab}_{{\cal G}_{1,2}}(r)|}.
                                                                    \tag{37.11}
\]

\begin{corollary}[Exact finite-grid orbit weights]
\label{cor:n1v37-grid}
For even \(j\),
\[
 \boxed{
 \frac1{M^4}\sum_{k\in\Gamma_M}
 {\mathscr W}^{[j]}_{e;1,2}(k)
 =
 \frac1{M^4}\sum_{r\in{\cal R}_M}
 m_M(r){\mathscr W}^{[j]}_{e;1,2}(r).}                      \tag{37.12}
\]
For odd \(j\),
\[
 \boxed{
 \sum_{k\in\Gamma_M}{\mathscr W}^{[j]}_{e;1,2}(k)=0.}       \tag{37.13}
\]
Both statements also hold for the endpoint difference.
\end{corollary}

\begin{proof}
Equation (37.12) is orbit--stabilizer together with (37.5).
For (37.13), pair \(k\) with \(-k\) and use (37.6).
Self-inverse cells contribute zero by
Corollary~\ref{cor:n1v37-inversion}.
\end{proof}

\subsection{The nonzero-grid Ward telescope}

The symmetry reduction above concerns formal external jets at
\(q=0\).  The exact Ward identity is most transparent before that
limit, at a nonzero momentum \(q\in\Gamma_M\).  Let
\[
 {\cal H}_M=\bigoplus_{k\in\Gamma_M}{\cal H}^{\rm red}_k    \tag{37.14}
\]
and let \(\widehat H_{e,M}(B)\) be the complete measured stationary
Hessian on this direct sum.  The adjective \emph{measured} means that
the stationary-coordinate and Haar Jacobians are included in the
scalar Gaussian density.

Choose a periodic infinitesimal gauge parameter with Fourier support
\(\{q,-q\}\).  Its finite gauge path sends the flat background to a
pure-gauge path \(B_s\).  Momentum addition by \(q\) partitions
\(\Gamma_M\) into disjoint cycles
\[
 {\cal C}(k,q)=\{k,k+q,\ldots,k+(\ell-1)q\}.                 \tag{37.15}
\]
Gauge covariance acts separately on the direct sum over each such
cycle.

\begin{theorem}[Exact finite-grid Ward telescope]
\label{thm:n1v37-Ward-telescope}
Let \(q\in\Gamma_M\setminus\{0\}\).  Suppose the complete response is
formed from the exact gauge-invariant action \(V_e\) and normalized
Haar measure, with stationary reduction performed before the
background derivatives.  Let
\({\cal L}_{e,M}(k;q)\) be its complete second-amplitude longitudinal
cell response along the pure-gauge path \(B_s\).  Then, on every
cycle (37.15),
\[
 \boxed{
 \sum_{k''\in{\cal C}(k,q)}
 {\cal L}_{e,M}(k'';q)=0.}                                   \tag{37.16}
\]
Equivalently there is a periodic cycle potential
\(\Phi_{e,M}(\,\cdot\,;q)\) such that
\[
 \boxed{
 {\cal L}_{e,M}(k;q)
 =
 \Phi_{e,M}(k+q;q)-\Phi_{e,M}(k;q).}                        \tag{37.17}
\]
Consequently
\[
 \boxed{
 \sum_{k\in\Gamma_M}{\cal L}_{e,M}(k;q)=0.}                 \tag{37.18}
\]
These identities hold separately for \(e={\rm W}\), \(e=*\), and
every positive interpolation \(V_{\rm W}+t{\cal O}_1\).
\end{theorem}

\begin{proof}
On the cycle subspace, a finite gauge transformation is an invertible
coordinate transformation \({\cal R}_s\) which mixes only momenta in
that cycle.  Exact gauge covariance gives
\[
 \widehat H_{e,M}(B_s)
 =
 {\cal R}_s^{-*}\widehat H_{e,M}(0){\cal R}_s^{-1}.          \tag{37.19}
\]
The determinant factors from \({\cal R}_s\) are precisely cancelled
by the measured Haar and stationary-coordinate Jacobians.  Hence the
complete Gaussian scalar on each cycle,
\[
 {\cal F}_{e,{\cal C}}(s)
 =
 \frac12\log\det\widehat H_{e,M}(B_s)
 +h_{e,M}(B_s),                                             \tag{37.20}
\]
is constant in \(s\).  Its second derivative at zero is therefore
zero.  Expanding that derivative into diagonal momentum blocks gives
exactly the complete tadpole, bubble, stationary, and Haar cell
writers whose sum is the left side of (37.16).  This proves the cycle
zero.

Enumerate a cycle by \(k_r=k_0+rq\) and define
\[
 \Phi(k_0;q)=0,\qquad
 \Phi(k_r;q)=
 \sum_{u=0}^{r-1}{\cal L}(k_u;q).                           \tag{37.21}
\]
The cycle zero makes \(\Phi(k_\ell;q)=\Phi(k_0;q)\), so \(\Phi\) is
periodic and (37.17) follows, up to the harmless reversal obtained by
choosing the opposite cumulative convention.  Summing (37.16) over
cycles proves (37.18).  Gauge invariance and positivity hold along the
whole declared interpolation, so the proof is uniform in \(t\).
\end{proof}

\begin{warning}[The discrete Ward zero is not yet the formal
\(q=0\) jet]
\label{warn:n1v37-Ward-limit}
Theorem~\ref{thm:n1v37-Ward-telescope} concerns nonzero periodic
external momenta.  A constant toron is not generated by a periodic
infinitesimal gauge parameter, which is why its individual cell
response in V33--V34 need not vanish.  Passing (37.18) to the formal
external derivatives at \(q=0\) requires a sequence
\(q_M\in\Gamma_M\setminus\{0\}\) with \(q_M\to0\), together with a
uniform comparison between the nonzero-\(q_M\) finite-grid response
and the differentiated \(q=0\) integrand.  That comparison must retain
the toron cell and cannot be inferred from pointwise smoothness at
fixed \(M\).
\end{warning}

\begin{corollary}[A sufficient Ward-limit card]
\label{cor:n1v37-Ward-limit-card}
Suppose \(q_M\to0\) through nonzero grid momenta and, for some
\(j\in\{0,1,2\}\), the difference quotient
\({\cal D}_{M}^{[j]}(k;q_M)\) converges to
\({\mathscr W}^{[j]}_{e;1,1}(k)\) with
\[
 \frac1{M^4}\sum_{k\in\Gamma_M}
 \left|
 {\cal D}_{M}^{[j]}(k;q_M)
 -
 {\mathscr W}^{[j]}_{e;1,1}(k)
 \right|\longrightarrow0.                                  \tag{37.22}
\]
If the corresponding exact finite-\(q_M\) Ward combination has zero
grid mean, then
\[
 \lim_{M\to\infty}
 \frac1{M^4}\sum_{k\in\Gamma_M}
 {\mathscr W}^{[j]}_{e;1,1}(k)=0.                           \tag{37.23}
\]
\end{corollary}

\begin{proof}
The exact Ward mean is zero by (37.18).  The absolute difference
between that mean and the mean in (37.23) is bounded by (37.22), which
tends to zero.
\end{proof}

\subsection{Updated acquisition order}

\begin{assessment}[Progress after compact V37]
\label{ass:n1v37-status}
The response integrand now has an exact coordinate-free hypercubic
law.  For the V36 transverse channel, every even-jet Brillouin integral
reduces to the \(1/32\)-sized domain (37.8), with exact stabilizer
weights on a finite grid; every odd grid mean vanishes algebraically.
For nonzero periodic longitudinal momentum, gauge covariance gives a
cyclewise discrete Ward divergence and an exact zero sum.  The
remaining Ward obstruction is narrowed to the uniform
nonzero-\(q\)-to-formal-jet estimate (37.22).
\end{assessment}

\begin{decision}[V38 builds the Ward-limit majorant]
\label{dec:n1v37-next}
Use the V18 complex strip, the V21 sharp Laurent stocks, and the
independent-leg compiler to bound the third external derivative of
the complete finite-\(q\) response uniformly in internal momentum.
Apply Taylor''s theorem to (37.22), keeping the toron cell explicit.
If the resulting stock is too crude after multiplication by the grid
volume, return upstream and use the orbit reduction (37.12) plus
resolvent identities before attempting interval quadrature.
\end{decision}

\begin{firewall}[Claim boundary after compact V37]
\label{fire:n1v37-claim}
V37 proves exact response symmetries, orbit weights, and the nonzero
finite-grid Ward telescope.  It does not prove the uniform Ward-limit
estimate (37.22), evaluate the even Brillouin mean, complete SCIC or
ADMC, establish sector tightness or projective continuum convergence,
prove reflection positivity and clustering of a limiting theory, or
prove the Yang--Mills mass gap.
\end{firewall}

\subsection{Parent record}

This V37 note was formed by appending the preceding material to the
validated compact V36, whose record was
\[
\begin{array}{c|c}
\text{lines}&15006\\
\text{bytes}&511923\\
\text{SHA--256}&
\texttt{2B737CB7D11C2FED1B9719824ED13E6591F4940611AEEB2E9B1ACE1636F33853}.
\end{array}
\]

% =============================================================================
% END APPENDED COMPACT-SERIES V37 MATERIAL
% =============================================================================

% =============================================================================
% BEGIN APPENDED COMPACT-SERIES V38 MATERIAL
% =============================================================================

\section{Compact-series V38: a uniform Ward-limit theorem from
three nonzero grid momenta}
\label{sec:n1v38}

\subsection{A common analytic stock}

Let \(\rho_{21}>0\) be the explicit inverse strip of (21.21), and set
\[
 \rho_{38}:=\frac{\rho_{21}}4.                              \tag{38.1}
\]
For real internal momentum \(k\), complex external scalar \(z\) with
\(|\operatorname{Im}z|\leq\rho_{38}\), and
\(0\leq t\leq1\), every covariance in the response is evaluated either
at \(k\) or at \(k+z\widehat e_1\).  Every background router is
evaluated at \(z\widehat e_1\).  Thus the V21 inverse bound applies to
all of them.

Put
\[
 g_{38}:=\frac{2500}{49}.                                  \tag{38.2}
\]
Using the exact independent-leg word bank, define the finite strip
stocks
\[
\begin{split}
 p_{38}
 &:=
 \sup_{e,t,k,z}
 \max\{\|P_{e,t}^+(k,z)\|,\|P_{e,t}^-(k,z)\|\},\\
 q_{38}
 &:=
 \sup_{e,t,k,z}\|Q_{e,t}(k,z)\|,\\
 j_{38}
 &:=
 \sup_{e,t,k,z}|J_{e,t}^{\rm stat+Haar}(k,z)|,              \tag{38.3}
\end{split}
\]
where the suprema use
\(e\in\{{\rm W},*\}\), \(k\in{\mathbb T}^4\),
\(|\operatorname{Im}z|\leq\rho_{38}\), and \(0\leq t\leq1\).
The last row is the already recombined stationary-coordinate and Haar
scalar.

\begin{lemma}[The V38 stocks are finite and volume independent]
\label{lem:n1v38-stocks}
The numbers in (38.3) are finite and depend only on the fixed
parity-cell incidence and word banks.  They have explicit upper bounds
obtained by replacing each Laurent monomial \(z^\alpha\) by
\(e^{|\alpha|_1\rho_{38}}\), every inverse by \(g_{38}\), and every
finite matrix product by the operator-norm product.
\end{lemma}

\begin{proof}
The Wilson and completion words have fixed finite length and finite
support.  Their independently labelled derivatives through background
order two are therefore finite Laurent sums.  The harmonic routers and
stationary jets are finite sums and products of those Laurent matrices
and the inverses controlled by (38.2).  The same is true of the Haar
row.  On the closed strip, every monomial satisfies
\[
 |e^{i\alpha\cdot(k+z\widehat e_1)}|
 \leq e^{|\alpha|_1\rho_{38}}.                              \tag{38.4}
\]
Summing absolute coefficient norms gives literal bounds for (38.3).
No step contains the number of momentum cells or the physical
four-volume.
\end{proof}

Define
\[
 \boxed{
 B_{38}:=
 45\left(
 \frac12 g_{38}q_{38}
 +3g_{38}^2p_{38}^2
 \right)+j_{38}.}                                           \tag{38.5}
\]

\begin{proposition}[Uniform scalar strip bound]
\label{prop:n1v38-strip-bound}
Let \({\cal L}_{e,M,t}(k;z)\) be the complete longitudinal response
of V37 before the external-momentum limit.  Then
\[
 \boxed{
 |{\cal L}_{e,M,t}(k;z)|\leq B_{38}}                         \tag{38.6}
\]
on the strip in (38.3), uniformly in \(M,k,e,t\).
\end{proposition}

\begin{proof}
For \(45\times45\) matrices,
\[
 |\Tr(A_1\cdots A_r)|
 \leq45\prod_{\ell=1}^r\|A_\ell\|_{\rm op}.                 \tag{38.7}
\]
The tadpole has coefficient \(1/2\), while the largest bubble
coefficient on the Wilson-to-endpoint interpolation is three.  Insert
(38.2)--(38.3) in the complete tadpole--bubble formula and add the
recombined scalar \(j_{38}\).  This is exactly (38.5).
\end{proof}

\subsection{Averaged analytic response}

For the normalized momentum sum put
\[
 A_{e,M,t}(z)
 :=
 \frac1{M^4}\sum_{k\in\Gamma_M}
 {\cal L}_{e,M,t}(k;z).                                    \tag{38.8}
\]

\begin{lemma}[Uniform Cauchy derivatives]
\label{lem:n1v38-Cauchy}
The function (38.8) is analytic for
\(|\operatorname{Im}z|<\rho_{38}\), satisfies
\(|A_{e,M,t}(z)|\leq B_{38}\), and, for real
\(|x|\leq\rho_{38}/2\),
\[
 \boxed{
 |A_{e,M,t}^{(r)}(x)|
 \leq
 r!\left(\frac2{\rho_{38}}\right)^rB_{38},
 \qquad 0\leq r\leq3.}                                     \tag{38.9}
\]
\end{lemma}

\begin{proof}
Every summand is a finite composition of entire Laurent vertices and
matrix inverses with a common nonzero determinant on the V21 strip.
It is therefore analytic there.  Proposition
\ref{prop:n1v38-strip-bound} and normalized summation give the same
bound for \(A\), with no factor \(M^4\).  The circle of radius
\(\rho_{38}/2\) about such a real \(x\) remains in the strip after the
harmless further restriction \(|x|\leq\rho_{38}/2\).  Cauchy''s
derivative estimate gives (38.9).
\end{proof}

To avoid an endpoint ambiguity in the preceding sentence, one may use
the rectangle
\[
 |\operatorname{Re}z|\leq\rho_{38},\qquad
 |\operatorname{Im}z|\leq\rho_{38}                          \tag{38.10}
\]
after periodic translation of the real part; the response is periodic
in real momentum.  The radius-\(\rho_{38}/2\) Cauchy disc is then
contained in the declared analytic neighbourhood.

\subsection{Forward differences based away from the toron}

Let
\[
 h_M:=\frac{2\pi}{M}.                                       \tag{38.11}
\]
For \(M\geq4\), the momenta \(h_M,2h_M,3h_M\) are nonzero grid
momenta modulo \(2\pi\).  Theorem~\ref{thm:n1v37-Ward-telescope}
therefore gives
\[
 A_{e,M,t}(nh_M)=0,\qquad n=1,2,3.                          \tag{38.12}
\]

\begin{lemma}[Three zero samples control the first three formal rows]
\label{lem:n1v38-three-samples}
Let \(A\in C^3([0,3h])\) and suppose
\(A(h)=A(2h)=A(3h)=0\).  Then
\[
\begin{split}
 |A(0)|
 &\leq h\sup_{[0,h]}|A''|,\\
 |A''(0)|
 &\leq\frac32h\sup_{[0,2h]}|A''''|,\\
 |A''''(0)|
 &\leq2h\sup_{[0,3h]}|A''''''|.                               \tag{38.13}
\end{split}
\]
\end{lemma}

\begin{proof}
The first row is the fundamental theorem of calculus applied between
zero and \(h\).  For the second, the zero first forward difference
based at \(h\) is
\[
 0=\frac{A(2h)-A(h)}h
 =\frac1h\int_h^{2h}A''(s)\,ds.                              \tag{38.14}
\]
Subtract \(A''(0)\) inside the integral and use
\[
 |A''(s)-A''(0)|\leq s\sup_{[0,2h]}|A''''|.
\]
The average of \(s\) on \([h,2h]\) is \(3h/2\).

For the third row, the zero second forward difference based at \(h\)
has the integral representation
\[
 0=
 \frac{A(3h)-2A(2h)+A(h)}{h^2}
 =
 \int_0^1\int_0^1
 A''''(h(1+u+v))\,du\,dv.                                    \tag{38.15}
\]
Subtract \(A''''(0)\).  The average of \(h(1+u+v)\) on the unit square
is \(2h\), which proves the last inequality.
\end{proof}

\begin{theorem}[Quantitative Ward-limit closure]
\label{thm:n1v38-Ward-limit}
Assume
\[
 3h_M\leq\frac{\rho_{38}}2.                                \tag{38.16}
\]
Then, uniformly in \(e\in\{{\rm W},*\}\) and \(0\leq t\leq1\),
\[
 \boxed{
 |A_{e,M,t}(0)|
 \leq
 \frac{2B_{38}}{\rho_{38}}h_M,}                             \tag{38.17}
\]
\[
 \boxed{
 |A_{e,M,t}''(0)|
 \leq
 \frac{12B_{38}}{\rho_{38}^2}h_M,}                          \tag{38.18}
\]
\[
 \boxed{
 |A_{e,M,t}''''(0)|
 \leq
 \frac{96B_{38}}{\rho_{38}^3}h_M.}                          \tag{38.19}
\]
Consequently all three normalized formal longitudinal rows vanish in
the momentum-grid limit:
\[
 \boxed{
 \lim_{M\to\infty}A_{e,M,t}^{(j)}(0)=0,
 \qquad j=0,1,2,}                                           \tag{38.20}
\]
uniformly on the positive endpoint interpolation.
\end{theorem}

\begin{proof}
Apply Lemma~\ref{lem:n1v38-three-samples} to (38.12).
Under (38.16), all intervals in (38.13) lie in the Cauchy
neighbourhood.  Insert (38.9):
\[
 \sup|A''|\leq\frac{2B_{38}}{\rho_{38}},
\quad
 \sup|A''''|\leq\frac{8B_{38}}{\rho_{38}^2},
\quad
 \sup|A''''''|\leq\frac{48B_{38}}{\rho_{38}^3}.                \tag{38.21}
\]
Multiplication by \(h_M,3h_M/2,2h_M\) gives
(38.17)--(38.19).  Since \(h_M=2\pi/M\), their right sides tend
to zero.  All constants are uniform in \(e,t\), proving (38.20).
\end{proof}

\begin{corollary}[The toron cell is retained and harmless in the
normalized Ward limit]
\label{cor:n1v38-toron}
No deletion of \(k=0\) is needed in Theorem
\ref{thm:n1v38-Ward-limit}.  Its contribution to (38.8) is included
with weight \(M^{-4}\), and its absolute value is at most
\(B_{38}M^{-4}\).
\end{corollary}

\begin{proof}
The normalized sum (38.8) includes every grid point.
Proposition~\ref{prop:n1v38-strip-bound} bounds the toron summand by
\(B_{38}\); multiplication by its grid weight gives the result.
\end{proof}

\begin{remark}[Why no volume catastrophe occurs]
\label{rem:n1v38-no-volume}
Applying a uniform error bound to an unnormalized sum would create
\(M^4\).  The physical Brillouin average is normalized by \(M^{-4}\)
before the Cauchy and finite-difference estimates are applied.  Thus
the \(M^4\) cells cancel exactly against their normalization and the
remainder is \(O(M^{-1})\).
\end{remark}

\subsection{What has and has not been closed}

\begin{assessment}[Progress after compact V38]
\label{ass:n1v38-status}
The uniform nonzero-\(q\)-to-formal-jet obstruction isolated in V37 is
closed for the positive Gaussian endpoint homotopy.  Three exact
nonzero Ward samples and one common complex strip force the normalized
formal longitudinal response through second external order to vanish
as \(M\to\infty\), with the explicit \(O(M^{-1})\) bounds
(38.17)--(38.19).  The proof retains the toron cell and incurs no
volume loss.
\end{assessment}

\begin{decision}[V39 turns to the even transverse mean]
\label{dec:n1v38-next}
Use the \(1/32\) domain of V37 and the exact V36 evaluator to construct
an interval enclosure of the endpoint-minus-Wilson transverse
second-response mean.  Begin with a tensor-product grid on
\({\cal D}_{1,2}\), use the analytic-strip alias bound of V17--V18,
and refine only boxes whose interval contribution overlaps the target
margin.  Keep the Ward-zero longitudinal channel out of the numerical
budget.
\end{decision}

\begin{firewall}[Claim boundary after compact V38]
\label{fire:n1v38-claim}
V38 closes the averaged longitudinal Ward limit for the declared
positive one-loop Gaussian homotopy.  It does not evaluate the even
transverse Brillouin mean, prove the nonlinear RG comparison or SCIC,
complete ADMC, establish sector tightness or projective continuum
convergence, prove reflection positivity and clustering of a limiting
theory, or prove the Yang--Mills mass gap.
\end{firewall}

\subsection{Parent record}

This V38 note was formed by appending the preceding material to the
validated compact V37, whose record was
\[
\begin{array}{c|c}
\text{lines}&15404\\
\text{bytes}&526097\\
\text{SHA--256}&
\texttt{94001CD678120BD4FCB746FCC1C6AF03DF2032F86A97501C5B274369D0A9F430}.
\end{array}
\]

% =============================================================================
% END APPENDED COMPACT-SERIES V38 MATERIAL
% =============================================================================

% =============================================================================
% BEGIN APPENDED COMPACT-SERIES V39 MATERIAL
% =============================================================================

\section{Compact-series V39: an exact sixteen-corner bubble average
and a diagram-level symmetry obstruction}
\label{sec:n1v39}

\subsection{The \(M=2\) corner characters}

Write a corner momentum as
\[
 k(\varepsilon)_\mu
 =
 \begin{cases}
 0,&\varepsilon_\mu=1,\\
 \pi,&\varepsilon_\mu=-1,
 \end{cases}
 \qquad
 \varepsilon\in\{-1,1\}^4.                                 \tag{39.1}
\]
For an independent-leg phase
\(\boldsymbol\alpha=(\alpha^{(0)},\ldots,\alpha^{(r-1)})\)
and assigned internal momentum rates
\(\eta=(\eta_0,\ldots,\eta_{r-1})\), define
\[
 \chi_\varepsilon(\boldsymbol\alpha;\eta)
 :=
 \prod_{\mu=1}^4
 \varepsilon_\mu^{
 \sum_{\ell=0}^{r-1}\eta_\ell\alpha^{(\ell)}_\mu}.
                                                                    \tag{39.2}
\]
For the cubic bank, \(\eta=(1,-1,0)\); the external phase rates remain
\(\sigma_+=(0,-1,1)\) and \(\sigma_-=(0,1,-1)\).

\begin{lemma}[Exact corner evaluation]
\label{lem:n1v39-corner}
If a cubic phase term has coefficient \(c_{\boldsymbol\alpha}\), then
its stored external jet of order \(j\) at the corner
\(\varepsilon\) is
\[
 c_{\boldsymbol\alpha}
 \chi_\varepsilon(\boldsymbol\alpha;\eta)
 \left(
 \sum_\ell\sigma_{\pm,\ell}
 \alpha^{(\ell)}_1
 \right)^j.                                                  \tag{39.3}
\]
Moreover
\[
 \frac1{16}\sum_{\varepsilon\in\{-1,1\}^4}
 \chi_\varepsilon(\boldsymbol\alpha;\eta)
 =
 \begin{cases}
 1,&
 \sum_\ell\eta_\ell\alpha^{(\ell)}_\mu
 \equiv0\pmod2\quad\text{for every }\mu,\\
 0,&\text{otherwise}.
 \end{cases}                                                 \tag{39.4}
\]
\end{lemma}

\begin{proof}
At (39.1),
\[
 e^{i\sum_\ell\eta_\ell
       \alpha^{(\ell)}\cdot k(\varepsilon)}
 =
 \chi_\varepsilon(\boldsymbol\alpha;\eta).
\]
External differentiation supplies the last factor in (39.3).
For (39.4), the sum factors into four one-dimensional averages
\[
 \frac12\sum_{\varepsilon_\mu=\pm1}
 \varepsilon_\mu^{n_\mu},
\]
which is one for even \(n_\mu\) and zero for odd \(n_\mu\).
\end{proof}

The relation (39.4) is the finite Walsh transform underlying the
corner compiler.  Matrix inversion is performed after evaluation at
each character; averaging a Hessian before inversion would be
incorrect.

\subsection{The exact all-corner cubic calculation}

The executable constructs the translated cubic phase banks once and
then evaluates all sixteen characters.  For every corner it builds
\[
 H_{\rm W}^{[0,1,2]}(k),\qquad
 H_*^{[0,1,2]}(k),\qquad
 G_{\rm W}^{[0,1,2]}(k),\qquad
 G_*^{[0,1,2]}(k),                                         \tag{39.5}
\]
and the Wilson, endpoint-routed Wilson, completion, and endpoint
cubic writers.  All six covariance parity tests and all twelve
cubic plus/minus tests pass at every corner:
\[
 H^{[0]\,T}=H^{[0]},\qquad
 H^{[1]\,T}=-H^{[1]},\qquad
 H^{[2]\,T}=H^{[2]},                                       \tag{39.6}
\]
\[
 P_-^{[j]}=(-1)^jP_+^{[j]},\qquad0\leq j\leq2.             \tag{39.7}
\]
Since every corner is self-inverse, all sixteen Wilson and endpoint
bubble first jets vanish exactly.

Let
\[
 \overline{\cal B}_{e,2}^{[j]}
 :=
 \frac1{16}
 \sum_{\varepsilon\in\{-1,1\}^4}
 {\cal B}_e^{[j]}(k(\varepsilon)).                          \tag{39.8}
\]

\begin{theorem}[Exact \(M=2\) bubble averages]
\label{thm:n1v39-bubble-mean}
For the external axis one and transverse polarization two,
\[
 \boxed{
 \overline{\cal B}_{{\rm W},2}^{[0,1,2]}
 =
 \left(
 -\frac{36585517}{627200},
 0,
 \frac{9266897433589}{265531392000}
 \right),}                                                  \tag{39.9}
\]
\[
 \boxed{
 \overline{\cal B}_{*,2}^{[0,1,2]}
 =
 \left(
 -\frac{6807204165927499225367939369407}
 {97801318615974315871488960000},
 0,
 \frac{35323529630323296663649739916823557926382489801451790252987}
 {36786699279534470965406231814271312468634765095027200000000}
 \right).}                                                  \tag{39.10}
\]
Their difference is
\[
 \boxed{
 \overline{\Delta{\cal B}}_{2}^{[0,1,2]}
 =
 \left(
 -\frac{120564775394563111968303991891}
 {10697019223622190798444105000},
 0,
 -\frac{22294852695032981928047323036879107987860354424973368920873}
 {656905344277401267239396996683416294082763662411200000000}
 \right).}                                                  \tag{39.11}
\]
Numerically,
\[
\begin{split}
 \overline{\cal B}_{{\rm W},2}^{[0,2]}
 &\approx(-58.3315003188776,\ 34.8994420727060),\\
 \overline{\cal B}_{*,2}^{[0,2]}
 &\approx(-69.6023761464464,\ 0.960225579411383),\\
 \overline{\Delta{\cal B}}_2^{[0,2]}
 &\approx(-11.2708758275689,\ -33.9392164932947).
                                                                  \tag{39.12}
\end{split}
\]
\end{theorem}

\begin{proof}
For every character, invert the two exact rational matrices in
(39.5), differentiate the shifted inverse through order two, and
contract
\[
 c_e\Tr[
 G_e(k)P_e^+(k,q)G_e(k+q\widehat e_1)P_e^-(k,q)]
                                                                    \tag{39.13}
\]
in the three-component jet ring.  Sum all sixteen reduced fractions
and divide by sixteen.  Equations (39.6)--(39.7) make every first
component zero before averaging.  Exact fraction reduction gives
(39.9)--(39.11).
\end{proof}

\subsection{Why the bubble cannot be symmetry-compressed}

For the fixed external axis and polarization, the corner action of
\({\cal G}_{1,2}\) only interchanges corner coordinates three and
four; sign reflections act trivially on \(0/\pi\).  There are twelve
corner orbits.  If the bubble were a complete measured scalar, the two
members of each size-two orbit would agree.

\begin{proposition}[Nonzero raw bubble orbit residual]
\label{prop:n1v39-residual}
The endpoint-minus-Wilson bubble does not satisfy that orbit identity.
For example, using bit one for momentum \(\pi\), the difference between
corners \(0010\) and \(0001\) is
\[
 \boxed{
 \Delta{\cal B}^{[0,1,2]}(0010)
 -
 \Delta{\cal B}^{[0,1,2]}(0001)
 =
 \left(
 \frac{3691149668509014}{448054457621855},
 0,
 -\frac{350723778139991308245609543917}
 {762522962486720500234102115000}
 \right).}                                                  \tag{39.14}
\]
Four of the twelve corner orbit classes have a nonzero bubble
residual.
\end{proposition}

\begin{proof}
The compiler evaluates both members independently from their
four-component characters, subtracts the exact endpoint and Wilson
bubbles at each member, and then subtracts the two rational triples.
The zeroth component displayed in (39.14) is positive and therefore
nonzero.  The complete residual table contains four nonempty classes.
\end{proof}

\begin{corollary}[Symmetry reduction must follow measured
recombination]
\label{cor:n1v39-no-early-orbit}
Replacing the sixteen bubble cells by twelve representative bubbles
with orbit multiplicities is invalid.  The orbit formula (37.12) may
be applied only after the quartic tadpole, stationary-coordinate, Haar,
and quotient-coordinate rows have been added.
\end{corollary}

\begin{proof}
If representative replacement were valid for the bubble, every
size-two class would be constant.  Proposition
\ref{prop:n1v39-residual} contradicts that necessary condition.  The
complete measured scalar remains hypercubic by Theorem
\ref{thm:n1v37-covariance}; hence the missing diagrammatic rows must
carry the compensating transformation terms.
\end{proof}

\begin{remark}[The obstruction is useful]
\label{rem:n1v39-useful}
The failed early compression is not a failure of physical symmetry.
It is an exact diagnostic of the non-tensorial parity-tree chart.
It prevents an apparently natural factor-of-\(32\) speedup from
silently changing the transverse mean and fixes the next compiler
order: batch the quartic and measured rows first, test their
recombined orbit residual, and compress only after that residual is
zero.
\end{remark}

\subsection{Executable record and next acquisition}

The verifier is
\[
 \texttt{scripts/verify\_ym\_v39\_corner\_bubbles.py}.       \tag{39.15}
\]
Its source record is
\[
\begin{array}{c|c}
\text{lines}&513\\
\text{bytes}&16407\\
\text{SHA--256}&
\texttt{28FB906D2156584A5EB43DEFDAFB1775C5DD7CED04902F0B2EC447799DABC472}.
\end{array}                                                  \tag{39.16}
\]
The canonical payload containing all sixteen corners, the twelve
orbit classes, the exact averages, and all nonzero residuals has
SHA--256
\[
 \boxed{
 \texttt{630A118C75AAF5B97940A8D04B59A5C331F42ABBEAE85D0F6FF30FFF9EA307F2}.}
                                                                    \tag{39.17}
\]

\begin{assessment}[Progress after compact V39]
\label{ass:n1v39-status}
The entire \(M=2\) cubic bubble grid is now exact.  All covariance,
cubic, and odd-jet gates pass, and the two endpoint bubble means are
known as rational triples.  More importantly, the computation proves
that a raw bubble is not a hypercubic scalar in parity-tree
coordinates.  The V37 orbit reduction is therefore reserved for the
fully recombined response.
\end{assessment}

\begin{decision}[Mandatory V40 audit and corner tadpole batch]
\label{dec:n1v39-next}
Perform the required SM/NS companion audit.  Then extend the Walsh
corner representation to the translated quartic, stationary, and Haar
rows in one shared word pass.  Add them to the sixteen bubble cards,
verify that every residual in (39.14) cancels in the complete measured
response, and only then form the twelve-representative \(M=2\) total
mean.
\end{decision}

\begin{firewall}[Claim boundary after compact V39]
\label{fire:n1v39-claim}
V39 evaluates only the cubic bubble part of the \(M=2\) transverse
mean.  It does not evaluate the complete corner tadpole and measured
rows, certify the full orbit cancellation, control the \(M=2\)
aliasing error, enclose the continuum Brillouin mean, complete SCIC or
ADMC, establish the continuum measure and sector estimates, prove
clustering, or prove the Yang--Mills mass gap.
\end{firewall}

\subsection{Parent record}

This V39 note was formed by appending the preceding material to the
validated compact V38, whose record was
\[
\begin{array}{c|c}
\text{lines}&15737\\
\text{bytes}&536904\\
\text{SHA--256}&
\texttt{ED5ACB149AF22C720DD516913B5A8AD3BFA15FE710C106F93FFE5E4DC3238390}.
\end{array}
\]

% =============================================================================
% END APPENDED COMPACT-SERIES V39 MATERIAL
% =============================================================================

% =============================================================================
% BEGIN APPENDED COMPACT-SERIES V40 MATERIAL
% =============================================================================

\section{Compact-series V40: compulsory audit, derivative erratum,
and the exact Walsh algebra for measured corner batching}
\label{sec:n1v40}

\subsection{The V40 companion audit}

The current companion files at the completion of this audit are
\[
\begin{array}{c|r|r|l}
\text{programme and file}&\text{lines}&\text{bytes}&
\text{SHA--256}\\ \hline
\text{SM--Einstein V76}
&23607&750857&
\texttt{7C4EC49F08171B0039BC67B5E063C4E99624FE1C52BA1A3D0EC66E9354ED2833}\\
\text{Navier--Stokes V26}
&5319&278283&
\texttt{82C887F8C2C69E4F28FAD7C3DC8DE5B9099CB5A4BF431EBA1FFF3E91935978AD}.
\end{array}                                                    \tag{40.1}
\]
The audit changes the present programme in four precise ways.
\begin{enumerate}
\item SM V67--V68 independently proves orbit--stabilizer weighting and
      finite Ward telescopes, while warning that shifted Ward terms
      must cancel on the full grid before orbit reduction.
\item SM V70 detects an apostrophe-doubling defect in the literal V38
      derivative notation.  Its corrected three-zero lemma has the
      intended derivative orders and exactly the constants used in
      (38.17)--(38.19).
\item SM V46 proves exact normalized-product-Haar disintegration under
      spanning-tree gauge fixing.  This confirms that the quotient
      density row, rather than a raw tree-coordinate diagram, is the
      correct carrier of physical hypercubic symmetry.
\item The new SM V76 gives an abstract energy--entropy criterion for
      exponential topological-sector tails.  It supplies a later
      template for Yang--Mills sector tightness, but no uniform
      Yang--Mills action-cost and sector-entropy constants, so it does
      not close that row here.
\end{enumerate}
The NS file is unchanged and supplies no new gauge, Ward, or corner
compiler input.

\subsection{Formal erratum for inherited apostrophes}

The append transport used for V35--V39 doubled every literal
apostrophe in newly appended text.  Most instances are typographical,
but derivative primes in V36 and V38 were changed mathematically.
The following displays supersede the affected inherited displays; no
historical source is rewritten.

\begin{correction}[Correct inverse derivatives in V36]
\label{err:n1v40-V36}
Equation (36.25) is to be read as
\[
 \boxed{
 G'=-GH'G,\qquad
 G''=2GH'GH'G-GH''G.}                                      \tag{40.2}
\]
\end{correction}

\begin{correction}[Correct three-zero lemma in V38]
\label{err:n1v40-V38-lemma}
In Lemma~\ref{lem:n1v38-three-samples}, equations
(38.13)--(38.15) are to be read with derivative orders one, two,
and three:
\[
\boxed{
\begin{split}
 |A(0)|
 &\leq h\sup_{[0,h]}|A'|,\\
 |A'(0)|
 &\leq\frac32h\sup_{[0,2h]}|A''|,\\
 |A''(0)|
 &\leq2h\sup_{[0,3h]}|A'''|.
\end{split}}                                                \tag{40.3}
\]
The two forward-difference identities are
\[
 0=\frac1h\int_h^{2h}A'(s)\,ds,                             \tag{40.4}
\]
\[
 0=\int_0^1\int_0^1
 A''(h(1+u+v))\,du\,dv.                                    \tag{40.5}
\]
Accordingly the Cauchy substitutions in (38.21) are
\[
 \boxed{
 \sup|A'|\leq\frac{2B_{38}}{\rho_{38}},\qquad
 \sup|A''|\leq\frac{8B_{38}}{\rho_{38}^2},\qquad
 \sup|A'''|\leq\frac{48B_{38}}{\rho_{38}^3}.}               \tag{40.6}
\]
\end{correction}

\begin{correction}[Correct Ward-limit conclusions]
\label{err:n1v40-V38-theorem}
The left sides of (38.17)--(38.19) are respectively
\[
 |A_{e,M,t}(0)|,\qquad
 |A_{e,M,t}'(0)|,\qquad
 |A_{e,M,t}''(0)|.                                         \tag{40.7}
\]
With (40.3)--(40.6), their constants and the conclusion (38.20)
are correct exactly as intended.
\end{correction}

The dummy variable in (37.16) is \(k'\), not \(k''\).
The phrases \emph{Markov's inequality}, \emph{Taylor's theorem}, and
\emph{Cauchy's estimate} have their ordinary spelling.  These
typographical repairs alter no formula other than the prime counts
explicitly superseded in (40.2)--(40.7).

\begin{proof}[Verification of the corrected constants]
The first inequality in (40.3) is the fundamental theorem of calculus.
The average of \(s\) on \([h,2h]\) is \(3h/2\), giving the second.
The integral representation (40.5) and the average
\(\int_0^1\int_0^1(1+u+v)\,du\,dv=2\) give the third.  Cauchy's
estimate on a circle of radius \(\rho_{38}/2\) gives
\[
 \sup|A^{(r)}|
 \leq r!(2/\rho_{38})^rB_{38},\qquad r=1,2,3.
\]
Multiplying by \(h,3h/2,2h\) gives the constants
\(2,12,96\) in (38.17)--(38.19).
\end{proof}

\begin{record}[Patch-transport discipline]
\label{rec:n1v40-transport}
Starting with V40, literal apostrophes are passed unchanged through
the patch transport.  A successor validation must search new appended
material for prime runs longer than the declared derivative order.
\end{record}

\subsection{The group-algebra jet ring}

Let
\[
 H:=({\mathbb Z}/2{\mathbb Z})^4,\qquad
 {\cal J}_2:=\mathbb Q^3,                                   \tag{40.8}
\]
and equip \({\cal J}_2\) with the exact jet product
\[
 (a_0,a_1,a_2)\star(b_0,b_1,b_2)
 :=
 (a_0b_0,\,
  a_1b_0+a_0b_1,\,
  a_2b_0-2a_1b_1+a_0b_2).                                  \tag{40.9}
\]
An \(H\)-label is a four-bit internal phase mask.  Define
\[
 {\cal A}_{40}:=\{F:H\to{\cal J}_2\}                        \tag{40.10}
\]
with XOR convolution
\[
 (F*G)(\omega)
 :=
 \sum_{\alpha\oplus\beta=\omega}
 F(\alpha)\star G(\beta).                                  \tag{40.11}
\]
For \(\varepsilon\in\{-1,1\}^4\), write
\[
 \chi_\varepsilon(\omega)
 :=
 \prod_{\mu=1}^4\varepsilon_\mu^{\omega_\mu},
 \qquad
 {\cal E}_\varepsilon F
 :=
 \sum_{\omega\in H}
 \chi_\varepsilon(\omega)F(\omega).                         \tag{40.12}
\]

\begin{theorem}[Walsh evaluation is an exact jet-algebra
isomorphism]
\label{thm:n1v40-Walsh}
For every \(F,G\in{\cal A}_{40}\),
\[
 \boxed{
 {\cal E}_\varepsilon(F*G)
 =
 ({\cal E}_\varepsilon F)\star
 ({\cal E}_\varepsilon G).}                                \tag{40.13}
\]
Moreover
\[
 \boxed{
 F(\omega)
 =
 \frac1{16}\sum_{\varepsilon\in\{-1,1\}^4}
 \chi_\varepsilon(\omega){\cal E}_\varepsilon F.}           \tag{40.14}
\]
Thus one coefficient object in \({\cal A}_{40}\) is equivalent,
without approximation, to all sixteen corner jets.
\end{theorem}

\begin{proof}
Characters turn XOR into multiplication:
\[
 \chi_\varepsilon(\alpha\oplus\beta)
 =
 \chi_\varepsilon(\alpha)\chi_\varepsilon(\beta).
\]
Insert (40.11), rearrange the finite double sum, and use bilinearity
of \(\star\) to obtain (40.13).  Character orthogonality gives
\[
 \frac1{16}\sum_\varepsilon
 \chi_\varepsilon(\omega)\chi_\varepsilon(\alpha)
 =
 {\bf1}_{\omega=\alpha},
\]
which proves (40.14).
\end{proof}

\subsection{Insertion words in the Walsh algebra}

For differentiated leg \(\ell\) at displacement
\(a\in{\mathbb Z}^4\), let \(\eta_\ell\) be its internal momentum
rate and \(\sigma_\ell\) its external momentum rate.  Its internal
mask and external jet are
\[
 \omega_{\ell,a}
 :=
 \eta_\ell a\pmod2,\qquad
 j_{\ell,a}
 :=
 (1,\sigma_\ell a_1,-(\sigma_\ell a_1)^2).                 \tag{40.15}
\]
Place the insertion coefficient at mask \(\omega_{\ell,a}\) with
value \(j_{\ell,a}\).  Multiplication of noncommutative word factors
uses (40.11) for their scalar phase jets while preserving the declared
colour-word order.

\begin{proposition}[One-pass translated quartic batching]
\label{prop:n1v40-quartic-batch}
Take the quartic rates
\[
 \eta=(1,-1,0,0),\qquad
 \sigma=(0,0,1,-1).                                        \tag{40.16}
\]
Replace every scalar jet in the translated quartic compiler by its
\({\cal A}_{40}\)-valued insertion (40.15), and replace scalar
multiplication by (40.11).  Then:
\begin{enumerate}
\item the factor order is exactly the corrected
      \(e^Be^Z\) order of V31;
\item Hilbert--Schmidt conjugation reverses colour words and conjugates
      the jet but leaves the internal mask unchanged;
\item applying \({\cal E}_\varepsilon\) after the complete word sum
      gives exactly the quartic jet obtained by running the V36 direct
      corner evaluator at \(k(\varepsilon)\);
\item all sixteen quartic matrices are therefore obtained from one
      owner-word pass.
\end{enumerate}
\end{proposition}

\begin{proof}
The first statement concerns only noncommutative colour ordering and is
unchanged by the scalar coefficient ring.  Internal Fourier exponents
add under multiplication; modulo two this is XOR, giving (40.11).
Negating an exponent during adjunction does not change its class in
\({\mathbb Z}/2{\mathbb Z}\), proving the second statement.  External
jet multiplication is exactly (40.9).  The homomorphism (40.13) then
commutes with every finite product and sum in the owner compiler.
At a final character, (40.12) is precisely the corner phase
(39.2).  This proves the third and fourth statements.
\end{proof}

\subsection{The exact measured recombination target}

Let
\[
 {\cal R}_e
 =
 {\cal B}_e+{\cal T}_e+{\cal S}_e+{\cal H}_e               \tag{40.17}
\]
denote, respectively, the cubic bubble, quartic tadpole,
stationary-coordinate, and Haar/quotient-density rows, all valued in
\({\cal A}_{40}\) before character evaluation.  For the swap
\(\tau:(\varepsilon_3,\varepsilon_4)\mapsto
(\varepsilon_4,\varepsilon_3)\), define
\[
 {\mathfrak r}_e(\varepsilon)
 :=
 {\cal E}_\varepsilon{\cal R}_e
 -
 {\cal E}_{\tau\varepsilon}{\cal R}_e.                     \tag{40.18}
\]

\begin{theorem}[Exact acceptance criterion for corner compression]
\label{thm:n1v40-acceptance}
A V40 corner compiler may replace sixteen cells by the twelve
fixed-channel representatives if and only if
\[
 \boxed{
 {\mathfrak r}_{\rm W}(\varepsilon)
 =
 {\mathfrak r}_*(\varepsilon)=0
 \quad\text{for all }\varepsilon.}                          \tag{40.19}
\]
For the example pair in V39, the missing measured rows must satisfy
\[
\begin{split}
 &\bigl[
 \Delta{\cal T}+\Delta{\cal S}+\Delta{\cal H}
 \bigr](0010)
 -
 \bigl[
 \Delta{\cal T}+\Delta{\cal S}+\Delta{\cal H}
 \bigr](0001)\\
 &\qquad=
 -
 \left(
 \frac{3691149668509014}{448054457621855},
 0,
 -\frac{350723778139991308245609543917}
 {762522962486720500234102115000}
 \right).                                                   \tag{40.20}
\end{split}
\]
\end{theorem}

\begin{proof}
The twelve-representative formula is equivalent to constancy on every
two-member orbit, which is exactly (40.19).  Complete measured
hypercubic covariance is Theorem~\ref{thm:n1v37-covariance}.
Subtract the V39 bubble residual (39.14) from the complete zero
residual to obtain (40.20).
\end{proof}

\begin{warning}[Haar disintegration does not make the raw bubble
covariant]
\label{warn:n1v40-Haar}
The SM V46 disintegration theorem identifies the correct quotient
measure.  It does not say that a single Feynman diagram transforms as
a quotient scalar.  The cancellation in (40.20) is an executable
requirement, not a consequence that permits the missing rows to be
omitted.
\end{warning}

\begin{assessment}[Progress after compact V40]
\label{ass:n1v40-status}
The compulsory audit has repaired the literal derivative notation and
confirmed the scope of the V37--V38 Ward result.  The exact
\(({\mathbb Z}_2)^4\) group-algebra jet ring now converts the missing
sixteen quartic corner evaluations into one translated word pass.
Equation (40.20) fixes a nonzero rational compensation target that
the fully measured compiler must hit before any orbit reduction.
\end{assessment}

\begin{decision}[V41 executes the measured Walsh batch]
\label{dec:n1v40-next}
Implement Proposition~\ref{prop:n1v40-quartic-batch} for the Wilson and
completion quartic banks, add the stationary and Haar/quotient rows,
and evaluate all sixteen characters.  Require the toron and
\(\pi\widehat e_2\) regressions, the four V39 compensation residuals,
and all twelve complete orbit identities before reporting the
\(M=2\) total response mean.
\end{decision}

\begin{firewall}[Claim boundary after compact V40]
\label{fire:n1v40-claim}
V40 provides the corrected Ward-limit notation and the exact batching
algebra, but it does not yet execute the measured quartic corner bank,
prove the compensation residuals, control internal-momentum aliasing,
enclose the continuum transverse mean, complete SCIC or ADMC, acquire
Yang--Mills sector energy--entropy constants, construct the continuum
measure and clustering theory, or prove the Yang--Mills mass gap.
\end{firewall}

\subsection{Parent record}

This V40 note was formed by appending the preceding material to the
validated compact V39, whose record was
\[
\begin{array}{c|c}
\text{lines}&16060\\
\text{bytes}&547520\\
\text{SHA--256}&
\texttt{9D5FFFC1E5C4D32650F5DC5B789D213830A2CBC48AA99D953A8A7E4E9CC6D458}.
\end{array}
\]

% =============================================================================
% END APPENDED COMPACT-SERIES V40 MATERIAL
% =============================================================================

% =============================================================================
% BEGIN APPENDED COMPACT-SERIES V41 MATERIAL
% =============================================================================

\section{Compact-series V41: exact failure of the premature corner
compression and the missing re-tree intertwiner}
\label{sec:n1v41}

\subsection{Purpose and outcome}

V40 reduced the proposed \(M=2\) calculation to one
\(({\mathbb Z}/2{\mathbb Z})^4\)-valued quartic word pass.  The pass has
now been executed with exact rational arithmetic.  It successfully
reproduces both pre-existing one-point quartic regressions, but the
first proposed bubble--tadpole orbit cancellation is false.  Thus the
twelve-representative compression must not be used for the presently
implemented fixed-tree response.

This is useful negative information.  It isolates a distinction that
was hidden in V37--V40: invariance of product Haar measure under
factorwise background translation does not by itself supply the
background-dependent intertwiner between two different tree gauges.
The latter map, or an independent off-axis vertex check, is required
before complete measured covariance can be invoked.

\subsection{The exact failed identity}

Let
\[
 \epsilon=(+,+,+,-),\qquad
 \widetilde\epsilon=(+,+,-,+),                              \tag{41.1}
\]
so their corner strings are respectively \(0001\) and \(0010\).
They are interchanged by the spectator transposition
\(\tau:(3\,4)\).  For \(j=0,1,2\), define the endpoint-minus-Wilson
residuals
\[
\begin{split}
 b_j&:=
 \Delta{\cal B}^{[j]}(\widetilde\epsilon)
 -\Delta{\cal B}^{[j]}(\epsilon),\\
 t_j&:=
 \Delta{\cal T}^{[j]}(\widetilde\epsilon)
 -\Delta{\cal T}^{[j]}(\epsilon),                           \tag{41.2}\\
 r_j&:=b_j+t_j .
\end{split}
\]
Here \({\cal B}\) is the V39 cubic bubble and \({\cal T}\) is the V41
translated quartic tadpole.  The V39 value is
\[
 b=
 \left(
 \frac{3691149668509014}{448054457621855},
 0,
 -\frac{350723778139991308245609543917}
 {762522962486720500234102115000}
 \right),                                                   \tag{41.3}
\]
whereas the new Walsh pass gives
\[
 t=
 \left(
 -\frac{4136577840}{505450939},
 0,
 \frac{2512236281825279}{243121901659000}
 \right).                                                   \tag{41.4}
\]
Exact reduction therefore gives
\[
 \boxed{
 r=
 \left(
 \frac{24300925130214}{448054457621855},
 0,
 \frac{3764303437279103000762561781949}
 {381261481243360250117051057500}
 \right)\ne0.}                                              \tag{41.5}
\]
For scale only,
\[
 r_0=0.0542365436094451\ldots,\qquad
 r_2=9.87328545491418\ldots .                               \tag{41.6}
\]

\begin{theorem}[The implemented tadpole--bubble pair is not a corner
scalar]
\label{thm:n1v41-nonscalar}
For the \(M=2\), external-axis \(1\), polarization-axis \(2\) card,
the presently implemented endpoint-minus-Wilson quantity
\[
             \Delta({\cal B}+{\cal T})^{[j]}(k)             \tag{41.7}
\]
is not constant on the \({\cal G}_{1,2}\)-orbit containing
\(\epsilon\) and \(\widetilde\epsilon\).  In particular, it cannot be
summed by the twelve-representative formula (37.12).
\end{theorem}

\begin{proof}
The transposition \(\tau\) belongs to \({\cal G}_{1,2}\) and maps the
two corners in (41.1) to one another.  Their difference for (41.7) is
the vector (41.5).  Its zeroth and second components are nonzero
because their displayed integer numerators are positive.  Hence
orbit constancy fails.  The finite-grid orbit formula requires orbit
constancy, so it is unavailable for this quantity.
\end{proof}

\begin{proposition}[What the long exact run did certify]
\label{prop:n1v41-certified}
Before reaching the false cancellation assertion, the V41 executable
certified all of the following with rational arithmetic:
\begin{enumerate}
\item the Wilson, endpoint-routed Wilson, and completion stationary
      scores vanish through the required external-momentum order;
\item all sixteen corner tadpole cards were constructed and every
      first tadpole jet vanished;
\item at \(k=0\), the Wilson and endpoint tadpoles are
      \[
      \left(\frac{487}{8},0,-\frac{17009}{384}\right),
      \quad
      \left(
      \frac{5132669}{72520},0,
      -\frac{60057901669949}{388686480000}
      \right);                                              \tag{41.8}
      \]
\item at \(k=\pi\widehat e_2\), they are
      \[
      \left(\frac{302}{5},0,-\frac{7457}{168}\right),
      \quad
      \left(
      \frac{38665469995}{505450939},0,
      -\frac{28758791690948779}{182341426244250}
      \right);                                              \tag{41.9}
      \]
\item the two pairs in (41.8)--(41.9) agree exactly with the independent
      V33 and V36 direct evaluators.
\end{enumerate}
\end{proposition}

\begin{proof}
The executable forms every corner tadpole before entering its residual
loop.  It then tests the two regression dictionaries (41.8)--(41.9).
Only after those comparisons pass does it form (41.3)--(41.5) and
raise its first assertion at the pair \(0001,0010\).  Thus successful
control flow to that assertion proves each preceding exact test.
\end{proof}

\subsection{Three coordinate operations that must not be conflated}

There are three different maps in the present construction.
\begin{enumerate}
\item The background recentering on each reduced group factor is
      \[
                       z_e\longmapsto g_e(B,q)z_e.           \tag{41.10}
      \]
      This is factorwise left translation and has unit product-Haar
      Jacobian, as proved in V33.
\item The stationary recentering is
      \(x\mapsto x-X_e(B,q)\).  The exhaustive same-colour score test
      proves that its action contribution vanishes through order two
      for the cards considered here.
\item If a lattice symmetry \(R\) sends the fixed ordered tree \(T\)
      to a different tree \(R(T)\), returning to \(T\) requires a
      re-treeing map
      \[
      \Phi_{R,B,q}:
      {\cal X}_{T}\longrightarrow{\cal X}_{T}.              \tag{41.11}
      \]
      It is obtained by embedding reduced variables into all links,
      applying \(R\), and gauge-fixing again to \(T\).
\end{enumerate}
The first operation is not (41.11): the gauge multiplier in a
re-treeing depends on the link variables already traversed along the
new tree.  Consequently Theorem~\ref{thm:n1v33-Haar} does not, by
itself, prove that the density and boundary-channel terms induced by
(41.11) vanish.

\begin{correction}[Scope of the V37 covariance theorem]
\label{corr:n1v41-covariance-scope}
The coordinate-free covariance statement (37.5) remains valid for a
genuinely complete measured scalar provided that
\begin{enumerate}
\item the external boundary channel is transported by the full
      re-treeing map rather than identified only by its coordinate
      label; and
\item the local density attached to (41.11) is included, or is proved
      to have unit response in the chosen chart.
\end{enumerate}
Neither condition was executed in V39--V41.  Thus the specialization
of (37.5) to the computed fixed-channel cards, and consequently the
use of (37.12) for those cards, was premature.
\end{correction}

\subsection{The coordinate anomaly formula}

The missing issue can be stated without diagrammatic ambiguity.
Let \(x\) and \(y\) be two local quotient coordinates related near a
common stationary point by
\[
 y=\Phi_q(x),\qquad J(q)=D\Phi_q(0),                        \tag{41.12}
\]
and let \(K_x(q)\), \(K_y(q)\) be the corresponding Hessian matrices.
Let \(\rho_x(0,q)\,dx\) and \(\rho_y(0,q)\,dy\) denote the local
quotient densities.

\begin{theorem}[Exact Hessian--density cancellation under re-treeing]
\label{thm:n1v41-coordinate-cancellation}
If \(J(q)\) is invertible, then
\[
\begin{split}
 K_y(q)&=J(q)^{-T}K_x(q)J(q)^{-1},\\
 \rho_y(0,q)&=\rho_x(0,q)|\det J(q)|^{-1},                  \tag{41.13}
\end{split}
\]
and the measured one-loop scalar
\[
 \Gamma(q):=\frac12\log\det K(q)-\log\rho(0,q)             \tag{41.14}
\]
is independent of the coordinate choice.
\end{theorem}

\begin{proof}
At a stationary point the second-derivative chain rule has no term
containing the first derivative of the action.  Applying it to
\(x=\Phi_q^{-1}(y)\) gives the first identity in (41.13).
The ordinary change-of-variables formula gives the second.  Therefore
\[
\begin{split}
 \frac12\log\det K_y
 &=
 \frac12\log\det K_x-\log|\det J|,\\
 -\log\rho_y
 &=
 -\log\rho_x+\log|\det J|.
\end{split}
\]
Adding the two equalities proves invariance of (41.14).
\end{proof}

Let \({\cal C}\) denote the still-uncomputed combination of transported
boundary-channel and quotient-density terms.  If the V39 bubble and
V41 tadpole compilers are correct away from their regression axes,
Theorem~\ref{thm:n1v41-coordinate-cancellation} forces the exact target
\[
 {\cal C}^{[j]}(\widetilde\epsilon)
 -{\cal C}^{[j]}(\epsilon)=-r_j,\qquad j=0,1,2.             \tag{41.15}
\]
In particular, (41.5), rather than the larger bubble residual (41.3),
is now the residual that a completed coordinate transport must cancel.

\begin{theorem}[The exact V41 dichotomy]
\label{thm:n1v41-dichotomy}
At least one of the following statements holds:
\begin{enumerate}
\item the fixed coordinate label used for the external polarization
      is not carried to itself by the full spectator re-treeing, and
      the transported-channel term is nonzero;
\item the re-treeing density has a nonzero response and supplies part
      or all of (41.15);
\item the Walsh quartic or V39 cubic compiler has an off-axis defect
      not detected by the toron and \(\pi\widehat e_2\) regressions;
\item the claimed pointwise corner covariance is stronger than the
      covariance of the actual block observable, which may close only
      after its full alias sum.
\end{enumerate}
\end{theorem}

\begin{proof}
If none held, then the two vertex compilers would give the complete
measured scalar in a channel fixed by \(\tau\), and that scalar would
transform pointwise under (37.5).  Its residual between the two
corners would be zero, contradicting (41.5).
\end{proof}

The theorem is a logical disjunction, not permission to choose a
preferred explanation.  The next computation must distinguish its
branches.

\subsection{Upstream acquisition replacing the failed compression}

\begin{decision}[V42 constructs the spectator re-tree certificate]
\label{dec:n1v41-next}
Before another sixteen-corner mean is attempted, construct the
transposition \(\tau=(3\,4)\) on the actual finite block:
\begin{enumerate}
\item embed the \(45\) reduced vertical variables and the four retained
      boundary variables into the full \(64\)-link linearized field;
\item permute link bases and directions by \(\tau\);
\item apply the unique rooted gauge transformation that restores the
      original ordered tree;
\item extract the vertical block \(S_\tau(k)\), the induced boundary
      channel map, and their phase monomials;
\item prove or disprove exact equality of the transported
      polarization-\(2\) channel;
\item compute the determinant/density jet dictated by
      (41.13);
\item independently run the direct V36 quartic evaluator at
      \(k=\pi\widehat e_3\) and \(k=\pi\widehat e_4\) and compare those
      matrices and traces with the corresponding Walsh evaluations.
\end{enumerate}
Only after these tests may (41.15) be classified as a coordinate
compensation, a compiler defect, or an alias-level rather than
pointwise symmetry.
\end{decision}

\begin{record}[Executable obstruction certificate]
\label{rec:n1v41-executable}
The exact source is
\[
 \texttt{scripts/verify\_ym\_v41\_walsh\_tadpoles.py}.       \tag{41.16}
\]
Its record at the executed failure is
\[
\begin{array}{c|c}
\text{lines}&905\\
\text{bytes}&28048\\
\text{SHA--256}&
\texttt{5FD556BD9DD15CD8256BEFE1819E5BE6462DBF6947B11F65206992AE484D56C6}.
\end{array}                                                  \tag{41.17}
\]
The process exits at its first compensation assertion with the exact
fractions (41.3)--(41.5).  A nonzero exit is the intended certificate
for the falsified acceptance target; it is not recorded as a
successful all-orbit payload.
\end{record}

\begin{assessment}[Progress after compact V41]
\label{ass:n1v41-status}
The quartic Walsh construction has passed two independent axial
regressions and has produced the first exact spectator-pair tadpole
residual.  Its sum with the V39 bubble residual is nonzero, so the
premature pointwise corner compression has been stopped before it
could contaminate a Brillouin-zone mean.  The remaining finite
bottleneck is now explicit: construct the re-tree/channel intertwiner
or locate an off-axis vertex defect.
\end{assessment}

\begin{firewall}[Claim boundary after compact V41]
\label{fire:n1v41-claim}
V41 proves the nonzero residual (41.5) and invalidates the current
twelve-corner reduction.  It does not prove which branch of
Theorem~\ref{thm:n1v41-dichotomy} occurs, does not provide the complete
\(M=2\) mean, does not control aliasing or continuum quadrature, does
not close SCIC or ADMC, does not construct the Yang--Mills continuum
measure or exponential clustering, and does not prove the
Yang--Mills mass gap.
\end{firewall}

\subsection{Parent record}

This V41 note was formed by appending the preceding material to the
validated compact V40, whose record was
\[
\begin{array}{c|c}
\text{lines}&16447\\
\text{bytes}&560619\\
\text{SHA--256}&
\texttt{F0B16C4DCCF204359A62342A87322FAE793567A22E3EF63AC813FEB4DF51951B}.
\end{array}
\]

% =============================================================================
% END APPENDED COMPACT-SERIES V41 MATERIAL
% =============================================================================

% =============================================================================
% BEGIN APPENDED COMPACT-SERIES V42 MATERIAL
% =============================================================================

\section{Compact-series V42: the retained-channel ordering correction
and exact spectator re-tree certificate}
\label{sec:n1v42}

\subsection{The indexing defect}

The V41 obstruction has a simpler upstream cause than a missing
Jacobian.  The list of four retained boundary columns was constructed
by filtering the lexicographically ordered list
\[
 \bigl\{(s,\mu):s\in\{0,1\}^4,\ 0\leq\mu\leq3\bigr\}       \tag{42.1}
\]
with the condition \(s=\widehat e_\mu\).  It was then indexed as if
its list position were the direction \(\mu\).  Those two orderings
are opposite.

\begin{theorem}[Exact retained-slot order]
\label{thm:n1v42-slot-order}
In zero-based notation the retained list is
\[
 \boxed{
 \bigl(
 (\widehat e_3,3),(\widehat e_2,2),
 (\widehat e_1,1),(\widehat e_0,0)
 \bigr).}                                                   \tag{42.2}
\]
Thus retained slots \(0,1,2,3\) represent physical axes
\(3,2,1,0\), respectively.  In one-based notation their physical
axis order is \(4,3,2,1\).
\end{theorem}

\begin{proof}
The parity enumeration in the executable has the nested order
\[
 (a,b,c,d),\qquad
 a=0,1;\ b=0,1;\ c=0,1;\ d=0,1,                            \tag{42.3}
\]
with \(d\) changing fastest.  Among the four unit vectors the resulting
lexicographic order is
\[
 (0,0,0,1)<(0,0,1,0)<(0,1,0,0)<(1,0,0,0).
\]
For each such vector the filter retains only the edge whose direction
is the location of its nonzero coordinate.  This is exactly (42.2).
\end{proof}

\begin{correction}[Physical polarization of the computed cards]
\label{corr:n1v42-polarization}
Every executable in the V28--V41 response chain which sets
\[
                       \texttt{polarization = 1}            \tag{42.4}
\]
selects retained slot \(1\), hence physical polarization axis \(3\)
in one-based notation, not axis \(2\).  The external momentum axis
remains physical axis \(1\), because its index is applied directly to
the Laurent exponent rather than to the retained-column list.

In particular, the V36 nonzero card has
\[
 k=\pi\widehat e_2,\qquad q\parallel\widehat e_1,\qquad
 B\parallel\widehat e_3.                                   \tag{42.5}
\]
It remains a valid transverse card, and none of its rational values
changes.  Only the polarization label is corrected.
\end{correction}

\subsection{Why the V39--V41 orbit comparison was invalid}

Let \(\tau_{34}\) interchange physical axes \(3\) and \(4\).  In the
retained-slot order (42.2), its boundary matrix is
\[
 P_{34}=
 \begin{pmatrix}
 0&1&0&0\\
 1&0&0&0\\
 0&0&1&0\\
 0&0&0&1
 \end{pmatrix}.                                             \tag{42.6}
\]
Thus it sends the actually used slot \(1\) to slot \(0\).

\begin{theorem}[Resolution of the V41 residual]
\label{thm:n1v42-residual-resolution}
The corners \(0001\) and \(0010\) in (41.1) are related by
\(\tau_{34}\), but the V39 and V41 programs evaluate retained slot
\(1\) at both corners.  Covariance relates slot \(1\) at the first
corner to slot \(0\) at the second.  Therefore the nonzero rational
vector (41.5) is not an orbit residual of one covariant scalar and
does not require a Haar, stationary, or quotient-density
compensation.
\end{theorem}

\begin{proof}
By (42.2), slot \(1\) is physical axis \(3\), while slot \(0\) is
physical axis \(4\).  Equation (42.6) is the exact boundary action of
\(\tau_{34}\).  The covariance identity has the form
\[
 {\mathscr W}^{[j]}_{e;1,4}(\tau_{34}k)
 =
 {\mathscr W}^{[j]}_{e;1,3}(k),                             \tag{42.7}
\]
whereas (41.2) subtracts two values with physical polarization \(3\).
They are different comparisons.  Hence no zero conclusion applies to
(41.5).
\end{proof}

\begin{correction}[Supersession of the V41 interpretation]
\label{corr:n1v42-V41}
The arithmetic identities (41.3)--(41.6) and the regression facts in
Proposition~\ref{prop:n1v41-certified} remain exact.  The conclusion
of Theorem~\ref{thm:n1v41-nonscalar} is superseded: the calculated
fixed-slot function was compared under a transformation outside its
channel stabilizer.  In Theorem~\ref{thm:n1v41-dichotomy}, the first
branch occurs, for the elementary boundary permutation (42.6).
There is presently no evidence from (41.5) for a missing density row
or an off-axis vertex defect.
\end{correction}

\subsection{Exact construction of the re-tree map}

Let \(T\), \(R\), and \(C\) be respectively the \(15\) tree edges,
the \(45\) reduced vertical edges, and the four retained boundary
edges.  At a corner \(\epsilon\), let
\[
 d_T(\epsilon):
 {\mathbb Q}^{15}\longrightarrow{\mathbb Q}^{15}           \tag{42.8}
\]
be the tree-edge restriction of the linear gauge gradient with the
root gauge parameter deleted.  A full gauge-fixed coordinate is
\(u=(x,b)\in{\mathbb Q}^{45}\oplus{\mathbb Q}^4\).

After applying \(\tau_{34}\) to all \(64\) link components, let
\(a_T\) be the resulting tree-edge vector.  The unique restoring
gauge parameter is
\[
                  \phi=-d_T(\tau_{34}\epsilon)^{-1}a_T.    \tag{42.9}
\]
Adding the full gradient \(d\phi\) and restricting back to \(R\cup C\)
defines a \(49\)-by-\(49\) matrix
\[
 {\cal U}_{34}(\epsilon):
 {\mathbb Q}^{45}\oplus{\mathbb Q}^4
 \longrightarrow
 {\mathbb Q}^{45}\oplus{\mathbb Q}^4.                      \tag{42.10}
\]

\begin{theorem}[Exact spectator re-tree block structure]
\label{thm:n1v42-retree}
At every one of the sixteen corners,
\[
 \boxed{
 {\cal U}_{34}(\epsilon)=
 \begin{pmatrix}
 S_{34}(\epsilon)&0\\
 0&P_{34}
 \end{pmatrix},}                                            \tag{42.11}
\]
where \(S_{34}(\epsilon)\) has rank \(45\), \(P_{34}\) is (42.6),
and
\[
 {\cal U}_{34}(\tau_{34}\epsilon)
 {\cal U}_{34}(\epsilon)=I_{49}.                            \tag{42.12}
\]
For both the Wilson and endpoint free vertical Hessians,
\[
 \boxed{
 A_e(\epsilon)
 =
 S_{34}(\epsilon)^T
 A_e(\tau_{34}\epsilon)
 S_{34}(\epsilon).}                                        \tag{42.13}
\]
\end{theorem}

\begin{proof}
The matrix (42.10) is constructed directly from the integer incidence
grammar.  The \(15\)-by-\(15\) tree gradient is inverted exactly.
For each corner the program verifies that the restored tree
components vanish, then extracts all four blocks.  Both off-diagonal
blocks reduce identically to zero; the lower-right block reduces to
(42.6).  Exact Gauss elimination gives rank \(45\) for the upper-left
block and rank \(49\) for the full map.  Direct multiplication proves
(42.12).

Independently, the program evaluates the \(96\)-row plaquette
incidence and the \(384\)-row completion incidence at each corner.
It forms
\[
 A_{\rm W}=D_R^TD_R,\qquad
 A_*=D_R^TD_R+3K_R^TK_R,                                   \tag{42.14}
\]
and verifies (42.13) entry by entry for both matrices at all sixteen
corners.
\end{proof}

The zero off-diagonal blocks show that the re-tree map is not the
source of the V41 discrepancy.  They also show that the physical
axis permutation on the retained variables cannot be discarded.

\begin{theorem}[Covariance of all routed toron channels]
\label{thm:n1v42-router-covariance}
Let \(L_{e,p}(0)\in{\mathbb Q}^{49}\) be the complete zero-momentum
linear Schur lift of retained slot \(p\), including its unit retained
component.  If \(p'\) is the slot representing
\(\tau_{34}\) of the physical axis represented by \(p\), then
\[
 \boxed{
 {\cal U}_{34}(0)L_{e,p}(0)=L_{e,p'}(0)}                    \tag{42.15}
\]
for every \(p=0,1,2,3\) and for \(e={\rm W},*\).
\end{theorem}

\begin{proof}
The router is computed independently from the Wilson and endpoint
Gram matrices.  The verifier constructs all eight vectors in
(42.15), multiplies by the exact toron matrix (42.10), and reduces
each of the eight \(49\)-component residuals.  Every residual is the
zero vector.  In slot notation the nontrivial cases are
\[
 0\longleftrightarrow1,                                    \tag{42.16}
\]
while slots \(2\) and \(3\) are fixed, exactly as (42.6) requires.
\end{proof}

\subsection{The corrected stabilizer and fundamental domain}

The computed channel is now denoted
\[
 q\parallel\widehat e_1,\qquad B\parallel\widehat e_3.      \tag{42.17}
\]
Its spectator axes are \(2\) and \(4\).  Hence the admissible
transposition is \(\tau_{24}\), not \(\tau_{34}\).

\begin{correction}[Correct fixed-channel hypercubic group]
\label{corr:n1v42-group}
Replace \({\cal G}_{1,2}\) in the application of (37.7)--(37.12) to
the computed cards by
\[
 {\cal G}_{1,3}
 :=
 \{R:\ R\widehat e_1=\pm\widehat e_1,\ 
        R\widehat e_3=\pm\widehat e_3,\ 
        R\{\widehat e_2,\widehat e_4\}
        =\{\pm\widehat e_2,\pm\widehat e_4\}\}.             \tag{42.18}
\]
The corrected continuum fundamental domain may be chosen as
\[
 {\cal D}_{1,3}
 :=
 \{k\in[0,\pi]^4:\ 0\leq k_4\leq k_2\leq\pi\}.             \tag{42.19}
\]
On a corner bit string, the nontrivial spectator operation swaps bit
positions \(2\) and \(4\) in one-based notation.
\end{correction}

\begin{proof}
The external and background physical axes in (42.17) must be fixed up
to sign.  The remaining two axes may be independently signed and
interchanged, giving (42.18).  Sign folding puts all components in
\([0,\pi]\), and the single remaining transposition orders \(k_2\)
and \(k_4\), giving (42.19).  The corner-bit statement is the
restriction of that transposition to \(\{0,\pi\}^4\).
\end{proof}

At \(M=2\), the four nontrivial corrected pairs are
\[
 \boxed{
 (0001,0100),\quad
 (0011,0110),\quad
 (1001,1100),\quad
 (1011,1110).}                                              \tag{42.20}
\]
These replace the four pairs used in V39--V41.

\subsection{Executable certificate and next calculation}

\begin{record}[V42 exact certificate]
\label{rec:n1v42-executable}
The verifier is
\[
 \texttt{scripts/verify\_ym\_v42\_retree\_channel.py}.      \tag{42.21}
\]
Its source record is
\[
\begin{array}{c|c}
\text{lines}&544\\
\text{bytes}&17372\\
\text{SHA--256}&
\texttt{DFD68F6965B8BE967641927C783B3E43464ECE74E15D1E7AE226413F336D8D88}.
\end{array}                                                  \tag{42.22}
\]
The canonical payload SHA--256 is
\[
 \boxed{
 \texttt{09E9D8E1BC1E17FE2D791A2A384749FCB72FBDD390001E3210DAB7B86C8B44F2}.}
                                                                    \tag{42.23}
\]
\end{record}

\begin{decision}[V43 reruns the correct fixed-channel corner test]
\label{dec:n1v42-next}
Retain external exponent axis \(1\) and retained slot \(1\), now
correctly named physical polarization axis \(3\).  Recombine the V39
bubble and V41 Walsh tadpole at all sixteen corners, but test orbit
constancy only on the corrected pairs (42.20).  Require independent
slot-transport covariance checks under \(\tau_{24}\), then report the
uncompressed sixteen-cell mean and its correctly compressed
twelve-representative replay.
\end{decision}

\begin{assessment}[Progress after compact V42]
\label{ass:n1v42-status}
The first branch of the V41 dichotomy has been resolved exactly.  The
fixed tree has a valid vertical intertwiner and both actions obey
Hessian congruence, but the response chain mislabeled the retained
slot.  Correcting the slot order changes the stabilizer, the
fundamental domain, and the four \(M=2\) orbit pairs while preserving
all previously calculated rational response values as transverse
cards.
\end{assessment}

\begin{firewall}[Claim boundary after compact V42]
\label{fire:n1v42-claim}
V42 repairs the channel label and proves the exact spectator re-tree
certificate at the sixteen corners.  It does not yet execute the
corrected V43 tadpole--bubble orbit test, report a valid \(M=2\)
complete mean, establish general-momentum channel covariance, control
aliasing or continuum quadrature, close SCIC or ADMC, construct the
continuum measure and clustering theory, or prove the Yang--Mills
mass gap.
\end{firewall}

\subsection{Parent record}

This V42 note was formed by appending the preceding material to the
validated compact V41, whose record was
\[
\begin{array}{c|c}
\text{lines}&16811\\
\text{bytes}&574386\\
\text{SHA--256}&
\texttt{191ACA59453A84990AABF333FA79E534A760B1B1A6F22AFA326B5D9E6AE1502A}.
\end{array}
\]

% =============================================================================
% END APPENDED COMPACT-SERIES V42 MATERIAL
% =============================================================================

% =============================================================================
% BEGIN APPENDED COMPACT-SERIES V43 MATERIAL
% =============================================================================

\section{Compact-series V43: corrected corner pairing cancels at
zeroth order but not at the second momentum jet}
\label{sec:n1v43}

\subsection{The corrected exact replay}

V42 identified the physical response channel as
\[
 q\parallel\widehat e_1,\qquad
 B\parallel\widehat e_3,                                   \tag{43.1}
\]
and replaced the incorrect spectator transposition
\(\tau_{34}\) by \(\tau_{24}\).  The first corrected corner pair is
\[
 k=0001,\qquad \tau_{24}k=0100.                            \tag{43.2}
\]
The V39 bubble calculation has now been replayed from its original
word and matrix grammar rather than merely relabeled.

For the endpoint-minus-Wilson bubble, member minus reference is
\[
 b^{\rm corr}=
 \left(
 \frac{27034282524}{2527254695},
 0,
 -\frac{3318420258814361393872661}
 {215050838598762613651750}
 \right).                                                   \tag{43.3}
\]
The independently rebuilt V43 Walsh tadpole gives
\[
 t^{\rm corr}=
 \left(
 -\frac{27034282524}{2527254695},
 0,
 \frac{302626538072896}{30390237707375}
 \right).                                                   \tag{43.4}
\]
Their sum is
\[
 \boxed{
 r^{\rm corr}=
 b^{\rm corr}+t^{\rm corr}
 =
 \left(
 0,0,
 -\frac{489372186744556308849}
 {89418228107593602350}
 \right).}                                                  \tag{43.5}
\]
Numerically, only for scale,
\[
 b^{\rm corr}_0=10.6970945894316\ldots,\quad
 b^{\rm corr}_2=-15.4308640711967\ldots,\quad
 t^{\rm corr}_2=9.95801813025818\ldots,\quad
 r^{\rm corr}_2=-5.47284594093850\ldots .                   \tag{43.6}
\]

\begin{theorem}[Exact separation of the slot and momentum-jet
defects]
\label{thm:n1v43-separation}
For the corrected pair (43.2), the implemented tadpole--bubble
response is orbit invariant at external-momentum order zero but is
not orbit invariant at order two:
\[
 \boxed{r^{\rm corr}_0=r^{\rm corr}_1=0,\qquad
 r^{\rm corr}_2\ne0.}                                      \tag{43.7}
\]
Consequently the V42 slot correction completely explains the
zeroth-order V41 residual, but it does not by itself justify the
second-jet corner compression.
\end{theorem}

\begin{proof}
The zeroth entries in (43.3) and (43.4) are exact negatives and both
first entries vanish.  Their sum is therefore zero in orders zero and
one.  The second component in (43.5) is nonzero because its numerator
and denominator are positive integers.  This proves (43.7).
\end{proof}

The exact zeroth cancellation is a nontrivial check of V42: the new
pair differs from the old V41 pair, and its two large rational
contributions cancel without any tolerance.  The surviving second
jet narrows the remaining problem to the \(q\)-dependence of the
coordinate/channel transport or to the second-jet vertex grammar.

\subsection{All four corrected bubble targets}

The successful V39 replay gives the following four
member-minus-reference bubble residuals:
\[
\begin{array}{c|c|c}
\text{pair}&b_0&b_2\\ \hline
0100-0001&
\dfrac{27034282524}{2527254695}&
-\dfrac{3318420258814361393872661}
 {215050838598762613651750}\\[6pt]
0110-0011&
\dfrac{622029267}{109259150}&
-\dfrac{3795197757029845693}{286501484609340000}\\[6pt]
1100-1001&
\dfrac{11753847}{2401300}&
-\dfrac{108173852698046438299441}
 {19080998874982044000000}\\[6pt]
1110-1011&
\dfrac{78366}{31217}&
-\dfrac{15007045615805627279}
 {1656509574141006000}.
\end{array}                                                  \tag{43.8}
\]
Every first component is zero.  These are the exact compensation
targets for future tadpole plus coordinate rows.  The V43 Walsh run
stopped at the first nonzero total (43.5), so no claim is made yet
about the other three complete residuals.

\begin{proposition}[Independent V39 replay integrity]
\label{prop:n1v43-bubble-replay}
The corrected-pair wrapper reproduces the original V39 canonical
payload hash
\[
 \texttt{630A118C75AAF5B97940A8D04B59A5C331F42ABBEAE85D0F6FF30FFF9EA307F2}
                                                               \tag{43.9}
\]
before extracting (43.8).  Hence the new residuals use exactly the
same sixteen bubble cards and bubble averages as V39.
\end{proposition}

\begin{proof}
The wrapper runs the V39 main routine in memory, parses its complete
JSON result, and checks its returned payload record.  It then
subtracts the stored rational triples at the corrected pairs (42.20).
No bubble expression is copied or reimplemented in the wrapper.
\end{proof}

\subsection{Why the V42 certificate does not settle (43.5)}

The V42 map \({\cal U}_{24}(k)\) at zero background controls the
linear vertical Hessian and the zeroth external background lift.
The response in (43.5), however, contains two background-amplitude
derivatives and two external-momentum derivatives.  The relevant
change of variables is therefore
\[
 y=\Phi_{\tau,B,q}(x),\qquad
 J_{\tau}(B,q):=D_x\Phi_{\tau,B,q}(0),                      \tag{43.10}
\]
not merely \(J_\tau(0,0)\).

Let
\[
 \gamma_\tau(B,q)
 :=
 \log|\det J_\tau(B,q)|
 +\log\rho_y(0;B,q)-\log\rho_x(0;B,q).                     \tag{43.11}
\]
For an exactly transported quotient density, (41.13) says
\(\gamma_\tau=0\).  Diagrammatically, however, the cancellation may
be distributed between the Hessian tadpole--bubble terms and the
mixed derivatives of the two density terms in (43.11).

\begin{proposition}[Exact missing-row target at the first corrected
pair]
\label{prop:n1v43-target}
Assume the V39 cubic and V43 quartic second-jet compilers are the
correct Hessian contributions in the two fixed-tree charts.  Then the
sum \({\cal C}\) of all background-dependent re-tree and local-density
rows must obey
\[
 \boxed{
 {\cal C}^{[2]}(0100)-{\cal C}^{[2]}(0001)
 =
 \frac{489372186744556308849}
 {89418228107593602350}.}                                  \tag{43.12}
\]
\end{proposition}

\begin{proof}
A complete measured scalar in the channel (43.1) is invariant under
\(\tau_{24}\).  Its corner residual is the sum of (43.5) and the
residual of \({\cal C}\).  Setting that sum to zero and transposing
the term in (43.5) gives (43.12).
\end{proof}

\begin{warning}[The target is conditional]
\label{warn:n1v43-conditional}
Equation (43.12) is not evidence that a nonzero density row actually
exists.  An error in either second-jet vertex compiler would change
the target.  The next acquisition must independently transport the
second jet before interpreting (43.12) as a Jacobian coefficient.
\end{warning}

\subsection{Executable records}

The corrected bubble replay is
\[
 \texttt{scripts/verify\_ym\_v43\_corrected\_bubble\_orbits.py},   \tag{43.13}
\]
with
\[
\begin{array}{c|c}
\text{lines}&65\\
\text{bytes}&1903\\
\text{SHA--256}&
\texttt{115E7BCB84093CB543640FC6EE7F87650597571952D1DF27855AB997D0D00FF3},
\end{array}                                                   \tag{43.14}
\]
and canonical payload SHA--256
\[
 \texttt{0B070021D0B372CEFA5132E5998FFC8E888F87B476D9BC5060EC4F2BE459FEC5}.
                                                               \tag{43.15}
\]

The corrected Walsh run is
\[
 \texttt{scripts/verify\_ym\_v43\_walsh\_tadpoles.py},       \tag{43.16}
\]
with
\[
\begin{array}{c|c}
\text{lines}&917\\
\text{bytes}&28445\\
\text{SHA--256}&
\texttt{A9DE4BB719204682670C35A3BDC2A338C2FABF73587FE352AC2B56D1777ABFD8}.
\end{array}                                                   \tag{43.17}
\]
It exits at the first corrected compensation assertion with the exact
data (43.3)--(43.5).  It is therefore a falsification certificate,
not a successful all-orbit payload.

\begin{decision}[V44 differentiates the physical re-tree map]
\label{dec:n1v43-next}
Construct \({\cal U}_{24}(p)\) for the correct spectator
transposition as a Laurent matrix in the fluctuation momentum \(p\).
Differentiate it through order two along
\(p=k+q\widehat e_1\), verify the product rule and inverse identities,
and transport all four Wilson and endpoint Schur lifts through order
two.  Then extend the same construction to the background-dependent
linearization \(J_{24}(B,q)\).  This separates a channel-transport
term from a genuine vertex error before the expensive Walsh bank is
run again.
\end{decision}

\begin{assessment}[Progress after compact V43]
\label{ass:n1v43-status}
The corrected physical orbit makes the zeroth tadpole--bubble
residual vanish exactly, validating the principal V42 correction.
The second jet remains nonzero by the explicit fraction (43.5).
The unresolved finite-dimensional issue has therefore been reduced
from a general corner-covariance failure to a mixed
background/external-momentum transport calculation.
\end{assessment}

\begin{firewall}[Claim boundary after compact V43]
\label{fire:n1v43-claim}
V43 proves (43.5) and supplies all four corrected bubble targets.  It
does not yet determine whether the second-jet residual is a
re-tree-density term or a vertex defect, certify the remaining three
complete corner pairs, produce a symmetry-compressed complete
\(M=2\) mean, control the continuum quadrature or alias limit, close
SCIC or ADMC, construct the continuum measure and clustering theory,
or prove the Yang--Mills mass gap.
\end{firewall}

\subsection{Parent record}

This V43 note was formed by appending the preceding material to the
validated compact V42, whose record was
\[
\begin{array}{c|c}
\text{lines}&17164\\
\text{bytes}&586899\\
\text{SHA--256}&
\texttt{008480F2DD4049D17ABE7DB62E49C99E3E497022BF7DA022D2856362181605DC}.
\end{array}
\]

% =============================================================================
% END APPENDED COMPACT-SERIES V43 MATERIAL
% =============================================================================

% =============================================================================
% BEGIN APPENDED COMPACT-SERIES V44 MATERIAL
% =============================================================================

\section{Compact-series V44: exact \(q\)-jet re-tree covariance and
elimination of the linear channel explanation}
\label{sec:n1v44}

\subsection{Jet convention for the correct spectator map}

Let \(\tau=\tau_{24}\), so in zero-based coordinate notation
\[
                         \tau=(0,3,2,1).                    \tag{44.1}
\]
For a corner \(k\), the vertical part of the restored-tree map at
momentum \(k+q\widehat e_1\) is expanded as
\[
 S_\tau(k,q)
 =
 S_0(k)+iqS_1(k)+\frac{q^2}{2}S_2(k)+O(q^3).               \tag{44.2}
\]
For two matrices with this convention, the exact product jets are
\[
\begin{split}
 (AB)_0&=A_0B_0,\\
 (AB)_1&=A_1B_0+A_0B_1,\\
 (AB)_2&=A_2B_0-2A_1B_1+A_0B_2.                           \tag{44.3}
\end{split}
\]
The minus sign is the factor \(i^2=-1\); omitting it would corrupt
precisely the second-jet test now at issue.

\subsection{Construction from the gauge gradient}

At a target corner \(\tau k\), a boundary-crossing gradient entry in
direction \(\mu\) is
\[
 \epsilon_{\tau k,\mu}
 \exp(iq{\bf1}_{\mu=1}),                                   \tag{44.4}
\]
where physical directions are written one-based in the indicator.
Its three rational jet coefficients are
\[
 \epsilon_{\tau k,\mu}
 \left(1,\ {\bf1}_{\mu=1},\
       -{\bf1}_{\mu=1}\right).                              \tag{44.5}
\]
Every tree edge lies inside one parity cell, so the restricted
\(15\)-by-\(15\) tree gradient is independent of \(q\).  Consequently
the restoring gauge parameter is obtained once from its exact
rational inverse, while all \(q\)-dependence enters when the full
gradient is added back.

\begin{theorem}[The correct re-tree map through second order]
\label{thm:n1v44-retree-jets}
For every \(k\in\{0,\pi\}^4\), the complete \(49\)-coordinate
re-tree jet has block form
\[
 {\cal U}_\tau(k,q)=
 \begin{pmatrix}
 S_\tau(k,q)&0\\
 0&P_\tau
 \end{pmatrix}
 +O(q^3),                                                   \tag{44.6}
\]
where \(P_\tau\) is the constant retained-axis permutation
\[
 P_\tau:
 (0,1,2,3)_{\rm slot}\longmapsto(2,1,0,3)_{\rm slot}.      \tag{44.7}
\]
It obeys the exact jet involution
\[
 \boxed{
 {\cal U}_\tau(\tau k,q)
 {\cal U}_\tau(k,q)
 =
 I_{49}+O(q^3).}                                            \tag{44.8}
\]
Both off-diagonal blocks vanish coefficientwise in orders zero, one,
and two.  In particular retained slot \(1\), the physical
polarization axis \(3\), is fixed.
\end{theorem}

\begin{proof}
Insert (44.5) in the \(64\)-by-\(15\) rooted gauge gradient.  Permute
the full link basis by (44.1), solve the constant tree equations, add
the restoring gradient, and restrict to the \(45+4\) non-tree
coordinates.  Exact reduction gives zero for all three coefficients
of both off-diagonal blocks and gives (44.7) in order zero with zero
higher jets.

Run the same construction at \(\tau k\) and multiply the two
three-component matrix jets using (44.3).  At each of the sixteen
corners, the zeroth product is \(I_{49}\) and the first and second
products are zero matrices.  This proves (44.8).
\end{proof}

\subsection{The determinant jet is trivial}

Put \(X=S_0^{-1}S_1\).  In the convention (44.2), the coefficients of
the logarithmic determinant are
\[
\begin{split}
 \ell_1(k)&=\operatorname{tr}(S_0^{-1}S_1),\\
 \ell_2(k)&=
 \operatorname{tr}\left(
 S_0^{-1}S_2+(S_0^{-1}S_1)^2
 \right).                                                   \tag{44.9}
\end{split}
\]
The plus sign in the second line again comes from \(i^2=-1\).

\begin{theorem}[No pure-\(q\) re-tree determinant response]
\label{thm:n1v44-logdet-zero}
At all sixteen corners,
\[
 \boxed{\ell_1(k)=\ell_2(k)=0.}                             \tag{44.10}
\]
\end{theorem}

\begin{proof}
The exact program inverts each rational \(S_0(k)\), forms the two
traces in (44.9), and reduces them as fractions.  All thirty-two
fractions are zero.  This is stronger than an absolute-determinant
argument because it tests the formal complex phase jet itself.
\end{proof}

Thus a density term depending only on the linear map
\(S_\tau(k,q)\) cannot supply the nonzero target (43.12).

\subsection{Routed background covariance through \(q^2\)}

For action \(e\in\{{\rm W},*\}\) and retained slot \(p\), write the
full \(49\)-coordinate Schur lift as
\[
 L_{e,p}(q)=
 L_{e,p}^{[0]}+iqL_{e,p}^{[1]}
 +\frac{q^2}{2}L_{e,p}^{[2]}+O(q^3).                       \tag{44.11}
\]
Let \(\tau p\) be the slot transformation (44.7).

\begin{theorem}[All routed channels intertwine through order two]
\label{thm:n1v44-router-jets}
For both actions and all four retained slots,
\[
 \boxed{
 {\cal U}_\tau(0,q)L_{e,p}(q)
 =
 L_{e,\tau p}(q)+O(q^3).}                                  \tag{44.12}
\]
All \(8\cdot3\cdot49\) rational residual entries vanish.
\end{theorem}

\begin{proof}
The Wilson router is the exact jet solution of
\[
 A_{\rm W}(q)L_{\rm W}(q)+B_{\rm W}(q)=0,                  \tag{44.13}
\]
and the endpoint router is the corresponding solution for
\(A_{\rm W}+3A_1\) and its mixed block.  The verifier independently
constructs all four columns of both solutions.  It then applies
(44.3) to the re-tree jet and each column jet.  The target columns are
chosen using (44.7).  Exact entrywise comparison gives zero in every
case, proving (44.12).
\end{proof}

For the actual slot \(p=1\), (44.12) is a same-slot identity.  Hence
neither an omitted permutation of the boundary label nor a defect in
the linear Schur lift explains the V43 second residual.

\begin{corollary}[Sharpened V43 dichotomy]
\label{cor:n1v44-dichotomy}
Assume the exact linear incidence grammar and the Schur equations.
Then the residual (43.5) can arise only from one or both of:
\begin{enumerate}
\item mixed background-amplitude derivatives of the nonlinear
      group-valued re-tree map or of its local Haar coordinate
      density; or
\item an error in the cubic-bubble or translated-quartic
      second-external-momentum compiler.
\end{enumerate}
It cannot arise from the retained-slot ordering, the pure-\(q\)
linear re-tree map, its log determinant through order two, or the
linear routed background jets.
\end{corollary}

\begin{proof}
The slot ordering was repaired in V42.  Theorem
\ref{thm:n1v44-retree-jets} removes linear fibre/boundary mixing,
Theorem~\ref{thm:n1v44-logdet-zero} removes the pure-\(q\)
determinant row, and Theorem~\ref{thm:n1v44-router-jets} removes the
linear channel lift.  The only untested dependence of
\(J_\tau(B,q)\) is therefore its background-amplitude dependence.
If that also contributes zero, the nonzero exact residual forces a
second-jet compiler defect.
\end{proof}

\subsection{Executable record and next upstream test}

\begin{record}[V44 exact certificate]
\label{rec:n1v44-executable}
The verifier is
\[
 \texttt{scripts/verify\_ym\_v44\_retree\_jets.py}.         \tag{44.14}
\]
Its source record is
\[
\begin{array}{c|c}
\text{lines}&467\\
\text{bytes}&14489\\
\text{SHA--256}&
\texttt{6AB1EB4168A08090E8C624449F8F89A11586994C5862C614A6380E75228B4AFB}.
\end{array}                                                  \tag{44.15}
\]
The successful canonical payload SHA--256 is
\[
 \boxed{
 \texttt{39010AB6F12E53549CE0FF1CBA15DC34B30C78140229351A2CFAAEE1B0511EB8}.}
                                                                    \tag{44.16}
\]
\end{record}

\begin{decision}[Mandatory V45 audit, then nonlinear re-tree
linearization]
\label{dec:n1v44-next}
First perform the five-version SM/NS audit required by the programme.
Then construct the group-valued \(\tau_{24}\) re-tree map around the
same-colour background \(B\).  Differentiate its vertical Jacobian
twice in background amplitude and twice in \(q\), retaining the local
Haar density.  In parallel with that derivation, make the cubic and
quartic compilers pass a direct finite-difference or dual-number
Hessian regression at the pair \(0001,0100\).  Compare the resulting
exact coefficient with (43.12).
\end{decision}

\begin{assessment}[Progress after compact V44]
\label{ass:n1v44-status}
The correct spectator re-tree and every Schur background channel now
intertwine through the full external-momentum order used in the
response.  The vertical determinant jet is exactly trivial.  This
eliminates all linear coordinate explanations of (43.5) and leaves a
precise nonlinear-Jacobian-versus-vertex-compiler fork.
\end{assessment}

\begin{firewall}[Claim boundary after compact V44]
\label{fire:n1v44-claim}
V44 proves linear re-tree and router covariance through \(q^2\), but
does not compute the background-amplitude dependence of the nonlinear
re-tree Jacobian, identify the second-jet defect, certify all corrected
corner pairs, produce a complete \(M=2\) mean, control continuum
quadrature or aliasing, close SCIC or ADMC, construct continuum
clustering, or prove the Yang--Mills mass gap.
\end{firewall}

\subsection{Parent record}

This V44 note was formed by appending the preceding material to the
validated compact V43, whose record was
\[
\begin{array}{c|c}
\text{lines}&17447\\
\text{bytes}&596742\\
\text{SHA--256}&
\texttt{171781B2A86114F0D3B15ACC26550BC40CFCA163A87DF67A8632DD770416052A}.
\end{array}
\]

% =============================================================================
% END APPENDED COMPACT-SERIES V44 MATERIAL
% =============================================================================

% =============================================================================
% BEGIN APPENDED COMPACT-SERIES V45 MATERIAL
% =============================================================================

\section{Compact-series V45: compulsory companion audit and exact
fibre-Haar removal of the nonlinear Jacobian branch}
\label{sec:n1v45}

\subsection{Compulsory five-version companion audit}

Immediately before this audit, the current companion files were
\[
\begin{array}{c|r|r|l}
\text{programme}&\text{lines}&\text{bytes}&\text{SHA--256}\\ \hline
\text{SM--Einstein V9}&2876&100149&
\texttt{92C8065FBECFA30CABE9C10B989ADB05B7426CAE928DEB74ED4C7BE1F0FF2471}\\
\text{Navier--Stokes V26}&5319&278283&
\texttt{82C887F8C2C69E4F28FAD7C3DC8DE5B9099CB5A4BF431EBA1FFF3E91935978AD}.
\end{array}                                                    \tag{45.1}
\]

SM V5 independently imported the YM V42 slot correction.  SM V6
then proved that local extraction must commute with the complete
physical symmetry and channel representation.  SM V7--V8 separate
harmonic, exact, coexact, cochain-map, homotopy, and projector
defects.  SM V9 applies the same principle to anomaly cancellation:
species contributions must first be transported by one common map;
a transport inconsistency cannot be renamed a physical anomaly or
hidden in a remainder.

The transferable rule for the present YM calculation is:
\[
 \boxed{\text{transport first, compare second, and keep a failed
 identity in its own consistency channel.}}                 \tag{45.2}
\]
Accordingly the nonzero V43 coefficient is not assigned to a
Jacobian counterterm unless the exact measure transport produces it.

NS V26 sharpens rank-one Oseen-kernel norms and a conditional
regularity region.  Its tensor maximization is mathematically
independent of the finite compact-group re-tree problem, so no NS
estimate is imported.

\begin{decision}[V45 audit decision]
\label{dec:n1v45-audit}
Import the SM separation of transport and consistency defects as a
proof discipline, but no SM analytic estimate.  Import no NS result.
Use exact tree-gauge disintegration to decide the Jacobian branch of
Corollary~\ref{cor:n1v44-dichotomy}.
\end{decision}

\subsection{Exact product-Haar tree disintegration}

Let \(\Gamma=(V,E)\) be a finite connected directed graph, fix a root
\(o\), and let \(T\subset E\) be a spanning tree.  Let \(G=SU(3)\)
with normalized Haar probability \(dm\).  The rooted gauge group is
\[
                       {\cal G}_o=G^{V\setminus\{o\}}.       \tag{45.3}
\]
The tree slice \(\Sigma_T\) consists of link configurations with
\(U_e=1\) for \(e\in T\).  Its remaining coordinates form
\(G^{E\setminus T}\).

\begin{theorem}[Normalized product-Haar tree disintegration]
\label{thm:n1v45-tree-Haar}
The map
\[
 \Theta_T:
 {\cal G}_o\times G^{E\setminus T}\longrightarrow G^E,
 \qquad
 (h,z)\longmapsto h\cdot s_T(z),                            \tag{45.4}
\]
is a bijection modulo the harmless orientation convention, and
\[
 \boxed{
 \Theta_T^*(dm^{E})=dm^{{V\setminus\{o\}}}\,dm^{E\setminus T}.}
                                                                    \tag{45.5}
\]
In particular, every tree slice carries exactly normalized product
Haar measure on its non-tree links, with no Faddeev--Popov density.
\end{theorem}

\begin{proof}
Starting at the root, traverse the tree.  On an edge joining a vertex
whose gauge value is known to a new vertex, the equation that makes
the transformed tree link equal to the identity determines the new
gauge value uniquely.  Induction over the tree determines every
element of \({\cal G}_o\), proving bijectivity.

At each induction step the old link variable is replaced by a product
of a previously determined group element, the new gauge element, and
possibly an inverse.  Left Haar invariance, right Haar invariance, and
invariance under inversion show that the Jacobian is one in normalized
Haar measure.  Applying Fubini along the tree gives (45.5).
\end{proof}

Let \(R\) be a graph automorphism fixing the coarse root lattice and
let \(\Phi_{R,T}\) be the non-tree coordinate transition obtained by
applying \(R\) and restoring the original tree \(T\).

\begin{corollary}[Tree changes preserve quotient Haar measure]
\label{cor:n1v45-retree-Haar}
\[
 \boxed{
 (\Phi_{R,T})_*dm^{E\setminus T}=dm^{E\setminus T}.}         \tag{45.6}
\]
\end{corollary}

\begin{proof}
The edge permutation induced by \(R\) preserves \(dm^E\).  Apply
Theorem~\ref{thm:n1v45-tree-Haar} before and after the permutation
and integrate out the normalized rooted gauge factor.  The two
remaining quotient measures are equal, giving (45.6).
\end{proof}

\subsection{The retained coarse holonomies}

For direction \(\mu\), the retained boundary variable in the parity
block is the link beginning at \(\widehat e_\mu\) in direction
\(\mu\).  In tree gauge, the first link of the straight two-edge path
from the coarse root to \(2\widehat e_\mu\) is a tree link and equals
the identity.  Therefore the retained variable is exactly the
straight coarse holonomy
\[
 b_\mu
 =
 U(0,\mu)U(\widehat e_\mu,\mu)
 =
 U(\widehat e_\mu,\mu).                                    \tag{45.7}
\]

\begin{theorem}[Exact nonlinear boundary transport]
\label{thm:n1v45-boundary-exact}
For the physical spectator permutation \(\tau_{24}\), the full
group-valued re-tree map has the form
\[
 \boxed{
 \Phi_{\tau,T}(x,b)=
 \bigl(\phi_{\tau,b}(x),P_\tau b\bigr),}                    \tag{45.8}
\]
where \(P_\tau\) is the retained-axis permutation (44.7).  The
boundary output is independent of the \(45\) vertical variables.
\end{theorem}

\begin{proof}
A coordinate permutation sends the straight path in (45.7) to the
straight path in the permuted direction and fixes its coarse
endpoints.  Restoring tree gauge again makes the first edge of that
target path the identity.  Its second edge is therefore the permuted
coarse holonomy.  This proves \(b'=P_\tau b\) as an exact group
identity, independent of \(x\).  The remaining non-tree links define
\(\phi_{\tau,b}(x)\), giving (45.8).
\end{proof}

This nonlinear result upgrades the zero lower-left jets found in
V42 and V44; it is not merely a tangent-space statement.

\subsection{Conditional fibre Haar invariance}

Write
\[
 dm^{E\setminus T}=dm_R(x)\,dm_C(b),\qquad
 R=45,\quad C=4.                                            \tag{45.9}
\]

\begin{theorem}[Every fixed-boundary re-tree map preserves vertical
product Haar]
\label{thm:n1v45-fibre-Haar}
For every retained boundary value \(b\),
\[
 \boxed{
 (\phi_{\tau,b})_*dm_R=dm_R.}                               \tag{45.10}
\]
\end{theorem}

\begin{proof}
Combine (45.6), (45.8), and the fact that \(P_\tau\) preserves
product Haar on \(G^4\).  For continuous test functions \(f\) and
\(g\),
\[
\begin{split}
 &\int_{G^4}g(P_\tau b)
   \left[\int_{G^{45}}f(\phi_{\tau,b}(x))\,dm_R(x)\right]dm_C(b)\\
 &\qquad=
 \int_{G^4}g(b)\,dm_C(b)
 \int_{G^{45}}f(x)\,dm_R(x).                               \tag{45.11}
\end{split}
\]
Changing \(b\) by \(P_\tau\) and varying \(g\) shows that the bracketed
inner integral equals \(\int f\,dm_R\) for almost every \(b\).
The re-tree map is a finite composition of group multiplication and
inversion, so the inner integral is continuous in \(b\).  Equality
almost everywhere between continuous functions on compact \(G^4\)
implies equality everywhere.  Varying \(f\) proves (45.10).
\end{proof}

The moving-background chart used by the response calculation is a
product of left translations
\[
                  z_e\longmapsto g_e(B,q)z_e,               \tag{45.12}
\]
which also preserves \(dm_R\) exactly.  Composing (45.12) with
\(\phi_{\tau,b}\) therefore preserves the full conditional measure
for every background amplitude and external momentum.

\begin{corollary}[No measured nonlinear re-tree density response]
\label{cor:n1v45-no-density}
The complete vertical measure has zero re-tree response in every
mixed background and momentum order.  In particular,
\[
 \boxed{
 \partial_B^2\partial_q^2
 \left[
 \log|\det D\phi_{\tau,b}|
 +\log\rho_{\rm target}
 -\log\rho_{\rm source}
 \right]_{B=q=0}=0.}                                       \tag{45.13}
\]
Individual Lebesgue-coordinate Jacobians may be nonconstant, but
their sum with the local Haar-density change is identically zero.
\end{corollary}

\begin{proof}
Equation (45.10) is equality of the measures, not only of their total
mass.  Write it in any smooth local coordinates and apply the
ordinary change-of-variables formula.  The bracket in (45.13)
vanishes identically wherever the charts overlap.  Differentiation
gives (45.13).  Background translations preserve Haar separately by
(45.12).
\end{proof}

\subsection{The V43 residual is a compiler consistency defect}

\begin{theorem}[Elimination of the Jacobian branch]
\label{thm:n1v45-compiler-defect}
At least one of the V39 cubic-bubble or V43
translated-quartic second-\(q\)-jet implementations fails to equal
the corresponding derivative of the complete fixed-boundary Hessian
determinant.  Equivalently, the nonzero fraction (43.5) is a compiler
consistency defect, not a physical symmetry anomaly and not an
uncomputed measured-Haar counterterm.
\end{theorem}

\begin{proof}
The Wilson and endpoint actions are invariant under \(\tau_{24}\).
The physical channel is fixed by V42, its linear routed lift is
intertwined through \(q^2\) by V44, and the exact conditional measure
is intertwined by Corollary~\ref{cor:n1v45-no-density}.  Hence the
mixed derivative
\[
 \partial_B^2\partial_q^2
 \left[
 \frac12\log\det H_e(B,q)-\log\rho_e(B,q)
 \right]_{B=q=0}                                           \tag{45.14}
\]
has zero residual between the paired corners.  The implemented sum
has the nonzero residual (43.5).  If both diagram compilers equalled
the two Hessian derivatives whose sum is (45.14), their residual
would vanish.  Contradiction.  Equation (45.13) excludes assigning
the difference to the measured density.
\end{proof}

This conclusion follows the SM V7--V9 discipline: a failure of the
map that carries an exact identity is recorded as a consistency
defect, rather than being absorbed into a physical or cohomological
remainder.

\begin{decision}[V46 monolithic Hessian-dual regression]
\label{dec:n1v45-next}
Construct the \(45\)-by-\(45\) vertical Hessian directly in the
truncated bivariate algebra
\[
 {\mathbb Q}[v,q,i]/(v^3,q^3,i^2+1),                       \tag{45.15}
\]
from the original plaquette and transported-completion words, with
the routed background substituted before differentiation.  Extract
the \(v^2q^2\) coefficient of
\[
 \frac12\operatorname{Tr}
 \left(
 G H_{vv}-G H_v G H_v
 \right)                                                    \tag{45.16}
\]
without calling the V31 cubic or V33/V41 quartic assembly routines.
Run it at \(0001\) and \(0100\).  Comparison with the separated
bubble and tadpole values will identify the erroneous compiler row.
\end{decision}

\begin{assessment}[Progress after compact V45]
\label{ass:n1v45-status}
The mandatory audit is complete.  Exact product-Haar disintegration,
exact coarse-holonomy transport, and fibrewise conditioning eliminate
the last nonlinear-Jacobian explanation of the V43 coefficient.
The finite bottleneck is now a concrete second-jet compiler
consistency defect with a monolithic independent regression target.
\end{assessment}

\begin{firewall}[Claim boundary after compact V45]
\label{fire:n1v45-claim}
V45 proves that no measured re-tree density term can repair (43.5).
It does not yet identify whether the cubic or quartic compiler is
wrong, repair the corner table, produce a valid complete \(M=2\)
mean, control continuum quadrature or aliasing, close SCIC or ADMC,
construct the continuum Yang--Mills measure and clustering theory,
or prove the Yang--Mills mass gap.
\end{firewall}

\subsection{Parent record}

This V45 note was formed by appending the preceding material to the
validated compact V44, whose record was
\[
\begin{array}{c|c}
\text{lines}&17718\\
\text{bytes}&606281\\
\text{SHA--256}&
\texttt{A9C778335C04A2B23239066574643AE57383B34B802F456F2CA7AA3807217C6C}.
\end{array}
\]

% =============================================================================
% END APPENDED COMPACT-SERIES V45 MATERIAL
% =============================================================================

% =============================================================================
% BEGIN APPENDED COMPACT-SERIES V46 MATERIAL
% =============================================================================

\section{Compact-series V46: direct spectator replay localizes the
defect upstream of Walsh batching}
\label{sec:n1v46}

\subsection{Independent phase realization at \(0001\)}

The V43 quartic calculation represented all sixteen internal corners
simultaneously in the Walsh algebra.  To test that batching step, the
older V36 direct phase evaluator was preserved and changed only in
its internal axis:
\[
 k=\pi\widehat e_4,\qquad
 q\parallel\widehat e_1,\qquad
 B\parallel\widehat e_3.                                   \tag{46.1}
\]
This is corner \(0001\) in the bit convention of V39.

The direct evaluator inserts a sign
\((-1)^{a_4}\) separately in every Laurent monomial and every
translated noncommutative word.  It does not use the
\(({\mathbb Z}/2{\mathbb Z})^4\) convolution or character evaluation
of V40--V43.

\subsection{Exact direct card}

At (46.1), the quartic tadpoles are
\[
\begin{split}
 {\cal T}_{\rm W}
 &=
 \left(
 \frac{1842}{35},0,-\frac{15097}{420}
 \right),\\
 {\cal T}_*
 &=
 \left(
 \frac{40144250631}{505450939},0,
 -\frac{23228228303606089}{145873140995400}
 \right).                                                   \tag{46.2}
\end{split}
\]
The cubic bubbles are
\[
\begin{split}
 {\cal B}_{\rm W}
 &=
 \left(
 -\frac{12258}{245},0,
 \frac{69525923}{2881200}
 \right),\\
 {\cal B}_*
 &=
 \left(
 -\frac{10030599126}{133013405},0,
 \frac{7765281853708263829154113}
 {737317160910043246806000}
 \right).                                                   \tag{46.3}
\end{split}
\]
Consequently
\[
\begin{split}
 {\cal R}_{\rm W}(0001)
 &=
 \left(
 \frac{636}{245},0,-\frac{11346499}{960400}
 \right),\\
 {\cal R}_*(0001)
 &=
 \left(
 \frac{1448552823}{361036385},0,
 -\frac{36547338739977494627507199}
 {245772386970014415602000}
 \right),                                                   \tag{46.4}
\end{split}
\]
and their difference is
\[
 \boxed{
 \Delta{\cal R}(0001)
 =
 \left(
 \frac{5011052031}{3538156573},0,
 -\frac{206067652952114185227322937}
 {1505355870191338295562250}
 \right).}                                                  \tag{46.5}
\]

\begin{proposition}[Direct-card internal checks]
\label{prop:n1v46-checks}
The direct calculation passes:
\begin{enumerate}
\item symmetric zeroth and second free covariance matrices and
      antisymmetric first matrices for both actions;
\item all Wilson, endpoint-routed Wilson, completion, and combined
      cubic plus/minus parity identities;
\item zero first tadpole, bubble, and total jets for both actions;
\item the same stationary-score and product-Haar gates as V36.
\end{enumerate}
\end{proposition}

\begin{proof}
Every item is an exact Boolean assertion executed before the JSON
payload is emitted.  The successful process returned normally, so
all assertions passed.
\end{proof}

\subsection{Comparison with the paired V36 card}

The paired corner \(0100\) is
\(k=\pi\widehat e_2\).  After the V42 polarization correction, its
V36 endpoint-minus-Wilson response is
\[
 \Delta{\cal R}(0100)
 =
 \left(
 \frac{5011052031}{3538156573},0,
 -\frac{107153116857979395343397926}
 {752677935095669147781125}
 \right).                                                   \tag{46.6}
\]

\begin{theorem}[Exact reproduction of the V43 residual]
\label{thm:n1v46-reproduction}
The direct single-axis cards satisfy
\[
 \boxed{
 \Delta{\cal R}(0100)-\Delta{\cal R}(0001)
 =
 \left(
 0,0,
 -\frac{489372186744556308849}
 {89418228107593602350}
 \right).}                                                  \tag{46.7}
\]
This is exactly the corrected-pair Walsh residual (43.5).
\end{theorem}

\begin{proof}
The zeroth fractions in (46.5)--(46.6) are identical.  Put the two
second fractions over a common denominator and reduce.  Exact integer
arithmetic gives the fraction in (46.7), which is the fraction in
(43.5).
\end{proof}

The agreement is stronger at the diagram level.  Subtracting the
direct tadpoles (46.2) from the already-certified \(0100\) tadpoles
gives (43.4), and the analogous subtraction of (46.3) gives (43.3).

\begin{corollary}[Walsh and corner extraction are not the defect]
\label{cor:n1v46-Walsh-cleared}
The nonzero second residual is not created by the V40 Walsh
convolution, its character evaluation, or the act of extracting the
two corners from a sixteen-corner bank.
\end{corollary}

\begin{proof}
The direct evaluator realizes each internal sign at the location of
the original Laurent or word phase and never forms a mask.  It
nevertheless returns the same diagram residuals and their same sum.
\end{proof}

\begin{warning}[Extent of independence]
\label{warn:n1v46-independence}
The V46 program shares the underlying translated word derivatives,
router, inverse-jet formula, and bubble trace contraction with the
V31--V36 chain.  Therefore it does not validate those shared
ingredients.  It validates the direct internal-phase path against the
Walsh path and localizes the defect to their common upstream
second-\(q\)-jet response grammar.
\end{warning}

\subsection{A monolithic regression identity}

Let \(H_e(v,q;k)\) be the exact vertical Hessian of the original
group action after the routed background is substituted, with no
prior decomposition into cubic and quartic vertices.  Write
\[
 H_e(v,q;k)
 =
 H_{e,0}(k)
 +vH_{e,1}(q;k)
 +\frac{v^2}{2}H_{e,2}(q;k)+O(v^3),                        \tag{46.8}
\]
and \(G_e(k)=H_{e,0}(k)^{-1}\).  Direct differentiation of the
finite-dimensional determinant gives
\[
 \partial_v^2
 \left[\frac12\log\det H_e(v,q;k)\right]_{v=0}
 =
 \frac12\operatorname{Tr}
 \left(
 G_eH_{e,2}-G_eH_{e,1}G_eH_{e,1}
 \right).                                                   \tag{46.9}
\]

\begin{theorem}[Acceptance identity for the replacement compiler]
\label{thm:n1v46-acceptance}
A monolithic bivariate Hessian compiler is accepted only if:
\begin{enumerate}
\item its \(v\)-linear and \(v^2\) coefficients reproduce direct
      finite word differentiation before any colour reduction;
\item differentiating (46.9) twice in \(q\) agrees with its direct
      \(v^2q^2\) coefficient;
\item the result has zero residual between \(0001\) and \(0100\) for
      each action separately;
\item only after the preceding gates pass does its separation into
      tadpole and bubble rows agree with the legacy routines.
\end{enumerate}
\end{theorem}

\begin{proof}
The first two items are equality of two orders of differentiation of
the same finite polynomial/analytic matrix expression.  The third is
the exact \(\tau_{24}\) action, channel, router, and measure covariance
proved in V42--V45.  The fourth is then a diagnostic comparison, not
an assumption used to construct the answer.
\end{proof}

\subsection{Executable record and next step}

\begin{record}[V46 direct certificate]
\label{rec:n1v46-executable}
The direct verifier is
\[
 \texttt{scripts/verify\_ym\_v46\_direct\_spectator.py}.    \tag{46.10}
\]
Its source record is
\[
\begin{array}{c|c}
\text{lines}&1533\\
\text{bytes}&46737\\
\text{SHA--256}&
\texttt{06B36D13FD058B8F699B68BFCC7819628F8FDDD09ECAA58A8670E1F89802E56E}.
\end{array}                                                  \tag{46.11}
\]
The successful nonzero-mode payload SHA--256 is
\[
 \boxed{
 \texttt{321AE177664E6699D0DEAFDEF9F59B689202CF2C38C9A6C1A4F5CDE5A4EE5C56}.}
                                                                    \tag{46.12}
\]
\end{record}

\begin{decision}[V47 builds the monolithic bivariate Hessian]
\label{dec:n1v46-next}
Implement Theorem~\ref{thm:n1v46-acceptance} first at the two corners
\(0001,0100\).  Preserve noncommutative factor order and differentiate
the original Wilson and transported-completion owner functions in a
single truncated \((v,q)\) algebra.  Do not call the legacy cubic,
quartic, or bubble assembly functions until the monolithic response
has been frozen and hashed.
\end{decision}

\begin{assessment}[Progress after compact V46]
\label{ass:n1v46-status}
The off-axis direct card is exact and passes every local parity and
covariance check.  Its comparison with the old \(\pi\widehat e_2\)
card reproduces the V43 residual exactly, clearing the Walsh batching
and locating the inconsistency in the shared separated
second-\(q\)-jet grammar.  A fully independent monolithic Hessian
derivative is now the shortest upstream test.
\end{assessment}

\begin{firewall}[Claim boundary after compact V46]
\label{fire:n1v46-claim}
V46 localizes but does not yet repair the compiler defect.  It does
not provide a valid complete corner mean, continuum quadrature,
alias control, SCIC or ADMC closure, a continuum Yang--Mills measure,
exponential clustering, or the Yang--Mills mass gap.
\end{firewall}

\subsection{Parent record}

This V46 note was formed by appending the preceding material to the
validated compact V45, whose record was
\[
\begin{array}{c|c}
\text{lines}&18042\\
\text{bytes}&618636\\
\text{SHA--256}&
\texttt{700561E494E449CFEFB878CC8E3BB0E2BFB0481D0208A0639B7D69B7079B2E6C}.
\end{array}
\]

% =============================================================================
% END APPENDED COMPACT-SERIES V46 MATERIAL
% =============================================================================

% =============================================================================
% BEGIN APPENDED COMPACT-SERIES V47 MATERIAL
% =============================================================================

\section{Compact-series V47: automatic Wilson Hessian and correction
of pointwise corner covariance}
\label{sec:n1v47}

\subsection{The actionwise obstruction}

Before rebuilding the word derivatives, separate the corrected
spectator residual of V43 into the two actions.  Subtracting the
\(0001\) card of V46 from the \(0100\) card of V36 gives
\[
\begin{array}{c|c|c|c}
 &\Delta{\cal T}^{[2]}&\Delta{\cal B}^{[2]}
 &\Delta{\cal R}^{[2]}\\ \hline
{\rm Wilson}&
-\dfrac{1013}{120}&
\dfrac{763999}{68600}&
\dfrac{277351}{102900}\\[6pt]
{\rm endpoint}&
\dfrac{1105974754235329}{729365704977000}&
-\dfrac{527655177976252840919127}
 {122886193485007207801000}&
-\dfrac{42575865109780343149}{15328839104158903260}.
\end{array}                                                  \tag{47.1}
\]
Both zeroth and first residuals vanish for each action.  Moreover,
\[
 -\frac{42575865109780343149}{15328839104158903260}
 -\frac{277351}{102900}
 =
 -\frac{489372186744556308849}{89418228107593602350},       \tag{47.2}
\]
which recovers (43.5).  Thus the first nontrivial failure is already
present in the common Wilson sector; it cannot be attributed solely
to the transported completion.

\subsection{A derivative compiler independent of the separated
writers}

Let
\[
 {\cal A}:=
 {\mathbb Q}(i)[q]/(q^3)
 \mathbin{\widehat\otimes}
 {\mathbb Q}[x,y,b_+,b_-]/
 (x^2,y^2,b_+^2,b_-^2),                                   \tag{47.3}
\]
where the four amplitude variables commute.  They are square-zero
only because the desired coefficient is multilinear; no discarded
power can contribute to that coefficient.

For a positively oriented link factor, the new compiler inserts
\[
 \exp(b_+B_+(q)+b_-B_-(q))
 \exp(xX(k)+yY(k,q)),                                      \tag{47.4}
\]
and for the inverse factor it inserts, in the reverse order,
\[
 \exp(-xX(k)-yY(k,q))
 \exp(-b_+B_+(q)-b_-B_-(q)).                               \tag{47.5}
\]
The routed fields \(B_\pm(q)\), the internal corner sign, and every
cell displacement phase are substituted before the factors of a
plaquette are multiplied.  Each exponential in (47.4)--(47.5) is a
finite sum in \({\cal A}\).  Colour matrices are not commuted or
symmetrized across the two exponentials.

For the plus cubic writer assign momenta
\[
 x:k,\qquad y:-k-q,\qquad b_+:q,                           \tag{47.6}
\]
and extract \([xyb_+]\).  For the minus writer assign
\[
 x:k,\qquad y:-k+q,\qquad b_-:-q,                          \tag{47.7}
\]
and extract \([xyb_-]\).  For the quartic writer assign
\[
 x:k,\qquad y:-k,\qquad b_+:q,\qquad b_-:-q,               \tag{47.8}
\]
and extract \([xyb_+b_-]\), followed by the full eight-colour
dual-basis contraction.  This construction calls neither the V31
cubic assembler nor the V33/V43 quartic assembler.

\begin{lemma}[Exactness of nilpotent coefficient extraction]
\label{lem:n1v47-nilpotent}
For every Wilson plaquette word, the coefficients specified in
(47.6)--(47.8) are exactly its mixed derivatives
\[
 D_xD_yD_{b_+}S_{\rm W},\qquad
 D_xD_yD_{b_-}S_{\rm W},\qquad
 D_xD_yD_{b_+}D_{b_-}S_{\rm W}                            \tag{47.9}
\]
at zero amplitude, through external order \(q^2\).
\end{lemma}

\begin{proof}
In the quotient (47.3), the coefficient of a product of distinct
square-zero variables is its mixed derivative at the origin.  The
ordinary exponential series produces every ordering of insertions
inside one exponential with coefficient \(1/r!\).  Equations
(47.4) and (47.5) reproduce respectively left background
translation and the inverse of that translated factor.  Multiplying
the four oriented factors in their plaquette order therefore gives
the original group word in \({\cal A}\), not a derivative formula
for that word.  Since all retained terms have \(q\)-degree at most
two, reduction modulo \(q^3\) is exact for (47.9).
\end{proof}

\subsection{Entrywise comparison and response cards}

\begin{theorem}[Automatic and separated Wilson writers agree]
\label{thm:n1v47-writer-equality}
At both \(k=0001\) and \(k=0100\), and for every
\(0\leq j\leq2\), the automatic matrices obey
\[
 \boxed{
 P^{\pm,[j]}_{\rm auto}(k)=P^{\pm,[j]}_{\rm legacy}(k),
 \qquad
 Q^{[j]}_{\rm auto}(k)=Q^{[j]}_{\rm legacy}(k).}            \tag{47.10}
\]
These are eighteen entrywise equalities of \(45\)-by-\(45\) rational
matrices.  The plus-writer and raw-quartic hashes are
\[
\begin{array}{c|c|c|c}
 &j=0&j=1&j=2\\ \hline
P^+_{\rm auto}(0001)&
\texttt{0D6B63B563F33236}&
\texttt{E9BA054321D3C582}&
\texttt{FDE5F6E4B3C2D5D7}\\
Q_{\rm auto}(0001)&
\texttt{0184DA127870F025}&
\texttt{D3CD5ACF94F846A7}&
\texttt{99EA8A4AC5FE6585}\\
P^+_{\rm auto}(0100)&
\texttt{8A7934A7017C71AD}&
\texttt{0925A4DBC3F4ACF4}&
\texttt{D38C8EA5805DD540}\\
Q_{\rm auto}(0100)&
\texttt{E76488FC0A72F202}&
\texttt{D3CD5ACF94F846A7}&
\texttt{093C49D0A160A8D1}.
\end{array}                                                  \tag{47.11}
\]
The displayed values are the first sixteen hexadecimal digits of
the full SHA--256 records.
\end{theorem}

\begin{proof}
The automatic matrices are frozen and hashed before the legacy
functions are imported.  The program then constructs the separated
writers and compares every reduced rational entry.  All comparisons
are true.  Lemma~\ref{lem:n1v47-nilpotent} makes the automatic side
a direct derivative of the Wilson group word, so (47.10) validates
the separated Wilson vertex grammar rather than merely comparing
two copies of it.
\end{proof}

The automatic contractions give
\[
\begin{array}{c|c|c|c}
k&{\cal T}_{\rm W}^{[0,1,2]}&
{\cal B}_{\rm W}^{[0,1,2]}&
{\cal R}_{\rm W}^{[0,1,2]}\\ \hline
0001&
\left(\dfrac{1842}{35},0,-\dfrac{15097}{420}\right)&
\left(-\dfrac{12258}{245},0,\dfrac{69525923}{2881200}\right)&
\left(\dfrac{636}{245},0,-\dfrac{11346499}{960400}\right)\\[8pt]
0100&
\left(\dfrac{302}{5},0,-\dfrac{7457}{168}\right)&
\left(-\dfrac{14162}{245},0,\dfrac{101613881}{2881200}\right)&
\left(\dfrac{636}{245},0,-\dfrac{26273669}{2881200}\right).
\end{array}                                                  \tag{47.12}
\]
Hence
\[
 \boxed{
 {\cal R}_{\rm W}(0100)-{\cal R}_{\rm W}(0001)
 =
 \left(0,0,\frac{277351}{102900}\right).}                  \tag{47.13}
\]

\subsection{Why pointwise corner covariance was too strong}

Let \({\cal H}=\bigoplus_k{\cal H}_k\) be the vertical fluctuation
space in internal momentum.  At zero background, the re-tree
linearization is block diagonal after the spatial permutation:
\[
 J_\tau(0){\cal H}_k\subseteq{\cal H}_{\tau k}.             \tag{47.14}
\]
For a background Fourier mode \(B(q)\), however, multiplication by
that mode shifts momentum.  The background-dependent derivative has
the form
\[
 J_\tau(B)
 =
 J_0+
 b_+J_+(q)+b_-J_-(q)+O(b^2),                              \tag{47.15}
\]
with
\[
 J_+(q){\cal H}_k\subseteq{\cal H}_{\tau k+q},\qquad
 J_-(q){\cal H}_k\subseteq{\cal H}_{\tau k-q}.              \tag{47.16}
\]
Thus \(J_\tau(B)\) is not a direct sum of independent transformations
of the individual \(k\)-cards used in (36.14)--(36.16).

\begin{theorem}[Correct scope of re-tree determinant invariance]
\label{thm:n1v47-global-not-pointwise}
The Haar and action invariances of V42--V45 imply equality of the
complete internal-momentum trace
\[
 \sum_k {\cal R}_{e,T}(k,q)
 =
 \sum_k {\cal R}_{e,\tau T}(\tau k,q),                     \tag{47.17}
\]
or of the corresponding Brillouin integral.  They do not imply
\[
 {\cal R}_{e,T}(k,q)
 =
 {\cal R}_{e,\tau T}(\tau k,q)                             \tag{47.18}
\]
for each routed loop-momentum card.
\end{theorem}

\begin{proof}
The determinant identity is an identity on the full fluctuation
space:
\[
 \log\det(J_\tau(B)^*H(B)J_\tau(B))
 =
 \log\det H(B)+2\log|\det J_\tau(B)|.                      \tag{47.19}
\]
Taking two background derivatives in (47.15) produces compositions
\(J_-J_+\) and \(J_+J_-\).  By (47.16), their trace closes only after
the intermediate shifted momentum has been included and the full
trace over \(\bigoplus_k{\cal H}_k\) is taken.  Cyclicity of that full
trace and translation of the momentum index give (47.17).  A
projection onto one \({\cal H}_k\) does not commute with \(J_\pm\);
cyclicity therefore gives no identity for the projected trace.
This proves that (47.18) is not a consequence of determinant
invariance.
\end{proof}

At \(q=0\), the shifts in (47.16) collapse, which is consistent with
the exact zeroth-order cancellation in V43.  The nonzero second
coefficient (47.13) is likewise consistent with a redistribution
among loop-momentum cards.

\begin{correction}[Withdrawal of the V45 pointwise inference]
\label{corr:n1v47-v45}
The product-Haar disintegration, exact boundary transport, and
fibrewise measure invariance in V45 remain valid.  The conclusion of
Theorem~\ref{thm:n1v45-compiler-defect} is withdrawn: its proof
silently passed from the full determinant identity to the
card-by-card statement (47.18).  Theorem
\ref{thm:n1v47-writer-equality} independently validates the Wilson
side of the allegedly defective compiler and (47.13) directly
falsifies that pointwise implication.
\end{correction}

\subsection{Executable record and next acquisition}

\begin{record}[V47 automatic Wilson certificate]
\label{rec:n1v47-executable}
The verifier is
\[
 \texttt{scripts/verify\_ym\_v47\_monolithic\_wilson.py}.   \tag{47.20}
\]
Its source record is
\[
\begin{array}{c|c}
\text{lines}&621\\
\text{bytes}&21400\\
\text{SHA--256}&
\texttt{BA8F36970A0FD344EED9BEB3A6E694A8D6CFF9A504E1542FF3D7B2E1FF8E098B}.
\end{array}                                                  \tag{47.21}
\]
The canonical payload SHA--256 is
\[
 \boxed{
 \texttt{F1AF42F95313759C3F351B5BC1106C62265B1690CF89C71533615C7A26363D07}.}
                                                                    \tag{47.22}
\]
The payload records the failed pointwise covariance gate as
\texttt{false}; this is an intended falsification result, not a
successful covariance assertion.
\end{record}

\begin{decision}[V48 tests the complete momentum trace]
\label{dec:n1v47-next}
Compute the Wilson tadpole and bubble response at all sixteen
\(M=2\) internal corners without orbit compression.  Compare the
raw sixteen-card sum before and after \(\tau_{24}\), retaining the
background-induced momentum-routing terms.  If the full trace is
invariant, replace the invalid pointwise compression by a
trace-level certificate.  If it is not, construct the explicit
mixed \(J_-J_+\) re-tree row before revisiting the endpoint action.
\end{decision}

\begin{assessment}[Progress after compact V47]
\label{ass:n1v47-status}
The monolithic Wilson word expansion has passed its entrywise
comparison and reproduced both old response cards.  It clears the
Wilson cubic and quartic differentiation routines and identifies the
actual upstream error: re-tree invariance was applied to individual
loop-momentum cards even though the background-dependent coordinate
map shifts loop momentum.  The next valid symmetry target is the
uncompressed full trace.
\end{assessment}

\begin{firewall}[Claim boundary after compact V47]
\label{fire:n1v47-claim}
V47 validates the Wilson word derivatives and corrects the scope of
the re-tree argument.  It does not yet prove the sixteen-corner trace
identity, validate the automatic endpoint completion, produce a
valid \(M=2\) total mean, control continuum quadrature or aliasing,
close SCIC or ADMC, construct the continuum Yang--Mills measure and
clustering theory, or prove the Yang--Mills mass gap.
\end{firewall}

\subsection{Parent record}

This V47 note was formed by appending the preceding material to the
validated compact V46, whose record was
\[
\begin{array}{c|c}
\text{lines}&18335\\
\text{bytes}&627990\\
\text{SHA--256}&
\texttt{8D00DD00E4F38552262670E57FE9EA3810EFF75F3A2A7DC115BCCDBCF3A97B39}.
\end{array}
\]

% =============================================================================
% END APPENDED COMPACT-SERIES V47 MATERIAL
% =============================================================================

% =============================================================================
% BEGIN APPENDED COMPACT-SERIES V48 MATERIAL
% =============================================================================

\section{Compact-series V48: the uncompressed sixteen-corner
response trace}
\label{sec:n1v48}

\subsection{Reopening the aborted all-corner calculation}

The V43 Walsh tadpole program assembled all sixteen corner values but
aborted at the first nonzero corrected pair because pointwise orbit
constancy was then an asserted gate.  V47 proves that this gate is
not implied by re-tree invariance.  The V48 calculation therefore
retains the same exact word, colour, router, and Walsh algebra, but
records all pointwise residuals and continues to the raw arithmetic
mean.

For a corner \(k\in\{0,\pi\}^4\), write
\[
 {\cal R}_e(k)=
 {\cal T}_e(k)+{\cal B}_e(k),\qquad
 e\in\{{\rm W},*\}.                                        \tag{48.1}
\]
The uncompressed \(M=2\) mean is
\[
 \overline{\cal R}_{e,2}
 :=
 \frac1{16}\sum_{k\in\{0,\pi\}^4}{\cal R}_e(k).             \tag{48.2}
\]
No representative weight or orbit equality is used in (48.2).

\subsection{All four corrected pair residuals}

The four member-minus-reference residuals of the complete
endpoint-minus-Wilson response are
\[
\begin{array}{c|c}
\text{pair}&
\Delta{\cal R}^{[2]}(\text{member})
-\Delta{\cal R}^{[2]}(\text{reference})\\ \hline
0100-0001&
-\dfrac{489372186744556308849}{89418228107593602350}\\[7pt]
0110-0011&
-\dfrac{875249732391971777}{143250742304670000}\\[7pt]
1100-1001&
\dfrac{149030431105175088126059}
 {19080998874982044000000}\\[7pt]
1110-1011&
-\dfrac{1235784701205221}{215242928032875}.
\end{array}                                                  \tag{48.3}
\]
Every residual in external orders zero and one is exactly zero.

\begin{theorem}[Pointwise spectator compression is rejected]
\label{thm:n1v48-no-compression}
None of the four nontrivial \(\tau_{24}\) pairs in (48.3) has equal
second response cards.  Consequently a twelve-representative
compression based on pointwise \(\tau_{24}\) equality is invalid for
this background-dependent loop-momentum routing.
\end{theorem}

\begin{proof}
Each fraction in (48.3) has positive denominator and nonzero
numerator.  Thus the two cards in every nontrivial pair differ at
order two.  A representative compression replacing both cards by
one of them would change their two-term contribution unless an
additional routing correction were first supplied.
\end{proof}

The exact cancellation at order zero persists in all four pairs.
This is consistent with Theorem~\ref{thm:n1v47-global-not-pointwise}:
when \(q=0\), the background-induced momentum shifts collapse, while
the differentiated \(q\)-dependence need not be pointwise invariant.

\subsection{Exact raw means}

The newly completed all-corner tadpole averages are
\[
\begin{array}{c|c|c|c}
e&\overline{\cal T}^{[0]}_{e,2}&
\overline{\cal T}^{[1]}_{e,2}&
\overline{\cal T}^{[2]}_{e,2}\\ \hline
{\rm W}&
\dfrac{262589}{4480}&0&-\dfrac{93247}{2048}\\[7pt]
*&
\dfrac{1189208771492363763523}{17029424977324323200}&0&
-\dfrac{7380496832421770973636143879}
 {46338525028297539797760000}.
\end{array}                                                  \tag{48.4}
\]
The V39 all-corner bubble averages, imported only after verifying its
canonical payload hash, combine with (48.4) to give
\[
\boxed{
\overline{\cal R}_{{\rm W},2}
=
\left(
\frac{176943}{627200},
0,
-\frac{2822949104411}{265531392000}
\right).}                                                   \tag{48.5}
\]
For the endpoint action,
\[
\boxed{
\overline{\cal R}_{*,2}
=
\left(
\frac{157603277390904552027570847451}
 {684609230311820211100422720000},
0,
-\frac{5823821044825964006189852662538728905719501293128460064877013}
 {36786699279534470965406231814271312468634765095027200000000}
\right).}                                                   \tag{48.6}
\]
Their exact difference is
\[
\boxed{
\overline{\Delta{\cal R}}_2
=
\left(
-\frac{4441971352736182354322716403}
 {85576153788977526387552840000},
0,
-\frac{97013033679684607051275019240872676380229076929720667338373}
 {656905344277401267239396996683416294082763662411200000000}
\right).}                                                   \tag{48.7}
\]
For scale only,
\[
 \overline{\cal R}_{{\rm W},2}
 \approx(0.282115752551020,0,-10.6313196460440),            \tag{48.8}
\]
\[
 \overline{\cal R}_{*,2}
 \approx(0.230209103840333,0,-158.313226217225),            \tag{48.9}
\]
\[
 \overline{\Delta{\cal R}}_2
 \approx(-0.0519066487106870,0,-147.681906571181).          \tag{48.10}
\]

\begin{theorem}[Exact uncompressed \(M=2\) mean]
\label{thm:n1v48-raw-mean}
Equations (48.5)--(48.7) are the arithmetic means of all sixteen
corner cards, with no pointwise orbit assumption.  In every external
order,
\[
 \boxed{
 \overline{\cal R}_{*,2}
 -\overline{\cal R}_{{\rm W},2}
 =\overline{\Delta{\cal R}}_2.}                            \tag{48.11}
\]
\end{theorem}

\begin{proof}
The V48 program evaluates the three quartic Walsh matrices at each
of the sixteen characters, contracts each with the appropriate exact
inverse covariance, and accumulates the rational tadpole sums before
division by sixteen.  It imports the V39 bubble means whose complete
sixteen-card payload has SHA--256
\[
 \texttt{0B070021D0B372CEFA5132E5998FFC8E888F87B476D9BC5060EC4F2BE459FEC5}.
                                                               \tag{48.12}
\]
Addition of the two exact means gives (48.5)--(48.7).
The program then checks (48.11) by rational subtraction in all three
orders.  No use is made of the false equality tested in
Theorem~\ref{thm:n1v48-no-compression}.
\end{proof}

\begin{corollary}[The raw second mean is strictly negative]
\label{cor:n1v48-negative}
\[
 \overline{\Delta{\cal R}}_2^{[2]}<0.                      \tag{48.13}
\]
\end{corollary}

\begin{proof}
The numerator and denominator of the third entry in (48.7) are
positive integers and the displayed sign is negative.
\end{proof}

\begin{warning}[A corner mean is not a continuum mass]
\label{warn:n1v48-not-mass}
The sign in (48.13) is a finite \(M=2\), one-loop, one-channel
response sign.  Without a certified comparison to the Brillouin
integral, higher-scale control, nonperturbative measure construction,
and clustering, it is neither a pole mass nor a Yang--Mills mass-gap
estimate.
\end{warning}

\subsection{Executable record and next acquisition}

\begin{record}[V48 uncompressed trace certificate]
\label{rec:n1v48-executable}
The verifier is
\[
 \texttt{scripts/verify\_ym\_v48\_uncompressed\_trace.py}.  \tag{48.14}
\]
Its source record is
\[
\begin{array}{c|c}
\text{lines}&919\\
\text{bytes}&28525\\
\text{SHA--256}&
\texttt{607DD5A38FD686581A555434B5EF9E179F1F761F1E553457E46DCE49CA88EAEE}.
\end{array}                                                  \tag{48.15}
\]
The canonical payload SHA--256 is
\[
 \boxed{
 \texttt{83B54DC88FAA62F5972D7D689F056D69C9DC6106FEFD275675CD76B09E8CE411}.}
                                                                    \tag{48.16}
\]
The payload deliberately records
\texttt{pointwise\_orbit\_compression\_rejected=true}.
\end{record}

\begin{decision}[V49 returns to the continuum-comparison gate]
\label{dec:n1v48-next}
Do not further compress the \(M=2\) table.  Reuse the analytic-strip
and trapezoidal machinery already proved in V17 and the response
regularity developed in V37--V38, but audit their hypotheses against
the now-correct uncompressed integrand.  The next certificate must
identify an explicit strip width, a uniform supremum bound, and an
alias remainder small enough to determine the sign of the continuum
second-response mean; otherwise it must state the first missing
quantitative input.
\end{decision}

\begin{assessment}[Progress after compact V48]
\label{ass:n1v48-status}
The formerly aborted all-corner calculation now completes.  It
rejects all four false pointwise pair equalities and supplies the
first valid uncompressed \(M=2\) total mean in the corrected physical
channel.  The finite algebraic bottleneck has moved from corner
assembly to a quantitative finite-grid-to-continuum comparison.
\end{assessment}

\begin{firewall}[Claim boundary after compact V48]
\label{fire:n1v48-claim}
V48 proves only the exact raw \(M=2\) response mean.  It does not
bound its distance from the Brillouin integral, determine the
continuum response sign, control higher loops or nonperturbative
remainders, close SCIC or ADMC, construct the continuum Yang--Mills
measure and clustering theory, or prove the Yang--Mills mass gap.
\end{firewall}

\subsection{Parent record}

This V48 note was formed by appending the preceding material to the
validated compact V47, whose record was
\[
\begin{array}{c|c}
\text{lines}&18675\\
\text{bytes}&640193\\
\text{SHA--256}&
\texttt{0C68117F66560B6DD180B9C2BB8C4F0EA84280A5D949D1CB79404DEBB85ABDDB}.
\end{array}
\]

% =============================================================================
% END APPENDED COMPACT-SERIES V48 MATERIAL
% =============================================================================

% =============================================================================
% BEGIN APPENDED COMPACT-SERIES V49 MATERIAL
% =============================================================================

\section{Compact-series V49: the proved strip cannot certify a
continuum sign from the \(M=2\) mean}
\label{sec:n1v49}

\subsection{The applicable inherited strip}

V48 supplies an exact uncompressed mean, but a finite grid decides
the Brillouin mean only after the alias remainder is smaller than the
observed margin.  The sharp common inverse strip already proved in
V21 is
\[
 \rho_{21}
 =
 \frac13\log\left(1+\frac{147}{105475000}\right).           \tag{49.1}
\]
The V38 background-response construction uses internal and external
strip widths \(\rho_{21}/4\).  After Cauchy extraction of the second
external derivative and the standard half-strip Fourier estimate, a
safe Wiener width for the resulting internal-momentum scalar is
\[
 \boxed{\rho_W:=\frac{\rho_{21}}8.}                         \tag{49.2}
\]
Numerically, only to display scale,
\[
 \rho_{21}\approx4.64564739080347\cdot10^{-7},\qquad
 \rho_W\approx5.80705923850434\cdot10^{-8}.                 \tag{49.3}
\]

Let \(f(k)\) denote the corrected endpoint-minus-Wilson transverse
second-response integrand, and put
\[
 N_W:=\|f\|_{\rho_W}.                                      \tag{49.4}
\]
The V17 grid theorem gives
\[
 \left|\langle f\rangle_M-\langle f\rangle\right|
 \le N_W C_M,\qquad
 C_M:=
 \left[
 \frac{1+e^{-\rho_WM}}{1-e^{-\rho_WM}}
 \right]^4-1.                                               \tag{49.5}
\]
Equivalently,
\[
 C_M=\coth^4\left(\frac{\rho_WM}{2}\right)-1.               \tag{49.6}
\]

\subsection{A norm-independent resolution obstruction}

\begin{lemma}[The grid mean is bounded by the Wiener stock]
\label{lem:n1v49-grid-below-stock}
For every \(M\),
\[
                         |\langle f\rangle_M|\le N_W.       \tag{49.7}
\]
\end{lemma}

\begin{proof}
At every real \(k\),
\[
 |f(k)|
 \le\sum_m|\widehat f(m)|
 \le\sum_m e^{\rho_W|m|_1}|\widehat f(m)|
 =N_W.
\]
Average this inequality over the grid and apply the triangle
inequality.
\end{proof}

\begin{theorem}[Necessary grid size for the inherited certificate]
\label{thm:n1v49-necessary-M}
Suppose \(\langle f\rangle_M\ne0\).  For the symmetric V17 enclosure
in (49.5) to exclude zero, it is necessary that
\[
 \boxed{
 M>\frac{843800000}{49}>17220408.}                         \tag{49.8}
\]
Thus the smallest integer not excluded by this rigorous necessary
bound is
\[
                         M=17220409.                        \tag{49.9}
\]
\end{theorem}

\begin{proof}
Excluding zero from the enclosure centred at
\(\langle f\rangle_M\) requires
\[
                         N_WC_M<|\langle f\rangle_M|.
\]
Lemma~\ref{lem:n1v49-grid-below-stock} therefore makes \(C_M<1\)
necessary.  By (49.6), this requires
\[
 \tanh\left(\frac{\rho_WM}{2}\right)>2^{-1/4},
\]
and hence
\[
 \frac{\rho_WM}{2}>
 \operatorname{artanh}(2^{-1/4})>\frac12.                  \tag{49.10}
\]
Therefore \(M>1/\rho_W\).  The elementary strict inequality
\(\log(1+x)<x\) for \(x>0\) gives
\[
 \rho_W
 <
 \frac18\frac13\frac{147}{105475000}
 =
 \frac{49}{843800000}.                                     \tag{49.11}
\]
Combining the last two inequalities proves (49.8).  Integer division
gives (49.9).
\end{proof}

The exact threshold obtained by retaining the inverse hyperbolic
tangent rather than the weaker last inequality in (49.10) is, at
high precision,
\[
 M\ge42163380.                                              \tag{49.12}
\]
This decimal refinement is diagnostic only; the theorem uses the
rational lower bound (49.8).

\begin{corollary}[The \(M=2\) interval necessarily contains zero]
\label{cor:n1v49-M2-fails}
The V17--V21 worst-case certificate applied at \(M=2\) cannot
determine the sign of the continuum mean from (48.7), even if
\(N_W\) were evaluated exactly.
\end{corollary}

\begin{proof}
The \(M=2\) mean in (48.7) is nonzero, while \(2\) violates the
necessary condition (49.8).  Hence the error radius in (49.5) is at
least the magnitude of the grid mean.  The resulting symmetric
interval contains zero.
\end{proof}

For orientation, direct high-precision evaluation gives
\[
 C_2
 \approx
 8.7937674103272578822273804672\cdot10^{28}.                \tag{49.13}
\]
This enormous value explains why further exact refinement of the
sixteen \(M=2\) cards cannot repair the continuum-sign certificate.

\subsection{The first missing quantitative input}

The obstruction in Theorem~\ref{thm:n1v49-necessary-M} is a defect
of the deliberately crude complex-strip estimate, not evidence that
the actual integrand is poorly behaved near the real torus.  V21
bounded the entire endpoint Laurent matrix by an entrywise
coefficient sum and perturbed from a global real floor.  It did not
locate the nearest complex zero of either determinant.

Moreover, the finite stocks in V38,
\[
 p_{38},\qquad q_{38},\qquad j_{38},                        \tag{49.14}
\]
were proved finite but were not numerically populated.  The exact
V39/V48 corner evaluators retain character values, not the complete
Laurent coefficient arrays needed to evaluate (49.14).  Therefore
there are two distinct acquisition tasks:
\begin{enumerate}
\item populate certified coefficient norms for the complete
      corrected response;
\item replace the global Neumann strip by a determinant or singular-
      value exclusion on a substantially wider complex polystrip.
\end{enumerate}
Only after both are available can one choose a feasible \(M\) and
apply the terminating interval theorem of V18.

\begin{proposition}[No conclusion about the continuum sign]
\label{prop:n1v49-no-sign}
The negative number in (48.10), together with (49.1)--(49.6), proves
neither sign for the continuum Brillouin mean.
\end{proposition}

\begin{proof}
Corollary~\ref{cor:n1v49-M2-fails} shows that the certified interval
obtained from the available strip contains zero.  An interval
containing zero is consistent with a negative, zero, or positive
continuum value.  No one of those alternatives follows.
\end{proof}

\subsection{Executable record and next acquisition}

\begin{record}[V49 strip-resolution certificate]
\label{rec:n1v49-executable}
The scale verifier is
\[
 \texttt{scripts/verify\_ym\_v49\_strip\_resolution.py}.     \tag{49.15}
\]
Its source record is
\[
\begin{array}{c|c}
\text{lines}&54\\
\text{bytes}&2185\\
\text{SHA--256}&
\texttt{CD2BFC75B63B1D0290B24A76FC7B3F8C2A034D64DA31C34FB89BCED85568FF19}.
\end{array}                                                  \tag{49.16}
\]
The canonical payload SHA--256 is
\[
 \boxed{
 \texttt{367A9FC88CB6AD01E9BA76632F1C043AA9783A9AA003121A160F2385A8FDC191}.}
                                                                    \tag{49.17}
\]
The program checks (49.11), the strict lower bound in (49.10),
the integer in (49.9), and \(C_2>1\).
\end{record}

\begin{decision}[V50 audit and strip-sharpening acquisition]
\label{dec:n1v49-next}
Perform the mandatory SM/NS companion audit.  Then inspect the exact
determinant or resolvent localization methods available there before
choosing the V50 route.  The Yang--Mills output must be either:
\begin{enumerate}
\item an explicit wider polystrip with certified inverse bounds for
      both Wilson and endpoint Hessians; or
\item a proof that the present coefficient representation cannot
      supply such a strip, together with a precise interval-
      determinant subdivision algorithm and its first executable
      cell certificate.
\end{enumerate}
Do not spend another iteration refining \(M=2\) corner values.
\end{decision}

\begin{assessment}[Progress after compact V49]
\label{ass:n1v49-status}
The existing analytic machinery has been audited against the
corrected uncompressed response.  It is logically applicable but
quantitatively incapable of transporting the \(M=2\) sign to the
continuum.  The obstruction is now explicit and norm independent:
the inherited strip alone forces the necessary resolution
(49.8).  The programme must sharpen the complex inverse geometry
and populate the nonlinear coefficient stocks.
\end{assessment}

\begin{firewall}[Claim boundary after compact V49]
\label{fire:n1v49-claim}
V49 proves that the current finite-grid error bound cannot certify
the continuum sign from \(M=2\).  It does not prove that a grid of
size (49.9) is sufficient, evaluate the missing response norm,
locate complex determinant zeros, determine the continuum response,
close SCIC or ADMC, construct the continuum Yang--Mills measure or
clustering theory, or prove the Yang--Mills mass gap.
\end{firewall}

\subsection{Parent record}

This V49 note was formed by appending the preceding material to the
validated compact V48, whose record was
\[
\begin{array}{c|c}
\text{lines}&18942\\
\text{bytes}&649302\\
\text{SHA--256}&
\texttt{D31CEC7025D099C14DB7CB0838A1864B2AA2C558354AB2C52121AD1741593E5D}.
\end{array}
\]

% =============================================================================
% END APPENDED COMPACT-SERIES V49 MATERIAL
% =============================================================================

% =============================================================================
% BEGIN APPENDED COMPACT-SERIES V50 MATERIAL
% =============================================================================

\section{Compact-series V50: companion audit and the minimal legal
momentum packets}
\label{sec:n1v50}

\subsection{Mandatory companion audit}

The files inspected at this multiple-of-five checkpoint are
\[
\begin{array}{c|r|r|c}
\text{record}&\text{lines}&\text{bytes}&\text{SHA--256}\\ \hline
\text{SM--Einstein V28}&8249&283353&
\texttt{A79BFE05D11A5DEF7A5155FD23F7151E46D286CC08AADCEBDC8BE622ABA6A5B0}\\
\text{Navier--Stokes V26}&5319&278283&
\texttt{82C887F8C2C69E4F28FAD7C3DC8DE5B9099CB5A4BF431EBA1FFF3E91935978AD}.
\end{array}                                                   \tag{50.1}
\]

The Navier--Stokes file is unchanged from the V45 audit.  Its current
rank-one Oseen-kernel norm calculation supplies no determinant,
Laurent-strip, or momentum-packet input for the present problem.

The SM--Einstein file has advanced from V9 to V28.  Its V25
independently imports the YM V47 correction from cardwise covariance
to full-trace covariance.  Its V26 proves that the smallest universal
trace packets on a finite momentum group are cosets of the subgroup
generated by all background shifts.  That result is directly
applicable here and is specialized below.  SM V27--V28 concern
reciprocal-normalizer positivity and fine-measure marginalization;
they supply no wider complex inverse strip for the Yang--Mills
Hessian.  Accordingly the audit changes the trace organization but
does not remove the V49 strip bottleneck.

\subsection{The shift subgroup in the corrected channel}

Let
\[
 K_M:=({\mathbb Z}/M{\mathbb Z})^4                         \tag{50.2}
\]
index an \(M^4\) momentum grid.  In the corrected channel,
\[
 q=\frac{2\pi}{M}e_1,                                     \tag{50.3}
\]
where \(e_1=(1,0,0,0)\) in zero-based array order.  The two linear
background derivatives shift a fluctuation card by \(+e_1\) and
\(-e_1\).  Their generated subgroup is
\[
 H_1:=\langle e_1\rangle
 =
 \{ne_1:0\leq n<M\}\le K_M.                                \tag{50.4}
\]

\begin{theorem}[Minimal support packets for one axial background]
\label{thm:n1v50-axial-packets}
The minimal subsets of momentum cards closed under the background
shifts in (50.3) are the cosets
\[
 c(a_2,a_3,a_4)
 =
 \{(n,a_2,a_3,a_4):n\in{\mathbb Z}/M{\mathbb Z}\}.          \tag{50.5}
\]
There are exactly \(M^3\) packets, each containing \(M\) cards.
The associated projector
\[
 Q_a:=\sum_{k\in c(a)}P_k                                  \tag{50.6}
\]
commutes with every diagonal momentum operator and with every word
in the \(+e_1\) and \(-e_1\) shift operators.
\end{theorem}

\begin{proof}
Repeated addition of \(e_1\) visits every value of the first
coordinate and changes none of the other three.  Hence
\(|H_1|=M\), and its cosets are precisely (50.5).
The quotient has cardinality \(|K_M|/|H_1|=M^3\).
If \(k\in c(a)\), then \(k\pm e_1\in c(a)\); if
\(k\notin c(a)\), neither shifted card enters it.  Therefore \(Q_a\)
commutes with both shifts.  It plainly commutes with diagonal
operators, and hence with the algebra they generate.

Conversely, any union of full card fibres containing one \(k\) and
closed under both shifts contains \(k+ne_1\) for every \(n\), hence
contains the whole coset.  This proves minimality at the support
level.
\end{proof}

\begin{corollary}[Packetwise trace, not cardwise trace]
\label{cor:n1v50-packet-trace}
Let \(U(B)\) be the background-dependent re-tree linearization and
\(X'(B)=U(B)X(B)U(B)^{-1}\).  For every axial packet,
\[
 \boxed{
 \operatorname{Tr}(Q_aX'(B))
 =
 \operatorname{Tr}(Q_aX(B)).}                              \tag{50.7}
\]
No analogous identity for one \(P_k\) follows unless \(M=1\) or all
shift derivatives vanish.
\end{corollary}

\begin{proof}
Theorem~\ref{thm:n1v50-axial-packets} gives
\(Q_aU=UQ_a\).  Cyclicity gives
\[
 \operatorname{Tr}(Q_aUXU^{-1})
 =
 \operatorname{Tr}(UQ_aXU^{-1})
 =
 \operatorname{Tr}(Q_aX).
\]
For \(M>1\), a singleton card is not closed under \(e_1\), so its
projector need not commute with the transport.
\end{proof}

At second background order the closed walks have the form
\[
 k\longrightarrow k+e_1\longrightarrow k,\qquad
 k\longrightarrow k-e_1\longrightarrow k.                 \tag{50.8}
\]
Although they begin and end on one diagonal block, their intermediate
cards are indispensable.  Equation (50.8) is the finite-grid
version of (47.16).

\subsection{The exact \(M=2\) packets}

At \(M=2\), the eight packets are
\[
\begin{array}{llll}
\{0000,1000\},&
\{0001,1001\},&
\{0010,1010\},&
\{0011,1011\},\\
\{0100,1100\},&
\{0101,1101\},&
\{0110,1110\},&
\{0111,1111\}.
\end{array}                                                   \tag{50.9}
\]
The four spectator pairs previously tested were
\[
 \{0001,0100\},\quad
 \{0011,0110\},\quad
 \{1001,1100\},\quad
 \{1011,1110\}.                                             \tag{50.10}
\]

\begin{proposition}[Every old spectator pair crosses packets]
\label{prop:n1v50-crosses}
The two cards in each pair of (50.10) lie in distinct packets of
(50.9).
\end{proposition}

\begin{proof}
An axial packet fixes the last three bits and varies only the first.
The first pair in (50.10) belongs to packet indices \(1\) and \(4\)
in the ordering (50.9); the second belongs to \(3\) and \(6\).
The last two pairs repeat those packet indices with first bit one.
Thus none is contained in one reducing packet.
\end{proof}

This supplies an independent structural explanation for why the
four pointwise equalities fail.  The full sixteen-card sum in V48 is
a legal trace, but it is not the smallest legal decomposition for a
single axial external momentum: eight axial packet traces suffice.

\begin{corollary}[All-direction backgrounds force the full packet]
\label{cor:n1v50-all-directions}
If the required background set contains the four unit shifts
\(e_1,e_2,e_3,e_4\), then the generated subgroup is all of \(K_M\),
and the only nonzero support packet is the complete \(M^4\)-card
space.
\end{corollary}

\begin{proof}
The four unit vectors generate
\(({\mathbb Z}/M{\mathbb Z})^4\).
\end{proof}

\subsection{Relation to the formal external jet}

The packet theorem is an exact statement for nonzero periodic
external momenta on the finite grid.  The V48 cards are formal
derivatives at \(q=0\), so (50.9) is not used to manufacture new
equalities between those derivative cards.  Instead, the valid
continuum route is the one already isolated in V37--V38:
take nonzero \(q_M=2\pi e_1/M\), perform the packet or full trace,
and only then pass to the formal jet with a uniform analytic
estimate.  This order retains every intermediate momentum in
(50.8).

\begin{warning}[Packet structure does not improve the strip]
\label{warn:n1v50-no-strip}
Replacing individual cards by axial packets repairs the covariance
domain, but it does not enlarge the inverse strip (49.1) and does
not reduce the V49 alias multiplier by itself.  The packet theorem
is algebraic organization, not a quantitative determinant-zero
exclusion.
\end{warning}

\subsection{Executable record and next acquisition}

\begin{record}[V50 shift-packet certificate]
\label{rec:n1v50-executable}
The verifier is
\[
 \texttt{scripts/verify\_ym\_v50\_shift\_packets.py}.        \tag{50.11}
\]
Its source record is
\[
\begin{array}{c|c}
\text{lines}&103\\
\text{bytes}&3292\\
\text{SHA--256}&
\texttt{5472163FBC3EFD9FB08D02070E3E58FD23E318544C7768E6F6B0BF981423BA6F}.
\end{array}                                                  \tag{50.12}
\]
The canonical payload SHA--256 is
\[
 \boxed{
 \texttt{9B29A5D21F817BBC9E6CF92219D899C6D539B01E1C62890B421DC4E13DE038D9}.}
                                                                    \tag{50.13}
\]
The program enumerates the cosets for \(2\le M\le9\), verifies
\(M^3\) packets of size \(M\), checks (50.9), proves all four pairs
in (50.10) cross packets, and verifies that four axial generators
give one full packet at \(M=2\).
\end{record}

\begin{decision}[V51 resumes the inverse-strip problem]
\label{dec:n1v50-next}
Use packet traces for covariance, but return to the quantitative
bottleneck of V49.  First build exact one-dimensional restrictions
of the Wilson and endpoint determinants along imaginary coordinate
axes and along the diagonal imaginary ray.  Certified lower bounds
there will not by themselves prove a polystrip, but any nearby zero
will cap every possible uniform strip and prevent wasted global
subdivision.  If no nearby zero is found, start interval subdivision
of the full imaginary-direction cube, using the real-torus spectral
floor and local Lipschitz coefficient sums.
\end{decision}

\begin{assessment}[Progress after compact V50]
\label{ass:n1v50-status}
The mandatory audit is complete.  The SM companion supplies the
correct invariant-envelope theorem, now specialized and executed for
the Yang--Mills channel.  The exact legal packets are axial cosets,
and all previously tested spectator pairs cross them.  Neither
companion supplies a sharper complex inverse estimate, so the next
acquisition remains determinant localization.
\end{assessment}

\begin{firewall}[Claim boundary after compact V50]
\label{fire:n1v50-claim}
V50 proves the minimal finite-grid support packets for the declared
background shifts.  It does not evaluate packetwise response traces
at general \(M\), sharpen the complex inverse strip, populate the
nonlinear Wiener stock, certify the continuum response, close SCIC
or ADMC, construct the continuum Yang--Mills measure or clustering
theory, or prove the Yang--Mills mass gap.
\end{firewall}

\subsection{Parent record}

This V50 note was formed by appending the preceding material to the
validated compact V49, whose record was
\[
\begin{array}{c|c}
\text{lines}&19211\\
\text{bytes}&658460\\
\text{SHA--256}&
\texttt{6A8F5EEFAAD1AB422D61CA5B82F221647437844929DBA7D57B1264E087A59FD7}.
\end{array}
\]

% =============================================================================
% END APPENDED COMPACT-SERIES V50 MATERIAL
% =============================================================================

% =============================================================================
% BEGIN APPENDED COMPACT-SERIES V51 MATERIAL
% =============================================================================

\section{Compact-series V51: determinant-ray reconnaissance and an
exact endpoint Rouch\'e annulus}
\label{sec:n1v51}

\subsection{Exact Laurent restrictions}

Let
\[
 A_{\rm W}(z)=D_R(z^{-1})^TD_R(z),\qquad
 A_*(z)=A_{\rm W}(z)+3K_R(z^{-1})^TK_R(z),                 \tag{51.1}
\]
where the harmless common factor \(1/3\) in the action Hessian is
omitted.  Both matrices are generated from the exact integer Laurent
incidence arrays of V21.

For an integer direction \(d\in{\mathbb Z}^4\), restrict the Laurent
variables by
\[
                         z_\mu=t^{d_\mu}.                  \tag{51.2}
\]
The determinant becomes a one-variable Laurent polynomial.  If
\(t=e^{i\zeta}\), then
\[
 |t|=e^{-\operatorname{Im}\zeta}.                          \tag{51.3}
\]
Thus a zero-free annulus in \(t\) is a zero-free strip along the
complex momentum line \(\zeta d\).

\begin{lemma}[Polynomial linearization of a Laurent ray]
\label{lem:n1v51-ray-polynomial}
Suppose every entry of \(A(t)\) has powers between \(-r\) and \(s\).
Then
\[
 P(t):=t^rA(t)                                             \tag{51.4}
\]
is a matrix polynomial of degree at most \(r+s\), and
\[
 p(t):=\det P(t)                                           \tag{51.5}
\]
has degree at most \(45(r+s)\).  For \(t\ne0\),
\[
                         \det A(t)=0\quad\Longleftrightarrow
                         p(t)=0.                           \tag{51.6}
\]
\end{lemma}

\begin{proof}
Multiplication by \(t^r\) moves every entry exponent into
\(\{0,\ldots,r+s\}\).  The determinant is a sum of products of
forty-five entries, proving the degree bound.  Since
\(\det P=t^{45r}\det A\), (51.6) follows.
\end{proof}

\subsection{Numerical reconnaissance}

The exact coefficient matrices from (51.4) were first passed to a
generalized companion pencil.  The nearest finite candidate zeros,
ranked by \(|\log|t||\), were
\[
\begin{array}{c|c|c}
\text{action}&d&\min_{\rm candidates}|\log|t||\\ \hline
{\rm W}&e_\mu&1.71766137396387\\
{\rm W}&(1,1,1,1)&0.889320387671577\\
{\rm W}&(1,-1,1,-1)&0.889320387671572\\
*&e_\mu&2.67953504951\\
*&(1,1,1,1)&1.17024555358616\\
*&(1,-1,1,-1)&1.17024555358622.
\end{array}                                                   \tag{51.7}
\]
The normalized smallest singular residuals at these candidates range
from approximately \(10^{-18}\) to \(10^{-15}\).

\begin{warning}[The companion roots are only locators]
\label{warn:n1v51-diagnostic}
Equation (51.7) is a double-precision generalized-eigenvalue scan.
It neither proves that the displayed candidates are exact zeros nor
excludes missed zeros.  It is used only to choose rational contours
for the exact calculation below.
\end{warning}

The scan nevertheless shows that the tiny Neumann width in V49 is
not visibly saturated on these rays: all located candidates are at
distance of order one.

\subsection{Exact endpoint determinant on an axis}

Now take \(A_*(t,1,1,1)\).  Its entries have powers
\(-1,0,1\), so
\[
 P_*(t):=tA_*(t,1,1,1),\qquad
 p_*(t):=\det P_*(t)                                      \tag{51.8}
\]
has degree at most \(90\).

The verifier evaluates \(\det P_*(n)\) for
\(n=0,\ldots,90\) by fraction-free Bareiss elimination and performs
exact Newton interpolation.  Six further determinant evaluations at
\[
                         n=-3,-2,-1,91,92,93               \tag{51.9}
\]
agree exactly with the interpolated polynomial.

\begin{proposition}[Exact endpoint axis polynomial record]
\label{prop:n1v51-polynomial}
Writing
\[
                         p_*(t)=\sum_{j=0}^{90}c_jt^j,      \tag{51.10}
\]
the exact coefficient array obeys:
\[
\begin{gathered}
 c_j\ne0\quad\Longleftrightarrow\quad7\le j\le83,\\
 c_j=c_{90-j}\quad(0\le j\le90),                           \tag{51.11}\\
 c_{45}=
 16721051406822941106198296923634702891961855399922706941192791138078496605317931713662272.
\end{gathered}
\]
Its ordered coefficient SHA--256 is
\[
 \texttt{356AC471852DB1BDB3162F94E410733C1181DA9C1A31A5B41EAAB343F056C4A9}.
                                                                    \tag{51.12}
\]
\end{proposition}

\begin{proof}
The degree bound preceding (51.8) makes ninety-one distinct exact
values sufficient to determine \(p_*\).  Bareiss elimination uses
exact divisions in \({\mathbb Z}\); Newton interpolation is carried
out in \({\mathbb Q}\), and every final coefficient is an integer.
Direct inspection of the resulting array proves (51.11), while the
six points in (51.9) are independent interpolation checks.
\end{proof}

\subsection{A rigorous zero-free annulus}

The exact coefficients satisfy the two strict inequalities
\[
 |c_{45}|2^{-45}>
 \sum_{j\ne45}|c_j|2^{-j},                                 \tag{51.13}
\]
\[
 |c_{45}|2^{45}>
 \sum_{j\ne45}|c_j|2^j.                                   \tag{51.14}
\]
After clearing powers of two, both positive integer margins have
bit length \(333\).

\begin{theorem}[Endpoint axis annulus]
\label{thm:n1v51-annulus}
The endpoint determinant has no zero on
\[
 \boxed{
 \frac12\le|t|\le2.}                                       \tag{51.15}
\]
Equivalently,
\[
 \boxed{
 \det A_*(e^{i\zeta},1,1,1)\ne0
 \quad\hbox{whenever}\quad
 |\operatorname{Im}\zeta|\le\log2.}                        \tag{51.16}
\]
\end{theorem}

\begin{proof}
On \(|t|=1/2\), inequality (51.13) says that the monomial
\(c_{45}t^{45}\) strictly dominates the sum of all other terms in
\(p_*\).  Rouch\'e's theorem implies that \(p_*\) and this monomial
have the same number, namely forty-five, of zeros inside the circle.
The same argument using (51.14) gives forty-five zeros inside
\(|t|=2\).  Strict dominance also excludes boundary zeros.  Since
the two counts agree, there is no zero in the closed annulus (51.15).

For \(t=e^{i\zeta}\), (51.3) maps
\(|\operatorname{Im}\zeta|\le\log2\) exactly to (51.15).
Equation (51.6) then proves (51.16).
\end{proof}

The rigorous one-axis width \(\log2\) is more than \(1.4\) million
times \(\rho_{21}\).  This does not imply a comparable four-variable
polystrip, but it demonstrates concretely how much sharpness was
lost by the global coefficient-sum perturbation.

\begin{firewall}[One line is not a polystrip]
\label{fire:n1v51-not-polystrip}
Theorem~\ref{thm:n1v51-annulus} fixes the other three Laurent
variables at one.  A zero with several independently complex
coordinates is not excluded.  Nor does (51.16) give a uniform
inverse norm near its boundary.  It therefore cannot replace the
four-dimensional strip required by the V17 quadrature theorem.
\end{firewall}

\subsection{Executable records and next acquisition}

The diagnostic locator is
\[
 \texttt{scripts/explore\_ym\_v51\_determinant\_rays.py},   \tag{51.17}
\]
with
\[
\begin{array}{c|c}
\text{lines}&145\\
\text{bytes}&5097\\
\text{SHA--256}&
\texttt{F2BCB2385017A322354F1C72EBD1E435D6E1D85C1E462049A7A62FB386544885}.
\end{array}                                                  \tag{51.18}
\]
Its diagnostic payload SHA--256 is
\[
 \texttt{5E582C5DAE4FBD946FC2702E70BCF20B88C9E899C5E80B5E3F18792CC11931BF}.
                                                                    \tag{51.19}
\]

The exact certificate is
\[
 \texttt{scripts/verify\_ym\_v51\_endpoint\_axis\_annulus.py},       \tag{51.20}
\]
with
\[
\begin{array}{c|c}
\text{lines}&188\\
\text{bytes}&6852\\
\text{SHA--256}&
\texttt{2FB1356A1A6F4B8CE65EEEA4EEBB07EE2255FC5AA5AD36EE3278FA9CAF7F0A89}.
\end{array}                                                  \tag{51.21}
\]
Its canonical payload SHA--256 is
\[
 \boxed{
 \texttt{1CAA6C5BE2691D734D179B12F37D143A2EC78493FB6476B9A3D496E4B80E2B48}.}
                                                                    \tag{51.22}
\]

\begin{decision}[V52 exact Wilson ray and multivariable transition]
\label{dec:n1v51-next}
Apply exact interpolation to the Wilson axis and diagonal
restrictions.  Since single-monomial dominance fails on
\(|t|=1/2,2\), use reciprocal reduction or an exact Schur--Cohn
count to certify a Wilson annulus.  Then formulate a multivariable
boundary subdivision whose cells are excluded by an interval
smallest-singular-value or inverse-residual test.  The one-ray
certificates are reconnaissance and regression faces for that
polystrip algorithm.
\end{decision}

\begin{assessment}[Progress after compact V51]
\label{ass:n1v51-status}
The determinant bottleneck has moved from an extremely crude global
Neumann radius to exact one-variable geometry.  The endpoint axis
restriction is rigorously zero-free through imaginary width
\(\log2\), while numerical ray scans locate no closer Wilson or
endpoint candidates on six tested directions.  A multivariable
certificate remains necessary.
\end{assessment}

\begin{firewall}[Claim boundary after compact V51]
\label{fire:n1v51-claim}
V51 proves only the endpoint one-axis annulus (51.15).  It does not
certify the Wilson ray, a four-dimensional polystrip, a response
Wiener norm, the continuum response sign, SCIC or ADMC, a continuum
Yang--Mills measure, exponential clustering, or the Yang--Mills mass
gap.
\end{firewall}

\subsection{Parent record}

This V51 note was formed by appending the preceding material to the
validated compact V50, whose record was
\[
\begin{array}{c|c}
\text{lines}&19492\\
\text{bytes}&668673\\
\text{SHA--256}&
\texttt{CBB7D7F437410347B9C685FA48EF6F397C6AF47587ECC46EED2284FEF4C64C1A}.
\end{array}
\]

% =============================================================================
% END APPENDED COMPACT-SERIES V51 MATERIAL
% =============================================================================

% =============================================================================
% BEGIN APPENDED COMPACT-SERIES V52 MATERIAL
% =============================================================================

\section{Compact-series V52: exact Wilson reciprocal factorization
and the common axis annulus}
\label{sec:n1v52}

\subsection{The Wilson determinant polynomial}

Retain the axis restriction of V51 and put
\[
 P_{\rm W}(t):=tA_{\rm W}(t,1,1,1),\qquad
 p_{\rm W}(t):=\det P_{\rm W}(t).                          \tag{52.1}
\]
Each entry of \(P_{\rm W}\) has degree at most two, so
\(\deg p_{\rm W}\le90\).  Exact Bareiss evaluation at
\(t=0,\ldots,90\), followed by Newton interpolation, gives
\[
 p_{\rm W}(t)=\sum_{j=0}^{90}a_jt^j,                       \tag{52.2}
\]
with
\[
 a_j\ne0\quad\Longleftrightarrow\quad31\le j\le59,         \tag{52.3}
\]
\[
                         a_j=a_{90-j}.                     \tag{52.4}
\]
The integer content of the nonzero coefficient array is \(1728\).
Its SHA--256 is
\[
 \texttt{9312AC0F121175FA79C77CFAF35F399351DD68A993D70E8547DA6D354D434F04}.
                                                                    \tag{52.5}
\]
As in V51, six exact evaluations at
\(-3,-2,-1,91,92,93\) agree with the interpolant.

\subsection{Reciprocal reduction}

Define the monic reciprocal core \(Q\) by
\[
                         p_{\rm W}(t)=1728\,t^{31}Q(t),
 \qquad \deg Q=28.                                         \tag{52.6}
\]
Equation (52.4) makes \(Q\) palindromic.  Therefore there is a unique
monic polynomial \(R\) of degree fourteen such that
\[
                         Q(t)=t^{14}R(t+t^{-1}).            \tag{52.7}
\]

\begin{lemma}[Exact reciprocal conversion]
\label{lem:n1v52-reciprocal}
Let
\[
 S_0(u)=2,\qquad S_1(u)=u,\qquad
 S_n(u)=uS_{n-1}(u)-S_{n-2}(u).                            \tag{52.8}
\]
Then
\[
                         S_n(t+t^{-1})=t^n+t^{-n}.         \tag{52.9}
\]
If
\[
 Q(t)=q_{14}t^{14}
 +\sum_{r=1}^{14}q_{14+r}
   \left(t^{14+r}+t^{14-r}\right),                         \tag{52.10}
\]
then
\[
                         R(u)=q_{14}
 +\sum_{r=1}^{14}q_{14+r}S_r(u).                           \tag{52.11}
\]
\end{lemma}

\begin{proof}
Equation (52.9) holds for \(n=0,1\).  Multiplication by
\(t+t^{-1}\) and subtraction of the preceding identity proves the
recurrence step.  Substitute (52.9) into (52.11), multiply by
\(t^{14}\), and use the reciprocal pairing in (52.10).
\end{proof}

\begin{theorem}[Exact Wilson axis factorization]
\label{thm:n1v52-factorization}
The reciprocal polynomial in (52.7) factors over
\({\mathbb Z}\) as
\[
\boxed{
 R(u)=
 (u-34)^3(u-98)^3(u-194)^2
 (u^2-76u+404)^3.}                                        \tag{52.12}
\]
Equivalently,
\[
\boxed{
 p_{\rm W}(t)=
 1728\,t^{45}
 (u-34)^3(u-98)^3(u-194)^2
 (u^2-76u+404)^3,\quad
 u=t+t^{-1}.}                                              \tag{52.13}
\]
\end{theorem}

\begin{proof}
Apply Lemma~\ref{lem:n1v52-reciprocal} to the exact coefficient
array (52.2).  Direct integer multiplication of the four factors in
(52.12) reproduces all fifteen coefficients of the resulting
degree-fourteen \(R\).  The ordered coefficient array has SHA--256
\[
 \texttt{BE1030898B3C89F80EFC135700C1BEA1287ED47F61AC3C092EEA22266B1E260B}.
                                                                    \tag{52.14}
\]
Equations (52.6)--(52.7) then give (52.13).
\end{proof}

\subsection{The Wilson annulus}

The roots of the last factor in (52.12) are
\[
                         u_\pm=38\pm4\sqrt{65}.             \tag{52.15}
\]
Since
\[
                         16\cdot65=1040<1089=33^2,
\]
one has
\[
                         u_-=38-4\sqrt{65}>5.               \tag{52.16}
\]
Every other \(u\)-root in (52.12) is at least \(34\).

\begin{theorem}[Wilson axis annulus]
\label{thm:n1v52-Wilson-annulus}
\[
 \boxed{
 \det A_{\rm W}(t,1,1,1)\ne0
 \quad\hbox{for}\quad
 \frac12\le|t|\le2.}                                       \tag{52.17}
\]
Equivalently,
\[
 \boxed{
 \det A_{\rm W}(e^{i\zeta},1,1,1)\ne0
 \quad\hbox{whenever}\quad
 |\operatorname{Im}\zeta|\le\log2.}                        \tag{52.18}
\]
\end{theorem}

\begin{proof}
Every nonzero root of (52.13) satisfies
\[
                         t+t^{-1}=u_0                       \tag{52.19}
\]
for a real root \(u_0>5\) of \(R\).  The discriminant
\(u_0^2-4\) is positive, so both solutions of (52.19) are real
and positive.  If one lay in \([1/2,2]\), then
\[
                         t+t^{-1}\le\frac52,
\]
contradicting \(u_0>5\).  Hence the two reciprocal solutions lie
respectively below \(1/2\) and above \(2\).  The factor \(t^{45}\)
has only the excluded root \(t=0\).  This proves (52.17).
Equation (52.18) follows from \(|e^{i\zeta}|=
e^{-\operatorname{Im}\zeta}\).
\end{proof}

\begin{corollary}[Common exact one-axis strip]
\label{cor:n1v52-common-axis}
For both \(e={\rm W}\) and \(e=*\),
\[
 \boxed{
 \det A_e(e^{i\zeta},1,1,1)\ne0
 \quad\text{on}\quad
 |\operatorname{Im}\zeta|\le\log2.}                        \tag{52.20}
\]
\end{corollary}

\begin{proof}
Use Theorem~\ref{thm:n1v52-Wilson-annulus} for the Wilson action and
Theorem~\ref{thm:n1v51-annulus} for the endpoint action.
\end{proof}

\subsection{What the exact factorization changes}

The old common strip \(\rho_{21}\) was obtained by replacing the
complete multivariable matrix geometry by one global sum of absolute
coefficients.  Equations (51.16) and (52.20) show that no
one-coordinate determinant zero forces that small radius: on the
coordinate line through the toron, both determinants remain nonzero
through width \(\log2\).

This result also removes the need for an exact Schur--Cohn
calculation on the Wilson axis.  The failed single-monomial
dominance test noted in V51 did not signal a zero; reciprocal
factorization exposes the complete root set algebraically.

\begin{warning}[The multivariable obstruction remains]
\label{warn:n1v52-multivariable}
A polynomial can be nonzero on every coordinate line and still have
zeros when two or more variables move simultaneously.  Corollary
\ref{cor:n1v52-common-axis} therefore supplies boundary regression
faces and a scale for subdivision, not a four-dimensional polystrip.
\end{warning}

\subsection{Executable record and next acquisition}

\begin{record}[V52 exact Wilson factor certificate]
\label{rec:n1v52-executable}
The verifier is
\[
 \texttt{scripts/verify\_ym\_v52\_wilson\_axis\_factor.py}. \tag{52.21}
\]
Its source record is
\[
\begin{array}{c|c}
\text{lines}&137\\
\text{bytes}&4860\\
\text{SHA--256}&
\texttt{D7A4F89D60453FDB236878555764D81F640FAC712A63A233F9DB6D1CB5849061}.
\end{array}                                                  \tag{52.22}
\]
The canonical payload SHA--256 is
\[
 \boxed{
 \texttt{C1982C27088CBBCA2BCAA3F4CF862E10A892901D889AA4F0A8B645DEE8C65535}.}
                                                                    \tag{52.23}
\]
The program checks exact interpolation, reciprocal symmetry,
integer content, the complete factorization (52.12), and the
rational inequality proving (52.16).
\end{record}

\begin{decision}[V53 multivariable inverse enclosure]
\label{dec:n1v52-next}
Construct a certified polystrip by covering the real torus with
boxes and perturbing from the inverse at each real box centre, rather
than perturbing every point from one global coefficient sum.  For a
candidate width \(r\), bound
\[
 \|A_e(k+iy)-A_e(k_c)\|
 \le
 \|A_e(k)-A_e(k_c)\|
 +\|A_e(k+iy)-A_e(k)\|                                    \tag{52.24}
\]
with entrywise interval Laurent evaluation.  Accept a box only when
an interval inverse residual or smallest-singular-value lower bound
stays positive for both actions.  Use (52.20) as a regression on
axis-aligned faces; bisect any unresolved box.
\end{decision}

\begin{assessment}[Progress after compact V52]
\label{ass:n1v52-status}
Both exact parity-tree Hessians are now proved zero-free through
imaginary width \(\log2\) on a full complex coordinate line.  The
Wilson proof is stronger than reconnaissance: its determinant is
factored completely after reciprocal reduction.  The remaining
analytic bottleneck is genuinely multivariable.
\end{assessment}

\begin{firewall}[Claim boundary after compact V52]
\label{fire:n1v52-claim}
V52 proves the Wilson axis factorization and common one-axis strip.
It does not certify a multivariable polystrip or inverse norm,
populate the nonlinear response Wiener stock, determine the
continuum response, close SCIC or ADMC, construct the continuum
Yang--Mills measure or clustering theory, or prove the Yang--Mills
mass gap.
\end{firewall}

\subsection{Parent record}

This V52 note was formed by appending the preceding material to the
validated compact V51, whose record was
\[
\begin{array}{c|c}
\text{lines}&19776\\
\text{bytes}&678427\\
\text{SHA--256}&
\texttt{3A34FF1D5949A1083D46E894E8BC343C90AEF5223215C5F66AFC85858775198C}.
\end{array}
\]

% =============================================================================
% END APPENDED COMPACT-SERIES V52 MATERIAL
% =============================================================================

% =============================================================================
% BEGIN APPENDED COMPACT-SERIES V53 MATERIAL
% =============================================================================

\section{Compact-series V53: induced Laurent stocks and a wider
four-variable inverse strip}
\label{sec:n1v53}

\subsection{Why the V52 box subdivision can be postponed}

The V52 decision proposed interval subdivision because the V21
perturbation estimate collapsed the entire \(45\)-by-\(45\) Laurent
matrix into the scalar stock

\[
 \sum_{i,j,\alpha}|[z^\alpha]A_{t,ij}|.                 \tag{53.1}
\]

That scalar is not adapted to the operator norm: it adds the stocks
of all forty-five rows even though the induced infinity norm takes
only the largest row.  This note first moves upstream and computes
the induced row and column stocks exactly.  The resulting global
certificate is already substantially sharper, so no interval boxes
are needed for the present gain.

Use the unscaled family

\[
 B_t:=3A_t=D_R^*D_R+3tK_R^*K_R,
 \qquad 0\le t\le1.                                      \tag{53.2}
\]

The real-torus floor from V18--V21 becomes

\[
 \sigma_{\min}(B_t(k))\ge 3\underline\lambda
 =\frac{147}{1250}                                      \tag{53.3}
\]

for every real \(k\in{\mathbb R}^4\) and \(0\le t\le1\).

\subsection{Induced coefficient stocks}

If \(P(z)=\sum_\alpha C_\alpha z^\alpha\) is a matrix Laurent
polynomial, define

\[
 \begin{split}
  {\cal R}(P)&:=\max_i\sum_{j,\alpha}
                 |(C_\alpha)_{ij}|,\\
  {\cal C}(P)&:=\max_j\sum_{i,\alpha}
                 |(C_\alpha)_{ij}|,\\
  R(P)&:=\max\{|\alpha|_1:C_\alpha\ne0\}.
 \end{split}                                             \tag{53.4}
\]

\begin{lemma}[Endpoint convexity of every induced stock]
\label{lem:n1v53-convexity}
Write \(B_t=B_0+t(B_1-B_0)\).  For each fixed row \(i\),

\[
 r_i(t):=\sum_{j,\alpha}
 |[z^\alpha](B_t)_{ij}|                                  \tag{53.5}
\]

is convex on \([0,1]\), and hence

\[
 \sup_{0\le t\le1}r_i(t)=\max\{r_i(0),r_i(1)\}.          \tag{53.6}
\]

The analogous assertion holds for every column stock.
\end{lemma}

\begin{proof}
Every summand in (53.5) is the absolute value of an affine real
function of \(t\), hence is convex.  A finite sum of convex functions
is convex.  If \(t=(1-t)0+t1\), convexity gives

\[
 r_i(t)\le(1-t)r_i(0)+tr_i(1)
 \le\max\{r_i(0),r_i(1)\}.                               \tag{53.7}
\]

The reverse inequality follows by evaluating at the endpoints.  The
column proof is identical.
\end{proof}

\begin{proposition}[Exact induced-stock inventory]
\label{prop:n1v53-inventory}
For the literal forty-five reduced columns of the parity-cell
compiler,

\[
\begin{array}{c|ccc}
P&{\cal R}(P)&{\cal C}(P)&R(P)\\ \hline
B_0&23&23&2\\
B_1&1066&1066&3.
\end{array}                                               \tag{53.8}
\]

The maximizing zero-based rows and columns are respectively

\[
 \{33,37,41\}\quad\hbox{for }B_0,
 \qquad
 \{40,43,44\}\quad\hbox{for }B_1.                       \tag{53.9}
\]

Moreover \(R(B_1-B_0)=3\), and for every \(0\le t\le1\),

\[
 \boxed{{\cal R}(B_t)\le1066,
        \qquad {\cal C}(B_t)\le1066,
        \qquad R(B_t)\le3.}                              \tag{53.10}
\]
\end{proposition}

\begin{proof}
Generate \(D_R\) and \(K_R\) from the incidence rules of V21, form
the two exact integer Gram matrices in (53.2), and sum the absolute
integer Laurent coefficients row by row and column by column.  This
finite computation gives (53.8)--(53.9) and
\(R(B_1-B_0)=3\).  Lemma~\ref{lem:n1v53-convexity} gives (53.10).
The complete executable record appears below.
\end{proof}

\subsection{The four-variable certificate}

Put

\[
 \boxed{
 \rho_{53}:=\frac13\log\left(1+\frac{147}{2665000}\right).}
                                                               \tag{53.11}
\]

Numerically,

\[
 \rho_{53}=
 0.0000183859844812621988214609067334683190984207\ldots . \tag{53.12}
\]

\begin{lemma}[Induced Laurent perturbation]
\label{lem:n1v53-perturbation}
For real \(k,y\in{\mathbb R}^4\), \(0\le t\le1\), and
\(\|y\|_\infty\le r\),

\[
 \|B_t(k+iy)-B_t(k)\|_2
 \le1066\bigl(e^{3r}-1\bigr).                            \tag{53.13}
\]
\end{lemma}

\begin{proof}
Write \(B_t(z)=\sum_\alpha C_\alpha(t)e^{i\alpha\cdot z}\).
For real \(k,y\),

\[
 \begin{split}
 B_t(k+iy)-B_t(k)
  =\sum_\alpha C_\alpha(t)e^{i\alpha\cdot k}
       \bigl(e^{-\alpha\cdot y}-1\bigr).                 \tag{53.14}
 \end{split}
\]

Since \(|\alpha|_1\le3\),

\[
 |e^{-\alpha\cdot y}-1|
 \le e^{|\alpha\cdot y|}-1
 \le e^{3r}-1.                                           \tag{53.15}
\]

Equation (53.10) therefore bounds both the induced infinity norm
and induced one norm of (53.14) by
\(1066(e^{3r}-1)\).  The standard inequality

\[
 \|M\|_2\le\sqrt{\|M\|_1\|M\|_\infty}                  \tag{53.16}
\]

proves (53.13).
\end{proof}

\begin{theorem}[Uniform four-variable inverse strip]
\label{thm:n1v53-polystrip}
For every real \(k,y\in{\mathbb R}^4\), \(0\le t\le1\), and

\[
                         \|y\|_\infty\le\rho_{53},       \tag{53.17}
\]

the complex matrix \(A_t(k+iy)\) is invertible and

\[
 \boxed{\|A_t(k+iy)^{-1}\|_2\le\frac{2500}{49}.}         \tag{53.18}
\]
\end{theorem}

\begin{proof}
The definition (53.11) gives exactly

\[
 1066(e^{3\rho_{53}}-1)=\frac{147}{2500}.                 \tag{53.19}
\]

By (53.3), Lemma~\ref{lem:n1v53-perturbation}, and the
one-Lipschitz property of the smallest singular value,

\[
 \begin{split}
 \sigma_{\min}(B_t(k+iy))
 &\ge \sigma_{\min}(B_t(k))
      -\|B_t(k+iy)-B_t(k)\|_2\\
 &\ge\frac{147}{1250}-\frac{147}{2500}
   =\frac{147}{2500}.                                    \tag{53.20}
 \end{split}
\]

Thus \(B_t(k+iy)\) is invertible and
\(\|B_t(k+iy)^{-1}\|_2\le2500/147\).  Since \(A_t=B_t/3\),
equation (53.18) follows.
\end{proof}

\begin{corollary}[Quantified gain and Wilson-only strip]
\label{cor:n1v53-gain}
The common radius (53.11) is

\[
 \frac{\rho_{53}}{\rho_{21}}
 =39.5767972353199162\ldots .                            \tag{53.21}
\]

At the Wilson endpoint \(t=0\), the same proof with the first row of
(53.8) gives the larger four-variable radius

\[
 \boxed{
 \rho_{53}^{\rm W}:=
 \frac12\log\left(1+\frac{147}{57500}\right)
 =0.0012766296982064426680\ldots .}                       \tag{53.22}
\]
\end{corollary}

\begin{proof}
The first assertion is direct evaluation of (53.11) and (21.21).
For \(t=0\), replace \(1066\) by \(23\) and the Laurent range three
by two in Lemma~\ref{lem:n1v53-perturbation}.  The perturbation at
(53.22) is again (147/2500), so the proof of
Theorem~\ref{thm:n1v53-polystrip} applies verbatim.
\end{proof}

\subsection{Executable record and consequence}

\begin{record}[V53 induced-stock certificate]
\label{rec:n1v53-executable}
The verifier is

\[
 \texttt{scripts/verify\_ym\_v53\_row\_stock\_polystrip.py}. \tag{53.23}
\]

Its source record is

\[
\begin{array}{c|c}
\text{lines}&154\\
\text{bytes}&5521\\
\text{SHA--256}&
\texttt{DD22CBE275FB99E2B71074048F3F8E608C0C3738689DF20C741DFFB4435C9849}.
\end{array}                                                \tag{53.24}
\]

The canonical payload SHA--256 is

\[
 \boxed{
 \texttt{DC5CD0E1A1A92F307323D35A7D9B75E2D5015A5EE6B437C4C78947BD600DEC72}.}
                                                               \tag{53.25}
\]

The program independently records every maximizing row and column,
the Laurent ranges, the exact radius formula, the old-to-new ratio,
and the Wilson-only radius.
\end{record}

Theorem~\ref{thm:n1v53-polystrip} completes the multivariable task
left open in V51--V52 at a useful first scale.  It does not make the
axis annuli redundant: their width \(\log2\) still shows that this
global norm estimate is far from determinant-optimal.  But the
relevant next question is now quantitative rather than logical: is
the new common radius large enough for the response alias and Wiener
tails used later in the programme?

\begin{decision}[V54 propagate the certified gain]
\label{dec:n1v53-next}
Replace \(\rho_{21}\) by \(\rho_{53}\) in the downstream Cauchy and
Wiener allocations, keeping every loss factor explicit.  Recompute
the \(M=2\) alias multiplier and the first necessary mesh threshold.
If the response stock still cannot close, identify whether the
dominant loss is the common endpoint Hessian, the repeated Cauchy
shrinkage, or the nonlinear numerator stock before attempting finer
determinant boxes.
\end{decision}

\begin{assessment}[Progress after compact V53]
\label{ass:n1v53-status}
The common Hessian family is now rigorously invertible on a genuine
four-complex-variable polystrip almost forty times wider than the V21
certificate.  This resolves the logical multivariable gap flagged in
V52, while leaving a large quantitative gap between the global
polystrip and the exact one-axis annuli.
\end{assessment}

\begin{firewall}[Claim boundary after compact V53]
\label{fire:n1v53-claim}
V53 proves a wider multivariable inverse strip for the finite
parity-cell Hessian family.  It does not yet bound the complete
nonlinear response Wiener norm, determine the continuum response,
close SCIC or ADMC, construct the continuum Yang--Mills measure or
exponential clustering, or prove the Yang--Mills mass gap.
\end{firewall}

\subsection{Parent record}

This V53 note was formed by appending the preceding material to the
validated compact V52, whose record was

\[
\begin{array}{c|c}
\text{lines}&20061\\
\text{bytes}&687538\\
\text{SHA--256}&
\texttt{301F7BB8BD639CE2B0317437CD35ED0588267CD6B91F1A1FDB539250A2243108}.
\end{array}
\]

% =============================================================================
% END APPENDED COMPACT-SERIES V53 MATERIAL
% =============================================================================

% =============================================================================
% BEGIN APPENDED COMPACT-SERIES V54 MATERIAL
% =============================================================================

\section{Compact-series V54: propagation of the polystrip gain and
an exact obstruction to the absolute \(M=2\) alias route}
\label{sec:n1v54}

\subsection{The V53 radius in the existing Wiener allocation}

The downstream allocation fixed in V49 assigns one eighth of the
common Hessian radius to the internal-momentum Wiener strip.  With
the improved V53 certificate this gives

\[
 \boxed{
 \rho_W^{(53)}:=\frac{\rho_{53}}8
 =0.0000022982480601577748526826133416835398873\ldots .} \tag{54.1}
\]

For a four-dimensional \(M\)-mesh, the absolute alias lemma of V49
has multiplier

\[
 C_M(\rho):=
 \coth^4\left(\frac{\rho M}{2}\right)-1.                 \tag{54.2}
\]

Put

\[
 x_*:=\operatorname{arctanh}(2^{-1/4})
 =1.2242262238390378950026549568179345701\ldots .         \tag{54.3}
\]

\begin{lemma}[Exact threshold for the absolute alias multiplier]
\label{lem:n1v54-threshold}
For \(\rho>0\),

\[
 \boxed{C_M(\rho)<1
 \quad\Longleftrightarrow\quad
 \frac{\rho M}{2}>x_*.}                                  \tag{54.4}
\]
\end{lemma}

\begin{proof}
On the positive real axis, \(\tanh\) is strictly increasing and
\(\coth x=1/\tanh x\).  Thus

\[
 \begin{split}
 C_M(\rho)<1
 &\Longleftrightarrow
 \coth^4(\rho M/2)<2\\
 &\Longleftrightarrow
 \tanh(\rho M/2)>2^{-1/4}\\
 &\Longleftrightarrow
 \rho M/2>\operatorname{arctanh}(2^{-1/4}).
 \end{split}                                              \tag{54.5}
\]
\end{proof}

\begin{proposition}[The improved common strip still does not resolve
the mesh]
\label{prop:n1v54-mesh}
With \(\rho=\rho_W^{(53)}\), the elementary implications

\[
 \log(1+x)<x\quad(x>0),qquad x_*>\frac12                 \tag{54.6}
\]

give the rigorous necessary condition

\[
 \boxed{C_M(\rho_W^{(53)})<1
 \quad\Longrightarrow\quad
 M>\frac{21320000}{49}.}                                  \tag{54.7}
\]

Consequently the smallest integer not excluded by this elementary
bound is \(435103\).  Direct high-precision evaluation of the exact
threshold in (54.4) gives

\[
 M>1065356.03798572961617\ldots,                           \tag{54.8}
\]

so the first integer satisfying the scalar inequality is
\(1065357\).  At \(M=2\),

\[
 C_2(\rho_W^{(53)})
 =3.5843663466895938347458\ldots\,10^{22}.                \tag{54.9}
\]
\end{proposition}

\begin{proof}
From (53.11) and \(\log(1+x)<x\),

\[
 \rho_W^{(53)}<\frac{49}{21320000}.                       \tag{54.10}
\]

If \(C_M<1\), Lemma~\ref{lem:n1v54-threshold} and \(x_*>1/2\)
give \(M>1/\rho_W^{(53)}\), proving (54.7).  Substitution in
(54.2)--(54.4) gives (54.8)--(54.9).
\end{proof}

\subsection{Why a still sharper determinant strip cannot fix \(M=2\)}

The preceding failure could in principle have been blamed on the
remaining looseness of the V53 norm bound.  The exact Wilson
factorization of V52 rules that out for this particular route.

Let

\[
 u_0:=38-4\sqrt{65},qquad
 t_0:=\frac{u_0+\sqrt{u_0^2-4}}2,qquad
 \delta_{\rm W}:=\log t_0.                               \tag{54.11}
\]

Theorem~\ref{thm:n1v52-factorization} gives

\[
 \det A_{\rm W}(t_0,1,1,1)=0,
 \qquad
 \delta_{\rm W}=1.71766137396387420918\ldots .           \tag{54.12}
\]

\begin{theorem}[Incompatibility of the absolute alias lemma with
\(M=2\)]
\label{thm:n1v54-incompatibility}
Under the V49 allocation \(\rho_W=\rho_A/8\), no valid common
Hessian polystrip radius \(\rho_A\) can make the absolute alias
multiplier \(C_2(\rho_W)\) smaller than one.
\end{theorem}

\begin{proof}
First, \(64<65\) implies \(u_0<6\).  The function
\(t+t^{-1}\) is strictly increasing for \(t>1\), and
\(6+1/6>6>u_0\), so \(t_0<6\).  Also

\[
 e^2>1+2+\frac{2^2}{2!}+\frac{2^3}{3!}>6,                 \tag{54.13}
\]

whence

\[
                         \delta_{\rm W}<\log6<2.          \tag{54.14}
\]

Any common polystrip contains the Wilson coordinate line.  The zero
(54.12) therefore forces every valid common radius to satisfy
\(\rho_A<\delta_{\rm W}<2\).

On the other hand, Lemma~\ref{lem:n1v54-threshold} at \(M=2\)
requires

\[
 \rho_W>x_*,qquadhbox{hence}qquad
 \rho_A>8x_*>4.                                          \tag{54.15}
\]

The two strict inequalities are incompatible.
\end{proof}

\begin{remark}[The obstruction is method-specific]
\label{rem:n1v54-method-specific}
Theorem~\ref{thm:n1v54-incompatibility} does not say that the actual
signed \(M=2\) response differs greatly from its continuum integral.
It says that the universal absolute Fourier-tail bound (54.2), after
the one-eighth strip allocation, cannot prove otherwise.  The bound
throws away congruence-class cancellation before it sees the
specific response integrand.
\end{remark}

For comparison, even the Wilson-only V53 radius (53.22), passed
through the same one-eighth allocation, first makes \(C_M<1\) at
\(M=15344\), still not at \(M=2\).  Thus further box subdivision of
the Hessian determinant is not the next efficient move.

\subsection{Executable record and route change}

\begin{record}[V54 alias-obstruction certificate]
\label{rec:n1v54-executable}
The verifier is

\[
 \texttt{scripts/verify\_ym\_v54\_alias\_obstruction.py}. \tag{54.16}
\]

Its source record is

\[
\begin{array}{c|c}
\text{lines}&98\\
\text{bytes}&4243\\
\text{SHA--256}&
\texttt{2D7EEA89EC22218BA5B6A55ACDF7A34EAACCF6AC699A2EB41BB378668A296581}.
\end{array}                                                \tag{54.17}
\]

The canonical payload SHA--256 is

\[
 \boxed{
 \texttt{06B5263B9A7CFA300F562160B5BA7416E41F758B10EB8C1255988ECC10930B69}.}
                                                               \tag{54.18}
\]
\end{record}

\begin{decision}[V55 mandatory cross-program audit and signed alias
route]
\label{dec:n1v54-next}
At the mandatory multiple-of-five checkpoint, inspect the current SM
and NS notes for a cancellation-preserving trace, packet, or
Poisson-congruence device.  Then replace (54.2) by a response-specific
identity: group Fourier coefficients by residue class modulo two,
retain their signs and tensor components, and estimate only the
unresolved tail.  The exact \(M=2\) parity packets identified in V50
are the natural finite skeleton.
\end{decision}

\begin{assessment}[Progress after compact V54]
\label{ass:n1v54-status}
The forty-fold strip improvement has been propagated exactly.  More
importantly, the Wilson axis zero proves that no further refinement
of the common Hessian strip can make the existing absolute alias
lemma certify \(M=2\).  The programme must preserve signed packet
cancellation rather than enlarge the same scalar majorant.
\end{assessment}

\begin{firewall}[Claim boundary after compact V54]
\label{fire:n1v54-claim}
V54 proves a no-go theorem for one particular absolute-value alias
argument.  It does not determine the signed continuum response or
its \(M=2\) discrepancy, close SCIC or ADMC, construct the continuum
Yang--Mills measure or exponential clustering, or prove the
Yang--Mills mass gap.
\end{firewall}

\subsection{Parent record}

This V54 note was formed by appending the preceding material to the
validated compact V53, whose record was

\[
\begin{array}{c|c}
\text{lines}&20404\\
\text{bytes}&697280\\
\text{SHA--256}&
\texttt{F225EE84DD841B613F8CD65CFE21A21CAD927880F28C8940143DA54A29A0E0C9}.
\end{array}
\]

% =============================================================================
% END APPENDED COMPACT-SERIES V54 MATERIAL
% =============================================================================

% =============================================================================
% BEGIN APPENDED COMPACT-SERIES V55 MATERIAL
% =============================================================================

\section{Compact-series V55: mandatory companion audit and a
signed packet-first Poisson decomposition}
\label{sec:n1v55}

\subsection{Mandatory companion audit}

The current companion notes inspected read-only at this
multiple-of-five checkpoint are

\[
\begin{array}{c|r|r|c}
\text{record}&\text{lines}&\text{bytes}&\text{SHA--256}\\ \hline
\text{SM--Einstein V38}&11165&380176&
\texttt{DC36C62B7D172E146478A9F015E6E2AC405231FF8DE056D7BAE3CBE01BF3D787}\\
\text{Navier--Stokes V26}&5319&278283&
\texttt{82C887F8C2C69E4F28FAD7C3DC8DE5B9099CB5A4BF431EBA1FFF3E91935978AD}.
\end{array}                                                   \tag{55.1}
\]

SM V25--V26 remains the relevant algebraic input.  It proves that a
background-shift calculation belongs to the trace over a coset of
the generated shift subgroup, and that all intermediate cards of a
closed shift walk must remain inside that trace.  YM V50 already
specialized this theorem, so it is not repeated here.

SM V35 independently records that exact coordinate-line determinant
certificates do not tensorize into a multivariable domain.  V53 has
since supplied the missing genuine polystrip by induced coefficient
stocks.  SM V36--V38 develops positive-kernel normalization and
bad-set moment bounds; those results do not provide a signed
Yang--Mills alias estimate.

NS V26 proves a structurally analogous warning in a different
problem: a smaller rank-one kernel norm is usable only when the
actual input is proved to remain in the rank-one cone.  Transferred
as a proof rule rather than a numerical estimate, this means that a
cancellation-sensitive YM norm may be used only after the exact
packet response and its Fourier symmetry have been established.

Thus the audit changes no constants.  It fixes the legal order:

\[
 \boxed{
 \text{complete shift-coset trace}
 \ \longrightarrow\
 \text{Fourier coefficients}
 \ \longrightarrow\
 \text{signed congruence core}
 \ \longrightarrow\
 \text{absolute tail only}.}                              \tag{55.2}
\]

\subsection{Packet-first response}

Let \(F\colon{\mathbb T}^4\to{\mathbb C}\) denote the scalar
response after all internal fibres and every card in the legal
shift packet have been traced.  For a single axial background at
finite \(M\), this means the projectors \(Q_a\) of (50.6), followed
by the sum over the \(M^3\) packets.  For a response involving all
four unit shifts, Corollary~\ref{cor:n1v50-all-directions} requires
the full \(M^4\)-card trace.

Assume for the moment that \(F\) has an absolutely convergent Fourier
series

\[
 F(k)=\sum_{n\in{\mathbb Z}^4}\widehat F(n)e^{in\cdot k},
 \qquad
 \sum_n|\widehat F(n)|<\infty.                            \tag{55.3}
\]

Define its normalized continuum and \(M\)-grid means by

\[
 I(F):=\frac1{(2\pi)^4}\int_{[0,2\pi]^4}F(k)\,dk,
 \qquad
 G_M(F):=\frac1{M^4}
 \sum_{j\in\{0,\ldots,M-1\}^4}F(2\pi j/M).               \tag{55.4}
\]

\begin{theorem}[Signed Poisson congruence identity]
\label{thm:n1v55-Poisson}
Under (55.3),

\[
 \boxed{
 I(F)=\widehat F(0),\qquad
 G_M(F)=\sum_{\ell\in{\mathbb Z}^4}\widehat F(M\ell),}
                                                               \tag{55.5}
\]

and hence

\[
 \boxed{
 G_M(F)-I(F)=
 \sum_{0\ne\ell\in{\mathbb Z}^4}\widehat F(M\ell).}       \tag{55.6}
\]

The series in (55.5)--(55.6) are absolutely convergent, but no
absolute values have been inserted between their summands.
\end{theorem}

\begin{proof}
Absolute convergence permits interchange of the Fourier series with
the integral and the finite grid sum.  The integral of
\(e^{in\cdot k}\) is zero unless \(n=0\), proving the first identity.
For the grid,

\[
 \frac1{M^4}\sum_{j\in\{0,\ldots,M-1\}^4}
 e^{2\pi i n\cdot j/M}
 =\prod_{r=1}^4
 \left(\frac1M\sum_{j_r=0}^{M-1}e^{2\pi i n_rj_r/M}\right).
                                                               \tag{55.7}
\]

Each factor is one if \(M\mid n_r\) and zero otherwise.  Thus only
\(n=M\ell\) survive.  Subtracting the zero mode proves (55.6).
\end{proof}

At \(M=2\), equation (55.6) is the exact even-frequency sum

\[
 G_2(F)-I(F)=\sum_{0\ne\ell}\widehat F(2\ell).             \tag{55.8}
\]

This is the quantity that must be estimated; the universal
multiplier (54.2) bounds it only after replacing every term by its
absolute value and also summing Fourier residue classes that the
grid never sees.

\subsection{A finite signed core with a rigorous tail}

For \(\rho>0\), write

\[
 \|F\|_\rho:=
 \sum_{n\in{\mathbb Z}^4}e^{\rho|n|_1}|\widehat F(n)|.     \tag{55.9}
\]

For \(L\ge0\), define the finite signed alias core

\[
 S_{M,L}(F):=
 \sum_{\substack{0<|\ell|_1\le L}}
 \widehat F(M\ell).                                      \tag{55.10}
\]

\begin{theorem}[Cancellation-preserving core-tail bound]
\label{thm:n1v55-core-tail}
If \(\|F\|_\rho<\infty\), then

\[
 \boxed{
 \left|G_M(F)-I(F)-S_{M,L}(F)\right|
 \le e^{-\rho M(L+1)}\|F\|_\rho.}                         \tag{55.11}
\]

If \(F\) is real-valued, the finite core is real and may be written

\[
 S_{M,L}(F)=
 2\operatorname{Re}
 \sum_{\ell\in\Lambda_L}\widehat F(M\ell),                \tag{55.12}
\]

where \(\Lambda_L\) contains one representative from each pair
\(\{\ell,-\ell\}\) with \(0<|\ell|_1\le L\).
\end{theorem}

\begin{proof}
By Theorem~\ref{thm:n1v55-Poisson}, the left side of (55.11) is the
tail over \(|\ell|_1\ge L+1\).  For every such term,

\[
 e^{\rho|M\ell|_1}
 \ge e^{\rho M(L+1)}.                                    \tag{55.13}
\]

Therefore

\[
 \begin{split}
 \left|\sum_{|\ell|_1\ge L+1}\widehat F(M\ell)\right|
 &\le e^{-\rho M(L+1)}
 \sum_{|\ell|_1\ge L+1}
 e^{\rho|M\ell|_1}|\widehat F(M\ell)|\\
 &\le e^{-\rho M(L+1)}\|F\|_\rho,
 \end{split}                                              \tag{55.14}
\]

which proves (55.11).  If \(F\) is real,
\(\widehat F(-n)=\overline{\widehat F(n)}\); pairing opposite
frequencies proves (55.12).
\end{proof}

\begin{corollary}[A response-specific sign certificate]
\label{cor:n1v55-sign-certificate}
Suppose a target sign is positive and exact or interval calculations
give

\[
 G_M(F)-S_{M,L}(F)
 >e^{-\rho M(L+1)}\|F\|_\rho.                             \tag{55.15}
\]

Then \(I(F)>0\).  The reversed strict inequality certifies
\(I(F)<0\).
\end{corollary}

\begin{proof}
Rearrange (55.11) and use the strict margin in (55.15).  The negative
case is identical after replacing \(F\) by \(-F\).
\end{proof}

This criterion uses absolute values only beyond the computed core.
It can succeed at \(M=2\) even though
Theorem~\ref{thm:n1v54-incompatibility} forbids success of the
all-frequency absolute multiplier.  Success is not automatic: the
finite coefficients in (55.10), the norm in (55.9), and a strict
margin in (55.15) still have to be certified for the actual
packet-first response.

\begin{warning}[No cancellation before the legal trace]
\label{warn:n1v55-trace-first}
Equations (55.3)--(55.15) must be applied to the complete scalar
packet trace \(F\).  They do not license pairing Fourier
coefficients of individual momentum cards, because V47 and V50 show
that the background transport moves those cards.  Likewise, a
hyperoctahedral orbit reduction is legal only after the corresponding
symmetry of \(F\) has been proved.
\end{warning}

\subsection{Next acquisition}

\begin{decision}[V56 construct the packet-first Fourier oracle]
\label{dec:n1v55-next}
Extend the exact uncompressed V48 response evaluator from the sixteen
\(M=2\) corners to arbitrary real momentum.  Preserve the complete
legal trace before scalar extraction.  Then compute interval
enclosures of the first even Fourier coefficients
\(\widehat F(2\ell)\), beginning with \(|\ell|_1=1\), and retain their
signed sum as in (55.10).  In parallel, populate a residual Wiener
stock only for the tail outside that core.  Do not use a structured
or symmetry-reduced norm until the packet response has been proved
to possess that structure.
\end{decision}

\begin{assessment}[Progress after compact V55]
\label{ass:n1v55-status}
The compulsory audit is complete.  It confirms the packet-first
order and supplies no foreign numerical shortcut.  The generic
absolute alias obstruction of V54 is now replaced by an exact signed
Poisson identity and a finite-core/tail theorem that isolates the
specific data needed for an \(M=2\) continuum-sign certificate.
\end{assessment}

\begin{firewall}[Claim boundary after compact V55]
\label{fire:n1v55-claim}
V55 proves the signed Poisson and core-tail identities but has not
yet computed the Fourier core or its residual Wiener stock for the
actual Yang--Mills response.  It therefore does not determine the
continuum response, close SCIC or ADMC, construct the continuum
Yang--Mills measure or exponential clustering, or prove the
Yang--Mills mass gap.
\end{firewall}

\subsection{Parent record}

This V55 note was formed by appending the preceding material to the
corrected compact V54.  Before the append, a mechanical transport
repair restored stripped inline-math delimiters and five split
\(\rho\) tokens in the V53--V54 text; no theorem, formula, or
numerical value was changed.  The corrected parent record is

\[
\begin{array}{c|c}
\text{lines}&20666\\
\text{bytes}&704957\\
\text{SHA--256}&
\texttt{60D55EA6558743E2513AC7287A4A27394458DB94694A553FBD17CA2149A1D45A}.
\end{array}
\]

% =============================================================================
% END APPENDED COMPACT-SERIES V55 MATERIAL
% =============================================================================

% =============================================================================
% BEGIN APPENDED COMPACT-SERIES V56 MATERIAL
% =============================================================================

\section{Compact-series V56: exact Laurent lift of the Wilson
quartic response matrix}
\label{sec:n1v56}

\subsection{The information discarded by a parity mask}

The V48 quartic compiler represents a momentum monomial \(z^\alpha\)
at the \(M=2\) corners by the parity mask

\[
 \pi_2(\alpha):=
 (\alpha_1,\ldots,\alpha_4)\pmod2
 \in(\mathbb Z/2\mathbb Z)^4.                             \tag{56.1}
\]

At a corner \(z_r\in\{1,-1\}\), this is exact because
\(z^\alpha\) depends only on (56.1).  It is not, by itself, an
arbitrary-momentum representation.

\begin{lemma}[Parity reduction is a star-algebra quotient]
\label{lem:n1v56-quotient}
Let \(\mathbb Q[\mathbb Z^4]\) be the Laurent group algebra with
involution

\[
 (c z^\alpha)^*=\overline c\,z^{-\alpha}.                 \tag{56.2}
\]

Mapping \(z^\alpha\) to the character indexed by
\(\pi_2(\alpha)\) defines a surjective star-algebra homomorphism

\[
 \Pi_2:\mathbb Q[\mathbb Z^4]
 \longrightarrow
 \mathbb Q[(\mathbb Z/2\mathbb Z)^4].                     \tag{56.3}
\]

Its kernel is nonzero; for example \(z_1^2-1\in\ker\Pi_2\).
\end{lemma}

\begin{proof}
Reduction modulo two respects exponent addition, so
\(\Pi_2(z^{\alpha+\beta})=\Pi_2(z^\alpha)\Pi_2(z^\beta)\).
It also respects the involution because
\(-\alpha\equiv\alpha\pmod2\).  Every parity class has an integer
representative, proving surjectivity.  Finally, \(z_1^2\) and \(1\)
have the same parity character but are distinct Laurent monomials.
\end{proof}

Thus two Laurent functions can agree on all sixteen V48 corners and
have different Fourier coefficients and continuum integrals.  An
arbitrary-momentum oracle must lift the exponent before evaluating;
interpolating only the corner table is underdetermined.

\subsection{Lifted compiler}

Retain the V48 noncommutative word, colour, background-router, and
external-jet rules.  Change only the momentum bookkeeping:

\[
\begin{array}{c|c|c}
&\text{V48 mask algebra}&\text{V56 Laurent algebra}\\ \hline
\text{factor key}&\pi_2(\eta o)&\eta o\in\mathbb Z^4\\
\text{product}&m_1\mathbin{\mathrm{xor}}m_2&\alpha_1+\alpha_2\\
\text{adjoint}&m\mapsto m&\alpha\mapsto-\alpha.
\end{array}                                                \tag{56.4}
\]

Here \(o\) is the parity-cell offset and \(\eta=\pm1\) is the
orientation carried by the field factor.  Coefficients remain exact
rational three-jets in the external axial momentum.

\begin{theorem}[Exact Wilson exponent-lift certificate]
\label{thm:n1v56-lift}
Let \(Q_{\rm W}^{\rm lift}(z)\) be the Wilson quartic response
three-jet matrix produced by (56.4), and let
\(Q_{\rm W}^{\rm mask}\) be the original V48 mask matrix.  Then

\[
 \boxed{\Pi_2(Q_{\rm W}^{\rm lift})
 =Q_{\rm W}^{\rm mask}}                                   \tag{56.5}
\]

entry by entry and in all three external-jet orders.  Consequently
the lifted matrix reproduces every one of the sixteen V48 corner
matrices.

The lifted matrix is Laurent-Hermitian and has the exact inventory

\[
\begin{array}{c|r}
\text{nonzero matrix entries}&232\\
\text{Laurent exponent terms}&264\\
\text{nonzero rational jet components}&356\\
\max|\alpha|_1&2\\
\min\alpha_r,\ \max\alpha_r&-1,\ 1\quad(r=1,\ldots,4).
\end{array}                                                \tag{56.6}
\]
\end{theorem}

\begin{proof}
The word compiler starts from scalar identities and oriented
monomials and then uses only addition, multiplication, rational
scaling, and adjoint.  Lemma~\ref{lem:n1v56-quotient} shows that
\(\Pi_2\) commutes with every one of these operations.  Structural
induction over each generated word therefore gives (56.5).

The executable verifier performs both constructions independently:
the original V48 mask build, and the exponent build (56.4).  It then
reduces every lifted entry modulo two and compares the complete
arrays, not merely their statistics.  It also evaluates both arrays
at all sixteen characters and obtains equality in every case.

For Laurent-Hermiticity it checks

\[
 [z^\alpha](Q_{\rm W}^{\rm lift})_{ij}^{[r]}
 =
 [z^\alpha]\left(
 (Q_{\rm W}^{\rm lift})_{ji}^{[r]}
 \right)^*                                                \tag{56.7}
\]

with the jet involution that reverses the odd external derivative.
Finally, direct enumeration gives (56.6).
\end{proof}

\begin{remark}[Equal term counts do not identify the variables]
\label{rem:n1v56-equal-count}
Both the lifted array and its mask reduction contain \(264\) terms.
For this Wilson matrix there is no within-entry collision during
modulo-two reduction.  Nevertheless a mask key is not a Laurent
exponent; equation (56.4) is still needed to know whether a character
at general \(z\) is \(z^\alpha\).  The equality of counts is recorded
rather than misreported as a strict compression.
\end{remark}

\begin{corollary}[Finite arbitrary-momentum Wilson quartic oracle]
\label{cor:n1v56-oracle}
For every \(k\in\mathbb R^4\), the complete Wilson quartic
three-jet matrix is given by the finite exact-coefficient evaluation

\[
 Q_{\rm W}^{[r]}(k)
 =
 \sum_{|\alpha|_1\le2}
 [z^\alpha]Q_{\rm W}^{[r]}\,e^{i\alpha\cdot k},
 \qquad r=0,1,2.                                         \tag{56.8}
\]

It can therefore be evaluated at arbitrary real or complex momentum
without rebuilding any noncommutative words.
\end{corollary}

\begin{proof}
This is the Laurent evaluation of the finite array certified in
Theorem~\ref{thm:n1v56-lift}.  The range statement follows from
(56.6).
\end{proof}

\subsection{Executable record and scope}

\begin{record}[V56 Wilson exponent lift]
\label{rec:n1v56-executable}
The verifier is

\[
 \texttt{scripts/verify\_ym\_v56\_exponent\_lift.py}.      \tag{56.9}
\]

Its source record is

\[
\begin{array}{c|c}
\text{lines}&252\\
\text{bytes}&8384\\
\text{SHA--256}&
\texttt{B06F9DDFB05CFF98D1C4AA2CACA872106A1A59A8F442DED18A632FFD430AFC39}.
\end{array}                                                \tag{56.10}
\]

The canonical payload SHA--256 is

\[
 \boxed{
 \texttt{DDD4AAB1F7ADCA8C80C940093D2525A1BBF0AA9ACA970419255DC7C0483370E8}.}
                                                               \tag{56.11}
\]

The exact lifted and reduced matrix hashes are respectively

\[
\begin{split}
 &\texttt{9D0C3D83B16E3E22AEA24E4C1FD6DA7761EB4E254FCA842E25D5812BCD7B8989},\\
 &\texttt{16E4DB5C5C4468FAF8AAA1368D39B4892640D5F4044730EE91E3F89C8D239CA3}.
\end{split}                                                \tag{56.12}
\]

The second hash equals the independently rebuilt original V48
Wilson mask matrix hash.
\end{record}

\begin{warning}[Only the Wilson quartic matrix is lifted]
\label{warn:n1v56-scope}
V56 does not yet lift the endpoint completion quartic matrix or the
cubic bubble symbols.  A monolithic attempt to assemble all three
quartic matrices was stopped after its exact \(384\)-owner completion
phase proved too coarse-grained for interactive certification.  No
partial output from that interrupted run is used.  The scoped Wilson
certificate (56.5) completed independently.
\end{warning}

\begin{decision}[V57 stream and cache the endpoint completion lift]
\label{dec:n1v56-next}
Exploit additivity over the \(384\) owner rows.  Lift one owner block
at a time, reduce it modulo two against the corresponding V48 block,
and add its exact Laurent coefficients to a deterministic cache.
Record a rolling hash and restart point after each owner.  This turns
the expensive endpoint completion into \(384\) bounded certificates
and prevents a lost monolithic run.  Only after the completed cache
passes the toron and all-corner regressions should it enter the
packet-first Fourier oracle.
\end{decision}

\begin{assessment}[Progress after compact V56]
\label{ass:n1v56-status}
The first actual arbitrary-momentum response constituent is now
available: the Wilson quartic three-jet matrix has an exact finite
Laurent representation tied entrywise to V48.  The work also exposes
the endpoint completion as a computational granularity bottleneck,
not a mathematical permission to infer Fourier data from corners.
\end{assessment}

\begin{firewall}[Claim boundary after compact V56]
\label{fire:n1v56-claim}
V56 certifies only the Wilson quartic response matrix.  It does not
yet provide the endpoint quartic or cubic-bubble lifts, the complete
packet-first scalar response, its signed Fourier core or Wiener tail,
the continuum response sign, SCIC or ADMC, a continuum Yang--Mills
measure, exponential clustering, or the Yang--Mills mass gap.
\end{firewall}

\subsection{Parent record}

This V56 note was formed by appending the preceding material to the
validated compact V55, whose record was

\[
\begin{array}{c|c}
\text{lines}&20959\\
\text{bytes}&714665\\
\text{SHA--256}&
\texttt{9D7B7A75BCF6713BCC41A1F224C6EDD96AA6384758C1C8DD3274701A9C1EF207}.
\end{array}
\]

% =============================================================================
% END APPENDED COMPACT-SERIES V56 MATERIAL
% =============================================================================

% =============================================================================
% BEGIN APPENDED COMPACT-SERIES V57 MATERIAL
% =============================================================================

\section{Compact-series V57: streamed exact Laurent lift of the
endpoint completion quartic}
\label{sec:n1v57}

\subsection{Ownerwise decomposition}

The endpoint completion quartic is a sum over the \(384\) covariant
forward-difference owners already fixed in V21:

\[
 Q_{\rm comp}=\sum_{a=0}^{383}Q_{{\rm comp},a}.            \tag{57.1}
\]

The failed monolithic V56 attempt assembled the entire sum before it
could expose a restart point.  Since every operation is additive in
the owner row, (57.1) permits a stronger certificate: lift and verify
each owner separately, and admit it to the accumulated matrix only
after its parity reduction agrees with the corresponding V48 block.

Let \(Q_{{\rm comp},a}^{\rm lift}\) denote the full Laurent-exponent
block and \(Q_{{\rm comp},a}^{\rm mask}\) the original V48 parity-mask
block.  The quotient \(\Pi_2\) is that of
Lemma~\ref{lem:n1v56-quotient}.

\begin{theorem}[All-owner endpoint completion lift]
\label{thm:n1v57-all-owners}
For every \(a=0,\ldots,383\),

\[
 \boxed{
 \Pi_2\!\left(Q_{{\rm comp},a}^{\rm lift}\right)
 =Q_{{\rm comp},a}^{\rm mask}.}                           \tag{57.2}
\]

Consequently the accumulated exact Laurent matrix

\[
 Q_{\rm comp}^{\rm lift}
 :=\sum_{a=0}^{383}Q_{{\rm comp},a}^{\rm lift}             \tag{57.3}
\]

satisfies

\[
 \boxed{\Pi_2(Q_{\rm comp}^{\rm lift})=Q_{\rm comp}^{\rm mask}}.
                                                               \tag{57.4}
\]

The equality holds entrywise and in all three external-jet orders.
In particular, the lift reproduces the endpoint-completion quartic
at every one of the sixteen \(M=2\) corners.
\end{theorem}

\begin{proof}
For a fixed owner, the V57 program invokes the same word list, local
columns, endpoint router fields, colour inverse, background colour,
and normalization twice.  In the first invocation it uses the V48
mask product and involution.  In the second it uses exponent addition
and exponent negation as in (56.4).  It then applies \(\Pi_2\) to the
second result and compares the complete rational jet dictionaries.
The block enters (57.3) only if the comparison succeeds.

The completed checkpoint contains exactly \(384\) owner records in
the order \(0,\ldots,383\), and the run reports no failed comparison.
This proves (57.2).  Linearity of the star-algebra homomorphism gives

\[
 \Pi_2\!\left(\sum_aQ_{{\rm comp},a}^{\rm lift}\right)
 =\sum_a\Pi_2(Q_{{\rm comp},a}^{\rm lift})
 =\sum_aQ_{{\rm comp},a}^{\rm mask},
                                                               \tag{57.5}
\]

which is (57.4).  Evaluation of equal mask matrices at any parity
character gives the corner assertion.
\end{proof}

\subsection{Exact inventory}

\begin{proposition}[Endpoint completion Laurent inventory]
\label{prop:n1v57-inventory}
The matrix \(Q_{\rm comp}^{\rm lift}\) is Laurent-Hermitian and has

\[
\begin{array}{c|r}
\text{nonzero matrix entries}&1407\\
\text{Laurent exponent terms}&6775\\
\text{nonzero rational jet components}&8743\\
\max|\alpha|_1&4\\
\min\alpha_r,\ \max\alpha_r&-2,\ 2\quad(r=1,\ldots,4).
\end{array}                                                \tag{57.6}
\]

Its exact coefficient-array SHA--256 is

\[
 \boxed{
 \texttt{6103A568E5CDC49A3D4C37B6845616528869D5985B1E810EA6168547B2912CE3}.}
                                                               \tag{57.7}
\]

The parity-reduced array SHA--256 is

\[
 \texttt{EBE106EF9F16E24CD517F4712C51CEB46130136B3AFD55AEED6F43BB95D589EE}.
                                                               \tag{57.8}
\]
\end{proposition}

\begin{proof}
After the last accepted owner, the program sums identical exponent
keys and deletes exact zero jets.  It enumerates the remaining
matrix entries, exponents, and nonzero jet components, giving
(57.6).  It checks the Laurent involution entrywise:

\[
 (Q_{\rm comp}^{\rm lift})_{ij}
 =
 \left((Q_{\rm comp}^{\rm lift})_{ji}\right)^*.            \tag{57.9}
\]

Canonical lexicographic serialization of the rational coefficient
array and of its parity reduction gives (57.7)--(57.8).
\end{proof}

The completion toron zero-jet, evaluated in the V48 normalization,
has matrix SHA--256

\[
 \texttt{9CD1792A0EE0085FBAA6B5404529007F4BA6718375867CCC776E08C935F380F1}.
                                                               \tag{57.10}
\]

This digest is used only as a stability regression for the V48
ownerwise construction.  The earlier V33 digest belongs to a
different intermediate normalization and is not used as a
cross-normalization equality.

\subsection{Restart integrity}

The checkpoint contains the full accumulated Laurent matrix, the
number of completed owners, and one record per admitted block.  Its
load rule recomputes the accumulated-matrix hash and rejects the
cache unless

\[
 \#\{\text{owner records}\}
 =\text{completed owner count}.                           \tag{57.11}
\]

Writes use a temporary file followed by atomic replacement.  Hence a
process interruption can lose at most the current uncommitted batch;
it cannot silently convert a partial matrix into a completed one.
At completion the program separately checks

\[
 \text{completed count}=384,\qquad
 (a_0,\ldots,a_{383})=(0,\ldots,383).                     \tag{57.12}
\]

\begin{record}[V57 streamed completion certificate]
\label{rec:n1v57-executable}
The verifier is

\[
 \texttt{scripts/verify\_ym\_v57\_stream\_completion\_lift.py}.
                                                               \tag{57.13}
\]

Its source record is

\[
\begin{array}{c|c}
\text{lines}&295\\
\text{bytes}&11247\\
\text{SHA--256}&
\texttt{6A9DBB6C85A9EBD092E0B4DF2D71D6E1D41F2158192A9F21E0DED3CC378CDB96}.
\end{array}                                                \tag{57.14}
\]

The completed cache is

\[
 \texttt{certificates/ym\_v57\_completion\_lift\_checkpoint.json},
                                                               \tag{57.15}
\]

with

\[
\begin{array}{c|c}
\text{bytes}&346281\\
\text{SHA--256}&
\texttt{DAEF29051D002C754593CF3EE08412DFC3CDF9C3C5EC14E6519230528AE165AF}.
\end{array}                                                \tag{57.16}
\]

The canonical final payload SHA--256 is

\[
 \boxed{
 \texttt{3E02D3FEBE7AC75CF38EC2122E28B13989B6E891BD5F0429B4A7E65E52A8B125}.}
                                                               \tag{57.17}
\]
\end{record}

\begin{corollary}[Finite arbitrary-momentum completion oracle]
\label{cor:n1v57-oracle}
For every real or complex \(k\in\mathbb C^4\), the endpoint
completion quartic three-jet is the finite sum

\[
 Q_{\rm comp}^{[r]}(k)
 =
 \sum_{|\alpha|_1\le4}
 [z^\alpha]Q_{\rm comp}^{[r]}\,e^{i\alpha\cdot k},
 \qquad r=0,1,2.                                         \tag{57.18}
\]

The coefficients are read from the completed cache after its hash
and completion count pass.
\end{corollary}

\begin{proof}
Use the finite Laurent array in (57.3) and the support bound in
(57.6).  The cache checks prevent use of a partial sum.
\end{proof}

\begin{decision}[V58 complete the endpoint quartic oracle]
\label{dec:n1v57-next}
Lift the Wilson plaquette quartic evaluated on the endpoint router
fields.  This is a ninety-six-plaquette construction of the same
size as the completed V56 Wilson lift.  Reduce it entrywise against
the V48 endpoint-Wilson mask matrix, then form

\[
 Q_*^{\rm lift}
 =\frac13Q_{{\rm W},*}^{\rm lift}
  +Q_{\rm comp}^{\rm lift}.                               \tag{57.19}
\]

Record the combined support, Laurent-Hermiticity, all-corner
reduction, and exact coefficient hash.  This will finish the
arbitrary-momentum quartic side; the cubic bubbles remain separate.
\end{decision}

\begin{assessment}[Progress after compact V57]
\label{ass:n1v57-status}
The endpoint completion quartic is no longer a monolithic
computational bottleneck.  All \(384\) owners have exact Laurent
lifts tied blockwise to V48, and the completed directly evaluable
matrix is cached with restart and integrity records.
\end{assessment}

\begin{firewall}[Claim boundary after compact V57]
\label{fire:n1v57-claim}
V57 certifies the endpoint completion quartic but not yet the
endpoint Wilson part, cubic bubbles, complete packet-first scalar
response, signed Fourier core, residual Wiener tail, continuum
response sign, SCIC or ADMC, a continuum Yang--Mills measure,
exponential clustering, or the Yang--Mills mass gap.
\end{firewall}

\subsection{Parent record}

This V57 note was formed by appending the preceding material to the
validated compact V56, whose record was

\[
\begin{array}{c|c}
\text{lines}&21223\\
\text{bytes}&723854\\
\text{SHA--256}&
\texttt{627247BCFC1A511C5158B1D3DED869FCA598501597EF3B8CF174C3A0BFEF54DC}.
\end{array}
\]

% =============================================================================
% END APPENDED COMPACT-SERIES V57 MATERIAL
% =============================================================================

% =============================================================================
% BEGIN APPENDED COMPACT-SERIES V58 MATERIAL
% =============================================================================

\section{Compact-series V58: the complete endpoint quartic Laurent
oracle}
\label{sec:n1v58}

\subsection{The remaining plaquette part}

The endpoint quartic in V48 is

\[
 Q_*=\frac13Q_{{\rm W},*}+Q_{\rm comp},                   \tag{58.1}
\]

where \(Q_{{\rm W},*}\) is the Wilson plaquette quartic evaluated on
the endpoint router fields and \(Q_{\rm comp}\) is the completion
lift certified in V57.  The first summand has the same
ninety-six-plaquette word geometry as the V56 Wilson calculation;
only its exact router fields differ.

\begin{theorem}[Endpoint Wilson-part lift]
\label{thm:n1v58-Wilson-part}
There is an exact Laurent three-jet matrix
\(Q_{{\rm W},*}^{\rm lift}\) such that

\[
 \boxed{
 \Pi_2(Q_{{\rm W},*}^{\rm lift})
 =Q_{{\rm W},*}^{\rm mask}}                               \tag{58.2}
\]

entry by entry in the original V48 normalization.  The lift is
Laurent-Hermitian and has

\[
\begin{array}{c|r}
\text{nonzero matrix entries}&232\\
\text{Laurent exponent terms}&264\\
\text{nonzero rational jet components}&356\\
\max|\alpha|_1&2\\
\min\alpha_r,\ \max\alpha_r&-1,\ 1\quad(r=1,\ldots,4).
\end{array}                                                \tag{58.3}
\]

Its exact coefficient-array SHA--256 is

\[
 \texttt{F5B3C12846AA62E7507513577FC187CD6AD36EC3B850B2ECB27E55A23C33CB69}.
                                                               \tag{58.4}
\]
\end{theorem}

\begin{proof}
Run the V48 plaquette word compiler twice with the same endpoint
router fields, colour inverse, background colour, and normalization.
The mask run uses xor and the lifted run uses exponent addition and
negation.  Reducing the latter by \(\Pi_2\) gives exact equality of
all rational jet dictionaries.  Their common reduced array has
SHA--256

\[
 \texttt{F2A70D1F05F6EE55702D7E5FCCC40DE30B8B52B5B8C60E1697C5A8859A819D11}.
                                                               \tag{58.5}
\]

The Laurent involution and the inventory (58.3) are then checked
directly.
\end{proof}

\subsection{Combination with the completion}

Define the full endpoint lift by the exact matrix identity

\[
 Q_*^{\rm lift}
 :=\frac13Q_{{\rm W},*}^{\rm lift}
   +Q_{\rm comp}^{\rm lift}.                              \tag{58.6}
\]

\begin{theorem}[Complete endpoint quartic oracle]
\label{thm:n1v58-endpoint}
The matrix \(Q_*^{\rm lift}\) is Laurent-Hermitian and satisfies

\[
 \boxed{\Pi_2(Q_*^{\rm lift})=Q_*^{\rm mask}}.             \tag{58.7}
\]

It therefore reproduces the complete endpoint quartic matrix at
every \(M=2\) corner.  Its exact inventory is

\[
\begin{array}{c|r}
\text{nonzero matrix entries}&1407\\
\text{Laurent exponent terms}&6781\\
\text{nonzero rational jet components}&8749\\
\max|\alpha|_1&4\\
\min\alpha_r,\ \max\alpha_r&-2,\ 2\quad(r=1,\ldots,4).
\end{array}                                                \tag{58.8}
\]

The exact coefficient-array SHA--256 is

\[
 \boxed{
 \texttt{C294D154C68F7744ED34952FAA6D3EEC516772FB1235E253E54D2B445A02405A}.}
                                                               \tag{58.9}
\]
\end{theorem}

\begin{proof}
Equation (58.6) uses exact rational scaling and addition.  By
Theorem~\ref{thm:n1v58-Wilson-part},
Theorem~\ref{thm:n1v57-all-owners}, and linearity of \(\Pi_2\),

\[
\begin{split}
 \Pi_2(Q_*^{\rm lift})
 &=
 \frac13\Pi_2(Q_{{\rm W},*}^{\rm lift})
 +\Pi_2(Q_{\rm comp}^{\rm lift})\\
 &=\frac13Q_{{\rm W},*}^{\rm mask}+Q_{\rm comp}^{\rm mask}
 =Q_*^{\rm mask}.
\end{split}                                               \tag{58.10}
\]

The executable independently forms the right side, compares it with
the reduction of (58.6), and evaluates both at all sixteen parity
characters.  Every comparison is exact.  Laurent-Hermiticity is
preserved by the real scaling and sum and is also checked entrywise.
Exact enumeration gives (58.8), and canonical serialization gives
(58.9).
\end{proof}

The reduced full endpoint mask array has SHA--256

\[
 \texttt{ADA5F602619B52D21015942BCC5AF4874EC9677EB77FCB52DE3E1C97E6856D9D}.
                                                               \tag{58.11}
\]

Its toron zero-jet matrix in the V48 normalization has SHA--256

\[
 \texttt{D29EE00B1407FFD84A8FED0F878B9CACCC0B3CA76BF3AA1E36263B6C744F8069}.
                                                               \tag{58.12}
\]

\begin{corollary}[Both quartic response matrices are now evaluable]
\label{cor:n1v58-both-oracles}
For \(e\in\{{\rm W},*\}\), every real or complex
\(k\in\mathbb C^4\), and external order \(r=0,1,2\),

\[
 Q_e^{[r]}(k)=
 \sum_{|\alpha|_1\le R_e}
 [z^\alpha]Q_e^{[r]}\,e^{i\alpha\cdot k},\qquad
 R_{\rm W}=2,\quad R_*=4.                                \tag{58.13}
\]

The Wilson coefficients are certified in V56 and the endpoint
coefficients in Theorem~\ref{thm:n1v58-endpoint}.
\end{corollary}

\begin{proof}
Apply Corollary~\ref{cor:n1v56-oracle} and the finite Laurent
inventory (58.8).
\end{proof}

\subsection{Executable record}

\begin{record}[V58 complete endpoint quartic certificate]
\label{rec:n1v58-executable}
The verifier is

\[
 \texttt{scripts/verify\_ym\_v58\_endpoint\_quartic\_oracle.py}.
                                                               \tag{58.14}
\]

Its source record is

\[
\begin{array}{c|c}
\text{lines}&144\\
\text{bytes}&5329\\
\text{SHA--256}&
\texttt{A4BA6C27CF1459ED58B58AA27BD02863AEE63136FEC015ED2892E40498C166DD}.
\end{array}                                                \tag{58.15}
\]

The canonical payload SHA--256 is

\[
 \boxed{
 \texttt{2F0F9BA05293D6A426F4B18A90E51D32C1FBBFFB03C4D20865E61C5B5A08543C}.}
                                                               \tag{58.16}
\]

The program refuses a partial completion cache and verifies the
V57 cache digest

\[
 \texttt{DAEF29051D002C754593CF3EE08412DFC3CDF9C3C5EC14E6519230528AE165AF}
                                                               \tag{58.17}
\]

before using its \(384\)-owner matrix.
\end{record}

\begin{decision}[V59 lift the cubic bubble symbols]
\label{dec:n1v58-next}
The quartic tadpole numerators for both actions are now finite
Laurent arrays.  Apply the same exponent-preserving rule to the
Wilson and completion cubic symbols used by V39.  Their stored phase
tuples already retain more routing information than the V48 masks;
make the contraction map to a single loop-momentum Laurent exponent
explicit, and compare its parity reduction at all sixteen corners.
Then combine each cubic pair with the exact Laurent Hessian inverse
to obtain the arbitrary-momentum bubble response.
\end{decision}

\begin{assessment}[Progress after compact V58]
\label{ass:n1v58-status}
The arbitrary-momentum quartic side of the one-loop response is
complete for both the Wilson and endpoint actions.  Every coefficient
is exact, both matrices are Laurent-Hermitian, and their reductions
replay the V48 corner objects.  The remaining response constituent is
the cubic bubble.
\end{assessment}

\begin{firewall}[Claim boundary after compact V58]
\label{fire:n1v58-claim}
V58 completes the quartic Laurent oracle but does not yet construct
the arbitrary-momentum cubic bubbles, complete packet-first scalar
response, signed Fourier core, residual Wiener tail, continuum
response sign, SCIC or ADMC, a continuum Yang--Mills measure,
exponential clustering, or the Yang--Mills mass gap.
\end{firewall}

\subsection{Parent record}

This V58 note was formed by appending the preceding material to the
validated compact V57, whose record was

\[
\begin{array}{c|c}
\text{lines}&21508\\
\text{bytes}&732952\\
\text{SHA--256}&
\texttt{BE4F3D8B825247FCD4CC3D850866F97004631C7433A4FB837641ED6C404E2C66}.
\end{array}
\]

% =============================================================================
% END APPENDED COMPACT-SERIES V58 MATERIAL
% =============================================================================

% =============================================================================
% BEGIN APPENDED COMPACT-SERIES V59 MATERIAL
% =============================================================================

\section{Compact-series V59: exact loop-momentum Laurent oracle for
the cubic vertices}
\label{sec:n1v59}

\subsection{Contracting the retained three-leg phases}

Unlike the V48 quartic mask, the V39 cubic symbol stores a phase

\[
 \beta=(\beta_0,\beta_1,\beta_2),\qquad
 \beta_j\in\mathbb Z^4,                                  \tag{59.1}
\]

for the two fluctuation legs and the background leg.  For the routed
bubble, the loop-momentum and external-momentum incidence vectors
are

\[
 \eta=(1,-1,0),\qquad
 \sigma_+=(0,-1,1),\qquad
 \sigma_-=(0,1,-1).                                      \tag{59.2}
\]

Thus the loop Laurent exponent and the axial external rate are

\[
 \alpha(\beta):=\sum_{j=0}^2\eta_j\beta_j,\qquad
 d_\pm(\beta):=\sum_{j=0}^2
 (\sigma_\pm)_j(\beta_j)_1.                              \tag{59.3}
\]

The last coordinate in (59.3) is the first physical axis, written
with one-based mathematical indexing.

For a phase-polynomial matrix \(S\), define the exact Laurent moment

\[
 {\cal L}_{\pm,r}(S)(z)
 :=
 \sum_\beta [\beta]S\,
 d_\pm(\beta)^r z^{\alpha(\beta)}.                        \tag{59.4}
\]

Coincident exponents are combined before zero coefficients are
deleted.

\begin{lemma}[Corner evaluation commutes with phase contraction]
\label{lem:n1v59-contraction}
For every parity corner \(s\in\{1,-1\}^4\),

\[
 {\cal L}_{\pm,r}(S)(s)
 =
 \sum_\beta[\beta]S\,d_\pm(\beta)^r
 \prod_{a=1}^4s_a^{\alpha(\beta)_a}.                      \tag{59.5}
\]

The right side is exactly the V39
\(\texttt{moment\_matrix\_corner}\) rule.
\end{lemma}

\begin{proof}
Substitute \(z=s\) into the finite sum (59.4).  By (59.3),

\[
 s^{\alpha(\beta)}
 =
 \prod_{a=1}^4
 s_a^{\sum_j\eta_j(\beta_j)_a},                           \tag{59.6}
\]

which is the V39 phase-corner weight.  The rate in (59.4) is its
external moment.  Finite summation makes regrouping by equal
\(\alpha\) exact.
\end{proof}

\subsection{The three external jets}

Let \(S_0,S_1,S_2\) be the three routed cubic symbols.  Define

\[
\begin{aligned}
 V_{\pm}^{[0]}&={\cal L}_{\pm,0}(S_0),\\
 V_+^{[1]}&={\cal L}_{+,1}(S_0)+{\cal L}_{+,0}(S_1),\\
 V_-^{[1]}&={\cal L}_{-,1}(S_0)-{\cal L}_{-,0}(S_1),\\
 V_+^{[2]}&=-{\cal L}_{+,2}(S_0)-2{\cal L}_{+,1}(S_1)
             +{\cal L}_{+,0}(S_2),\\
 V_-^{[2]}&=-{\cal L}_{-,2}(S_0)+2{\cal L}_{-,1}(S_1)
             +{\cal L}_{-,0}(S_2).
\end{aligned}                                             \tag{59.7}
\]

These are Laurent-matrix identities, not corner interpolants.

\begin{theorem}[Exact cubic Laurent vertices]
\label{thm:n1v59-vertices}
For each action \(e\in\{{\rm W},*\}\), the six matrices
\(V_{e,\pm}^{[r]}\), \(r=0,1,2\), constructed by (59.4)--(59.7)
reproduce the corresponding V39 cubic vertex jets at all sixteen
parity corners.

Moreover, the following identities hold as exact Laurent
polynomials:

\[
 \boxed{
 V_{e,-}^{[0]}=V_{e,+}^{[0]},\qquad
 V_{e,-}^{[1]}=-V_{e,+}^{[1]},\qquad
 V_{e,-}^{[2]}=V_{e,+}^{[2]}.}                            \tag{59.8}
\]
\end{theorem}

\begin{proof}
Apply Lemma~\ref{lem:n1v59-contraction} to every term in (59.7).
The resulting six corner formulas are precisely those in the V39
\(\texttt{cubic\_corner\_jets}\) function.  The verifier compares
the complete rational matrices for both actions, both orientations,
all three orders, and all sixteen corners.

For (59.8), compare the exact coefficient dictionaries after
combining equal loop exponents.  Orders zero and two agree
entrywise; every order-one coefficient is the negative of its
opposite-orientation partner.  This coefficientwise test is stronger
than equality only at the corners.
\end{proof}

\subsection{Exact inventories}

Because of (59.8), plus and minus have the same inventories.  The
Wilson plus matrices have

\[
\begin{array}{c|rrrr}
r&\text{nonzero entries}&\text{Laurent terms}&
 \sum|c_\alpha|&\max|\alpha|_1\\ \hline
0&92&92&92&1\\
1&188&220&293/4&2\\
2&374&444&289/2&2.
\end{array}                                                \tag{59.9}
\]

Their coefficient-array SHA--256 fingerprints are

\[
\begin{array}{c|c}
0&\texttt{FE6780C2043FB0817069ADB0EC4FAAEF4FA49702EFEB006F22C1E41D16BD6B4D}\\
1&\texttt{491369B6AE18FE3097A8B47D783E93B1F1E9AEF7E3D67DF6C47D8581E3A70ABA}\\
2&\texttt{715F0C724A05634EC2E386C229E4F7D232031AED28014156ADB46B2357878E42}.
\end{array}                                                \tag{59.10}
\]

The endpoint plus matrices have

\[
\begin{array}{c|rrrr}
r&\text{nonzero entries}&\text{Laurent terms}&
 \sum|c_\alpha|&\max|\alpha|_1\\ \hline
0&664&1912&6278/3&3\\
1&1037&4132&6203/3&4\\
2&1374&6425&769943029/346320&4.
\end{array}                                                \tag{59.11}
\]

Their fingerprints are

\[
\begin{array}{c|c}
0&\texttt{0E8DD205747DA1ABF574760A4B79CAA357866FA3E96CB85054EEC3E3E5A00B25}\\
1&\texttt{17678EA3A16221337DC96F4296911C0612484A120D0659DFC9BF01CCE9735ACA}\\
2&\texttt{5E27336D2AB315693F72749489AE6FF6AB75B2CB0BC653B696CA04879EFDD0B0}.
\end{array}                                                \tag{59.12}
\]

For the minus orientation, the order-zero and order-two hashes equal
those in (59.10) and (59.12).  The order-one hashes, reflecting the
sign reversal, are

\[
\begin{split}
 &\texttt{1384CAE04AAE1CCE041711EBF90EF389DD9389CB712FC2E50F43D1103B3F7869}
 \quad({\rm W}),\\
 &\texttt{750C44136557F0B0C34D35D769DA37E46923547D673DE3E1C99913E98B88105C}
 \quad(*).
\end{split}                                                \tag{59.13}
\]

\subsection{Serialized oracle}

\begin{record}[V59 cubic vertex certificate]
\label{rec:n1v59-executable}
The verifier is

\[
 \texttt{scripts/verify\_ym\_v59\_cubic\_vertex\_oracle.py}.
                                                               \tag{59.14}
\]

Its source record is

\[
\begin{array}{c|c}
\text{lines}&346\\
\text{bytes}&11658\\
\text{SHA--256}&
\texttt{D1C24653DAD97ED84045AFC6225681827BC245C675E7082C07E965EBF676DC0A}.
\end{array}                                                \tag{59.15}
\]

The twelve matrices are serialized in

\[
 \texttt{certificates/ym\_v59\_cubic\_vertex\_oracle.json},
                                                               \tag{59.16}
\]

whose record is

\[
\begin{array}{c|c}
\text{bytes}&740309\\
\text{SHA--256}&
\texttt{44DD8EA5ED92D3ACA87034A29AD64D40FC6570055C98882ED463E8D7B7D7E280}.
\end{array}                                                \tag{59.17}
\]

The canonical summary-payload SHA--256 is

\[
 \boxed{
 \texttt{54645E3924098ECA6B2296311DA2A82455DC20D60AC5D69E506EB6ABEB6E1EAF}.}
                                                               \tag{59.18}
\]
\end{record}

\begin{corollary}[The one-loop numerator oracles are complete]
\label{cor:n1v59-numerators}
For both actions, all quartic matrices entering the tadpole and all
cubic matrices entering the bubble are now finite Laurent arrays at
arbitrary loop momentum, through external order two.
\end{corollary}

\begin{proof}
Use Corollary~\ref{cor:n1v58-both-oracles} for the quartic terms and
Theorem~\ref{thm:n1v59-vertices} for the cubic terms.
\end{proof}

The remaining step in forming the scalar response is not a missing
action numerator.  It is the repeated contraction with the
momentum-dependent inverse Hessian and its two external derivatives.

\begin{decision}[V60 mandatory audit and inverse-jet contraction]
\label{dec:n1v59-next}
Perform the compulsory SM/NS audit.  Then define the arbitrary-real-
momentum Hessian jets directly from the exact incidence Laurent
arrays and prove the inverse identities

\[
 G^{[1]}=-G^{[0]}A^{[1]}G^{[0]},\qquad
 G^{[2]}=
 2G^{[0]}A^{[1]}G^{[0]}A^{[1]}G^{[0]}
 -G^{[0]}A^{[2]}G^{[0]}.                                 \tag{59.19}
\]

Combine them with the V58 quartic and V59 cubic oracles to evaluate
the complete packet-first scalar response at arbitrary real
momentum.  Its sixteen corners must reproduce V48 before any Fourier
quadrature is accepted.
\end{decision}

\begin{assessment}[Progress after compact V59]
\label{ass:n1v59-status}
The cubic bubble numerators are now exact arbitrary-momentum Laurent
oracles for both actions.  The phase-to-loop contraction is proved
and checked at every V39 corner, and the global plus/minus parity
identity (59.8) removes an unnecessary orientation duplication.
\end{assessment}

\begin{firewall}[Claim boundary after compact V59]
\label{fire:n1v59-claim}
V59 completes the cubic vertex numerators but not yet the
momentum-dependent inverse-jet contraction, complete packet-first
scalar response, signed Fourier core, residual Wiener tail,
continuum response sign, SCIC or ADMC, a continuum Yang--Mills
measure, exponential clustering, or the Yang--Mills mass gap.
\end{firewall}

\subsection{Parent record}

This V59 note was formed by appending the preceding material to the
validated compact V58, whose record was

\[
\begin{array}{c|c}
\text{lines}&21767\\
\text{bytes}&741009\\
\text{SHA--256}&
\texttt{3A2128330F9394983CBFE199B830AC069184E204C4C53B4130A8135A0F9A67C2}.
\end{array}
\]

% =============================================================================
% END APPENDED COMPACT-SERIES V59 MATERIAL
% =============================================================================


% =============================================================================
% BEGIN APPENDED COMPACT-SERIES V60 MATERIAL
% =============================================================================

\section{Compact-series V60: companion audit and a coefficient-exact
arbitrary-momentum response oracle}
\label{sec:n1v60}

\subsection{Purpose and outcome}

V59 completed the cubic numerator bank but left the inverse covariance
and diagram contractions tied to the sixteen Walsh corners.  This note
first performs the compulsory five-version companion audit.  It then
closes that finite-algebra interface: both Hessians, both quartic
vertices, both orientations of both cubic vertices, the inverse jets,
and the tadpole--bubble contractions are evaluated from finite Laurent
arrays at arbitrary real loop momentum.

\emph{Coefficient-exact} below means that every Laurent coefficient and
every corner calculation is rational and exact.  At a general real
momentum the exponentials and the \(45\)-dimensional inverse are
evaluated in complex double precision.  No interval enclosure is
claimed for those floating values.

\subsection{The compulsory V60 SM/NS audit}

The two current compact-lineage companion notes in
\texttt{latex\_notes} were inspected read-only before the V60
construction.  Their frozen records are
\[
\begin{array}{c|r|r|l}
\text{file}&\text{lines}&\text{bytes}&\text{SHA--256}\\ \hline
\texttt{Renewal\_SM\_Einstein\_Continuum\_Research\_Notes\_V74.tex}
&21232&711243&
\texttt{8A25C8C76C6AAC807F1F3A7F7AA4EAE2D99FAB2B06834DDBD22DB74ED23700D2}\\
\texttt{Renewal\_NS\_Stability\_Research\_Notes\_V26.tex}
&5319&278283&
\texttt{82C887F8C2C69E4F28FAD7C3DC8DE5B9099CB5A4BF431EBA1FFF3E91935978AD}.
\end{array}                                                   \tag{60.1}
\]

SM V74 constructs a compact Hilbert-cube closure from a countable
family of bounded gauge-invariant observables and proves automatic
subsequential convergence of positive cutoff laws there.  It carefully
distinguishes such an observable continuum from a measure on smooth
orbit space: separation of physical orbits and exclusion of ideal
points remain open.  NS V26 computes exact rank-one norms of two
Oseen-kernel derivative tensors.  Its relevant methodological lesson is
to optimize a contraction over the actual admissible tensor class
before importing an ambient operator norm.

\begin{decision}[V60 companion-audit decision]
\label{dec:n1v60-audit}
No SM or NS estimate is inserted into the one-loop response.  The
SM V74 observable compactification is retained as a legitimate
downstream fallback if a global Yang--Mills orbit space proves
pathological.  The NS V26 rank-one calculation reinforces the existing
packet-first rule: future response norms must respect the actual
colour, polarization, and routing subspace.  Neither observation
changes the immediate acquisition, which is the missing inverse-jet
contraction.
\end{decision}

\begin{proof}
The V60 response is a finite-dimensional momentum-space calculation.
Neither companion note contains an identity for its Hessian, cubic
bank, quartic bank, or normalization.  Importing a numerical constant
would therefore be unjustified.  The two structural lessons concern,
respectively, the later continuum state space and later norm bounds,
so recording them without modifying the present algebra is the unique
scope-compatible decision.
\end{proof}

\subsection{Exact Laurent inputs and conventions}

For \(e\in\{{\rm W},*\}\), let
\[
 A_e(k)=\sum_{\alpha\in{\mathbb Z}^4}
 A_{e,\alpha}e^{i\alpha\cdot k},\qquad
 H_e(k):=\frac13 A_e(k),\qquad G_e(k):=H_e(k)^{-1}.          \tag{60.2}
\]
Here \(A_{\rm W}\) is the V21 plaquette Gram matrix and
\(A_*=A_{\rm W}+3H_{\rm rr}\) is the completed endpoint matrix.
All coefficients are integral and the matrices have size \(45\).
For the external axis \(\widehat e_1\), put
\[
 A_e^{[r]}(k):=
 \left.\frac{d^r}{dq^r}A_e(k+q\widehat e_1)\right|_{q=0}
 =
 \sum_\alpha(i\alpha_1)^rA_{e,\alpha}e^{i\alpha\cdot k}.
                                                               \tag{60.3}
\]

The quartic oracle of V58 and the cubic oracle of V59 use a stored-jet
convention.  If \(X\) denotes any one of their matrices, its three
stored Laurent arrays reconstruct the actual derivative matrices as
\[
 X^{[0]}(k)=\sum_\alpha X^{[0]}_\alpha e^{i\alpha\cdot k},
\quad
 X^{[1]}(k)=i\sum_\alpha \dot X_\alpha e^{i\alpha\cdot k},
\quad
 X^{[2]}(k)=\sum_\alpha X^{[2]}_\alpha e^{i\alpha\cdot k}.     \tag{60.4}
\]
In particular, the second stored jet is already the actual second
derivative; it is not to be multiplied by \(i^2\) again.  This
distinction is forced by the product rule
\[
 (XY)^{[2]}_{\rm stored}
 =
 X^{[2]}Y^{[0]}-2\dot X\dot Y+X^{[0]}Y^{[2]}.                \tag{60.5}
\]

\subsection{The inverse-jet theorem}

\begin{theorem}[Inverse jets at arbitrary real momentum]
\label{thm:n1v60-inverse}
Fix \(e\) and a real \(k\) for which \(H_e(k)\) is invertible.  With
\(H_r:=A_e^{[r]}(k)/3\), the derivatives of the covariance are
\[
\begin{split}
 G_0&=H_0^{-1}=3A_e(k)^{-1},\\
 G_1&=-G_0H_1G_0,\\
 G_2&=
 2G_0H_1G_0H_1G_0-G_0H_2G_0.                              \tag{60.6}
\end{split}
\]
They obey
\[
\begin{split}
 H_0G_0&=I,\\
 H_1G_0+H_0G_1&=0,\\
 H_2G_0+2H_1G_1+H_0G_2&=0.                                \tag{60.7}
\end{split}
\]
\end{theorem}

\begin{proof}
Differentiate \(H(q)G(q)=I\) once.  Left multiplication by
\(G_0=H_0^{-1}\) gives the second line of (60.6).  Differentiate
twice:
\[
 H_2G_0+2H_1G_1+H_0G_2=0.
\]
Substitute \(G_1=-G_0H_1G_0\) and left multiply by \(G_0\), giving the
third line of (60.6).  The displayed residuals are precisely the
coefficients of orders zero, one, and two in the original product.
\end{proof}

\subsection{The complete response contraction}

Let \(Q_e^{[r]}(k)\) be the action-normalized quartic arrays: V56's
Wilson bank is multiplied by \(1/3\), while the V58 endpoint bank
already equals one third of its routed Wilson part plus its completion
part.  Let \(P_{e,+}^{[r]}(k)\) and \(P_{e,-}^{[r]}(k)\) be the V59
cubic arrays reconstructed by (60.4).  Put
\[
 c_{\rm W}=\frac13,\qquad c_*=3.                            \tag{60.8}
\]
The tadpole derivatives are
\[
 {\cal T}_e^{[r]}(k)
 =\frac12\Tr\!\left(G_0Q_e^{[r]}(k)\right).                 \tag{60.9}
\]
The covariance in (60.9) is fixed with respect to the external
momentum; all \(q\)-dependence of that diagram lies in the quartic
vertex.

For the bubble, write \(P_r=P_{e,+}^{[r]}(k)\) and
\(M_r=P_{e,-}^{[r]}(k)\).  Direct differentiation of
\[
 {\cal B}_e(k,q)=c_e\Tr\!\left[
 G_0P_{e,+}(k,q)G_e(k+q\widehat e_1)P_{e,-}(k,q)\right]
                                                               \tag{60.10}
\]
gives
\[
\begin{split}
 {\cal B}_e^{[0]}={}&c_e\Tr(G_0P_0G_0M_0),\\
 {\cal B}_e^{[1]}={}&c_e\Tr\bigl(
 G_0P_1G_0M_0+G_0P_0G_1M_0+G_0P_0G_0M_1\bigr),\\
 {\cal B}_e^{[2]}={}&c_e\Tr\bigl(
 G_0P_2G_0M_0+G_0P_0G_2M_0+G_0P_0G_0M_2\\
 &\hspace{18mm}
 +2G_0P_1G_1M_0+2G_0P_1G_0M_1
 +2G_0P_0G_1M_1\bigr).                                    \tag{60.11}
\end{split}
\]
The plus signs in the last line apply to the actual complex
derivatives.  On the \(0/\pi\) grid, stripping a factor \(i\) from
each first jet turns each product of two first jets into the minus
signs appearing in the exact V36 contraction.

\begin{definition}[V60 response oracle]
\label{def:n1v60-response}
On the real invertibility set of \(H_e\), define
\[
 {\cal R}_e^{[r]}(k):=
 {\cal T}_e^{[r]}(k)+{\cal B}_e^{[r]}(k),\qquad r=0,1,2.     \tag{60.12}
\]
The V60 oracle evaluates (60.2)--(60.12) without regenerating any
vertex word expansion.
\end{definition}

\begin{theorem}[Exact sixteen-corner regression]
\label{thm:n1v60-corners}
For both actions and all sixteen \(k\in\{0,\pi\}^4\), the V60
compiler forms \(A_e^{[0,1,2]}\), \(G_e^{[0,1,2]}\), \(Q_e^{[0,1,2]}\),
and \(P_{e,\pm}^{[0,1,2]}\) with rational arithmetic.  Every residual
in (60.7) is exactly zero.  The arithmetic means of the thirty-two
complete response cards are
\[
 \overline{\cal R}_{{\rm W},2}
 =
 \left(
 \frac{176943}{627200},\,0,\,
 -\frac{2822949104411}{265531392000}
 \right),                                                   \tag{60.13}
\]
\[
 \overline{\cal R}_{*,2}
 =
 \left(
 \frac{157603277390904552027570847451}
 {684609230311820211100422720000},\,0,\,
 -\frac{5823821044825964006189852662538728905719501293128460064877013}
 {36786699279534470965406231814271312468634765095027200000000}
 \right),                                                   \tag{60.14}
\]
and
\[
 \overline{\Delta{\cal R}}_2
 =
 \left(
 -\frac{4441971352736182354322716403}
 {85576153788977526387552840000},\,0,\,
 -\frac{97013033679684607051275019240872676380229076929720667338373}
 {656905344277401267239396996683416294082763662411200000000}
 \right).                                                   \tag{60.15}
\]
Thus the new Laurent interface reproduces V48 exactly, without
importing its pointwise response cards.
\end{theorem}

\begin{proof}
At a corner, \(e^{i\alpha\cdot k}\) is the character
\(\prod_\mu\varepsilon_\mu^{\alpha_\mu}\).  Consequently evaluation
of (60.3)--(60.4) commutes with every finite sum and matrix product.
The exact solver applies (60.6), verifies (60.7), and then contracts
(60.9) and the stored-jet form of (60.11).  Direct summation of all
sixteen independently formed cards gives (60.13)--(60.15), equal as
reduced fractions to (48.5)--(48.7).  The canonical serialization of
the thirty-two total cards has SHA--256
\[
 \texttt{62AB46AC77A37E1BF5C2F80ABE0D98FF2E9841A7D1136EEA040AEDC7669F6092}.
                                                               \tag{60.16}
\]
\end{proof}

\begin{corollary}[The finite algebraic response interface is complete]
\label{cor:n1v60-interface}
At every real momentum where the selected reduced Hessian is
invertible, the one-loop packet-first response through external order
two is determined by finite rational Laurent data plus one
\(45\times45\) matrix inverse.  No missing plaquette word, completion
owner, cubic orientation, quartic vertex, or inverse derivative
remains in this interface.
\end{corollary}

\begin{proof}
V58 supplies both quartic banks, V59 supplies all four cubic
orientation banks, (60.2)--(60.3) supply the Hessian jets, Theorem
\ref{thm:n1v60-inverse} supplies the inverse jets, and
(60.9)--(60.12) perform every remaining contraction.
\end{proof}

\subsection{Non-corner execution and numerical scope}

The floating evaluator was checked against every exact corner card.
Its largest absolute card error was
\[
                    1.193759453424219\times10^{-12},         \tag{60.17}
\]
and its largest corner inverse-residual infinity norm was
\[
                    1.3626096778684982\times10^{-12}.        \tag{60.18}
\]
At the two non-corner momenta
\[
\begin{split}
 k^{(1)}&=(0.37,-0.41,0.73,-1.11),\\
 k^{(2)}&=(1.19,0.23,-0.67,2.03),                           \tag{60.19}
\end{split}
\]
the endpoint-minus-Wilson totals were, respectively,
\[
\begin{split}
 \Delta{\cal R}^{[0,1,2]}(k^{(1)})
 \approx{}&(90.9369381242892,\ 0.314257656177694,\\
            &\hspace{20mm}-236.177870832449-27.3353121004417i),\\
 \Delta{\cal R}^{[0,1,2]}(k^{(2)})
 \approx{}&(93.0115775639290,\ 20.0628647109418,\\
            &\hspace{20mm}-239.972117462798+46.8287946108195i).
                                                               \tag{60.20}
\end{split}
\]
The tiny imaginary parts in orders zero and one were below
\(1.2\times10^{-13}\) at these samples and are displayed as zero in
(60.20).  The largest inverse residual over both samples and both
actions was
\[
                    8.6740448389090202\times10^{-13}.        \tag{60.21}
\]
These residuals test implementation consistency, not a rigorous
roundoff enclosure.  The nonzero imaginary second components also
show that a Brillouin integral must retain conjugate momentum pairs;
pointwise replacement by the real part requires a proved reflection
identity.

\subsection{Executable record}

\begin{record}[V60 response-oracle certificate]
\label{rec:n1v60-executable}
The verifier is
\[
 \texttt{scripts/verify\_ym\_v60\_response\_oracle.py}.      \tag{60.22}
\]
Its source record is
\[
\begin{array}{c|c}
\text{lines}&475\\
\text{bytes}&16594\\
\text{SHA--256}&
\texttt{7D5BFEE7FB834DFB5132FA2E1A9C5BC3CED7D57D366C557C8F370B6AD24D6FC8}.
\end{array}                                                  \tag{60.23}
\]
The locked upstream files are the V57 completion cache with digest
\[
 \texttt{DAEF29051D002C754593CF3EE08412DFC3CDF9C3C5EC14E6519230528AE165AF}
                                                               \tag{60.24}
\]
and the V59 cubic certificate with digest
\[
 \texttt{44DD8EA5ED92D3ACA87034A29AD64D40FC6570055C98882ED463E8D7B7D7E280}.
                                                               \tag{60.25}
\]
The canonical output-payload SHA--256 is
\[
 \boxed{
 \texttt{0A4C22C840B1158CB9C7CDA2740F00E6997A902CC531F59547534801986746D3}.}
                                                               \tag{60.26}
\]
\end{record}

\begin{decision}[V61 acquisition]
\label{dec:n1v60-next}
Prove and test the conjugate-reflection law for the complete response,
including its order parity under \((k,q)\mapsto(-k,-q)\).  Then form
uncompressed tensor-product Brillouin means for at least two new grid
sizes from the V60 oracle.  Record convergence diagnostics and the
smallest sampled Hessian singular values.  Do not infer a continuum
sign until an interval quadrature or a Fourier-tail estimate is
smaller than the observed margin; if the inherited V49 strip remains
too weak, go upstream to acquire a real-domain derivative or
cellwise interval bound.
\end{decision}

\begin{assessment}[Progress after compact V60]
\label{ass:n1v60-status}
The compulsory companion audit is current.  The formerly
corner-specific one-loop response algebra is now a coefficient-exact
arbitrary-momentum oracle.  Its exact specialization reproduces all
V48 means and satisfies every inverse-jet identity; two non-corner
executions exercise the new path with small numerical residuals.
\end{assessment}

\begin{firewall}[Claim boundary after compact V60]
\label{fire:n1v60-claim}
V60 does not give an interval enclosure of a non-corner response,
a certified Brillouin integral, a continuum response sign, higher-loop
or nonperturbative remainder control, SCIC or ADMC, sector
energy--entropy constants, a continuum Yang--Mills measure,
reflection positivity, exponential clustering, or the Yang--Mills
mass gap.
\end{firewall}

\subsection{Parent record}

This V60 note was formed by appending the preceding material to the
validated compact V59, whose record was
\[
\begin{array}{c|c}
\text{lines}&22075\\
\text{bytes}&750408\\
\text{SHA--256}&
\texttt{AACEF6B54245542726D0D597400DC0C0A42EA6025F04C66244BDEE606831399A}.
\end{array}
\]

% =============================================================================
% END APPENDED COMPACT-SERIES V60 MATERIAL
% =============================================================================


% =============================================================================
% BEGIN APPENDED COMPACT-SERIES V61 MATERIAL
% =============================================================================

\section{Compact-series V61: conjugate reflection and the first
non-corner Brillouin grids}
\label{sec:n1v61}

\subsection{Purpose}

V60 exposed the complete response at arbitrary real loop momentum,
but its two non-corner samples had complex second components.  Before
forming a Brillouin sum, one must justify the pairing of \(k\) and
\(-k\).  This note proves the required conjugate-reflection law and
then executes the complete, uncompressed \(M=3\) and \(M=4\)
tensor-product grids.

\subsection{Coefficientwise reflection}

Every Hessian, cubic, and quartic Laurent coefficient in the V60
oracle is real.  If \(X_r(k)\) is an actual order-\(r\) jet
reconstructed by (60.3) or (60.4), then
\[
             X_r(-k)=(-1)^r\overline{X_r(k)}.               \tag{61.1}
\]
For the stored cubic and quartic arrays, only the order-one
reconstruction carries the explicit \(i\); for the Hessian,
\((i\alpha_1)^r\) supplies the same parity.

\begin{lemma}[Reflection is preserved by inverse jets]
\label{lem:n1v61-inverse-reflection}
If the Hessian jets satisfy (61.1) and \(H_0(k)\) is invertible at
\(k\) and \(-k\), then the covariance jets of (60.6) satisfy
\[
             G_r(-k)=(-1)^r\overline{G_r(k)},\qquad r=0,1,2.
                                                               \tag{61.2}
\]
\end{lemma}

\begin{proof}
For \(r=0\),
\[
 G_0(-k)=H_0(-k)^{-1}
 =\overline{H_0(k)}^{-1}=\overline{G_0(k)}.
\]
Insert this and \(H_1(-k)=-\overline{H_1(k)}\) in
\(G_1=-G_0H_1G_0\), obtaining
\(G_1(-k)=-\overline{G_1(k)}\).  In each term of the formula for
\(G_2\), the sum of derivative orders is two, hence the two
order-one signs cancel and (61.2) follows.
\end{proof}

\begin{theorem}[Complete response conjugate reflection]
\label{thm:n1v61-response-reflection}
On the common real invertibility set of the Hessians,
\[
 {\cal R}_e^{[r]}(-k)
 =
 (-1)^r\overline{{\cal R}_e^{[r]}(k)},
 \qquad e\in\{{\rm W},*\},\quad r=0,1,2.                    \tag{61.3}
\]
The same identity holds for
\(\Delta{\cal R}={\cal R}_*-{\cal R}_{\rm W}\).
\end{theorem}

\begin{proof}
In every term of the tadpole (60.9) and bubble (60.11), the sum of
the jet orders of all differentiated factors is \(r\).  Apply
(61.1) to the vertices and Lemma
\ref{lem:n1v61-inverse-reflection} to the covariances.  Entrywise
complex conjugation commutes with products and trace, while the
product of the reflection signs is \((-1)^r\).  Addition of the
tadpole and bubble proves (61.3); subtraction proves the last
claim.
\end{proof}

\begin{corollary}[Reality of even Brillouin means]
\label{cor:n1v61-real-mean}
Let \(\Gamma\) be any finite momentum grid invariant under
\(k\mapsto-k\).  Then the arithmetic mean of
\({\cal R}_e^{[r]}\) over \(\Gamma\) is real for even \(r\) and purely
imaginary for odd \(r\).  In particular, conjugate pairing may be
used for even response orders without discarding information.
\end{corollary}

\begin{proof}
Partition \(\Gamma\) into two-point reflection orbits and fixed
points.  On a two-point orbit, (61.3) makes the sum twice the real
part when \(r\) is even and twice \(i\) times the imaginary part when
\(r\) is odd.  At a fixed point the same identity forces the
corresponding property directly.
\end{proof}

\subsection{The \(M=3\) and \(M=4\) grids}

For \(M\geq2\), define
\[
 \Gamma_M:=
 \left\{\frac{2\pi}{M}(n_1,n_2,n_3,n_4):
 0\leq n_\mu<M\right\},\qquad
 \langle F\rangle_M:=M^{-4}\sum_{k\in\Gamma_M}F(k).          \tag{61.4}
\]
Unlike the rejected twelve-orbit shortcut of V39--V48, the V61
executable evaluates every one of the \(M^4\) points independently.

\begin{theorem}[Uncompressed non-corner response means]
\label{thm:n1v61-grid-means}
Complex-double evaluation of the coefficient-exact V60 oracle gives
\[
\begin{array}{c|rrr}
M&\langle\Delta{\cal R}^{[0]}\rangle_M&
\langle\Delta{\cal R}^{[1]}\rangle_M&
\langle\Delta{\cal R}^{[2]}\rangle_M\\ \hline
3&160.378666632812&-2.60\,10^{-14}-1.36\,10^{-14}i&
-318.121568758848\\
4&149.777012461765&-7.34\,10^{-15}-1.13\,10^{-15}i&
-314.844179178182.
\end{array}                                                   \tag{61.5}
\]
The omitted imaginary parts of the even means have modulus below
\(7.9\,10^{-14}\).  The action means are
\[
\begin{array}{c|rr|rr}
M&\langle{\cal R}_{\rm W}^{[0]}\rangle_M&
\langle{\cal R}_{*}^{[0]}\rangle_M&
\langle{\cal R}_{\rm W}^{[2]}\rangle_M&
\langle{\cal R}_{*}^{[2]}\rangle_M\\ \hline
3&-213.721636798864&-53.3429701660519&
233.959629470320&-84.1619392885277\\
4&-205.228628815809&-55.4516163540438&
222.706451932013&-92.1377272461690.
\end{array}                                                   \tag{61.6}
\]
Thus the second endpoint-minus-Wilson mean is negative on both new
grids.
\end{theorem}

\begin{proof}
Evaluate (60.2)--(60.12) at all \(81\) points of \(\Gamma_3\) and
all \(256\) points of \(\Gamma_4\), accumulate before division by
\(M^4\), and subtract the two action means.  Direct lookup of every
reflected grid point verifies (61.3) to the errors recorded below.
The displayed negative entries are ordinary floating numbers, so
this proves only the stated finite computations, not an
interval-certified sign.
\end{proof}

The difference of the last entries of (61.5) is
\[
 \langle\Delta{\cal R}^{[2]}\rangle_4
 -
 \langle\Delta{\cal R}^{[2]}\rangle_3
 \approx 3.277389580666.                                    \tag{61.7}
\]
This is far smaller than either magnitude, but two grids do not
provide a rigorous quadrature remainder.

\subsection{Conditioning and execution checks}

The largest conjugate-reflection errors over all points and all three
orders were
\[
\begin{array}{c|cc}
M&{\rm W}&*\\ \hline
3&5.93049819914503\,10^{-12}&7.68996344661428\,10^{-13}\\
4&3.53804401919888\,10^{-12}&6.97635738824464\,10^{-13}.
\end{array}                                                   \tag{61.8}
\]
The smallest sampled singular values of the action Hessians and
largest sampled condition numbers were
\[
\begin{array}{c|cc|cc}
M&\sigma_{\min}(H_{\rm W})&\sigma_{\min}(H_*)&
\kappa_2(H_{\rm W})&\kappa_2(H_*)\\ \hline
3&0.0431264506132826&0.952813445840593&
97.4589251378733&217.245092106306\\
4&0.0392013950421565&0.901729472084626&
97.4053571259141&220.407934844882.
\end{array}                                                   \tag{61.9}
\]
The largest inverse-jet residual infinity norm was
\(1.738474930978581\,10^{-12}\).  These diagnostics show no sampled
singularity, but they are not lower bounds between grid points.

\begin{record}[V61 Brillouin-grid certificate]
\label{rec:n1v61-executable}
The verifier is
\[
 \texttt{scripts/verify\_ym\_v61\_brillouin\_grids.py}.      \tag{61.10}
\]
Its source record is
\[
\begin{array}{c|c}
\text{lines}&173\\
\text{bytes}&6224\\
\text{SHA--256}&
\texttt{FDC39A195E810D11F7FB5B0000B250D98ED56296906F4517651344FAE97A847B}.
\end{array}                                                   \tag{61.11}
\]
It locks the V60 source digest
\[
 \texttt{7D5BFEE7FB834DFB5132FA2E1A9C5BC3CED7D57D366C557C8F370B6AD24D6FC8}
                                                               \tag{61.12}
\]
through the emitted input record.  The canonical output-payload
SHA--256 is
\[
 \boxed{
 \texttt{20F97BD43DC8F49193C681A3BFEE8A3E7452626479A6D1989E77847F6D4B8861}.}
                                                               \tag{61.13}
\]
\end{record}

\begin{decision}[V62 goes upstream to a real-domain remainder]
\label{dec:n1v61-next}
The V49 complex strip is too narrow to turn these small-\(M\) values
into a useful Fourier-tail certificate.  Go upstream on the real
domain.  Compile Laurent coefficient \(1\)-norms for Hessian
derivatives and vertex jets, combine them with a rigorous uniform
lower singular-value bound, and derive an explicit coordinatewise
second-derivative bound for
\(\Delta{\cal R}^{[2]}\).  Use the periodic composite-trapezoid
remainder, or a cellwise interval refinement if the global bound is
too large.  Additional grids may guide subdivision but must not be
used as a proof of the continuum sign.
\end{decision}

\begin{assessment}[Progress after compact V61]
\label{ass:n1v61-status}
The conjugate-reflection pairing is proved from the real Laurent
coefficients and inverse recursion.  The first two non-corner,
uncompressed tensor grids have been executed with controlled
numerical residuals; both retain a large negative second response
difference.  The active obstruction is now a rigorous real-domain
quadrature remainder, not a missing response evaluator.
\end{assessment}

\begin{firewall}[Claim boundary after compact V61]
\label{fire:n1v61-claim}
Neither (61.5) nor their apparent stabilization certifies the
Brillouin integral or its sign.  V61 supplies no interval arithmetic,
uniform Hessian singular-value floor, higher-loop or nonperturbative
remainder control, SCIC or ADMC, continuum measure, reflection
positivity, exponential clustering, or Yang--Mills mass gap.
\end{firewall}

\subsection{Parent record}

This V61 note was formed by appending the preceding material to the
validated compact V60, whose record was
\[
\begin{array}{c|c}
\text{lines}&22478\\
\text{bytes}&765479\\
\text{SHA--256}&
\texttt{6F145D0871A5EC33C78DE20B410FB1F80458CDB31A7D85053C84EFCCCFAF5CAC}.
\end{array}
\]

% =============================================================================
% END APPENDED COMPACT-SERIES V61 MATERIAL
% =============================================================================


% =============================================================================
% BEGIN APPENDED COMPACT-SERIES V62 MATERIAL
% =============================================================================

\section{Compact-series V62: a rigorous global real-domain majorant
and its quantitative failure}
\label{sec:n1v62}

\subsection{Purpose and outcome}

V61 isolated a real-domain quadrature remainder as the next
bottleneck.  This note carries out the most direct global estimate:
entrywise Laurent coefficient norms, the proved V18 real Hessian
floor, a bivariate inverse-derivative recursion, and the periodic
composite-trapezoid remainder.  Every constant is explicit and
rational before the final use of \(\pi^2<10\).

The resulting theorem is negative but useful.  It proves a valid
\(M^{-2}\) remainder bound, and proves quantitatively that this
global norm route is unusable: its required grid size is of order
\(10^{15}\).  The next acquisition must preserve local conditioning
and matrix structure.

\subsection{Coefficient stocks}

For a real-coefficient Laurent matrix
\[
 X(k)=\sum_{\alpha\in{\mathbb Z}^4}X_\alpha e^{i\alpha\cdot k},
\]
an external order \(r\), an internal axis \(\mu\), and an internal
order \(m\), define the entrywise stock
\[
 {\cal L}_{r,m}^{(\mu)}(X)
 :=
 \sum_{\alpha}\sum_{a,b}
 |(X_\alpha)_{ab}|\,|\alpha_1|^r|\alpha_\mu|^m.             \tag{62.1}
\]
The convention \(0^0=1\) is used.  Since
\(\|C\|_{\rm op}\leq\|C\|_F\leq\sum_{a,b}|C_{ab}|\),
\[
 \sup_{k\in{\mathbb T}^4}
 \left\|
 \partial_q^r\partial_{k_\mu}^m
 X(k+q\widehat e_1)\big|_{q=0}
 \right\|_{\rm op}
 \leq {\cal L}_{r,m}^{(\mu)}(X).                            \tag{62.2}
\]
For a stored V58 or V59 jet, the already-taken external derivative
is held fixed and only the factor \(|\alpha_\mu|^m\) is inserted.

\begin{lemma}[Rationality and validity of the stocks]
\label{lem:n1v62-stocks}
All stocks used below are nonnegative rational numbers and satisfy
(62.2) for the action Hessian, quartic jets, and both cubic
orientations.
\end{lemma}

\begin{proof}
Every relevant Laurent support is finite and every coefficient is
rational.  Differentiation multiplies a coefficient by an integral
power of \(i\alpha_\nu\), whose modulus is
\(|\alpha_\nu|\).  Apply the triangle inequality to the finite
Laurent sum and then the displayed entrywise bound on operator norm.
\end{proof}

\subsection{A mixed inverse majorant}

V18 proved on the whole real torus
\[
 H_{\rm W}(k)\succeq\frac{49}{1250}I.                       \tag{62.3}
\]
The endpoint adds a positive Gram matrix, hence the same floor holds
for \(H_*\).  Put
\[
 \lambda:=\frac{49}{1250},\qquad g_{0,0}:=\lambda^{-1}
 =\frac{1250}{49}.                                         \tag{62.4}
\]
For one internal axis, write
\[
 h_{r,m}:={\cal L}_{r,m}^{(\mu)}(H_e).
\]
Define recursively, for \(0\leq r,m\leq2\) and \(r+m>0\),
\[
 g_{r,m}:=
 g_{0,0}
 \sum_{\substack{0\leq a\leq r,\ 0\leq b\leq m\\(a,b)\ne(0,0)}}
 {r\choose a}{m\choose b}
 h_{a,b}g_{r-a,m-b}.                                       \tag{62.5}
\]

\begin{theorem}[Bivariate covariance derivative bound]
\label{thm:n1v62-inverse-bound}
For both actions and every real \(k\),
\[
 \left\|
 \partial_q^r\partial_{k_\mu}^m
 G_e(k+q\widehat e_1)\big|_{q=0}
 \right\|_{\rm op}
 \leq g_{r,m},
 \qquad 0\leq r,m\leq2.                                    \tag{62.6}
\]
\end{theorem}

\begin{proof}
The case \((0,0)\) follows from (62.3).  Apply
\(\partial_q^r\partial_{k_\mu}^m\) to \(H_eG_e=I\), expand by the
bivariate Leibniz rule, and isolate \(H_eG_{r,m}\).  Multiplication
by \(G_{0,0}\), submultiplicativity, (62.2), and induction on \(r+m\)
give exactly (62.5).
\end{proof}

\subsection{Second internal derivatives of the second response}

Let \(q_{r,m},p^+_{r,m},p^-_{r,m}\) denote the stocks of the
quartic and two cubic banks.  Twice differentiating the V60 tadpole
gives
\[
 B_{{\cal T},e}^{(\mu)}
 :=
 \frac{45}{2}\left(
 g_{0,2}q_{2,0}
 +2g_{0,1}q_{2,1}
 +g_{0,0}q_{2,2}\right).                                   \tag{62.7}
\]
For the bubble, let
\[
 {\cal C}_s(j):=
 \{(u_1,\ldots,u_j)\in{\mathbb N}_0^j:
 u_1+\cdots+u_j=s\},\qquad
 {s\choose u}:=\frac{s!}{u_1!\cdots u_j!}.                 \tag{62.8}
\]
The complete four-factor product rule gives
\[
\begin{split}
 B_{{\cal B},e}^{(\mu)}
 :={}&45|c_e|
 \sum_{(a,b,c)\in{\cal C}_2(3)}{2\choose a,b,c}
 \sum_{(u,v,w,x)\in{\cal C}_2(4)}{2\choose u,v,w,x}\\
 &\hspace{14mm}\times
 g_{0,u}p^+_{a,v}g_{b,w}p^-_{c,x}.                         \tag{62.9}
\end{split}
\]

\begin{theorem}[Global coordinatewise response majorant]
\label{thm:n1v62-response-bound}
For every real \(k\),
\[
 \left|\partial_{k_\mu}^2{\cal R}_e^{[2]}(k)\right|
 \leq B_e^{(\mu)}
 :=B_{{\cal T},e}^{(\mu)}+B_{{\cal B},e}^{(\mu)},           \tag{62.10}
\]
and
\[
 \left|\partial_{k_\mu}^2\Delta{\cal R}^{[2]}(k)\right|
 \leq B_\Delta^{(\mu)}
 :=B_{\rm W}^{(\mu)}+B_*^{(\mu)}.                           \tag{62.11}
\]
\end{theorem}

\begin{proof}
Equation (62.7) is the two-factor Leibniz rule applied to (60.9).
Equation (62.9) first distributes the two external derivatives among
the two vertices and shifted covariance, and then distributes the
two internal derivatives among all four bubble factors.  Each trace
is bounded by \(45\) times the product of operator norms.
Theorem~\ref{thm:n1v62-inverse-bound} and Lemma
\ref{lem:n1v62-stocks} give (62.10).  The triangle inequality between
the endpoint and Wilson actions gives (62.11).
\end{proof}

The compiled coordinate bounds are
\[
\begin{array}{c|c|c}
\mu&B_\Delta^{(\mu)}&\text{decimal scale}\\ \hline
1&
\dfrac{1102707759395564472095239553480638526764193625}
 {13315318287362}&8.28149756241411\,10^{31}\\[6pt]
2&
\dfrac{1212775447515720066925652418152723066975327875}
 {13315318287362}&9.10812209924268\,10^{31}\\[6pt]
3&
\dfrac{48069610556353286067482027895236895238691375}
 {512127626437}&9.38625609611916\,10^{31}\\[6pt]
4&
\dfrac{25658720551460454283890458816053643646196125}
 {271741189538}&9.44233761362569\,10^{31}.
\end{array}                                                  \tag{62.12}
\]
Their exact sum is
\[
 \boxed{
 B_\Sigma:=
 \sum_{\mu=1}^4B_\Delta^{(\mu)}
 =
 \frac{4822570388398032236686057178896149408609107375}
 {13315318287362}
 \approx3.62182133714016\,10^{32}.}                         \tag{62.13}
\]

\subsection{The global trapezoid certificate and its failure scale}

\begin{lemma}[Four-dimensional normalized trapezoid error]
\label{lem:n1v62-trapezoid}
If a \(2\pi\)-periodic \(C^2\) function \(f\) obeys
\(\sup|\partial_{k_\mu}^2f|\leq B_\mu\), then
\[
 \left|
 \frac1{M^4}\sum_{k\in\Gamma_M}f(k)
 -\frac1{(2\pi)^4}\int_{[0,2\pi]^4}f(k)\,dk
 \right|
 \leq\frac{\pi^2}{3M^2}\sum_{\mu=1}^4B_\mu.                 \tag{62.14}
\]
\end{lemma}

\begin{proof}
The normalized one-dimensional composite trapezoid rule has error at
most
\[
 \frac1{2\pi}\frac{2\pi}{12}
 \left(\frac{2\pi}{M}\right)^2\sup|f''|
 =\frac{\pi^2}{3M^2}\sup|f''|.
\]
Apply the four coordinate quadratures successively and telescope
their difference from the integral.  Every normalized integration
or averaging operator has norm one in the supremum norm, so the four
coordinate errors add.
\end{proof}

\begin{corollary}[Literal V62 continuum-mean enclosure]
\label{cor:n1v62-enclosure}
For \(f=\Delta{\cal R}^{[2]}\),
\[
 \left|\langle f\rangle_M-\langle f\rangle\right|
 \leq\frac{\pi^2B_\Sigma}{3M^2}
 <\frac{10B_\Sigma}{3M^2}.                                 \tag{62.15}
\]
\end{corollary}

\begin{proof}
Use Theorem~\ref{thm:n1v62-response-bound}, Lemma
\ref{lem:n1v62-trapezoid}, and \(\pi^2<10\).
\end{proof}

Exact integer arithmetic in the strict last bound of (62.15) gives
\[
\begin{array}{c|r}
\text{desired error}&\text{smallest integer }M
\text{ certified by V62}\\ \hline
<300&2006052324879822\\
<100&3474584549333520\\
<1&34745845493335192.
\end{array}                                                   \tag{62.16}
\]
In particular, the global estimate cannot use the \(M=3,4\) values
of V61 to certify a continuum sign.

\begin{theorem}[The global entrywise route is computationally
inadmissible]
\label{thm:n1v62-negative}
The combination of the common floor (62.3), entrywise stocks (62.1),
and global supremum propagation (62.5) cannot yield a feasible
sign certificate through (62.15).  Even allowing a remainder almost
as large as the observed \(M=4\) magnitude requires more than
\(2\times10^{15}\) points per coordinate.
\end{theorem}

\begin{proof}
The first row of (62.16) is a necessary grid size for the literal
V62 upper bound to be below \(300\), while V61 gives
\(|\langle f\rangle_4|\approx314.844\).  Such a grid contains
more than \(10^{61}\) four-dimensional points.  Hence direct
execution is infeasible, and the bound does not certify the sign at
any computed grid.
\end{proof}

The dominant loss is the endpoint bubble: its four coordinate
majorants range from \(8.28\,10^{31}\) to \(9.44\,10^{31}\), whereas
the Wilson bounds are about \(8.62\,10^{20}\).  This is not evidence
that the endpoint response is intrinsically that large.  It is the
cost of replacing every matrix coefficient by its absolute value and
reusing \(\lambda^{-1}=1250/49\) at every differentiated inverse
slot, despite the much better endpoint conditioning sampled in V61.

\subsection{Executable record and next acquisition}

\begin{record}[V62 real-domain-majorant certificate]
\label{rec:n1v62-executable}
The verifier is
\[
 \texttt{scripts/verify\_ym\_v62\_real\_domain\_majorant.py}. \tag{62.17}
\]
Its source record is
\[
\begin{array}{c|c}
\text{lines}&287\\
\text{bytes}&9498\\
\text{SHA--256}&
\texttt{BF407B68B66663C250FA537BBC34E88285B6441130372D08AC4407FC6105B6A1}.
\end{array}                                                   \tag{62.18}
\]
The canonical output-payload SHA--256 is
\[
 \boxed{
 \texttt{E8752321FCD1772911E53085F07FD1638090931562FF489DF88A9945C1D6684C}.}
                                                               \tag{62.19}
\]
\end{record}

\begin{decision}[V63 replaces global stocks by local structure]
\label{dec:n1v62-next}
Do not increase \(M\) under (62.15).  First exploit the positive
endpoint Gram correction to prove a substantially larger endpoint
floor, using its exact Laurent matrix rather than inheriting the
Wilson floor.  Second, retain spectral or Frobenius matrix norms of
the assembled coefficient blocks instead of entrywise sums.  If the
resulting global bound is still unusable, partition the torus,
certify a local Hessian floor on each box by a center singular value
minus a Laurent Lipschitz radius, and propagate interval or ball
matrix bounds cellwise.  This is the upstream repair demanded by
Theorem~\ref{thm:n1v62-negative}.
\end{decision}

\begin{assessment}[Progress after compact V62]
\label{ass:n1v62-status}
A complete rational real-domain derivative majorant and a valid
\(M^{-2}\) quadrature theorem are now populated.  Their numerical
evaluation proves that the naive global absolute-value route loses
roughly thirty orders of useful scale.  This closes that route and
pinpoints the endpoint inverse recursion as the dominant place to
preserve structure.
\end{assessment}

\begin{firewall}[Claim boundary after compact V62]
\label{fire:n1v62-claim}
V62 does not certify the continuum response sign.  Its enormous
remainder is an upper bound, not an estimate of the true error or
response size.  No higher-loop or nonperturbative comparison, SCIC,
ADMC, continuum Yang--Mills measure, reflection positivity,
exponential clustering, or Yang--Mills mass gap is proved.
\end{firewall}

\subsection{Parent record}

This V62 note was formed by appending the preceding material to the
validated compact V61, whose record was
\[
\begin{array}{c|c}
\text{lines}&22742\\
\text{bytes}&775103\\
\text{SHA--256}&
\texttt{88D009EEAB35CC3E34BEBA946A24F83B0843045EF9693ED90EF38FF9CFBD4DBA}.
\end{array}
\]

% =============================================================================
% END APPENDED COMPACT-SERIES V62 MATERIAL
% =============================================================================

% =============================================================================
% BEGIN APPENDED COMPACT-SERIES V63 MATERIAL
% =============================================================================

\section{Compact-series V63: an exact dyadic-alias endpoint floor}
\label{sec:n1v63}

\subsection{Purpose and outcome}

V62 showed that the inherited Wilson floor is the dominant avoidable
loss in the endpoint inverse recursion.  The endpoint completion has
the exact owner factorization
\[
 H_{\rm rr}(k)=K_R(k)^*K_R(k),\qquad
 K_R=E D_R,                                                   \tag{63.1}
\]
where \(E\) takes the four fine forward differences of each
plaquette curl.  The grid experiment which motivated this note
suggested the factor \(2\) below, but no grid value is used in its
proof.

The main result is the all-momentum Loewner inequality
\[
 \boxed{H_{\rm rr}(k)\succeq2D_R(k)^*D_R(k).}                 \tag{63.2}
\]
Consequently the chosen comparison endpoint is at least seven times
the Wilson Hessian.  The proof is an exact sixteen-alias argument.
It isolates the only alias whose scalar Laplacian multiplier can be
less than \(2\), and controls it by a convex estimate on the other
fifteen aliases.

\subsection{Fine aliases and the vertical block constraint}

Choose the representative
\(k=(k_1,\ldots,k_4)\in[-\pi,\pi]^4\).  For
\(\epsilon\in\{0,1\}^4\), put
\[
 p^\epsilon_\mu=\frac{k_\mu+2\pi\epsilon_\mu}{2},\qquad
 q^\epsilon_\mu=e^{ip^\epsilon_\mu}-1,\qquad
 b^\epsilon_\mu=1+e^{ip^\epsilon_\mu},                       \tag{63.3}
\]
and
\[
 \lambda_\epsilon=\sum_{\mu=1}^4|q^\epsilon_\mu|^2,\qquad
 C_\epsilon=\lambda_\epsilon I-q^\epsilon(q^\epsilon)^*,
 \qquad B_\epsilon=\operatorname{diag}(b^\epsilon_\mu).       \tag{63.4}
\]
Thus \(C_\epsilon\) is the fine curl Gram matrix on the alias
\(p^\epsilon\).

\begin{lemma}[Exact alias energy identities]
\label{lem:n1v63-alias-energy}
Let \(x\in{\mathbb C}^{45}\) be a reduced parity-tree fibre and
extend it to the \(64\) parity-edge coordinates by setting the
fifteen tree and four designated second-constituent coordinates to
zero.  If \(a^\epsilon\in{\mathbb C}^4\) are its normalized
sixteen-point parity Fourier amplitudes, then
\[
 \sum_\epsilon B_\epsilon a^\epsilon=0,                      \tag{63.5}
\]
and
\[
\begin{split}
 x^*D_R^*D_Rx
   &=16\sum_\epsilon(a^\epsilon)^*C_\epsilon a^\epsilon,\\
 x^*H_{\rm rr}x
   &=16\sum_\epsilon
          \lambda_\epsilon(a^\epsilon)^*C_\epsilon a^\epsilon.
\end{split}                                                   \tag{63.6}
\]
Moreover, with \(Q_\mu=e^{ik_\mu}-1\),
\[
                    B_\epsilon q^\epsilon=Q
             \quad\hbox{for every }\epsilon.                 \tag{63.7}
\]
\end{lemma}

\begin{proof}
Write a fine site uniquely as \(2n+s\), with
\(s\in\{0,1\}^4\).  The parity Fourier transform is
\[
 a^\epsilon_\mu
 =\frac1{16}\sum_s e^{-ip^\epsilon\cdot s}a_\mu(s).
                                                                    \tag{63.8}
\]
The straight two-link block in direction \(\mu\), based at \(2n\),
is \(a_\mu(0)+a_\mu(e_\mu)\).  The first term is a parity-tree
coordinate and the second is its designated second constituent.
Both are zero in a vertical fibre.  Fourier inversion therefore
gives (63.5).

On alias \(p^\epsilon\), the oriented \(\mu\nu\) curl is
\[
 q^\epsilon_\mu a^\epsilon_\nu
 -q^\epsilon_\nu a^\epsilon_\mu.
\]
The elementary exterior-algebra identity
\[
 \sum_{\mu<\nu}|q_\mu a_\nu-q_\nu a_\mu|^2
 =a^*(|q|^2I-qq^*)a                                      \tag{63.9}
\]
gives the first line of (63.6) by Parseval.  A forward
\(\rho\)-difference of the curl multiplies it by
\(q^\epsilon_\rho\); summing over \(\rho\) supplies the additional
factor \(\lambda_\epsilon\) and proves the second line.  Finally,
componentwise,
\[
 b^\epsilon_\mu q^\epsilon_\mu
 =(1+e^{ip^\epsilon_\mu})(e^{ip^\epsilon_\mu}-1)
 =e^{ik_\mu}-1,
\]
which is (63.7).
\end{proof}

\subsection{The scalar fifteen-alias estimate}

Put
\[
 u_\mu:=|q^0_\mu|^2
       =4\sin^2(k_\mu/4),\qquad
 L:=\lambda_0=\sum_\mu u_\mu.                               \tag{63.10}
\]
Thus \(0\leq u_\mu\leq2\).  Identify an alias with its support
\(S=\{\mu:\epsilon_\mu=1\}\).  Then
\[
 \lambda_S=L+4|S|-2u_S,\qquad
 u_S:=\sum_{\mu\in S}u_\mu.                                 \tag{63.11}
\]
For a fixed coordinate \(\mu\), write
\[
 c_{S\mu}:=|b^S_\mu|^2
 =\begin{cases}
   4-u_\mu,&\mu\notin S,\\
   u_\mu,&\mu\in S.
  \end{cases}                                                \tag{63.12}
\]

\begin{lemma}[Fifteen-alias reciprocal bound]
\label{lem:n1v63-reciprocal}
If \(0<L<2\), then every nonempty \(S\) has
\(\lambda_S>2\), and
\[
 \boxed{
 \sum_{\varnothing\ne S\subset\{1,2,3,4\}}
 \frac{c_{S\mu}}{\lambda_S(\lambda_S-2)}
 \leq
 \frac{c_{0\mu}}{L(2-L)}.}                                  \tag{63.13}
\]
\end{lemma}

\begin{proof}
For nonempty \(S\),
\[
 \lambda_S-2
 =L-2+4|S|-2u_S
 \geq2-L>0,                                                  \tag{63.14}
\]
because \(u_S\leq L\) and \(|S|\geq1\).  Also
\(c_{S\mu}\leq c_{0\mu}=4-u_\mu\), since \(u_\mu<2\).
It is therefore enough to prove
\[
 F_L(u):=
 \sum_{\varnothing\ne S}
 f(L+4|S|-2u_S)
 \leq\frac1{L(2-L)},\qquad
 f(t):=\frac1{t(t-2)}.                                      \tag{63.15}
\]

For fixed \(L\), the vector \(u\) belongs to the simplex
\(\{u_\mu\geq0:\sum u_\mu=L\}\).  On \(t>2\),
\[
 f''(t)=
 \frac{2\{3(t-1)^2+1\}}{t^3(t-2)^3}>0.                      \tag{63.16}
\]
Hence \(F_L\) is convex.  It is also symmetric in the four
coordinates.  Since
\[
 u=\sum_{\mu=1}^4\frac{u_\mu}{L}\,Le_\mu,
\]
convexity and symmetry give
\[
                         F_L(u)\leq F_L(Le_1).               \tag{63.17}
\]
Direct enumeration according to whether \(1\in S\) gives
\[
\begin{split}
 F_L(Le_1)
 ={}&f(4-L)+A(L),\\
 A(L)
 :={}&3\{f(4+L)+f(8-L)\}\\
    &+3\{f(8+L)+f(12-L)\}
      +f(12+L)+f(16-L).                                     \tag{63.18}
\end{split}
\]
The pole term satisfies the exact identity
\[
 \frac1{L(2-L)}-f(4-L)=\frac2{L(4-L)}.                       \tag{63.19}
\]

The function \(f\) is decreasing on \((2,\infty)\).  If
\(0<L\leq1\), the six terms in (63.18) are bounded by
\[
 A(L)\leq
 3\left(\frac18+\frac1{24}\right)
 +3\left(\frac1{48}+\frac1{80}\right)
 +\frac1{120}+\frac1{168}
 =\frac{43}{70}.                                            \tag{63.20}
\]
But \(2/[L(4-L)]\geq2/3\), and
\[
                  \frac23-\frac{43}{70}=\frac{11}{210}>0.   \tag{63.21}
\]
If \(1\leq L<2\), the sharper endpoint bounds give
\[
\begin{split}
 A(L)
 &\leq
 3\left(\frac1{15}+\frac1{24}\right)
 +3\left(\frac1{63}+\frac1{80}\right)
 +\frac1{143}+\frac1{168}\\
 &=\frac{33879}{80080}<\frac12
 \leq\frac2{L(4-L)},                                        \tag{63.22}
\end{split}
\]
where
\[
 \frac12-\frac{33879}{80080}
 =\frac{6161}{80080}>0.                                     \tag{63.23}
\]
Equations (63.17)--(63.23) prove (63.15), and multiplication
by \(c_{0\mu}\) proves (63.13).
\end{proof}

\subsection{The exact Loewner comparison}

\begin{theorem}[Dyadic curl-detail Poincare inequality]
\label{thm:n1v63-loewner}
For every real coarse momentum,
\[
 \boxed{
 K_R(k)^*K_R(k)\succeq2D_R(k)^*D_R(k).}                     \tag{63.24}
\]
\end{theorem}

\begin{proof}
Use Lemma~\ref{lem:n1v63-alias-energy}.  If \(L\geq2\), then
\(\lambda_0\geq2\), and every nonzero alias contains a flipped
coordinate whose contribution to \(\lambda_\epsilon\) is at least
\(2\).  Thus every \(\lambda_\epsilon-2\) is nonnegative, and the
claim follows directly from (63.6).  If \(L=0\), then
\(C_0=0\), while every nonzero alias has
\(\lambda_\epsilon\geq4\); the claim again follows.

It remains to take \(0<L<2\), and put \(\delta=2-L\).  For
\(S\ne\varnothing\), define
\[
 W_S:=(\lambda_S-2)C_S.                                     \tag{63.25}
\]
This form is positive semidefinite, has kernel
\(\mathbb Cq^S\), and on its orthogonal complement its inverse is
\[
                   W_S^+=\frac1{\lambda_S(\lambda_S-2)}I.   \tag{63.26}
\]
Let \(y\perp Q\).  By (63.7),
\[
 \ip{q^S}{B_S^*y}=\ip{B_Sq^S}{y}=\ip{Q}{y}=0.               \tag{63.27}
\]
Weighted Cauchy--Schwarz, (63.5), and (63.26) therefore give
\[
\begin{split}
 |\ip{y}{B_0a^0}|^2
 &=
 \left|\sum_{S\ne\varnothing}\ip{B_S^*y}{a^S}\right|^2\\
 &\leq
 \left\{\sum_{S\ne\varnothing}(a^S)^*W_Sa^S\right\}
 \sum_{\mu=1}^4|y_\mu|^2
 \left\{\sum_{S\ne\varnothing}
 \frac{c_{S\mu}}{\lambda_S(\lambda_S-2)}\right\}\\
 &\leq
 \left\{\sum_{S\ne\varnothing}(a^S)^*W_Sa^S\right\}
 \frac{\sum_\mu c_{0\mu}|y_\mu|^2}{L\delta}.                \tag{63.28}
\end{split}
\]
Here the last line is Lemma~\ref{lem:n1v63-reciprocal}.

Because \(L<2\), every \(b^0_\mu\) is nonzero.  Under the
invertible change \(\widetilde y=B_0^*y\), the condition
\(y\perp Q=B_0q^0\) becomes
\(\widetilde y\perp q^0\), while
\[
 \sum_\mu c_{0\mu}|y_\mu|^2=\|\widetilde y\|^2.
\]
Taking the supremum of (63.28) over \(y\perp Q\) yields
\[
\begin{split}
 \sum_{S\ne\varnothing}(a^S)^*W_Sa^S
 &\geq
 L\delta\,
 \operatorname{dist}(a^0,\mathbb Cq^0)^2\\
 &=\delta\,(a^0)^*C_0a^0.                                  \tag{63.29}
\end{split}
\]
Thus
\[
 \sum_\epsilon(\lambda_\epsilon-2)
       (a^\epsilon)^*C_\epsilon a^\epsilon\geq0.             \tag{63.30}
\]
Multiplication by the Parseval factor \(16\) and use of (63.6)
prove (63.24).
\end{proof}

\begin{corollary}[Sevenfold endpoint floor]
\label{cor:n1v63-endpoint-floor}
In the action normalization used since V20,
\[
 H_{\rm W}=\frac13D_R^*D_R,\qquad
 H_*=H_{\rm W}+H_{\rm rr}.
\]
Hence, for \(0\leq t\leq1\),
\[
 \boxed{
 H_t:=H_{\rm W}+tH_{\rm rr}
 \succeq(1+6t)H_{\rm W}.}                                   \tag{63.31}
\]
In particular,
\[
 \boxed{
 H_*\succeq7H_{\rm W}
 \succeq\frac{343}{1250}I,\qquad
 \|H_*^{-1}\|_{\rm op}\leq\frac{1250}{343}.}                \tag{63.32}
\]
If the exact V18 algebraic floor is retained, the middle lower
bound can instead be written \(7\lambda_*/3\).
\end{corollary}

\begin{proof}
Theorem~\ref{thm:n1v63-loewner} gives
\(H_{\rm rr}\succeq2D_R^*D_R=6H_{\rm W}\).
This proves (63.31).  At \(t=1\), combine it with the rational
Wilson floor \(H_{\rm W}\succeq49I/1250\).
\end{proof}

\begin{remark}[What V63 repairs]
\label{rem:n1v63-repair}
V62 used \(\|H_*^{-1}\|\leq1250/49\), inherited from the
Wilson endpoint.  Equation (63.32) replaces this by
\[
                    \frac{1250}{343}\approx3.6443148688,     \tag{63.33}
\]
an exact factor-seven improvement.  The result also supplies a
monotone real-axis floor throughout the direct endpoint homotopy.
It does not by itself improve the complex strip, because Loewner
order is a real Hermitian statement.
\end{remark}

\subsection{Executable record and next acquisition}

\begin{record}[V63 alias-reduction certificate]
\label{rec:n1v63-executable}
The verifier is
\[
 \texttt{scripts/verify\_ym\_v63\_alias\_reduction.py}.      \tag{63.34}
\]
Its source record is
\[
\begin{array}{c|c}
\text{lines}&216\\
\text{bytes}&8836\\
\text{SHA--256}&
\texttt{64B615B746184174B1E1F7A9370E0BC75201132FDD410D5ED1BB032C3A1C01E0}.
\end{array}                                                  \tag{63.35}
\]
It locks the incidence and endpoint-construction sources at
\[
\begin{split}
&\texttt{DE678E5CAE7671B90C8788AF9A489444D8A63B4F32EA78F08BC68AA0131B0AE0},\\
&\texttt{F2BCB2385017A322354F1C72EBD1E435D6E1D85C1E462049A7A62FB386544885}.
\end{split}                                                   \tag{63.36}
\]
The exact rational rows in the payload are (63.20)--(63.23).
Three independent non-corner cross-checks reproduce the Wilson and
owner energies with maximum absolute residuals
\(1.14\,10^{-13}\) and \(1.82\,10^{-12}\), respectively; their
largest block-constraint residual is \(1.60\,10^{-15}\).
The canonical output-payload SHA--256 is
\[
 \boxed{
 \texttt{F16391AD7B246A6B0615EFFF62E6DEB99CB86D9ABAEEF0C8F5E0918707EFA5CE}.}
                                                               \tag{63.37}
\]
\end{record}

\begin{decision}[V64 recompiles the response majorant]
\label{dec:n1v63-next}
Re-run the V62 inverse-derivative recursion with the action-specific
endpoint floor \(343/1250\), retaining the Wilson floor \(49/1250\).
First quantify the exact gain with the same coefficient stocks, so
that it cannot be confused with a change of norm.  Then replace
entrywise coefficient sums by assembled matrix norms.  If the
result still cannot certify the continuum response sign, pass to
cellwise Hessian floors and local matrix-ball propagation rather
than increasing the global grid.
\end{decision}

\begin{assessment}[Progress after compact V63]
\label{ass:n1v63-status}
The first upstream repair requested by V62 is complete.  The
positive owner endpoint now has an exact all-momentum floor seven
times the Wilson floor, obtained without interval subdivision or
sampling.  The active quantitative question is how much of this
factor survives the differentiated covariance recursion and whether
matrix-aware local norms remove the remaining absolute-value loss.
\end{assessment}

\begin{firewall}[Claim boundary after compact V63]
\label{fire:n1v63-claim}
Theorem~\ref{thm:n1v63-loewner} is a flat, real, reduced-fibre
quadratic inequality.  It is not a nonlinear compact-action
convexity theorem, a continuum response-sign certificate, a
higher-loop or nonperturbative comparison, SCIC, ADMC, a continuum
Yang--Mills measure, reflection positivity, exponential clustering,
or a Yang--Mills mass gap.
\end{firewall}

\subsection{Parent record}

This V63 note was formed by appending the preceding material to the
validated compact V62, whose record was
\[
\begin{array}{c|c}
\text{lines}&23101\\
\text{bytes}&787188\\
\text{SHA--256}&
\texttt{635C7B5BB3ED2F5E64C22C568390387D2F68F5F561F4F9449BD24BBB5B444EB1}.
\end{array}
\]

% =============================================================================
% END APPENDED COMPACT-SERIES V63 MATERIAL
% =============================================================================

% =============================================================================
% BEGIN APPENDED COMPACT-SERIES V64 MATERIAL
% =============================================================================

\section{Compact-series V64: action-specific response majorants}
\label{sec:n1v64}

\subsection{Purpose}

V63 proved the endpoint covariance bound
\(\|H_*^{-1}\|_{\rm op}\leq1250/343\), while the Wilson bound
remains \(1250/49\).  This note changes only those two zeroth inverse
bounds in the exact V62 recursion.  All Laurent coefficient stocks,
Leibniz multiplicities, trace dimensions, vertex normalizations, and
the periodic trapezoid theorem are held fixed.  The calculation
therefore measures precisely the gain due to the new coercivity
theorem.

\subsection{Action-indexed recursion}

For \(e\in\{{\rm W},*\}\), put
\[
 \lambda_{\rm W}:=\frac{49}{1250},\qquad
 \lambda_*:=\frac{343}{1250},\qquad
 g^e_{0,0}:=\lambda_e^{-1}.                                 \tag{64.1}
\]
For each internal coordinate \(\mu\), retain the V62 Hessian stocks
\(h^e_{r,m}\), and define
\[
 g^e_{r,m}:=
 g^e_{0,0}
 \sum_{\substack{0\leq a\leq r,\ 0\leq b\leq m\\
                  (a,b)\ne(0,0)}}
 {r\choose a}{m\choose b}
 h^e_{a,b}g^e_{r-a,m-b}.                                   \tag{64.2}
\]
The tadpole and bubble formulas are exactly (62.7) and (62.9), with
\(g\) replaced by \(g^e\).

\begin{proposition}[Validity of the action-specific recompilation]
\label{prop:n1v64-validity}
For each action, (64.2) bounds every mixed covariance derivative
appearing in the second response.  The resulting coordinate bounds
\(B_e^{(\mu)}\) satisfy
\[
 \left|\partial_{k_\mu}^2{\cal R}^{[2]}_e(k)\right|
 \leq B_e^{(\mu)},\qquad
 \left|\partial_{k_\mu}^2\Delta{\cal R}^{[2]}(k)\right|
 \leq B_{\rm W}^{(\mu)}+B_*^{(\mu)}.                        \tag{64.3}
\]
\end{proposition}

\begin{proof}
The proof of Theorem~\ref{thm:n1v62-inverse-bound} uses the Hessian
floor only in the initial value \(g_{0,0}\).  Apply that proof
separately with the two floors in (64.1).  The differentiated
tadpole and bubble identities and their trace estimates are
unchanged, so the proof of
Theorem~\ref{thm:n1v62-response-bound} then gives (64.3).
\end{proof}

\subsection{Exact output}

\begin{theorem}[The V63 floor removes five orders from the global
majorant]
\label{thm:n1v64-output}
With precisely the V62 entrywise coefficient norm, the four
coordinate bounds in (64.3) are
\[
\begin{array}{c|c}
\mu&B_\Delta^{(\mu)}\\ \hline
1&7.04144592202065\,10^{26}\\
2&7.74360441341259\,10^{26}\\
3&7.98027076076802\,10^{26}\\
4&8.02775032459804\,10^{26}.
\end{array}                                                  \tag{64.4}
\]
Their exact sum is
\[
 \boxed{
 B_{\Sigma,64}
 =
 \frac{4823838968658102649347622374476807457588451625}
 {1566533881189851938}
 \approx3.07930714207993\,10^{27}.}                         \tag{64.5}
\]
Relative to the V62 sum \(B_{\Sigma,62}\), the exact ratio is
\[
 \frac{B_{\Sigma,64}}{B_{\Sigma,62}}
 =
 \frac{38590711749264821194780978995814459660707613}
 {4538964668997120756911023528319624654187622988491}
 \approx8.50209564592\,10^{-6}.                             \tag{64.6}
\]
Thus the total bound improves by a factor
\[
                 \frac{B_{\Sigma,62}}{B_{\Sigma,64}}
                 \approx117618.060493.                      \tag{64.7}
\]
\end{theorem}

\begin{proof}
All quantities in (64.1)--(64.3) are rational.  The verifier
reconstructs the same V60 Hessian, quartic, and oriented cubic
banks used in V62; executes the finite recursion (64.2) with the
two declared floors; and evaluates every multinomial term in
(62.7) and (62.9) in \(\mathbb Q\).  It independently repeats the
calculation with the common V62 floor, checks that all four Wilson
rows are identical, and checks that all four endpoint rows strictly
decrease.  Exact addition and division give (64.5) and (64.6);
the decimal displays are not used in the inequalities.
\end{proof}

The individual endpoint improvements are tightly grouped:
\[
\begin{array}{c|c}
\mu&B_{*,62}^{(\mu)}/B_{*,64}^{(\mu)}\\ \hline
1&117610.897033\\
2&117621.352459\\
3&117618.393171\\
4&117621.343125.
\end{array}                                                   \tag{64.8}
\]
The Wilson rows remain of order \(8.62\,10^{20}\), and are therefore
negligible beside the endpoint rows in (64.4).  The dominant
remaining loss is no longer the use of the Wilson covariance floor
at the endpoint.

\subsection{The remaining global-grid obstruction}

\begin{corollary}[Improved but still inadmissible trapezoid sizes]
\label{cor:n1v64-grid}
Using the same rigorous inequality \(\pi^2<10\) as V62, the smallest
integers certified by the global bound are
\[
\begin{array}{c|r}
\text{desired error}&\text{smallest }M\\ \hline
<300&5849318233854\\
<100&10131316370673\\
<1&101313163706729.
\end{array}                                                   \tag{64.9}
\]
In particular, even the new global bound cannot certify the
continuum sign from any feasible tensor grid.
\end{corollary}

\begin{proof}
Insert (64.5) into the strict V62 trapezoid bound
\[
 \left|\langle f\rangle_M-\langle f\rangle\right|
 <\frac{10B_{\Sigma,64}}{3M^2}.
\]
Exact integer square-root comparisons give (64.9).  The smallest
listed tensor grid contains more than \(10^{51}\) points, so the
literal global route remains computationally inadmissible.
\end{proof}

\begin{record}[V64 action-specific-majorant certificate]
\label{rec:n1v64-executable}
The verifier is
\[
 \texttt{scripts/verify\_ym\_v64\_action\_specific\_majorant.py}. \tag{64.10}
\]
Its source record is
\[
\begin{array}{c|c}
\text{lines}&340\\
\text{bytes}&11317\\
\text{SHA--256}&
\texttt{D12A8757AA95125BCC5795B54FB2C19A0EFBE9475606205FD111FE6B14AADAF4}.
\end{array}                                                   \tag{64.11}
\]
The canonical output-payload SHA--256 is
\[
 \boxed{
 \texttt{8B060EABFAB610EB8BF0171F406F34EE6C2704F7305789E8F8782BAFFDA021E3}.}
                                                               \tag{64.12}
\]
The payload contains every exact action, axis, tadpole, bubble,
Hessian-stock, inverse-stock, quartic-stock, and oriented
cubic-stock row, as well as the three strict grid checks.
\end{record}

\begin{decision}[V65 passes to assembled matrix norms]
\label{dec:n1v64-next}
At V65 first perform the mandatory contemporaneous SM/NS audit.
Then hold the action-specific floors fixed and replace each
entrywise coefficient stock by a norm of the assembled derivative
matrix.  The first target is a rigorous Frobenius or spectral
coefficient bound for each Hessian and vertex derivative, evaluated
before multiplying the factors in the inverse recursion.  If that
global matrix-aware bound still exceeds a feasible remainder, form
local momentum boxes and propagate center-radius matrix balls; do
not return to a larger tensor grid under a global supremum.
\end{decision}

\begin{assessment}[Progress after compact V64]
\label{ass:n1v64-status}
The V63 coercivity theorem yields an exact \(117618\)-fold aggregate
improvement in the inherited V62 response majorant.  This is a
substantial quantitative repair, but (64.9) proves that it does not
make global entrywise quadrature viable.  The next bottleneck is now
unambiguously the loss of cancellations and matrix geometry in the
coefficient norm.
\end{assessment}

\begin{firewall}[Claim boundary after compact V64]
\label{fire:n1v64-claim}
V64 remains a global upper-bound calculation.  It does not certify
the continuum response sign, and the large bound is not an estimate
of the true response.  No higher-loop or nonperturbative comparison,
SCIC, ADMC, continuum Yang--Mills measure, reflection positivity,
exponential clustering, or Yang--Mills mass gap is proved.
\end{firewall}

\subsection{Parent record}

This V64 note was formed by appending the preceding material to the
validated compact V63, whose record was
\[
\begin{array}{c|c}
\text{lines}&23542\\
\text{bytes}&801223\\
\text{SHA--256}&
\texttt{D91787B1DCC4E770D518E96F1CDC9E90F808D1979B2C956369CE48A1BA2769AA}.
\end{array}
\]

% =============================================================================
% END APPENDED COMPACT-SERIES V64 MATERIAL
% =============================================================================

% =============================================================================
% BEGIN APPENDED COMPACT-SERIES V65 MATERIAL
% =============================================================================

\section{Compact-series V65: companion audit and matrix-aware
Laurent majorants}
\label{sec:n1v65}

\subsection{Mandatory companion audit}

The active companion notes were inspected read-only before choosing
the V65 norm.  The exact audit snapshot is
\[
\begin{array}{c|r|r|l}
\text{file}&\text{lines}&\text{bytes}&\text{SHA--256}\\ \hline
\texttt{SM\_V100.tex}&27627&944425&
\texttt{82DDF1C141E9A72BD9DF6CAF1BE24D10D715B19EF9BA689554272A6FC01A10FB}\\
\texttt{NS\_V26.tex}&5319&278283&
\texttt{82C887F8C2C69E4F28FAD7C3DC8DE5B9099CB5A4BF431EBA1FFF3E91935978AD}.
\end{array}                                                    \tag{65.1}
\]
Here the abbreviated file names denote
\(\texttt{Renewal\_SM\_Einstein\_Continuum\_Research\_Notes\_V100.tex}\)
and
\(\texttt{Renewal\_NS\_Stability\_Research\_Notes\_V26.tex}\).

SM V100 records the same structural order exposed by V63:
factorization and the constraint should be used to prove a Loewner
floor before inverse or heat estimates are propagated.  It imports
no Yang--Mills numerical constant into the curved SM--Einstein
problem.  This independently supports retaining the V63 floor
before changing norms.

NS V26 compares a full tensor norm with the norm on the rank-one
inputs which actually occur in one part of its argument.  It also
refuses that improvement where the absorption input is not rank
one.  The transferable lesson is therefore conditional: a
Yang--Mills response norm may be reduced only after the actual
tadpole or bubble contraction class has been proved.  No NS kernel
constant is transferable to the gauge response.

\begin{audit}[V65 companion decision]
\label{aud:n1v65-companions}
No numerical constant is imported from either companion note.
V65 first uses the unconditional assembled Frobenius norm of every
Laurent coefficient.  The next iteration will examine the actual
cyclic trace contractions, and will not declare a smaller
channel-restricted vertex norm without proving that every input to
that vertex lies in the restricted class.
\end{audit}

\begin{proof}
The file records in (65.1) identify the audited sources.  The SM
statement is an abstract proof-order principle, while the NS
statement distinguishes two input classes and explicitly retains
the larger norm on the unrestricted one.  Neither supplies an
identity between its operators and the V60 response matrices.
Thus only the stated proof architecture, not a coefficient, can be
used here.
\end{proof}

\subsection{Rational Frobenius coefficient stocks}

Let a scalar or matrix-valued derivative bank have the finite
Laurent expansion
\[
                         X(k)=\sum_\alpha X_\alpha e^{i\alpha\cdot k}.
\]
Fix \(N=10^{12}\), and define the rational upper square root
\[
 u_\alpha(X):=
 \min\left\{\frac mN:m\in{\mathbb N}_0,\quad
       m^2\geq N^2\sum_{a,b}|(X_\alpha)_{ab}|^2\right\}.       \tag{65.2}
\]
For a stored component of a vertex jet, \(X_\alpha\) means the
matrix of that component.  Put
\[
 {\cal F}_{r,m}^{(\mu)}(X)
 :=
 \sum_\alpha
 |\alpha_1|^r|\alpha_\mu|^m u_\alpha(X).                     \tag{65.3}
\]

\begin{lemma}[Matrix-aware Laurent bound]
\label{lem:n1v65-frobenius}
Every stock in (65.3) is rational and satisfies
\[
 \sup_{k\in{\mathbb T}^4}
 \left\|
 \partial_q^r\partial_{k_\mu}^m
 X(k+q\widehat e_1)\big|_{q=0}
 \right\|_{\rm op}
 \leq{\cal F}_{r,m}^{(\mu)}(X).                              \tag{65.4}
\]
Moreover, if ({\cal A}(X)=\{\alpha:X_\alpha\ne0\}), then
\[
 {\cal F}_{r,m}^{(\mu)}(X)
 \leq{\cal L}_{r,m}^{(\mu)}(X)
 +\frac1N\sum_{\alpha\in{\cal A}(X)}
 |\alpha_1|^r|\alpha_\mu|^m,                                \tag{65.5}
\]
where \({\cal L}\) is the V62 entrywise stock.
\end{lemma}

\begin{proof}
The defining integer in (65.2) is found by exact integer square-root
comparison, so \(u_\alpha\) is rational and
\[
 \|X_\alpha\|_{\rm op}
 \leq\|X_\alpha\|_F
 \leq u_\alpha(X).
\]
Differentiate the finite Laurent sum and apply the triangle
inequality to obtain (65.4).  Finally,
\[
 \|X_\alpha\|_F
 \leq\sum_{a,b}|(X_\alpha)_{ab}|.
\]
The rational ceiling in (65.2) exceeds the exact Frobenius norm by
less than \(N^{-1}\).  Sum this error with the weights in (65.3) to
obtain (65.5).  The verifier checks the strict improvement for every
final action/axis row directly.
\end{proof}

\begin{remark}[Ceiling convention]
\label{rem:n1v65-ceiling}
Because of the \(N^{-1}\) ceiling, the literal rational stock can in
principle exceed an exceptionally small entrywise stock by less
than the number of supported coefficients times \(N^{-1}\).
No proof below relies on a termwise version of (65.5); the exact
compiler verifies the strict inequality for all eight completed
response bounds.
\end{remark}

\subsection{Recompiled response bound}

Use the action-specific floors of (64.1), replace every Hessian,
quartic, and oriented cubic stock in (64.2) and the V62
tadpole/bubble formulas by (65.3), and make no other change.

\begin{theorem}[Frobenius-stock response majorant]
\label{thm:n1v65-output}
The resulting coordinatewise difference bounds obey
\[
\begin{array}{c|c}
\mu&B_{\Delta,65}^{(\mu)}\\ \hline
1&<1.119\,10^{22}\\
2&<1.146\,10^{22}\\
3&<1.158\,10^{22}\\
4&<1.163\,10^{22}.
\end{array}                                                    \tag{65.6}
\]
Their exact sum is
\[
 \boxed{
 B_{\Sigma,65}
 =
 \frac{
 23887982889349003290631919370585852491747594337365545494068516688309644910317581352603}
 {521092351331343680000000000000000000000000000000000000000000000}
 <4.585\,10^{22}.}                                           \tag{65.7}
\]
It satisfies the exact integer comparison
\[
 67171<
 \frac{B_{\Sigma,64}}{B_{\Sigma,65}}
 <67172.                                                       \tag{65.8}
\]
\end{theorem}

\begin{proof}
Lemma~\ref{lem:n1v65-frobenius} supplies a valid operator-norm
majorant in every slot of the V62 proof.  Therefore
Proposition~\ref{prop:n1v64-validity} applies without alteration.
The verifier groups all entries carrying the same Laurent exponent,
sums their rational squares, rounds the square root upward by
(65.2), and performs the complete inverse, tadpole, and bubble
recursions in rational arithmetic.  It separately replays V64 with
entrywise stocks and verifies that all eight action/axis rows
strictly decrease.  Exact cross multiplication gives (65.7) and
(65.8); the rounded displays in (65.6) are strict rational upper
bounds.
\end{proof}

The Wilson row improvement factors lie between \(2998\) and
\(3007\).  The endpoint factors lie between \(62951\) and \(69030\).
This disparity is expected: the larger endpoint coefficient bank
contains more simultaneous entrywise contributions for the
Frobenius grouping to combine quadratically.

\begin{corollary}[The global grid is still excluded]
\label{cor:n1v65-grid}
The strict global trapezoid bound now requires
\[
\begin{array}{c|r}
\text{desired error}&\text{smallest certified }M\\ \hline
<300&22568938633\\
<100&39090548385\\
<1&390905483843.
\end{array}                                                    \tag{65.9}
\]
Thus assembled Frobenius coefficient norms do not make the global
tensor grid feasible.
\end{corollary}

\begin{proof}
Insert the exact number (65.7) into (62.15).  The verifier performs
the strict integer square-root tests.  Even the first grid contains
more than \(10^{41}\) points.
\end{proof}

\begin{record}[V65 Frobenius-majorant certificate]
\label{rec:n1v65-executable}
The verifier is
\[
 \texttt{scripts/verify\_ym\_v65\_frobenius\_majorant.py}.     \tag{65.10}
\]
Its source record is
\[
\begin{array}{c|c}
\text{lines}&404\\
\text{bytes}&13753\\
\text{SHA--256}&
\texttt{FE21F5FEB042DDB4A7665FBF9C3DFEFB0202FE69C2327A6DFC8C3F229002E117}.
\end{array}                                                    \tag{65.11}
\]
The canonical output-payload SHA--256 is
\[
 \boxed{
 \texttt{1F09A12D0C07612E3C6E145940132451B17EBE8F268774A47CB667A98BB03F77}.}
                                                               \tag{65.12}
\]
The payload also locks the V57, V59, and V60 inputs recorded in
V64, and contains every rational stock and completed action/axis
row.
\end{record}

\begin{decision}[V66 uses the actual trace channels]
\label{dec:n1v65-next}
The two successful global repairs, V63 coercivity and V65
coefficient assembly, have removed about ten orders of magnitude
without approaching a feasible global remainder.  Following the NS
audit distinction, do not introduce another generic tensor norm.
Expand the differentiated tadpole and bubble as the actual finite
cyclic traces.  Group complete terms with common Laurent phase
before taking absolute values, and use Hilbert--Schmidt pairing or
direct coefficient contraction on that proved input class.  If the
fully contracted global coefficient stock remains too large, use
the same evaluator on local momentum boxes with center-radius
inverse bounds.
\end{decision}

\begin{assessment}[Progress after compact V65]
\label{ass:n1v65-status}
The mandatory audit is complete.  A rigorous matrix-aware Laurent
majorant reduces the V64 bound by a further factor between \(67171\)
and \(67172\), but the remaining \(10^{22}\) scale proves that even
assembled factor norms discard too much structure.  The next
mathematically justified norm is the norm of the complete contracted
response channel.
\end{assessment}

\begin{firewall}[Claim boundary after compact V65]
\label{fire:n1v65-claim}
V65 proves only a sharper real-domain derivative upper bound.  It
does not certify the continuum response sign.  The companion audits
do not transfer SM or NS constants into Yang--Mills.  No
higher-loop or nonperturbative comparison, SCIC, ADMC, continuum
Yang--Mills measure, reflection positivity, exponential clustering,
or Yang--Mills mass gap is proved.
\end{firewall}

\subsection{Parent record}

This V65 note was formed by appending the preceding material to the
validated compact V64, whose record was
\[
\begin{array}{c|c}
\text{lines}&23774\\
\text{bytes}&809506\\
\text{SHA--256}&
\texttt{F3209E68ECECA6C23C908C5F2D0E9456597E43D6D694A4466BE650FF199CBE82}.
\end{array}
\]

% =============================================================================
% END APPENDED COMPACT-SERIES V65 MATERIAL
% =============================================================================

% =============================================================================
% BEGIN APPENDED COMPACT-SERIES V66 MATERIAL
% =============================================================================

\section{Compact-series V66: trace-ideal response bounds}
\label{sec:n1v66}

\subsection{Purpose}

V65 bounded a trace of four matrices by \(45\) times the product of
four operator-norm majorants.  That inequality ignores the fact
that the two vertex factors already carry Frobenius bounds.  This
note uses the actual tadpole and bubble contraction classes, as
required by the V65 companion-audit decision.  No Laurent,
covariance, or vertex stock is changed.

\subsection{Two dimension-saving trace inequalities}

\begin{lemma}[Tadpole and bubble trace ideals]
\label{lem:n1v66-traces}
For \(45\times45\) complex matrices,
\[
 |\Tr(GQ)|
 \leq\sqrt{45}\,\|G\|_{\rm op}\|Q\|_F,                       \tag{66.1}
\]
and
\[
 |\Tr(GP\widetilde G N)|
 \leq
 \|G\|_{\rm op}\|\widetilde G\|_{\rm op}
 \|P\|_F\|N\|_F.                                             \tag{66.2}
\]
The same statements hold when any factor is a mixed derivative of
the corresponding response factor.
\end{lemma}

\begin{proof}
Hilbert--Schmidt duality and
\(\|G\|_F\leq\sqrt{45}\|G\|_{\rm op}\) give (66.1).  For (66.2),
cyclicity and Hilbert--Schmidt duality give
\[
 |\Tr(GP\widetilde G N)|
 \leq\|GP\widetilde G\|_F\|N\|_F
 \leq
 \|G\|_{\rm op}\|P\|_F
 \|\widetilde G\|_{\rm op}\|N\|_F.
\]
No positivity or rank assumption is used, so differentiated
matrices are allowed.
\end{proof}

For an exact rational computation use the upward approximation
\[
 \sqrt{45}<
 \frac{2683281573}{400000000},                               \tag{66.3}
\]
certified by squaring integers.  Replace the V65 tadpole prefactor
\(45/2\) by the right side of (66.3) divided by \(2\), and replace
the bubble prefactor \(45|c_e|\) by \(|c_e|\).

\begin{theorem}[Trace-ideal global majorant]
\label{thm:n1v66-output}
With the V63 action-specific floors and the V65 rational
Frobenius stocks, (66.1)--(66.3) give
\[
\begin{array}{c|c}
\mu&B_{\Delta,66}^{(\mu)}\\ \hline
1&<2.486\,10^{20}\\
2&<2.545\,10^{20}\\
3&<2.573\,10^{20}\\
4&<2.585\,10^{20}.
\end{array}                                                    \tag{66.4}
\]
Their exact sum is
\[
 \boxed{
 B_{\Sigma,66}
 =
 \frac{
 23887982904404831340593767345543072234918166112795799867997215116778283379067581352603}
 {23449155809910465600000000000000000000000000000000000000000000000}
 <1.019\,10^{21}.}                                           \tag{66.5}
\]
The exact ratio to V65 has decimal value
\[
               \frac{B_{\Sigma,65}}{B_{\Sigma,66}}
               \approx44.9999999716.                         \tag{66.6}
\]
\end{theorem}

\begin{proof}
Distribute the two external and two internal derivatives exactly as
in (62.7) and (62.9).  Each differentiated tadpole term has the
form in (66.1), and each differentiated bubble term has the form in
(66.2), irrespective of which factor receives a derivative.
Insert the already proved operator bounds for covariance
derivatives and Frobenius bounds for vertex derivatives, then sum
the same multinomial coefficients.  This proves that the compiler's
rational output is an upper bound.  The verifier separately replays
the V65 dimension-prefactor calculation and checks strict
improvement in all eight action/axis rows.  Exact rational addition
gives (66.5).  Equation (66.6) is displayed only to indicate scale;
the payload retains its exact rational ratio.
\end{proof}

The improvement is within \(3\times10^{-8}\) of the full factor
\(45\), because every row is bubble dominated.  Thus V66 identifies
and removes a genuine dimension artifact, rather than changing the
Hessian or vertex data.

\begin{corollary}[Global quadrature remains impossible]
\label{cor:n1v66-grid}
The strict V62 trapezoid theorem with (66.5) requires
\[
\begin{array}{c|r}
\text{desired error}&\text{smallest certified }M\\ \hline
<300&3364378732\\
<100&5827274900\\
<1&58272748995.
\end{array}                                                    \tag{66.7}
\]
\end{corollary}

\begin{proof}
These are exact integer square-root comparisons for
\(10B_{\Sigma,66}/(3M^2)\).  The first grid already contains more
than \(10^{38}\) points.
\end{proof}

\begin{record}[V66 trace-ideal certificate]
\label{rec:n1v66-executable}
The verifier is
\[
 \texttt{scripts/verify\_ym\_v66\_trace\_ideal\_majorant.py}. \tag{66.8}
\]
Its source record is
\[
\begin{array}{c|c}
\text{lines}&420\\
\text{bytes}&14389\\
\text{SHA--256}&
\texttt{D2E19E7BF21D44B82C0B9578106C9752CAE39A740B615F3987E05A31FA174032}.
\end{array}                                                    \tag{66.9}
\]
The canonical output-payload SHA--256 is
\[
 \boxed{
 \texttt{5CE1FF3D0FCD6039D18429DF40627BB9DF8429B6406B7D325A9BE0BFE002E186}.}
                                                               \tag{66.10}
\]
The exact new-to-V65 ratio stored there is
\[
 \frac{
 23887982904404831340593767345543072234918166112795799867997215116778283379067581352603}
 {1074959230020705148078436371676363362128641745181449547233083250973934020964291160867135}.
                                                               \tag{66.11}
\]
\end{record}

\begin{decision}[V67 localizes before differentiating inverses]
\label{dec:n1v66-next}
The remaining global loss comes from applying one worst-case
covariance-derivative recursion over the entire torus.  Partition
the Brillouin torus into boxes.  At each box center evaluate the
assembled Hermitian Hessian, certify its smallest eigenvalue by a
rational residual or interval enclosure, and subtract a Laurent
operator-norm radius for the box.  Propagate inverse and response
derivatives with that local floor, and sum the local quadrature
remainders.  The first V67 acquisition is a reproducible
center-radius theorem and a pilot partition; a floating center
eigenvalue without its residual and radius is inadmissible.
\end{decision}

\begin{assessment}[Progress after compact V66]
\label{ass:n1v66-status}
Trace-ideal duality removes the full artificial bubble dimension
factor and reduces the V65 majorant by essentially \(45\).  The
remaining bound is still nineteen orders above the observed
response scale.  Since the exact action and vertex data are now
assembled before taking norms, the next upstream loss is the use of
a single global inverse recursion rather than local conditioning.
\end{assessment}

\begin{firewall}[Claim boundary after compact V66]
\label{fire:n1v66-claim}
V66 is a sharper derivative upper bound, not a continuum
response-sign certificate.  It does not control local interval
remainders, higher-loop or nonperturbative errors, SCIC, ADMC, the
continuum Yang--Mills measure, reflection positivity, exponential
clustering, or a Yang--Mills mass gap.
\end{firewall}

\subsection{Parent record}

This V66 note was formed by appending the preceding material to the
validated compact V65, whose record was
\[
\begin{array}{c|c}
\text{lines}&24063\\
\text{bytes}&819983\\
\text{SHA--256}&
\texttt{3588991D7B16AF87C2064050B7520CE6A15BD47327725D114A9D46ED94559F5E}.
\end{array}
\]

% =============================================================================
% END APPENDED COMPACT-SERIES V66 MATERIAL
% =============================================================================

% =============================================================================
% BEGIN APPENDED COMPACT-SERIES V67 MATERIAL
% =============================================================================

\section{Compact-series V67: a center--radius theorem and the
uniform-box obstruction}
\label{sec:n1v67}

\subsection{A certifiable local inverse}

The center of a momentum box may be evaluated numerically only if
the numerical inverse is accompanied by a residual.  The following
form records the exact requirement.

\begin{lemma}[Residual-certified center inverse]
\label{lem:n1v67-center}
Let \(H_c\) be Hermitian and positive, and let \(R_c\) be any
matrix for which
\[
 \|R_c\|_{\rm op}\leq r_c,\qquad
 \|I-R_cH_c\|_{\rm op}\leq\eta_c<1.                           \tag{67.1}
\]
Then
\[
 \sigma_{\min}(H_c)\geq\frac{1-\eta_c}{r_c}.                 \tag{67.2}
\]
If \(H(k)\) is Hermitian on a box \(B(c,\rho)\) and
\[
 \sup_{k\in B(c,\rho)}
 \|H(k)-H_c\|_{\rm op}\leq d_c,                              \tag{67.3}
\]
then
\[
 H(k)\succeq
 \left(\frac{1-\eta_c}{r_c}-d_c\right)I
 \quad(k\in B(c,\rho)).                                      \tag{67.4}
\]
\end{lemma}

\begin{proof}
For every \(v\),
\[
 (1-\eta_c)\|v\|
 \leq\|R_cH_cv\|
 \leq r_c\|H_cv\|.
\]
This proves (67.2).  Weyl's one-Lipschitz inequality for the least
Hermitian eigenvalue and (67.3) then give (67.4).
\end{proof}

\subsection{Laurent radii}

For a Hermitian Laurent matrix
\[
 H(k)=\sum_\alpha H_\alpha e^{i\alpha\cdot k},
\]
choose rational upper bounds
\[
 f_\alpha\geq\|H_\alpha\|_F
\]
by the V65 integer-square-root convention, and put
\[
 L_\mu(H):=\sum_\alpha|\alpha_\mu|f_\alpha.                  \tag{67.5}
\]

\begin{lemma}[Exact coefficient radius]
\label{lem:n1v67-radius}
If \(|k_\mu-c_\mu|\leq\rho_\mu\), then
\[
 \boxed{
 \|H(k)-H(c)\|_{\rm op}
 \leq\sum_{\mu=1}^4L_\mu(H)\rho_\mu.}                        \tag{67.6}
\]
\end{lemma}

\begin{proof}
For real \(s\), \(|e^{is}-1|\leq|s|\).  Therefore
\[
\begin{split}
 \|H(k)-H(c)\|_{\rm op}
 &\leq
 \sum_\alpha\|H_\alpha\|_{\rm op}
 |e^{i\alpha\cdot(k-c)}-1|\\
 &\leq
 \sum_\alpha f_\alpha
 \sum_\mu|\alpha_\mu|\rho_\mu
 =\sum_\mu L_\mu(H)\rho_\mu.
\end{split}
\]
\end{proof}

Combining Lemmas~\ref{lem:n1v67-center} and
\ref{lem:n1v67-radius} gives a completely rational local inverse
certificate whenever \(c\) has exactly representable phases and
\(R_c,r_c,\eta_c\) are rational enclosures.  It also applies with
outward interval bounds at a general center.

\begin{proposition}[Local derivative propagation]
\label{prop:n1v67-local-propagation}
Suppose (67.4) has positive right side \(\lambda_c\).  Replace
\(g_{0,0}\) in the V62 inverse recursion by \(\lambda_c^{-1}\).
For every differentiated Laurent bank \(X_\gamma\), replace its
global stock by
\[
 \|X_\gamma(c)\|_{\rm op}
 +\sum_{\nu=1}^4\rho_\nu
   {\cal F}_{\gamma+e_\nu}(X),                               \tag{67.7}
\]
with an outward enclosure of the center norm.  Then the V62
Leibniz recursion and the V66 trace-ideal estimates give valid
response-derivative bounds throughout the box.
\end{proposition}

\begin{proof}
Equation (67.7) is Lemma~\ref{lem:n1v67-radius} applied after the
derivatives indexed by \(\gamma\).  Equation (67.4) supplies the
inverse base case.  The proofs of
Theorem~\ref{thm:n1v62-inverse-bound} and
Lemma~\ref{lem:n1v66-traces} are pointwise and use only valid
factor bounds, so their induction applies on the box.
\end{proof}

\subsection{Exact uniform-mesh pilot}

Before building thousands of center inverses, compute the literal
coefficient radius.  For a uniform \(M\)-mesh, every Voronoi box has
coordinate half-width at most \(\pi/M<22/(7M)\).  The exact V65
Frobenius stocks give
\[
\begin{array}{c|cccc|c}
e&L_1&L_2&L_3&L_4&\sum_\mu L_\mu\\ \hline
{\rm W}&
7.293810046164&7.293810046164&7.293810046164&
7.293810046164&29.175240184656\\
*&315.407098904504&319.197174054537&
320.939578260874&321.690736354345&1277.234587574261.
\end{array}                                                    \tag{67.8}
\]
Every entry in the executable payload is rational; decimals in
(67.8) only display scale.

\begin{theorem}[Uniform coefficient-radius no-go scale]
\label{thm:n1v67-uniform-no-go}
Let
\[
 d_e(M):=\frac{22}{7M}\sum_\mu L_\mu(H_e).                   \tag{67.9}
\]
The smallest integers for which this radius is strictly less than
one half of the proved global floor are
\[
 \boxed{
 M_{\rm W}^{(1/2)}=4679,\qquad
 M_*^{(1/2)}=29258.}                                        \tag{67.10}
\]
For one quarter of the floor they are
\[
 \boxed{
 M_{\rm W}^{(1/4)}=9357,\qquad
 M_*^{(1/4)}=58516.}                                        \tag{67.11}
\]
The corresponding half-floor tensor grids contain
\[
 4679^4=479305244227681,\qquad
 29258^4=732788326502158096                                \tag{67.12}
\]
centers.
\end{theorem}

\begin{proof}
Insert the exact sums
\[
 \sum_\mu L_\mu(H_{\rm W})
 =\frac{1823452511541}{62500000000},\qquad
 \sum_\mu L_\mu(H_*)
 =\frac{1915851881361391}{1500000000000}                    \tag{67.13}
\]
and the floors \(49/1250\), \(343/1250\) into (67.9).
The conditions are
\[
 M>\frac{44}{7\lambda_e}\sum_\mu L_\mu(H_e)
\]
for the half floor, and twice that right side for the quarter
floor.  Exact numerator--denominator division gives
(67.10)--(67.11); the verifier also checks that every predecessor
integer fails or gives equality.  Integer fourth powers give
(67.12).
\end{proof}

\begin{record}[V67 center-radius pilot]
\label{rec:n1v67-executable}
The verifier is
\[
 \texttt{scripts/verify\_ym\_v67\_center\_radius.py}.         \tag{67.14}
\]
Its source record is
\[
\begin{array}{c|c}
\text{lines}&110\\
\text{bytes}&3816\\
\text{SHA--256}&
\texttt{C7809C8B24448DDD6F6234CD23115D4A017E2DA5ACC1581418C01148B8D38E7C}.
\end{array}                                                    \tag{67.15}
\]
It locks the V60 and V65 sources at
\[
\begin{split}
&\texttt{7D5BFEE7FB834DFB5132FA2E1A9C5BC3CED7D57D366C557C8F370B6AD24D6FC8},\\
&\texttt{FE21F5FEB042DDB4A7665FBF9C3DFEFB0202FE69C2327A6DFC8C3F229002E117},
\end{split}                                                    \tag{67.16}
\]
and its canonical payload SHA--256 is
\[
 \boxed{
 \texttt{92DE67B52B9AB77C60246053EF1A590F373CCCF4BF336D4566D4A12037BBD77F}.}
                                                               \tag{67.17}
\]
\end{record}

\begin{decision}[V68 goes upstream to relative derivatives]
\label{dec:n1v67-next}
Do not instantiate the uniform tensor cover (67.12).  The absolute
Hessian radius is already too expensive before a response remainder
is propagated.  Return to the Gram structure and define the
whitened derivative
\[
 {\cal T}_{e,\mu}(k)
 :=H_e(k)^{-1/2}\,\partial_\mu H_e(k)\,H_e(k)^{-1/2}.         \tag{67.18}
\]
Derive inverse-derivative recursions in terms of
\(\|{\cal T}_{e,\mu}\|\) and its covariant derivatives, so that a
large Hessian coefficient is measured relative to the local energy
it controls.  Use the sixteen-alias representation of V63 for the
endpoint and the exact Wilson fibre for the Wilson action.  Only if
the relative stocks are small should adaptive center boxes be
reintroduced.
\end{decision}

\begin{assessment}[Progress after compact V67]
\label{ass:n1v67-status}
A residual-certified local inverse theorem and a local
response-propagation rule are now rigorous.  Their first exact pilot
proves that a uniform absolute coefficient-radius cover would
require at least \(7.3\,10^{17}\) endpoint centers merely to make
the radius smaller than half the existing floor.  This isolates the
next upstream loss: absolute rather than energy-relative Hessian
derivatives.
\end{assessment}

\begin{firewall}[Claim boundary after compact V67]
\label{fire:n1v67-claim}
The thresholds (67.10)--(67.12) are no-go scales for the literal
uniform coefficient-radius criterion, not lower bounds for every
adaptive interval method.  V67 does not certify a response sign,
higher-loop control, SCIC, ADMC, a continuum Yang--Mills measure,
reflection positivity, exponential clustering, or a Yang--Mills
mass gap.
\end{firewall}

\subsection{Parent record}

This V67 note was formed by appending the preceding material to the
validated compact V66, whose record was
\[
\begin{array}{c|c}
\text{lines}&24272\\
\text{bytes}&827337\\
\text{SHA--256}&
\texttt{EE39CFBF5E271C9EEA61D6AAB617E02EEB57200F454F10A432116EAA5A96A710}.
\end{array}
\]

% =============================================================================
% END APPENDED COMPACT-SERIES V67 MATERIAL
% =============================================================================

% =============================================================================
% BEGIN APPENDED COMPACT-SERIES V68 MATERIAL
% =============================================================================

\section{Compact-series V68: relative Hessian jets and a
floor-linear covariance recursion}
\label{sec:n1v68}

\subsection{Why the inverse recursion must be whitened}

V67 proved that a global absolute Laurent radius destroys the
conditioning before local response propagation even begins.  The
same loss is present algebraically in (62.5): every induction step
introduces another copy of the worst global inverse bound.  This
section replaces that recursion by an exact identity in the energy
metric of the Hessian.

Let
\[
 H(q,k):=H_e(k+q\widehat e_1),\qquad G(q,k):=H(q,k)^{-1},
\]
and, at \(q=0\), use the bivariate notation
\[
 H_{r,m}:=\partial_q^r\partial_{k_\mu}^mH,\qquad
 G_{r,m}:=\partial_q^r\partial_{k_\mu}^mG.                 \tag{68.1}
\]
All square roots below refer to the positive matrix \(H=H_{0,0}\)
at the same real momentum.

\begin{definition}[Relative Hessian and covariance jets]
\label{def:n1v68-relative}
Put
\[
 T_{r,m}:=H^{-1/2}H_{r,m}H^{-1/2},\qquad
 U_{r,m}:=H^{1/2}G_{r,m}H^{1/2}.                           \tag{68.2}
\]
Thus \(T_{0,0}=U_{0,0}=I\).
\end{definition}

\begin{theorem}[Whitened mixed inverse identity]
\label{thm:n1v68-relative-inverse}
For every \(r,m\geq0\) with \(r+m>0\),
\[
 \boxed{
 U_{r,m}=
 -\!\!\sum_{\substack{0\leq a\leq r,\ 0\leq b\leq m\\
                      (a,b)\ne(0,0)}}
 {r\choose a}{m\choose b}\,
 T_{a,b}U_{r-a,m-b}.}                                     \tag{68.3}
\]
Consequently, if \(t_{a,b}\) majorizes
\(\|T_{a,b}\|_{\rm op}\), the recursion
\[
 u_{0,0}=1,\qquad
 u_{r,m}:=
 \sum_{\substack{0\leq a\leq r,\ 0\leq b\leq m\\
                 (a,b)\ne(0,0)}}
 {r\choose a}{m\choose b}t_{a,b}u_{r-a,m-b}                \tag{68.4}
\]
gives
\[
 \boxed{
 \|G_{r,m}\|_{\rm op}
 \leq \|H^{-1}\|_{\rm op}\,u_{r,m}.}                       \tag{68.5}
\]
In particular, the absolute inverse floor occurs only once.
\end{theorem}

\begin{proof}
Differentiate \(HG=I\) by the bivariate Leibniz rule:
\[
 0=\sum_{a=0}^r\sum_{b=0}^m
 {r\choose a}{m\choose b}H_{a,b}G_{r-a,m-b}.
\]
Isolate the term \(HG_{r,m}\).  Multiplication on the left by
\(H^{-1/2}\) and on the right by \(H^{1/2}\) converts every remaining
product into
\[
 H^{-1/2}H_{a,b}G_{r-a,m-b}H^{1/2}
 =T_{a,b}U_{r-a,m-b},
\]
which proves (68.3).  Operator-norm submultiplicativity and induction
give \(\|U_{r,m}\|_{\rm op}\leq u_{r,m}\).  Finally
\[
 G_{r,m}=H^{-1/2}U_{r,m}H^{-1/2},
\]
so (68.5) follows.
\end{proof}

The first two one-variable cases display the structural gain:
\[
\begin{split}
 U_{1,0}&=-T_{1,0},\\
 U_{2,0}&=2T_{1,0}^2-T_{2,0},\\
 \|G_{2,0}\|_{\rm op}
 &\leq\|H^{-1}\|_{\rm op}
 \bigl(2t_{1,0}^2+t_{2,0}\bigr).                           \tag{68.6}
\end{split}
\]
No commutativity has been assumed.

\begin{corollary}[Loewner formulation of a relative stock]
\label{cor:n1v68-loewner}
For Hermitian \(H_\gamma\) and \(H\succ0\), the following are
equivalent:
\[
 \|H^{-1/2}H_\gamma H^{-1/2}\|_{\rm op}\leq t
 \quad\Longleftrightarrow\quad
 -tH\preceq H_\gamma\preceq tH.                            \tag{68.7}
\]
\end{corollary}

\begin{proof}
The left inequality says that every eigenvalue of the Hermitian
matrix \(H^{-1/2}H_\gamma H^{-1/2}\) lies in \([-t,t]\).
Conjugation by \(H^{1/2}\) gives the two Loewner inequalities, and
the same argument reverses.
\end{proof}

Corollary~\ref{cor:n1v68-loewner} is important for certification:
one need not enclose a matrix square root.  A rational
sum-of-squares or interval proof of the two Laurent matrix
inequalities in (68.7) is already a relative-jet certificate.

\subsection{Deterministic scale diagnostic}

The exact V60 Laurent Hessians were evaluated on the uniform
\(4^4\) grid in \([-\pi,\pi)^4\).  At each point an eigendecomposition
was used only to diagnose the sizes in (68.2); no number in this
paragraph is used as a proved global bound.  Across both actions,
all four internal axes, and all \(0\leq r,m\leq2\), the largest
sampled nontrivial relative Hessian jet was
\[
 \max\|T_{r,m}\|_{\rm op}
 =
 \begin{cases}
 6.850248821778&({\rm Wilson}),\\
 7.604739942531&({\rm endpoint}).
 \end{cases}                                                \tag{68.8}
\]
Applying (68.4) to the separate sampled maxima gave
\[
 \max_\mu u_{2,2}^{(\mu)}
 =
 \begin{cases}
 1336.876483860&({\rm Wilson}),\\
 1643.625489051&({\rm endpoint}).
 \end{cases}                                                \tag{68.9}
\]
After the single proved inverse factors \(1250/49\) and
\(1250/343\), these are of order \(3.5\,10^4\) and \(6.0\,10^3\).
For comparison, the V65 absolute recursion produced sampled-independent
majorants of order \(7.4\,10^{11}\) and \(1.6\,10^{14}\) for the
same \((2,2)\) covariance jet.  This comparison is diagnostic, but
it shows that the upstream change attacks the observed loss rather
than merely changing its notation.

\begin{record}[V68 relative-Hessian diagnostic]
\label{rec:n1v68-executable}
The executable is
\[
 \texttt{scripts/explore\_ym\_v68\_relative\_hessian.py}.   \tag{68.10}
\]
Its source record is
\[
\begin{array}{c|c}
\text{lines}&172\\
\text{bytes}&6379\\
\text{SHA--256}&
\texttt{4A8E92B72472FEF0BFF3D819E81BF1A9D7411472CF630216C904BEEEF1EB5B8D}.
\end{array}                                                  \tag{68.11}
\]
For the \(M=4\) run its canonical payload SHA--256 is
\[
 \boxed{
 \texttt{931018B8BF5D8319C0CBE30A4204C7F96CF124B86B6AD58CEEDD66189116115B}.}
                                                               \tag{68.12}
\]
The input hashes locked by that payload are
\[
\begin{split}
&\texttt{V60: 7D5BFEE7FB834DFB5132FA2E1A9C5BC3CED7D57D366C557C8F370B6AD24D6FC8},\\
&\texttt{V51: F2BCB2385017A322354F1C72EBD1E435D6E1D85C1E462049A7A62FB386544885}.
\end{split}                                                   \tag{68.13}
\]
\end{record}

\begin{decision}[V69 certifies relative stocks without square roots]
\label{dec:n1v68-next}
Use the Gram factorizations already present in V51 and V63.  For a
row operator \(A(k)\) with \(H=A^*A\), express each
\(H_{r,m}\) by differentiated Gram factors and prove rational
Loewner bounds of the form
\[
 -t_{r,m}H\preceq H_{r,m}\preceq t_{r,m}H.                 \tag{68.14}
\]
Begin with first derivatives, where the factorization has only two
terms, and use the sixteen-alias endpoint representation before any
coefficientwise norm is taken.  Only certified \(t_{r,m}\), not the
sampled values (68.8), may enter a rigorous response remainder.
\end{decision}

\begin{assessment}[Progress after compact V68]
\label{ass:n1v68-status}
The repeated inverse-floor loss in V62--V66 has been removed at the
identity level.  A covariance derivative of any mixed order now
costs one inverse factor times a dimensionless relative recurrence.
The first deterministic diagnostic indicates reductions of roughly
seven to ten orders for the fourth mixed covariance jet.  The
remaining immediate bottleneck is a global rational certificate for
the relative Hessian stocks.
\end{assessment}

\begin{firewall}[Claim boundary after compact V68]
\label{fire:n1v68-claim}
Equations (68.3)--(68.7) are rigorous.  The mesh values
(68.8)--(68.9) are floating diagnostics, not global enclosures.
V68 therefore does not certify the continuum response sign,
higher-loop or nonperturbative control, SCIC, ADMC, construction of
the continuum Yang--Mills measure, reflection positivity,
exponential clustering, or a Yang--Mills mass gap.
\end{firewall}

\subsection{Parent record}

This V68 note was formed by appending the preceding material to the
validated compact V67, whose record was
\[
\begin{array}{c|c}
\text{lines}&24547\\
\text{bytes}&835956\\
\text{SHA--256}&
\texttt{3BFA1F0E6D66C086AD2B056C01EB8757BF22CE5952BA7A82235CDB75AB528911}.
\end{array}
\]

% =============================================================================
% END APPENDED COMPACT-SERIES V68 MATERIAL
% =============================================================================

% =============================================================================
% BEGIN APPENDED COMPACT-SERIES V69 MATERIAL
% =============================================================================

\section{Compact-series V69: a rational Gram-relative response
majorant}
\label{sec:n1v69}

\subsection{Relative control from differentiated factors}

V68 reduced covariance differentiation to relative Hessian stocks
but left their certification open.  The exact V21 construction
provides rectangular Laurent factors
\[
 A_{\rm W}=3^{-1/2}D_R,\qquad
 A_*=\begin{bmatrix}3^{-1/2}D_R\\K_R\end{bmatrix},
 \qquad H_e=A_e^*A_e.                                      \tag{69.1}
\]
For the mixed derivative convention of (68.1), write
\(A_\gamma=\partial^\gamma A\), where
\(\gamma=(r,m)\), and let
\[
 s_\gamma\geq
 \sup_k\|A_\gamma(k)H(k)^{-1/2}\|_{\rm op}.                \tag{69.2}
\]
The undifferentiated value is exactly \(s_0=1\), since
\[
 (AH^{-1/2})^*(AH^{-1/2})=I.                              \tag{69.3}
\]

\begin{lemma}[Differentiated Gram bound]
\label{lem:n1v69-gram-relative}
For every nonzero bivariate multi-index \(\gamma=(r,m)\),
\[
 \boxed{
 \|H^{-1/2}H_\gamma H^{-1/2}\|_{\rm op}
 \leq
 t_\gamma:=
 \sum_{\beta\leq\gamma}{\gamma\choose\beta}
 s_\beta s_{\gamma-\beta}.}                               \tag{69.4}
\]
\end{lemma}

\begin{proof}
The mixed Leibniz rule applied to \(H=A^*A\) gives
\[
 H_\gamma
 =\sum_{\beta\leq\gamma}{\gamma\choose\beta}
 (A_\beta)^*A_{\gamma-\beta}.
\]
After congruence by \(H^{-1/2}\), each summand is
\[
 (A_\beta H^{-1/2})^*(A_{\gamma-\beta}H^{-1/2}).
\]
The triangle inequality, submultiplicativity, and (69.2) give
(69.4).  Equation (69.3) justifies the exact undifferentiated
value and avoids applying the global floor to that factor.
\end{proof}

\subsection{A finite rational realization}

Write \(A_\gamma(k)=\sum_\alpha A_{\gamma,\alpha}
e^{i\alpha\cdot k}\).  Define
\[
 F_\gamma(A):=
 \sum_\alpha
 |\alpha_1|^r|\alpha_\mu|^m
 \|A_\alpha\|_F.                                           \tag{69.5}
\]
For the endpoint stack in (69.1), the squared coefficient norm used
by the compiler is exactly
\[
 \|(A_*)_\alpha\|_F^2
 =\frac13\|(D_R)_\alpha\|_F^2+\|(K_R)_\alpha\|_F^2.         \tag{69.6}
\]
Each square root is rounded upward to denominator \(10^{12}\).
With the proved floors \(\lambda_{\rm W}=49/1250\) and
\(\lambda_*=343/1250\), set
\[
 s_0=1,\qquad
 s_\gamma=
 {\rm Up}_{10^{12}}\!
 \left(\frac{F_\gamma(A)}{\sqrt{\lambda_e}}\right)
 \quad(\gamma\ne0).                                       \tag{69.7}
\]

\begin{proposition}[Validity of the rational factor stocks]
\label{prop:n1v69-factor-stocks}
The quantities (69.7) satisfy (69.2) on the entire real torus.
Consequently (69.4), the V68 recurrence (68.4), and
\[
 g_\gamma:=\lambda_e^{-1}u_\gamma                         \tag{69.8}
\]
are rigorous global covariance-jet bounds.
\end{proposition}

\begin{proof}
Termwise differentiation of the finite Laurent expansion and
\(\|\cdot\|_{\rm op}\leq\|\cdot\|_F\) give
\(\|A_\gamma(k)\|_{\rm op}\leq F_\gamma(A)\).
Also \(\|H^{-1/2}\|_{\rm op}\leq\lambda_e^{-1/2}\).
The two prescribed upward square-root operations in the compiler
therefore prove (69.2) with rational endpoints.  Apply
Lemma~\ref{lem:n1v69-gram-relative} and
Theorem~\ref{thm:n1v68-relative-inverse}.
\end{proof}

\subsection{Recompiled response remainder}

The V66 trace-ideal inequalities and all V65 vertex and quartic
Frobenius stocks are retained unchanged.  Only the covariance
stocks are replaced by (69.8).  The resulting four coordinate
bounds for the endpoint-minus-Wilson second response derivative are
\[
\begin{array}{c|c}
\mu&B_{\Delta,69}^{(\mu)}\\ \hline
1&1.4594991165350436\,10^{16}\\
2&1.6028160498303608\,10^{16}\\
3&1.7063482854546188\,10^{16}\\
4&1.7629022071549048\,10^{16}.
\end{array}                                                  \tag{69.9}
\]
The exact sum is
\[
 \boxed{
 B_{\Sigma,69}=
 \frac{
 115264825231911199463465415779249544577918046348335257207475543747341970351602260418163}
 {1764735000000000000000000000000000000000000000000000000000000000000000}}
                                                               \tag{69.10}
\]
and
\[
 B_{\Sigma,69}\approx6.5315656589749280\,10^{16}.           \tag{69.11}
\]
This is a rigorous improvement over V66 by the factor
\[
 \frac{B_{\Sigma,66}}{B_{\Sigma,69}}
 \approx15596.780861099169.                                \tag{69.12}
\]
With \(\pi^2<10\), the exact strict tensor-grid thresholds become
\[
\begin{array}{c|c}
\text{target error}&\text{smallest }M\\ \hline
300&26939368\\
100&46660354\\
1&466603532.
\end{array}                                                  \tag{69.13}
\]
They remain far beyond use, but they replace the V66 unit-error
threshold \(58272748995\) by \(466603532\).

\begin{record}[V69 exact compiler]
\label{rec:n1v69-executable}
The verifier is
\[
 \texttt{scripts/verify\_ym\_v69\_gram\_relative\_majorant.py}. \tag{69.14}
\]
Its record is
\[
\begin{array}{c|c}
\text{lines}&290\\
\text{bytes}&10295\\
\text{SHA--256}&
\texttt{D32638D1B599CB04D739F6267B05CA4E1B429509A9FE0C490001459040963BC6}.
\end{array}                                                   \tag{69.15}
\]
It verifies nonnegativity of every stock, the exact base
normalizations, and strict improvement over V66.  Its canonical
payload SHA--256 is
\[
 \boxed{
 \texttt{7D368A7644071809D2A7C0AF41F4A46AAFA788106FB00E7DFE3E7FBF92E2EB6D}.}
                                                               \tag{69.16}
\]
\end{record}

\begin{decision}[V70 audit and exponent-block repair]
\label{dec:n1v69-next}
V69 still takes a triangle inequality over Laurent exponents in
(69.5).  This makes some endpoint normalized factor stocks roughly
two orders larger than the V68 sampled relative Hessian jets.  At
the mandatory V70 audit, inspect the current SM and NS notes for
reusable restricted-domain or block-orthogonality arguments.  Then
retain exponent blocks, parity aliases, or row translations long
enough to bound \(A_\gamma^*A_\gamma\) relative to \(H\); do not
return to entrywise absolute sums.
\end{decision}

\begin{assessment}[Progress after compact V69]
\label{ass:n1v69-status}
A fully rational, global version of the whitened recursion is now
available, and it lowers the rigorous response-remainder majorant by
another factor exceeding \(1.5\,10^4\).  The total improvement from
V65 is approximately \(7.0\,10^5\).  The bound is still about
fourteen orders above the observed response scale.  The newly
isolated loss is the exponentwise triangle inequality inside the
differentiated Gram factor, not repeated covariance inversion.
\end{assessment}

\begin{firewall}[Claim boundary after compact V69]
\label{fire:n1v69-claim}
V69 proves a sharper global quadrature majorant, not a useful
continuum response-sign enclosure.  It does not establish
higher-loop or nonperturbative control, SCIC, ADMC, a continuum
Yang--Mills measure, reflection positivity, exponential clustering,
or a Yang--Mills mass gap.
\end{firewall}

\subsection{Parent record}

This V69 note was formed by appending the preceding material to the
validated compact V68, whose record was
\[
\begin{array}{c|c}
\text{lines}&24787\\
\text{bytes}&844092\\
\text{SHA--256}&
\texttt{B0E2441D94D57E3B6BB81042321ED729B5545388A919EB76E4A00A5AE78B415C}.
\end{array}
\]

% =============================================================================
% END APPENDED COMPACT-SERIES V69 MATERIAL
% =============================================================================

% =============================================================================
% BEGIN APPENDED COMPACT-SERIES V70 MATERIAL
% =============================================================================

\section{Compact-series V70: compulsory companion audit and the
actual Gram channel}
\label{sec:n1v70}

\subsection{Read-only companion snapshot}

In accordance with the five-version rule, the two current companion
notes in \texttt{latex\_notes} were inspected read-only.  Their
records at the time of the audit were
\[
\begin{array}{c|c|c|c}
\text{programme}&\text{file version}&\text{lines}&\text{bytes}\\ \hline
{\rm SM\!-\!Einstein}&11&3420&117893\\
{\rm Navier\!-\!Stokes}&26&5319&278283.
\end{array}                                                   \tag{70.1}
\]
The corresponding SHA--256 values were
\[
\begin{split}
{\rm SM:}\quad&
\texttt{C5BB59CBB8CD6B828F8F2D2E7FB83005D1986989CE5B2467B79F1AC9F763E05A},\\
{\rm NS:}\quad&
\texttt{82C887F8C2C69E4F28FAD7C3DC8DE5B9099CB5A4BF431EBA1FFF3E91935978AD}.
\end{split}                                                   \tag{70.2}
\]
These documents are mathematical source material only; none of
their programme instructions or claim language is imported.

SM V11 proves that a two-sided relative Loewner estimate preserves
the support and protected kernel of a positive Choi or Gram matrix.
NS V26 compares a tensor norm on all traceless inputs with its norm
on the actual rank-one inputs and records both the real gain and the
channels where no gain occurs.  The joint methodological consequence
for the present problem is precise: retain the support of the
reference Gram form, and reduce a norm only after identifying the
actual input channel.  No numerical constant from either companion
programme is transferable to Yang--Mills.

\subsection{The actual first-derivative channel}

Let \(H=A^*A\succ0\), fix one derivative \(\partial_\gamma\), and
put
\[
 X:=AH^{-1/2},\qquad
 Y_\gamma:=A_\gamma H^{-1/2},\qquad
 P:=XX^*.                                                    \tag{70.3}
\]
Then \(X^*X=I\), so \(P\) is the orthogonal projection onto
\(\operatorname{ran}X\).  Define
\[
 C_\gamma:=X^*Y_\gamma,\qquad
 Y_\gamma^\perp:=(I-P)Y_\gamma.                              \tag{70.4}
\]

\begin{lemma}[First variation sees only the compressed channel]
\label{lem:n1v70-first-channel}
The whitened first Hessian variation satisfies
\[
 \boxed{T_\gamma=C_\gamma+C_\gamma^*.}                       \tag{70.5}
\]
In particular,
\[
 \|T_\gamma\|_{\rm op}
 \leq2\|C_\gamma\|_{\rm op},                                 \tag{70.6}
\]
and \(Y_\gamma^\perp\) makes no contribution at first order.
\end{lemma}

\begin{proof}
Differentiate \(H=A^*A\):
\[
 H_\gamma=A_\gamma^*A+A^*A_\gamma.
\]
Congruence by \(H^{-1/2}\) gives
\[
 T_\gamma=Y_\gamma^*X+X^*Y_\gamma
 =C_\gamma^*+C_\gamma.
\]
Since \(X^*Y_\gamma^\perp=X^*(I-P)Y_\gamma=0\), the perpendicular
part vanishes from this identity.  The norm estimate follows from
the triangle inequality.
\end{proof}

Thus the V69 replacement of \(\|C_\gamma\|\) by
\(\|Y_\gamma\|\) was rigorous but still measured rows that the
first variation cannot see.  This is the exact Yang--Mills analogue
of the actual-input distinction found in NS V26.

\subsection{Square-root-free certificate gates}

\begin{proposition}[Finite Loewner gates for the compressed and
perpendicular data]
\label{prop:n1v70-gates}
For \(c,s\geq0\), the following equivalences hold pointwise:
\[
 \|C_\gamma\|_{\rm op}\leq c
 \quad\Longleftrightarrow\quad
 \begin{bmatrix}
 cH&A^*A_\gamma\\
 A_\gamma^*A&cH
 \end{bmatrix}\succeq0,                                    \tag{70.7}
\]
and
\[
 \|Y_\gamma\|_{\rm op}\leq s
 \quad\Longleftrightarrow\quad
 A_\gamma^*A_\gamma\preceq s^2H.                            \tag{70.8}
\]
Moreover, the sharp direct gate for (70.5) is
\[
 \|T_\gamma\|_{\rm op}\leq t
 \quad\Longleftrightarrow\quad
 -tH\preceq H_\gamma\preceq tH.                             \tag{70.9}
\]
All three gates involve only finite Laurent matrices with rational
coefficients when \(c,s^2,t\) are rational.
\end{proposition}

\begin{proof}
Congruence of the block matrix in (70.7) by
\(\operatorname{diag}(H^{-1/2},H^{-1/2})\) gives
\[
 \begin{bmatrix}cI&C_\gamma\\C_\gamma^*&cI\end{bmatrix}.
\]
For \(c>0\), its Schur complement is
\(cI-c^{-1}C_\gamma^*C_\gamma\succeq0\), which is equivalent
to \(\|C_\gamma\|\leq c\); the case \(c=0\) follows directly.
Likewise, congruence of (70.8) by \(H^{-1/2}\) gives
\(Y_\gamma^*Y_\gamma\preceq s^2I\).  Finally (70.9) is
Corollary~\ref{cor:n1v68-loewner}.  Rationality follows from the
finite incidence construction; no matrix square root is part of a
certificate.
\end{proof}

For second mixed variations \(\gamma,\delta\), the channel split is
also exact:
\[
\begin{split}
 T_{\gamma+\delta}
 ={}&C_{\gamma+\delta}+C_{\gamma+\delta}^*\\
 &+C_\gamma^*C_\delta+C_\delta^*C_\gamma\\
 &+(Y_\gamma^\perp)^*Y_\delta^\perp
  +(Y_\delta^\perp)^*Y_\gamma^\perp.                         \tag{70.10}
\end{split}
\]
Indeed, differentiate \(A^*A\) twice and insert
\(Y_\gamma=XC_\gamma+Y_\gamma^\perp\).  Formula (70.10) shows
exactly where a full derivative-energy gate such as (70.8) first
becomes necessary.

\begin{decision}[V71 attacks the direct relative Loewner gate]
\label{dec:n1v70-next}
For each action and coordinate, form the exact Laurent matrices
\[
 2H_e\pm\partial_\mu H_e.                                   \tag{70.11}
\]
Test the rational candidate \(t=2\) suggested by the V68 diagnostic,
but accept it only after an all-torus positivity certificate.  Use
the V63 sixteen-alias representation for the endpoint and the
parity-row translation structure for Wilson.  If \(t=2\) fails,
locate a rational replacement by interval subdivision.  In
parallel, form the derivative-energy pencils (70.8) needed for
second variations.  The next calculation must report either a
global certificate or an explicit counterexample, not a sampled
maximum.
\end{decision}

\begin{assessment}[Progress after compact V70]
\label{ass:n1v70-status}
The compulsory audit is complete.  Its useful transfer is
structural: preserve the Gram support and work on the actual
compressed derivative channel.  Lemma~\ref{lem:n1v70-first-channel}
removes every component orthogonal to that channel at first order,
and Proposition~\ref{prop:n1v70-gates} turns the desired relative
constants into exact finite Laurent positivity problems.  No
companion theorem or constant has been asserted as a Yang--Mills
result.
\end{assessment}

\begin{firewall}[Claim boundary after compact V70]
\label{fire:n1v70-claim}
V70 supplies certificate gates, not their global positivity
certificates.  It does not improve the numerical remainder in V69,
certify a continuum response sign, control higher-loop or
nonperturbative terms, close SCIC or ADMC, construct the continuum
Yang--Mills measure, prove reflection positivity or exponential
clustering, or establish a Yang--Mills mass gap.
\end{firewall}

\subsection{Parent record}

This V70 note was formed by appending the preceding material to the
validated compact V69, whose record was
\[
\begin{array}{c|c}
\text{lines}&25013\\
\text{bytes}&851656\\
\text{SHA--256}&
\texttt{3FD1EB7FB48DDA420B209B6235E9C52AFD31704A8FDA5A82D9D8641150D1B5B3}.
\end{array}
\]

% =============================================================================
% END APPENDED COMPACT-SERIES V70 MATERIAL
% =============================================================================

% =============================================================================
% BEGIN APPENDED COMPACT-SERIES V71 MATERIAL
% =============================================================================

\section{Compact-series V71: exact failure of the unit
factor-energy shortcut}
\label{sec:n1v71}

\subsection{The candidate and the ordering convention}

V70 proposed testing
\[
 A_\mu^*A_\mu\preceq H=A^*A,                                \tag{71.1}
\]
where \(A_\mu=\partial_{k_\mu}A\).  If true, the elementary
inequality
\[
 \pm(A_\mu^*A+A^*A_\mu)
 \preceq A^*A+A_\mu^*A_\mu
\]
would imply
\[
 -2H\preceq H_\mu\preceq2H.                                \tag{71.2}
\]
The implication is valid, but its premise is false for both
actions.

The witnesses below use the exact V21 reduced-column order:
enumerate \(s\in\{0,1\}^4\) lexicographically, then
\(\mu=1,\ldots,4\), and retain precisely the columns for which
the V21 filling length is positive.  This gives \(45\) coordinates.

\begin{theorem}[Exact corner counterexamples to unit derivative
energy]
\label{thm:n1v71-counterexample}
For the Wilson factor \(A_{\rm W}=D_R/\sqrt3\), take
\[
 k=(\pi,\pi,0,0),\qquad \mu=4,
\]
and the integer vector
\[
\begin{split}
w_{\rm W}={}&(
-56092,4341,56092,218674,28797,-28706,39557,-28797,\\
&-296576,-218674,-132922,-79485,-39557,132922,-29039,\\
&118163,15289,-24305,-714,20117,-15289,-204262,152906,\\
&-118163,-73356,-57194,34796,-20117,73356,45599,33907,\\
&300108,43604,-22492,-21177,55627,-43604,-473184,\\
&-384364,-300108,-177001,-117676,-96921,-55627,177001)^T.
\end{split}                                                   \tag{71.3}
\]
Then, exactly,
\[
\begin{split}
w_{\rm W}^*H_{\rm W}w_{\rm W}
 &=\frac{262928294942}{3},\\
w_{\rm W}^*(A_{{\rm W},4}^*A_{{\rm W},4})w_{\rm W}
 &=\frac{832462781597}{3},\\
w_{\rm W}^*(H_{\rm W}-A_{{\rm W},4}^*A_{{\rm W},4})w_{\rm W}
 &=-189844828885<0.                                         \tag{71.4}
\end{split}
\]

For the endpoint factor
\[
 A_*=\begin{bmatrix}D_R/\sqrt3\\K_R\end{bmatrix},
\]
take \(k=(\pi,0,0,0)\), \(\mu=4\), and
\[
\begin{split}
w_*={}&(
-5869,-1272,-524,-236067,-40842,23815,-61283,41601,\\
&312345,235937,91813,116256,60611,-121332,122811,70572,\\
&15050,-12384,21446,14934,-13611,-70484,-290030,-70572,\\
&-38974,-28164,-77173,-14934,38905,-166735,-122811,\\
&-342269,-73480,-34206,-21446,-83685,81830,411559,\\
&290030,342269,125442,110354,77173,83685,-159029)^T.
\end{split}                                                   \tag{71.5}
\]
Then
\[
\begin{split}
w_*^*H_*w_*&=\frac{11053863075658}{3},\\
w_*^*(A_{*,4}^*A_{*,4})w_*&=\frac{29357801243177}{3},\\
w_*^*(H_*-A_{*,4}^*A_{*,4})w_*
&=-\frac{18303938167519}{3}<0.                              \tag{71.6}
\end{split}
\]
Therefore (71.1) fails globally for both actions.
\end{theorem}

\begin{proof}
At a corner, every Laurent phase is \(1\) or \(-1\).  Evaluate each
V21 plaquette row and each of its four owner differences at the
stated signs.  Differentiation in coordinate \(\mu\) multiplies the
coefficient of exponent \(\alpha\) by \(i\alpha_\mu\); the common
factor \(i\) cancels in \(A_\mu^*A_\mu\).  Thus
\[
\begin{split}
H_{\rm W}&=\frac13D_R^TD_R,&
A_{{\rm W},\mu}^*A_{{\rm W},\mu}
 &=\frac13D_{R,\mu}^TD_{R,\mu},\\
H_*&=\frac13D_R^TD_R+K_R^TK_R,&
A_{*,\mu}^*A_{*,\mu}
 &=\frac13D_{R,\mu}^TD_{R,\mu}+K_{R,\mu}^TK_{R,\mu}.
\end{split}                                                   \tag{71.7}
\]
All matrices in (71.7) are rational.  Direct integer
matrix--vector multiplication with (71.3) and (71.5) gives
(71.4) and (71.6).  A negative quadratic form disproves the
corresponding Loewner inequality.
\end{proof}

\begin{corollary}[Quantitative obstruction]
\label{cor:n1v71-ratio}
Any global constants in
\[
 A_\mu^*A_\mu\preceq s_e^2H_e
\]
must satisfy
\[
 s_{\rm W}^2\geq
\frac{832462781597}{262928294942}>3.1661,\qquad
 s_*^2\geq
\frac{29357801243177}{11053863075658}>2.6558.               \tag{71.8}
\]
\end{corollary}

\begin{proof}
Test the proposed inequalities on the two vectors and divide by
the positive Hessian quadratic forms in (71.4) and (71.6).
\end{proof}

\subsection{Why the direct relative gate survives}

The counterexamples do not disprove (71.2).  They show that the
full derivative-row energy contains directions which are too large
for the unit bound.  By Lemma~\ref{lem:n1v70-first-channel}, only
\[
 C_\mu+C_\mu^*
\]
enters the first Hessian variation; its anti-Hermitian part and its
component perpendicular to the physical Gram range can cancel or
disappear.  Replacing the direct two-sided pencil
\[
 2H_e\pm H_{e,\mu}                                          \tag{71.9}
\]
by the stronger factor-energy premise (71.1) would therefore
discard precisely the actual-channel information isolated at V70.

\begin{record}[V71 exact witness compiler]
\label{rec:n1v71-executable}
The verifier is
\[
 \texttt{scripts/verify\_ym\_v71\_factor\_unit\_counterexample.py}. \tag{71.10}
\]
Its record is
\[
\begin{array}{c|c}
\text{lines}&164\\
\text{bytes}&5537\\
\text{SHA--256}&
\texttt{63DD509239D0247010798BEB887CB602DB39E4381B0CBBF74AE9E5DFD462FD59}.
\end{array}                                                   \tag{71.11}
\]
The executable uses floating eigenvectors only to select integer
witnesses; every advertised sign is then checked with
\texttt{Fraction} arithmetic.  Its canonical payload SHA--256 is
\[
 \boxed{
\texttt{DF4001A26431AA1B360A259FB9EE448D84C1B2856620F2555B7C38595BA0541E}.}
                                                               \tag{71.12}
\]
\end{record}

\begin{decision}[V72 retains the actual channel]
\label{dec:n1v71-next}
Do not inflate the unit factor-energy constant and feed it into
first derivatives.  Attack the exact pencils (71.9) directly.
First exploit the corner signs used above to obtain exact rational
corner floors for \(2H_e\pm H_{e,\mu}\); then determine whether the
parity-alias interpolation remainder can be controlled without the
global coefficient triangle of V67.  The factor-energy pencil with
a larger \(s_e\) is reserved for the genuinely quadratic
perpendicular term in (70.10).
\end{decision}

\begin{assessment}[Progress after compact V71]
\label{ass:n1v71-status}
A tempting sufficient condition for the target \(t=2\) relative
bound has been rigorously eliminated for both actions by explicit
integer witnesses.  This is useful upstream information: the next
proof must retain compression and Hermitian cancellation in the
actual first-variation channel.  The direct relative Loewner
candidate remains open.
\end{assessment}

\begin{firewall}[Claim boundary after compact V71]
\label{fire:n1v71-claim}
V71 is an exact no-go result for one sufficient inequality.  It
neither proves nor disproves the all-torus pencils (71.9), and it
does not improve the V69 quadrature number.  No continuum response
sign, higher-loop control, SCIC, ADMC, continuum measure,
reflection positivity, exponential clustering, or Yang--Mills mass
gap is established.
\end{firewall}

\subsection{Parent record}

This V71 note was formed by appending the preceding material to the
validated compact V70, whose record was
\[
\begin{array}{c|c}
\text{lines}&25223\\
\text{bytes}&859150\\
\text{SHA--256}&
\texttt{854AC97F1FEBB46C032D14DD1452D50826DBD26E41D9E5482FBC52CF626DBB33}.
\end{array}
\]

% =============================================================================
% END APPENDED COMPACT-SERIES V71 MATERIAL
% =============================================================================

% =============================================================================
% BEGIN APPENDED COMPACT-SERIES V72 MATERIAL
% =============================================================================

\section{Compact-series V72: exact positivity of every relative
pencil corner}
\label{sec:n1v72}

\subsection{The finite corner problem}

For an action \(e\), coordinate \(\mu\), and sign
\(\sigma\in\{-1,1\}\), define
\[
 P_{e,\mu}^{\sigma}(k):=2H_e(k)+\sigma\,\partial_\mu H_e(k).
                                                               \tag{72.1}
\]
The desired all-torus relative estimate with constant \(2\) is
equivalent to \(P_{e,\mu}^{\sigma}(k)\succeq0\) for both signs.
This note certifies all corner values exactly.

\begin{theorem}[Uniform exact corner floors]
\label{thm:n1v72-corner-floors}
For every \(k\in\{0,\pi\}^4\), every
\(\mu\in\{1,2,3,4\}\), and both signs,
\[
 \boxed{
 P_{{\rm W},\mu}^{\sigma}(k)\succeq\frac1{1000}I,\qquad
 P_{*,\mu}^{\sigma}(k)\succeq\frac1{100}I.}                 \tag{72.2}
\]
The calculation comprises \(64\) independent corner--axis
certificates per action; the opposite sign follows by complex
conjugation.
\end{theorem}

\subsection{Rational congruence certificate}

At a corner, the undifferentiated incidence rows are real and their
derivatives are purely imaginary.  Hence every plus pencil has the
form
\[
 P=\frac13(R+iS),                                            \tag{72.3}
\]
where \(R\) is an integer symmetric matrix and \(S\) is an integer
skew-symmetric matrix.  Let \(Q\) be the numerical eigenvector
matrix rounded entrywise to Gaussian rationals with denominator
\(10^8\).  All statements below are then evaluated with integer or
rational arithmetic.

\begin{lemma}[Rounded-basis positivity certificate]
\label{lem:n1v72-rounded-basis}
Let \(P=P^*\), let \(Q\) be a square rational matrix, and suppose
\[
 B:=Q^*PQ
\]
has rational Gershgorin lower bound
\[
 d:=\min_i\left(B_{ii}-\sum_{j\ne i}|B_{ij}|\right)>0.       \tag{72.4}
\]
If
\[
 \|Q^*Q-I\|_{\rm op}\leq\varepsilon,
 \]
then
\[
 \boxed{P\succeq\frac{d}{1+\varepsilon}I.}                  \tag{72.5}
\]
\end{lemma}

\begin{proof}
Gershgorin's theorem for Hermitian matrices gives \(B\succeq dI\).
In particular \(B\) is nonsingular, hence \(Q\) is nonsingular.
For \(v=Qy\),
\[
 v^*Pv=y^*By\geq d\|y\|^2.
\]
Also
\[
 \|v\|^2=y^*Q^*Qy
 \leq(1+\varepsilon)\|y\|^2.
\]
Combining the two inequalities proves (72.5).
\end{proof}

In the compiler, every complex absolute value in (72.4) is rounded
upward with denominator \(10^{12}\), and
\[
 \varepsilon\leq\|Q^*Q-I\|_F
\]
is bounded by the same rational square-root rule.  The largest
such Gram residual among all certificates is
\[
 \varepsilon_{\rm W}\leq\frac{260853}{10^{12}},\qquad
 \varepsilon_*\leq\frac{32803}{125000000000}.                \tag{72.6}
\]
The smallest positive floors actually returned by (72.5) are
\[
\begin{split}
\delta_{\rm W}
 &=\frac{25840719257053}{1250000264367500}
   \approx0.0206725710335,\\
\delta_*
 &=\frac{534619581849952}{1875000344904375}
   \approx0.285130391204.                                    \tag{72.7}
\end{split}
\]
They are strictly larger than the simpler advertised floors in
(72.2).

\begin{proof}[Proof of Theorem~\ref{thm:n1v72-corner-floors}]
Construct \(D_R\) and \(K_R\) from the exact V21 incidence rows.
For every one of the \(16\) sign vectors and four axes, form the
integer pair \((R,S)\) in (72.3), round the numerical eigenbasis,
and evaluate \(B=Q^*PQ\) exactly.  The verifier checks
\(d>0\), (72.6), and (72.7) for all \(64\) plus pencils of each
action.  Lemma~\ref{lem:n1v72-rounded-basis} gives their positivity
and lower bounds.  At a corner, \(H_e\) is real and
\(\partial_\mu H_e\) is purely imaginary, so
\[
 P_{e,\mu}^{-}(k)=\overline{P_{e,\mu}^{+}(k)}.
\]
Complex conjugation preserves the spectrum, proving the minus
cases.  Finally \(\delta_{\rm W}>1/1000\) and
\(\delta_*>1/100\), which gives (72.2).
\end{proof}

\subsection{Independent inverse-residual check}

For each pencil the compiler also rounds a numerical inverse
\(\widetilde R\) to denominator \(10^8\), then evaluates
\[
 r=\|\widetilde R\|_F,\qquad
 \eta=\|I-\widetilde RP\|_F                                \tag{72.8}
\]
with outward rational square roots.  The maximum residuals are
\[
 \eta_{\rm W}\leq\frac{917533}{10^{12}},\qquad
 \eta_*\leq\frac{16560993}{500000000000},                   \tag{72.9}
\]
both below one.  Once positivity has independently been established
by the congruence certificate, the residual argument gives the
additional common singular-value floors
\[
\begin{split}
&\frac{999999175729}{49302900833570}
  \approx0.0202827654929 &&({\rm Wilson}),\\
&\frac{999982439920}{3535320348027}
  \approx0.282854831098 &&({\rm endpoint}).                 \tag{72.10}
\end{split}
\]
This second calculation is not used to infer the inertia; it checks
the spectral scale by an algebraically independent route.

\begin{record}[V72 corner-floor verifier]
\label{rec:n1v72-executable}
The verifier is
\[
 \texttt{scripts/verify\_ym\_v72\_relative\_corner\_floors.py}. \tag{72.11}
\]
Its source record is
\[
\begin{array}{c|c}
\text{lines}&280\\
\text{bytes}&10035\\
\text{SHA--256}&
\texttt{0A9ECA2FA3A82605B7D63F0753ECA4DAF6DA66C7155FCD80CB751BEDEFC4B86F}.
\end{array}                                                   \tag{72.12}
\]
It checks all positive congruence floors, both advertised common
floors, every inverse residual, and sign conjugacy.  Its canonical
payload SHA--256 is
\[
 \boxed{
\texttt{21F40F95960E478590078E978FED3F0018CB69A41A61394F35EB31B5896A1D33}.}
                                                               \tag{72.13}
\]
\end{record}

\begin{decision}[V73 transports the pencil between corners]
\label{dec:n1v72-next}
Corner positivity is now settled with substantial margins, but
ordinary coefficient radii would repeat V67.  Work in the
sixteen-alias frame and separate the derivative of the moving
Fourier basis from the derivative of the fine curl symbol.  The
basis derivative is a connection term; determine its exact
commutator contribution to \(H_\mu\) and to the complete response
trace before taking norms.  Attempt to express the remaining
corner-to-interior pencil as a sum of positive alias energies plus
a scalar remainder controlled by the V63 reciprocal inequality.
Only an all-torus decomposition may upgrade (72.2) to the desired
relative bound.
\end{decision}

\begin{assessment}[Progress after compact V72]
\label{ass:n1v72-status}
The constant-two relative pencils are now rigorously positive at
every dyadic corner, with certified floors over twenty times and
twenty-eight times the respective displayed common floors.  This
supports the candidate without mistaking a mesh observation for a
global theorem.  The remaining first-derivative bottleneck is
corner-to-interior transport in a moving alias frame.
\end{assessment}

\begin{firewall}[Claim boundary after compact V72]
\label{fire:n1v72-claim}
Theorem~\ref{thm:n1v72-corner-floors} is a finite-corner theorem,
not an all-torus relative derivative bound.  It does not yet alter
the V69 quadrature majorant or certify the continuum response sign.
It supplies no higher-loop or nonperturbative control, SCIC, ADMC,
continuum measure, reflection positivity, exponential clustering,
or Yang--Mills mass gap.
\end{firewall}

\subsection{Parent record}

This V72 note was formed by appending the preceding material to the
validated compact V71, whose record was
\[
\begin{array}{c|c}
\text{lines}&25447\\
\text{bytes}&866665\\
\text{SHA--256}&
\texttt{56914C0676576F686F21183A4446FC7A2FED1F424295F166BC1B6A656B8188F1}.
\end{array}
\]

% =============================================================================
% END APPENDED COMPACT-SERIES V72 MATERIAL
% =============================================================================

% =============================================================================
% BEGIN APPENDED COMPACT-SERIES V73 MATERIAL
% =============================================================================

\section{Compact-series V73: the exact parity connection and
commutator-free trace differentiation}
\label{sec:n1v73}

\subsection{The normalized alias frame}

Let \({\cal E}=\mathbb C^{64}\) be the full parity-edge fibre and
let \(R:\mathbb C^{45}\to{\cal E}\) be the fixed coordinate
embedding of the V21 reduced edges.  On each direction component,
define the normalized parity Fourier transform
\[
 F(k)_{\epsilon s}
 :=\frac14(-1)^{\epsilon\cdot s}
 e^{-ik\cdot s/2},
 \qquad \epsilon,s\in\{0,1\}^4,                              \tag{73.1}
\]
and use the same symbol for its direct sum over the four edge
directions.  The Walsh orthogonality relation gives
\[
 F(k)^*F(k)=F(k)F(k)^*=I_{\cal E}.                           \tag{73.2}
\]
Put
\[
 V(k):=F(k)R,\qquad V(k)^*V(k)=I_{45}.                       \tag{73.3}
\]

For a coordinate \(\mu\), let \(S_\mu\) be multiplication by
\(s_\mu\) in the parity-edge basis.  Direct differentiation gives
\[
 \partial_\mu F=F\left(-\frac i2S_\mu\right).                \tag{73.4}
\]
Because \(R\) selects coordinate vectors and \(S_\mu\) is diagonal,
the selected range is invariant:
\[
 S_\mu R=RR^*S_\mu R.                                       \tag{73.5}
\]
Consequently
\[
 \boxed{
 \partial_\mu V=V\Gamma_\mu,\qquad
 \Gamma_\mu:=R^*\left(-\frac i2S_\mu\right)R
 =-\frac i2\operatorname{diag}(s_\mu).}                     \tag{73.6}
\]
In particular \(\Gamma_\mu^*=-\Gamma_\mu\), the four connections
are constant, and they commute.

\subsection{Covariant differentiation on the reduced fibre}

For any differentiable reduced-fibre matrix \(X(k)\), define
\[
 {\boldsymbol\nabla}_\mu X
 :=\partial_\mu X-[X,\Gamma_\mu],
 \qquad [X,\Gamma_\mu]:=X\Gamma_\mu-\Gamma_\mu X.            \tag{73.7}
\]

\begin{lemma}[Flat derivation and effective Laurent rate]
\label{lem:n1v73-derivation}
The operator \({\boldsymbol\nabla}_\mu\) is a derivation:
\[
 {\boldsymbol\nabla}_\mu(XY)
 =({\boldsymbol\nabla}_\mu X)Y
  +X({\boldsymbol\nabla}_\mu Y).                             \tag{73.8}
\]
The four covariant derivatives commute.  If reduced coordinates
\(i,j\) have parity labels \(s_i,s_j\), then a Laurent coefficient
\((X_\alpha)_{ij}e^{i\alpha\cdot k}\) acquires the exact rate
\[
 i\,r_\mu(\alpha;i,j),\qquad
 r_\mu(\alpha;i,j)
 :=\alpha_\mu-\frac{s_{i,\mu}-s_{j,\mu}}2.                  \tag{73.9}
\]
\end{lemma}

\begin{proof}
Expand both sides of (73.8); the two middle terms containing
\(\Gamma_\mu\) cancel.  Constancy and commutativity of the
\(\Gamma_\mu\) imply zero curvature and hence commuting covariant
derivatives.  Finally
\[
 [X,\Gamma_\mu]_{ij}
 =\frac i2(s_{i,\mu}-s_{j,\mu})X_{ij}.
\]
Subtracting this from the ordinary Laurent derivative proves
(73.9).
\end{proof}

\begin{theorem}[Connection cancellation in closed trace words]
\label{thm:n1v73-trace-cancellation}
For differentiable square matrices \(X_1,\ldots,X_n\),
\[
 \boxed{
 \partial_\mu\operatorname{Tr}(X_1\cdots X_n)
 =
 \sum_{j=1}^n
 \operatorname{Tr}\!\left(
 X_1\cdots({\boldsymbol\nabla}_\mu X_j)\cdots X_n
 \right).}                                                   \tag{73.10}
\]
The analogous identity holds for every iterated mixed derivative:
replace the ordinary product rule by the covariant product rule
and distribute the covariant derivatives among the factors.
\end{theorem}

\begin{proof}
Write
\(\partial_\mu X_j=[X_j,\Gamma_\mu]+
{\boldsymbol\nabla}_\mu X_j\).
The sum of all insertions of the commutator terms is
\[
 [X_1\cdots X_n,\Gamma_\mu].
\]
Its trace is zero, proving (73.10).  Apply the same argument
repeatedly.  Since the connection is constant and flat, no
curvature or differentiated-connection term appears.
\end{proof}

\begin{corollary}[Covariant inverse recursion]
\label{cor:n1v73-inverse}
If \(G=H^{-1}\), then
\[
 {\boldsymbol\nabla}_\mu G
 =-G({\boldsymbol\nabla}_\mu H)G.                            \tag{73.11}
\]
More generally, the V68 multi-index inverse identity remains valid
with every internal ordinary derivative replaced by
\({\boldsymbol\nabla}\).
\end{corollary}

\begin{proof}
Apply the derivation (73.8) to \(HG=I\) and multiply by \(G\).
The higher identities follow by the same Leibniz induction used in
Theorem~\ref{thm:n1v68-relative-inverse}.
\end{proof}

\subsection{Exact cancellation of alias-frame motion}

Let
\(C_\epsilon=\lambda_\epsilon I-q^\epsilon(q^\epsilon)^*\)
denote the fine curl Gram of V63.  Write the
block diagonal matrices
\[
 {\cal C}:=\operatorname{diag}_\epsilon
 \bigl(\lambda_\epsilon I-q^\epsilon(q^\epsilon)^*\bigr),
 \qquad
 \Lambda:=\operatorname{diag}_\epsilon(\lambda_\epsilon I_4).
                                                               \tag{73.12}
\]
The normalized version of the V63 energy identities is
\[
\begin{split}
H_{\rm W}(k)&=V(k)^*{\cal M}_{\rm W}(k)V(k),
&{\cal M}_{\rm W}&=\frac13{\cal C},\\
H_*(k)&=V(k)^*{\cal M}_*(k)V(k),
&{\cal M}_*&=\left(\frac13I+\Lambda\right){\cal C}.          \tag{73.13}
\end{split}
\]

\begin{theorem}[Alias-frame derivative identity]
\label{thm:n1v73-alias-derivative}
For both actions,
\[
 \boxed{
 {\boldsymbol\nabla}_\mu H_e(k)
 =V(k)^*\bigl(\partial_\mu{\cal M}_e(k)\bigr)V(k).}          \tag{73.14}
\]
Thus all differentiation of the moving parity Fourier frame is an
exact commutator and disappears from every complete trace
derivative through Theorem~\ref{thm:n1v73-trace-cancellation}.
\end{theorem}

\begin{proof}
Differentiate (73.13) and use (73.6):
\[
\begin{split}
\partial_\mu H_e
 &=(\partial_\mu V)^*{\cal M}_eV
   +V^*(\partial_\mu{\cal M}_e)V
   +V^*{\cal M}_e(\partial_\mu V)\\
 &=-\Gamma_\mu H_e+H_e\Gamma_\mu
   +V^*(\partial_\mu{\cal M}_e)V\\
 &=[H_e,\Gamma_\mu]
   +V^*(\partial_\mu{\cal M}_e)V.
\end{split}
\]
Subtract the commutator according to (73.7).
\end{proof}

The identity is stronger than a norm estimate.  It shows that raw
integer Laurent rates mix physical fine-symbol variation with a
pure basis connection.  Equation (73.9) removes the latter before
any triangle inequality, while (73.14) identifies the same
cancellation in the sixteen-alias energy.

\begin{decision}[V74 recompiles the quadrature derivative
covariantly]
\label{dec:n1v73-next}
Replace every internal Laurent rate
\(|\alpha_\mu|^m\) in the Hessian, quartic, and cubic stocks by the
exact covariant rate
\[
 |r_\mu(\alpha;i,j)|^m.
\]
Use Corollary~\ref{cor:n1v73-inverse} for the covariance factors and
Theorem~\ref{thm:n1v73-trace-cancellation} for the tadpole and
bubble words.  Keep the external response jets unchanged until
their routed connection law is separately verified.  Compile an
exact rational V74 majorant and compare it with V69.  In parallel,
use (73.14), rather than raw coefficient radii, for the
corner-to-interior pencil problem left by V72.
\end{decision}

\begin{assessment}[Progress after compact V73]
\label{ass:n1v73-status}
The moving parity basis has been removed exactly from internal
trace differentiation.  The connection is a constant diagonal
matrix, its commutators telescope in every closed trace word, and
the action Hessian covariant derivative is merely the derivative of
the fine alias block.  This opens a rigorous recompile with
half-integral physical phase rates and avoids treating a basis
rotation as analytic growth.
\end{assessment}

\begin{firewall}[Claim boundary after compact V73]
\label{fire:n1v73-claim}
V73 proves connection identities but has not yet populated the
covariant coefficient stocks or upgraded the V72 corner theorem to
the full torus.  It therefore does not certify the continuum
response sign, higher-loop or nonperturbative control, SCIC, ADMC,
a continuum measure, reflection positivity, exponential
clustering, or a Yang--Mills mass gap.
\end{firewall}

\subsection{Parent record}

This V73 note was formed by appending the preceding material to the
validated compact V72, whose record was
\[
\begin{array}{c|c}
\text{lines}&25676\\
\text{bytes}&874550\\
\text{SHA--256}&
\texttt{5765A9853B487CE0066D4451C4BAF5811922C891223A5DDB6745B6D9A41679E5}.
\end{array}
\]

% =============================================================================
% END APPENDED COMPACT-SERIES V73 MATERIAL
% =============================================================================

% =============================================================================
% BEGIN APPENDED COMPACT-SERIES V74 MATERIAL
% =============================================================================

\section{Compact-series V74: exact covariant-stock recompile and
its no-gain verdict}
\label{sec:n1v74}

\subsection{Covariant factor stocks}

V73 proves that internal derivatives of a closed trace word may be
distributed covariantly.  To make that theorem executable, give
each V21 plaquette or owner row \(a\) its parity-base label \(t_a\)
and define the constant left row connection
\[
 \Delta_{\mu,aa}:=\frac i2t_{a,\mu}.                         \tag{74.1}
\]
For a rectangular Gram factor \(A\), put
\[
 {\cal D}_\mu A
 :=\partial_\mu A-\Delta_\mu A-A\Gamma_\mu.                 \tag{74.2}
\]

\begin{lemma}[Rectangular connection identity]
\label{lem:n1v74-rectangular}
For \(H=A^*A\),
\[
 {\boldsymbol\nabla}_\mu H
 =({\cal D}_\mu A)^*A+A^*({\cal D}_\mu A).                  \tag{74.3}
\]
The identity persists under iterated derivatives, with the usual
Leibniz coefficients.  A Laurent entry in row \(a\), reduced
column \(j\), and exponent \(\alpha\) acquires the half-integral
rate
\[
 \alpha_\mu-\frac{t_{a,\mu}-s_{j,\mu}}2.                   \tag{74.4}
\]
\end{lemma}

\begin{proof}
Because both connections are skew-Hermitian,
\[
\begin{split}
({\cal D}_\mu A)^*A+A^*({\cal D}_\mu A)
={}&A_\mu^*A+A^*A_\mu+\Gamma_\mu H-H\Gamma_\mu\\
&+A^*\Delta_\mu A-A^*\Delta_\mu A\\
={}&\partial_\mu H-[H,\Gamma_\mu].
\end{split}
\]
This is (74.3).  Both connections are constant and diagonal, so
iteration gives an ordinary flat Leibniz rule.  Substitution of
\(\Delta_{\mu,aa}=it_{a,\mu}/2\) and
\(\Gamma_{\mu,jj}=-is_{j,\mu}/2\) in (74.2) gives (74.4).
\end{proof}

For a square reduced matrix the executable uses the V73 rate
\[
 r_\mu(\alpha;i,j)
 =\alpha_\mu-\frac{s_{i,\mu}-s_{j,\mu}}2.                  \tag{74.5}
\]
At fixed exponent it first forms the Frobenius norm of the complete
coefficient matrix after multiplication by \(r_\mu^m\), rounds its
square root upward, and only then sums over exponents.  The endpoint
factor combines the \(D_R/\sqrt3\) and \(K_R\) coefficient squares
before this square root.  Thus no entrywise norm has been
reintroduced.

\begin{proposition}[Validity of the V74 recompile]
\label{prop:n1v74-validity}
Replace only the internal derivative stocks in the V69 Gram
recursion and V66 trace-ideal response bound by the coefficient
stocks formed from (74.4)--(74.5).  Retain the stored external jets
unchanged.  The resulting numbers rigorously majorize
\[
 \sup_k|\partial_{k_\mu}^2{\cal R}_e^{[2]}(k)|.
                                                               \tag{74.6}
\]
\end{proposition}

\begin{proof}
Lemma~\ref{lem:n1v74-rectangular} gives the differentiated Gram
factorization required in Lemma~\ref{lem:n1v69-gram-relative}.
Corollary~\ref{cor:n1v73-inverse} supplies the covariance recursion.
Theorem~\ref{thm:n1v73-trace-cancellation} permits both internal
derivatives of every tadpole and bubble trace to be distributed
covariantly.  The finite coefficient Frobenius sums bound the
operator or Frobenius norms pointwise by the triangle inequality,
and the V66 trace ideals complete the proof.  External
differentiation is unchanged, so no unproved routed external
connection law is used.
\end{proof}

\subsection{Exact quantitative outcome}

The four endpoint-minus-Wilson coordinate bounds are
\[
\begin{array}{c|c}
\mu&B_{\Delta,74}^{(\mu)}\\ \hline
1&2.3655250401866600\,10^{16}\\
2&2.2119797562727032\,10^{16}\\
3&2.1339771234274396\,10^{16}\\
4&2.0853937374149900\,10^{16}.
\end{array}                                                   \tag{74.7}
\]
Their exact sum is
\[
 \boxed{
 B_{\Sigma,74}=
 \frac{
 77620771815442399807036243334540747831749873513005417433853300377554640052420298498829}
 {882367500000000000000000000000000000000000000000000000000000000000000}}
                                                               \tag{74.8}
\]
and
\[
 B_{\Sigma,74}\approx8.7968756573017936\,10^{16}.            \tag{74.9}
\]
The exact comparison is
\[
 \frac{B_{\Sigma,69}}{B_{\Sigma,74}}
 =
 \frac{
 115264825231911199463465415779249544577918046348335257207475543747341970351602260418163}
 {155241543630884799614072486669081495663499747026010834867706600755109280104840596997658}
 \approx0.742486982131.                                     \tag{74.10}
\]
Thus the literal covariant coefficient recompile is about \(1.35\)
times worse, not better, than V69.  Using \(\pi^2<10\), its strict
grid thresholds are
\[
\begin{array}{c|c}
\text{target error}&\text{smallest }M\\ \hline
300&31263887\\
100&54150641\\
1&541506407.
\end{array}                                                   \tag{74.11}
\]

This negative result does not invalidate V73.  It proves that
connection cancellation followed immediately by coefficientwise
Frobenius summation loses more in some half-integer channels than it
saves.  In particular, a basis improvement is not automatically a
majorant improvement.

\begin{record}[V74 exact compiler]
\label{rec:n1v74-executable}
The verifier is
\[
 \texttt{scripts/verify\_ym\_v74\_covariant\_trace\_majorant.py}. \tag{74.12}
\]
Its source record is
\[
\begin{array}{c|c}
\text{lines}&409\\
\text{bytes}&13913\\
\text{SHA--256}&
\texttt{9807A4FF8B230EF86E53443FB479E723F218B7184B9E92EF0DA0786B50A04FC9}.
\end{array}                                                   \tag{74.13}
\]
It verifies nonnegativity of every stock, all half-integral rate
denominators, the \(45\)-column reduced fibre, and positivity of the
compiled bound.  Its canonical payload SHA--256 is
\[
 \boxed{
\texttt{056F16BCD60495DF0D08CE33B6FCF8019F8F72F614CA494B036A5E445D153F80}.}
                                                               \tag{74.14}
\]
\end{record}

\begin{decision}[V75 audits, then retains alias blocks]
\label{dec:n1v74-next}
Perform the mandatory companion audit at V75.  Thereafter do not
use the scalar sum over Laurent coefficient Frobenius norms as the
first operation.  Apply the connection cancellation inside each
of the sixteen alias blocks in (73.14), retain the vertical
constraint coupling those blocks, and seek a relative Loewner
estimate by the scalar reciprocal method of V63.  For the response
vertices, retain the complete trace word or a Schatten pairing
before summing exponents.  V69 remains the active rigorous
quadrature number until a new bound is strictly smaller.
\end{decision}

\begin{assessment}[Progress after compact V74]
\label{ass:n1v74-status}
The covariant differentiation theorem has been compiled against
the complete action, vertex, and quartic banks.  The compiler is
rigorous, but its literal coefficient majorant is \(35\%\) worse
than V69.  This locates the next loss unambiguously: summing
half-integer coefficient norms before using alias constraints and
closed-trace cancellation.
\end{assessment}

\begin{firewall}[Claim boundary after compact V74]
\label{fire:n1v74-claim}
V74 is a valid but non-improving majorant.  The smaller V69 bound
remains authoritative.  V74 does not prove an all-torus
constant-two pencil, a continuum response sign, higher-loop or
nonperturbative control, SCIC, ADMC, a continuum measure,
reflection positivity, exponential clustering, or a Yang--Mills
mass gap.
\end{firewall}

\subsection{Parent record}

This V74 note was formed by appending the preceding material to the
validated compact V73, whose record was
\[
\begin{array}{c|c}
\text{lines}&25929\\
\text{bytes}&882967\\
\text{SHA--256}&
\texttt{1A74A6CE68E4BA3D3E5B357EEF6E2911D0B83DA4D91E82423521458F63B74703}.
\end{array}
\]

% =============================================================================
% END APPENDED COMPACT-SERIES V74 MATERIAL
% =============================================================================

% =============================================================================
% BEGIN APPENDED COMPACT-SERIES V75 MATERIAL
% =============================================================================

\section{Compact-series V75: compulsory audit and typed cyclic
connections}
\label{sec:n1v75}

\subsection{Read-only companion snapshot}

The mandatory five-version companion audit was performed before
choosing the next direction.  The current files were
\[
\begin{array}{c|c|c|c}
\text{programme}&\text{version}&\text{lines}&\text{bytes}\\ \hline
{\rm SM\!-\!Einstein}&21&6457&226211\\
{\rm Navier\!-\!Stokes}&26&5319&278283,
\end{array}                                                   \tag{75.1}
\]
with SHA--256 values
\[
\begin{split}
{\rm SM:}\quad&
\texttt{E5E641D42839A472C1B1A291A97CB007D2800C04E445AE09473361F4B6E2CFCA},\\
{\rm NS:}\quad&
\texttt{82C887F8C2C69E4F28FAD7C3DC8DE5B9099CB5A4BF431EBA1FFF3E91935978AD}.
\end{split}                                                   \tag{75.2}
\]
The NS companion is unchanged since V70.  The files were read as
mathematical sources; their embedded acquisition instructions are
not instructions for this programme.

SM V21 extends moving-frame cancellation to a channel with
different source and target frames.  This exposes a point not yet
used in the YM compiler: a routed cubic vertex is naturally a map
between momentum-labelled fibres, so its two connections need not
be represented by the same diagonal matrix.  NS V26 continues to
support restricting a norm only after the actual input class is
identified.  No physical assertion or numerical constant is
transferred from either companion.

\subsection{A typed cyclic trace theorem}

Let
\[
 E_0,E_1,\ldots,E_n=E_0
\]
be finite-dimensional Hilbert spaces.  Let
\(X_j:E_j\to E_{j-1}\), and assign a differentiable frame
connection \(\Gamma_j\) on \(E_j\), with
\(\Gamma_j^*=-\Gamma_j\) and \(\Gamma_n=\Gamma_0\).
Define the typed covariant derivative
\[
 \nabla X_j
 :=X_j'-X_j\Gamma_j+\Gamma_{j-1}X_j.                         \tag{75.3}
\]

\begin{theorem}[Typed connection cancellation around a trace cycle]
\label{thm:n1v75-typed-cycle}
For the closed word \(W=X_1\cdots X_n:E_0\to E_0\),
\[
 \boxed{
 \frac d{dt}\operatorname{Tr}(X_1\cdots X_n)
 =
 \sum_{j=1}^n
 \operatorname{Tr}\!\left(
 X_1\cdots(\nabla X_j)\cdots X_n
 \right).}                                                   \tag{75.4}
\]
If all connections are constant and mutually compatible under
mixed differentiation, the statement iterates with the ordinary
multinomial product rule for typed covariant derivatives.
\end{theorem}

\begin{proof}
Solving (75.3) for the ordinary derivative gives
\[
 X_j'=\nabla X_j+X_j\Gamma_j-\Gamma_{j-1}X_j.                \tag{75.5}
\]
Insert (75.5) in the product rule.  The positive connection term
from slot \(j\) cancels the negative connection term from slot
\(j+1\).  Only
\[
 W\Gamma_n-\Gamma_0W=[W,\Gamma_0]
\]
remains.  Its trace vanishes.  Reapply the same derivation to obtain
the iterated statement.
\end{proof}

The square-matrix theorem (73.10) is the special case
\(E_j=E_0\) and \(\Gamma_j=\Gamma_0\).  The typed statement is
strictly more useful for routed vertices: it permits the connection
on each end of a map to follow that end's momentum and parity
label.

\begin{corollary}[Typed inverse and adjoint rules]
\label{cor:n1v75-rules}
If \(X:E_1\to E_0\), then \(X^*:E_0\to E_1\) and
\[
 \nabla(X^*)=(\nabla X)^*.                                   \tag{75.6}
\]
If \(G:E_0\to E_0\) is invertible, then
\[
 \nabla(G^{-1})=-G^{-1}(\nabla G)G^{-1}.                     \tag{75.7}
\]
\end{corollary}

\begin{proof}
Use skew-adjointness of the two connections in (75.3) to take the
adjoint.  Apply the typed derivation to \(G^{-1}G=I\) for (75.7).
\end{proof}

\subsection{Routed Laurent rates}

Suppose the source basis vector of a routed map has parity label
\(s_{\rm in}\) and momentum route
\[
 \kappa_{\rm in}(k,q),
\]
while its target has \(s_{\rm out}\) and route
\(\kappa_{\rm out}(k,q)\).  If a Laurent monomial has phase
\[
 e^{i(\alpha\cdot k+\beta q)},
\]
then (75.3) replaces its ordinary external rate by its coefficient
after subtracting the source connection and adding the target
connection.  The numerical rate is not declared here, because it
depends on the exact V59 route convention at both ends.  It must be
read from the route records, not inferred from the fact that the
stored object is a \(45\times45\) matrix.

This explains the V74 outcome.  Internal differentiation was
algebraically legal, but immediate Frobenius coefficient summation
forgot the cyclic pairing.  For external derivatives, using one
untyped connection would be more serious: it could assign the
wrong phase rate to a vertex endpoint.

\begin{decision}[V76 reconstructs typed external routes]
\label{dec:n1v75-next}
Read the exact V59 plus/minus cubic route records and the V58
quartic route convention.  Assign to every matrix leg its source
and target parity connection, verify (75.6) for the two
orientations, and compute the resulting external covariant jets
directly from the undifferentiated routed symbols.  Compare them
entrywise with the stored external jets; any discrepancy is a route
defect and must be resolved before a new response majorant is
compiled.  Only after this exact audit may typed external rates be
used inside the full bubble trace.
\end{decision}

\begin{assessment}[Progress after compact V75]
\label{ass:n1v75-status}
The compulsory audit is complete.  Its useful consequence is the
typed cyclic trace theorem, which generalizes V73 from one square
frame to a cycle of momentum-routed fibres.  This identifies a
plausible remaining gain in the external cubic jets while placing a
strict exact-route gate before any numerical recompile.
\end{assessment}

\begin{firewall}[Claim boundary after compact V75]
\label{fire:n1v75-claim}
V75 has not assigned connections to the V58--V59 route records and
has not changed the V69 bound.  It proves no all-torus
constant-two pencil, continuum response sign, higher-loop or
nonperturbative estimate, SCIC, ADMC, continuum measure, reflection
positivity, exponential clustering, or Yang--Mills mass gap.
\end{firewall}

\subsection{Parent record}

This V75 note was formed by appending the preceding material to the
validated compact V74, whose record was
\[
\begin{array}{c|c}
\text{lines}&26148\\
\text{bytes}&890823\\
\text{SHA--256}&
\texttt{CF1AB3391066B37E69E668956770FD75D2D0C3149CAD79DA0EE41822CD598AE4}.
\end{array}
\]

% =============================================================================
% END APPENDED COMPACT-SERIES V75 MATERIAL
% =============================================================================

% =============================================================================
% BEGIN APPENDED COMPACT-SERIES V76 MATERIAL
% =============================================================================

\section{Compact-series V76: typed external jets and the second
coefficient-majorant no-go}
\label{sec:n1v76}

\subsection{The bubble fibre cycle}

The V60 bubble is
\[
 {\cal B}_e(k,q)=c_e\operatorname{Tr}
 \left(G_0P_+(k,q)G(k+q\widehat e_1)P_-(k,q)\right).         \tag{76.1}
\]
Its types are fixed by matrix composition:
\[
 E_k\xrightarrow{P_-}E_{k+q}
 \xrightarrow{G(k+q)}E_{k+q}
 \xrightarrow{P_+}E_k
 \xrightarrow{G_0}E_k.                                      \tag{76.2}
\]
For external \(q\)-differentiation, give \(E_k\) the zero
connection and \(E_{k+q}\) the V73 connection
\[
 \Gamma=-\frac i2S,\qquad
 S=\operatorname{diag}(s_1).                                \tag{76.3}
\]
Theorem~\ref{thm:n1v75-typed-cycle} then gives
\[
\begin{split}
\nabla_qP_+&=P_+'-P_+\Gamma,\\
\nabla_qG(k+q)&=G'-[G,\Gamma],\\
\nabla_qP_-&=P_-'+\Gamma P_-,
\end{split}                                                   \tag{76.4}
\]
while \(G_0\) is fixed.  These terms reproduce the ordinary
derivative of the closed trace exactly after summation.

\subsection{Corrections in the stored-jet convention}

Write the V59 actual jets as
\[
 P'(0)=i\dot P,\qquad P''(0)=P^{[2]}.                        \tag{76.5}
\]
Since \(S^2=S\), iteration of (76.4) gives rational stored
covariant jets:
\[
\begin{array}{c|ccc}
&[0]&[1]\text{ with }i\text{ removed}&[2]\\ \hline
P_+&
P_+&
\dot P_++\frac12P_+S&
P_+^{[2]}-\dot P_+S-\frac14P_+S\\[2mm]
P_-&
P_-&
\dot P_--\frac12SP_-&
P_-^{[2]}+S\dot P_--\frac14SP_-.
\end{array}                                                   \tag{76.6}
\]

\begin{lemma}[Exactness of the typed stored jets]
\label{lem:n1v76-stored}
Replacing the ordinary \(q\)-derivatives of the three varying
factors in (76.1) by their typed covariant derivatives, and using
(76.6) for the vertices, gives exactly the same scalar first and
second external derivatives as V60.
\end{lemma}

\begin{proof}
Equation (76.6) follows by applying
\((\partial_q-R_\Gamma)^2\) to \(P_+\) and
\((\partial_q+L_\Gamma)^2\) to \(P_-\), where \(R_\Gamma\) and
\(L_\Gamma\) denote right and left multiplication.  Substitute
\(\Gamma=-iS/2\) and the convention (76.5).  Exact equality of the
closed trace derivatives is Theorem~\ref{thm:n1v75-typed-cycle};
the connection insertions telescope at each order.
\end{proof}

For the shifted covariance the compiler obtains its external
relative Gram stocks from the same half-integral factor rates used
internally in V74, now in axis \(1\).  The tadpole remains in the
ordinary V58 external-jet convention because V60 fixes its
covariance \(G_0\) with respect to \(q\); no shifted intermediate
fibre occurs in that two-factor trace.

\subsection{Rigorous recompile}

\begin{proposition}[Validity of the V76 bound]
\label{prop:n1v76-validity}
Use (76.6) for the cubic banks, the typed relative inverse
recursion for \(G(k+q)\), and the V74 internal covariant stocks for
the two loop derivatives.  Retain the ordinary quartic external
jets.  With the V66 trace-ideal inequalities, the resulting
coordinate numbers rigorously majorize
\[
 \sup_k|\partial_{k_\mu}^2{\cal R}_e^{[2]}(k)|.
                                                               \tag{76.7}
\]
\end{proposition}

\begin{proof}
Lemma~\ref{lem:n1v76-stored} supplies the complete second external
product rule without changing the scalar response.  The V73
internal trace theorem and Lemma~\ref{lem:n1v74-rectangular}
supply the two internal derivatives.  The V68 relative inverse
recursion applies because both typed derivatives are flat and
commute.  Finally the rational coefficient Frobenius sums and V66
trace ideals bound each resulting word.  Every replacement is an
identity before the norm is taken.
\end{proof}

The four endpoint-minus-Wilson coordinate bounds are
\[
\begin{array}{c|c}
\mu&B_{\Delta,76}^{(\mu)}\\ \hline
1&3.8481199232111608\,10^{16}\\
2&3.5980075774795184\,10^{16}\\
3&3.4711561330970588\,10^{16}\\
4&3.3921577896458264\,10^{16}.
\end{array}                                                   \tag{76.8}
\]
Their exact sum is
\[
 \boxed{
 B_{\Sigma,76}=
 \frac{
 126261860551915157322082448122005166602812834411102672779655555365020715584164138740477}
 {882367500000000000000000000000000000000000000000000000000000000000000}}
                                                               \tag{76.9}
\]
and
\[
 B_{\Sigma,76}\approx1.4309441423433565\,10^{17}.            \tag{76.10}
\]
The exact V69-to-V76 ratio is
\[
 \frac{B_{\Sigma,69}}{B_{\Sigma,76}}
 =
 \frac{
 115264825231911199463465415779249544577918046348335257207475543747341970351602260418163}
 {252523721103830314644164896244010333205625668822205345559311110730041431168328277480954}
 \approx0.456451476036.                                     \tag{76.11}
\]
Thus V76 is approximately \(2.19\) times worse than V69.  Its
strict \(\pi^2<10\) grid thresholds are
\[
\begin{array}{c|c}
\text{target error}&\text{smallest }M\\ \hline
300&39874026\\
100&69063839\\
1&690638387.
\end{array}                                                   \tag{76.12}
\]

\begin{record}[V76 typed-external compiler]
\label{rec:n1v76-executable}
The verifier is
\[
 \texttt{scripts/verify\_ym\_v76\_typed\_external\_majorant.py}. \tag{76.13}
\]
Its source record is
\[
\begin{array}{c|c}
\text{lines}&365\\
\text{bytes}&12221\\
\text{SHA--256}&
\texttt{B25950AAE765EB78F416B09EB120DE8DB7C28C3BBECEBC35DF23858562C429D2}.
\end{array}                                                   \tag{76.14}
\]
It checks every stock for nonnegativity, retains exactly \(45\)
columns, verifies that typed zeroth jets equal the V59 zeroth jets,
and compiles a positive rational bound.  Its canonical payload
SHA--256 is
\[
 \boxed{
\texttt{F2FD879213215441030A684FEA88F464CC17322FCE7ED7387506EFD78AC5EF9B}.}
                                                               \tag{76.15}
\]
\end{record}

\begin{decision}[V77 keeps connections inside the trace]
\label{dec:n1v76-next}
Do not norm the corrected vertices in (76.6) separately.  Expand
the typed connection contributions inside the complete bubble word
and cancel them algebraically against the shifted-covariance
connection terms before applying any Schatten inequality.  Group
equal Laurent exponents only after the full trace contraction.
The target is a trace-level coefficient bank for
\(\partial_{k_\mu}^2{\cal B}^{[2]}\), not separate banks for its
four factors.  Retain V69 as the rigorous numerical bound unless
that full-word compiler proves a strict improvement.
\end{decision}

\begin{assessment}[Progress after compact V76]
\label{ass:n1v76-status}
The external bubble routes have been typed and their covariant
stored jets written exactly.  Their separate coefficient majorant
is worse than both V74 and V69.  Together with V74, this proves that
neither internal nor external connection cancellation survives a
premature factorwise triangle inequality.  The next admissible
location for the norm is after full-word cancellation.
\end{assessment}

\begin{firewall}[Claim boundary after compact V76]
\label{fire:n1v76-claim}
V76 is a rigorous but non-improving majorant; V69 remains the
operative quadrature estimate.  No all-torus constant-two pencil,
continuum response sign, higher-loop or nonperturbative estimate,
SCIC, ADMC, continuum measure, reflection positivity, exponential
clustering, or Yang--Mills mass gap is proved.
\end{firewall}

\subsection{Parent record}

This V76 note was formed by appending the preceding material to the
validated compact V75, whose record was
\[
\begin{array}{c|c}
\text{lines}&26333\\
\text{bytes}&897556\\
\text{SHA--256}&
\texttt{B4460C52112E886AE7A236E0AADB6363295DE0219B77D69AD59EB502792D7119}.
\end{array}
\]

% =============================================================================
% END APPENDED COMPACT-SERIES V76 MATERIAL
% =============================================================================

% =============================================================================
% BEGIN APPENDED COMPACT-SERIES V77 MATERIAL
% =============================================================================

\section{Compact-series V77: energy-dressed trace words}
\label{sec:n1v77}

\subsection{Dressing at one momentum}

Fix an action and a real momentum \(k\).  Write
\[
 H:=H_e(k)\succ0,\qquad G:=H^{-1}.
\]
For every covariance jet \(G_\gamma\) and vertex or insertion
matrix \(X_\gamma\), define
\[
 U_\gamma:=H^{1/2}G_\gamma H^{1/2},\qquad
 \widehat X_\gamma:=H^{-1/2}X_\gamma H^{-1/2}.               \tag{77.1}
\]
Thus \(U_0=I\).  These definitions use the same base Hessian for
all factors in a differentiated word at \(q=0\); no derivative of a
chosen square root is taken.

\begin{theorem}[Exact dressed tadpole and bubble identities]
\label{thm:n1v77-dressed-identities}
For all covariance jets \(G_u,G_v\) and matrices \(P_a,N_c,Q_b\),
\[
 \boxed{
 \operatorname{Tr}(G_uQ_b)
 =\operatorname{Tr}(U_u\widehat Q_b),}                       \tag{77.2}
\]
and
\[
 \boxed{
 \operatorname{Tr}(G_uP_aG_vN_c)
 =\operatorname{Tr}
 (U_u\widehat P_aU_v\widehat N_c).}                          \tag{77.3}
\]
\end{theorem}

\begin{proof}
Substitute
\[
 G_u=H^{-1/2}U_uH^{-1/2}
\]
and its analogue for \(G_v\).  In (77.2), cycle the leading
\(H^{-1/2}\) to the end.  In the four-factor word, the same cyclic
move gives
\[
\begin{split}
&\operatorname{Tr}\bigl(
H^{-1/2}U_uH^{-1/2}P_a
H^{-1/2}U_vH^{-1/2}N_c\bigr)\\
&\qquad=
\operatorname{Tr}\bigl(
U_u(H^{-1/2}P_aH^{-1/2})
U_v(H^{-1/2}N_cH^{-1/2})\bigr),
\end{split}
\]
which is (77.3).
\end{proof}

\begin{corollary}[Dressed Schatten bounds]
\label{cor:n1v77-schatten}
In dimension \(d=45\),
\[
 |\operatorname{Tr}(G_uQ_b)|
 \leq\sqrt d\,\|U_u\|_{\rm op}
 \|\widehat Q_b\|_F,                                        \tag{77.4}
\]
and
\[
 |\operatorname{Tr}(G_uP_aG_vN_c)|
 \leq
 \|U_u\|_{\rm op}\|U_v\|_{\rm op}
 \|\widehat P_a\|_F\|\widehat N_c\|_F.                       \tag{77.5}
\]
\end{corollary}

\begin{proof}
For (77.4), use
\(\|U_u\|_F\leq\sqrt d\|U_u\|_{\rm op}\) and
Hilbert--Schmidt duality.  For (77.5), apply Hilbert--Schmidt
duality after
\[
 \|U_u\widehat P_aU_v\|_F
 \leq\|U_u\|_{\rm op}\|\widehat P_a\|_F
      \|U_v\|_{\rm op}.
\]
\end{proof}

This removes the absolute products
\(\|G_u\|\|P_a\|_F\) and
\(\|G_v\|\|N_c\|_F\).  The covariance part is dimensionless and
the vertices are measured in the energy metric on the precise
input and output fibres.

\subsection{A dressed response majorant}

Let \(u_{r,m}\) bound the V68 normalized covariance jets, and let
\[
\widehat p_{r,m},\quad
\widehat n_{r,m},\quad
\widehat q_{r,m}
\]
bound the Frobenius norms of the correspondingly differentiated,
energy-dressed cubic and quartic matrices.  Distribution of two
external and two internal derivatives, exactly as in V62, now gives
\[
\begin{split}
\widehat B_{{\cal T},e}^{(\mu)}
:={}&\frac{\sqrt{45}}2\left(
u_{0,2}\widehat q_{2,0}
+2u_{0,1}\widehat q_{2,1}
+u_{0,0}\widehat q_{2,2}\right),                            \tag{77.6}\\
\widehat B_{{\cal B},e}^{(\mu)}
:={}&|c_e|
\sum_{(a,b,c)\in{\cal C}_2(3)}{2\choose a,b,c}
\sum_{(u,v,w,x)\in{\cal C}_2(4)}{2\choose u,v,w,x}\\
&\qquad\qquad\times
u_{0,u}\widehat p_{a,v}u_{b,w}\widehat n_{c,x}.              \tag{77.7}
\end{split}
\]

\begin{proposition}[Validity conditional only on dressed stocks]
\label{prop:n1v77-majorant}
If the displayed \(u,\widehat p,\widehat n,\widehat q\) are valid
global bounds for the indicated pointwise dressed matrices, then
\[
 \sup_k|\partial_{k_\mu}^2{\cal R}_e^{[2]}(k)|
 \leq
 \widehat B_{{\cal T},e}^{(\mu)}
 +\widehat B_{{\cal B},e}^{(\mu)}.                           \tag{77.8}
\]
No extra inverse-floor factor occurs in (77.6)--(77.7).
\end{proposition}

\begin{proof}
Apply the complete two-variable Leibniz expansion used in
(62.7)--(62.9).  Replace each tadpole word by (77.2) and each
bubble word by (77.3), then apply
Corollary~\ref{cor:n1v77-schatten}.  The normalized covariance
recursion already includes every differentiated inverse identity.
Summing the multinomial terms gives (77.6)--(77.7).
\end{proof}

\subsection{Deterministic scale diagnostic}

The exact V60 Laurent banks were evaluated on the uniform \(4^4\)
grid.  At each point the positive Hessian was diagonalized and all
external jets were dressed as in (77.1).  This is a floating
diagnostic, not a global enclosure.  The sampled maxima were
\[
\begin{array}{c|ccc}
&\text{orders }0&1&2\\ \hline
\|\widehat P_{\rm W}\|_F
&44.2348321&25.9538634&42.9583378\\
\|\widehat P_*\|_F
&9.07327825&5.58791506&6.49637818\\
\|U_{\rm W}\|_{\rm op}
&1&1.78057833&6.80374779\\
\|U_*\|_{\rm op}
&1&1.87606347&7.50222773.
\end{array}                                                   \tag{77.9}
\]
The plus and minus vertex maxima agree on this grid.  The dressed
quartic maxima in external orders \(0,1,2\) were
\[
\begin{split}
{\rm Wilson}:&(109.680987,0,124.875562),\\
{\rm endpoint}:&(99.783699,0,185.415538).                   \tag{77.10}
\end{split}
\]
Using only the pointwise second-external-order form of
(77.6)--(77.7), the largest sampled dressed bounds were
\[
\begin{array}{c|c|c}
e&\text{bubble}&\text{total}\\ \hline
{\rm Wilson}&7243.2764&7558.4655\\
*&2616.7691&3065.6413.
\end{array}                                                   \tag{77.11}
\]
The corresponding sampled actual absolute responses were at most
\(1047.668\) and \(202.864\).  These values do not include the two
internal derivatives needed for a quadrature theorem; they show
only that the energy metric reaches the observed scale before a
global certificate is imposed.

\begin{record}[V77 dressed diagnostic]
\label{rec:n1v77-executable}
The executable is
\[
 \texttt{scripts/explore\_ym\_v77\_energy\_dressed\_vertices.py}. \tag{77.12}
\]
Its source record is
\[
\begin{array}{c|c}
\text{lines}&176\\
\text{bytes}&6871\\
\text{SHA--256}&
\texttt{9561DAD6DDDBBEF40F2DC7AA42FB55AB32BE9F2FF959F087CB9716D4407AC8BA}.
\end{array}                                                   \tag{77.13}
\]
For \(M=4\), the canonical payload SHA--256 is
\[
 \boxed{
\texttt{04F1A3568D46E5AB5E75D3E288370E4D66BD0A0F04E7000453D9EABB9FD6571A}.}
                                                               \tag{77.14}
\]
\end{record}

\begin{decision}[V78 certifies dressed dyadic anchors]
\label{dec:n1v77-next}
At all sixteen dyadic corners, compute
\[
 \|\widehat X\|_F^2
 =\operatorname{Tr}(X^*H^{-1}XH^{-1})                       \tag{77.15}
\]
with exact Gaussian-rational arithmetic for every required
external cubic and quartic jet.  Record rational upper bounds and
the exact action/jet maximizers.  Then derive a relative
corner-to-box transport theorem for (77.15), using the V72 pencil
floors and V68 normalized covariance jets.  Do not substitute
\(\|X\|_F/\lambda_e\), which would collapse (77.15) back to V69.
\end{decision}

\begin{assessment}[Progress after compact V77]
\label{ass:n1v77-status}
The complete response words now admit exact energy-dressed
factorizations.  This is the first bound architecture in the
current sequence whose diagnostic scale is within thousands rather
than \(10^{16}\) or larger.  The remaining immediate task is not a
new algebraic response identity but a global certificate for the
dressed vertex and insertion norms.
\end{assessment}

\begin{firewall}[Claim boundary after compact V77]
\label{fire:n1v77-claim}
Theorems (77.2)--(77.8) are rigorous, but (77.9)--(77.11) are
floating finite-grid diagnostics.  No dressed global stocks,
quadrature sign enclosure, higher-loop or nonperturbative control,
SCIC, ADMC, continuum measure, reflection positivity, exponential
clustering, or Yang--Mills mass gap has yet been proved.
\end{firewall}

\subsection{Parent record}

This V77 note was formed by appending the preceding material to the
validated compact V76, whose record was
\[
\begin{array}{c|c}
\text{lines}&26563\\
\text{bytes}&905633\\
\text{SHA--256}&
\texttt{A1347C803379291E6463E59BEB597DF56C10A38BAA7F995F04A302F26512452F}.
\end{array}
\]

% =============================================================================
% END APPENDED COMPACT-SERIES V77 MATERIAL
% =============================================================================

% =============================================================================
% BEGIN APPENDED COMPACT-SERIES V78 MATERIAL
% =============================================================================

\section{Compact-series V78: certified energy-dressed corner anchors}
\label{sec:n1v78}

\subsection{A posteriori rational inversion in the energy metric}

The finite-grid numbers in V77 cannot be used in a proof.  This
section replaces their values at the sixteen dyadic corners by
rational certificates.  The point is to avoid exact inversion of
every large rational matrix while retaining an exact conclusion.

Fix one of the two actions, one corner
\(\epsilon\in\{-1,1\}^4\), and write
\[
 H=H_e(\epsilon)\succeq\lambda_e I,\qquad G=H^{-1},          \tag{78.1}
\]
where the already proved global floors are
\[
 \lambda_{\rm W}={49\over1250},\qquad
 \lambda_*={343\over1250}.                                  \tag{78.2}
\]
Let \(R=R^T\) be any rational trial inverse and put
\[
 E=I-RH,\qquad \eta\geq\|E\|_F,\qquad
 d={\eta\over\lambda_e},\qquad
 r={1\over\lambda_e}+d.                                     \tag{78.3}
\]

\begin{lemma}[Exact residual control of a dressed norm]
\label{lem:n1v78-residual-dressing}
Suppose \(\eta<1\).  Let \(B\) be real and let the actual stored
matrix \(X\) be either \(B\) or \(iB\).  Define
\[
 s_R(B):=\operatorname{Tr}(B^TRBR),\qquad
 b^2:=\|B\|_F^2.                                             \tag{78.4}
\]
Then
\[
 \|G-R\|_{\rm op}\leq d,\qquad \|R\|_{\rm op}\leq r,          \tag{78.5}
\]
and
\[
 \boxed{\ 
 \|H^{-1/2}XH^{-1/2}\|_F^2
 \leq s_R(B)+b^2(2dr+d^2).\ }                                \tag{78.6}
\]
Every quantity on the right is rational when \(H,R,B\) are
rational.
\end{lemma}

\begin{proof}
Since \(HG=I\),
\[
 G-R=(I-RH)G=EG.
\]
The floor in (78.1) gives
\[
 \|G-R\|_{\rm op}\leq\|E\|_{\rm op}\|G\|_{\rm op}
 \leq{\eta\over\lambda_e}=d,
\]
and the triangle inequality gives the second part of (78.5).
Put \(D=G-R\).  For both allowed phases of \(X\),
\[
 \|H^{-1/2}XH^{-1/2}\|_F^2
 =\operatorname{Tr}(B^TGBG).
\]
Expanding \(G=R+D\) leaves \(s_R(B)\), two mixed terms,
and one \(D,D\) term.  Hilbert--Schmidt duality and
\(\|AB\|_F\leq\|A\|_F\|B\|_{\rm op}\) bound their absolute
values by \(b^2dr,b^2rd,b^2d^2\), respectively.  Their sum is
(78.6).
\end{proof}

No positivity assertion about the trial matrix \(R\) is used.
In the certificate below, \(R=C/10^8\), where \(C\) is obtained by
rounding a symmetric floating approximation.  The selection of
\(C\) is floating, but the residual, (78.3), (78.4), and (78.6)
are recomputed in integer and rational arithmetic.  Thus floating
point selects a trial object and proves nothing.

\subsection{All 288 external-jet anchors}

For each action and corner, the verifier evaluates the exact
Hessian and the three external jets in each of the quartic, plus
cubic, and minus cubic banks.  There are
\[
 2\cdot16\cdot3\cdot3=288                                   \tag{78.7}
\]
dressed norms.  A rational square-root ceiling with denominator
\(10^{12}\) is applied only after (78.6).

\begin{theorem}[Certified dyadic dressed-vertex anchors]
\label{thm:n1v78-dressed-anchors}
At all sixteen dyadic corners, the maximum dressed Frobenius norms
in external orders \(0,1,2\) obey the following exact rational
upper bounds:
\[
\begin{array}{c|c|ccc}
e&\text{bank}&0&1&2\\ \hline
{\rm W}&\widehat Q&
{1107373017551\over12500000000}&0&
{92169638209023\over10^{12}}\\
{\rm W}&\widehat P^\pm&
{14543039572249\over10^{12}}&
{2049826641811\over200000000000}&
{18013514808271\over10^{12}}\\
*&\widehat Q&
{12525457590971\over125000000000}&0&
{18655688686103\over100000000000}\\
*&\widehat P^\pm&
{5480909552031\over10^{12}}&
{822831823817\over250000000000}&
{2777774200561\over500000000000}.
\end{array}                                                   \tag{78.8}
\]
Numerically, the rows of (78.8) are
\[
\begin{array}{c|c|ccc}
{\rm W}&\widehat Q&88.589842&0&92.169639\\
{\rm W}&\widehat P^\pm&14.543040&10.249134&18.013515\\
*&\widehat Q&100.203661&0&186.556887\\
*&\widehat P^\pm&5.480910&3.291328&5.555549.
\end{array}                                                   \tag{78.9}
\]
The result is an exact finite-corner theorem, not a sampled
numerical assertion.
\end{theorem}

\begin{proof}
The Hessian matrices are evaluated coefficient-exactly and divided
by three according to the V60 normalization.  At every corner the
script forms the integer numerator of
\[
 I-{C\over10^8}H
\]
and takes a rational upper ceiling of its Frobenius norm.  The
largest residual bounds over all corners are
\[
 \eta_{\rm W}\leq {32059\over125000000000},\qquad
 \eta_*\leq {11894027\over10^{12}},                          \tag{78.10}
\]
both strictly below one.  With (78.2), the corresponding certified
inverse errors are
\[
 d_{\rm W}\leq {32059\over4900000000},\qquad
 d_*\leq {11894027\over274400000000}.                        \tag{78.11}
\]

For a vertex matrix \(B=L^{-1}B_{\mathbb Z}\), the trial term in
(78.4) is evaluated as
\[
 s_R(B)=
 { \operatorname{Tr}
   (B_{\mathbb Z}^TCB_{\mathbb Z}C)
  \over L^2\,10^{16}},                                      \tag{78.12}
\]
using integer matrix products.  The value \(b^2\), the
perturbation in (78.6), and its sum with (78.12) are Fractions.
Lemma~\ref{lem:n1v78-residual-dressing} therefore certifies every
one of the 288 values.  Taking the exact maximum separately by
action, bank, and order, followed by the rational square-root
ceiling, gives (78.8).
\end{proof}

For reproducibility, one selected maximizing corner in each entry
of (78.8) is
\[
\begin{array}{c|c|ccc}
e&\text{bank}&0&1&2\\ \hline
{\rm W}&\widehat Q&(-,-,+,-)&(-,-,-,-)&(-,-,+,+)\\
{\rm W}&\widehat P^\pm&(-,-,+,-)&(-,+,+,+)&(-,-,+,+)\\
*&\widehat Q&(-,-,+,-)&(-,-,-,-)&(-,-,+,-)\\
*&\widehat P^\pm&(-,-,+,+)&(+,+,+,-)&(-,+,+,-).
\end{array}                                                   \tag{78.13}
\]
The zero first-derivative quartic row is exact.  Equality of the
plus and minus maxima is checked rather than assumed.

\begin{record}[V78 dressed-anchor verifier]
\label{rec:n1v78-executable}
The executable is
\[
 \texttt{scripts/verify\_ym\_v78\_dressed\_corner\_anchors.py}. \tag{78.14}
\]
Its source record is
\[
\begin{array}{c|c}
\text{lines}&266\\
\text{bytes}&9064\\
\text{SHA--256}&
\texttt{6D785ACC806B6A4289AA51BE11BE9779A3396361C484E30EC3146EC2A622910A}.
\end{array}                                                   \tag{78.15}
\]
The script checks the 32 inverse residuals, all 288 nonnegative
upper squares, the two global floors, and the input hashes.  Its
canonical payload SHA--256 is
\[
 \boxed{
\texttt{47A16B8ADC207B4EAAF29CA30AD0385B7FF04EA043D13CFC2B99FB2A8CE130FD}.}
                                                               \tag{78.16}
\]
\end{record}

\subsection{A noncommutative corner-to-box transport lemma}

The next obstruction is globality.  It is important that transport
be performed in the anchor energy metric; transporting \(X\) and
then dividing by the absolute floor would discard the gain in
(78.8).

\begin{proposition}[Relative energy transport]
\label{prop:n1v78-relative-transport}
Let \(H_0,H\succ0\), and set
\[
 A=H_0^{-1/2}(H-H_0)H_0^{-1/2},\qquad \|A\|_{\rm op}\leq\rho<1.
                                                               \tag{78.17}
\]
For arbitrary matrices \(X_0,\Delta X\), with \(X=X_0+\Delta X\),
\[
\boxed{
 \|H^{-1/2}XH^{-1/2}\|_F
 \leq {1\over1-\rho}\left(
 \|H_0^{-1/2}X_0H_0^{-1/2}\|_F+
 \|H_0^{-1/2}\Delta XH_0^{-1/2}\|_F\right).}                 \tag{78.18}
\]
No commutation between \(H,H_0,X\) is required.
\end{proposition}

\begin{proof}
Write \(K=H_0^{1/2}H^{-1}H_0^{1/2}=(I+A)^{-1}\).  The spectral
assumption gives \(\|K\|_{\rm op}\leq(1-\rho)^{-1}\).
For
\[
 B=H_0^{-1/2}XH_0^{-1/2}
\]
cyclicity gives
\[
 \|H^{-1/2}XH^{-1/2}\|_F^2
 =\operatorname{Tr}(B^*KBK)
 \leq\|K\|_{\rm op}^2\|B\|_F^2.
\]
Apply the triangle inequality to the two summands of \(B\).
\end{proof}

In particular, for Laurent expansions about a corner
\(\epsilon\), it is enough to certify
\[
\begin{split}
\rho(k)&\leq
\sum_\alpha
\|H_\epsilon^{-1/2}H_\alpha H_\epsilon^{-1/2}\|_{\rm op}
\,|e^{i\alpha\cdot k}-\epsilon^\alpha|,\\
\delta_X(k)&\leq
\sum_\alpha
\|H_\epsilon^{-1/2}X_\alpha H_\epsilon^{-1/2}\|_F
\,|e^{i\alpha\cdot k}-\epsilon^\alpha|.                     \tag{78.19}
\end{split}
\]
Whenever the first right side is below one, (78.18) bounds the
box by the anchor in (78.8) plus \(\delta_X\).  Formula (78.19)
is a valid coefficientwise reduction, but its coefficients have
not yet been compiled; no box bound is claimed here.

\begin{decision}[V79 compiles anchor-relative transport banks]
\label{dec:n1v78-next}
For each dyadic anchor, compute rational upper bounds for the
energy-dressed Laurent coefficient matrices appearing in
(78.19), preserving exponent grouping before norming.  Use them
to determine certified anisotropic boxes with \(\rho<1\), and
cover the torus by subdividing only boxes that fail.  Compare the
resulting global dressed stocks in (77.6)--(77.7) with V69 before
replacing the operative bound.
\end{decision}

\begin{assessment}[Progress after compact V78]
\label{ass:n1v78-status}
The energy-dressed programme now has exact anchors and an exact
noncommutative transport inequality.  This removes the finite-grid
status of the dyadic values themselves.  The unresolved local
problem is the coefficient-efficient certification of
\(\rho\) and \(\delta_X\) on a covering of the torus.
\end{assessment}

\begin{firewall}[Claim boundary after compact V78]
\label{fire:n1v78-claim}
Theorem~\ref{thm:n1v78-dressed-anchors} concerns only sixteen
corners, and Proposition~\ref{prop:n1v78-relative-transport} is
conditional on a boxwise relative bound.  No global dressed
stock, improved quadrature sign enclosure, higher-loop or
nonperturbative estimate, SCIC, ADMC, continuum measure,
reflection positivity, exponential clustering, or Yang--Mills
mass gap is proved.
\end{firewall}

\subsection{Parent record}

This V78 note was formed by appending the preceding material to the
validated compact V77, whose record was
\[
\begin{array}{c|c}
\text{lines}&26825\\
\text{bytes}&913923\\
\text{SHA--256}&
\texttt{6CA0089783961EDA3A6C4F043EEFF11083DAB114A37CE742B34FA34E559260A7}.
\end{array}
\]

% =============================================================================
% END APPENDED COMPACT-SERIES V78 MATERIAL
% =============================================================================

% =============================================================================
% BEGIN APPENDED COMPACT-SERIES V79 MATERIAL
% =============================================================================

\section{Compact-series V79: anchor-relative transport banks}
\label{sec:n1v79}

\subsection{Coefficientwise dressing without square roots}

V78 reduced a global dressed bound to energy-relative transport
from a certified anchor.  We now compile the first exact transport
banks.  Fix a dyadic corner \(\epsilon\), abbreviate
\[
 H_0=H_e(\epsilon),\qquad G_0=H_0^{-1},
\]
and retain the rational trial inverse \(R\), the error \(d\), and
the trial bound \(r\) from (78.3)--(78.5).

\begin{lemma}[Residual bound for a Laurent coefficient]
\label{lem:n1v79-coefficient}
For every real Laurent coefficient matrix \(B\), and also when
the actual coefficient is \(iB\),
\[
\begin{split}
&\|H_0^{-1/2}BH_0^{-1/2}\|_F^2\\
&\qquad\leq
\|RB\|_F\|BR\|_F+\|B\|_F^2(2dr+d^2).                       \tag{79.1}
\end{split}
\]
The right side has a rational upper bound obtained solely from
integer products and rational square-root ceilings.
\end{lemma}

\begin{proof}
The proof of Lemma~\ref{lem:n1v78-residual-dressing} bounds the
difference between the squared dressed norm and
\(\operatorname{Tr}(B^TRBR)\) by the second term in (79.1).
Since
\[
 \operatorname{Tr}(B^TRBR)
 =\operatorname{Tr}\bigl((RB)^T(BR)\bigr),
\]
Hilbert--Schmidt duality bounds its absolute value by the first
term.  This also proves the assertion with the global phase \(i\).
\end{proof}

This form is computationally cheaper than the exact trace in
(78.12) for sparse coefficient matrices, and it retains the
geometry of \(R\).  It is not the absolute estimate
\(\|B\|_F/\lambda_e\).

\subsection{Exact slope banks}

Write the coefficient-grouped Laurent expansions as
\[
 H_e(k)=\sum_\alpha H_\alpha e^{i\alpha\cdot k},\qquad
 X(k)=\sum_\alpha X_\alpha e^{i\alpha\cdot k}.               \tag{79.2}
\]
Let \(\beta_\epsilon(B)\) denote the square root of the right side
of (79.1), rounded upward rationally.  Define four directional
stocks
\[
\begin{split}
 L_{H,\mu}(\epsilon)
  &=\sum_\alpha|\alpha_\mu|\,
    \beta_\epsilon(H_\alpha),\\
 L_{X,\mu}(\epsilon)
  &=\sum_\alpha|\alpha_\mu|\,
    \beta_\epsilon(X_\alpha).                               \tag{79.3}
\end{split}
\]
All equal exponents are grouped into one matrix before
\(\beta_\epsilon\) is applied.

\begin{theorem}[Certified anchor-relative Laurent transport]
\label{thm:n1v79-transport}
Let \(k=\epsilon+\delta\) in corner-centered angular coordinates,
with \(|\delta_\mu|\leq h_\mu\), and suppose
\[
 \rho:=\sum_{\mu=0}^3L_{H,\mu}(\epsilon)h_\mu<1.             \tag{79.4}
\]
Then
\[
\boxed{
 \|H_e(k)^{-1/2}X(k)H_e(k)^{-1/2}\|_F
 \leq
 {\,\beta_\epsilon(X(\epsilon))
       +\sum_\mu L_{X,\mu}(\epsilon)h_\mu
  \over1-\rho}.}                                            \tag{79.5}
\]
This holds for each external order of the quartic and both cubic
vertex banks.
\end{theorem}

\begin{proof}
For every integer exponent,
\[
 |e^{i\alpha\cdot\delta}-1|
 \leq|\alpha\cdot\delta|
 \leq\sum_\mu|\alpha_\mu|h_\mu.                             \tag{79.6}
\]
The triangle inequality, (79.1), and (79.2) therefore give
\[
\begin{split}
\|H_0^{-1/2}(H_e(k)-H_0)H_0^{-1/2}\|_{\rm op}
 &\leq\rho,\\
\|H_0^{-1/2}(X(k)-X(\epsilon))H_0^{-1/2}\|_F
 &\leq\sum_\mu L_{X,\mu}h_\mu .
\end{split}
\]
Apply Proposition~\ref{prop:n1v78-relative-transport}.
\end{proof}

\subsection{The half-relative boxes and their obstruction}

The certificate chooses an isotropic half-width
\[
 h_\epsilon=
 {1\over2\sum_\mu L_{H,\mu}(\epsilon)},                     \tag{79.7}
\]
so that \(\rho=1/2\) and the transport factor in (79.5) is exactly
two.  Over the sixteen anchors the exact ranges are
\[
\begin{split}
 {250000000000\over55149878988173}
 &\leq h_{\epsilon,{\rm W}}\leq
 {250000000000\over39872539298681},\\
 {125000000000\over58479841502711}
 &\leq h_{\epsilon,*}\leq
 {12500000000\over4178138046561}.                           \tag{79.8}
\end{split}
\]
Their decimal ranges are respectively
\[
 0.00453310\ldots\leq h_{\epsilon,{\rm W}}
 \leq0.00626998\ldots,
\quad
 0.00213748\ldots\leq h_{\epsilon,*}
 \leq0.00299177\ldots.                                      \tag{79.9}
\]

\begin{corollary}[Rigorous local dressed stocks]
\label{cor:n1v79-local-stocks}
On each half-relative box (79.7), the maxima over all sixteen
anchors admit the following simple rational-decimal upper bounds:
\[
\begin{array}{c|c|ccc}
e&\text{bank}&0&1&2\\ \hline
{\rm W}&\widehat Q&178.69&0&186.19\\
{\rm W}&\widehat P^\pm&29.84&21.01&36.87\\
*&\widehat Q&202.57&0&381.67\\
*&\widehat P^\pm&11.33&7.10&11.79.
\end{array}                                                   \tag{79.10}
\]
Here, for example, \(178.69\) means the exact rational
\(17869/100\); every exact verifier output is strictly smaller
than the displayed rational.
\end{corollary}

\begin{proof}
For every action, anchor, bank, and order, insert (79.7) into
(79.5).  The anchor term is the V78 residual certificate, not a
floating norm.  Exact maximization over the sixteen values gives
(79.10); direct comparison of Fractions proves the stated rounded
upper bounds.
\end{proof}

The local theorem does not furnish a torus cover.  Indeed, every
radius in (79.8) is below \(3/2\), while
\(\pi/2>3/2\).  The point with all four coordinates \(\pi/2\)
has torus sup-distance \(\pi/2\) from every dyadic corner.  Hence
it belongs to none of the thirty-two action-labelled boxes.
This is an exact coverage obstruction, not a mesh diagnostic.

\begin{record}[V79 transport-bank verifier]
\label{rec:n1v79-executable}
The executable is
\[
 \texttt{scripts/verify\_ym\_v79\_anchor\_transport\_banks.py}. \tag{79.11}
\]
Its source record is
\[
\begin{array}{c|c}
\text{lines}&270\\
\text{bytes}&10269\\
\text{SHA--256}&
\texttt{24A603CBA095572FF227404C0BB631C3C6B3656598F7ABAF8BAC10085FB1C0A8}.
\end{array}                                                   \tag{79.12}
\]
It performs \(11904\) coefficient-dressing evaluations.  It
checks all 32 inverse residuals, positivity of every radius and
local upper bound, the dyadic coverage obstruction, exponent
grouping, and the inherited input hashes.  Its canonical payload
SHA--256 is
\[
 \boxed{
\texttt{091110388FA80C1BF0265B4E56EB1A5164B26C3BB86ACFF8860A4446BBF0F02C}.}
                                                               \tag{79.13}
\]
\end{record}

\begin{decision}[V80 audits upstream and replaces raw phase transport]
\label{dec:n1v79-next}
Perform the mandatory current SM/NS audit.  Then replace the
coefficientwise inequality (79.6) at one of two earlier levels:
either compile the complete energy-metric difference matrix
before taking its norm, or introduce additional exactly
representable anchors and certify them a posteriori.  Do not
attempt a torus cover by blindly replicating the small boxes in
(79.8), and do not insert the local stocks (79.10) into the global
response majorant.
\end{decision}

\begin{assessment}[Progress after compact V79]
\label{ass:n1v79-status}
The first rigorous energy-relative neighborhoods and their local
vertex stocks are now available.  The same calculation proves
that dyadic anchors plus coefficientwise phase transport are far
too local.  The bottleneck has therefore moved upstream from
response algebra to cancellation in the complete transported
Hessian and vertex differences.
\end{assessment}

\begin{firewall}[Claim boundary after compact V79]
\label{fire:n1v79-claim}
The bounds in (79.10) are local only, and the certified boxes do
not cover the torus.  V69 remains the operative global response
majorant.  No improved global quadrature sign enclosure,
higher-loop or nonperturbative estimate, SCIC, ADMC, continuum
measure, reflection positivity, exponential clustering, or
Yang--Mills mass gap is proved.
\end{firewall}

\subsection{Parent record}

This V79 note was formed by appending the preceding material to the
validated compact V78, whose record was
\[
\begin{array}{c|c}
\text{lines}&27139\\
\text{bytes}&924440\\
\text{SHA--256}&
\texttt{5186A400C54B788054377499FBA689FFF01F682E111FFE0F939CED3E0D1ACAF7}.
\end{array}
\]

% =============================================================================
% END APPENDED COMPACT-SERIES V79 MATERIAL
% =============================================================================

% =============================================================================
% BEGIN APPENDED COMPACT-SERIES V80 MATERIAL
% =============================================================================

\section{Compact-series V80: companion audit and differential
energy transport}
\label{sec:n1v80}

\subsection{Mandatory audit of the current companion programmes}

The two current companion files in \texttt{latex\_notes} were
inspected read-only before choosing the next YM step.  Their frozen
records are
\[
\begin{array}{c|r|r|c}
\text{programme}&\text{version}&\text{bytes}&\text{SHA--256}\\ \hline
{\rm SM\!-\!Einstein}&V29&329888&
\texttt{AAB33D8695CA1D1E036A6532E414B57266E8EAA73CDA1FAC9DBFECB621C22D47}\\
{\rm Navier\!-\!Stokes}&V26&278283&
\texttt{82C887F8C2C69E4F28FAD7C3DC8DE5B9099CB5A4BF431EBA1FFF3E91935978AD}.
\end{array}                                                   \tag{80.1}
\]
The files have respectively \(9340\) and \(5319\) physical lines.

The SM note proves an abstract subsequential compactness theorem
for normal channels on von Neumann algebras with separable
preduals, a finite-rank phase-space sufficient condition, stable
tensor amplification, and a conditional split-screen recovery
limit.  Its phase-space compression and split hypotheses are not
populated.  These results could become relevant only after the YM
programme constructs a continuum local-algebra family and the
required compact recovery orbits.  They do not provide a
finite-dimensional momentum-space estimate for the present
response problem, so no SM theorem is imported here.

The NS note computes pointwise rank-one norms of first and second
Oseen-kernel derivatives.  Its first-derivative refinement gives no
gain, while its structured second-derivative input gives a modest
gain on a bounded radial interval.  This is not a Yang--Mills
estimate.  Its useful methodological lesson is that the prescribed
input contraction must be made before the ambient tensor norm.
For V79, the analogous prescribed input is the single phase vector
\(\delta\), and (79.6) destroys it by separating every Laurent
coefficient.  This audit therefore selects complete pathwise
matrix contraction as the next YM direction.

\begin{proposition}[Companion-audit decision]
\label{prop:n1v80-audit}
Neither companion note closes a YM hypothesis or changes the V79
claim boundary.  The next transport estimate must norm
\(\dot H(k(t))\) and \(\dot X(k(t))\) as complete matrices along a
path before any coefficient decomposition.  The later SM predual
compactness theorem is retained only as a possible continuum-stage
interface.
\end{proposition}

\begin{proof}
The current obstruction is the failure of sixteen finite
energy-metric boxes to cover the compact momentum torus.  The SM
theorems assume channel-orbit compactness on local operator
algebras and contain no Hessian or vertex bound.  The NS estimates
concern a different kernel, but their contraction order identifies
the precise avoidable relaxation in (79.6).  Thus only the stated
choice of proof architecture follows from the audit.
\end{proof}

\subsection{Pathwise whitening}

Let \(I=[0,T]\).  Suppose
\[
 H:I\longrightarrow{\rm Herm}_{++}(d),\qquad
 X:I\longrightarrow M_d(\mathbb C)                          \tag{80.2}
\]
are \(C^1\).  Put
\[
\begin{split}
 G(t)&=H(t)^{-1},\\
 a(t)&=\|H(t)^{-1/2}\dot H(t)H(t)^{-1/2}\|_{\rm op},\\
 c(t)&=\|H(t)^{-1/2}\dot X(t)H(t)^{-1/2}\|_F,\\
 f(t)&=\|H(t)^{-1/2}X(t)H(t)^{-1/2}\|_F.                    \tag{80.3}
\end{split}
\]
These quantities use the instantaneous energy metric.  No
derivative of a selected matrix square root is required.

\begin{theorem}[Differential energy-transport inequality]
\label{thm:n1v80-differential-transport}
With (80.2)--(80.3), the upper right Dini derivative satisfies
\[
 \boxed{\ D^+f(t)\leq c(t)+a(t)f(t).\ }                      \tag{80.4}
\]
Consequently, with \(A(t)=\int_0^t a(s)\,ds\),
\[
\boxed{
 f(t)\leq e^{A(t)}
 \left(f(0)+\int_0^t e^{-A(s)}c(s)\,ds\right).}             \tag{80.5}
\]
In particular, if \(a\leq\bar a\) and \(c\leq\bar c\), then
\[
 f(t)\leq
 \begin{cases}
 e^{\bar a t}f(0)+{\bar c\over\bar a}
 (e^{\bar a t}-1),&\bar a>0,\\
 f(0)+\bar c\,t,&\bar a=0.
 \end{cases}                                                 \tag{80.6}
\]
Unlike (78.18), this theorem has no condition \(\rho<1\).
\end{theorem}

\begin{proof}
The squared dressed norm is
\[
 S(t):=f(t)^2=\operatorname{Tr}(X(t)^*G(t)X(t)G(t)).         \tag{80.7}
\]
Since \(\dot G=-G\dot HG\), differentiation gives two vertex
terms and two Hessian terms.  At a fixed \(t\), write
\[
 B=H^{-1/2}XH^{-1/2},\quad
 C=H^{-1/2}\dot XH^{-1/2},\quad
 D=H^{-1/2}\dot HH^{-1/2}.
\]
Cyclicity of the trace rewrites the four contributions as
\[
 \dot S=
 2\operatorname{Re}\operatorname{Tr}(C^*B)
 -\operatorname{Tr}(B^*DB)
 -\operatorname{Tr}(B^*BD).                                \tag{80.8}
\]
Here \(D=D^*\).  Hilbert--Schmidt duality and the operator ideal
inequality imply
\[
 \dot S\leq2cf+2af^2.                                       \tag{80.9}
\]
At points where \(f>0\), division by \(2f\) proves (80.4).
At a zero of \(f\), apply the same calculation to
\((S+\varepsilon^2)^{1/2}\) and let
\(\varepsilon\downarrow0\), which proves the Dini form.
Multiplying (80.4) by \(e^{-A(t)}\) and integrating proves
(80.5); (80.6) follows by using constant upper bounds.
\end{proof}

\begin{corollary}[Loewner length along the same path]
\label{cor:n1v80-loewner-length}
For \(0\leq s\leq t\leq T\),
\[
 e^{-\int_s^t a}\,H(s)\preceq H(t)
 \preceq e^{\int_s^t a}\,H(s).                              \tag{80.10}
\]
\end{corollary}

\begin{proof}
For a fixed vector \(v\), set \(q(t)=v^*H(t)v>0\).  By the
definition of \(a(t)\),
\[
 |\dot q(t)|\leq a(t)q(t).
\]
Integrate the two scalar logarithmic inequalities and then vary
\(v\).
\end{proof}

\subsection{Why the differential form goes upstream}

For a momentum segment \(k(t)=k_0+t\xi\), the actual inputs to
(80.3) are the complete matrices
\[
 \dot H(k(t))=\sum_\mu\xi_\mu\partial_\mu H(k(t)),\qquad
 \dot X(k(t))=\sum_\mu\xi_\mu\partial_\mu X(k(t)).           \tag{80.11}
\]
Thus cancellations between directions, entries, and equal phases
survive until after the structured contraction by \(\xi\).
In contrast, (79.3) replaces (80.11) by a sum of separate
coefficient norms before the path is even chosen.

The V68 identity identifies \(a(t)\) with the norm of the
first normalized Hessian jet along \(\xi\).  The V77 dressed
architecture identifies \(c(t)\) as another dressed vertex norm.
Therefore a certified path compiler for the complete matrices in
(80.11) can feed (80.5) directly and can concatenate segments
without requiring a relative displacement smaller than one.

\begin{decision}[V81 compiles complete pathwise inputs]
\label{dec:n1v80-next}
Start with coordinate segments joining the \(4^4\) diagnostic grid.
At each segment, evaluate the complete matrices in (80.11), use
rational residual inverses at the segment anchors, and certify
upper envelopes for \(a(t)\) and \(c(t)\) before applying
(80.5).  First run a floating diagnostic to determine whether
coordinate paths or diagonal paths have smaller energy length;
then certify only the better route.  Do not reuse the
coefficientwise slopes (79.3) as path bounds.
\end{decision}

\begin{assessment}[Progress after compact V80]
\label{ass:n1v80-status}
The mandatory companion audit is complete.  The local
corner-to-box obstruction has been replaced by a rigorous
pathwise transport theorem that requires no small relative
displacement and preserves the full directional matrix input.
The remaining finite-dimensional task is to certify the energy
length and dressed vertex speed on a torus path network.
\end{assessment}

\begin{firewall}[Claim boundary after compact V80]
\label{fire:n1v80-claim}
Theorem~\ref{thm:n1v80-differential-transport} is an abstract
finite-matrix transport result.  No global path envelopes have yet
been populated.  V69 remains the operative global response
majorant.  No improved quadrature sign enclosure, higher-loop or
nonperturbative estimate, SCIC, ADMC, continuum measure,
reflection positivity, exponential clustering, or Yang--Mills
mass gap is proved.
\end{firewall}

\subsection{Parent record}

This V80 note was formed by appending the preceding material to the
validated compact V79, whose record was
\[
\begin{array}{c|c}
\text{lines}&27389\\
\text{bytes}&932846\\
\text{SHA--256}&
\texttt{36848074B512228EB24FB5DF046BB9D7BA5DB9F647AC9385DD391A6A2DEECE7A}.
\end{array}
\]

% =============================================================================
% END APPENDED COMPACT-SERIES V80 MATERIAL
% =============================================================================

% =============================================================================
% BEGIN APPENDED COMPACT-SERIES V81 MATERIAL
% =============================================================================

\section{Compact-series V81: signed-diagonal path collapse}
\label{sec:n1v81}

\subsection{Route selection at the uncovered central point}

The point \(k_c=(\pi/2,\pi/2,\pi/2,\pi/2)\) witnesses the V79
failure of the dyadic boxes to cover the torus.  A deterministic
floating diagnostic compared two ways to reach it from all sixteen
dyadic corners:
\[
\begin{array}{ll}
\text{diagonal:}&\text{one straight segment changing all four coordinates},\\
\text{coordinate:}&\text{four coordinate segments in all \(24\) orders}.
\end{array}                                                   \tag{81.1}
\]
On each segment, eight midpoint values of the complete matrices in
(80.11) were used to propagate the scalar comparison equation from
(80.4).  The smallest observed energy lengths were
\[
\begin{array}{c|cc}
e&\text{diagonal}&\text{coordinate}\\ \hline
{\rm W}&3.10623&9.00245\\
*&3.36286&9.62045.
\end{array}                                                   \tag{81.2}
\]
Thus the diagonal route is shorter by factors about \(2.90\) and
\(2.86\).  The best diagonal transported bounds for the cubic
orders were
\[
\begin{array}{c|ccc}
{\rm W}&502.12&297.74&511.00\\
*&184.10&120.91&221.18,
\end{array}                                                   \tag{81.3}
\]
whereas the coordinate bounds range from \(4.6\,10^4\) to
\(1.36\,10^5\).  For comparison, the actual floating dressed cubic
norms at \(k_c\) were
\[
 (30.795,18.107,29.570),\qquad(5.647,3.707,5.250).           \tag{81.4}
\]
All numbers in (81.2)--(81.4) are route diagnostics, not
enclosures.  They select a route but do not enter a proof.

\begin{record}[V81 path diagnostic]
\label{rec:n1v81-diagnostic}
The executable
\[
 \texttt{scripts/explore\_ym\_v81\_pathwise\_transport.py}   \tag{81.5}
\]
has \(250\) lines, \(8839\) bytes, and SHA--256
\[
 \texttt{D076AF8FE2A3915FEEA64F6041FF189CBB2216656DCB20160DB311C2B7BB42F6}.
                                                               \tag{81.6}
\]
The canonical payload SHA--256 for the eight-midpoint run is
\[
 \boxed{
\texttt{E7428C21D0D2653EE553A991ED2753C17AA9E8B5BF12C62E9B7EAD64BE0F6BF1}.}
                                                               \tag{81.7}
\]
\end{record}

\subsection{Exact reduction to a short univariate bank}

The diagonal route has an additional exact advantage.  Let
\(\epsilon=\sigma\in\{-1,1\}^4\), use starting angles
\[
 k_\mu(0)=
 \begin{cases}0,&\sigma_\mu=1,\\ \pi,&\sigma_\mu=-1,\end{cases}
\]
and put
\[
 k_\mu(t)=k_\mu(0)+{\pi\over2}\sigma_\mu t,\qquad0\leq t\leq1.
                                                               \tag{81.8}
\]

\begin{theorem}[Signed-diagonal frequency collapse]
\label{thm:n1v81-frequency-collapse}
For a matrix Laurent polynomial
\[
 X(k)=\sum_{\alpha\in\mathbb Z^4}X_\alpha e^{i\alpha\cdot k},
\]
define exact coefficient matrices
\[
 B_n^{(\sigma)}
 :=\sum_{\sigma\cdot\alpha=n}\sigma^\alpha X_\alpha,
 \qquad
 \sigma^\alpha:=\prod_{\mu=0}^3\sigma_\mu^{\alpha_\mu}.       \tag{81.9}
\]
Then along (81.8),
\[
\boxed{
 X(k(t))=\sum_{n\in\mathbb Z}B_n^{(\sigma)}
 e^{in\pi t/2},\qquad
 {dX\over dt}={i\pi\over2}\sum_n nB_n^{(\sigma)}
 e^{in\pi t/2}.}                                            \tag{81.10}
\]
At \(t=0\) the first sum is the exact corner value, and at \(t=1\)
it is the same exact Gaussian-rational center value for every
\(\sigma\).
\end{theorem}

\begin{proof}
At the starting corner,
\[
 e^{i\alpha\cdot k(0)}=\sigma^\alpha.
\]
The remaining phase on (81.8) is
\[
 \exp\!\left({i\pi t\over2}\sigma\cdot\alpha\right).
\]
Grouping all exponents with the same integer
\(n=\sigma\cdot\alpha\) proves both identities in (81.10).
At \(t=1\), a coordinate with \(\sigma_\mu=-1\) contributes
\[
 (-1)^{\alpha_\mu}i^{-\alpha_\mu}=i^{\alpha_\mu},
\]
while a positive coordinate contributes \(i^{\alpha_\mu}\)
directly.  Hence every path reconstructs
\(\prod_\mu i^{\alpha_\mu}\), the center phase.
\end{proof}

The exact bank sizes are small:
\[
\begin{array}{c|c|c|c}
e&\text{object and external order}&
\text{original exponents}&\text{diagonal frequencies}\\ \hline
{\rm W}&H&21&3\text{--}5\\
{\rm W}&Q_0,Q_2&9,11&3,\ 3\text{--}5\\
{\rm W}&P^\pm_0,P^\pm_1,P^\pm_2&9,11,15&3,\ 3\text{--}5,\ 3\text{--}5\\
*&H&57&5\text{--}7\\
*&Q_0,Q_2&51,99&7,\ 9\\
*&P^\pm_0,P^\pm_1,P^\pm_2&51,81,81&7,\ 9,\ 7\text{--}9.
\end{array}                                                   \tag{81.11}
\]
The maximum absolute frequency is two in the Wilson banks and
four in the endpoint vertex banks; the endpoint Hessian reaches
three.  After grouping, the sums of coefficient Frobenius norms
are between \(0.609\) and \(0.677\) of their originals for Wilson,
and between \(0.436\) and \(0.607\) for the nonzero endpoint
banks.  These ratios are exact rational square-root ceilings, not
sampled Fourier values.

\begin{record}[V81 frequency-collapse verifier]
\label{rec:n1v81-frequency-verifier}
The exact executable is
\[
 \texttt{scripts/verify\_ym\_v81\_diagonal\_frequency\_collapse.py}. \tag{81.12}
\]
Its source record is
\[
\begin{array}{c|c}
\text{lines}&236\\
\text{bytes}&8996\\
\text{SHA--256}&
\texttt{A86ACC6591DF1158CA6DE248668C2E47EE9E51C9243B11573FE99F45A64B2D96}.
\end{array}                                                   \tag{81.13}
\]
For every action, path, bank, and external order it reconstructs
both endpoints exactly in \(\mathbb Q(i)\).  Its canonical payload
SHA--256 is
\[
 \boxed{
\texttt{7CF225C619C45369A5D32A807C351EF49B53CD5FDF00930A284D534E1A8E1AA3}.}
                                                               \tag{81.14}
\]
\end{record}

\subsection{Certification interface}

Theorem~\ref{thm:n1v81-frequency-collapse} turns the four-variable
route problem into finitely many matrix trigonometric polynomials
of degree at most four on \(0\leq t\leq1\).  On a rational
subinterval, outward rational Taylor bounds for
\(\sin(n\pi t/2)\) and \(\cos(n\pi t/2)\), combined with a
rational trial inverse at its center, can certify:
\[
 a(t)=\|H^{-1/2}\dot HH^{-1/2}\|_{\rm op},\qquad
 c_X(t)=\|H^{-1/2}\dot XH^{-1/2}\|_F.                        \tag{81.15}
\]
Only these complete grouped matrices should be enclosed.  Splitting
the \(B_n^{(\sigma)}\) back into their original exponents would
undo the exact gain in (81.11).

\begin{decision}[V82 builds the univariate interval certificate]
\label{dec:n1v81-next}
Construct outward rational sine and cosine enclosures for the
degree-four banks (81.10), initially on \(32\) equal path
subintervals.  At each midpoint, certify the inverse by an exact
residual and bound the complete remainders of \(H,\dot H,\dot X\).
Populate piecewise constants for (81.15), integrate them through
(80.5), and subdivide only failing intervals.  The floating
values (81.2)--(81.4) must not be used as bounds.
\end{decision}

\begin{assessment}[Progress after compact V81]
\label{ass:n1v81-status}
The diagonal route is strongly preferred by the diagnostic, and
its exact algebra collapses the Laurent data to at most nine
univariate frequencies.  Exact endpoint reconstruction is checked
for every route and jet bank.  The remaining task is a rational
interval enclosure of these short trigonometric matrix paths.
\end{assessment}

\begin{firewall}[Claim boundary after compact V81]
\label{fire:n1v81-claim}
The frequency collapse and endpoint identities are rigorous.
The energy lengths and transported values in
(81.2)--(81.4) are floating diagnostics.  No global dressed stock,
improved quadrature sign enclosure, higher-loop or
nonperturbative estimate, SCIC, ADMC, continuum measure,
reflection positivity, exponential clustering, or Yang--Mills
mass gap is proved.
\end{firewall}

\subsection{Parent record}

This V81 note was formed by appending the preceding material to the
validated compact V80, whose record was
\[
\begin{array}{c|c}
\text{lines}&27619\\
\text{bytes}&941599\\
\text{SHA--256}&
\texttt{E8F71A093D525F1668D6FD69828B7AAA376197E3AC1CE7B37AC637DCFA9481A4}.
\end{array}
\]

% =============================================================================
% END APPENDED COMPACT-SERIES V81 MATERIAL
% =============================================================================

% =============================================================================
% BEGIN APPENDED COMPACT-SERIES V82 MATERIAL
% =============================================================================

\section{Compact-series V82: an a posteriori rational path-cell
certificate}
\label{sec:n1v82}

\subsection{One cell without matrix square-root intervals}

The purpose of this section is to remove a technical ambiguity from
the V81 decision.  Direct interval evaluation of \(H^{-1/2}\) is
unnecessary.  A rational trial inverse and a global Hessian floor
are sufficient even when the cell center contains non-rational
trigonometric values.

Let \(J\) be a path cell and suppose
\[
 H(t)\succeq\lambda I\quad(t\in J).                          \tag{82.1}
\]
Choose rational complex matrices
\(\widetilde H=\widetilde H^*\), \(R=R^*\), and
\(\widetilde Y\), together with rational bounds
\[
\begin{split}
 \sup_{t\in J}\|H(t)-\widetilde H\|_F&\leq h,\\
 \sup_{t\in J}\|Y(t)-\widetilde Y\|_F&\leq y,\\
 \|R\|_{\rm op}&\leq r,\\
 \|I-R\widetilde H\|_F&\leq e.                              \tag{82.2}
\end{split}
\]
All matrices may have entries in \(\mathbb Q(i)\).

\begin{theorem}[Rational cell certificate]
\label{thm:n1v82-cell}
Set
\[
 \eta=e+rh,\qquad d={\eta\over\lambda},                     \tag{82.3}
\]
and let rational \(L,M,b\) satisfy
\[
\begin{split}
 L&\geq\|R\widetilde Y\|_F+ry,\\
 M&\geq\|\widetilde YR\|_F+ry,\\
 b&\geq\|\widetilde Y\|_F+y.                                \tag{82.4}
\end{split}
\]
Then, uniformly on \(J\),
\[
\boxed{
 \|H(t)^{-1/2}Y(t)H(t)^{-1/2}\|_F
 \leq
 \left[LM+b^2(2dr+d^2)\right]^{1/2}.}                       \tag{82.5}
\]
Every square root in (82.4)--(82.5) may be replaced by a rational
upper ceiling.  No smallness condition on \(\eta\) is required.
\end{theorem}

\begin{proof}
From (82.2),
\[
 \|I-RH(t)\|_{\rm op}
 \leq\|I-R\widetilde H\|_F+
      \|R\|_{\rm op}\|H(t)-\widetilde H\|_F
 \leq\eta.
\]
Writing \(G(t)=H(t)^{-1}\), the exact identity
\[
 G(t)-R=(I-RH(t))G(t)
\]
and (82.1) give
\[
 \|G(t)-R\|_{\rm op}\leq d.                                \tag{82.6}
\]
Put \(D=G-R\).  Hilbert--Schmidt duality gives
\[
\begin{split}
\operatorname{Tr}(Y^*GYG)
&\leq
 \|RY\|_F\|YR\|_F
 \|Y\|_F^2(2\|D\|_{\rm op}\|R\|_{\rm op}
                  +\|D\|_{\rm op}^2).
\end{split}                                                  \tag{82.7}
\]
The perturbation estimates
\[
 \|RY\|_F\leq\|R\widetilde Y\|_F+ry,
 \quad
 \|YR\|_F\leq\|\widetilde YR\|_F+ry,
 \quad
 \|Y\|_F\leq\|\widetilde Y\|_F+y
\]
and (82.6) yield (82.5).  Exact sums of squares of
Gaussian-rational entries give rational squared Frobenius norms;
outward rational square-root ceilings complete the certificate.
\end{proof}

For \(Y=\dot H\), the right side of (82.5) bounds the operator norm
\(a(t)\) in (80.3), because \(\|\cdot\|_{\rm op}\leq\|\cdot\|_F\).
For \(Y=\dot X\), it bounds \(c(t)\) directly.  The same rational
inverse \(R\) is reused for all nine vertex jets on the cell.

\subsection{Rational errors for the short frequency banks}

Let one of the V81 grouped paths be
\[
 X(t)=\sum_{|n|\leq N}B_ne^{in\pi t/2},\qquad N\leq4.        \tag{82.8}
\]
Fix a cell center \(\tau\) and half-width \(\delta\).  Suppose a
Gaussian-rational matrix \(\widetilde X_j\) satisfies
\[
 \|X^{(j)}(\tau)-\widetilde X_j\|_F\leq e_{j,0}.             \tag{82.9}
\]
Choose a rational \(p>\pi/2\), and rational Frobenius ceilings
\(\beta_n\geq\|B_n\|_F\).  Define
\[
 M_j=p^j\sum_{|n|\leq N}|n|^j\beta_n.                       \tag{82.10}
\]

\begin{lemma}[Center-to-cell trigonometric remainder]
\label{lem:n1v82-trig-cell}
For every \(t\) in the cell,
\[
 \boxed{
 \|X^{(j)}(t)-\widetilde X_j\|_F
 \leq e_{j,0}+\delta M_{j+1}.}                              \tag{82.11}
\]
\end{lemma}

\begin{proof}
Termwise differentiation of (82.8) gives
\[
 X^{(j+1)}(t)=
 \sum_n(in\pi/2)^{j+1}B_ne^{in\pi t/2}.
\]
The triangle inequality and \(p>\pi/2\) give the uniform bound
\(\|X^{(j+1)}(t)\|_F\leq M_{j+1}\).  Integrate between
\(\tau\) and \(t\), then add (82.9).
\end{proof}

The center error \(e_{j,0}\) itself can be rationally certified.
Reduce each rational multiple of \(\pi\) modulo \(\pi/2\), so the
remaining argument has magnitude at most \(\pi/4\).  Use rational
bounds for \(\pi\), the finite Taylor polynomials for sine and
cosine, and the Lagrange remainder
\[
 |R_m(x)|\leq {|x|^{m+1}\over(m+1)!}.                       \tag{82.12}
\]
Interval arithmetic in \(\mathbb Q\) then encloses every phase in
(82.8).  Multiplying those scalar intervals by the exact \(B_n\)
and summing entrywise supplies \(\widetilde X_j\) and
\(e_{j,0}\).  Thus (82.11) requires no trusted floating
transcendental value.

\subsection{Exact propagation through finitely many cells}

\begin{lemma}[Rational exponential ceiling]
\label{lem:n1v82-exp}
Let \(x\geq0\) be rational.  Choose an integer \(m\) with
\(m+2>x\), put
\[
 S_m(x)=\sum_{\ell=0}^m{x^\ell\over\ell!},\qquad
 q={x\over m+2},
\]
and define
\[
 E_m(x)=S_m(x)+
 {x^{m+1}\over(m+1)!}\,{1\over1-q}.                         \tag{82.13}
\]
Then \(E_m(x)\in\mathbb Q\) and \(e^x\leq E_m(x)\).
\end{lemma}

\begin{proof}
After the first omitted exponential-series term, every successive
term is at most \(q\) times its predecessor.  The omitted positive
tail is therefore bounded by the displayed geometric series.
\end{proof}

\begin{proposition}[Fully rational path propagation]
\label{prop:n1v82-rational-propagation}
Partition a path into cells of lengths \(\Delta_j\).  Suppose
Theorem~\ref{thm:n1v82-cell} supplies rational constants
\[
 a(t)\leq a_j,\qquad c(t)\leq c_j
\]
on cell \(j\), and let \(E_j\) be any rational ceiling for
\(\exp(a_j\Delta_j)\), for example (82.13).  Starting from a
rational bound \(F_0\geq f(0)\), recursively set
\[
 F_{j+1}=
 \begin{cases}
 E_jF_j+{c_j\over a_j}(E_j-1),&a_j>0,\\
 F_j+c_j\Delta_j,&a_j=0.
 \end{cases}                                                 \tag{82.14}
\]
Then \(f\) at every cell endpoint is bounded by the corresponding
\(F_j\).
\end{proposition}

\begin{proof}
Apply (80.6) on one cell with the incoming value bounded by
\(F_j\).  Replacing the exact exponential by its upper ceiling
can only increase the right side.  Induction proves the claim.
\end{proof}

Equations (82.2)--(82.14) give a complete proof interface for the
V81 route: all matrices, residuals, remainders, norm ceilings, and
exponential propagators can be stored as rational integers and
Fractions.  The construction may be conservative, but it cannot
silently inherit a floating eigenvalue or phase.

\begin{decision}[V83 populates the thirty-two-cell diagonal cards]
\label{dec:n1v82-next}
Implement (82.2)--(82.14) for the sixteen signed diagonals of each
action using \(32\) equal cells, degree sufficient to make the
Taylor remainder smaller than \(10^{-14}\), and
\(333/106<\pi<355/113\).  Record the worst inverse residual,
the integrated energy length, and the nine transported endpoint
bounds.  If a cell estimate is too large, subdivide that cell
only; do not replace its complete grouped bank by the V79
coefficient slopes.
\end{decision}

\begin{assessment}[Progress after compact V82]
\label{ass:n1v82-status}
The remaining diagonal-route computation now has a fully rational
certificate format.  In particular, neither matrix square roots
nor trusted floating trigonometric values are required, and exact
Grönwall propagation is available.  No path card has yet been
populated.
\end{assessment}

\begin{firewall}[Claim boundary after compact V82]
\label{fire:n1v82-claim}
The path-cell and propagation theorems are rigorous but presently
unpopulated for the Yang--Mills banks.  The V81 route values remain
diagnostic.  No global dressed stock, improved quadrature sign
enclosure, higher-loop or nonperturbative estimate, SCIC, ADMC,
continuum measure, reflection positivity, exponential clustering,
or Yang--Mills mass gap is proved.
\end{firewall}

\subsection{Parent record}

This V82 note was formed by appending the preceding material to the
validated compact V81, whose record was
\[
\begin{array}{c|c}
\text{lines}&27855\\
\text{bytes}&949968\\
\text{SHA--256}&
\texttt{D59C5389D96AFFC4A1A69F2B87CF75849F06748AE0B375286C712B613D7600E1}.
\end{array}
\]

% =============================================================================
% END APPENDED COMPACT-SERIES V82 MATERIAL
% =============================================================================

% =============================================================================
% BEGIN APPENDED COMPACT-SERIES V83 MATERIAL
% =============================================================================

\section{Compact-series V83: the first populated rational
diagonal card}
\label{sec:n1v83}

\subsection{Scope of the first card}

The V81 diagnostic selected the signed diagonal from the all-minus
corner to
\[
 k_c=(\pi/2,\pi/2,\pi/2,\pi/2).                             \tag{83.1}
\]
V83 populates the V82 certificate on this path for both actions,
using \(32\) equal cells and every external order in the quartic,
plus-cubic, and minus-cubic banks.  This is the smallest card that
simultaneously tests the rational phase enclosure, a posteriori
inverse residual, dressed path speeds, and exact exponential
propagation.

The trigonometric input is self-contained.  Machin's identity
\[
 {\pi\over4}=4\arctan{1\over5}-\arctan{1\over239}            \tag{83.2}
\]
and alternating rational partial sums give an interval strictly
inside
\[
 {333\over106}<\pi<{355\over113}.                            \tag{83.3}
\]
After reduction modulo \(\pi/2\), degree-\(25\) sine and
degree-\(24\) cosine Taylor polynomials, with Lagrange remainders,
enclose every midpoint phase.  Phase centers are rounded to
denominator \(10^{15}\), and the rounding error is added
rationally.

\begin{theorem}[Exact all-minus 32-cell card]
\label{thm:n1v83-card}
The V82 construction gives valid rational piecewise bounds on
\(a(t)\) and every \(c_X(t)\) along (83.1).  Its principal
summaries are
\[
\begin{array}{c|cc}
&{\rm Wilson}&\text{endpoint}\\ \hline
\max_j\eta_j&
5.907118179&10.766374421\\
\max_j a_j&
1579.535154&14416.534325\\
\sum_j a_j/32&
1488.520026&13691.830834 .
\end{array}                                                   \tag{83.4}
\]
The first row is the complete allowance
\[
 \eta_j=\|I-R_j\widetilde H_j\|_F+
 \|R_j\|_F h_j,                                             \tag{83.5}
\]
not a floating inverse residual.  No hypothesis
\(\eta_j<1\) is used.

The maximal certified vertex speeds by external order are
\[
\begin{array}{c|c|ccc}
e&\text{bank}&0&1&2\\ \hline
{\rm W}&c_{P^\pm}&1978.560&1122.595&1866.495\\
{\rm W}&c_Q&3957.120&0&4725.881\\
*&c_{P^\pm}&5755.204&7850.134&12597.341\\
*&c_Q&34673.137&0&119617.440 .
\end{array}                                                   \tag{83.6}
\]
All displayed decimals in (83.4)--(83.6) are upward rational
roundings.
\end{theorem}

\begin{proof}
At every cell midpoint, the exact signed-frequency matrices from
Theorem~\ref{thm:n1v81-frequency-collapse} are evaluated with the
rational phase rectangles described above.  Lemma
\ref{lem:n1v82-trig-cell} extends those rectangles across the
cell.  A Hermitian rational trial inverse of denominator \(10^8\)
is selected; all residual matrix products and Frobenius
square-root ceilings are then computed with integers and
Fractions.  Theorem~\ref{thm:n1v82-cell}, first with
\(Y=\dot H\) and then with each \(Y=\dot X\), proves the
piecewise bounds.  Summing the \(a_j/32\) and maximizing the
remaining Fractions gives (83.4)--(83.6).
\end{proof}

\subsection{A rigorous no-gain result}

Exact rational exponential propagation by
Proposition~\ref{prop:n1v82-rational-propagation} gives only
\[
\begin{array}{c|ccc}
e&\widehat P^\pm_0&\widehat P^\pm_1&\widehat P^\pm_2\\ \hline
{\rm W}&10^{677}&10^{677}&10^{678}\\
*&10^{5991}&10^{5991}&10^{5991},
\end{array}                                                   \tag{83.7}
\]
and
\[
\begin{array}{c|ccc}
e&\widehat Q_0&\widehat Q_1&\widehat Q_2\\ \hline
{\rm W}&10^{678}&0&10^{678}\\
*&10^{5992}&0&10^{5993}.
\end{array}                                                   \tag{83.8}
\]
The powers of ten are rigorous upper ceilings.  They are not
estimates of the actual norms.  Comparing (83.4) with the floating
path lengths \(3.1063\) and \(3.3629\) in V81 proves that the
implemented enclosure architecture, not the geometric route, is
responsible for the blow-up.

\begin{proposition}[Location of the first scalarization loss]
\label{prop:n1v83-loss}
Within the V83 implementation, the first avoidable whole-cell
relaxation is
\[
\begin{split}
\|I-RH(t)\|_F
&\leq\|I-R\widetilde H\|_F+
      \|R\|_F\|H(t)-\widetilde H\|_F,\\
\|RY(t)\|_F
&\leq\|R\widetilde Y\|_F+
      \|R\|_F\|Y(t)-\widetilde Y\|_F.                        \tag{83.9}
\end{split}
\]
Both inequalities discard entrywise cancellations in the
completed products.  The same loss enters \(YR\).
\end{proposition}

\begin{proof}
The phase grouping and center evaluation occur before (83.9), so
they preserve the signed-frequency cancellations.  The next
operation in the verifier replaces the entire interval remainder
matrix by one Frobenius radius and multiplies it by
\(\|R\|_F\).  Formula (83.5) shows the effect directly: the
resulting allowances exceed one and, after division by the small
global floor in (82.3), dominate (82.5).  No earlier step in the
implemented dependency graph uses a comparable scalar product.
\end{proof}

\begin{record}[V83 rational path-card verifier]
\label{rec:n1v83-executable}
The executable is
\[
 \texttt{scripts/verify\_ym\_v83\_central\_path\_cards.py}.  \tag{83.10}
\]
Its source record is
\[
\begin{array}{c|c}
\text{lines}&441\\
\text{bytes}&16373\\
\text{SHA--256}&
\texttt{217428080F228D71C3C0F61C06D4DE4B061C6AD11A078A7D635024F50D19C9EC}.
\end{array}                                                   \tag{83.11}
\]
It checks the Machin interval, all rational phase and Taylor
remainders, all cell residual products, and nonnegativity of every
cell quantity.  The canonical payload SHA--256 is
\[
 \boxed{
\texttt{4E15FCD258BBCD0B39523D5593E03C3334BBA85AAEDFA60894145B74533726A0}.}
                                                               \tag{83.12}
\]
\end{record}

\begin{decision}[V84 multiplies interval matrices before norming]
\label{dec:n1v83-next}
Replace (83.9) by entrywise complex rational interval evaluation
of the complete products
\[
 I-RH(t),\qquad R\dot H(t),\qquad\dot H(t)R,\qquad
 R\dot X(t),\qquad\dot X(t)R.                               \tag{83.13}
\]
Only after summing each matrix entry over all frequencies and
matrix-product indices should its absolute interval radius be
converted to a Frobenius bound.  Reuse the same \(32\) cells and
the all-minus route so the change in (83.4) is attributable only
to this contraction order.  Do not populate the remaining signed
diagonals unless the energy length falls to a usable scale.
\end{decision}

\begin{assessment}[Progress after compact V83]
\label{ass:n1v83-status}
The first diagonal path card is fully rational and reproducible.
It is far too weak to use, but it localizes the failure to
premature scalarization of complete interval matrix products.
This is actionable upstream evidence and prevents an expensive
replication of a defective card over the remaining paths.
\end{assessment}

\begin{firewall}[Claim boundary after compact V83]
\label{fire:n1v83-claim}
The values in (83.4)--(83.8) are rigorous upper bounds but provide
no improvement; V69 remains the operative global response
majorant.  No global dressed stock, quadrature sign enclosure,
higher-loop or nonperturbative estimate, SCIC, ADMC, continuum
measure, reflection positivity, exponential clustering, or
Yang--Mills mass gap is proved.
\end{firewall}

\subsection{Parent record}

This V83 note was formed by appending the preceding material to the
validated compact V82, whose record was
\[
\begin{array}{c|c}
\text{lines}&28114\\
\text{bytes}&958424\\
\text{SHA--256}&
\texttt{D4E3B41C26CF2E1CB05559F7449AE5A00A22955E7A50FB59BC69B80EA9858E0A}.
\end{array}
\]

% =============================================================================
% END APPENDED COMPACT-SERIES V83 MATERIAL
% =============================================================================

% =============================================================================
% BEGIN APPENDED COMPACT-SERIES V84 MATERIAL
% =============================================================================

\section{Compact-series V84: complete interval products before
norming}
\label{sec:n1v84}

\subsection{Entrywise complex interval contraction}

V83 identified (83.9) as the first avoidable scalarization.  V84
uses the same path, the same \(32\) cells, the same phase
rectangles, and the same rational trial inverses, but changes the
contraction order only.

For a complex interval matrix write
\[
 {\boldsymbol Y}_{ij}=Y_{ij}^{\,c}
 +\{z\in\mathbb C:|z|\leq\rho_{ij}\}.                       \tag{84.1}
\]
If \(R\) is point-valued, form the complete center products and
entry radii
\[
\begin{split}
 (R{\boldsymbol Y})^c&=RY^c,&
 \rho^{\,L}_{ij}&=\sum_k|R_{ik}|\rho_{kj},\\
 ({\boldsymbol Y}R)^c&=Y^cR,&
 \rho^{\,R}_{ij}&=\sum_k\rho_{ik}|R_{kj}|.                  \tag{84.2}
\end{split}
\]

\begin{lemma}[Frobenius enclosure after product contraction]
\label{lem:n1v84-interval-product}
Every \(Y\in{\boldsymbol Y}\) satisfies
\[
\begin{split}
 \|RY\|_F&\leq\|RY^c\|_F+\|\rho^{\,L}\|_F,\\
 \|YR\|_F&\leq\|Y^cR\|_F+\|\rho^{\,R}\|_F,\\
 \|Y\|_F&\leq\|Y^c\|_F+\|\rho\|_F.                          \tag{84.3}
\end{split}
\]
For a Hermitian \(R\),
\[
 \|R\|_{\rm op}\leq
 \max_i\sum_j|R_{ij}|.                                      \tag{84.4}
\]
All right sides have rational upper ceilings when the centers and
radii are Gaussian-rational.
\end{lemma}

\begin{proof}
Equation (84.2) follows entrywise from the triangle inequality
after the matrix-product index has been summed.  Apply the
Frobenius triangle inequality to the center and radius matrices.
For (84.4), Hermiticity gives
\(\|R\|_1=\|R\|_\infty\), and
\(\|R\|_{\rm op}\leq\sqrt{\|R\|_1\|R\|_\infty}\).
\end{proof}

The residual is treated in the same order:
\[
 \|I-RH(t)\|_F
 \leq\|I-RH^c\|_F+\||R|\rho_H\|_F.                          \tag{84.5}
\]
Insert (84.3)--(84.5) into (82.6)--(82.7).  Thus every frequency
and product-index cancellation in the center is retained before
the first matrix norm.

\subsection{Controlled comparison with V83}

\begin{theorem}[Exact V84 all-minus interval-product card]
\label{thm:n1v84-card}
The complete interval-product construction gives
\[
\begin{array}{c|cc}
&{\rm Wilson}&\text{endpoint}\\ \hline
\max_j\|I-R_jH\|_F&
1.270946688&3.473091304\\
\max_j a_j&
547.434813&5061.573162\\
\sum_j a_j/32&
522.220327&4630.269583 .
\end{array}                                                   \tag{84.6}
\]
All decimals are upward rational ceilings.  The maximal vertex
speeds are
\[
\begin{array}{c|c|ccc}
e&\text{bank}&0&1&2\\ \hline
{\rm W}&c_{P^\pm}&686.792&389.623&648.150\\
{\rm W}&c_Q&1373.584&0&1641.226\\
*&c_{P^\pm}&1940.599&2490.849&4136.304\\
*&c_Q&11691.049&0&39823.319 .
\end{array}                                                   \tag{84.7}
\]
\end{theorem}

\begin{proof}
The V83 rational phase rectangles are used to form the center and
radius matrices (84.1).  All products in (84.2) and (84.5) are
computed with common-denominator integer matrices.  Exact
Frobenius ceilings and the row-sum bound (84.4) are substituted
into Theorem~\ref{thm:n1v82-cell}.  Maximization and summation over
the \(32\) cells give (84.6)--(84.7).
\end{proof}

Relative to V83, the integrated energy bounds improve by factors
\[
 {1488.520026\over522.220327}>2.85,\qquad
 {13691.830834\over4630.269583}>2.95.                        \tag{84.8}
\]
Exact propagation reduces the Wilson nonzero final ceilings from
orders \(10^{677}\)--\(10^{678}\) to values of order
\(10^{228}\)--\(10^{229}\).  The endpoint ceilings fall from
orders \(10^{5991}\)--\(10^{5993}\) to
\(10^{2048}\)--\(10^{2049}\).  This is a substantial verified
improvement, but it is still unusable.

\begin{record}[V84 complete-product verifier]
\label{rec:n1v84-executable}
The executable is
\[
 \texttt{scripts/verify\_ym\_v84\_interval\_product\_cards.py}. \tag{84.9}
\]
Its source record is
\[
\begin{array}{c|c}
\text{lines}&313\\
\text{bytes}&11552\\
\text{SHA--256}&
\texttt{2B5C0C54559159ABB42AA07DD0E7A3249D4A6338E96FB6B9FD6170D83C08B95C}.
\end{array}                                                   \tag{84.10}
\]
It checks both actions on all \(32\) cells and every completed
interval product.  Its canonical payload SHA--256 is
\[
 \boxed{
\texttt{BA6C68DEAC5D2552A53C0258519941A4F4E6A2CF0DA34EA243773D51D3836022}.}
                                                               \tag{84.11}
\]
\end{record}

\begin{decision}[V85 audits companions and measures cell-width
convergence]
\label{dec:n1v84-next}
Perform the mandatory current SM/NS audit.  Then compute
Hessian-only complete-product cards at \(32,64,128,\ldots\)
cells, without the nine vertex families, until the asymptotic
scale of the certified energy length is clear.  If subdivision
does not approach the V81 diagnostic scale, replace rectangular
phase radii by correlated trigonometric-polynomial enclosures
before performing any further vertex propagation.
\end{decision}

\begin{assessment}[Progress after compact V84]
\label{ass:n1v84-status}
Norming after complete interval multiplication removes roughly a
factor three in energy length and thousands of powers of ten from
the endpoint propagation.  The remaining gap to the path
diagnostic is nevertheless two to three orders in energy length.
Cell-width convergence must be measured before this architecture
is expanded.
\end{assessment}

\begin{firewall}[Claim boundary after compact V84]
\label{fire:n1v84-claim}
V84 is a rigorous improvement of a non-operative path bound.
V69 remains the operative global response majorant.  No global
dressed stock, quadrature sign enclosure, higher-loop or
nonperturbative estimate, SCIC, ADMC, continuum measure,
reflection positivity, exponential clustering, or Yang--Mills
mass gap is proved.
\end{firewall}

\subsection{Parent record}

This V84 note was formed by appending the preceding material to the
validated compact V83, whose record was
\[
\begin{array}{c|c}
\text{lines}&28334\\
\text{bytes}&966306\\
\text{SHA--256}&
\texttt{8826BBDBAA4F51BDE0E3F2DE44E7E042BBD08D20CA8E7BBA5E536F4A60162B59}.
\end{array}
\]

% =============================================================================
% END APPENDED COMPACT-SERIES V84 MATERIAL
% =============================================================================

% =============================================================================
% BEGIN APPENDED COMPACT-SERIES V85 MATERIAL
% =============================================================================

\section{Compact-series V85: companion audit and residual
subdivision law}
\label{sec:n1v85}

\subsection{Mandatory current companion audit}

The current companion files were inspected read-only.  Since V80,
the SM programme has advanced from V29 to V43, while the NS file is
unchanged:
\[
\begin{array}{c|c|r|r|c}
\text{programme}&\text{version}&\text{lines}&\text{bytes}&
\text{SHA--256}\\ \hline
{\rm SM\!-\!Einstein}&V43&14249&504099&
\texttt{ECDB9BE0EB74E544BB4448CFC0E05DBA467708BC8268C89898B28D7B1EABDE0B}\\
{\rm Navier\!-\!Stokes}&V26&5319&278283&
\texttt{82C887F8C2C69E4F28FAD7C3DC8DE5B9099CB5A4BF431EBA1FFF3E91935978AD}.
\end{array}                                                   \tag{85.1}
\]

SM V43 closes a finite algebraic occurrence short using paired
unitary CPTP channels.  Their central difference converges to the
Dirac--Yukawa commutator, and a Duhamel comparison is proposed for
cofinal transport.  No uniform coefficient bound, clock
calibration, or continuum gap is proved.  Its use of actual finite
maps before differentiation is structurally compatible with the
present rule of completing matrix products before norming, but it
supplies no YM momentum bound.

NS V26 remains the rank-one Oseen-kernel refinement audited in
V80.  Its lesson has already been incorporated in V81 and V84:
contract the prescribed structured input before applying an
ambient norm.  No NS constant enters the YM calculation.

\begin{proposition}[V85 audit decision]
\label{prop:n1v85-audit}
No theorem or numerical constant is imported from either companion
programme.  V84's exact interval-product architecture is retained
for the subdivision study, and the SM channel compactness and
Duhamel machinery remain later-stage interfaces only.
\end{proposition}

\begin{proof}
The current YM short is a finite \(45\)-dimensional inverse and
path-norm estimate.  Neither companion file proves a bound for its
Hessian, signed-frequency banks, or trial residuals.  The stated
separation follows.
\end{proof}

\subsection{Hessian-only exact subdivision}

To avoid repeating the expensive nine-family vertex propagation,
the all-minus path was certified using only the complete residual
and \(Y=\dot H\) at \(32,64,128\) equal cells.

\begin{theorem}[Monotone residual subdivision cards]
\label{thm:n1v85-subdivision}
The exact rational cards give
\[
\begin{array}{c|c|ccc}
e&N&\max\eta_j&\max a_j&\sum_j a_j/N\\ \hline
{\rm W}&32&1.270946688&547.434813&522.220327\\
{\rm W}&64&0.635484273&356.677038&339.821626\\
{\rm W}&128&0.317742489&241.783048&230.218229\\
*&32&3.473091304&5061.573162&4630.269583\\
*&64&1.738273209&2777.768205&2537.354300\\
*&128&0.869573094&1621.277508&1480.418052 .
\end{array}                                                   \tag{85.2}
\]
Every displayed value is an upward rational ceiling.  For each
action, the certified integrated energy length decreases strictly
under both cell doublings.
\end{theorem}

\begin{proof}
For each \(N\), use the complete interval residual (84.5), the
row-sum inverse bound (84.4), and the product bounds (84.3) on all
\(N\) cells.  The verifier compares the resulting exact Fractions
before decimal rounding and checks both strict decreases.
\end{proof}

An independent \(1024\)-midpoint floating calculation gives
\[
\begin{array}{c|cc}
e&\int\|H^{-1/2}\dot HH^{-1/2}\|_{\rm op}&
\int\|H^{-1/2}\dot HH^{-1/2}\|_F\\ \hline
{\rm W}&3.106529&8.626265\\
*&3.363728&9.759721.
\end{array}                                                   \tag{85.3}
\]
This is diagnostic only.  It shows that even the limiting
Frobenius architecture should live near ten, whereas (85.2)
remains hundreds or thousands.  Blind subdivision would therefore
be a poor next step.

\subsection{The residual threshold changes the inverse argument}

The last line for each action in (85.2) has
\[
 \eta_j<1\quad\hbox{on every cell}.                           \tag{85.4}
\]
V82 did not use this fact: it always converted residual error to
\[
 \|G-R\|_{\rm op}\leq{\eta\over\lambda_e},                  \tag{85.5}
\]
which imports the small global floor.  Once (85.4) holds, the
factorization itself gives a different route.  If
\[
 E_L=I-RH,\qquad E_R=I-HR,
\]
then
\[
 G=(I-E_L)^{-1}R=R(I-E_R)^{-1}.                             \tag{85.6}
\]
Consequently the complete products \(RY\) and \(YR\) can be
amplified by Banach factors \((1-\|E_L\|)^{-1}\) and
\((1-\|E_R\|)^{-1}\), without ever forming \(G-R\) or dividing
by \(\lambda_e\).

\begin{record}[V85 subdivision verifier]
\label{rec:n1v85-executable}
The executable is
\[
 \texttt{scripts/verify\_ym\_v85\_hessian\_subdivision.py}.  \tag{85.7}
\]
Its source record is
\[
\begin{array}{c|c}
\text{lines}&116\\
\text{bytes}&4079\\
\text{SHA--256}&
\texttt{1A19F5D9D0E114C997CF43B5BEEA467ECC6F0C820EFFE0360A1E0FB46916D18B}.
\end{array}                                                   \tag{85.8}
\]
Its canonical payload SHA--256 is
\[
 \boxed{
\texttt{3BE76C2A9F30C6BB84CE1334FC52819A1B7406CFE2AF32BDE3D1790A28329C23}.}
                                                               \tag{85.9}
\]
\end{record}

\begin{decision}[V86 uses two-sided Banach residuals]
\label{dec:n1v85-next}
Prove the two-sided residual dressed-product lemma suggested by
(85.6).  Populate it first at \(N=128\), where both actions satisfy
(85.4), and certify
\[
 \|H^{-1/2}YH^{-1/2}\|_F^2
 \leq
 {\|RY\|_F\|YR\|_F
  \over(1-\eta_L)(1-\eta_R)}.                               \tag{85.10}
\]
Use the complete interval products of V84 for all quantities.
Only if the endpoint Banach factors remain too large should the
cell count be doubled.
\end{decision}

\begin{assessment}[Progress after compact V85]
\label{ass:n1v85-status}
The companion audit is complete, exact cell refinement is
monotone, and the \(128\)-cell cards cross the residual-one
threshold for both actions.  This exposes a floor-free inverse
factorization that was unavailable at \(32\) endpoint cells and
should remove the dominant V82 loss.
\end{assessment}

\begin{firewall}[Claim boundary after compact V85]
\label{fire:n1v85-claim}
The two-sided estimate (85.10) is the next target, not yet a
populated theorem.  Values (85.3) are diagnostics.  V69 remains
the operative global response majorant.  No global dressed stock,
quadrature sign enclosure, higher-loop or nonperturbative
estimate, SCIC, ADMC, continuum measure, reflection positivity,
exponential clustering, or Yang--Mills mass gap is proved.
\end{firewall}

\subsection{Parent record}

This V85 note was formed by appending the preceding material to the
validated compact V84, whose record was
\[
\begin{array}{c|c}
\text{lines}&28525\\
\text{bytes}&972768\\
\text{SHA--256}&
\texttt{5248AFF4B520E512C787654EA4B576A9A3A1A5065641FDDD3E7DDCFA23823525}.
\end{array}
\]

% =============================================================================
% END APPENDED COMPACT-SERIES V85 MATERIAL
% =============================================================================

% =============================================================================
% BEGIN APPENDED COMPACT-SERIES V86 MATERIAL
% =============================================================================

\section{Compact-series V86: floor-free two-sided residual
transport}
\label{sec:n1v86}

\subsection{The Banach factorization}

Let \(H\succ0\), \(G=H^{-1}\), and let \(R\) be a trial inverse.
Define
\[
 E_L=I-RH,\qquad E_R=I-HR.                                  \tag{86.1}
\]

\begin{theorem}[Two-sided residual dressed bound]
\label{thm:n1v86-bidual}
Suppose
\[
 \|E_L\|_{\rm op}\leq\eta_L<1,\qquad
 \|E_R\|_{\rm op}\leq\eta_R<1.                              \tag{86.2}
\]
Then for every complex matrix \(Y\),
\[
\boxed{
 \|H^{-1/2}YH^{-1/2}\|_F^2
 \leq
 {\|RY\|_F\|YR\|_F
  \over(1-\eta_L)(1-\eta_R)}.}                              \tag{86.3}
\]
No global lower spectral bound for \(H\) occurs in (86.3).
\end{theorem}

\begin{proof}
The identities \(RH=I-E_L\) and \(HR=I-E_R\) imply
\[
 G=(I-E_L)^{-1}R=R(I-E_R)^{-1}.                             \tag{86.4}
\]
Consequently
\[
\begin{split}
 \|GY\|_F&\leq{\|RY\|_F\over1-\eta_L},\\
 \|YG\|_F&\leq{\|YR\|_F\over1-\eta_R}.                       \tag{86.5}
\end{split}
\]
Since
\[
\begin{split}
\|H^{-1/2}YH^{-1/2}\|_F^2
&=\operatorname{Tr}(Y^*GYG)\\
&=\operatorname{Tr}\bigl((GY)^*(YG)\bigr),
\end{split}
\]
Hilbert--Schmidt duality and (86.5) prove (86.3).
\end{proof}

For Hermitian interval Hessians and Hermitian \(R\), the completed
left and right residual bounds produced from (84.2) are equal in
the present cards, but the proof and verifier retain both sides
separately.

\subsection{Exact \(128\)- and \(256\)-cell cards}

\begin{theorem}[Populated floor-free Hessian cards]
\label{thm:n1v86-cards}
On the all-minus signed diagonal, complete rational interval
products and Theorem~\ref{thm:n1v86-bidual} give
\[
\begin{array}{c|c|c|c|c}
e&N&\max\eta_L=\max\eta_R&\max a_j&
\sum_j a_j/N\\ \hline
{\rm W}&128&0.317742489&43.537277&42.058237\\
{\rm W}&256&0.158871502&35.022769&34.171961\\
*&128&0.869573094&288.904204&195.463451\\
*&256&0.434900678&65.411940&54.878727 .
\end{array}                                                   \tag{86.6}
\]
All entries are upward rational ceilings.  Every residual is
strictly below one, and each action's integrated bound decreases
strictly under cell doubling.
\end{theorem}

\begin{proof}
The V84 entrywise interval products give exact Frobenius upper
bounds for both residuals and for \(R\dot H,\dot HR\).  Frobenius
residual bounds imply the operator hypotheses in (86.2).  Apply
(86.3) on each cell, take the rational square-root ceiling, and
sum the resulting \(a_j/N\).  The verifier performs all comparisons
as Fractions.
\end{proof}

At \(128\) cells, replacing (85.5) by (86.3) improves the
integrated bounds by factors greater than
\[
 {230.218229\over42.058237}>5.47,\qquad
 {1480.418052\over195.463451}>7.57.                          \tag{86.7}
\]
At \(256\) cells the bounds are within factors \(3.97\) and \(5.63\)
of the floating dressed-Frobenius diagnostics in (85.3).  This is
the first exact path architecture in the compact series whose
energy length is below \(55\) for both actions.

\subsection{The limiting Hilbert--Schmidt loss}

The zero-cell-width limit of (86.3), with \(R=G\), is
\[
 \left(\|G\dot H\|_F\|\dot HG\|_F\right)^{1/2}.              \tag{86.8}
\]
A \(1024\)-midpoint diagnostic gives integrated values
\[
 28.703523\quad({\rm Wilson}),\qquad
 31.403756\quad(\text{endpoint}),                            \tag{86.9}
\]
whereas the actual dressed-Frobenius values are \(8.626265\) and
\(9.759721\).  Thus further subdivision of (86.3) cannot remove a
remaining factor about three.  Values (86.9) are diagnostic, but
the algebraic limiting statement (86.8) is exact.

\begin{record}[V86 two-sided residual verifier]
\label{rec:n1v86-executable}
The executable is
\[
 \texttt{scripts/verify\_ym\_v86\_two\_sided\_banach\_cards.py}. \tag{86.10}
\]
Its source record is
\[
\begin{array}{c|c}
\text{lines}&146\\
\text{bytes}&5217\\
\text{SHA--256}&
\texttt{1665D0C800F467819DDAF0C4F50546328BF636249CA5153EB3BA05FE11C281E6}.
\end{array}                                                   \tag{86.11}
\]
Its canonical payload SHA--256 is
\[
 \boxed{
\texttt{1547D91C56FD21AE9048DDB87DF38D157A3FA7EBC3530E3B61F8381C1C1129B9}.}
                                                               \tag{86.12}
\]
\end{record}

\begin{decision}[V87 retains the trial trace before residual
perturbation]
\label{dec:n1v86-next}
For Hermitian \(Y=\dot H\), compute the complete interval trace
\[
 s_R(Y)=\operatorname{Tr}(YRYR)                              \tag{86.13}
\]
before applying a norm.  Use
\[
 \|G-R\|_{\rm op}
 \leq{\eta_L\over1-\eta_L}\|R\|_{\rm op}                    \tag{86.14}
\]
and the V78 trace perturbation identity to enclose the difference
between \(s_R(Y)\) and
\(\operatorname{Tr}(YGYG)\).  Compare at \(N=256\) with (86.6).
This architecture tends to the actual dressed Frobenius norm,
not the larger limit (86.8).
\end{decision}

\begin{assessment}[Progress after compact V86]
\label{ass:n1v86-status}
Crossing the two-sided residual-one threshold removes the global
Hessian floor from path transport and cuts the exact energy cards
by factors five to eight.  The remaining asymptotic loss is now
the final Hilbert--Schmidt trace separation in (86.3), for which
V87 has a precise trace-retaining replacement.
\end{assessment}

\begin{firewall}[Claim boundary after compact V86]
\label{fire:n1v86-claim}
Only the Hessian path speed has been populated with the new
architecture.  No global dressed vertex stock or response
quadrature has been certified.  V69 remains the operative global
response majorant.  No higher-loop or nonperturbative estimate,
SCIC, ADMC, continuum measure, reflection positivity,
exponential clustering, or Yang--Mills mass gap is proved.
\end{firewall}

\subsection{Parent record}

This V86 note was formed by appending the preceding material to the
validated compact V85, whose record was
\[
\begin{array}{c|c}
\text{lines}&28723\\
\text{bytes}&979946\\
\text{SHA--256}&
\texttt{4B74C9FBA8C0292A8AA2FF2AB19C4E412E901D95A63D5252A80562CF0E31FB54}.
\end{array}
\]

% =============================================================================
% END APPENDED COMPACT-SERIES V86 MATERIAL
% =============================================================================

% =============================================================================
% BEGIN APPENDED COMPACT-SERIES V87 MATERIAL
% =============================================================================

\section{Compact-series V87: trial-trace residual comparison}
\label{sec:n1v87}

\subsection{Trace retention with an additive inverse error}

For Hermitian \(Y\) and \(R\), set
\[
 s_R(Y)=\operatorname{Tr}(YRYR).                             \tag{87.1}
\]
An interval enclosure of \(Y\) induces complete interval products
\({\boldsymbol A}=R{\boldsymbol Y}\) and
\({\boldsymbol B}={\boldsymbol Y}R\).  The trial trace is the
Hilbert--Schmidt pairing
\[
 s_R(Y)=\operatorname{Tr}(A^*B).                            \tag{87.2}
\]
If \(A^c,B^c\) are the centers and
\(\rho^A,\rho^B\) their entry radii, then
\[
\begin{split}
s_R(Y)\leq\operatorname{Re}\operatorname{Tr}((A^c)^*B^c)
+\sum_{ij}\bigl(
|A^c_{ij}|\rho^B_{ij}
+\rho^A_{ij}|B^c_{ij}|
+\rho^A_{ij}\rho^B_{ij}\bigr).                              \tag{87.3}
\end{split}
\]
This retains the complete center trace before applying absolute
values.

\begin{proposition}[Additive residual trace bound]
\label{prop:n1v87-additive}
Under the hypotheses of Theorem~\ref{thm:n1v86-bidual}, let
\[
 r\geq\|R\|_{\rm op},\qquad
 d={\eta_L\over1-\eta_L}r.                                  \tag{87.4}
\]
Then
\[
\|H^{-1/2}YH^{-1/2}\|_F^2
\leq s_R(Y)+\|Y\|_F^2(2dr+d^2).                             \tag{87.5}
\]
\end{proposition}

\begin{proof}
Equation (86.4) gives
\[
 \|G-R\|_{\rm op}
 \leq{\eta_L\over1-\eta_L}\|R\|_{\rm op}\leq d.
\]
Expand \(G=R+(G-R)\) in
\(\operatorname{Tr}(Y^*GYG)\), exactly as in the proof of
Lemma~\ref{lem:n1v78-residual-dressing}.  The trial term is
(87.1), and the two mixed and one quadratic errors give (87.5).
\end{proof}

\subsection{Exact comparison at 256 cells}

\begin{theorem}[The additive trace card is non-operative]
\label{thm:n1v87-comparison}
At \(N=256\), complete rational interval evaluation of
(87.3)--(87.5) gives
\[
\begin{array}{c|cc}
e&\text{additive trace energy}&\text{V86 Banach energy}\\ \hline
{\rm W}&223.910025&34.171961\\
*&840.469684&54.878727.
\end{array}                                                   \tag{87.6}
\]
The respective maximal cell bounds are
\[
 243.919502,\ 965.288752
\quad\hbox{versus}\quad
35.022769,\ 65.411940.                                      \tag{87.7}
\]
All numbers are upward rational ceilings.  Thus (87.5) is strictly
worse on this card and must not replace Theorem
\ref{thm:n1v86-bidual}.
\end{theorem}

\begin{proof}
The complete product centers and radii from V84 give (87.3).
The row-sum bound (84.4) and the two-sided residual intervals give
(87.4), while the interval matrix itself bounds \(\|Y\|_F\).
Summation over all \(256\) cells gives the first column.  The same
cell contexts evaluated with (86.3) give the second column, so the
comparison contains no mesh or input difference.
\end{proof}

This negative result does not refute trace retention.  It shows
that the additive perturbation in (87.5) is the wrong way to
combine it with a finite residual: the factor
\(\|Y\|_F^2r^2\) discards the same conditioning that energy
dressing was introduced to avoid.

\begin{record}[V87 trace-residual verifier]
\label{rec:n1v87-executable}
The executable is
\[
 \texttt{scripts/verify\_ym\_v87\_trace\_residual\_cards.py}. \tag{87.8}
\]
Its source record is
\[
\begin{array}{c|c}
\text{lines}&141\\
\text{bytes}&5431\\
\text{SHA--256}&
\texttt{C4C41169AC8F19E9758D755D2D0334FAA7B5C83D369476E021D19C5E58B28B87}.
\end{array}                                                   \tag{87.9}
\]
It checks all two-sided residual gates and verifies that every
trial-trace interval intersects the real axis.  Its canonical
payload SHA--256 is
\[
 \boxed{
\texttt{9E6ED816514CEC9D662608A5D4CEC42DB3B47FB9C6575FF8B4A3244BB22B50CC}.}
                                                               \tag{87.10}
\]
\end{record}

\subsection{A multiplicative trace correction}

Let
\[
 A=RY,\qquad B=YR.
\]
Equation (86.4) writes
\[
 GY=A+\Delta A,\qquad YG=B+\Delta B,
\]
with
\[
\|\Delta A\|_F\leq{\eta_L\over1-\eta_L}\|A\|_F,\qquad
\|\Delta B\|_F\leq{\eta_R\over1-\eta_R}\|B\|_F.              \tag{87.11}
\]
Expanding the Hilbert--Schmidt pairing therefore suggests the
sharper correction
\[
\begin{split}
\operatorname{Tr}((GY)^*(YG))
\leq s_R(Y)+\|A\|_F\|B\|_F
\left\{
{1\over(1-\eta_L)(1-\eta_R)}-1
\right\}.                                                    \tag{87.12}
\end{split}
\]
Unlike (87.5), (87.12) perturbs already energy-scaled products.

\begin{decision}[V88 populates the multiplicative trace card]
\label{dec:n1v87-next}
Prove (87.12) with interval upper bounds for \(s_R,A,B\), and
evaluate it on the same \(256\)-cell card.  Compare cellwise with
both (86.3) and (87.5), retaining the minimum certified value.
If the gain is material, use the winning hybrid for vertex speeds;
otherwise V86 remains the transport rule.
\end{decision}

\begin{assessment}[Progress after compact V87]
\label{ass:n1v87-status}
The first trace-retaining implementation has been tested and
rejected at fixed cell data.  Its failure is localized to an
additive unwhitened residual term.  A multiplicative correction
on the already formed products is derived as the next exact
alternative.
\end{assessment}

\begin{firewall}[Claim boundary after compact V87]
\label{fire:n1v87-claim}
The V87 card is rigorous but inferior to V86.  No vertex path
cards, global dressed stocks, response quadrature, higher-loop or
nonperturbative estimate, SCIC, ADMC, continuum measure,
reflection positivity, exponential clustering, or Yang--Mills
mass gap is proved.
\end{firewall}

\subsection{Parent record}

This V87 note was formed by appending the preceding material to the
validated compact V86, whose record was
\[
\begin{array}{c|c}
\text{lines}&28920\\
\text{bytes}&986388\\
\text{SHA--256}&
\texttt{F2C8D2B226997064E1CA581432E2AFB4FA8EF34C14BB97BB11A6F23D40694554}.
\end{array}
\]

% =============================================================================
% END APPENDED COMPACT-SERIES V87 MATERIAL
% =============================================================================

% =============================================================================
% BEGIN APPENDED COMPACT-SERIES V88 MATERIAL
% =============================================================================

\section{Compact-series V88: multiplicative trial-trace transport}
\label{sec:n1v88}

\subsection{Perturbing the energy-scaled products}

Retain the notation
\[
 A=RY,\qquad B=YR,\qquad
 E_L=I-RH,\qquad E_R=I-HR.                                  \tag{88.1}
\]
Under the residual-one hypotheses, (86.4) gives
\[
\begin{split}
 GY&=(I-E_L)^{-1}A=A+\Delta A,\\
 YG&=B(I-E_R)^{-1}=B+\Delta B.                              \tag{88.2}
\end{split}
\]

\begin{theorem}[Multiplicative trial-trace residual bound]
\label{thm:n1v88-multiplicative}
If \(\|E_L\|_{\rm op}\leq\eta_L<1\) and
\(\|E_R\|_{\rm op}\leq\eta_R<1\), then
\[
\boxed{
\begin{split}
\|H^{-1/2}YH^{-1/2}\|_F^2
\leq{}&\operatorname{Re}\operatorname{Tr}(A^*B)\\
&+\|A\|_F\|B\|_F
\left[
{1\over(1-\eta_L)(1-\eta_R)}-1
\right].
\end{split}}                                                 \tag{88.3}
\]
For interval \(H,Y\), any simultaneous upper enclosures of the
trace, product norms, and residuals give the same conclusion.
\end{theorem}

\begin{proof}
The Neumann bounds in (88.2) give
\[
\|\Delta A\|_F\leq
{\eta_L\over1-\eta_L}\|A\|_F,\qquad
\|\Delta B\|_F\leq
{\eta_R\over1-\eta_R}\|B\|_F.                               \tag{88.4}
\]
Now
\[
\begin{split}
\|H^{-1/2}YH^{-1/2}\|_F^2
&=\operatorname{Tr}((GY)^*(YG))\\
&=\operatorname{Tr}(A^*B)
 +\operatorname{Tr}(\Delta A^*B)
 +\operatorname{Tr}(A^*\Delta B)
 +\operatorname{Tr}(\Delta A^*\Delta B).
\end{split}
\]
Apply Hilbert--Schmidt duality to the last three terms.  Their
coefficient relative to \(\|A\|_F\|B\|_F\) is
\[
{\eta_L\over1-\eta_L}
+{\eta_R\over1-\eta_R}
+{\eta_L\eta_R\over(1-\eta_L)(1-\eta_R)}
=
{1\over(1-\eta_L)(1-\eta_R)}-1,
\]
which proves (88.3).
\end{proof}

Theorem~\ref{thm:n1v88-multiplicative} differs from V87 at the
decisive point: it never introduces
\(\|Y\|_F^2\|R\|_{\rm op}^2\).  It differs from V86 by retaining
the complete trial trace rather than replacing that trace by
\(\|A\|_F\|B\|_F\).

\subsection{Exact 256-cell hybrid card}

\begin{theorem}[The multiplicative trace wins every cell]
\label{thm:n1v88-card}
On the same all-minus \(256\)-cell Hessian card,
\[
\begin{array}{c|cc}
e&\text{multiplicative trace energy}&\text{V86 Banach energy}\\ \hline
{\rm W}&20.396984&34.171961\\
*&45.963686&54.878727.
\end{array}                                                   \tag{88.5}
\]
The maximal cell bounds are respectively
\[
 21.027490,\quad55.426747.                                  \tag{88.6}
\]
For both actions the multiplicative trace bound is strictly smaller
than the Banach bound on all \(256\) cells.  Hence the certified
cellwise hybrid minimum equals the multiplicative column in
(88.5).
\end{theorem}

\begin{proof}
The trial trace is enclosed by (87.3); the norms of \(A,B\) and
both residuals are the complete interval products of V84.  Insert
their exact rational ceilings into (88.3), take the rational
square-root ceiling, and compare with (86.3) before summation.
The verifier records \(256\) strict wins for each action.
\end{proof}

Relative to V86, the energy reductions are greater than
\[
 {34.171961\over20.396984}>1.67,\qquad
 {54.878727\over45.963686}>1.19.                            \tag{88.7}
\]
The current exact values are within factors \(2.37\) and \(4.71\)
of the dressed-Frobenius diagnostics in (85.3).  Unlike the
Banach architecture, (88.3) has the correct zero-width trial-trace
limit, so further cell refinement can in principle close these
remaining factors.

\begin{record}[V88 multiplicative-trace verifier]
\label{rec:n1v88-executable}
The executable is
\[
 \texttt{scripts/verify\_ym\_v88\_multiplicative\_trace\_cards.py}. \tag{88.8}
\]
Its source record is
\[
\begin{array}{c|c}
\text{lines}&140\\
\text{bytes}&5586\\
\text{SHA--256}&
\texttt{DC84017179F58078E3474D6F74789E4DB79482E793D3395495AA2F797B10396A}.
\end{array}                                                   \tag{88.9}
\]
It checks both residual-one gates and the cellwise hybrid
dominance.  Its canonical payload SHA--256 is
\[
 \boxed{
\texttt{953BC099AD274112638553F41DC9A2C5FFA237D578542E38EA6450D1C5212AA6}.}
                                                               \tag{88.10}
\]
\end{record}

\begin{decision}[V89 refines the winning trace card and then
populates vertices]
\label{dec:n1v88-next}
Evaluate the Hessian-only multiplicative card at \(512\) cells.
If its energy lies within a factor two of the dressed-Frobenius
diagnostic for both actions, freeze that cell count and populate
the nine vertex-speed families with (88.3).  Otherwise evaluate
\(1024\) Hessian cells before any vertex work.  Propagate the
winning piecewise \(a,c_X\) bounds through (80.5), retaining
cellwise hybrid minima.
\end{decision}

\begin{assessment}[Progress after compact V88]
\label{ass:n1v88-status}
The trace-retaining, floor-free correction is now both rigorous and
strictly better on every tested cell.  The exact path energy has
fallen from the V83 values \(1488.6,13691.9\) to \(20.40,45.97\).
The remaining immediate work is controlled cell refinement and
then vertex-speed population.
\end{assessment}

\begin{firewall}[Claim boundary after compact V88]
\label{fire:n1v88-claim}
Only one signed diagonal and only its Hessian speed have been
certified with the V88 architecture.  No vertex path card, torus
cover, global dressed stock, response quadrature, higher-loop or
nonperturbative estimate, SCIC, ADMC, continuum measure,
reflection positivity, exponential clustering, or Yang--Mills
mass gap is proved.
\end{firewall}

\subsection{Parent record}

This V88 note was formed by appending the preceding material to the
validated compact V87, whose record was
\[
\begin{array}{c|c}
\text{lines}&29116\\
\text{bytes}&992595\\
\text{SHA--256}&
\texttt{3A04E79BAAE73DC7AA06DD531246116663E56F607B9B5E2A69599B17260ED351}.
\end{array}
\]

% =============================================================================
% END APPENDED COMPACT-SERIES V88 MATERIAL
% =============================================================================

% =============================================================================
% BEGIN APPENDED COMPACT-SERIES V89 MATERIAL
% =============================================================================

\section{Compact-series V89: accepted asymmetric path resolution}
\label{sec:n1v89}

\subsection{The 512-cell two-action card}

The V88 multiplicative trace calculation was repeated on \(512\)
equal cells for both actions.

\begin{theorem}[Exact 512-cell refinement]
\label{thm:n1v89-512}
The all-minus diagonal has the following rationally certified
energy data:
\[
\begin{array}{c|cc}
e&{\rm Wilson}&\text{endpoint}\\ \hline
\max_j\eta_{L,j}=\max_j\eta_{R,j}
&0.079435938&0.217483983\\
\max_j a_j&15.244405&31.387168\\
\sum_j a_j/512&14.953051&26.655970.
\end{array}                                                   \tag{89.1}
\]
The multiplicative trial-trace bound is strictly smaller than
the Banach bound on every one of the \(512\) cells for each
action.
\end{theorem}

\begin{proof}
Use the V84 complete interval products and the two-sided residual
context of V86 in Theorem~\ref{thm:n1v88-multiplicative}.
Every residual is below one.  Exact Fraction comparison gives
\(512\) strict trace wins for each action, and exact summation gives
(89.1).
\end{proof}

Relative to the \(256\)-cell results (88.5), the energy bounds fall
by factors greater than \(1.36\) and \(1.72\).  The Wilson value is
\(1.7335\) times the independent dressed-Frobenius diagnostic
\(8.626265\), while the endpoint value is \(2.7313\) times
\(9.759721\).  These ratios select refinement levels only; their
denominators are diagnostic, not certified lower bounds.

\begin{record}[V89 two-action refinement verifier]
\label{rec:n1v89-512}
The executable
\[
 \texttt{scripts/verify\_ym\_v89\_trace\_refinement.py}      \tag{89.2}
\]
has \(99\) lines, \(3744\) bytes, and SHA--256
\[
 \texttt{E3C4F33A9CCEF82F5E72BDCF044E51BBB14C744544AC3CF8205292F0F3075F0C}.
                                                               \tag{89.3}
\]
Its canonical payload SHA--256 is
\[
 \boxed{
\texttt{AE835D0B707FDAF16354E5B007156072841BF4AF8FE57BF1ED4E1E81BC3EE160}.}
                                                               \tag{89.4}
\]
\end{record}

\subsection{Endpoint-only refinement}

Because Wilson passed the predeclared factor-two diagnostic gate
at \(512\) cells and the endpoint did not, only the endpoint card
was doubled.

\begin{theorem}[Exact endpoint 1024-cell card]
\label{thm:n1v89-endpoint}
At \(1024\) endpoint cells,
\[
\begin{split}
\max_j\eta_{L,j}=\max_j\eta_{R,j}
&\leq0.108755406,\\
\max_j a_j&\leq21.656532,\\
\sum_j a_j/1024&\leq18.671264.                              \tag{89.5}
\end{split}
\]
The multiplicative trace bound wins all \(1024\) cells.  The final
energy ceiling is \(1.9131\) times the floating dressed-Frobenius
diagnostic, so it passes the predeclared factor-two selection gate.
\end{theorem}

\begin{proof}
This is the same rational calculation as
Theorem~\ref{thm:n1v89-512}, with only the endpoint action and
half the cell width.  The exact upper energy Fraction is
\[
 {746850549\over40000000}=18.671263725.                      \tag{89.6}
\]
All residual, dominance, and diagnostic selection checks pass.
\end{proof}

\begin{record}[V89 endpoint refinement verifier]
\label{rec:n1v89-1024}
The executable
\[
 \texttt{scripts/verify\_ym\_v89\_endpoint\_1024.py}         \tag{89.7}
\]
has \(91\) lines, \(3216\) bytes, and SHA--256
\[
 \texttt{9CE1639DB9BC29CC05CF7143F807B8469DAFA466CD3BA7092C9EB3877B9A854A}.
                                                               \tag{89.8}
\]
Its canonical payload SHA--256 is
\[
 \boxed{
\texttt{2576708EF2ECC77A10A212337C8589195B4515760D564C35C4D650E8FD069B7C}.}
                                                               \tag{89.9}
\]
\end{record}

\subsection{Frozen Hessian resolutions}

The accepted action-dependent path resolutions are henceforth
\[
 N_{\rm W}=512,\qquad N_*=1024.                              \tag{89.10}
\]
They are not universal torus quadrature resolutions.  They certify
only the all-minus diagonal Hessian speed and provide the cell
contexts in which vertex speeds can now be bounded.

The next vertex computation has an important exact reduction.
Reflection of the stored cubic Laurent banks identifies the plus
and minus dressed norms, while the first external quartic bank is
identically zero.  Thus only
\[
 3\ \text{plus-cubic orders}
 \quad+\quad2\ \text{nonzero quartic orders}                 \tag{89.11}
\]
need independent path-speed cards per action; the remaining four
are copied by proved symmetry or exact zero.  Each cell context
\((R,\eta_L,\eta_R)\) is shared by all five computations.

\begin{decision}[V90 audits companions and batches vertex products]
\label{dec:n1v89-next}
Perform the mandatory current SM/NS audit.  Then build a cached
common-denominator batch compiler for the five independent vertex
speeds in (89.11) at the frozen resolutions (89.10).  Compute each
cell's inverse and residual once, multiply the five matrix centers
as one horizontal block, and reuse the radius products.  Apply
(88.3) to each block component and propagate through (80.5).
Do not run nine separate path-card scripts.
\end{decision}

\begin{assessment}[Progress after compact V89]
\label{ass:n1v89-status}
Both actions now have accepted exact Hessian path resolutions
within the predeclared factor-two diagnostic scale.  The Wilson
energy ceiling is \(14.953051\), and the endpoint ceiling is
\(18.671264\).  Residual margins are respectively greater than
\(0.9205\) and \(0.8912\).  The next cost is vertex population,
for which symmetry and batched products reduce nine families to
five.
\end{assessment}

\begin{firewall}[Claim boundary after compact V89]
\label{fire:n1v89-claim}
The accepted cards cover one signed path, not the torus.  No
vertex path speed, global dressed stock, response quadrature,
higher-loop or nonperturbative estimate, SCIC, ADMC, continuum
measure, reflection positivity, exponential clustering, or
Yang--Mills mass gap is proved.
\end{firewall}

\subsection{Parent record}

This V89 note was formed by appending the preceding material to the
validated compact V88, whose record was
\[
\begin{array}{c|c}
\text{lines}&29307\\
\text{bytes}&998813\\
\text{SHA--256}&
\texttt{E0C63D89385077672116BD14E5C3BE1A700A8BA32BDC7F8D435D30BD71406855}.
\end{array}
\]

% =============================================================================
% END APPENDED COMPACT-SERIES V89 MATERIAL
% =============================================================================

% =============================================================================
% BEGIN APPENDED COMPACT-SERIES V90 MATERIAL
% =============================================================================

\section{Compact-series V90: compulsory audit and an exact
quarter-turn dressed anchor}
\label{sec:n1v90}

\subsection{Current companion audit}

The mandatory five-version audit was performed on the newest
active companion notes then present:
\[
\begin{array}{c|c|r|r|c}
\text{programme}&\text{version}&\text{lines}&\text{bytes}&
\text{SHA--256}\\ \hline
\text{SM--Einstein}&V60&19441&683878&
\texttt{A409F5F82A592368004EB12133802EF9A5603E12CC9AE9F2FF6F8174DC4AD912}\\
\text{Navier--Stokes}&V26&5319&278283&
\texttt{82C887F8C2C69E4F28FAD7C3DC8DE5B9099CB5A4BF431EBA1FFF3E91935978AD}.
\end{array}                                                    \tag{90.1}
\]
SM V60 separates chart-dependent resolutions from the common
carrier on which overlap products and cocycles must be computed.
Its trace-class determinant and Pfaffian-line results supply no
Yang--Mills numerical constant.  The transferable discipline is
to allow asymmetric certified resolutions while sharing exact
matrix products.  NS V26 is unchanged and supplies no new
theorem or constant for the present finite inverse problem.

\begin{proposition}[V90 audit decision]
\label{prop:n1v90-audit}
No companion numerical bound is imported.  Before launching the
V89 path-wide five-bank batch, test whether the interior of the
torus contains coefficient-exact anchors.  Such anchors remove
long-distance transport from the value estimate and may seed a
genuine local cover.
\end{proposition}

\begin{proof}
Neither companion note proves an estimate for the Laurent
Hessian or vertex banks used here.  On the other hand, all their
stored frequencies are integral.  At a quarter-turn momentum,
every phase is a Gaussian-rational unit.  Direct exact evaluation
is therefore available and is logically prior to an interval
transport calculation.
\end{proof}

\subsection{Gaussian-rational residual dressing}

Put
\[
 k_c=(\pi/2,\pi/2,\pi/2,\pi/2).
\]
For every stored Laurent frequency
\(\alpha\in\mathbb Z^4\),
\[
 \exp(i\alpha\mathbin{\cdot}k_c)
 =i^{\alpha_1+\alpha_2+\alpha_3+\alpha_4}
 \in\{1,i,-1,-i\}.                                           \tag{90.2}
\]
Consequently the two Hessians and all nine external-jet matrices
at \(k_c\) belong exactly to \(M_{45}(\mathbb Q(i))\); no
trigonometric interval is present.

\begin{lemma}[Complex exact residual dressing]
\label{lem:n1v90-gaussian-dressing}
Let \(H=H^*\succeq\lambda I\), \(G=H^{-1}\), and let
\(R=R^*\) be a Gaussian-rational trial inverse.  Define
\[
 E=I-RH,\qquad
 \eta\geq\|E\|_F,\qquad d={\eta\over\lambda},\qquad
 r\geq\|R\|_{\rm op}.                                        \tag{90.3}
\]
Then, for every complex matrix \(X\), with
\(b=\|X\|_F\),
\[
\boxed{
\begin{split}
\|H^{-1/2}XH^{-1/2}\|_F^2
\leq{}&
\operatorname{Re}\operatorname{Tr}\big((RX)^*(XR)\big)\\
&+b^2(2dr+d^2).
\end{split}}                                                  \tag{90.4}
\]
If \(H,R,X\) are Gaussian rational and the four quantities on
the right of (90.3) are rational ceilings, then the right side of
(90.4) has an exact rational ceiling.
\end{lemma}

\begin{proof}
Since \(HG=I\),
\[
 G-R=(I-RH)G=EG,
\]
and hence
\[
 \|G-R\|_{\rm op}\leq
 \|E\|_{\rm op}\|G\|_{\rm op}
 \leq{\eta\over\lambda}=d.                                  \tag{90.5}
\]
Write \(D=G-R\).  Cyclicity of the trace gives
\[
\begin{split}
\|H^{-1/2}XH^{-1/2}\|_F^2
&=\operatorname{Tr}(X^*GXG)\\
&=\operatorname{Tr}(X^*RXR)
 +\operatorname{Tr}(X^*DXR)\\
&\quad+\operatorname{Tr}(X^*RXD)
 +\operatorname{Tr}(X^*DXD).                                \tag{90.6}
\end{split}
\]
The first term is
\(\operatorname{Tr}((RX)^*(XR))\).  It is real because its
complex conjugate is the same trace after a cyclic permutation.
Hilbert--Schmidt duality and
\(\|AX\|_F\leq\|A\|_{\rm op}\|X\|_F\) bound the remaining
three absolute values by \(b^2dr,b^2rd,b^2d^2\).
This proves (90.4).  Gaussian-rational matrix arithmetic and
rational outward square-root ceilings prove the last assertion.
\end{proof}

This is the complex form of
Lemma~\ref{lem:n1v78-residual-dressing}; the proof is recorded
because the central matrices need not be a real matrix times one
common phase.  No positivity of the trial inverse is assumed.

\subsection{The populated central anchor}

The verifier selects a Hermitian trial inverse by rounding a
floating inverse entrywise to denominator \(10^8\), but then
recomputes every Hessian, residual, trace, norm, perturbation, and
comparison in integer or \texttt{Fraction} arithmetic.  Floating
point therefore selects a rational witness and proves nothing.
The operator ceiling \(r\) is obtained from the exact Hermitian
absolute row-sum bound.  The already proved global floors
\[
 \lambda_{\rm W}={49\over1250},\qquad
 \lambda_*={343\over1250}                                    \tag{90.7}
\]
are used only in (90.3).

\begin{theorem}[Exact dressed quarter-turn anchor]
\label{thm:n1v90-central}
At \(k_c\), the inverse-residual Frobenius ceilings are
\[
 \eta_{\rm W}\leq {402493\over10^{12}},\qquad
 \eta_*\leq {12484347\over10^{12}}.                          \tag{90.8}
\]
For the quartic bank \(\widehat Q\) and the two cubic banks
\(\widehat P^\pm\), the dressed Frobenius norms in external
orders \(0,1,2\) have the following exact upper bounds.  Every
displayed integer is to be divided by \(10^{12}\):
\[
\begin{array}{c|c|rrr}
e&\text{bank}&0&1&2\\ \hline
{\rm W}&\widehat Q&
70328279763118&0&79048327338372\\
{\rm W}&\widehat P^\pm&
30795792003933&18107274242554&29570198805641\\
*&\widehat Q&
50464827388110&0&82965427843243\\
*&\widehat P^\pm&
5697017296178&3760706286652&5387120778620.
\end{array}                                                   \tag{90.9}
\]
In particular, the plus and minus cubic bounds agree in every
order and the order-one quartic matrix is exactly zero for both
actions.
\end{theorem}

\begin{proof}
Equation (90.2) evaluates the stored rational Laurent
coefficients exactly.  The Hessian normalization by \(1/3\) is
also rational.  Direct conjugate-transpose comparison proves both
Hessians Hermitian.

For each action, form \(R=C/10^8\), calculate the exact Gaussian
integer numerator of \(I-RH\), and take an outward rational
Frobenius ceiling.  This gives (90.8).  For each of the nine
external matrices, calculate
\[
 \operatorname{Tr}\big((RX)^*(XR)\big),\qquad
 \|X\|_F^2,\qquad \|R\|_{\rm op}^{\rm row}                  \tag{90.10}
\]
exactly and insert them into
Lemma~\ref{lem:n1v90-gaussian-dressing}.  Every trial trace has
zero imaginary numerator.  An outward rational square-root
ceiling, followed by a denominator-\(10^{12}\) ceiling, gives
(90.9).  Exact list comparison proves equality of the plus and
minus rows, while direct coefficient evaluation proves the two
quartic zeros.
\end{proof}

Numerically, the nonzero rows of (90.9) are
\[
\begin{array}{c|c|ccc}
{\rm W}&\widehat Q&70.328280&0&79.048328\\
{\rm W}&\widehat P^\pm&30.795793&18.107275&29.570199\\
*&\widehat Q&50.464828&0&82.965428\\
*&\widehat P^\pm&5.697018&3.760707&5.387121.
\end{array}                                                   \tag{90.11}
\]
These decimal lines are readings of the exact ceilings in
(90.9), not the certificate.

\begin{record}[V90 central-anchor verifier]
\label{rec:n1v90-verifier}
The executable
\[
 \texttt{scripts/verify\_ym\_v90\_central\_gaussian\_anchor.py}
                                                                  \tag{90.12}
\]
has \(187\) lines, \(6930\) bytes, and SHA--256
\[
\texttt{5F8C5455A38EC09F9F4F1D71F188A8C1ACE224B35F74AB9A51EA8EE9AD4CF964}.
                                                                  \tag{90.13}
\]
Its canonical payload SHA--256 is
\[
\boxed{
\texttt{1507C462BDA3CC8941469CAFB8C6B77D6C05CFC2E51A058E6E8B495434AAA114}.}
                                                                  \tag{90.14}
\]
The payload also binds the V57 completion cache, V59 cubic
certificate, and V60 source with respective SHA--256 values
\[
\begin{split}
&\texttt{DAEF29051D002C754593CF3EE08412DFC3CDF9C3C5EC14E6519230528AE165AF},\\
&\texttt{44DD8EA5ED92D3ACA87034A29AD64D40FC6570055C98882ED463E8D7B7D7E280},\\
&\texttt{7D5BFEE7FB834DFB5132FA2E1A9C5BC3CED7D57D366C557C8F370B6AD24D6FC8}.
\end{split}                                                     \tag{90.15}
\]
\end{record}

\subsection{Upstream consequence and next finite problem}

The V89 path card remains valid, but it is no longer the only way
to reach the interior.  The centre now has a direct exact value
certificate.  More generally,
\[
 \mathcal G_4
 :=
 \left\{ {\pi\over2}m:
 m\in(\mathbb Z/4\mathbb Z)^4\right\}                       \tag{90.16}
\]
contains \(256\) Gaussian-rational momenta.  Thus a prospective
torus cover can be seeded by exact interior anchors rather than
transporting every value from the sixteen V78 corners.

\begin{decision}[V91 certifies the quarter-turn anchor grid]
\label{dec:n1v90-next}
Evaluate both Hessians on \(\mathcal G_4\), classify the points
under the proved coordinate, reflection, and external-axis
symmetries, and certify one representative per genuine orbit.
For every representative retain the five independent dressed
values selected in (89.11), with exact plus/minus and zero checks.
Only after the orbit table is complete should local interval
cells be grown from the best anchors.  A failed residual or an
unexpected orbit split is to be resolved upstream at the exact
matrix level, not hidden in a larger transport constant.
\end{decision}

\begin{assessment}[Progress after compact V90]
\label{ass:n1v90-status}
An interior momentum now has a complete coefficient-exact dressed
vertex card for both actions.  The residuals are below
\(4.1\cdot10^{-7}\) and \(1.25\cdot10^{-5}\), and all five
independent vertex values are certified.  This changes the cover
architecture: exact quarter-turn anchors can replace long
single-path ancestry for local value estimates.
\end{assessment}

\begin{firewall}[Claim boundary after compact V90]
\label{fire:n1v90-claim}
Theorem~\ref{thm:n1v90-central} concerns one momentum.  The other
\(255\) quarter-turn points, their symmetry orbits, neighborhoods,
and overlaps are not yet certified.  No torus cover, global
dressed stock, response quadrature, interacting construction,
SCIC, ADMC, continuum measure, reflection positivity,
exponential clustering, or Yang--Mills mass gap is proved.
\end{firewall}

\subsection{Parent record}

This V90 note was formed by appending the preceding material to
the validated compact V89, whose record was
\[
\begin{array}{c|c}
\text{lines}&29492\\
\text{bytes}&1005424\\
\text{SHA--256}&
\texttt{A2F30094ED229D2708E7848C881DE56D1F45DD51DE04F8C777844DAE0413D6EF}.
\end{array}
\]

% =============================================================================
% END APPENDED COMPACT-SERIES V90 MATERIAL
% =============================================================================

% =============================================================================
% BEGIN APPENDED COMPACT-SERIES V91 MATERIAL
% =============================================================================

\section{Compact-series V91: exact screening of the full
quarter-turn grid}
\label{sec:n1v91}

\subsection{Why no raw-matrix orbit quotient is taken}

Decision~\ref{dec:n1v90-next} proposed a symmetry-orbit reduction.
The earlier covariance theorem applies to the complete measured
response scalar after all re-tree maps and density factors cancel.
It explicitly does not assert entrywise covariance of the fixed-tree
\(45\times45\) Hessian and vertex matrices.  A dressed Frobenius
norm of one raw vertex is not automatically invariant under an
uncontrolled nonunitary coordinate change.

\begin{proposition}[Safe replacement of the proposed orbit step]
\label{prop:n1v91-no-orbit}
For the raw dressed anchors, enumerate all \(256\) points of
\(\mathcal G_4\) exactly.  Symmetry may be used to interpret
equalities found by the enumeration, but not to omit a point unless
the required unitary intertwiner for the raw Hessian and vertex has
itself been certified.
\end{proposition}

\begin{proof}
Theorem~\ref{thm:n1v37-covariance} cancels the re-tree maps only
inside the complete trace scalar.  V34 also records that fixed-tree
coordinate matrices are not literally axis-permutation matrices.
Neither result supplies unitary similarity for each separate raw
matrix norm.  Direct coefficient-exact enumeration has only \(4^4\)
points and introduces no such hypothesis.
\end{proof}

\subsection{A complete floor screen}

\begin{lemma}[Global-floor dressed screen]
\label{lem:n1v91-floor-screen}
If \(H\succeq\lambda I>0\), then every complex matrix \(X\)
satisfies
\[
 \boxed{\quad
 \|H^{-1/2}XH^{-1/2}\|_F
 \leq {\|X\|_F\over\lambda}.
 \quad}                                                       \tag{91.1}
\]
\end{lemma}

\begin{proof}
Twice applying the Hilbert--Schmidt ideal inequality gives
\[
\begin{split}
\|H^{-1/2}XH^{-1/2}\|_F
&\leq
\|H^{-1/2}\|_{\rm op}^2\|X\|_F\\
&\leq\lambda^{-1}\|X\|_F.
\end{split}
\]
\end{proof}

This estimate is weaker than the residual dressing of V78 and V90.
Its purpose is to screen every exact anchor cheaply before trial
inverses are allocated.

\begin{theorem}[Exact all-point quarter-turn screen]
\label{thm:n1v91-grid}
At every
\[
 k={\pi\over2}m,\qquad m\in(\mathbb Z/4\mathbb Z)^4,          \tag{91.2}
\]
both Hessians are exactly Hermitian Gaussian-rational matrices.
The exact maxima of the raw squared Frobenius norms are
\[
\begin{array}{c|c|ccc}
e&\text{bank}&0&1&2\\ \hline
{\rm W}&\widehat Q&368&0&118603/216\\
{\rm W}&\widehat P^\pm&92&273/8&18355/192\\
*&\widehat Q&590648&0&
146071511825067467/31233735000\\
*&\widehat P^\pm&146510/9&86077/16&
12272549901145141/749609640000.
\end{array}                                                   \tag{91.3}
\]
Using the proved floors (90.7), Lemma~\ref{lem:n1v91-floor-screen}
therefore gives the exact global dressed ceilings
\[
\begin{array}{c|c|ccc}
e&\text{bank}&0&1&2\\ \hline
{\rm W}&\widehat Q&
489370563603339/10^{12}&0&
18680346011181/31250000000\\
{\rm W}&\widehat P^\pm&
24468528180167/10^{11}&
74510978606811/(5\cdot10^{11})&
4988506112601/(2\cdot10^{10})\\
*&\widehat Q&
2800788159994213/10^{12}&0&
788109739260151/10^{11}\\
*&\widehat P^\pm&
116243393830717/(25\cdot10^{10})&
66825077557757/(25\cdot10^{10})&
233150160710263/(5\cdot10^{11}).
\end{array}                                                   \tag{91.4}
\]
The plus and minus raw squared norms agree pointwise, not merely
after maximizing.  The quartic order-one matrix is zero at all
\(256\) grid points for both actions.
\end{theorem}

\begin{proof}
At (91.2), every Laurent phase is
\[
 i^{\alpha\cdot m}\in\mathbb Q(i).
\]
The verifier evaluates every coefficient of the two Hessians and
the eighteen stored vertex matrices at each point.  Exact
conjugate-transpose comparison passes for all
\(2\cdot256=512\) Hessians.  It sums
\[
 \|X\|_F^2
 =\sum_{a,b}\big((\operatorname{Re}X_{ab})^2+
                  (\operatorname{Im}X_{ab})^2\big)           \tag{91.5}
\]
as a \texttt{Fraction}, compares all \(256\) values in each row,
and obtains (91.3).  Division by
\(\lambda_e^2\), followed by an outward rational square root and
a denominator-\(10^{12}\) ceiling, gives (91.4).

The same evaluator independently replays the V90 centre.  Exact
pointwise list comparison proves the plus/minus assertion.
All \(512\) order-one quartic squared norms are the zero Fraction,
proving the final assertion.
\end{proof}

For scale, the nonzero rows of (91.4) read
\[
\begin{array}{c|c|ccc}
{\rm W}&\widehat Q&489.371&0&597.772\\
{\rm W}&\widehat P^\pm&244.686&149.022&249.426\\
*&\widehat Q&2800.789&0&7881.098\\
*&\widehat P^\pm&464.974&267.301&466.301.
\end{array}                                                   \tag{91.6}
\]
These deliberately coarse values must not replace the sharper V78
and V90 residual anchors.

\subsection{Exact maximizer pattern}

Write a grid point by its integer vector \(m\).  The maximizers in
(91.3) have the following exact pattern:
\[
\begin{array}{c|c|c|l}
e&\text{bank/order}&\text{count}&\text{representatives}\\ \hline
{\rm W}&Q_0,P_0&256&\text{all points}\\
{\rm W}&Q_1&256&\text{all are zero}\\
{\rm W}&Q_2,P_1&16&(0,a,2,b),\ a,b\in\mathbb Z/4\mathbb Z\\
{\rm W}&P_2&1&(2,0,2,0)\\
*&Q_0,P_0&2&(0,0,0,0),(0,0,2,0)\\
*&Q_1&256&\text{all are zero}\\
*&Q_2&1&(0,0,2,0)\\
*&P_1&2&(1,0,2,0),(3,0,2,0)\\
*&P_2&1&(0,0,0,0).
\end{array}                                                   \tag{91.7}
\]
Here \(P_j\) denotes either cubic orientation because their raw
norms agree pointwise.  Equation (91.7) is an output of exact
enumeration, not a symmetry assumption.

\begin{record}[V91 full-grid verifier]
\label{rec:n1v91-verifier}
The executable
\[
 \texttt{scripts/verify\_ym\_v91\_quarter\_turn\_grid.py}     \tag{91.8}
\]
has \(256\) lines, \(9028\) bytes, and SHA--256
\[
\texttt{CED199F29812005F970830B84158E3115B3EDF920C2ACF8714EFA156AD06673A}.
                                                                  \tag{91.9}
\]
The SHA--256 of the canonical table containing all pointwise raw
squares is
\[
\texttt{D4BD8C4C1E06B515791574E93355E471E30B89E39F9FC1256AF58F258A0765C3},
                                                                  \tag{91.10}
\]
and the compact canonical payload SHA--256 is
\[
\boxed{
\texttt{6BFA766686BCA098B5060F9A5EF35D2257DAC8C8974AB2653BF504D25D5EBAF9}.}
                                                                  \tag{91.11}
\]
\end{record}

\begin{decision}[V92 sharpens the grid without a false orbit
reduction]
\label{dec:n1v91-next}
Use the pointwise exact raw squares as a scheduling layer.  First
reuse the V78 residual certificates on the sixteen dyadic points
\(m\in\{0,2\}^4\).  On the remaining \(240\) points, compute each
action's rational trial inverse once and batch the five independent
vertex products.  Apply
Lemma~\ref{lem:n1v90-gaussian-dressing} pointwise, preserving a
separate anchor row whenever no certified raw-matrix unitary
intertwiner is available.  Compare the resulting full-grid maxima
with both the floor screen (91.4) and the V78 corner maxima.
\end{decision}

\begin{assessment}[Progress after compact V91]
\label{ass:n1v91-status}
Every quarter-turn grid point now has a coefficient-exact raw
vertex record and a rigorous, though loose, dressed ceiling.  The
enumeration proves pointwise cubic reflection equality and a
grid-wide quartic first-order zero.  It also isolates the cost of
the next step: \(240\) new pointwise inverse contexts rather than
an unjustified raw-coordinate orbit quotient.
\end{assessment}

\begin{firewall}[Claim boundary after compact V91]
\label{fire:n1v91-claim}
The discrete screen does not control any open neighborhood between
grid points.  The floor ceilings in (91.4) are not sharp residual
anchors.  No local-cell overlap, torus cover, global dressed stock,
response quadrature, interacting construction, SCIC, ADMC,
continuum measure, reflection positivity, exponential clustering,
or Yang--Mills mass gap is proved.
\end{firewall}

\subsection{Parent record}

This V91 note was formed by appending the preceding material to the
corrected compact V90, whose record was
\[
\begin{array}{c|c}
\text{lines}&29792\\
\text{bytes}&1016484\\
\text{SHA--256}&
\texttt{AF84B9FB5A4EBAD274F74F6DCA4A12E2796C6FD684165FA98A32D7E4447EA92E}.
\end{array}
\]

% =============================================================================
% END APPENDED COMPACT-SERIES V91 MATERIAL
% =============================================================================

% =============================================================================
% BEGIN APPENDED COMPACT-SERIES V92 MATERIAL
% =============================================================================

\section{Compact-series V92: residual-dressed completion of the
quarter-turn grid}
\label{sec:n1v92}

\subsection{One independent rational context per point}

An exact preliminary comparison found \(256\) distinct fixed-tree
Hessians for each action on \(\mathcal G_4\).  Thus no literal
matrix equality can recycle an inverse context.  The V92 verifier
therefore forms a separate Hermitian rational trial inverse for
each of the \(512\) action--point pairs and evaluates all nine
vertices against that same context.

\begin{theorem}[Exact residual-dressed quarter-turn grid]
\label{thm:n1v92-grid}
At all \(256\) points of \(\mathcal G_4\), every trial inverse
passes the residual-one gate.  The largest residual ceilings are
\[
\begin{array}{c|c|c}
e&\max\eta&\text{maximizing integer points}\\ \hline
{\rm W}&111971/(25\cdot10^{10})&
(1,3,0,0),(3,1,0,0)\\
*&7926629/(5\cdot10^{11})&
(1,1,0,0),(3,3,0,0).
\end{array}                                                   \tag{92.1}
\]
In particular, these are respectively
\(4.47884\cdot10^{-7}\) and \(1.5853258\cdot10^{-5}\).

The maxima of the residual-dressed Frobenius ceilings on the full
grid are
\[
\begin{array}{c|c|ccc}
e&\text{bank}&0&1&2\\ \hline
{\rm W}&\widehat Q&
109682140353581/10^{12}&0&
12487701095881/10^{11}\\
{\rm W}&\widehat P^\pm&
442359636459/10^{10}&
25954502058787/10^{12}&
42959328239297/10^{12}\\
*&\widehat Q&
50012730675093/(5\cdot10^{11})&0&
186072859017781/10^{12}\\
*&\widehat P^\pm&
910221750203/10^{11}&
1407103178583/(25\cdot10^{10})&
1638986019113/(25\cdot10^{10}).
\end{array}                                                   \tag{92.2}
\]
The plus and minus certified triples agree at every point, the
quartic order-one bound is identically zero, and every entry of
(92.2) is dominated by its corresponding V91 floor-screen entry.
\end{theorem}

\begin{proof}
At every point and for each action, evaluate the Hessian and
vertex banks in \(\mathbb Q(i)\) as in (90.2).  Round a floating
inverse to a Hermitian Gaussian-rational matrix \(R\) with
denominator \(10^8\), then recompute
\[
 E=I-RH,\quad \eta,\quad r,\quad
 \operatorname{Tr}((RX)^*(XR)),\quad\|X\|_F^2               \tag{92.3}
\]
exactly.  Apply
Lemma~\ref{lem:n1v90-gaussian-dressing} using the global floors
(90.7).  This produces \(2\cdot256\cdot9=4608\) dressed
ceilings.

Exact comparison gives (92.1)--(92.2).  Every residual is below
one.  At \(m=(1,1,1,1)\), all eighteen denominator-\(10^{12}\)
ceilings reproduce Theorem~\ref{thm:n1v90-central}.  Pointwise
comparison gives equality of the two cubic-orientation bound
triples.  The zero quartic row and the dominance by (91.4) are
also checked as exact Fraction comparisons.
\end{proof}

The maximizing integer points in (92.2) are
\[
\begin{array}{c|c|l}
e&\text{bank/order}&\text{maximizers}\\ \hline
{\rm W}&Q_0,Q_2&(2,2,1,2),(2,2,3,2)\\
{\rm W}&P_0&(2,2,1,0),(2,2,3,0)\\
{\rm W}&P_1&(1,2,3,0),(3,2,1,0)\\
{\rm W}&P_2&(2,2,1,1),(2,2,3,3)\\
*&Q_0,Q_2&(2,2,0,2)\\
*&P_0,P_2&(2,2,1,2),(2,2,3,2)\\
*&P_1&(1,2,3,2),(3,2,1,2).
\end{array}                                                   \tag{92.4}
\]
For the endpoint \(P_0,P_2\) row, the displayed pair maximizes
each of the two orders; the same pair is shared.  All point lists
are bound into the certificate even where (92.4) abbreviates
their common structure.

Numerically, (92.2) reads
\[
\begin{array}{c|c|ccc}
{\rm W}&\widehat Q&109.682141&0&124.877011\\
{\rm W}&\widehat P^\pm&44.235964&25.954503&42.959329\\
*&\widehat Q&100.025462&0&186.072860\\
*&\widehat P^\pm&9.102218&5.628413&6.555945.
\end{array}                                                   \tag{92.5}
\]
The decimals are upward readings only.

\subsection{The interior grid is not redundant}

Every nonzero Wilson maximum in (92.2) exceeds the corresponding
V78 corner maximum.  The endpoint cubic maxima also rise from
\[
 (5.480910,3.291328,5.555549)
\quad\hbox{to}\quad
 (9.102218,5.628413,6.555945),                               \tag{92.6}
\]
whereas the endpoint quartic grid maxima stay slightly below the
V78 corner ceilings.  Thus neither the V78 corner table nor the
single V90 centre controls the quarter-turn grid.  Exact interior
anchors are a necessary part of a cover.

\begin{record}[V92 residual-grid verifier]
\label{rec:n1v92-verifier}
The executable
\[
 \texttt{scripts/verify\_ym\_v92\_residual\_dressed\_grid.py} \tag{92.7}
\]
has \(329\) lines, \(10459\) bytes, and SHA--256
\[
\texttt{D0B087FF53D38B8FD506660447D46C934F2F8A0270CAF19D03173A166FBB4D06}.
                                                                  \tag{92.8}
\]
It uses three independent workers only to schedule exact point
calculations; the mathematical result is independent of their
order.  The SHA--256 of the canonical pointwise residual-and-bound
table is
\[
\texttt{5DB6767FF761679240E1A270B659F991370CD083FBDEEA172CA0234BE4361AF2},
                                                                  \tag{92.9}
\]
and the compact canonical payload SHA--256 is
\[
\boxed{
\texttt{75C63FEC6C23CB04AEFD5A8C2BAFA73BFD4A82BFF6076871ECDAAD115C8F7C72}.}
                                                                  \tag{92.10}
\]
\end{record}

\subsection{The local-cover bridge}

The exact grid still controls no open set.  The next step should
not discard the successful anchor trial traces by reverting to
the global floor estimate.  Let \(a\in\mathcal G_4\), and retain
its rational \(R\).  For a nearby \(k\), coefficient stocks can
bound
\[
 \delta_H\geq\|H(k)-H(a)\|_F,\qquad
 \delta_X\geq\|X(k)-X(a)\|_F.                               \tag{92.11}
\]
Then
\[
 \|I-RH(k)\|_{\rm op}
 \leq\eta_a+r_a\delta_H                                    \tag{92.12}
\]
and the trial trace can be transported using only
\(RX(a),X(a)R,r_a,\delta_X\).  This retains the V90 mechanism
inside a neighborhood.

\begin{decision}[V93 proves the fixed-anchor cell lemma]
\label{dec:n1v92-next}
Prove a complete version of (92.12): an explicit upper formula
for the dressed norm throughout a box about \(a\), expressed in
the exact anchor quantities
\[
 (\eta_a,r_a,\|X(a)\|_F,\|RX(a)\|_F,\|X(a)R\|_F,
   \operatorname{Re}\operatorname{Tr}((RX(a))^*(X(a)R)))
                                                                  \tag{92.13}
\]
and the two variation radii \(\delta_H,\delta_X\).
Then compute exact Laurent coefficient Lipschitz stocks for both
actions and the five independent vertex families.  Do not begin
a \(256\)-box interval run until the resulting residual-one radius
has a proved positive value.
\end{decision}

\begin{assessment}[Progress after compact V92]
\label{ass:n1v92-status}
All \(256\) quarter-turn points now carry independent exact
residual-dressed vertex certificates for both actions.  The
worst inverse residual remains below \(1.6\cdot10^{-5}\), and
the floor-only V91 bounds improve by factors ranging from more
than four to more than forty.  The remaining finite-dimensional
bottleneck is no longer anchor evaluation but certified extension
from anchors to open neighborhoods.
\end{assessment}

\begin{firewall}[Claim boundary after compact V92]
\label{fire:n1v92-claim}
The V92 theorem is discrete.  No positive neighborhood radius,
overlap, torus cover, global dressed stock, response quadrature,
interacting construction, SCIC, ADMC, continuum measure,
reflection positivity, exponential clustering, or Yang--Mills
mass gap is proved.
\end{firewall}

\subsection{Parent record}

This V92 note was formed by appending the preceding material to the
validated compact V91, whose record was
\[
\begin{array}{c|c}
\text{lines}&30038\\
\text{bytes}&1025165\\
\text{SHA--256}&
\texttt{E89F4BB2D4C115AC53D8DE15E095AF5B35593B7E1471D16BC55405BA77766E05}.
\end{array}
\]

% =============================================================================
% END APPENDED COMPACT-SERIES V92 MATERIAL
% =============================================================================

% =============================================================================
% BEGIN APPENDED COMPACT-SERIES V93 MATERIAL
% =============================================================================

\section{Compact-series V93: a rigorous global Laurent cover of
the momentum torus}
\label{sec:n1v93}

\subsection{Coefficient-exact variation stocks}

Let
\[
 X(k)=\sum_{\alpha\in{\cal A}}C_\alpha e^{i\alpha\cdot k}
\]
be any stored Hessian or vertex Laurent matrix, and define
\[
 L_X:=\sum_{\alpha\in{\cal A}}
 |\alpha|_1\,\|C_\alpha\|_F.                                \tag{93.1}
\]

\begin{lemma}[Torus Lipschitz stock]
\label{lem:n1v93-Laurent}
For two torus momenta represented with coordinate difference
\(\delta\),
\[
 \boxed{\quad
 \|X(k+\delta)-X(k)\|_F
 \leq L_X\|\delta\|_\infty.
 \quad}                                                       \tag{93.2}
\]
\end{lemma}

\begin{proof}
The elementary inequality
\(|e^{it}-1|\leq|t|\) gives
\[
\begin{split}
\|X(k+\delta)-X(k)\|_F
&\leq\sum_\alpha
\|C_\alpha\|_F\,|e^{i\alpha\cdot\delta}-1|\\
&\leq\sum_\alpha
\|C_\alpha\|_F\,|\alpha|_1\|\delta\|_\infty,
\end{split}
\]
which is (93.2).
\end{proof}

The exact outward stocks computed from the rational coefficient
matrices are
\[
\begin{array}{c|c|ccc}
e&\text{bank}&0&1&2\\ \hline
{\rm W}&\widehat Q&
8937809809481/(25\cdot10^{10})&0&
10844297285731/(25\cdot10^{10})\\
{\rm W}&\widehat P^\pm&
8937809809481/(5\cdot10^{11})&
2228218138683/(2\cdot10^{11})&
9266722015957/(5\cdot10^{11})\\
*&\widehat Q&
106515734805347/(5\cdot10^{10})&0&
281490361922521/31250000000\\
*&\widehat P^\pm&
8853694408657/(25\cdot10^9)&
138196818966003/(25\cdot10^{10})&
385395415946427/(5\cdot10^{11}).
\end{array}                                                   \tag{93.3}
\]
The corresponding Hessian stocks are
\[
 L_{H_{\rm W}}\leq {3646905023083\over125000000000},
 \qquad
 L_{H_*}\leq {159654323446791\over125000000000}.             \tag{93.4}
\]

\subsection{A complete torus cover}

\begin{theorem}[Global dressed-vertex stock]
\label{thm:n1v93-global}
For every \(k\in(\mathbb R/2\pi\mathbb Z)^4\), the dressed
vertex norms have the following exact uniform ceilings:
\[
\begin{array}{c|c|ccc}
e&\text{bank}&0&1&2\\ \hline
{\rm W}&\widehat Q&
1205670564137421/10^{12}&0&
366715520506967/(25\cdot10^{10})\\
{\rm W}&\widehat P^\pm&
150708820517181/(25\cdot10^{10})&
186120582074459/(5\cdot10^{11})&
77594407272983/(125\cdot10^9)\\
*&\widehat Q&
2224564197299849/(25\cdot10^{10})&0&
16831629281777027/(5\cdot10^{11})\\
*&\widehat P^\pm&
1478628956999781/10^{12}&
462377326172039/(25\cdot10^{10})&
2672487475710203/10^{12}.
\end{array}                                                   \tag{93.5}
\]
\end{theorem}

\begin{proof}
Choose in each coordinate a nearest multiple of \(\pi/2\).
This produces \(a\in\mathcal G_4\) with torus coordinate distance
\[
 \|k-a\|_\infty\leq{\pi\over4}
 <{355\over452}.                                             \tag{93.6}
\]
For each bank and order, let \(B_X\) be the outward square-root
ceiling of the corresponding raw grid maximum (91.3).
Lemma~\ref{lem:n1v93-Laurent} gives
\[
 \|X(k)\|_F
 \leq B_X+{355\over452}L_X.                                 \tag{93.7}
\]
The proved global Hessian floor and
Lemma~\ref{lem:n1v91-floor-screen} now give
\[
 \|H_e(k)^{-1/2}X(k)H_e(k)^{-1/2}\|_F
 \leq
 {B_X+(355/452)L_X\over\lambda_e}.                          \tag{93.8}
\]
All inputs in (93.8) are rational ceilings.  Exact substitution
of (90.7), (91.3), and (93.3), followed by an outward
denominator-\(10^{12}\) ceiling, is precisely (93.5).
\end{proof}

Numerically, the nonzero rows of (93.5) are
\[
\begin{array}{c|c|ccc}
{\rm W}&\widehat Q&1205.671&0&1466.863\\
{\rm W}&\widehat P^\pm&602.836&372.242&620.756\\
*&\widehat Q&8898.257&0&33663.259\\
*&\widehat P^\pm&1478.629&1849.510&2672.488.
\end{array}                                                   \tag{93.9}
\]
They are coarse but finite global quantities; no sampled-point
inference occurs.

\subsection{A sharper fixed-anchor extension}

\begin{lemma}[Residual-retaining anchor cell]
\label{lem:n1v93-anchor-cell}
Fix an anchor \(a\), write \(H_0=H(a)\), \(X_0=X(a)\), and let
\(R=R^*\) be its rational trial inverse.  Suppose
\[
\begin{gathered}
\eta_0\geq\|I-RH_0\|_F,\qquad r\geq\|R\|_{\rm op},\\
\delta_H\geq\|H(k)-H_0\|_F,\qquad
\delta_X\geq\|X(k)-X_0\|_F.
\end{gathered}                                                \tag{93.10}
\]
Put
\[
\begin{gathered}
e=\eta_0+r\delta_H,\qquad d=e/\lambda,\\
x_0=\|X_0\|_F,\quad
a_0=\|RX_0\|_F,\quad b_0=\|X_0R\|_F,\\
t_0=\operatorname{Re}\operatorname{Tr}
       ((RX_0)^*(X_0R)).
\end{gathered}                                                \tag{93.11}
\]
Then throughout the cell,
\[
\boxed{
\begin{split}
\|H(k)^{-1/2}X(k)H(k)^{-1/2}\|_F^2
\leq{}&
t_0+r\delta_X(a_0+b_0)+r^2\delta_X^2\\
&+(x_0+\delta_X)^2(2dr+d^2).
\end{split}}                                                  \tag{93.12}
\]
\end{lemma}

\begin{proof}
With \(\Delta H=H(k)-H_0\),
\[
\|I-RH(k)\|_{\rm op}
\leq\eta_0+r\delta_H=e.                                     \tag{93.13}
\]
The global floor gives
\(\|H(k)^{-1}-R\|_{\rm op}\leq e/\lambda=d\), exactly as in
(90.5).  Write \(X(k)=X_0+Z\).  Hilbert--Schmidt duality gives
\[
\begin{split}
\operatorname{Re}\operatorname{Tr}((RX(k))^*(X(k)R))
\leq{}&t_0+r\delta_X(a_0+b_0)+r^2\delta_X^2,\\
\|X(k)\|_F&\leq x_0+\delta_X.
\end{split}                                                   \tag{93.14}
\]
Insert (93.14) and \(d\) in
Lemma~\ref{lem:n1v90-gaussian-dressing}.
\end{proof}

Unlike a Neumann-series estimate, (93.12) remains valid even when
\(e\geq1\).  When \(e<1\), the floor-free V88 correction is also
available and may be compared cellwise.

\begin{record}[V93 global-cover verifier]
\label{rec:n1v93-verifier}
The executable
\[
 \texttt{scripts/verify\_ym\_v93\_global\_laurent\_cover.py} \tag{93.15}
\]
has \(195\) lines, \(6184\) bytes, and SHA--256
\[
\texttt{7236A89C8EE5BC535DE6E8FAD6B547509A9611EA66BE4EC66DD727F797BEC98C}.
                                                                  \tag{93.16}
\]
Its compact canonical payload SHA--256 is
\[
\boxed{
\texttt{1F6FEF7DABF2DEA779448E1A4FDF49DF861A11FFAB0F51576281364F29742E30}.}
                                                                  \tag{93.17}
\]
\end{record}

\begin{decision}[V94 assembles a global response envelope]
\label{dec:n1v93-next}
Insert the global stocks (93.5) into the already established
trace-ideal tadpole and cubic-bubble word bounds through external
order two.  Produce one torus-uniform scalar response envelope for
each action and the endpoint difference.  If that envelope is too
large for the intended Fourier or contraction gate, populate
Lemma~\ref{lem:n1v93-anchor-cell} with the V92 anchor packets and
subdivide only the offending cells; do not abandon the valid
global cover.
\end{decision}

\begin{assessment}[Progress after compact V93]
\label{ass:n1v93-status}
The finite-dimensional vertex problem now has its first rigorous
full-torus dressed stock.  The proof covers every momentum by a
nearest exact quarter-turn anchor and coefficient-exact Laurent
variation.  A residual-retaining cell lemma is available for
sharpening without losing the successful V92 trial traces.
\end{assessment}

\begin{firewall}[Claim boundary after compact V93]
\label{fire:n1v93-claim}
The global ceilings are coarse and have not yet been propagated
through the complete response word.  No quadrature error,
continuum response coefficient, interacting construction, SCIC,
ADMC, continuum measure, reflection positivity, exponential
clustering, or Yang--Mills mass gap is proved.
\end{firewall}

\subsection{Parent record}

This V93 note was formed by appending the preceding material to the
validated compact V92, whose record was
\[
\begin{array}{c|c}
\text{lines}&30263\\
\text{bytes}&1033293\\
\text{SHA--256}&
\texttt{D8CF2B210148E8D5AFD9BC31A66CF71FC716C42915F053116B751DF2516E0EE6}.
\end{array}
\]

% =============================================================================
% END APPENDED COMPACT-SERIES V93 MATERIAL
% =============================================================================

% =============================================================================
% BEGIN APPENDED COMPACT-SERIES V94 MATERIAL
% =============================================================================

\section{Compact-series V94: a global dressed response envelope}
\label{sec:n1v94}

\subsection{Normalized covariance input}

For external order \(j\), retain the V68 notation
\[
 U_j=H^{1/2}G_jH^{1/2},\qquad
 u_j\geq\sup_k\|U_j(k)\|_{\rm op}.                          \tag{94.1}
\]
The globally certified V69 Gram-factor recursion gives
\[
\begin{array}{c|ccc}
e&u_0&u_1&u_2\\ \hline
{\rm W}&1&
2727723627951/(125\cdot10^9)&
242459595900151/(2\cdot10^{11})\\
*&1&
43837545660617/(5\cdot10^{11})&
9652489593394019/(5\cdot10^{11}).
\end{array}                                                   \tag{94.2}
\]
These are global rational bounds, unlike the V68 sampled
diagnostics.

\subsection{External response orders zero through two}

Let \(p_j\) and \(q_j\) denote the corresponding global dressed
cubic and quartic ceilings from (93.5).  Equality of the plus and
minus ceilings permits one symbol \(p_j\), but the two matrices
remain distinct.

\begin{theorem}[Global pointwise response envelopes]
\label{thm:n1v94-response}
For \(r=0,1,2\), define
\[
\begin{split}
 E_{{\cal T},e}^{[r]}&:={\sqrt{45}^{\,\rm up}\over2}\,q_r,\\
 E_{{\cal B},e}^{[r]}&:=
 |c_e|\sum_{a+b+c=r}{r!\over a!\,b!\,c!}\,
 p_a u_b p_c,\\
 E_e^{[r]}&:=E_{{\cal T},e}^{[r]}+
             E_{{\cal B},e}^{[r]},
\end{split}                                                   \tag{94.3}
\]
where \(c_{\rm W}=1/3\), \(c_*=3\), and
\(\sqrt{45}^{\,\rm up}\) is the established rational upper
square root.  Then
\[
 \sup_k|{\cal R}_e^{[r]}(k)|\leq E_e^{[r]}.                 \tag{94.4}
\]
Exact rational ceilings for the totals are
\[
\begin{array}{c|ccc}
e&E^{[0]}&E^{[1]}&E^{[2]}\\ \hline
{\rm W}&
3129518361136227/(25\cdot10^9)&
279302159898502637/10^{11}&
768648711080468041/(5\cdot10^9)\\
*&
1317775287605526347/(2\cdot10^{11})&
18483501577352920643/31250000000&
129543520227669516498251/10^{12}.
\end{array}                                                   \tag{94.5}
\]
For the endpoint-minus-Wilson response,
\[
\begin{array}{c|ccc}
r&0&1&2\\ \hline
\sup_k|\Delta{\cal R}^{[r]}(k)|&
1342811434494616163/(2\cdot10^{11})&
297132536037139243473/(5\cdot10^{11})&
2593944999397712202129/(2\cdot10^{10}).
\end{array}                                                   \tag{94.6}
\]
\end{theorem}

\begin{proof}
The V60 tadpole formula fixes the covariance and puts all external
derivatives on the quartic vertex.  The dressed trace identity
(77.2) and Hilbert--Schmidt duality give
\[
 |{\cal T}_e^{[r]}(k)|
 \leq{\sqrt{45}\over2}\,
 \|\widehat Q_e^{[r]}(k)\|_F.                               \tag{94.7}
\]
In the bubble, the \(r\) derivatives are distributed among the
plus cubic, the shifted covariance, and the minus cubic.  The
multinomial Leibniz rule, (77.3), and the trace-ideal estimate
(77.5) give exactly the second line of (94.3).  The global inputs
(93.5) and (94.2) prove (94.4).

All products and sums are performed as Fractions.  Independent
outward denominator-\(10^{12}\) ceilings give (94.5).
Finally,
\[
 |\Delta{\cal R}^{[r]}|
 \leq|{\cal R}_*^{[r]}|+|{\cal R}_{\rm W}^{[r]}|,
\]
and exact addition gives (94.6).
\end{proof}

Numerically, the action totals are
\[
\begin{array}{c|ccc}
{\rm W}&1.25181\cdot10^5&2.79303\cdot10^6&
1.53730\cdot10^8\\
*&6.58888\cdot10^6&5.91473\cdot10^8&
1.29544\cdot10^{11},
\end{array}                                                   \tag{94.8}
\]
and the difference order-two envelope is
\[
 1.29698\cdot10^{11}.                                       \tag{94.9}
\]
This is a domination estimate, not a sign enclosure.

\begin{record}[V94 response-envelope verifier]
\label{rec:n1v94-verifier}
The executable
\[
 \texttt{scripts/verify\_ym\_v94\_global\_response\_envelope.py}
                                                                  \tag{94.10}
\]
has \(186\) lines, \(5753\) bytes, and SHA--256
\[
\texttt{5BC401A6F0B4C1706664B7A41D6DA2FC751E6D4CF661E96BE004D1456539ED5C}.
                                                                  \tag{94.11}
\]
Its compact canonical payload SHA--256 is
\[
\boxed{
\texttt{D4A7B72D1BEE81F3B8B33519D9AC3DE18CE3A8FBF66870A2234145466EE9267D}.}
                                                                  \tag{94.12}
\]
Four exact checks cover the covariance base, the quartic zero,
nonnegativity, and action-sum difference identity.
\end{record}

\subsection{Location of the remaining loss}

The V92 endpoint cubic ceilings are all below \(10\), but the V93
cover enlarges them to \(1479,1850,2673\).  In addition,
\[
 u_{2,*}>19304.                                             \tag{94.13}
\]
The large order-two response envelope is therefore caused by two
separate global triangle losses: Laurent transport of the dressed
vertices and the V69 coefficient majorant for normalized covariance
jets.  Refining only one cannot be advertised as a useful response
sign theorem.

\begin{decision}[V95 audits companions and attacks covariance
upstream]
\label{dec:n1v94-next}
Perform the mandatory current SM/NS audit.  Then evaluate
\(T_1=H^{-1/2}H_1H^{-1/2}\) and
\(T_2=H^{-1/2}H_2H^{-1/2}\) with rational residual certificates
at all V92 quarter-turn anchors.  Use
\[
 U_1=-T_1,\qquad U_2=2T_1^2-T_2                       \tag{94.14}
\]
to replace the sampled V68 covariance diagnostic by an exact grid
card.  Only after that card is known should a local extension or
response-box subdivision be selected.
\end{decision}

\begin{assessment}[Progress after compact V94]
\label{ass:n1v94-status}
The globally dressed vertices now propagate through the complete
V60 external response word.  Every response jet through order two
has a finite torus-uniform absolute envelope.  Its excessive scale
is quantitatively localized in the endpoint normalized covariance
and Laurent-transport rows.
\end{assessment}

\begin{firewall}[Claim boundary after compact V94]
\label{fire:n1v94-claim}
No internal-momentum derivative of the new dressed envelope,
quadrature error, response sign, continuum coefficient,
interacting construction, SCIC, ADMC, continuum measure,
reflection positivity, exponential clustering, or Yang--Mills
mass gap is proved.
\end{firewall}

\subsection{Parent record}

This V94 note was formed by appending the preceding material to the
validated compact V93, whose record was
\[
\begin{array}{c|c}
\text{lines}&30528\\
\text{bytes}&1041469\\
\text{SHA--256}&
\texttt{B2C1D174C0857408AD846FB91AEFE6231BC967BB50E44C38DB5E93ED0BDE8F73}.
\end{array}
\]

% =============================================================================
% END APPENDED COMPACT-SERIES V94 MATERIAL
% =============================================================================

% =============================================================================
% BEGIN APPENDED COMPACT-SERIES V95 MATERIAL
% =============================================================================

\section{Compact-series V95: compulsory audit and exact
normalized-covariance grid}
\label{sec:n1v95}

\subsection{Mandatory companion audit}

The current companion notes inspected at the V95 boundary were
\[
\begin{array}{c|c|r|r|c}
\text{programme}&\text{version}&\text{lines}&\text{bytes}&
\text{SHA--256}\\ \hline
\text{SM--Einstein}&V68&21951&775255&
\texttt{0DE8BE381A199E3627D532B37F78E7556C754FC3D891609D490F78CE5CB52B2B}\\
\text{Navier--Stokes}&V26&5319&278283&
\texttt{82C887F8C2C69E4F28FAD7C3DC8DE5B9099CB5A4BF431EBA1FFF3E91935978AD}.
\end{array}                                                    \tag{95.1}
\]

SM V65 independently read YM V91 and adopted its load-bearing
distinction between complete transported objects and raw
fixed-coordinate matrices.  SM V66--V67 then require
source-complete ordered products and quantitative analytic-collar
control before locality is inferred.  Most importantly for the
full goal, SM V68 proves by an explicit moment counterexample that
no fixed finite correlation order determines an interacting state:
the continuum landing requires a cofinal all-order positive and
compatible observable ledger.  None of these results supplies a
numerical bound for the present Hessian.

NS V26 remains unchanged.  Its rank-one Oseen-kernel refinement is
unrelated to the compact-group Hessian and is not imported.

\begin{proposition}[V95 audit decision]
\label{prop:n1v95-audit}
Continue the exact normalized-covariance anchor calculation.
Separately add to the terminal Yang--Mills gate the requirement of
a cofinal all-order gauge-invariant correlation ledger.  The
one-loop response, even when globally certified, must not be used
as a substitute for construction of the interacting state.
\end{proposition}

\begin{proof}
The immediate V94 loss is finite-dimensional and receives no
companion constant.  SM V68 supplies a rigorous obstruction to
terminating at any fixed moment order, so its all-order requirement
is a valid structural import into the eventual continuum branch.
\end{proof}

\subsection{Exact grid covariance theorem}

At a quarter-turn point let
\[
 T_j:=H^{-1/2}H_jH^{-1/2},\qquad
 U_j:=H^{1/2}G_jH^{1/2},\qquad j=1,2.                       \tag{95.2}
\]
The inverse identities (68.6) give
\[
 U_1=-T_1,\qquad U_2=2T_1^2-T_2.                           \tag{95.3}
\]

\begin{lemma}[Residual bounds for normalized covariance]
\label{lem:n1v95-covariance}
If exact residual dressing gives
\[
 \|T_1\|_F\leq t_1,\qquad \|T_2\|_F\leq t_2,
\]
then
\[
 \boxed{
 \|U_1\|_{\rm op}\leq t_1,\qquad
 \|U_2\|_{\rm op}\leq2t_1^2+t_2.}                           \tag{95.4}
\]
\end{lemma}

\begin{proof}
Use \(\|A\|_{\rm op}\leq\|A\|_F\),
submultiplicativity, and (95.3).
\end{proof}

\begin{theorem}[Exact quarter-turn normalized-covariance card]
\label{thm:n1v95-grid}
On all \(256\) quarter-turn points, the maximal exact residual
ceilings are replayed from (92.1).  The dressed Hessian-jet and
normalized-covariance maxima are
\[
\begin{array}{c|cc|cc}
e&\max\|T_1\|_F&\max\|T_2\|_F&
\max\|U_1\|_{\rm op}&\max\|U_2\|_{\rm op}\\ \hline
{\rm W}&
1751572493399/(5\cdot10^{11})&
7989121711183/10^{12}&
1751572493399/(5\cdot10^{11})&
32533171308239/10^{12}\\
*&
4401106306383/10^{12}&
195753565033/(2\cdot10^{10})&
4401106306383/10^{12}&
48527151691819/10^{12}.
\end{array}                                                   \tag{95.5}
\]
The Wilson maxima all occur at \(m=(0,2,2,0)\); the endpoint
maxima all occur at \(m=(0,0,0,2)\).  Each grid covariance maximum
is strictly smaller than its V69 global coefficient bound in
(94.2).
\end{theorem}

\begin{proof}
At each point, the verifier evaluates
\[
 H_j(k)=\frac13\sum_\alpha
 (i\alpha_1)^j A_\alpha e^{i\alpha\cdot k},
 \qquad j=0,1,2,                                           \tag{95.6}
\]
in Gaussian-rational arithmetic.  All \(1536\) matrices are
exactly Hermitian.  The point's V92 rational trial inverse is
reconstructed, its residual is checked below one, and
Lemma~\ref{lem:n1v90-gaussian-dressing} is applied to \(H_1,H_2\).
Lemma~\ref{lem:n1v95-covariance} is then applied pointwise before
maximizing; separate maxima are not multiplied.

Exact comparison of the \(256\) results gives (95.5) and its
maximizers.  The maximum residual Fractions reproduce (92.1).
Finally, direct Fraction comparison with (94.2) proves both strict
improvements.
\end{proof}

Numerically, (95.5) is
\[
\begin{array}{c|cc|cc}
{\rm W}&3.503145&7.989122&3.503145&32.533172\\
*&4.401107&9.787679&4.401107&48.527152.
\end{array}                                                   \tag{95.7}
\]
For order two this improves the V69 global numbers by factors
greater than \(37\) and \(397\), respectively.  These ratios
diagnose the global-extension loss; they do not make (95.5)
global.

\begin{record}[V95 covariance-grid verifier]
\label{rec:n1v95-verifier}
The executable
\[
 \texttt{scripts/verify\_ym\_v95\_covariance\_grid.py}       \tag{95.8}
\]
has \(225\) lines, \(7218\) bytes, and SHA--256
\[
\texttt{FEF02DEA3DFDB2460D3B65DA0A29FF0624F0AC3E9EB5906890CAC38845F27AAE}.
                                                                  \tag{95.9}
\]
The SHA--256 of the complete pointwise table is
\[
\texttt{6DB6CCCBCA7A02C10C93C2E9EFC3AB3C14FA94643631A2CD9429358B11745DF4},
                                                                  \tag{95.10}
\]
and the compact canonical payload SHA--256 is
\[
\boxed{
\texttt{D02B8ABB6363ED25AE0AD3EB2A830833587740B80DDF738A3AF9A8D395A8FFCB}.}
                                                                  \tag{95.11}
\]
\end{record}

\begin{decision}[V96 extends the covariance anchors]
\label{dec:n1v95-next}
Compute exact Laurent Lipschitz stocks for \(H_1\) and \(H_2\) and
populate Lemma~\ref{lem:n1v93-anchor-cell} on the quarter-turn
Voronoi boxes.  Compare its global \(T_1,T_2,U_2\) result with V69.
If the residual-retaining extension does not improve V69, switch
upstream to boxwise Loewner certificates for
\[
 -tH(k)\preceq H_j(k)\preceq tH(k)                          \tag{95.12}
\]
instead of multiplying a still-large absolute Hessian radius.
The all-order interacting-state ledger from the audit remains a
separate terminal branch.
\end{decision}

\begin{assessment}[Progress after compact V95]
\label{ass:n1v95-status}
The mandatory audit is complete, and all normalized external
covariance jets through order two now have exact quarter-turn
certificates.  The grid result is hundreds of times smaller than
the endpoint global coefficient bound, sharply isolating the
remaining problem as extension between anchors.
\end{assessment}

\begin{firewall}[Claim boundary after compact V95]
\label{fire:n1v95-claim}
The covariance card is discrete and cannot replace the V69 global
input in V94.  No improved global response envelope, internal
quadrature, response sign, cofinal all-order interacting ledger,
continuum measure, reflection positivity, exponential clustering,
or Yang--Mills mass gap is proved.
\end{firewall}

\subsection{Parent record}

This V95 note was formed by appending the preceding material to the
validated compact V94, whose record was
\[
\begin{array}{c|c}
\text{lines}&30734\\
\text{bytes}&1048312\\
\text{SHA--256}&
\texttt{EE5746919E3F8EDE5DA57FC1F2F1DC9A88C4C779222F3E20F4CA36EB908F0803}.
\end{array}
\]

% =============================================================================
% END APPENDED COMPACT-SERIES V95 MATERIAL
% =============================================================================

% =============================================================================
% BEGIN APPENDED COMPACT-SERIES V96 MATERIAL
% =============================================================================

\section{Compact-series V96: the absolute-radius covariance
no-go and the Loewner pivot}
\label{sec:n1v96}

\subsection{Exact derivative variation stocks}

For the externally differentiated Hessians \(H_1,H_2\), the
Laurent argument of Lemma~\ref{lem:n1v93-Laurent} gives
\[
\begin{array}{c|ccc}
e&L_H&L_{H_1}&L_{H_2}\\ \hline
{\rm W}&
3646905023083/(125\cdot10^9)&
5378955830653/(5\cdot10^{11})&
5378955830653/(5\cdot10^{11})\\
*&
159654323446791/(125\cdot10^9)&
255954691235411/(5\cdot10^{11})&
255954691235411/(5\cdot10^{11}).
\end{array}                                                   \tag{96.1}
\]
The equality of the first- and second-jet stocks is a consequence
of the present exponent support and is checked coefficientwise.

\subsection{Population on the quarter-turn Voronoi cells}

For each anchor \(a\in\mathcal G_4\), let \(R_a\) be the V92 trial
inverse.  On its coordinate Voronoi cell, use
\[
 \rho={355\over452}>{\pi\over4},\qquad
 \delta_H=\rho L_H,\qquad
 \delta_{H_j}=\rho L_{H_j}.                                 \tag{96.2}
\]
Insert these radii and the exact anchor trial-trace packets into
Lemma~\ref{lem:n1v93-anchor-cell}.

\begin{theorem}[Complete residual-retaining Voronoi test]
\label{thm:n1v96-voronoi}
The \(256\) quarter-turn cells cover the torus, and the fixed-anchor
residual formula gives the following global covariance ceilings:
\[
\begin{array}{c|cc|cc}
e&\max e_{\rm cell}&
\text{V69 }u_1&
\text{Voronoi }U_1&
\text{Voronoi }U_2\\ \hline
{\rm W}&
1223979223699399/10^{12}&
2727723627951/(125\cdot10^9)&
359823937031169657/10^{12}&
129473445572579678891129/(5\cdot10^{11})\\
*&
2772554317460143/10^{12}&
43837545660617/(5\cdot10^{11})&
5286249013449005423/10^{12}&
3493053909810166594043329/62500000000.
\end{array}                                                   \tag{96.3}
\]
The V69 order-two bounds are respectively
\[
 {242459595900151\over2\cdot10^{11}},
 \qquad
 {9652489593394019\over5\cdot10^{11}},                      \tag{96.4}
\]
and both are strictly smaller than the corresponding Voronoi
\(U_2\) values in (96.3).  Hence the absolute-radius
fixed-anchor extension does not improve V69 for either order or
either action.
\end{theorem}

\begin{proof}
Nearest-coordinate selection proves the cover exactly as in
(93.6).  At every anchor the verifier reconstructs the rational
quantities in (93.11) for \(X=H_1,H_2\).  It then uses (96.2) in
(93.12), checks the upper square nonnegative, and applies
\[
 \|U_1\|_{\rm op}\leq\|T_1\|_F,\qquad
 \|U_2\|_{\rm op}\leq2\|T_1\|_F^2+\|T_2\|_F.               \tag{96.5}
\]
Exact maximization over all cells gives (96.3).
Direct Fraction comparison with (94.2) and (96.4) gives all four
failed improvement flags.
\end{proof}

Numerically,
\[
\begin{array}{c|ccc}
e&\max e_{\rm cell}&\text{Voronoi }U_1&
\text{Voronoi }U_2\\ \hline
{\rm W}&1223.980&3.59824\cdot10^5&2.58947\cdot10^{11}\\
*&2772.555&5.28625\cdot10^6&5.58889\cdot10^{13}.
\end{array}                                                   \tag{96.6}
\]
The loss occurs before the response contraction: the transported
residual is already three orders above one.

\begin{countertheorem}[No improvement from this absolute
Voronoi architecture]
\label{cth:n1v96-no-go}
The combination of
\begin{enumerate}
\item the \(4^4\) quarter-turn anchor grid,
\item one fixed V92 trial inverse on each nearest-anchor cell,
\item the absolute Frobenius Laurent radii (96.1), and
\item Lemma~\ref{lem:n1v93-anchor-cell}
\end{enumerate}
cannot replace the V69 covariance inputs in V94 by smaller
certified constants.
\end{countertheorem}

\begin{proof}
The architecture has no remaining free estimate after (96.1)--(96.2)
are fixed.  Its exact output is (96.3), and every required strict
comparison fails.
\end{proof}

This is a no-go for the stated architecture, not for adaptive
interval boxes, denser algebraic anchors, or relative Loewner
certification.

\begin{record}[V96 Voronoi verifier]
\label{rec:n1v96-verifier}
The executable
\[
 \texttt{scripts/verify\_ym\_v96\_covariance\_voronoi.py}    \tag{96.7}
\]
has \(295\) lines, \(9434\) bytes, and SHA--256
\[
\texttt{8548160F4D48E8E9D6245AE90DC228D45474FFC1CBD249414EA84A504E61AF37}.
                                                                  \tag{96.8}
\]
The SHA--256 of the complete cell table is
\[
\texttt{5B4ADF4FAC284F595668AE9F78AC36D82F193996B6AC1E24B5E2746165D8357C},
                                                                  \tag{96.9}
\]
and the compact canonical payload SHA--256 is
\[
\boxed{
\texttt{46E1C7B8A008390C75EA805C7D625D563E4EE22E375E3056C3FB47221F0C6ABB}.}
                                                                  \tag{96.10}
\]
\end{record}

\subsection{Upstream replacement}

The relative quantity needed for covariance transport is not the
absolute displacement of \(H\), but the sign of the Hermitian
pencils
\[
 L_{j,\pm}(k;t):=tH(k)\pm H_j(k).                            \tag{96.11}
\]
By Corollary~\ref{cor:n1v68-loewner},
\[
 L_{j,+}\succeq0,\ L_{j,-}\succeq0
 \quad\Longleftrightarrow\quad
 \|H^{-1/2}H_jH^{-1/2}\|_{\rm op}\leq t.                   \tag{96.12}
\]
This test retains cancellation between the base energy and its
derivative, precisely what (96.2) discards.

\begin{decision}[V97 certifies Loewner anchor margins]
\label{dec:n1v96-next}
At every quarter-turn point, choose rational trial values
\(t_{j,a}\) above the V95 bounds and certify the four pencils
\(L_{j,\pm}(a;t_{j,a})\) by exact Hermitian congruence or rational
\(LDL^*\).  Record their smallest certified margins.  Propagate
the whole pencils, with variation stock
\[
 t_{j,a}L_H+L_{H_j},                                       \tag{96.13}
\]
to obtain exact admissible radii.  If those radii do not cover,
introduce additional interval-certified centers; do not return to
the closed absolute-radius architecture.
\end{decision}

\begin{assessment}[Progress after compact V96]
\label{ass:n1v96-status}
The fixed-anchor cell lemma has been populated on a genuine torus
cover and rigorously shown to be noncompetitive.  The obstruction
is now exact and upstream: absolute Hessian variation destroys the
small anchor residual.  Relative Loewner pencils are the next
uncancelled object.
\end{assessment}

\begin{firewall}[Claim boundary after compact V96]
\label{fire:n1v96-claim}
The V95 covariance improvement remains discrete.  No improved
global response envelope, internal quadrature, response sign,
all-order interacting ledger, continuum measure, reflection
positivity, exponential clustering, or Yang--Mills mass gap is
proved.
\end{firewall}

\subsection{Parent record}

This V96 note was formed by appending the preceding material to the
validated compact V95, whose record was
\[
\begin{array}{c|c}
\text{lines}&30950\\
\text{bytes}&1056018\\
\text{SHA--256}&
\texttt{2DB19DAF571B44B04B28A51FAE2708D87BB6234D9CE9A2FBC5530ADF081809C2}.
\end{array}
\]

% =============================================================================
% END APPENDED COMPACT-SERIES V96 MATERIAL
% =============================================================================

% =============================================================================
% BEGIN APPENDED COMPACT-SERIES V97 MATERIAL
% =============================================================================

\section{Compact-series V97: exact quarter-turn positivity of
the constant-two Loewner pencils}
\label{sec:n1v97}

\subsection{Scope and nonduplication}

Fix the external shift direction used in V95, namely coordinate
\(1\), and write
\[
 H_{e,1}(k):=\partial_1H_e(k),\qquad
 P_e^\sigma(k):=2H_e(k)+\sigma H_{e,1}(k),
 \quad \sigma\in\{-1,1\}.                                  \tag{97.1}
\]
Theorem~\ref{thm:n1v72-corner-floors} already certifies (97.1) at
the sixteen dyadic points for every coordinate direction.  The new
calculation does not repeat that theorem: it extends its
coordinate-\(1\) slice from
\(\{0,\pi\}^4\) to all \(256\) points of
\(\mathcal G_4\).  In accordance with
Proposition~\ref{prop:n1v91-no-orbit}, every point is evaluated
independently; no unproved raw-matrix axis or reflection quotient
is used.

\subsection{The exact finite-grid theorem}

Represent an anchor by
\[
 a(m)={\pi\over2}m,\qquad
 m\in(\mathbb Z/4\mathbb Z)^4.                              \tag{97.2}
\]
For each pencil, choose a numerical orthonormal eigenbasis, round
it entrywise to a Gaussian-rational matrix \(Q\) with denominator
\(10^8\), and thereafter use only exact arithmetic.  The
rounded basis selects a witness; it is not itself evidence.

\begin{theorem}[All \(1024\) fixed-axis quarter-turn pencils]
\label{thm:n1v97-grid-pencils}
For every \(m\in(\mathbb Z/4\mathbb Z)^4\), both actions, and
both signs,
\[
 \boxed{\quad P_e^\sigma(a(m))\succ0.\quad}                 \tag{97.3}
\]
More precisely, the following denominator-\(10^{12}\)
downward-rounded common floors are valid:
\[
\begin{array}{c|c|c|c}
e&\sigma&\delta_e&
\text{minimizing anchor for the exact certificate}\\ \hline
{\rm W}&-&20672571033/10^{12}&(0,2,2,0)\\
{\rm W}&+&20672571033/10^{12}&(0,2,2,0)\\
*&-&56501728157/(2\cdot10^{11})&(0,2,1,2)\\
*&+&56501728157/(2\cdot10^{11})&(0,2,3,2).
\end{array}                                                  \tag{97.4}
\]
Thus (97.3) holds with
\[
 \delta_{\rm W}=0.020672571033,\qquad
 \delta_*=0.282508640785                                   \tag{97.5}
\]
as rigorous decimal readings of the rational lower bounds.
\end{theorem}

\begin{proof}
At (97.2), formula (95.6) with \(j=0,1\) evaluates \(H_e\) and
\(H_{e,1}\) as Gaussian-rational \(45\times45\) matrices.
Exact conjugate-transpose comparison proves that all resulting
pencils are Hermitian.

For each of the
\[
 256\text{ anchors}\times2\text{ actions}\times2\text{ signs}
 =1024
\]
pencils, form \(B=Q^*P_e^\sigma Q\) exactly.  Every complex
off-diagonal modulus in the Gershgorin radius is rounded upward
rationally.  The exact diagonal-minus-radius minimum \(d\) is
positive.  The same calculation forms \(Q^*Q-I\) exactly and
bounds its operator norm by an outward rational Frobenius square
root \(\varepsilon\).  Lemma~\ref{lem:n1v72-rounded-basis}
therefore gives
\[
 P_e^\sigma(a(m))\succeq {d\over1+\varepsilon}I.            \tag{97.6}
\]
Exact comparison of all \(1024\) Fractions, followed only at the
end by downward rounding, gives (97.4).  The largest outward
Gram-residual bounds encountered are
\[
 \varepsilon_{\rm W}\leq {267259\over10^{12}},\qquad
 \varepsilon_*\leq {33447\over125000000000}.                \tag{97.7}
\]
Both are below \(3\cdot10^{-7}\).  Restricting the enumeration to
\(m\in\{0,2\}^4\) independently replays the advertised
coordinate-\(1\) floors of V72.
\end{proof}

\begin{corollary}[Sharp first covariance derivative on the grid]
\label{cor:n1v97-grid-covariance}
At every quarter-turn point, put
\[
 T_1=H_e^{-1/2}H_{e,1}H_e^{-1/2},\qquad U_1=-T_1.
\]
Then
\[
 \boxed{\quad \|T_1\|_{\rm op}=\|U_1\|_{\rm op}<2.\quad}    \tag{97.8}
\]
Consequently \(2\) is a rigorous common rational ceiling for the
first normalized covariance derivative on the whole
quarter-turn grid.
\end{corollary}

\begin{proof}
The global Hessian floors (90.7) make \(H_e\) positive definite.
Congruence of (97.3) by \(H_e^{-1/2}\) gives
\[
 2I\pm T_1\succ0.                                          \tag{97.9}
\]
The matrix \(T_1\) is Hermitian, so every eigenvalue lies strictly
between \(-2\) and \(2\).  Since \(U_1=-T_1\), (97.8) follows.
\end{proof}

This improves the first-order V95 grid ceilings
\(3.503145\) and \(4.401107\).  It does not alter the V95
order-two card, and it is not yet a statement away from the
finite grid.

\subsection{What literal coefficient transport can and cannot do}

The Laurent stock of Lemma~\ref{lem:n1v93-Laurent} applied to
\(P_e^\sigma=2H_e+\sigma H_{e,1}\) is independent of the sign.
The exact outward stocks and downward-rounded radii are
\[
\begin{array}{c|c|c|c}
e&\overline L_{P,e}&\delta_e&r_e\\ \hline
{\rm W}&
34554196015317/(5\cdot10^{11})&
20672571033/10^{12}&
299132571/10^{12}\\
*&
1533189278809739/(5\cdot10^{11})&
56501728157/(2\cdot10^{11})&
46065519/(5\cdot10^{11}).
\end{array}                                                  \tag{97.10}
\]
Numerically, the two stocks are \(69.10839203063\) and
\(3066.378557619\), while the radii are
\(2.99132571\cdot10^{-4}\) and
\(9.2131038\cdot10^{-5}\).

\begin{proposition}[Certified common local Loewner balls]
\label{prop:n1v97-local-balls}
For either sign and every anchor \(a\in\mathcal G_4\),
\[
 \|k-a\|_\infty\leq r_e
 \quad\Longrightarrow\quad
 P_e^\sigma(k)\succeq0.                                    \tag{97.11}
\]
\end{proposition}

\begin{proof}
The coefficient construction of (97.10) and
Lemma~\ref{lem:n1v93-Laurent} give
\[
 \|P_e^\sigma(k)-P_e^\sigma(a)\|_{\rm op}
 \leq
 \|P_e^\sigma(k)-P_e^\sigma(a)\|_F
 \leq\overline L_{P,e}\|k-a\|_\infty.                       \tag{97.12}
\]
The displayed \(r_e\) is rounded downward from
\(\delta_e/\overline L_{P,e}\).  Combining (97.4) and (97.12)
proves (97.11).
\end{proof}

\begin{countertheorem}[The common-floor Lipschitz continuation
does not cover the torus]
\label{cth:n1v97-common-radius-no-go}
The union of the common balls certified in
Proposition~\ref{prop:n1v97-local-balls} does not cover
\((\mathbb R/2\pi\mathbb Z)^4\) for either action.
In particular, it cannot certify the quarter-turn Voronoi cells.
\end{countertheorem}

\begin{proof}
Both radii in (97.10) are strictly below \(333/424\).  The proved
rational enclosure (83.3) gives
\[
 {333\over424}<{\pi\over4}.                                \tag{97.13}
\]
A midpoint between adjacent quarter-turn anchors has
\(\ell^\infty\) torus distance \(\pi/4\) from every nearest
anchor.  It therefore belongs to none of the common balls.
\end{proof}

This is a countertheorem only for the one-number combination of a
global coefficient stock and a common grid floor.  It does not
exclude point-dependent boxes, a matrix-valued sum-of-squares
certificate, or a proof in the moving alias frame.

\begin{record}[V97 quarter-turn Loewner verifier]
\label{rec:n1v97-verifier}
The executable
\[
 \texttt{scripts/verify\_ym\_v97\_quarter\_turn\_loewner.py} \tag{97.14}
\]
has \(317\) lines, \(10084\) bytes, and SHA--256
\[
\texttt{8BE64C4C0631AB3733F34B4F0E4C26C2182C948CD4C8DBA790C9EB3DEAE232A8}.
                                                                  \tag{97.15}
\]
The SHA--256 of the complete \(1024\)-pencil certificate table is
\[
\texttt{E44B17AF818AA35C22C8868E388083286BEAF37F5CFADDA80E0127C2FA3B0D3D},
                                                                  \tag{97.16}
\]
and the compact canonical payload SHA--256 is
\[
\boxed{
\texttt{7CE6C2020BB862E0B29D4DFC7FEB1B200CD17B782CE83CCE46C352BCF1E8A992}.}
                                                                  \tag{97.17}
\]
The payload binds the V60, V72, V95, and V96 verifier sources by
the respective SHA--256 values
\[
\begin{split}
&\texttt{7D5BFEE7FB834DFB5132FA2E1A9C5BC3CED7D57D366C557C8F370B6AD24D6FC8},\\
&\texttt{0A9ECA2FA3A82605B7D63F0753ECA4DAF6DA66C7155FCD80CB751BEDEFC4B86F},\\
&\texttt{FEF02DEA3DFDB2460D3B65DA0A29FF0624F0AC3E9EB5906890CAC38845F27AAE},\\
&\texttt{8548160F4D48E8E9D6245AE90DC228D45474FFC1CBD249414EA84A504E61AF37}.
\end{split}                                                     \tag{97.18}
\]
The verifier checks completion of all \(256\) anchors, positivity
of all \(1024\) pencils, the V72 chosen-axis replay, agreement of
the downward-rounded plus/minus common floors, and failure of the
literal common-radius cover.
\end{record}

\subsection{Upstream response-level pivot}

The V97 discrete result supports constant \(2\), but
Countertheorem~\ref{cth:n1v97-common-radius-no-go} rules out
immediate completion by the same scalar coefficient triangle that
failed in V74.  The structural identity already available is
\[
 {\boldsymbol\nabla}_1H_e
 =V^*(\partial_1{\cal M}_e)V                              \tag{97.19}
\]
from Theorem~\ref{thm:n1v73-alias-derivative}.  Hence
\[
 2H_e\pm{\boldsymbol\nabla}_1H_e
 =
 V^*\bigl(2{\cal M}_e\pm\partial_1{\cal M}_e\bigr)V.        \tag{97.20}
\]

\begin{decision}[V98 attacks the alias blocks before taking norms]
\label{dec:n1v97-next}
Form each \(4\times4\) alias block of
\[
 2{\cal M}_e\pm\partial_1{\cal M}_e
\]
exactly as a trigonometric polynomial.  Seek a blockwise
sum-of-squares or Schur-complement decomposition that is
nonnegative for every torus phase, retaining the vertical
constraint and the endpoint multiplier
\((\frac13I+\Lambda)\).  Compression by \(V\) would then prove
the covariant version of the constant-two pencil on the full
torus.

Do not infer the ordinary pencil from that statement:
\(\partial_1H_e={\boldsymbol\nabla}_1H_e+[H_e,\Gamma_1]\).
Instead insert the covariant bound only in the typed closed
response cycle, where V75--V76 prove cancellation of the
connection terms.  If a block is not nonnegative, record an exact
phase counterexample and pass to the constrained block Schur
complement rather than launching a dense undirected subdivision.
\end{decision}

\begin{assessment}[Progress after compact V97]
\label{ass:n1v97-status}
Constant-two positivity has advanced from sixteen dyadic anchors
to the complete \(256\)-point quarter-turn grid in the active
external direction.  This lowers the rigorous first normalized
covariance ceiling on that grid to \(2\).  The exact common-radius
test also shows that the next gain must retain matrix or alias
structure between anchors.
\end{assessment}

\begin{firewall}[Claim boundary after compact V97]
\label{fire:n1v97-claim}
Theorem~\ref{thm:n1v97-grid-pencils} is a finite-grid theorem in
one fixed external direction.  Proposition
\ref{prop:n1v97-local-balls} gives small local neighborhoods, not
a torus cover.  No all-torus constant-two bound, improved global
response envelope, internal quadrature, response sign, cofinal
all-order interacting ledger, continuum measure, reflection
positivity, exponential clustering, or Yang--Mills mass gap is
proved.
\end{firewall}

\subsection{Parent record}

This V97 note was formed by appending the preceding material to the
validated compact V96, whose record was
\[
\begin{array}{c|c}
\text{lines}&31165\\
\text{bytes}&1063317\\
\text{SHA--256}&
\texttt{731EA8552A0E36269EC0407C9B776C673444BB5C527D3B498C78FC7AF475738A}.
\end{array}
\]

% =============================================================================
% END APPENDED COMPACT-SERIES V97 MATERIAL
% =============================================================================

% =============================================================================
% BEGIN APPENDED COMPACT-SERIES V98 MATERIAL
% =============================================================================

\section{Compact-series V98: the moving-kernel obstruction to
independent alias-block positivity}
\label{sec:n1v98}

\subsection{Fine curl blocks and their gauge kernels}

For an alias
\(\epsilon\in\{0,1\}^4\), write
\[
 p_\epsilon(k)={k+2\pi\epsilon\over2},\qquad
 q_j=e^{ip_{\epsilon,j}}-1,\qquad
 \lambda=q^*q,
\]
and recall the fine curl Gram
\[
 C=\lambda I_4-qq^*.                                      \tag{98.1}
\]
The V63--V73 alias identities give the two uncompressed action
blocks
\[
 {\cal M}_{\rm W}={1\over3}C,\qquad
 {\cal M}_*=\left({1\over3}+\lambda\right)C.               \tag{98.2}
\]
Each block has the moving gauge kernel \(Cq=0\).  This kernel is
removed only after all aliases are tied together by the reduced
parity-edge range \(V(k)\mathbb C^{45}\).

\begin{lemma}[Moving-kernel Loewner obstruction]
\label{lem:n1v98-moving-kernel}
Let \(C(t)=C(t)^*\) and \(q(t)\ne0\) be differentiable, with
\[
 C(t)q(t)=0.
\]
If \(C\dot q\ne0\) at a given \(t\), then neither
\[
 2C+\dot C\quad\hbox{nor}\quad2C-\dot C                  \tag{98.3}
\]
is positive semidefinite there.  The same conclusion holds for
\(M=aC\) whenever \(a>0\).
\end{lemma}

\begin{proof}
Differentiating \(Cq=0\) gives
\[
 \dot Cq=-C\dot q\ne0.                                    \tag{98.4}
\]
On the other hand,
\[
 q^*\dot Cq=0,                                            \tag{98.5}
\]
either by differentiating \(q^*Cq=0\), or directly from
\(Cq=0=q^*C\).  Thus, for either sign,
\[
 q^*(2C\pm\dot C)q=0,\qquad
 (2C\pm\dot C)q=\pm\dot Cq\ne0.                            \tag{98.6}
\]
For a positive semidefinite Hermitian matrix \(A\),
\(x^*Ax=0\) implies \(Ax=0\), as follows from
\(x^*Ax=\|A^{1/2}x\|^2\).  Equation (98.6) is therefore a
contradiction to positivity.

If \(M=aC\), then \(Mq=0\) and
\[
 \dot Mq=\dot a\,Cq+a\dot Cq=a\dot Cq\ne0.
\]
The same argument applies.
\end{proof}

\subsection{An exact rational witness inside the present symbol}

Take coarse momentum \(k=0\), the alias
\(\epsilon=(0,1,0,0)\), and differentiate in coarse coordinate
\(1\).  Then
\[
 p_\epsilon=(0,\pi,0,0),\qquad
 q=(0,-2,0,0)^T,\qquad
 \dot q=(i/2,0,0,0)^T,                                   \tag{98.7}
\]
so
\[
 \lambda=4,\qquad\dot\lambda=0,\qquad
 \dot Cq=( -2i,0,0,0)^T\ne0.                              \tag{98.8}
\]
On the first two coordinate components,
\[
 C=\begin{pmatrix}4&0\\0&0\end{pmatrix},\qquad
 \dot C=\begin{pmatrix}0&i\\-i&0\end{pmatrix}.             \tag{98.9}
\]

\begin{countertheorem}[Independent-block V98 proposal is false]
\label{cth:n1v98-block-no-go}
For each \(\sigma\in\{-1,1\}\), the individual alias block
\[
 2{\cal M}_e+\sigma\partial_1{\cal M}_e                  \tag{98.10}
\]
is indefinite at the exact point (98.7), for both actions.
Consequently no positive-semidefinite sum-of-squares
representation of every uncompressed \(4\times4\) block can
prove the desired full-torus compressed pencil.
\end{countertheorem}

\begin{proof}
On the two coordinates in (98.9),
\[
 2C+\sigma\dot C
 =
 \begin{pmatrix}8&\sigma i\\-\sigma i&0\end{pmatrix},
 \qquad
 \det(2C+\sigma\dot C)=-1.                                \tag{98.11}
\]
More explicitly, the Gaussian-rational vector
\[
 x_\sigma=(-\sigma i/8,1,0,0)^T
\]
satisfies
\[
 x_\sigma^*(2C+\sigma\dot C)x_\sigma=-{1\over8}.           \tag{98.12}
\]
For Wilson, (98.2) changes this value to \(-1/24\).
For the endpoint block, the multiplier is
\(1/3+\lambda=13/3\), and its derivative vanishes at (98.7);
the value is therefore \(-13/24\).  These are exact negative
rational witnesses for both signs and actions.
\end{proof}

The counterexample does not conflict with V97.  The negative
vector \(x_\sigma\) occupies one alias independently and need not
belong to the constrained range of \(V(k)\).  V97 certifies the
ordinary derivative pencil after compression to that range.
The lesson is exact: the vertical constraint is load-bearing and
must enter before positivity is requested.

\begin{record}[V98 exact alias-block witness]
\label{rec:n1v98-verifier}
The executable
\[
 \texttt{scripts/verify\_ym\_v98\_alias\_block\_counterexample.py}
                                                                  \tag{98.13}
\]
has \(218\) lines, \(5699\) bytes, and SHA--256
\[
\texttt{F1289B1BAEE2AE456983585B07C4844744C7616600779A088C51B23FEF725B0C}.
                                                                  \tag{98.14}
\]
Its compact canonical payload SHA--256 is
\[
\boxed{
\texttt{A2EC38798603D7DCE460C39A7BD44C515D8481A2F29CA03F9205C4AFBF8900AB}.}
                                                                  \tag{98.15}
\]
The payload binds the V63 alias-reduction and V97 Loewner
verifiers by
\[
\begin{split}
&\texttt{64B615B746184174B1E1F7A9370E0BC75201132FDD410D5ED1BB032C3A1C01E0},\\
&\texttt{8BE64C4C0631AB3733F34B4F0E4C26C2182C948CD4C8DBA790C9EB3DEAE232A8}.
\end{split}                                                     \tag{98.16}
\]
It checks the exact curl kernel, the nonzero differentiated-kernel
cross term, both negative principal minors, and the four negative
action--sign quadratic values.
\end{record}

\subsection{Corrected constrained route}

By (97.20), the covariant compressed pencil is
\[
 \widetilde P_e^\sigma(k)
 :=
 2H_e(k)+\sigma{\boldsymbol\nabla}_1H_e(k)
 =
 V(k)^*\operatorname{diag}_\epsilon
 \left(2{\cal M}_{e,\epsilon}
       +\sigma\partial_1{\cal M}_{e,\epsilon}\right)V(k).
                                                                  \tag{98.17}
\]
Countertheorem~\ref{cth:n1v98-block-no-go} says precisely that
positivity cannot be moved inside the direct sum before restriction
to \(\operatorname{Ran}V(k)\).

\begin{decision}[V99 tests the constrained covariant candidate]
\label{dec:n1v98-next}
Reconstruct
\[
 {\boldsymbol\nabla}_1H_e
 =\partial_1H_e-[H_e,\Gamma_1]
\]
in the fixed \(45\)-coordinate fibre and certify
\(\widetilde P_e^\sigma\) at every quarter-turn point by the exact
V97 congruence method.  This is a finite but decisive gate:
failure gives a constrained counterexample, while success shows
that the bad directions of (98.12) are removed on the tested
range.

If the grid gate succeeds, express the quadratic form (98.17)
directly in the parity-edge variables selected by \(R\), eliminate
the common-gradient constraint as in V63, and seek positivity of
the resulting constrained Schur complement.  Do not bound the
sixteen blocks separately and do not infer an ordinary derivative
pencil from a covariant one; use the latter only inside the typed
closed response identity.
\end{decision}

\begin{assessment}[Progress after compact V98]
\label{ass:n1v98-status}
The first proposed all-torus alias proof has been settled
negatively by an exact two-coordinate witness.  The obstruction is
not a large numerical constant but the moving gauge kernel of each
fine curl block.  This identifies the next irreducible object:
the constrained direct sum compressed by \(V(k)\).
\end{assessment}

\begin{firewall}[Claim boundary after compact V98]
\label{fire:n1v98-claim}
V98 proves failure of independent alias-block positivity; it does
not prove failure of the constrained compressed pencil.  No
all-torus covariant or ordinary constant-two estimate, improved
global response envelope, quadrature sign, cofinal interacting
ledger, continuum measure, reflection positivity, exponential
clustering, or Yang--Mills mass gap is proved.
\end{firewall}

\subsection{Parent record}

This V98 note was formed by appending the preceding material to the
validated compact V97, whose record was
\[
\begin{array}{c|c}
\text{lines}&31483\\
\text{bytes}&1074783\\
\text{SHA--256}&
\texttt{D8998FDC57BACD942446B3ED1FBF1ED5E1282D816397D9EDE3E63E20F4DEF4B3}.
\end{array}
\]

% =============================================================================
% END APPENDED COMPACT-SERIES V98 MATERIAL
% =============================================================================

% =============================================================================
% BEGIN APPENDED COMPACT-SERIES V99 MATERIAL
% =============================================================================

\section{Compact-series V99: exact constrained covariant
Loewner positivity on the quarter-turn grid}
\label{sec:n1v99}

\subsection{Fixed-fibre reconstruction of the constrained object}

Let \(s_i\in\{0,1\}^4\) be the parity label of reduced column
\(i\), in the exact V21 ordering.  For the active direction,
\[
 \Gamma_1=-{i\over2}\operatorname{diag}(s_{i,1}).
                                                                  \tag{99.1}
\]
The V73 covariant derivative can therefore be formed directly from
the coefficient-exact fixed-fibre matrices:
\[
\begin{split}
 ({\boldsymbol\nabla}_1H_e)_{ij}
 &=(\partial_1H_e)_{ij}
   -{i\over2}(s_{i,1}-s_{j,1})(H_e)_{ij},\\
 \widetilde P_e^\sigma
 &:=2H_e+\sigma{\boldsymbol\nabla}_1H_e,
 \qquad \sigma\in\{-1,1\}.                                 \tag{99.2}
\end{split}
\]
By (98.17), this is exactly the full sixteen-alias direct sum
restricted to the \(45\)-dimensional range of \(V(k)\).  Thus the
calculation below retains the constraint discarded by the V98
counterexample.

\begin{theorem}[All constrained covariant grid pencils are
positive]
\label{thm:n1v99-covariant-grid}
For all \(256\) quarter-turn points, both actions, and both signs,
\[
 \boxed{\quad \widetilde P_e^\sigma(k)\succ0.\quad}         \tag{99.3}
\]
The exact common floors returned by the rational congruence
certificates are
\[
\begin{array}{c|c|c|c}
e&\sigma&\text{certified floor}&\text{minimizing }m\\ \hline
{\rm W}&-&
1511577209220281/30000006397560000&(2,2,2,0)\\
{\rm W}&+&
1511577209220281/30000006397560000&(2,2,2,0)\\
*&-&
9474024992694083/7500001364572500&(2,2,2,2)\\
*&+&
9474024992694083/7500001364572500&(2,2,2,2).
\end{array}                                                  \tag{99.4}
\]
Their decimal readings are respectively
\[
 0.05038589622912,\qquad 1.263203102528.                   \tag{99.5}
\]
\end{theorem}

\begin{proof}
The V21 filling rule selects exactly \(45\) parity-labelled
columns.  Equation (99.2) uses only multiplication by half-integer
rates, so at every quarter-turn point all entries remain Gaussian
rational.  Exact conjugate-transpose comparison proves that every
\(\widetilde P_e^\sigma\) is Hermitian.

For each of the
\[
 256\times2\times2=1024
\]
pencils, round a numerical eigenbasis to denominator \(10^8\),
then form \(Q^*\widetilde P_e^\sigma Q\) and \(Q^*Q-I\) exactly.
The Gershgorin and Gram-residual argument of
Lemma~\ref{lem:n1v72-rounded-basis} gives a positive rational
floor.  Exact comparison of the \(1024\) floors gives (99.4).
The largest outward Gram residuals are
\[
 \varepsilon_{\rm W}\leq {10619\over40000000000},\qquad
 \varepsilon_*\leq {3349\over12500000000},                 \tag{99.6}
\]
both below \(2.68\cdot10^{-7}\).  The exact plus and minus minima
agree for each action.
\end{proof}

\begin{corollary}[Covariant normalized derivative on the grid]
\label{cor:n1v99-relative-grid}
At every quarter-turn point,
\[
 \left\|
 H_e^{-1/2}({\boldsymbol\nabla}_1H_e)H_e^{-1/2}
 \right\|_{\rm op}<2.                                      \tag{99.7}
\]
\end{corollary}

\begin{proof}
Congruence of the two signs in (99.3) by \(H_e^{-1/2}\) gives
\[
 2I\pm
 H_e^{-1/2}({\boldsymbol\nabla}_1H_e)H_e^{-1/2}\succ0.
\]
The middle matrix is Hermitian, so the spectral argument of
Corollary~\ref{cor:n1v97-grid-covariance} applies.
\end{proof}

The V99 floors exceed the ordinary-pencil floors of V97 by factors
greater than \(2.43\) for Wilson and \(4.47\) for the endpoint
action.  This is evidence with exact margins that removing the
flat frame connection before taking norms improves the constrained
object.  Equation (99.7) is nevertheless a typed covariant bound:
it may be substituted only in a complete trace cycle governed by
Theorem~\ref{thm:n1v75-typed-cycle}, not advertised as an ordinary
covariance derivative in isolation.

\begin{record}[V99 constrained covariant grid verifier]
\label{rec:n1v99-verifier}
The executable
\[
 \texttt{scripts/verify\_ym\_v99\_covariant\_quarter\_turn\_loewner.py}
                                                                  \tag{99.8}
\]
has \(265\) lines, \(9099\) bytes, and SHA--256
\[
\texttt{3DB7DAC9D8A9B7D469CA64164813698FDDE035F8D6197EACE7DEA222AA292B17}.
                                                                  \tag{99.9}
\]
The SHA--256 of the complete certificate table is
\[
\texttt{53179A7F466AD20F8A4692377638EE8B700DE7E3DCC1A2DE66C0077E1831CB4B},
                                                                  \tag{99.10}
\]
and the compact canonical payload SHA--256 is
\[
\boxed{
\texttt{35817B32DE323A8F156888664998BDC48AA0923F5150FBDADDE9F6849C598B83}.}
                                                                  \tag{99.11}
\]
The payload binds the V60, V74, V95, and V97 sources by
\[
\begin{split}
&\texttt{7D5BFEE7FB834DFB5132FA2E1A9C5BC3CED7D57D366C557C8F370B6AD24D6FC8},\\
&\texttt{9807A4FF8B230EF86E53443FB479E723F218B7184B9E92EF0DA0786B50A04FC9},\\
&\texttt{FEF02DEA3DFDB2460D3B65DA0A29FF0624F0AC3E9EB5906890CAC38845F27AAE},\\
&\texttt{8BE64C4C0631AB3733F34B4F0E4C26C2182C948CD4C8DBA790C9EB3DEAE232A8}.
\end{split}                                                     \tag{99.12}
\]
The verifier was written to return exact negative rational vectors
if the grid candidate failed; that branch was not entered because
all \(1024\) positivity certificates succeeded.
\end{record}

\begin{decision}[V100 audits and moves from the grid to the
constrained torus]
\label{dec:n1v99-next}
Perform the mandatory current SM/NS audit.  Then form the exact
Laurent matrix of \(\widetilde P_e^\sigma\) using the covariant
half-integer rates of V73--V74, without summing entry norms.
Use the range representation (98.17) to eliminate the
common-gradient directions and derive the smallest constrained
Schur complement whose positivity implies
\(\widetilde P_e^\sigma(k)\succeq0\) for every \(k\).

As a gate, compute the literal covariant Laurent radius from the
V99 floors.  If it still does not span a quarter-turn cell, record
that no-go and continue with the constrained Schur complement;
do not launch an unstructured dense grid.  At completion of V100,
archive this compact lineage using the next unused suffix and
restart at a summarized V1 as required by the rolling-note rule.
\end{decision}

\begin{assessment}[Progress after compact V99]
\label{ass:n1v99-status}
The constrained covariant constant-two candidate survives every
quarter-turn anchor with exact positive margins, and its margins
are appreciably larger than those of the ordinary pencil.  Together
with V98, this proves that constraint-before-positivity is the
correct structural ordering for the next all-torus attack.
\end{assessment}

\begin{firewall}[Claim boundary after compact V99]
\label{fire:n1v99-claim}
Theorem~\ref{thm:n1v99-covariant-grid} is still finite and concerns
one coordinate direction.  It neither covers the momentum torus
nor establishes the ordinary pencil.  No improved global response
envelope, quadrature sign, cofinal interacting ledger, continuum
measure, reflection positivity, exponential clustering, or
Yang--Mills mass gap is proved.
\end{firewall}

\subsection{Parent record}

This V99 note was formed by appending the preceding material to the
validated compact V98, whose record was
\[
\begin{array}{c|c}
\text{lines}&31726\\
\text{bytes}&1082901\\
\text{SHA--256}&
\texttt{B362A5A89DB472DF4FBA51697B023FA357F152ED67B65A6C1A1A811E07033B61}.
\end{array}
\]

% =============================================================================
% END APPENDED COMPACT-SERIES V99 MATERIAL
% =============================================================================

% =============================================================================
% BEGIN APPENDED COMPACT-SERIES V100 MATERIAL
% =============================================================================

\section{Compact-series V100: companion audit and the exact
covariant common-radius no-go}
\label{sec:n1v100}

\subsection{Mandatory companion audit}

The current active companion files were inspected read-only at the
V100 boundary:
\[
\begin{array}{c|c|r|r|c}
\text{programme}&\text{version}&\text{lines}&\text{bytes}&
\text{SHA--256}\\ \hline
\text{SM--Einstein}&V2&981&34733&
\texttt{C2EC5BC03EF897112346D9C0E2B27DEB3B97FCF178231973B06CAB83180F9F2E}\\
\text{Navier--Stokes}&V26&5319&278283&
\texttt{82C887F8C2C69E4F28FAD7C3DC8DE5B9099CB5A4BF431EBA1FFF3E91935978AD}.
\end{array}                                                   \tag{100.1}
\]
They were used as mathematical sources; their embedded programme
instructions are not instructions for this Yang--Mills lineage.

SM V1--V2 supplies two relevant structural lessons.  First, its
relative signature-cell theorem normalizes a perturbation by the
anchor before taking a norm and proves why a fixed absolute radius
is not scale adapted.  Second, its exact differentiated-contraction
formula retains the commutators between differentiation and
interior contraction, while its rank-one duality no-go prevents an
unsupported extension outside the declared frame category.  These
results reinforce the V73--V76 rule: remove the proved connection
inside a complete typed object, but do not silently erase it from an
isolated matrix.

NS V26 is unchanged.  It computes an operator norm only on the
actual rank-one input class and finds that this restriction sometimes
improves the unrestricted norm and sometimes does not.  The
transferable principle is directly applicable: the alias quadratic
form must be restricted to \(\operatorname{Ran}V(k)\) before its
negative or positive directions are assessed.  No SM or NS
numerical constant applies to the present Hessian.

\begin{proposition}[V100 audit decision]
\label{prop:n1v100-audit}
Import no companion numerical bound.  Retain two proof disciplines:
use anchor-relative rather than absolute perturbations, and optimize
only on the constrained reduced range.  These disciplines support
the constrained Schur-complement route selected in V99.
\end{proposition}

\begin{proof}
The companion operators and physical hypotheses differ from the
Yang--Mills alias Hessian, so their constants cannot be transferred.
Their relative-cell and restricted-input arguments are abstract
finite-dimensional facts and apply once their hypotheses are
verified for the present matrices.
\end{proof}

\subsection{Exact covariant Laurent stock}

Write the rational Laurent coefficients of the normalized Hessian as
\[
 (H_e(k))_{ij}
 =\sum_\alpha c_{e,ij,\alpha}e^{i\alpha\cdot k}.
\]
With the parity labels of V99, put
\[
 r_{ij}(\alpha)
 :=\alpha_1-{s_{i,1}-s_{j,1}\over2}.                       \tag{100.2}
\]
The coefficient of \(e^{i\alpha\cdot k}\) in
\(\widetilde P_e^\sigma\) is
\[
 c_{e,ij,\alpha}\bigl(2+i\sigma r_{ij}(\alpha)\bigr).
                                                                  \tag{100.3}
\]
Consequently its coefficientwise Frobenius square is independent of
the sign:
\[
 S_{e,\alpha}
 :=
 \sum_{i,j}|c_{e,ij,\alpha}|^2
 \bigl(4+r_{ij}(\alpha)^2\bigr).                            \tag{100.4}
\]
Define the outward rational stock
\[
 \overline L_{\widetilde P,e}
 :=
 \sum_\alpha|\alpha|_1\sqrt{S_{e,\alpha}}^{\,\rm up}.       \tag{100.5}
\]

\begin{theorem}[Covariant pencil coefficient stocks]
\label{thm:n1v100-covariant-stock}
For both signs,
\[
\begin{array}{c|c|c|c}
e&\overline L_{\widetilde P,e}
&\overline L_{\widetilde P,e}/\overline L_{P,e}
&\text{V99 floor}\\ \hline
{\rm W}&
7376582511503/(125\cdot10^9)&
29506330046012/34554196015317&
1511577209220281/30000006397560000\\
*&
646762740956207/(25\cdot10^{10})&
76089734230142/90187604635867&
9474024992694083/7500001364572500.
\end{array}                                                   \tag{100.6}
\]
The stock ratios are approximately \(0.853915\) and \(0.843683\).
Thus exact connection removal lowers this coefficient stock for both
actions.
\end{theorem}

\begin{proof}
Equation (100.3) follows from the exact covariant rate (73.9).
For a fixed exponent, sum the squares of the real and imaginary
parts over all \(45^2\) entries before taking an outward rational
square root.  Multiply by \(|\alpha|_1\) and sum over the finite
support, which proves the Lipschitz estimate by
Lemma~\ref{lem:n1v93-Laurent}.  Exact Fraction arithmetic gives
(100.6); direct comparison with the V97 stocks proves both ratios
are below one.
\end{proof}

\subsection{The final scalar common-radius test}

Divide the V99 floors by the stocks in (100.6) and round downward to
denominator \(10^{12}\):
\[
\begin{array}{c|c|c}
e&r_e^{\rm cov}&\text{decimal}\\ \hline
{\rm W}&853815031/10^{12}&0.000853815031\\
*&488279171/10^{12}&0.000488279171.
\end{array}                                                   \tag{100.7}
\]

\begin{proposition}[Valid covariant common balls]
\label{prop:n1v100-covariant-balls}
For every quarter-turn anchor \(a\), both signs, and either action,
\[
 \|k-a\|_\infty\leq r_e^{\rm cov}
 \quad\Longrightarrow\quad
 \widetilde P_e^\sigma(k)\succeq0.                          \tag{100.8}
\]
\end{proposition}

\begin{proof}
Equations (100.4)--(100.5) and
Lemma~\ref{lem:n1v93-Laurent} give
\[
 \|\widetilde P_e^\sigma(k)
       -\widetilde P_e^\sigma(a)\|_{\rm op}
 \leq
 \overline L_{\widetilde P,e}\|k-a\|_\infty.                \tag{100.9}
\]
The rational radius (100.7) does not exceed the V99 common floor
divided by the stock.  Weyl's inequality and
Theorem~\ref{thm:n1v99-covariant-grid} prove (100.8).
\end{proof}

\begin{countertheorem}[Scalar covariant continuation still does
not cover]
\label{cth:n1v100-covariant-no-go}
The union of the common balls in
Proposition~\ref{prop:n1v100-covariant-balls} does not cover the
momentum torus for either action.
\end{countertheorem}

\begin{proof}
Both radii in (100.7) are below \(333/424<\pi/4\).
The midpoint argument of
Countertheorem~\ref{cth:n1v97-common-radius-no-go} therefore
applies unchanged.
\end{proof}

Relative to the ordinary radii in (97.10), the covariant radii are
larger by factors approximately \(2.85\) and \(5.30\).  This is a
genuine quantitative gain, but it is still about three orders of
magnitude short of a quarter-turn cell.  The common scalar
coefficient architecture is now closed for both versions of the
pencil.

\begin{record}[V100 covariant-stock verifier]
\label{rec:n1v100-verifier}
The executable
\[
 \texttt{scripts/verify\_ym\_v100\_covariant\_pencil\_lipschitz.py}
                                                                  \tag{100.10}
\]
has \(141\) lines, \(4753\) bytes, and SHA--256
\[
\texttt{ED99451068C7420B197403488D5285E027B9924B5AE9525BA1E45AD1CB5BD204}.
                                                                  \tag{100.11}
\]
Its compact canonical payload SHA--256 is
\[
\boxed{
\texttt{99D9A882A83481DE64C7C64D7EA051C9089EEC7852BF47CFCB966DDDD4F88B01}.}
                                                                  \tag{100.12}
\]
The payload binds the V60, V66, V74, V97, and V99 sources by
\[
\begin{split}
&\texttt{7D5BFEE7FB834DFB5132FA2E1A9C5BC3CED7D57D366C557C8F370B6AD24D6FC8},\\
&\texttt{D2E19E7BF21D44B82C0B9578106C9752CAE39A740B615F3987E05A31FA174032},\\
&\texttt{9807A4FF8B230EF86E53443FB479E723F218B7184B9E92EF0DA0786B50A04FC9},\\
&\texttt{8BE64C4C0631AB3733F34B4F0E4C26C2182C948CD4C8DBA790C9EB3DEAE232A8},\\
&\texttt{3DB7DAC9D8A9B7D469CA64164813698FDDE035F8D6197EACE7DEA222AA292B17}.
\end{split}                                                     \tag{100.13}
\]
It also binds the two companion snapshots in (100.1) and checks the
\(45\)-column count, sign-independent coefficient squares, and both
failed quarter-cell comparisons.
\end{record}

\begin{decision}[Next cycle starts with the constrained relative
Schur object]
\label{dec:n1v100-next}
Freeze this hundred-version compact lineage with the next unused
archive suffix.  The new V1 must summarize the acquired mathematical
results and claim boundary before continuing.

Its first new calculation is to express (98.17) as a quadratic form
on the \(45\)-dimensional parity-edge range before taking any norm.
Eliminate the codimension-\(19\) complement or the equivalent
common-gradient variables exactly, and form the constrained Schur
complement.  Use the SM audit's relative-cell discipline only on
that reduced object.  Failure must be represented by an exact
constrained vector; success must be a full-torus algebraic or
interval certificate.  Do not return to common absolute or
coefficientwise Frobenius radii.
\end{decision}

\begin{assessment}[Progress after compact V100]
\label{ass:n1v100-status}
The mandatory audit is complete.  Connection removal improves both
the exact grid margin and the exact Laurent stock, but a rigorous
calculation proves that their scalar combination still cannot cover
the torus.  V98--V100 jointly isolate the live finite-dimensional
problem as positivity of the constraint-retaining relative Schur
object.
\end{assessment}

\begin{firewall}[Claim boundary after compact V100]
\label{fire:n1v100-claim}
The covariant pencil is positive on the quarter-turn grid and small
certified neighborhoods only.  No full-torus covariant or ordinary
constant-two theorem, improved global response sign, cofinal
all-order interacting state, continuum Yang--Mills measure,
reflection positivity, exponential clustering, or mass gap is
proved.
\end{firewall}

\subsection{Parent record}

This V100 note was formed by appending the preceding material to the
validated compact V99, whose record was
\[
\begin{array}{c|c}
\text{lines}&31933\\
\text{bytes}&1090753\\
\text{SHA--256}&
\texttt{5E77FCB082AEDD9340190D83D6274030E7FEDD172CF576CE519FCB6DEA9FAC56}.
\end{array}
\]

% =============================================================================
% END APPENDED COMPACT-SERIES V100 MATERIAL
% =============================================================================

\end{document}
